\documentclass[10pt,a4paper]{article}
\pdfoutput=1
\usepackage[T1]{fontenc}
\usepackage[a4paper,top=9mm,bottom=8mm,left=9mm,right=9mm,includefoot,footskip=5mm]{geometry}
\usepackage{newpxtext}
\usepackage{microtype}
\usepackage{fancyhdr}
\usepackage{needspace}
\usepackage[compact]{titlesec}
\usepackage{enumitem}
\usepackage{longtable}
\usepackage[latin1]{inputenc}
\usepackage{mystyle}
\let\otmltransp\transp
\let\transp\undefined
\usepackage{newpxmath}
\setcounter{localmathalphabets}{0}
\let\newpxmathtransp\transp
\let\transp\otmltransp
\usepackage{wrapfig}
\usepackage{notations_ot}
\usepackage{tikz}

\definecolor{compactblue}{HTML}{234F70}
\definecolor{compactparagraph}{HTML}{315F83}
\definecolor{compactlink}{HTML}{7A3B2E}
\definecolor{compactgray}{HTML}{606971}
\definecolor{compactwash}{HTML}{F2F5F7}
\usepackage[bookmarks,bookmarksdepth=2,colorlinks=true,
    linkcolor=compactlink,urlcolor=compactlink,
    pdftitle={Optimal Transport for Machine Learners: Compact Teaching Notes},
    pdfauthor={Gabriel Peyr{\'e}},
    pdfsubject={Compact mathematical notes on optimal transport}]{hyperref}
\renewcommand{\a}{\VectMode{a}}
\renewcommand{\b}{\VectMode{b}}
\renewcommand{\c}{c}
\renewcommand{\d}{\ins{d}}
\renewcommand{\H}{H}
\renewcommand{\P}{\VectMode{P}}
\renewcommand{\S}{\VectMode{S}}
\renewcommand{\SS}{\VectMode{S}}
\renewcommand{\th}{\theta}

\setlength{\parindent}{0pt}
\setlength{\parskip}{.6pt plus .2pt}
\linespread{0.92}
\AtBeginDocument{%
    \setlength{\abovedisplayskip}{3pt plus .7pt minus .5pt}%
    \setlength{\belowdisplayskip}{3pt plus .7pt minus .5pt}%
    \setlength{\abovedisplayshortskip}{1pt plus .5pt minus .5pt}%
    \setlength{\belowdisplayshortskip}{1pt plus .5pt minus .5pt}%
}
\setlength{\jot}{1pt}
\setlength{\topsep}{0pt}
\setlength{\partopsep}{0pt}
\setlength{\parsep}{0pt}
\setlength{\itemsep}{0pt}
\setlist{nosep,leftmargin=*}
\titlespacing*{\section}{0pt}{2.2ex plus .5ex}{.65ex}
\titlespacing*{\subsection}{0pt}{1.45ex plus .3ex}{.35ex}
\titlespacing*{\subsubsection}{0pt}{1ex plus .2ex}{.2ex}
\titlespacing*{\paragraph}{0pt}{1.05ex plus .25ex}{.55em}
\titleformat{\section}{\Large\bfseries\color{compactblue}}{\thesection}{.55em}{}[{\vspace{.15ex}\titlerule[.35pt]}]
\titleformat{\subsection}{\large\bfseries\color{compactblue}}{\thesubsection}{.5em}{}
\titleformat{\subsubsection}{\normalsize\bfseries\color{compactgray}}{\thesubsubsection}{.45em}{}
\titleformat{\paragraph}[runin]{\normalfont\bfseries\color{compactparagraph}}{}{0pt}{}
\pretocmd{\paragraph}{\Needspace{5\baselineskip}}{}{}
\mdfdefinestyle{definitionbox}{
    backgroundcolor=compactwash,linecolor=compactparagraph,
    linewidth=.65pt,topline=false,bottomline=false,rightline=false,
    innertopmargin=1.5pt,innerbottommargin=1.5pt,
    innerleftmargin=5pt,innerrightmargin=5pt,
    skipabove=.18\baselineskip,skipbelow=.18\baselineskip,
    nobreak=true,needspace=2.5\baselineskip}
\pagestyle{fancy}
\fancyhf{}
\fancyfoot[L]{\scriptsize\color{compactgray}Optimal Transport\enspace /\enspace Compact Notes}
\fancyfoot[R]{\small\color{compactblue}\thepage}
\renewcommand{\headrulewidth}{0pt}
\renewcommand{\footrulewidth}{0pt}
\allowdisplaybreaks[2]
\emergencystretch=2em
\sloppy

\newcommand{\todo}[1]{}
\newcommand{\blue}[1]{#1}
\newcommand{\removed}[1]{}
\newcommand{\dims}{d}
\providecommand{\mathscr}{\mathcal}

\begin{document}
\begin{center}
{\fontsize{21}{23}\selectfont\bfseries\color{compactblue}Optimal Transport for Machine Learners}\\[.35em]
{\large\itshape Compact teaching notes}\\[.4em]
{\small Gabriel Peyr{\'e}\quad\textperiodcentered\quad\today}
\end{center}
\vspace{-.3em}

\section{Optimal Matching between Point Clouds}
\label{sec-matching}
\subsection{Monge Problem for Discrete Points}
\label{sec-monge-pbm}
\paragraph{Assignment problem.}
\begin{defn}[Optimal assignment problem]\label{def-optimal-assignment}
Let $n\geq1$. For a cost matrix $\C\in\RR^{n\times n}$ and the set $\Perm(n)$ of bijections of $\range{n}=\{1,\ldots,n\}$, the optimal assignment value is
\eql{\label{eq-optimal-assignment}
\umin{\sigma \in \Perm(n)} \frac{1}{n}\sum_{i=1}^n \C_{i,\sigma(i)}.
}
A minimizer is an optimal assignment, or optimal permutation.
\end{defn}
When $\C_{i,j}=c(x_i,y_j)$, every permutation induces the empirical coupling $\pi_\sigma=n^{-1}\sum_{i=1}^n\delta_{(x_i,y_{\sigma(i)})}$ between $\alpha=n^{-1}\sum_{i=1}^n\delta_{x_i}$ and $\beta=n^{-1}\sum_{j=1}^n\delta_{y_j}$. If the source locations are pairwise distinct, this coupling is the graph of the map $T_\sigma(x_i)=y_{\sigma(i)}$, and the assignment problem is exactly the Monge problem between $\alpha$ and $\beta$.

\paragraph{Convex costs on the line.}
\begin{proposition}[Monotone matching on the line]\label{prop-matching-1d-monotone}
Assume that $x_1,\ldots,x_n$ are pairwise distinct and that $y_1,\ldots,y_n$ are pairwise distinct. If $\C_{i,j}=h(x_i-y_j)$ for a strictly convex function $h:\RR\to\RR$, then the unique optimizer is the order-preserving permutation, characterized by
\((x_i-x_{i'})(y_{\sigma(i)}-y_{\sigma(i')})>0 \qquad\text{for every }i\neq i'.\)
In particular, the result applies to $h(s)=|s|^p$ for $p>1$.
\end{proposition}
\begin{proof}
A permutation that does not preserve order contains an inversion: after relabeling, $x<x'$ are matched to $y'>y$. Set $d=y'-y>0$ and
\(D(s)=\frac{h(s)-h(s-d)}{d}\).
Strict convexity of $h$ makes $D$ strictly increasing. Therefore
\(h(x-y')+h(x'-y)-h(x-y)-h(x'-y') =d\bigl(D(x'-y)-D(x-y)\bigr)>0.\)
Swapping the two inverted targets strictly lowers the cost. Repeating this exchange eliminates every inversion; the only order-preserving bijection between the two sorted clouds pairs equal ranks.
\end{proof}
If $h$ is convex but not strictly convex, the same exchange inequality is non-strict. The equal-rank matching remains optimal, but other, non-monotone optimizers can coexist. To construct a monotone optimizer, choose sorting permutations $\sigma_X,\sigma_Y$ such that
\(x_{\sigma_X(1)} \leq x_{\sigma_X(2)} \leq \ldots \qandq y_{\sigma_Y(1)} \leq y_{\sigma_Y(2)} \leq \ldots\)
and match $x_{\sigma_X(k)}$ with $y_{\sigma_Y(k)}$. Thus $\sigma=\sigma_Y\circ\sigma_X^{-1}$ is optimal.

\paragraph{Histogram equalization.}
A typical application is grayscale histogram equalization. Let $I_\ell$, $\ell\in\range{N}$, be the labeled pixel intensities, write $\alpha_N=N^{-1}\sum_{\ell=1}^N\delta_{I_\ell}$, and choose a stable sorting permutation $\iota$ such that $I_{\iota(1)}\leq\cdots\leq I_{\iota(N)}$. For a target law $\beta$, assign the pixel of rank $k$ the intensity $\cumul{\beta}^{-1}((k-1/2)/N)$, using the quantile function of Definition~\ref{def-cdf-quantile}. Stable sorting breaks ties between equal source intensities.

\paragraph{Two-dimensional geometry.}
\begin{proposition}[No proper crossings in Euclidean-distance matchings]\label{prop-planar-no-proper-crossing}
Let $x_1,\ldots,x_n,y_1,\ldots,y_n\in\RR^2$ and let $c(x,y)=\norm{x-y}$. If $\sigma$ is optimal, then two matched segments cannot cross properly: their relative interiors cannot meet at a point where their supporting lines are distinct.
\end{proposition}
\begin{proof}
Suppose that $[x_i,y_{\sigma(i)}]$ and $[x_j,y_{\sigma(j)}]$ cross properly at $z$. The triangle inequality along the two reconnected paths gives
\(\norm{x_i-y_{\sigma(j)}} \leq \norm{x_i-z}+\norm{z-y_{\sigma(j)}} ,\qquad \norm{x_j-y_{\sigma(i)}} \leq \norm{x_j-z}+\norm{z-y_{\sigma(i)}}.\)
The sum of the two right-hand sides is exactly the cost of the two original segments. At least one inequality is strict because a proper crossing is non-collinear. Swapping $\sigma(i)$ and $\sigma(j)$ therefore strictly decreases the assignment cost, contradicting optimality. Collinear overlaps are excluded from the statement precisely because the exchange can then have equal cost.
\end{proof}
For an alternately colored convex $2n$-gon, the number of non-crossing perfect matchings is the Catalan number
\[
C_n=\frac{1}{n+1}\binom{2n}{n}\sim \frac{4^n}{\sqrt{\pi}n^{3/2}}.
\]

The Hungarian algorithm solves the assignment problem exactly in $O(n^3)$ operations by maintaining feasible dual potentials and augmenting a matching along edges where the dual constraint is tight.

\section{Monge Problem between Measures}
\label{sec-monge}
\subsection{Measures}
\label{sec-measures}
\paragraph{Histograms.}
\begin{defn}[Probability simplex]\label{def-probability-simplex}
The probability simplex of length $n$ is \(\simplex_n \eqdef \enscond{\a \in \RR_+^n}{ \sum_{i=1}^n \a_i = 1 }.\)
Its elements are also called probability vectors or histograms.
\end{defn}
\paragraph{Discrete measure, empirical measure.}
\begin{defn}[Discrete measure]\label{def-discrete-measure}
A discrete measure with weights $\a$ and locations $x_1,\dots,x_n\in\X$ is
\eql{\label{eq-discr-meas}
\al = \sum_{i=1}^n \a_i \de_{x_i},
}
where $\de_x$ is the Dirac mass at position $x$. It is a probability measure if $\a\in\simplex_n$, and a positive measure if all weights $\a_i$ are nonnegative.
\end{defn}
\paragraph{General measures.} We write $\Mm(\X)$ for finite signed Borel measures on a metric space $(\X,d)$.
A Dirac measure $\de_x$ is then defined as $\de_x(A)=1$ if $x \in A$ and $0$ otherwise, and this extends by linearity for discrete measures of the form~\eqref{eq-discr-meas} as \(\al(A) = \sum_{x_i \in A} \a_i\)
\paragraph{Polish metric spaces.}
\begin{defn}[Polish metric space]\label{def-polish-metric-space}
A metric space $(\X,d)$ is Polish if it is complete and separable: every Cauchy sequence converges in $\X$, and $\X$ contains a countable dense subset. More generally, a topological space is called Polish if its topology can be induced by some complete separable metric.
\end{defn}
A Polish path space with its uniform metric is \(\Ss=C([0,1];\RR^d),\qquad d_\infty(\gamma,\eta) \eqdef \sup_{t\in[0,1]}\norm{\gamma(t)-\eta(t)}.\)
\begin{defn}[Support of a measure]\label{def:support}
Let $\al$ be a positive Borel measure on a metric space $\Xx$. Its support is \(\supp(\al) \eqdef \enscond{x\in\Xx}{\al(B(x,r))>0\text{ for every }r>0}.\)
If $\Xx$ is separable, this is the smallest closed set of full $\al$-mass.
\end{defn}
\paragraph{Radon measures.}
\begin{defn}[Radon and probability measures]\label{def-radon-probability-measures}
A locally finite positive Borel measure $\al$ on $\Xx$ is Radon if
\(\al(A)=\sup\enscond{\al(K)}{K\subset A\text{ compact}}\)
for every Borel set $A$. We write $\Mm_+(\Xx)$ for finite positive Radon measures and
$\Mm_+^1(\Xx)=\enscond{\al\in\Mm_+(\Xx)}{\al(\Xx)=1}$ for probability measures.
\end{defn}
For an integrable function $f$, we use the pairing notation
\(\dotp{f}{\al}\eqdef\int_\Xx f(x)\d\al(x).\)
Integration of such a measurable $f$ against a discrete measure $\al$ computes a sum
\(\int_\X f(x) \d\al(x) = \sum_{i=1}^n \a_i f(x_i).\)
\paragraph{Relative densities.}
\begin{defn}[Relative density]\label{def-relative-density}
If $\al$ is absolutely continuous with respect to a positive reference measure $\lambda$, its relative density is the Radon--Nikodym derivative
\(\density{\al}\eqdef\frac{\d\al}{\d\lambda}, \qquad \d\al(x)=\density{\al}(x)\d\lambda(x).\)
Equivalently, for every bounded measurable function $h$, and more generally whenever $h$ is integrable with respect to the total variation of $\al$, \(\int_{\Xx} h(x) \d\al(x) = \int_{\Xx} h(x) \density{\al}(x) \d\lambda(x).\)
If $\al$ is positive, the same identity holds in $[0,+\infty]$ for every nonnegative measurable $h$.
\end{defn}
\paragraph{Total variation norm.}
\begin{defn}[Total variation]\label{defn-total-variation}
For a finite signed Radon measure $\al$ on a compact space $\Xx$, \(\norm{\al}_{\TV} \eqdef \usup{f \in \Cc(\Xx)} \enscond{\dotp{f}{\al}}{\norm{f}_\infty \leq 1}.\)
\end{defn}
The absolute value of a signed measure is \(|\al|(A) \eqdef \usup{A=\cup_i B_i} \sum_i |\al(B_i)|,\)
where the supremum is over finite or countable measurable partitions of $A$. Thus, if $\al=\sum_i \a_i \de_{x_i}$ with distinct atoms, $|\al|=\sum_i |\a_i|\de_{x_i}$; if $\d\al(x)=\rho(x)\d\lambda(x)$, then $\d|\al|(x)=|\rho(x)|\d\lambda(x)$.

\begin{prop}[Dual and measure definitions of total variation]\label{prop-tv-dual-measure}
For a finite signed Radon measure $\al$ on a compact space $\Xx$,
one has $\norm{\al}_{\TV}=|\al|(\Xx)$.
\end{prop}
\begin{proof}
The inequality $\norm{\al}_{\TV}\leq |\al|(\Xx)$ follows from
\(\abs{\int f\,\d\al}\leq \int |f|\,\d|\al|\leq \norm{f}_\infty |\al|(\Xx).\)
For the reverse inequality, write the Jordan decomposition $\al=\al^+-\al^-$, so that $|\al|=\al^++\al^-$. The measurable sign $s=\frac{\d\al}{\d|\al|}$ takes values in $\{-1,1\}$ outside a $|\al|$-null set and satisfies $\d\al=s\,\d|\al|$. By regularity of Radon measures on compact spaces, $s$ can be approximated in $L^1(|\al|)$ by continuous functions $f_k$ with $\norm{f_k}_\infty\leq1$. Hence $\int f_k\,\d\al\to\int s\,\d\al=|\al|(\Xx)$, which proves the equality. The final statement is the Riesz--Markov--Kakutani representation theorem with this norm.
\end{proof}
For two absolutely continuous measures $\d\al=\density{\al}\d\lambda$ and $\d\be=\density{\be}\d\lambda$, this gives the concrete formula
\(\norm{\al-\be}_{\TV} = \int_\Xx |\density{\al}(x)-\density{\be}(x)|\,\d\lambda(x).\)
If
\(\al=\sum_{i=1}^n \tilde a_i\delta_{x_i}, \qquad \be=\sum_{j=1}^m \tilde b_j\delta_{y_j},\)
let $z_1,\ldots,z_r$ be the distinct points in the union $\{x_i\}_i\cup\{y_j\}_j$ and set
\(a_k\eqdef\sum_{i:x_i=z_k}\tilde a_i, \qquad b_k\eqdef\sum_{j:y_j=z_k}\tilde b_j.\)
Then $\al=\sum_{k=1}^r a_k\delta_{z_k}$ and $\be=\sum_{k=1}^r b_k\delta_{z_k}$. This convention merges masses located at the same point before comparing the two measures. With this common-support notation,
\(\norm{\al-\be}_{\TV} = \sum_{k=1}^r |a_k-b_k|.\)
\paragraph{Probabilistic interpretation.}
\begin{defn}[Random variable and law]\label{def-random-variable-law}
Let $(\Omega,\Ff,\PP)$ be a probability space and let $\Xx$ be Polish. An $\Xx$-valued random variable is a measurable map $X:(\Omega,\Ff)\to(\Xx,\Bb(\Xx))$. Its law is
\(\operatorname{Law}(X)\eqdef X_\sharp\PP, \qquad \operatorname{Law}(X)(A)=\PP(X\in A)\)
for every Borel set $A\subset\Xx$.
\end{defn}
For every integrable function $f$, integration with respect to this law is expectation: \(\int_\X f(x)\,\d\operatorname{Law}(X)(x)=\EE[f(X)].\)
\subsection{Push-Forward}
\label{sec-push-forward}
For a measurable map $\T: \X \rightarrow \Y$, the push-forward operator $\T_\sharp: \Mm(\X) \rightarrow \Mm(\Y)$ records the distribution of the image points.

For discrete measures~\eqref{eq-discr-meas}, the push-forward operation consists simply in moving the positions of all the points in the support of the measure
\(\T_{\sharp} \al \eqdef \sum_i \a_i \de_{\T(x_i)}.\)
\begin{defn}[Push-forward]\label{defn-pushfwd}
For a measurable map $\T: \X \rightarrow \Y$ and $\al \in \Mm(\X)$, the push-forward measure $\be = \T_\sharp \al \in \Mm(\Y)$ is defined, for every measurable set $B\subset\Y$, by
\eql{\label{eq-equiv-pushfwd}
\be(B) = \al\bigl(\enscond{x \in \X}{\T(x) \in B}\bigr).
}
Equivalently, for every bounded measurable function $h:\Y\to\RR$,
\eql{\label{eq-push-fwd}
\int_\Y h(y) \d \be(y) = \int_\X h(\T(x)) \d\al(x).
}
Note that $\T_\sharp$ preserves positivity and total mass, so that if $\al \in \Mm_+^1(\X)$ then $\T_\sharp \al \in \Mm_+^1(\Y)$.
\end{defn}
\begin{prop}[Push-forward formula for densities]\label{prop-push-forward-densities}
Let $U,V\subset\RR^\dims$ be open, let $\al$ and $\be$ have Lebesgue densities $\density{\al}$ and $\density{\be}$ on $U$ and $V$, and let $\T:U\to V$ be a $C^1$ diffeomorphism such that $\be=\T_\sharp\al$. Then, for Lebesgue-almost every $x\in U$,
\eql{\label{eq-pfwd-density}
\density{\al}(x) = |\det(\T'(x))| \density{\be}(\T(x)).
}
Equivalently, for Lebesgue-almost every $y\in V$, \(\density{\be}(y)=\density{\al}(\T^{-1}y)\, \frac{1}{|\det(\T'(\T^{-1}y))|}.\)
\end{prop}
\begin{proof}
Apply~\eqref{eq-push-fwd} to an arbitrary compactly supported continuous test function $h$ and use the change of variables $y=\T(x)$, for which $\d y = |\det(\T'(x))| \d x$.
\begin{align*}
\int_V h(y)\density{\be}(y) \d y &= \int_V h(y) \d \be(y) = \int_U h(\T(x)) \d \al(x) = \int_U h(\T(x)) \density{\al}(x) \d x \\
&= \int_V h(y) \density{\al}(\T^{-1}y) \frac{\d y}{|\det(\T'(\T^{-1} y))|},
\end{align*}
Since the identity holds for every such $h$, uniqueness of Lebesgue densities gives the target-variable formula in the statement. Substituting $y=\T(x)$ then yields~\eqref{eq-pfwd-density}, where $\T'(x)\in\RR^{\dims\times\dims}$ is the Jacobian matrix of $\T$.
\end{proof}
\subsection{Monge's Formulation}
\label{sec-monge-formulation}
\label{sec-continuous-monge}
\paragraph{Monge problem.}
\begin{defn}[Monge problem and Monge map]\label{def-monge-problem}
For $\al\in\Mm_+^1(\Xx)$, $\be\in\Mm_+^1(\Yy)$ and measurable $c:\Xx\times\Yy\to[0,+\infty]$, define
\eql{\label{eq-monge-continuous}
\Monge_c(\al,\be)
\eqdef
\inf_{\T:\Xx\to\Yy\,\text{ measurable}}
\enscond{\int_\Xx c(x,\T(x))\d\al(x)}{\T_\sharp\al=\be}.
}
A feasible $\T$ is a Monge map; a minimizer is an optimal Monge map. The value is $+\infty$ if no feasible map exists.
\end{defn}
On the real line, for the quadratic cost and finite second moments, the optimal map from an atomless source to an arbitrary target is the monotone rearrangement
\(\T=\cumul{\be}^{-1}\circ\cumul{\al},\)
\begin{prop}[Empirical Monge maps and matchings]\label{prop-empirical-monge-matching}
Assume that the source atoms $x_1,\ldots,x_n$ are distinct and that
\(\al=\frac1n\sum_{i=1}^n\de_{x_i}, \qquad \be=\frac1n\sum_{j=1}^n\de_{y_j}.\)
If $\T_\sharp\al=\be$, then for each distinct target value $z$ in the support of $\be$, exactly $n\be(\{z\})$ source atoms are mapped to $z$. In particular, if the $y_j$ are distinct, then there exists a permutation $\sigma\in\Perm(n)$ such that $\T(x_i)=y_{\sigma(i)}$ for all $i$. Conversely, every such assignment of source atoms to target atoms with the correct masses defines a feasible Monge map on the support of $\al$, and in the distinct-target case
\(\int_\Xx c(x,\T(x))\d\al(x)=\frac1n\sum_{i=1}^n c(x_i,y_{\sigma(i)}).\)
The values of a feasible map away from $\{x_1,\ldots,x_n\}$ may be chosen arbitrarily, provided the resulting extension is measurable. If source locations are repeated, they should first be merged into atoms with larger masses; such atoms cannot be split by a Monge map.
\end{prop}
\begin{proof}
Since $\T_\sharp\al=\be$, all images $\T(x_i)$ must belong to the support of $\be$; otherwise the push-forward would give positive mass to a point outside that support. For any target atom $z$,
\(\be(\{z\})=\al(\T^{-1}(\{z\}))=\frac1n\#\enscond{i}{\T(x_i)=z}.\)
This proves the counting statement. If all target atoms have mass $1/n$, each target receives exactly one source atom, which is a permutation. The converse and the cost identity follow by direct substitution.
\end{proof}
\begin{prop}[Existence of transport maps from atomless sources]\label{prop-existence-transport-map-atomless}
Let $\al$ and $\be$ be Borel probability measures on Polish spaces, and assume that $\al$ is atomless. Then there exists a measurable map $\T$ such that $\T_\sharp\al=\be$.
\end{prop}
\begin{proof}
A standard measure-isomorphism theorem identifies the atomless standard probability space $(\Xx,\al)$ with $([0,1],\mathrm{Leb})$, modulo null sets. It is therefore enough to construct a map from $[0,1]$ to the target law $\be$. Since the target is Polish, it is a standard Borel space; choose a Borel isomorphism $i$ from a full-$\be$ Borel subset of $\Yy$ onto a Borel subset of $[0,1]$. The generalized quantile of $i_\sharp\be$ pushes Lebesgue measure to $i_\sharp\be$. Composing it with $i^{-1}$ and with the source isomorphism gives the desired measurable map, after arbitrary definitions on the discarded null sets.
\end{proof}
\paragraph{Monge distance.}
\begin{defn}[Directed Monge distance]\label{def-directed-monge-distance}
Let $(\Xx,d)$ be a metric space and $1\leq p<+\infty$. For probability measures with finite $p$th moments, define
\eql{\label{eq-monge-distance}
\tilde\Wass_p(\al,\be)^p
\eqdef
\inf_{\T:\Xx\to\Xx\,\text{ measurable}}
\enscond{\int_\Xx d(x,\T(x))^p\d\al(x)}{\T_\sharp\al=\be}.
}
The value is $+\infty$ when the constraint set is empty.
\end{defn}
\begin{prop}[Directed Monge distance]\label{prop-directed-monge-distance}
Let $\Xx=\Yy$ be a Polish metric space and $1\leq p<+\infty$. On probability measures with finite $p$-th moment, the quantity $\tilde\Wass_p$ is nonnegative, vanishes only on the diagonal, and satisfies the triangle inequality. It is therefore a directed extended distance: it need not be symmetric and may take the value $+\infty$.
\end{prop}
\begin{proof}
Nonnegativity is immediate. If $\tilde\Wass_p(\al,\be)=0$, choose feasible maps $\T_k$ with $\int d(x,\T_k(x))^p\d\al(x)\to0$. For every bounded $1$-Lipschitz function $h$,
\[
\left|\int h\d\be-\int h\d\al\right|
=
\left|\int h(\T_k(x))-h(x)\d\al(x)\right|
\leq
\left(\int d(x,\T_k(x))^p\d\al(x)\right)^{1/p}
\longrightarrow0.
\]
Since bounded Lipschitz functions separate probability measures on Polish metric spaces, $\al=\be$. Conversely, the identity map shows that $\tilde\Wass_p(\al,\al)=0$.

If either $\tilde\Wass_p(\al,\ga)$ or $\tilde\Wass_p(\ga,\be)$ is infinite, the claim is automatic. If both are finite, feasible maps $S_\sharp\al=\ga$ and $T_\sharp\ga=\be$ exist, so their composition is feasible from $\al$ to $\be$; in particular, the left-hand side is finite. Fix $\epsilon>0$ and choose $\epsilon$-minimizers satisfying
\(\Ee_\al(S)^{1/p}\leq\tilde\Wass_p(\al,\ga)+\epsilon, \qquad \Ee_\ga(T)^{1/p}\leq\tilde\Wass_p(\ga,\be)+\epsilon.\)
The composed map is feasible from $\al$ to $\be$, and Minkowski's inequality gives
\[
\begin{aligned}
\tilde\Wass_p(\al,\be)
&\leq
\left(\int d(x,T(S(x)))^p\d\al(x)\right)^{1/p} \\
&\leq
\left(\int d(x,S(x))^p\d\al(x)\right)^{1/p}
+
\left(\int d(S(x),T(S(x)))^p\d\al(x)\right)^{1/p} \\
&=
\Ee_\al(S)^{1/p}+\Ee_\ga(T)^{1/p}
\leq
\tilde\Wass_p(\al,\ga)+\tilde\Wass_p(\ga,\be)+2\epsilon.
\end{aligned}
\]
Letting $\epsilon\to0$ gives the result.
\end{proof}
\subsection{Existence and Uniqueness of the Monge Map}
\label{sec-monge-existence-uniqueness}
\paragraph{Brenier's theorem.}
\begin{thm}[Brenier]\label{thm-brenier}
Let $\al,\be\in\Mm_+^1(\RR^d)$ have finite second moments, and assume that $\al$ is absolutely continuous with respect to Lebesgue measure. For the quadratic cost $c(x,y)=\norm{x-y}^2$, there exists a convex function $\phi:\RR^d\to\RR\cup\{+\infty\}$ such that
\(\T=\nabla\phi, \qquad \T_\sharp\al=\be,\)
and $\T$ is the unique optimal Monge map $\al$-almost everywhere. The optimal Kantorovich plan is $(\Id,\T)_\sharp\al$.
\end{thm}
\begin{proof}
Apply Proposition~\ref{prop-quadratic-dual-attainment} to the normalized cost and multiply its dual potentials by two. This supplies optimal extended $c$-concave potentials $(f,g)$ for the cost $\norm{x-y}^2$, finite almost everywhere and integrable against their respective marginals. Define
\(\phi(x)\eqdef\frac12\norm{x}^2-\frac12 f(x), \qquad \psi(y)\eqdef\frac12\norm{y}^2-\frac12 g(y).\)
The $c$-concavity of $f$ makes $\phi$ convex. Dual feasibility is equivalent to $\phi(x)+\psi(y)\geq\dotp{x}{y}$, and equality holds for almost every pair under any optimal plan. Replacing $\psi$ by the convex conjugate $\phi^*$ therefore yields the Fenchel equality $\phi(x)+\phi^*(y)=\dotp{x}{y}$ on the optimal relation, hence $y\in\partial\phi(x)$. Since $\al$ has a density and convex functions are differentiable Lebesgue-almost everywhere, $\partial\phi(x)=\{\nabla\phi(x)\}$ for $\al$-almost every $x$. Every optimal plan is therefore concentrated on the same graph, which proves existence and uniqueness of both the optimal plan and the Monge map.
\end{proof}
Gradients of convex functions are monotone: \(\dotp{\nabla\phi(x)-\nabla\phi(x')}{x-x'}\geq0.\)
\begin{defn}[Measures not charging hypersurfaces]\label{defn-not-charging-hypersurfaces}
A Borel measure $\al$ on $\RR^d$ does not charge hypersurfaces if $\al(S)=0$ for every countably $(d-1)$-rectifiable set $S$, i.e.\ every set that can be covered, up to an $\mathcal H^{d-1}$-null set, by countably many $C^1$ hypersurfaces.
\end{defn}
\paragraph{Radial measures.}
Radial symmetry gives a useful higher-dimensional case where the Brenier map reduces to a one-dimensional monotone rearrangement. A measure $\al$ on $\RR^d$ is radial if $Q_\sharp\al=\al$ for every orthogonal map $Q$.

\begin{prop}[Optimal transport between radial measures]\label{prop-radial-transport}
Let $\al,\be\in\Mm_+^1(\RR^d)$ be radial probability measures with finite second moments, and assume that $\al$ is absolutely continuous with respect to Lebesgue measure. Define their radial distribution functions
\(F_\al(r)\eqdef \al(\{x:\norm{x}\leq r\}), \qquad F_\be(r)\eqdef \be(\{y:\norm{y}\leq r\}), \qquad r\geq0,\)
and let $F_\be^{-1}(u)\eqdef\inf\{r\geq0:F_\be(r)\geq u\}$ be the generalized inverse.
Set $\tau(r)\eqdef F_\be^{-1}(F_\al(r))$.
Then the quadratic Brenier map from $\al$ to $\be$ is the radial map
\(\T(0)=0, \qquad \T(x)=\frac{\tau(\norm{x})}{\norm{x}}x \quad\text{for }x\neq0.\)
Moreover, if $\al_R=(x\mapsto\norm{x})_\sharp\al$ is the law of the source radius, then
\(\Wass_2^2(\al,\be) = \int_0^\infty |r-\tau(r)|^2\,\d\al_R(r).\)
\end{prop}
\begin{proof}
Let $\al_R=(\norm{\cdot})_\sharp\al$ and $\be_R=(\norm{\cdot})_\sharp\be$ be the radial laws. Since $\al$ is absolutely continuous, $\al_R$ has no atoms. The map $\tau=F_\be^{-1}\circ F_\al$ is therefore the monotone one-dimensional transport from $\al_R$ to $\be_R$. If $X\sim\al$ and $R=\norm{X}$, then $\tau(R)$ has law $\be_R$. Moreover, by radiality, the conditional direction $X/\norm{X}$ given $R=r$ is uniform on the sphere for $\al_R$-a.e.\ $r>0$. Thus $\T$ changes only the radius, from $R$ to $\tau(R)$, and keeps the uniform angular distribution. It follows that $\T_\sharp\al$ is radial with radial law $\be_R$, hence $\T_\sharp\al=\be$.

The function $\tau$ is nondecreasing and nonnegative. Define $\psi(r)\eqdef\int_0^r\tau(s)\d s$ on $\RR_+$. Then $\psi$ is convex and nondecreasing, so $\phi(x)\eqdef\psi(\norm{x})$ is convex on $\RR^d$. Since $\al$ is absolutely continuous, it does not charge the origin nor the radii where the monotone function $\tau$ is discontinuous. Thus
$\nabla\phi(x)=\tau(\norm{x})x/\norm{x}=\T(x)$ for $\al$-a.e.\ $x$.
The map $\T$ is therefore a gradient of a convex potential and pushes $\al$ to $\be$. Brenier's theorem identifies it as the unique quadratic optimal Monge map. The displayed cost formula follows by substituting $\T(x)$ in $\int\norm{x-\T(x)}^2\d\al(x)$.
\end{proof}
If $\al$ is not absolutely continuous, the same radial idea is still useful but has to be interpreted at the level of couplings of the radial variables: atoms on spheres may have to be split, so a Monge map need not exist.

\paragraph{Displacement interpolation.}
\label{sec-monge-interpolation}
\begin{defn}[Monge and McCann displacement interpolation]\label{def-monge-mccann-interpolation}
If $\T_\sharp\al=\be$, the Monge interpolation generated by $\T$ is the curve
\(\T_t(x)\eqdef(1-t)x+t\T(x), \qquad \al_t\eqdef(\T_t)_\sharp\al, \qquad t\in[0,1].\)
For the quadratic cost, when $\T$ is the optimal Brenier map, this curve is called McCann's displacement interpolation.
\end{defn}
\begin{prop}[Directed Monge displacement geodesics]\label{prop-monge-displacement-geodesic}
Let $1\leq p<+\infty$, let $\al,\be\in\Mm_+^1(\RR^d)$ have finite $p$-th moments, and let $\T$ be an optimal map for $\tilde\Wass_p(\al,\be)$ in~\eqref{eq-monge-distance}. Set $\al_t=(\T_t)_\sharp\al$ with $\T_t=(1-t)\Id+t\T$. Assume that, for every $t<1$, $\T_t$ is one-to-one on a full $\al$-measure Borel set. Its inverse on the image of that set is then measurable and defined $\al_t$-almost everywhere. For $0\leq s\leq t\leq1$,
\(\tilde\Wass_p(\al_s,\al_t)=(t-s)\tilde\Wass_p(\al,\be).\)
Thus $t\mapsto\al_t$ is an oriented constant-speed geodesic for the directed Monge distance. In particular, for $p=2$, this applies to the Brenier map under the hypotheses of Theorem~\ref{thm-brenier}.
\end{prop}
\begin{proof}
The case $s=t$ is trivial, so assume $s<t$. Since $s<1$, the inverse $\T_s^{-1}$ is defined $\al_s$-almost everywhere along the transported particles. Define
$S_{s,t}\eqdef\T_t\circ\T_s^{-1}$ $\al_s$-a.e.
Then $(S_{s,t})_\sharp\al_s=\al_t$, and, using the optimality of $\T$, \(\tilde\Wass_p(\al_s,\al_t)^p \leq \int \norm{S_{s,t}(z)-z}^p\d\al_s(z) = \int \norm{\T_t(x)-\T_s(x)}^p\d\al(x) = (t-s)^p\tilde\Wass_p(\al,\be)^p.\)
The same particle construction gives $\tilde\Wass_p(\al,\al_s)\leq s\tilde\Wass_p(\al,\be)$. If $t<1$, the map $\T\circ\T_t^{-1}$ sends $\al_t$ to $\be$ and gives $\tilde\Wass_p(\al_t,\be)\leq(1-t)\tilde\Wass_p(\al,\be)$; if $t=1$, this latter distance is zero. Using the triangle inequality from Proposition~\ref{prop-directed-monge-distance},
\[
\tilde\Wass_p(\al,\be)
\leq
\tilde\Wass_p(\al,\al_s)+\tilde\Wass_p(\al_s,\al_t)+\tilde\Wass_p(\al_t,\be)
\leq
s\tilde\Wass_p(\al,\be)+\tilde\Wass_p(\al_s,\al_t)+(1-t)\tilde\Wass_p(\al,\be).
\]
Hence $\tilde\Wass_p(\al_s,\al_t)\geq(t-s)\tilde\Wass_p(\al,\be)$, which proves equality. For a Brenier map $\T=\nabla\phi$, the map $\T_t$ is the gradient of
\(x\mapsto (1-t)\frac{\norm{x}^2}{2}+t\phi(x)\),
which is $(1-t)$-strongly convex for every $t<1$. On the full-measure set where $\phi$ is differentiable, this gives
$\dotp{\T_t(x)-\T_t(y)}{x-y}\geq(1-t)\norm{x-y}^2$, so $\T_t$ is injective there.
\end{proof}
\paragraph{Regularity and the Monge--Amp\`ere equation.}
\begin{prop}[Caffarelli regularity]\label{prop-caffarelli-regularity}
Let $\Omega,\Lambda\subset\RR^d$ be bounded uniformly convex domains with $C^2$ boundaries. Let $\al=\density{\al}(x)\d x$ and $\be=\density{\be}(y)\d y$ be supported on $\Omega$ and $\Lambda$, respectively, with
\(0<m\leq\density{\al},\density{\be}\leq M<+\infty.\)
If $\density{\al}\in C^\alpha(\overline\Omega)$ and $\density{\be}\in C^\alpha(\overline\Lambda)$ for some $\alpha\in(0,1)$, then the Brenier potential $\phi$ is strictly convex in $\Omega$ and belongs to $C^{2,\alpha}_{\mathrm{loc}}(\Omega)$. Under the standard compatibility and higher boundary-regularity assumptions, the same regularity extends to the boundary.
\end{prop}
\begin{proof}
The potential solves the Monge--Amp\`ere equation
\(\det(\nabla^2\phi(x)) =\frac{\density{\al}(x)}{\density{\be}(\nabla\phi(x))}\)
in the Alexandrov sense, with second boundary condition $\nabla\phi(\Omega)=\Lambda$. The density bounds and convexity of the domains give strict convexity and localization of sections. Caffarelli's interior theory then yields $C^{2,\alpha}_{\mathrm{loc}}$ estimates for $\phi$; the boundary statement follows from the boundary regularity theory under uniform convexity and compatibility assumptions.
\end{proof}
For sufficiently smooth positive densities and a classical Brenier map, the change-of-variables formula~\eqref{eq-pfwd-density} gives the Monge--Amp\`ere equation
\eql{\label{eq-monge-ampere}
\det(\nabla^2\phi(x))\density{\be}(\nabla\phi(x))=\density{\al}(x).
}
\begin{prop}[Linearization of the Monge--Amp\`ere equation]\label{prop-linearized-monge-ampere}
Let $\rho_\epsilon=\rho_0+\epsilon r+o(\epsilon)$ be a smooth perturbation of a positive reference density $\rho_0$ on a smooth bounded domain $\Omega$, with $\int_\Omega r\d x=0$. Assume that the expansions hold in a topology strong enough to differentiate the change-of-variables identity. If $\T_\epsilon(x)=x+\epsilon\nabla u(x)+o(\epsilon)$ transports $\rho_0\d x$ to $\rho_\epsilon\d x$, then, to first order,
\(-\nabla\cdot(\rho_0\nabla u)=r.\)
In particular, when $\rho_0$ is constant, the linearized equation is $-\Delta u=r/\rho_0$. If the source and target occupy the same fixed domain and no mass crosses its boundary, the accompanying first-order boundary condition is $\rho_0\partial_n u=0$ on $\partial\Omega$; together with $\int_\Omega u\d x=0$, it fixes the additive constant.
\end{prop}
\begin{proof}
The change-of-variables equation for $\T_\epsilon$ is \(\rho_0(x)=\rho_\epsilon(\T_\epsilon(x))\det(\nabla \T_\epsilon(x)).\)
Expanding $\rho_\epsilon(x+\epsilon\nabla u)=\rho_0(x)+\epsilon r(x)+\epsilon\dotp{\nabla\rho_0}{\nabla u}+o(\epsilon)$ and
$\det(\Id+\epsilon\nabla^2u)=1+\epsilon\Delta u+o(\epsilon)$ gives
\[
\rho_0
=
\rho_0+
\epsilon\bigl(r+\dotp{\nabla\rho_0}{\nabla u}+\rho_0\Delta u\bigr)+o(\epsilon)
=
\rho_0+
\epsilon\bigl(r+\nabla\cdot(\rho_0\nabla u)\bigr)+o(\epsilon).
\]
The first-order term must vanish.
\end{proof}
\subsection{Beyond the Quadratic Euclidean Cost}
\label{sec-beyond-quadratic-euclidean-cost}
\paragraph{\texorpdfstring{$\Wass_p$ costs.}{Wp costs}}
The quadratic cost is special because it identifies the optimal map with the Euclidean gradient of an ordinary convex potential. For the Wasserstein cost, normalized as \(c(x,y)=\norm{x-y}^p/p\) with \(p>1\), the same Monge-map picture survives, but the potential is adapted to the cost. At differentiability points of \(f\),
\(\nabla f(x)=\nabla_x c(x,\T(x)).\)
For \(c(x,y)=\norm{x-y}^p/p\), this gives the explicit relation
\(\T(x)=x-\norm{\nabla f(x)}^{q-2}\nabla f(x), \qquad \frac1p+\frac1q=1.\)
Here $\norm{z}^{q-2}z$ denotes the continuous duality map, defined to be $0$ at $z=0$. This convention includes the case $q=2$ and removes the apparent singularity when $1<q<2$.

\paragraph{Squared geodesic distance.}
Let $M$ be a complete connected Riemannian manifold, let $\al,\be\in\Mm_+^1(M)$ have finite second moments, and assume that $\al$ is absolutely continuous with respect to Riemannian volume. For $c(x,y)=d_M(x,y)^2/2$, the optimal map is expressed using the exponential map:
\(\T(x)=\exp_x(-\nabla\phi(x)),\)
where \(\phi\) is \(c\)-convex and the formula is understood $\al$-almost everywhere, with $\T(x)$ outside the cut locus of $x$. This is the intrinsic version of the formula \(T=x-\nabla\phi\) in normal coordinates. The associated McCann interpolation is
\[
\T_t(x)
=
\exp_x\!\left(t\exp_x^{-1}(\T(x))\right)
=
\exp_x(-t\nabla\phi(x)),
\qquad
\al_t=(\T_t)_\sharp\al,
\qquad t\in[0,1].
\]
More generally, for an arbitrary optimal coupling $\pi$, choose a minimizing geodesic from $x$ to $y$ measurably for $\pi$-almost every $(x,y)$ and let $G_t(x,y)$ be its point at time $t$. Then $\al_t=(G_t)_\sharp\pi$.

\subsection{One-Dimensional Transport and Quantiles}
\label{sec-1d-transport-quantiles}
\paragraph{Cumulative and quantile functions.}
\begin{defn}[Cumulative and quantile functions]\label{def-cdf-quantile}
For $\al\in\Mm_+^1(\RR)$, its cumulative distribution function is
\eql{\label{eq-cumul-defn}
\cumul{\al}(x)\eqdef\al((-\infty,x]).
}
Its generalized inverse, or quantile function, is
\eql{\label{eq-OT-map-1d}
\cumul{\al}^{-1}(r)
\eqdef
\inf\enscond{x\in\RR}{\cumul{\al}(x)\geq r},
\qquad r\in(0,1).
}
\end{defn}
\begin{prop}[Quantile push-forward]\label{prop-quantile-pushforward}
One has $(\cumul{\al}^{-1})_\sharp\mathrm{Leb}_{(0,1)}=\al$. If $\al$ has no atoms, then $(\cumul{\al})_\sharp\al=\mathrm{Leb}_{(0,1)}$.
\end{prop}
\begin{proof}
For simplicity, assume first that $\al$ has a strictly positive density, so that $\cumul{\al}$ is strictly increasing and continuous. Denote $\ga\eqdef(\cumul{\al}^{-1})_\sharp\mathrm{Leb}_{(0,1)}$. It is enough to prove that $\cumul{\ga}=\cumul{\al}$. For every $x$,
\[
\begin{aligned}
\cumul{\ga}(x)
&=
\int_0^1 \mathbf 1_{(-\infty,x]}(\cumul{\al}^{-1}(z))\d z
=
\int_0^1 \mathbf 1_{[0,\cumul{\al}(x)]}(z)\d z
=
\cumul{\al}(x),
\end{aligned}
\]
where we used $\cumul{\al}^{-1}(z)\leq x$ if and only if $z\leq\cumul{\al}(x)$. General measures follow from the same argument with generalized inverses and right-continuity of the cumulative distribution function. If $\al$ has no atoms, the probability integral transform gives $(\cumul{\al})_\sharp\al=\mathrm{Leb}_{(0,1)}$.
\end{proof}
\paragraph{1-D Monge solutions.}\label{par-1d-monge-solution}
\begin{thm}[One-dimensional Monge solution]\label{prop-1d-quantile-map}
Let $\al,\be\in\Mm_+^1(\RR)$, assume that $\al$ has no atoms, and let $h:\RR\to[0,+\infty)$ be convex. Consider the cost $c(x,y)=h(x-y)$. Write $q_\al\eqdef\cumul{\al}^{-1}$ and $q_\be\eqdef\cumul{\be}^{-1}$, and assume that $h(q_\al-q_\be)\in L^1(0,1)$. Then
\eql{\label{eq-1d-monge-map}
\T=q_\be\circ\cumul{\al}
=
\cumul{\be}^{-1}\circ\cumul{\al}
}
satisfies $\T_\sharp\al=\be$ and minimizes
\(\umin{S_\sharp\al=\be}\int h(x-S(x))\d\al(x).\)
In particular, for any $1\leq p<+\infty$, if $\al$ and $\be$ have finite $p$th moments and $\al$ has no atoms, then $h(t)=|t|^p$ gives the usual one-dimensional Monge optimizer.
\end{thm}
\begin{proof}
By Proposition~\ref{prop-quantile-pushforward}, atomlessness of $\al$ gives $(\cumul{\al})_\sharp\al=\mathrm{Leb}_{(0,1)}$, while $(q_\be)_\sharp\mathrm{Leb}_{(0,1)}=\be$. Hence $\T_\sharp\al=\be$.

Let $S$ be any admissible map and set $\pi_S=(\Id,S)_\sharp\al$. The one-dimensional uncrossing theorem for relaxed couplings, proved later as Theorem~\ref{prop-1d-kantorovich-quantile-coupling}, shows that the quantile coupling $\pi^\star=(q_\al,q_\be)_\sharp\mathrm{Leb}_{(0,1)}$ is optimal among all couplings. Since $\al$ has no atoms, $(\Id,\T)_\sharp\al=\pi^\star$, and therefore
\(\int h(x-\T(x))\d\al(x) = \int h(x-y)\d\pi^\star(x,y) \leq \int h(x-y)\d\pi_S(x,y) = \int h(x-S(x))\d\al(x),\)
where every integral is understood in $[0,+\infty]$. This proves optimality of $\T$ among Monge maps.
\end{proof}
For discrete measures, one cannot directly apply the map formula when the source has atoms, but if the measures are uniform on the same number of Dirac masses, then it is exactly the sorting formula of Proposition~\ref{prop-matching-1d-monotone}.

\begin{prop}[One-dimensional Wasserstein formulas]\label{prop-wass-quantile-1d}
Let $1\leq p<+\infty$ and let $\al,\be\in\Mm_+^1(\RR)$ have finite $p$-th moments. Then
\eql{\label{eq-wass-cumul}
\Wass_p(\al,\be)^p
=
\int_0^1\abs{\cumul{\al}^{-1}(r)-\cumul{\be}^{-1}(r)}^p\d r
=
\norm{\cumul{\al}^{-1}-\cumul{\be}^{-1}}_{L^p([0,1])}^p.
}
For $p=1$, this is equivalently
\begin{equation}\label{eq-w1-1d}
\Wass_1(\al,\be)
=
\norm{\cumul{\al}-\cumul{\be}}_{L^1(\RR)}
=
\int_\RR \abs{\cumul{\al}(x)-\cumul{\be}(x)}\d x
=
\int_\RR \abs{\int_{-\infty}^x\d(\al-\be)}\d x.
\end{equation}
\end{prop}
\begin{proof}
The first formula follows from Theorem~\ref{prop-1d-kantorovich-quantile-coupling}: the optimal coupling is obtained by taking the same quantile level $r$ for both measures. For $p=1$, use the layer-cake identity. If $q_\al$ and $q_\be$ are the two quantile functions, then
\(\int_0^1 |q_\al(r)-q_\be(r)|\d r = \int_\RR \lambda\bigl(\{r:q_\al(r)\leq x<q_\be(r)\}\cup\{r:q_\be(r)\leq x<q_\al(r)\}\bigr)\d x.\)
The measure of the displayed set is exactly $|\cumul{\al}(x)-\cumul{\be}(x)|$ for almost every $x$.
\end{proof}
At time $t\in[0,1]$, interpolate the quantiles: \(\cumul{\al_t}^{-1}(r) = (1-t)\cumul{\al}^{-1}(r)+t\cumul{\be}^{-1}(r), \qquad r\in(0,1).\)
\subsection{Gaussian Measures and the Bures Metric}
\label{sec-gaussian-bures}
\paragraph{One-dimensional Gaussians.}
Let $\al=\Gaussian(m_\al,\sigma_\al^2)$ and $\be=\Gaussian(m_\be,\sigma_\be^2)$ be two nondegenerate Gaussians on $\RR$. Then
\(\T(x)=\frac{\sigma_\be}{\sigma_\al}(x-m_\al)+m_\be\)
satisfies $\T_\sharp\al=\be$. It is the derivative of the convex function
\(\phi(x)=\frac{\sigma_\be}{2\sigma_\al}(x-m_\al)^2+m_\be x,\)
The associated distance is
\[
\Wass_2(\al,\be)^2
=
\int_\RR\left(\frac{\sigma_\be}{\sigma_\al}(x-m_\al)+m_\be-x\right)^2\d\al(x)
=
(m_\al-m_\be)^2+(\sigma_\al-\sigma_\be)^2.
\]
The formula extends by continuity to Dirac masses, although the affine Monge map itself only pushes a Dirac source to another Dirac. Thus the OT geometry of one-dimensional Gaussians is the Euclidean geometry of the closed half-plane $(m,\sigma)\in\RR\times\RR_+$.

\paragraph{Multivariate Gaussians.} Consider two Gaussian laws and an affine transport candidate:
\eql{\label{eq-gauss-pf}
\al=\Gaussian(\mean_\al,\cov_\al),
\qquad
\be=\Gaussian(\mean_\be,\cov_\be),
\qquad
\T(x)=\mean_\be+A(x-\mean_\al),
}
\begin{prop}[Affine push-forward of Gaussians]\label{prop-gaussian-affine-push-forward}
One has $\T_\sharp\al=\be$ if and only if
\eql{\label{eq-gauss-covariance-pushforward}
A\cov_\al A^\top=\cov_\be.
}
\end{prop}
\begin{proof}
An affine function maps a Gaussian to a Gaussian, so the law of $\T(X)$ is determined by its mean and covariance. If $X\sim\al$ and $Y=\T(X)$, then
\(\EE(Y)=\mean_\be+A\EE(X-\mean_\al)=\mean_\be\),
and
\(\EE((Y-\mean_\be)(Y-\mean_\be)^\top) = A\EE((X-\mean_\al)(X-\mean_\al)^\top)A^\top = A\cov_\al A^\top.\)
Thus $A\cov_\al A^\top=\cov_\be$ is necessary and sufficient for $Y$ to have the same mean and covariance as $\be$; since both laws are Gaussian, this is equivalent to equality in distribution.
\end{proof}
\begin{defn}[Bures metric]\label{def-bures-metric}
For positive semidefinite covariance matrices $\Sigma$ and $\Lambda$, the Bures metric is
\eql{\label{eq-bures-defn}
\Bb(\Sigma,\Lambda)^2
\eqdef
\tr\pa{\Sigma+\Lambda-2(\Sigma^{1/2}\Lambda\Sigma^{1/2})^{1/2}}.
}
\end{defn}
\begin{prop}[Gaussian $\Wass_2$ formula and Bures covariance term]\label{prop-gaussian-w2-bures}
Assume that $\cov_\al$ and $\cov_\be$ are positive definite. The covariance equation and its unique symmetric positive-definite solution are
\eql{\label{eq-bures-map}
A\cov_\al A=\cov_\be,\quad A\succ0,
\qquad\Longleftrightarrow\qquad
A=
\cov_\al^{-1/2}
\pa{\cov_\al^{1/2}\cov_\be\cov_\al^{1/2}}^{1/2}
\cov_\al^{-1/2}.
}
The affine map $\T(x)=\mean_\be+A(x-\mean_\al)$ is the optimal quadratic-cost transport from $\Gaussian(\mean_\al,\cov_\al)$ to $\Gaussian(\mean_\be,\cov_\be)$, and
\eql{\label{eq-dist-gauss}
\Wass_2(\al,\be)^2
=
\norm{\mean_\al-\mean_\be}^2+\Bb(\cov_\al,\cov_\be)^2,
}
where $\Bb$ is the Bures metric of Definition~\ref{def-bures-metric}.
\end{prop}
\begin{proof}
Multiplying $A\cov_\al A=\cov_\be$ on the left and right by $\cov_\al^{1/2}$ gives
$(\cov_\al^{1/2}A\cov_\al^{1/2})^2
=\cov_\al^{1/2}\cov_\be\cov_\al^{1/2}$.
Since $A$ is symmetric positive, $\cov_\al^{1/2}A\cov_\al^{1/2}$ is symmetric positive and is therefore the unique positive square root of the right-hand side. Conversely, the matrix in~\eqref{eq-bures-map} is symmetric positive and satisfies the covariance equation.

By Proposition~\ref{prop-gaussian-affine-push-forward}, this affine map pushes $\al$ to $\be$. It is the gradient of a convex quadratic potential, so Brenier's theorem implies optimality. If $X\sim\al$, then
\begin{align*}
\EE\norm{X-\T(X)}^2
&=
\norm{\mean_\al-\mean_\be}^2
+
\EE\norm{(\Id-A)(X-\mean_\al)}^2 \\
&=
\norm{\mean_\al-\mean_\be}^2
+
\tr((\Id-A)\cov_\al(\Id-A)^\top) \\
&=
\norm{\mean_\al-\mean_\be}^2
+
\tr(\cov_\al)+\tr(A\cov_\al A)-2\tr(A\cov_\al) \\
&=
\norm{\mean_\al-\mean_\be}^2
+
\tr(\cov_\al)+\tr(\cov_\be)
-2\tr((\cov_\al^{1/2}\cov_\be\cov_\al^{1/2})^{1/2}),
\end{align*}
which is the desired expression. The formula for singular covariance matrices follows by adding $\eta\Id$ and letting $\eta\downarrow0$.
\end{proof}
\paragraph{Comparison with the Fisher--Rao metric for zero-mean Gaussians.}
For $\Sigma,\Sigma'\in\mathbb S_{++}^d$, write
\[
\KL(\Sigma\mid\Sigma')
\eqdef
\KL\big(\Gaussian(0,\Sigma)\mid\Gaussian(0,\Sigma')\big)
=
\frac12\left(
\tr((\Sigma')^{-1}\Sigma)-d-\log\det((\Sigma')^{-1}\Sigma)
\right).
\]
If $D=D^\top$ and $\Sigma+\varepsilon D$ remains positive definite, then
\[
\KL(\Sigma+\varepsilon D\mid\Sigma)
=
\frac{\varepsilon^2}{2}\langle Q_\Sigma(D),D\rangle_F+o(\varepsilon^2),
\qquad
Q_\Sigma(D)=\frac12\Sigma^{-1}D\Sigma^{-1},
\]
where $\langle A,B\rangle_F=\tr(A^\top B)$ is the Frobenius pairing. Thus the local quadratic form is \(g_\Sigma(D,E)=\frac12\tr(\Sigma^{-1}D\Sigma^{-1}E).\)
With this normalization, the Fisher--Rao, or affine-invariant, geodesic distance on $\mathbb S_{++}^d$ is \(d_{\mathrm{FR}}(\Sigma,\Sigma')^2 = \frac12 \norm{\log(\Sigma^{-1/2}\Sigma'\Sigma^{-1/2})}_F^2,\)
where $\log$ denotes the principal matrix logarithm, and the corresponding geodesic is
\[
\Sigma_t^{\mathrm{FR}}
=
\Sigma^{1/2}
\left(\Sigma^{-1/2}\Sigma'\Sigma^{-1/2}\right)^t
\Sigma^{1/2}.
\]
Rank-deficient covariances are at finite Bures distance, but at infinite Fisher--Rao distance.

For $2\times2$ covariances, use $z=(t,u,v)$ with $\Sigma(z)=2^{-1/2}\left(\begin{smallmatrix}t+u&v\\v&t-u\end{smallmatrix}\right)$ and $t\geq\sqrt{u^2+v^2}$, and set
\(\Delta_z=t^2-u^2-v^2,\qquad \langle z,z'\rangle_{\mathrm E}=tt'+uu'+vv', \qquad \langle z,z'\rangle_{\mathrm L}=tt'-uu'-vv'.\)
Then the two distances admit the side-by-side expressions
\begin{align}
\Bb(\Sigma,\Sigma')^2
&=
\sqrt2(t+t')
-2\sqrt{\langle z,z'\rangle_{\mathrm E}
+\sqrt{\Delta_z\Delta_{z'}}},
\label{eq-2x2-bures-cone-distance}\\
d_{\mathrm{FR}}(\Sigma,\Sigma')^2
&=
\frac12\left((\log\lambda_+)^2+(\log\lambda_-)^2\right),
\qquad
\lambda_\pm
=
\frac{\langle z,z'\rangle_{\mathrm L}
\pm\sqrt{\langle z,z'\rangle_{\mathrm L}^2-\Delta_z\Delta_{z'}}}
{\Delta_z},
\label{eq-2x2-fr-cone-distance}
\end{align}
where the Fisher--Rao formula requires $\Delta_z,\Delta_{z'}>0$ and $\lambda_\pm$ are the eigenvalues of $\Sigma^{-1}\Sigma'$.

\section{Kantorovich Relaxation}
\label{sec-kantorovich}
\subsection{Discrete Relaxation}
\label{sec-discrete-relaxation}
\paragraph{Mass splitting and couplings.}
The relaxation is encoded, in place of a permutation $\sigma$ or a map $T$, by a coupling matrix $\P\in\RR_+^{n\times m}$, where $\P_{i,j}$ describes the amount of mass flowing from $x_i$ to $y_j$ in the formalism of discrete measures
\(\al=\sum_i \a_i\de_{x_i}, \qquad \be=\sum_j \b_j\de_{y_j}.\)
\begin{defn}[Discrete couplings and mass conservation]\label{def-discrete-couplings}
Admissible couplings are only constrained to satisfy conservation of mass:
\eql{\label{eq-discr-couplings}
\CouplingsD(\a,\b)
\eqdef
\enscond{\P\in\RR_+^{n\times m}}{\P\ones_m=\a \qandq \transp{\P}\ones_n=\b}.
}
Equivalently, \(\P\ones_m=\bigl(\sum_j \P_{i,j}\bigr)_i\in\RR^n, \qquad \transp{\P}\ones_n=\bigl(\sum_i \P_{i,j}\bigr)_j\in\RR^m.\)
\end{defn}
\begin{defn}[Discrete product coupling]\label{def-discrete-product-coupling}
Given weights $\a\in\simplex_n$ and $\b\in\simplex_m$, the discrete product, or trivial, coupling is \((\a\otimes\b)_{i,j}\eqdef \a_i\b_j.\)
It belongs to $\CouplingsD(\a,\b)$ and corresponds to choosing the source and target labels independently.
\end{defn}
\begin{prop}[Discrete product optimality is degenerate]\label{prop-discrete-product-coupling-degenerate}
Assume that the zero-mass rows and columns have been removed, so that $\a_i>0$ and $\b_j>0$, and let $\C$ be a finite cost matrix. The product plan $\a\otimes\b$ minimizes $\P\mapsto\dotp{\C}{\P}$ over $\CouplingsD(\a,\b)$ if and only if every coupling $\P\in\CouplingsD(\a,\b)$ minimizes it.
\end{prop}
\begin{proof}
The reverse implication is immediate. Assume conversely that $\a\otimes\b$ is optimal and let $\Q\in\CouplingsD(\a,\b)$ be arbitrary. Since all entries of $\a\otimes\b$ are positive, there exists $t>0$ small enough that
\(\R\eqdef (1+t)(\a\otimes\b)-t\Q\)
has nonnegative entries. Its row and column sums are $\R\ones_m=(1+t)\a-t\a=\a$ and $\R^\top\ones_n=(1+t)\b-t\b=\b$, so $\R\in\CouplingsD(\a,\b)$. Moreover,
\(\a\otimes\b=\frac{1}{1+t}\R+\frac{t}{1+t}\Q\).
By optimality of $\a\otimes\b$, both $\dotp{\C}{\R}$ and $\dotp{\C}{\Q}$ are at least $\dotp{\C}{\a\otimes\b}$. Taking the scalar product of the convex-combination identity with $\C$ gives
\(\dotp{\C}{\a\otimes\b} = \frac{1}{1+t}\dotp{\C}{\R} +\frac{t}{1+t}\dotp{\C}{\Q}.\)
A convex average of two numbers not smaller than $\dotp{\C}{\a\otimes\b}$ can equal $\dotp{\C}{\a\otimes\b}$ only if both numbers are equal to it. Hence $\dotp{\C}{\Q}=\dotp{\C}{\a\otimes\b}$, and $\Q$ is optimal. Since $\Q$ was arbitrary, all couplings are optimal.
\end{proof}
Thus the product plan is mainly a feasibility witness. Except in the degenerate situation where the linear cost is constant on the whole transportation polytope, it is not expected to solve optimal transport.

\paragraph{Linear-programming structure.}
\begin{defn}[Discrete Kantorovich problem]\label{def-discrete-kantorovich-problem}
For $\a\in\simplex_n$, $\b\in\simplex_m$, and $\C\in\RR^{n\times m}$, define
\eql{\label{eq-kanto-discr}
\MKD_{\C}(\a,\b)
\eqdef
\umin{\P\in\CouplingsD(\a,\b)}
\dotp{\C}{\P}
=
\umin{\P\in\CouplingsD(\a,\b)}
\sum_{i,j}\C_{i,j}\P_{i,j}.
}
A minimizer $\P$ is an optimal coupling, or optimal transport plan.
\end{defn}
\begin{prop}[Convexity in the marginals and concavity in the cost]\label{prop-discrete-kantorovich-joint-convexity}
For every fixed cost matrix $\C\in\RR^{n\times m}$, the map
\((\a,\b)\longmapsto\MKD_{\C}(\a,\b)\)
is jointly convex on $\simplex_n\times\simplex_m$. For every fixed $(\a,\b)\in\simplex_n\times\simplex_m$, the map
\(\C\longmapsto\MKD_{\C}(\a,\b)\)
is concave on $\RR^{n\times m}$.
\end{prop}
\begin{proof}
For $r\in\{0,1\}$, let $\P_r$ be optimal between $\a_r$ and $\b_r$. For $t\in[0,1]$, the matrix $\P_t=(1-t)\P_0+t\P_1$ belongs to $\CouplingsD((1-t)\a_0+t\a_1,(1-t)\b_0+t\b_1)$. Therefore
\(\MKD_{\C}\big((1-t)\a_0+t\a_1,(1-t)\b_0+t\b_1\big) \leq \dotp{\C}{\P_t} = (1-t)\MKD_{\C}(\a_0,\b_0)+t\MKD_{\C}(\a_1,\b_1).\)
For fixed marginals, every $\P\in\CouplingsD(\a,\b)$ satisfies \(\dotp{(1-t)\C_0+t\C_1}{\P} \geq (1-t)\MKD_{\C_0}(\a,\b) +t\MKD_{\C_1}(\a,\b).\)
Minimizing the left-hand side over $\P$ proves concavity in $\C$.
\end{proof}
\begin{prop}[Rank-controlled sparse minimizers]\label{prop-lp-rank-sparsity}
Let $\mathcal A:\RR^{n\times m}\to\RR^q$ be linear, let $\mathbf d\in\RR^q$ and $\C\in\RR^{n\times m}$, and consider
\begin{equation}\label{eq-lp-rank-sparsity}
\inf_{\P\geq 0,\;\mathcal A(\P)\leq\mathbf d}
\dotp{\C}{\P},
\end{equation}
where the inequality is componentwise. If the feasible set is nonempty and the objective is bounded below on it, then~\eqref{eq-lp-rank-sparsity} admits a minimizer $\P^\star$ satisfying
\begin{equation}\label{eq-lp-rank-support-bound}
\#\enscond{(i,j)}{\P^\star_{i,j}>0}
\leq \rank(\mathcal A).
\end{equation}
Here $\rank(\mathcal A)=\dim(\operatorname{Im}\mathcal A)$. The same conclusion holds when $\mathcal A(\P)\leq\mathbf d$ is replaced by $\mathcal A(\P)=\mathbf d$.
\end{prop}
\begin{proof}
The feasible region is a nonempty polyhedron, and the objective is bounded below on it. The attainment theorem for linear programs therefore provides a minimizer. Among all minimizers, choose $\P^\star$ with the smallest number of positive entries, and denote its support by $S$.

Suppose that $\#S>\rank(\mathcal A)$. The restriction of $\mathcal A$ to the $\#S$-dimensional space of matrices supported on $S$ then has a nontrivial kernel. Hence there is a nonzero matrix $H$, supported on $S$, such that $\mathcal A(H)=0$. Since $\P^\star$ is strictly positive on $S$, both $\P^\star+tH$ and $\P^\star-tH$ remain nonnegative for all sufficiently small $t>0$. Moreover, $\mathcal A(\P^\star\pm tH)=\mathcal A(\P^\star)$, so both perturbations satisfy exactly the same linear inequalities, or equalities, as $\P^\star$. Optimality forces $\dotp{\C}{H}=0$; otherwise one of the two signs would strictly decrease the objective.

Choose $\sigma\in\{-1,1\}$ such that $\sigma H$ has a negative entry, and set $t_\star\eqdef\min_{(i,j):\,\sigma H_{i,j}<0}\P^\star_{i,j}/(-\sigma H_{i,j})>0$.
Then $\P^\star+t_\star\sigma H$ is nonnegative, feasible, and has the same cost as $\P^\star$. At least one of its positive entries has become zero, while no new entry has appeared outside $S$. This contradicts the minimality of $S$ and proves~\eqref{eq-lp-rank-support-bound}.
\end{proof}
\begin{prop}[Sparse optimal plans]\label{prop-sparse-optimal-plans}
Assume $\a_i>0$, $\b_j>0$ and $\sum_i \a_i=\sum_j \b_j=1$. The linear program~\eqref{eq-kanto-discr} admits an optimal coupling with at most $n+m-1$ nonzero entries.
\end{prop}
\begin{proof}
The transportation polytope is nonempty and compact, so the minimum is finite. Apply the equality case of Proposition~\ref{prop-lp-rank-sparsity} to the marginal operator $\mathcal A_{\mathrm{marg}}(\P)\eqdef(\P\ones_m,\transp{\P}\ones_n)\in\RR^{n+m}$.
Its rank is $n+m-1$. Indeed, every matrix satisfies
$\sum_i(\P\ones_m)_i=\sum_j(\transp{\P}\ones_n)_j$, so there is one linear dependence among the $n+m$ marginal coordinates. It is the only one: if $u\in\RR^n$ and $v\in\RR^m$ satisfy
\(\dotp{u}{\P\ones_m} + \dotp{v}{\transp{\P}\ones_n}=0 \qquad\text{for every }\P,\)
then testing this identity on each elementary matrix gives $u_i+v_j=0$ for every $(i,j)$. Thus all $u_i$ equal one scalar and all $v_j$ equal its opposite, so the annihilator of $\operatorname{Im}(\mathcal A_{\mathrm{marg}})$ is one-dimensional. It follows that $\rank(\mathcal A_{\mathrm{marg}})=n+m-1$, and Proposition~\ref{prop-lp-rank-sparsity} supplies the claimed optimal coupling.
\end{proof}
\begin{prop}[North-west corner feasible plan]\label{prop-northwest-corner}
Let $\a\in\RR_+^n$ and $\b\in\RR_+^m$ have the same positive total mass. Starting from $(i,j)=(1,1)$ with residual masses $r_i=\a_i$ and $s_j=\b_j$, skip zero residuals, set $\P_{i,j}=\min(r_i,s_j)$,
subtract this value from both residuals, and advance every index whose residual has become zero. Repeat until all residual masses are exhausted.
\end{prop}
\begin{proof}
All assignments are nonnegative. At each step, the mass placed in entry $(i,j)$ is subtracted from exactly one current row residual and one current column residual, so no row or column can receive more mass than prescribed. Conversely, an index is advanced only when its residual has been fully filled. When the algorithm stops, the total assigned mass is $\sum_i \a_i=\sum_j \b_j$, hence all row and column sums are exactly $\a$ and $\b$.

Each positive assignment exhausts at least one current row or one current column. Before the final assignment, at most $n-1$ row advances and $m-1$ column advances can occur without terminating the construction. Hence the number of positive entries is at most $(n-1)+(m-1)+1=n+m-1$. For acyclicity, view the positive support as a bipartite graph. Once a row or column index is advanced, it never appears again, so each new positive edge either starts a new component or attaches at least one new vertex to the component currently being swept. No edge is ever added between two old vertices of the same component, so no cycle can be created.
\end{proof}
\paragraph{One-dimensional cases.}
\begin{prop}[One-dimensional weighted sweep]\label{prop-1d-weighted-sweep}
Let $x_1\leq\cdots\leq x_n$ and $y_1\leq\cdots\leq y_m$ be points on the line, and let $c(x,y)=h(x-y)$ with $h$ convex. The north-west corner plan between the sorted weighted atoms is optimal for~\eqref{eq-kanto-discr}.
\end{prop}
\begin{proof}
The north-west plan is monotone: if $i<i'$ and $j>j'$, it cannot put positive mass on both $(i,j)$ and $(i',j')$, because the sweep exhausts rows and columns in increasing order. To prove optimality, start from an arbitrary feasible plan $\P$. If
$\P_{1,1}<\min(\a_1,\b_1)$, then row $1$ sends positive mass to some column $j>1$, while column $1$ receives positive mass from some row $i>1$. Moving
\(\eta=\min(\P_{1,j},\P_{i,1})\)
from $(1,j)$ and $(i,1)$ to $(1,1)$ and $(i,j)$ preserves both marginals. Convexity of $h$ gives the Monge inequality
\(h(x_1-y_j)+h(x_i-y_1) \geq h(x_1-y_1)+h(x_i-y_j),\)
so this exchange does not increase the cost. Each exchange either removes the chosen entry in row $1$ or removes the chosen entry in column $1$. Repeating finitely many times forces
$\P_{1,1}=\min(\a_1,\b_1)$, exactly the first north-west assignment. It exhausts row $1$, column $1$, or both. Removing every exhausted index reduces the problem to a smaller ordered transportation problem. Induction on $n+m$ produces the complete north-west plan without increasing cost. Sorting costs $O(n\log n+m\log m)$ and the sweep uses at most $n+m-1$ positive assignments.
\end{proof}
\paragraph{Permutation matrices as couplings.}
\begin{defn}[Permutation matrices]\label{def-permutation-matrices}
For a permutation $\sigma\in\Perm(n)$, its permutation matrix $\P_\sigma$ is
\[
(\P_\sigma)_{i,j}=\choice{1 \qifq j=\sigma(i)\\ 0 \quad\text{otherwise}.}
\]
The set of all permutation matrices is \(\mathcal P_n^{\mathrm{perm}} \eqdef \enscond{\P_\sigma}{\sigma\in\Perm(n)}.\)
\end{defn}
If the matching cost matrix is $\C$, then
\(\dotp{\C}{\P_\sigma/n}=\frac1n\sum_{i=1}^n \C_{i,\sigma(i)}.\)
\begin{defn}[Birkhoff polytope]\label{def-birkhoff-polytope}
The Birkhoff polytope is the convex set of bistochastic matrices
\(\mathcal B_n \eqdef \enscond{\P\in\RR_+^{n\times n}}{\P\ones_n=\ones_n \qandq \transp{\P}\ones_n=\ones_n}.\)
\end{defn}
Then $\CouplingsD(\ones_n/n,\ones_n/n)=\mathcal B_n/n$, and permutation couplings are included in this convex relaxation. More precisely, \(\mathcal P_n^{\mathrm{perm}} = \mathcal B_n\cap\{0,1\}^{n\times n} \subset \mathcal B_n,\)
so before using the structure of $\mathcal B_n$ one first obtains only the relaxed inequality
\(\min_{\P\in\mathcal B_n}\dotp{\C}{\P} \leq \min_{\P\in\mathcal P_n^{\mathrm{perm}}}\dotp{\C}{\P}.\)
\begin{defn}[Extreme points]\label{def-extreme-points}
For a convex set $\mathcal C$ in a finite-dimensional vector space, \(\operatorname{Extr}(\mathcal C) \eqdef \enscond{x\in\mathcal C}{x=(y+z)/2,\ y,z\in\mathcal C \Rightarrow y=z=x}.\)
\end{defn}
\begin{prop}[Existence of extreme points]\label{prop-extreme-point-existence}
If $\mathcal C$ is a non-empty compact convex subset of a finite-dimensional vector space, then $\operatorname{Extr}(\mathcal C)$ is non-empty.
\end{prop}
\begin{proof}
Among all non-empty faces of $\mathcal C$, choose one of minimal affine dimension. If this face contained two distinct points, maximizing a linear functional that is not constant on the face would produce a non-empty proper exposed subface, contradicting minimality. Hence the minimal face is a singleton, and its point is extreme.
\end{proof}
\begin{prop}[Linear programs have extreme minimizers]\label{prop-linear-program-extreme-minimizer}
Let $\mathcal C$ be a non-empty compact convex set. For every linear form $\ell$, \(\operatorname{Extr}(\mathcal C)\cap\operatorname*{argmin}_{x\in\mathcal C}\ell(x) \neq\emptyset.\)
\end{prop}
\begin{proof}
The set $S=\operatorname*{argmin}_{x\in\mathcal C}\ell(x)$ is non-empty, compact and convex. By Proposition~\ref{prop-extreme-point-existence}, it has an extreme point $x$. If $x=(y+z)/2$ with $y,z\in\mathcal C$, then by linearity and optimality of $x$, both $y$ and $z$ also minimize $\ell$ on $\mathcal C$, hence $y,z\in S$. Since $x$ is extreme in $S$, $y=z=x$. Thus $x$ is extreme in $\mathcal C$.
\end{proof}
\begin{thm}[Birkhoff--von Neumann]\label{thm-birkhoff-von-neumann}
The extreme points of $\mathcal B_n$ are exactly the permutation matrices.
\end{thm}
\begin{proof}
We first prove that permutation matrices are extreme. Let $\P_\sigma\in\mathcal P_n^{\mathrm{perm}}$ and assume that $\P_\sigma=(\Q+\R)/2$ with $\Q,\R\in\mathcal B_n$.
Every bistochastic matrix has entries in $[0,1]$. Since the only extreme points of $[0,1]$ are $0$ and $1$, each entry of $\P_\sigma$ fixes the corresponding entries of $\Q$ and $\R$: if $(\P_\sigma)_{i,j}=0$, then $\Q_{i,j}=\R_{i,j}=0$, while if $(\P_\sigma)_{i,j}=1$, then $\Q_{i,j}=\R_{i,j}=1$. Hence $\Q=\R=\P_\sigma$, so $\P_\sigma$ is extreme.

Pick $\P\in\mathcal B_n\setminus\mathcal P_n^{\mathrm{perm}}$. Since an integral bistochastic matrix is necessarily a permutation matrix, $\P$ has at least one fractional entry. We shall split
\(\P=\frac{\Q+\R}{2}\)
with $\Q,\R\in\mathcal B_n$ and $\Q\neq \R$, proving that $\P$ is not extreme.

Associate with $\P$ the bipartite graph whose left vertices are the rows, whose right vertices are the columns, and whose edges are the fractional entries
\(0<\P_{i,j}<1\).
An entry equal to $1$ uses the whole mass of its row and column, so it is isolated in the positive support and does not appear in this fractional graph. If a left vertex $i$ is incident to a fractional edge $(i,j_1)$, then it must be incident to at least one other fractional edge. Indeed, the row sum is one; after the contribution $\P_{i,j_1}\in(0,1)$, a positive amount $1-\P_{i,j_1}$ remains in the same row, and it cannot be carried by an entry equal to $1$. The same argument applies to right vertices, using the column sums. Thus every non-isolated vertex of the fractional graph has degree at least two.

Starting from any fractional edge, one may therefore walk through adjacent fractional edges without immediately backtracking and without getting stuck. Since the graph is finite, some vertex is eventually visited twice; the portion of the walk between the two visits contains a cycle. Choose a shortest such cycle and write it in alternating form $(i_1,j_1,i_2,j_2,\ldots,i_p,j_p)$, with $i_{p+1}=i_1$, where both $(i_s,j_s)$ and $(i_{s+1},j_s)$ are fractional for every $s$. The minimality of the cycle implies that the vertices $i_s$ are all distinct and that the vertices $j_s$ are all distinct. In particular,
\(0<\P_{i_s,j_s}<1 \qandq 0<\P_{i_{s+1},j_s}<1.\)
Define
\(\epsilon \eqdef \min_{1\leq s\leq p} \bigl\{ \P_{i_s,j_s},\, \P_{i_{s+1},j_s},\, 1-\P_{i_s,j_s},\, 1-\P_{i_{s+1},j_s} \bigr\}.\)
All these numbers are positive, so $\epsilon>0$. Split the cycle edges into the two alternating families
\(A\eqdef \{(i_s,j_s)\}_{s=1}^p, \qquad B\eqdef \{(i_{s+1},j_s)\}_{s=1}^p.\)
\[
\Q_{i,j}\eqdef
\begin{cases}
\P_{i,j}, & (i,j)\notin A\cup B,\\
\P_{i,j}+\epsilon/2, & (i,j)\in A,\\
\P_{i,j}-\epsilon/2, & (i,j)\in B,
\end{cases}
\qquad
\R_{i,j}\eqdef
\begin{cases}
\P_{i,j}, & (i,j)\notin A\cup B,\\
\P_{i,j}-\epsilon/2, & (i,j)\in A,\\
\P_{i,j}+\epsilon/2, & (i,j)\in B.
\end{cases}
\]
By the definition of $\epsilon$, all modified entries stay in $[0,1]$, so $\Q$ and $\R$ are nonnegative. Each row vertex $i_s$ of the cycle is incident to exactly one edge of $A$ and one edge of $B$; the $+\epsilon/2$ and $-\epsilon/2$ perturbations therefore cancel in that row. The same cancellation holds in each column vertex $j_s$, and all other rows and columns are unchanged. Finally, $\Q\neq \R$ because $\epsilon>0$ and the cycle is non-empty, while by construction $\P=(\Q+\R)/2$. Thus $\P$ is not extreme.
\end{proof}
\begin{cor}[Kantorovich for matching]\label{cor-kantorovich-matching}
If $m=n$ and $\a=\b=\ones_n/n$, then the discrete Kantorovich problem~\eqref{eq-kanto-discr} admits an optimal solution of the form $\P_\sigma/n$. The associated permutation $\sigma$ solves the assignment problem of Section~\ref{sec-monge-pbm}.
\end{cor}
\begin{proof}
The feasible set is $\mathcal B_n/n$. By Proposition~\ref{prop-linear-program-extreme-minimizer}, the linear objective has an optimal extreme point. Since scaling preserves extreme points and Theorem~\ref{thm-birkhoff-von-neumann} identifies the extreme points of $\mathcal B_n$, this optimizer is $\P_\sigma/n$ for some permutation $\sigma$. Its cost is exactly $n^{-1}\sum_i \C_{i,\sigma(i)}$, so $\sigma$ is an optimal assignment.
\end{proof}
\paragraph{Constructive Birkhoff--von Neumann decomposition.}
\begin{cor}[Birkhoff--von Neumann decomposition]\label{cor-birkhoff-von-neumann-decomposition}
Every $\P\in\mathcal B_n$ admits a representation
\(\P=\sum_{r=1}^{L}\lambda_r\P_{\sigma_r}, \qquad \lambda_r>0, \qquad \sum_{r=1}^{L}\lambda_r=1, \qquad L\leq (n-1)^2+1,\)
where each $\P_{\sigma_r}$ is a permutation matrix.
\end{cor}
\begin{proof}
Let $\mathcal K$ be the convex hull of the permutation matrices. Clearly $\mathcal K\subset\mathcal B_n$. Conversely, if some $\P\in\mathcal B_n$ did not belong to the compact convex set $\mathcal K$, strict separation would provide a linear form $\ell$ such that $\ell(\P)<\min_{\Q\in\mathcal K}\ell(\Q)$. Proposition~\ref{prop-linear-program-extreme-minimizer} and Theorem~\ref{thm-birkhoff-von-neumann} give
\(\min_{\Q\in\mathcal B_n}\ell(\Q) =\min_{\sigma\in\Perm(n)}\ell(\P_\sigma) =\min_{\Q\in\mathcal K}\ell(\Q),\)
contradicting $\P\in\mathcal B_n$. Hence $\mathcal B_n=\mathcal K$.

The $2n$ row- and column-sum equations have rank $2n-1$, and the strictly positive matrix $\ones_n\transp{\ones_n}/n$ lies in their affine space. Therefore $\mathcal B_n$ has affine dimension $n^2-(2n-1)=(n-1)^2$. Carath\'eodory's theorem in this affine hull now reduces any convex decomposition to at most $(n-1)^2+1$ permutation matrices; zero coefficients can be discarded.
\end{proof}
If $N(I)$ denotes the neighbor set of $I$, then
\(s|I| =\sum_{i\in I}\sum_{j\in N(I)}\R_{i,j} \leq\sum_{j\in N(I)}\sum_i\R_{i,j} =s|N(I)|.\)
Thus $|N(I)|\geq |I|$, so $G_{\R}$ contains a perfect matching.

If $\lambda<s$, then
\(\widehat\P' \eqdef \frac{\R-\lambda\P_\sigma}{s-\lambda}\)
If the current support graph has $E$ edges, the elementary matching subroutine performs at most $n$ breadth-first searches and costs $O(nE)$. Thus a dense implementation costs $O(n^3)$ per extracted permutation and, using the term bound above, at most $O(n^5)$ arithmetic and graph operations overall.

\paragraph{Rational weights.}
\begin{proposition}[Rational weights as duplicated uniform matching]\label{prop-rational-weights-duplicated-matching}
Let
\(\al=\sum_{i=1}^n \frac{k_i}{N}\delta_{x_i}, \qquad \be=\sum_{j=1}^m \frac{\ell_j}{N}\delta_{y_j}, \qquad \sum_i k_i=\sum_j\ell_j=N,\)
with positive integers $k_i,\ell_j$. Replace each $x_i$ by $k_i$ identical copies and each $y_j$ by $\ell_j$ identical copies, producing two uniform $N$-point clouds. Then the duplicated assignment problem and the discrete Kantorovich problem between $\al$ and $\be$ have the same optimal value. Moreover, there exists an optimal coupling of the form
\(\P_{i,j}=\frac{n_{i,j}}{N}, \qquad n_{i,j}\in\NN, \qquad \sum_j n_{i,j}=k_i, \quad \sum_i n_{i,j}=\ell_j,\)
and couplings whose scaled entries $N\P_{i,j}$ are integral are exactly the collapsed assignments between duplicated clouds. Fractional optimal couplings may nevertheless coexist when the optimum is degenerate.
\end{proposition}
\begin{proof}
Any assignment between the duplicated source and target clouds defines integers $n_{i,j}$ counting how many copied particles of type $x_i$ are matched to copied particles of type $y_j$. These counts satisfy $\sum_j n_{i,j}=k_i$ and $\sum_i n_{i,j}=\ell_j$, and the associated coupling $\P_{i,j}=n_{i,j}/N$ has marginals $k_i/N$ and $\ell_j/N$. The assignment cost is $\frac1N\sum_{i,j} n_{i,j}c(x_i,y_j)=\sum_{i,j}\P_{i,j}c(x_i,y_j)$.
Conversely, any nonnegative integer count matrix with those row and column sums can be realized by allocating the $k_i$ copies of each $x_i$ among the target copies according to $(n_{i,j})_j$. Scale a feasible coupling by $N$ and write $\Q=N\P$. The constraints on $\Q$ have integer right-hand sides $(k_1,\ldots,k_n,\ell_1,\ldots,\ell_m)$. After multiplying the target rows by $-1$, their coefficient matrix is the oriented node-edge incidence matrix of a bipartite graph and is therefore totally unimodular. Every vertex of this transportation polytope is integral. Since a linear objective attains its minimum at a vertex, an integral optimal $\Q$ exists, proving equality of the two optimal values. The argument establishes existence, not integrality of every optimizer; convex combinations of distinct integral optima give fractional optima.
\end{proof}
\subsection{Linear-Programming Algorithms}
\label{sec-kantorovich-lp-algorithms}
\paragraph{Interior-point methods.}
The logarithmic-barrier problem on the resulting transport polytope is
\eql{\label{eq-transport-log-barrier}
\P_\epsilon
\eqdef
\uargmin{\substack{\P\ones_m=\a,\ \P^\top\ones_n=\b\\ \P_{i,j}>0}}
\dotp{\C}{\P}
-
\epsilon\sum_{i,j}\log \P_{i,j},
}
where $\epsilon>0$ is decreased along the algorithm. The barrier is singular at the boundary, so each iterate stays strictly inside the transportation polytope; as $\epsilon\downarrow0$, the central path approaches the set of LP minimizers.

\subsection{Relaxation for Arbitrary Measures}
\label{sec-kantorovich-continuous}
\paragraph{Continuous couplings.}
\begin{defn}[Marginals of a joint measure]\label{def-joint-marginals}
Let $\pi\in\Mm_+^1(\Xx\times\Yy)$ and let $P_\Xx(x,y)=x$, $P_\Yy(x,y)=y$ be the coordinate projections. The marginals of $\pi$ are
\(\pi_1\eqdef(P_\Xx)_\sharp\pi\in\Mm_+^1(\Xx), \qquad \pi_2\eqdef(P_\Yy)_\sharp\pi\in\Mm_+^1(\Yy).\)
Equivalently, for all bounded continuous test functions $f$ on $\Xx$ and $g$ on $\Yy$, \(\int_{\Xx\times\Yy} f(x)\d\pi(x,y)=\int_\Xx f\d\pi_1, \qquad \int_{\Xx\times\Yy} g(y)\d\pi(x,y)=\int_\Yy g\d\pi_2.\)
\end{defn}
A useful mnemonic for the marginal constraint $\pi_1=\al$ and $\pi_2=\be$ is the formal row-and-column notation
\(\int_\Yy \d\pi(x,y)=\d\al(x), \qquad \int_\Xx \d\pi(x,y)=\d\be(y),\)
which is made rigorous by Definition~\ref{def-joint-marginals}, or equivalently by the identities $\pi(A\times\Yy)=\al(A)$ and $\pi(\Xx\times B)=\be(B)$ for measurable sets $A\subset\Xx$ and $B\subset\Yy$.

\begin{defn}[Couplings]\label{def-continuous-couplings}
Given $\al\in\Mm_+^1(\Xx)$ and $\be\in\Mm_+^1(\Yy)$, the set of couplings between $\al$ and $\be$ is
\eql{\label{eq-coupling-generic}
\Couplings(\al,\be)
\eqdef
\enscond{\pi\in\Mm_+^1(\Xx\times\Yy)}{\pi_1=\al \qandq \pi_2=\be}.
}
This is the continuous analogue of the transportation polytope~\eqref{eq-discr-couplings}.
\end{defn}
\begin{defn}[Tensor product and trivial coupling]\label{def-tensor-product-coupling}
Given $\al\in\Mm_+^1(\Xx)$ and $\be\in\Mm_+^1(\Yy)$, the tensor product coupling $\al\otimes\be$ is the probability measure on $\Xx\times\Yy$ defined by
\[
\int_{\Xx\times\Yy} h(x,y)\d(\al\otimes\be)(x,y)
=
\int_\Xx\left(\int_\Yy h(x,y)\d\be(y)\right)\d\al(x)
\]
for every bounded measurable $h$. It is also called the trivial coupling because it makes the two coordinates independent.
\end{defn}
Indeed, for every $f\in\Cc_b(\Xx)$,
\[
\int_{\Xx\times\Yy} f(x)\d(\al\otimes\be)(x,y)
=
\left(\int_\Xx f(x)\d\al(x)\right)\left(\int_\Yy\d\be(y)\right)
=
\int f\d\al,
\]
\begin{prop}[Product optimality is degenerate]\label{prop-product-coupling-degenerate}
Assume that $\Xx$ and $\Yy$ are compact metric spaces and that $c\in\Cc(\Xx\times\Yy)$.
\(\al\otimes\be\in\uargmin{\pi\in\Couplings(\al,\be)}\int c\,\d\pi, \qquad \text{every coupling in }\Couplings(\al,\be)\text{ is optimal}.\)
They are also equivalent to the additive decomposition of the cost on the product support, \(c(x,y)=u(x)+v(y), \qquad u\in\Cc(\supp(\al)),\quad v\in\Cc(\supp(\be)).\)
\end{prop}
\begin{proof}
If every coupling is optimal, then $\al\otimes\be$ is optimal. Conversely, assume that $\al\otimes\be$ is optimal. We first show that, for every $x_0,x_1\in\supp(\al)$ and $y_0,y_1\in\supp(\be)$,
\(c(x_0,y_0)+c(x_1,y_1) = c(x_0,y_1)+c(x_1,y_0).\)
Indeed, if this equality failed, after exchanging $y_0$ and $y_1$ if necessary one would have a strict inequality
\(c(x_0,y_0)+c(x_1,y_1) > c(x_0,y_1)+c(x_1,y_0).\)
Failure of the equality forces $x_0\neq x_1$ and $y_0\neq y_1$. By continuity, the strict inequality persists with a uniform margin on small neighborhoods $U_0,U_1$ of $x_0,x_1$ and $V_0,V_1$ of $y_0,y_1$, chosen disjoint within each pair. Since the four points lie in the supports, these neighborhoods have positive marginal mass. Denote by $\al_i$ and $\be_i$ the normalized restrictions of $\al$ to $U_i$ and of $\be$ to $V_i$, and choose
\(0<\lambda\leq \min\{\al(U_0)\be(V_0),\al(U_1)\be(V_1)\}.\)
The exchanged measure
\(\tilde\pi = \al\otimes\be -\lambda\,\al_0\otimes\be_0 -\lambda\,\al_1\otimes\be_1 +\lambda\,\al_0\otimes\be_1 +\lambda\,\al_1\otimes\be_0\)
is nonnegative and has the same two marginals as $\al\otimes\be$. The uniform strict inequality on the neighborhoods implies that $\int c\,\d\tilde\pi<\int c\,\d(\al\otimes\be)$, contradicting optimality.

Fixing any $x_\star\in\supp(\al)$ and $y_\star\in\supp(\be)$, the equality of cross differences gives, for all $(x,y)\in\supp(\al)\times\supp(\be)$, \(c(x,y)=c(x,y_\star)+c(x_\star,y)-c(x_\star,y_\star).\)
Thus $c=u+v$ on the product support, with $u(x)=c(x,y_\star)-c(x_\star,y_\star)$ and $v(y)=c(x_\star,y)$. These functions are continuous on the corresponding supports. Every coupling is concentrated on this product support, so for any $\pi\in\Couplings(\al,\be)$,
\(\int c\,\d\pi = \int u\,\d\al+\int v\,\d\be\),
which depends only on the marginals. Hence all couplings are optimal.
\end{proof}
The tensor product is therefore a trivial feasible coupling, not a typical optimizer. Product optimality means that the cost cannot distinguish between dependences once the marginals are fixed.

If there exists a map $T:\Xx\to\Yy$ such that $T_\sharp\al=\be$, then the Monge map induces the graph coupling $\pi=(\Id,T)_\sharp\al\in\Couplings(\al,\be)$, characterized by
\(\int_{\Xx\times\Yy} h(x,y)\d\pi(x,y) = \int_\Xx h(x,T(x))\d\al(x).\)
Applying this identity to $h(x,y)=f(x)$ or $h(x,y)=g(y)$ gives respectively $\pi_1=\al$ and $\pi_2=\be$. Thus graph couplings are precisely the Kantorovich representation of deterministic Monge maps.

\paragraph{Continuous Kantorovich problem.}
\begin{defn}[Continuous Kantorovich problem]\label{def-continuous-kantorovich-problem}
For $\al\in\Mm_+^1(\Xx)$, $\be\in\Mm_+^1(\Yy)$, and a nonnegative Borel cost $c:\Xx\times\Yy\to[0,+\infty]$, define
\eql{\label{eq-mk-generic}
\MK_c(\al,\be)
\eqdef
\inf_{\pi\in\Couplings(\al,\be)}
\int_{\Xx\times\Yy} c(x,y)\d\pi(x,y).
}
A minimizer is an optimal coupling, or optimal transport plan.
The value is understood in $[0,+\infty]$.
\end{defn}
\begin{prop}[Convexity in the marginals and concavity in the cost]\label{prop-kantorovich-value-curvature}
Fix a nonnegative Borel cost $c$. The extended-valued map
\((\al,\be)\longmapsto\MK_c(\al,\be)\)
is jointly convex on pairs of probability measures. Conversely, for fixed marginals $(\al,\be)$, the map $c\mapsto\MK_c(\al,\be)$ is concave on the cone of nonnegative Borel costs: for any such $c_0,c_1$ and $t\in[0,1]$,
\(\MK_{(1-t)c_0+tc_1}(\al,\be) \geq (1-t)\MK_{c_0}(\al,\be)+t\MK_{c_1}(\al,\be).\)
The inequality is understood in $[0,+\infty]$; in particular, the map is concave on every convex family of costs on which it is finite.
\end{prop}
\begin{proof}
For joint convexity, let $(\al_0,\be_0)$ and $(\al_1,\be_1)$ be two pairs of probability measures. The claim is immediate at $t\in\{0,1\}$ or if one of the two values on the right-hand side is infinite. Otherwise, for $\eta>0$, choose $\pi_i\in\Couplings(\al_i,\be_i)$ such that $\int c\d\pi_i\leq\MK_c(\al_i,\be_i)+\eta$ for $i\in\{0,1\}$.
Then $(1-t)\pi_0+t\pi_1$ couples $(1-t)\al_0+t\al_1$ and $(1-t)\be_0+t\be_1$, and hence
\(\MK_c\big((1-t)\al_0+t\al_1,(1-t)\be_0+t\be_1\big) \leq (1-t)\MK_c(\al_0,\be_0)+t\MK_c(\al_1,\be_1)+\eta.\)
Letting $\eta\to0$ proves joint convexity.

For concavity, the endpoint cases are immediate, so assume $t\in(0,1)$ and fix $\pi\in\Couplings(\al,\be)$. Linearity of integration for nonnegative functions gives
\(\int\big((1-t)c_0+tc_1\big)\d\pi \geq (1-t)\MK_{c_0}(\al,\be)+t\MK_{c_1}(\al,\be).\)
Taking the infimum over $\pi$ proves the stated inequality.
\end{proof}
\begin{prop}[Existence for lower-semicontinuous costs]\label{prop-kantorovich-existence-compact}
Assume that $\Xx$ and $\Yy$ are Polish spaces and that $c:\Xx\times\Yy\to[0,+\infty]$ is lower semicontinuous. Then the Kantorovich problem~\eqref{eq-mk-generic} admits at least one minimizer. In particular, this applies to every continuous nonnegative cost on compact metric spaces.
\end{prop}
\begin{proof}
The constraint set is non-empty because it contains the product coupling $\al\otimes\be$. It is uniformly tight: given $\varepsilon>0$, choose compact sets $K_\Xx\subset\Xx$ and $K_\Yy\subset\Yy$ such that $\al(K_\Xx)\geq1-\varepsilon/2$ and $\be(K_\Yy)\geq1-\varepsilon/2$. Every $\pi\in\Couplings(\al,\be)$ then satisfies
\(\pi(K_\Xx\times K_\Yy)\geq1-\varepsilon\).
By Prokhorov's theorem, the set of couplings is relatively compact for weak convergence; the weak topology is reviewed in Section~\ref{sec-wasserstein-topology-applications}. It is also weakly closed, because the coordinate projections are continuous and therefore preserve the marginal identities in the limit. Hence $\Couplings(\al,\be)$ is weakly compact. Finally, Portmanteau's theorem makes $\pi\mapsto\int c\d\pi$ weakly lower semicontinuous. The direct method gives a minimizer.
\end{proof}
\paragraph{Axiomatic characterization of the relaxation.}
\begin{thm}[Kantorovich cost as the maximal convex relaxation]\label{thm-kantorovich-maximal-convex-relaxation}
Assume that $\Xx$ and $\Yy$ are compact metric spaces and that $c:\Xx\times\Yy\to[0,+\infty)$ is continuous. Then $\MK_c$ is the pointwise largest proper functional
\(\mathsf F:\Mm_+^1(\Xx)\times\Mm_+^1(\Yy)\longrightarrow(-\infty,+\infty]\)
that is jointly convex, weakly lower semicontinuous, and satisfies \(\mathsf F(\de_x,\de_y)\leq c(x,y) \qquad\text{for every }(x,y)\in\Xx\times\Yy.\)
More precisely, every such $\mathsf F$ satisfies $\mathsf F(\al,\be)\leq\MK_c(\al,\be)$ for all $(\al,\be)$,
while $\MK_c$ itself has these properties and obeys $\MK_c(\de_x,\de_y)=c(x,y)$.
\end{thm}
\begin{proof}
Joint convexity follows from Proposition~\ref{prop-kantorovich-value-curvature}. For weak lower semicontinuity, consider converging marginals and optimal plans along a subsequence realizing the lower limit. Compactness gives a further weakly convergent subsequence; its limit has the limiting marginals, and lower semicontinuity of the integral gives the claim, exactly as in Proposition~\ref{prop-kantorovich-existence-compact}. The only coupling of $(\de_x,\de_y)$ is $\de_{(x,y)}$, which proves the identity on Dirac pairs.

Fix $\pi\in\Couplings(\al,\be)$. Its marginal constraints are exactly the barycentric identity $(\al,\be)=\int(\de_x,\de_y)\d\pi$.
Jensen's inequality for lower-semicontinuous convex functionals therefore gives \(\mathsf F(\al,\be) \leq\int\mathsf F(\de_x,\de_y)\d\pi(x,y) \leq\int c(x,y)\d\pi(x,y).\)
Taking the infimum over $\pi\in\Couplings(\al,\be)$ proves maximality.
\end{proof}
\paragraph{Monge--Kantorovich equivalence.}
\begin{cor}[Monge--Kantorovich equivalence under Brenier]\label{cor-monge-kantorovich-brenier}
Let $\al,\be\in\Pp_2(\RR^d)$, assume that $\al$ is absolutely continuous with respect to Lebesgue measure, and let $c(x,y)=\norm{x-y}^2$. If $T$ is the Brenier map solving Monge's problem, then $\pi=(\Id,T)_\sharp\al$ is the unique optimal coupling solving the Kantorovich problem. In particular, Monge and Kantorovich costs are the same.
\end{cor}
\begin{proof}
The proof of Brenier's theorem shows that the support of any optimal Kantorovich plan is contained in the subdifferential $\partial\phi$ of a convex function $\phi$. When $\al$ has a density, $\phi$ is differentiable $\al$-almost everywhere, so $\partial\phi(x)=\{\nabla\phi(x)\}$ for $\al$-almost every $x$. Thus every optimal coupling is concentrated on the graph of $T=\nabla\phi$ and must equal $(\Id,T)_\sharp\al$. The graph coupling is feasible and optimal, and the two formulations have the same value.
\end{proof}
\paragraph{Kantorovich solution in 1-D.}\label{sec-1d-kantorovich-solution}
\begin{thm}[One-dimensional Kantorovich solution]\label{prop-1d-kantorovich-quantile-coupling}
Let $\al,\be\in\Mm_+^1(\RR)$ and let $h:\RR\to[0,+\infty)$ be convex. Consider the cost $c(x,y)=h(x-y)$. Write $q_\al\eqdef\cumul{\al}^{-1}$ and $q_\be\eqdef\cumul{\be}^{-1}$, and assume that $h(q_\al-q_\be)\in L^1(0,1)$. Then the quantile coupling
\(\pi^\star=(q_\al,q_\be)_\sharp\mathrm{Leb}_{[0,1]} = (\cumul{\al}^{-1},\cumul{\be}^{-1})_\sharp\mathrm{Leb}_{[0,1]}\)
solves the Kantorovich problem~\eqref{eq-mk-generic} for $c(x,y)=h(x-y)$.
In particular, $h(t)=|t|^p$ gives the usual one-dimensional optimal coupling for every $p\geq1$.
\end{thm}
\begin{proof}
The push-forward statement $\pi^\star\in\Couplings(\al,\be)$ follows from Proposition~\ref{prop-quantile-pushforward}.

The key point is the one-dimensional uncrossing inequality. If $x<x'$ and $y>y'$, set $a=x-y$, $\delta=x'-x>0$ and $\eta=y-y'>0$. Convexity of $h$ implies that increments are monotone, hence
\(h(a+\eta)-h(a) \leq h(a+\delta+\eta)-h(a+\delta),\)
which is exactly
\(h(x-y)+h(x'-y') \geq h(x-y')+h(x'-y)\).
Thus removing a crossing never increases the cost. For a finite transport matrix on two ordered grids, if $i<i'$ and $j>j'$ carry crossed masses $\P_{i,j}$ and $\P_{i',j'}$, one may move $\theta=\min(\P_{i,j},\P_{i',j'})$ units of mass from the crossed entries $(i,j)$, $(i',j')$ to the uncrossed entries $(i,j')$, $(i',j)$. The row and column sums are unchanged, and the preceding inequality shows that the cost does not increase. Repeating this elementary uncrossing step yields an ordered plan; on an ordered uniform quantile grid, this is the diagonal plan.

For general measures, apply the same argument in quantile coordinates. Let $\pi\in\Couplings(\al,\be)$. Let $U_\al$ and $U_\be$ be uniform random variables on $(0,1)$, and denote by $K_\al(x,\d r)$ and $K_\be(y,\d s)$ regular conditional laws of $U_\al$ given $q_\al(U_\al)=x$ and of $U_\be$ given $q_\be(U_\be)=y$. Given $(X,Y)\sim\pi$, draw conditionally $R\sim K_\al(X,\cdot)$ and $S\sim K_\be(Y,\cdot)$. The law $\gamma$ of $(R,S)$ has uniform marginals and satisfies $\pi=(q_\al,q_\be)_\sharp\gamma$.

Let $\kappa_M(r)=\max(-M,\min(r,M))$ and set $q_{\al,M}=\kappa_M\circ q_\al$ and $q_{\be,M}=\kappa_M\circ q_\be$. These functions are bounded and nondecreasing. Approximate them almost everywhere by nondecreasing step functions $q_{\al,M}^N$ and $q_{\be,M}^N$ that are constant on the uniform intervals $I_k=((k-1)/N,k/N]$. The matrix
\(G_{k,\ell}^N=\gamma(I_k\times I_\ell)\)
is a coupling of the two uniform histograms. Proposition~\ref{prop-1d-weighted-sweep}, applied to the ordered step values, therefore gives
\(\int_0^1 h(q_{\al,M}^N(r)-q_{\be,M}^N(r))\d r \leq \int_{(0,1)^2}h(q_{\al,M}^N(r)-q_{\be,M}^N(s))\d\gamma(r,s).\)
For fixed $M$, the arguments of $h$ remain in the compact interval $[-2M,2M]$. Dominated convergence as $N\to\infty$ yields the same inequality for $q_{\al,M}$ and $q_{\be,M}$.

Finally, $\kappa_M$ is nondecreasing and $1$-Lipschitz. Hence, for every $x,y\in\RR$, there exists $t_M(x,y)\in[0,1]$ such that
\(\kappa_M(x)-\kappa_M(y)=t_M(x,y)(x-y)\).
Convexity and nonnegativity imply
\(h(t_M(x,y)(x-y)) \leq (1-t_M(x,y))h(0)+t_M(x,y)h(x-y) \leq h(0)+h(x-y).\)
The assumed integrability controls the diagonal term. If the cost of $\pi$ is finite, the same bound controls the right-hand side; if it is infinite, the desired comparison is immediate. Dominated convergence as $M\to\infty$ therefore gives
\(\int_0^1 h(q_\al(r)-q_\be(r))\d r \leq \int h(x-y)\d\pi(x,y)\)
for every $\pi\in\Couplings(\al,\be)$.
\end{proof}
\subsection{Cone of Convex Functions and Monge Gap}
\label{sec-monge-gap}
\paragraph{Definition and optimality certificate.}
\begin{defn}[Monge gap]\label{def-monge-gap}
Let $\al\in\Mm_+^1(\Xx)$, let $T:\Xx\to\Yy$ be measurable, and assume that the following costs are finite. The $c$-Monge gap of $T$ relative to $\al$ is
\eql{\label{eq-monge-gap}
\mathcal M_\al^c(T)
\eqdef
\int c(x,T(x))\d\al(x)-\MK_c(\al,T_\sharp\al).
}
\end{defn}
\begin{prop}[Elementary properties of the Monge gap]\label{prop-monge-gap-properties}
Under the standing Polish-space assumptions, suppose that the relevant integrals are finite. Then $\mathcal M_\al^c(T)\geq0$, with equality if and only if $(\Id,T)_\sharp\al$ is an optimal coupling between $\al$ and $T_\sharp\al$. Moreover,
\eql{\label{eq-monge-gap-self-coupling}
\mathcal M_\al^c(T)
=
\sup_{\eta\in\Couplings(\al,\al)}
\int\big[c(x,T(x))-c(x,T(z))\big]\d\eta(x,z).
}
Finally, the gap is unchanged when $c(x,y)$ is replaced by $c(x,y)+r(x)+s(y)$, provided the modified cost remains admissible and all marginal terms are integrable.
\end{prop}
\begin{proof}
The graph plan $(\Id,T)_\sharp\al$ is feasible for $\MK_c(\al,T_\sharp\al)$, which proves nonnegativity and the equality criterion. Every $\eta\in\Couplings(\al,\al)$ induces the coupling $(x,z)\mapsto(x,T(z))$. Conversely, disintegrating $\al$ along $T$ lifts every coupling in $\Couplings(\al,T_\sharp\al)$ to such a self-coupling. Subtracting the resulting infimum from the graph cost gives~\eqref{eq-monge-gap-self-coupling}. Adding $r(x)+s(y)$ shifts both terms in~\eqref{eq-monge-gap} by $\int r\d\al+\int s\d(T_\sharp\al)$.
\end{proof}
\paragraph{Convex gaps and convex potentials.}
\begin{prop}[Convex Monge gaps for affine-difference costs]\label{prop-affine-difference-monge-gap}
Let $\Yy\subset\RR^m$ be convex and let the nonnegative cost $c$ have the form
\eql{\label{eq-affine-difference-cost}
c(x,y)=a(x)+b(y)-\dotp{A(x)}{y},
}
where $A,T\in L^2(\al;\RR^m)$ and all displayed terms are finite. Then
\eql{\label{eq-affine-difference-monge-gap}
\mathcal M_\al^c(T)
=
\sup_{\eta\in\Couplings(\al,\al)}
\int\dotp{A(x)}{T(z)-T(x)}\d\eta(x,z).
}
Moreover, $\mathcal M_\al^c(T)=0$ if and only if there is a proper lower-semicontinuous convex function $\varphi$ such that
\eql{\label{eq-affine-difference-subgradient-characterization}
T(x)\in\partial\varphi(A(x))
\qquad\text{for $\al$-almost every }x.
}
\end{prop}
\begin{proof}
The marginal terms $a$ and $b$ cancel in~\eqref{eq-monge-gap-self-coupling}, giving~\eqref{eq-affine-difference-monge-gap}; this is a supremum of linear functionals of $T$. Next, disintegrating $\al$ along $A$ lifts every coupling between $A_\sharp\al$ and $T_\sharp\al$. Hence zero gap is equivalent to optimality of $(A,T)_\sharp\al$ for the bilinear cost $(u,y)\mapsto-\dotp{u}{y}$. Up to marginal terms, this is the quadratic cost $\frac12\norm{u-y}^2$. By Rockafellar's characterization, recalled in Section~\ref{sec-cyclical-monotonicity}, its optimal plans are exactly those concentrated on the subdifferential of a proper lower-semicontinuous convex function.
\end{proof}
It also contains the Bregman divergence of Definition~\ref{def-bregman-divergence} with $x$ as base point, \(B_\Phi(y|x)=\Phi(y)-\Phi(x)-\dotp{\nabla\Phi(x)}{y-x},\)
where $\Phi$ is differentiable and convex: here $a(x)=\dotp{\nabla\Phi(x)}{x}-\Phi(x)$, $b(y)=\Phi(y)$ and $A(x)=\nabla\Phi(x)$. The argument order is essential: on a full-dimensional domain, the reversed divergence $B_\Phi(x|y)$ is generally not of the form~\eqref{eq-affine-difference-cost} unless $\Phi$ is quadratic.

\paragraph{One-dimensional order and isotonicity.}
\begin{defn}[Increasing rearrangement relative to a source]\label{def-increasing-rearrangement-map}
Let $\al\in\Mm_+^1(\RR)$ be atomless and let $T:\RR\to\RR$ be measurable. The increasing rearrangement of $T$ relative to $\al$ is
\eql{\label{eq-increasing-rearrangement-map}
T^{\uparrow,\al}
\eqdef
\cumul{T_\sharp\al}^{-1}\circ\cumul{\al}.
}
It is the unique $\al$-almost-everywhere nondecreasing map satisfying $(T^{\uparrow,\al})_\sharp\al=T_\sharp\al$.
\end{defn}
\begin{prop}[One-dimensional Monge-gap formulas]\label{prop-monge-gap-one-dimensional}
Let $\al\in\Mm_+^1(\RR)$ be atomless, let $c(x,y)=h(x-y)$ with $h$ convex, and assume the relevant terms are integrable. Then
\eql{\label{eq-monge-gap-increasing-rearrangement}
\mathcal M_\al^c(T)
=
\int\big[h(x-T(x))-h(x-T^{\uparrow,\al}(x))\big]\d\al(x).
}
If $h$ is strictly convex, this gap vanishes if and only if $T$ has a nondecreasing representative $\al$-almost everywhere.
For the equal-weight atomic measure $\al_n=\frac1n\sum_{i=1}^n\de_{x_i}$ on the increasing uniform grid $x_i=x_1+(i-1)\Delta$ with $\Delta>0$, let $t_i=T(x_i)$ and let $t_{(1)}\leq\cdots\leq t_{(n)}$ be their order statistics. For $c(x,y)=|x-y|^2$,
\eql{\label{eq-monge-gap-uniform-grid}
\mathcal M_{\al_n}^{|\cdot|^2}(T)
=
\frac1n\sum_{i=1}^n\big[(x_i-t_i)^2-(x_i-t_{(i)})^2\big]
=
\frac{2\Delta}{n}\sum_{1\leq i<j\leq n}(t_i-t_j)_+.
}
\end{prop}
\begin{proof}
The one-dimensional rearrangement result of Theorem~\ref{prop-1d-quantile-map}, together with its coupling extension in Theorem~\ref{prop-1d-kantorovich-quantile-coupling}, shows that $T^{\uparrow,\al}$ is optimal between $\al$ and $T_\sharp\al$. If $h$ is strictly convex, atomlessness makes this optimal map unique $\al$-almost everywhere, so zero gap is equivalent to $T=T^{\uparrow,\al}$ $\al$-almost everywhere. For the empirical formula, $\Delta>0$ orders the source points increasingly, so sorting $(t_i)_i$ gives the increasing optimal assignment and hence the first equality. An adjacent exchange of an inversion $u>v$ decreases the average squared cost by $2\Delta(u-v)/n$; bubble sort exchanges every inverted pair once, which gives the second equality.
\end{proof}
\paragraph{Monge-gap-regularized kernel regression.}\label{par-monge-gap-kernel-regression}
The hinge representation above turns transport optimality into a tractable soft monotonicity prior. Let $x_i=x_1+(i-1)\Delta$ with $\Delta>0$ and let $y_i\in\RR$ be noisy observations. For a positive-semidefinite kernel $k$ with reproducing-kernel Hilbert space $\RKHS_k$, consider
\eql{\label{eq-monge-gap-kernel-ridge}
\umin{T\in\RKHS_k}
\left\{
\frac1n\sum_{i=1}^n|T(x_i)-y_i|^2
+\lambda\norm{T}_{\RKHS_k}^2
+\gamma\mathcal M_{\al_n}^{|\cdot|^2}(T)
\right\},
\qquad \lambda>0,\quad \gamma\geq0,
}
where $\al_n=n^{-1}\sum_i\de_{x_i}$. The usual orthogonal-projection argument shows that a minimizer has the representer form $T_\q(\cdot)=\sum_{j=1}^n\q_j k(x_j,\cdot)$. Writing $\K_{i,j}=k(x_i,x_j)$ and $\mathrm y=(y_i)_i$, formula~\eqref{eq-monge-gap-uniform-grid} reduces~\eqref{eq-monge-gap-kernel-ridge} to
\eql{\label{eq-monge-gap-kernel-finite}
\umin{\q\in\RR^n}
\left\{
\frac1n\norm{\K\q-\mathrm y}^2
+\lambda\dotp{\q}{\K\q}
+\frac{2\gamma\Delta}{n}
\sum_{1\leq i<j\leq n}
\big[(\K\q)_i-(\K\q)_j\big]_+
\right\}.
}
This is a finite-dimensional convex program: the first two terms are convex quadratics because $\K\succeq0$, and the last term is a sum of hinges of affine functionals. Thus every point along the regularization path can be computed globally, without alternating between regression and transport.

\subsection{\texorpdfstring{$c$}{c}-Cyclical Monotonicity}
\label{sec-cyclical-monotonicity}
\paragraph{\texorpdfstring{Support and $c$-cyclical monotonicity.}{Support and c-cyclical monotonicity.}}
The support of a coupling is the topological support introduced in Definition~\ref{def:support}, now applied to a Radon measure on $\Xx\times\Yy$. Thus $(x,y)\in\supp(\pi)$ exactly when every open neighborhood of $(x,y)$ has positive $\pi$-mass.

\begin{defn}[$c$-cyclical monotonicity]\label{def:ccm}
A set $\Gamma\subset\Xx\times\Yy$ is $c$-cyclically monotone if, for every $k\geq2$, every finite family $(x_i,y_i)_{i=1}^k\subset\Gamma$ and every permutation $\sigma$ of $\{1,\ldots,k\}$,
\(\sum_{i=1}^k c(x_i,y_i) \leq \sum_{i=1}^k c(x_i,y_{\sigma(i)}).\)
\end{defn}
Any permutation is a product of cycles, so it suffices to verify the inequality for cyclic permutations, \(\sum_{i=1}^k c(x_i,y_i) \leq \sum_{i=1}^k c(x_i,y_{i+1}), \qquad y_{k+1}=y_1.\)
\paragraph{Optimal matching to optimal transport.}
Let the marginals be uniform on $n$ points, $\al=\frac1n\sum_{i=1}^n\delta_{x_i}$ and $\be=\frac1n\sum_{i=1}^n\delta_{y_i}$. By Corollary~\ref{cor-kantorovich-matching}, there exists an optimal plan induced by a permutation. Its support $\Gamma=\{(x_i,y_{\sigma(i)})\}_i$ is $c$-cyclically monotone: otherwise exchanging the finitely many targets along a violating cycle would lower the matching cost.

\begin{thm}[Optimal plans are $c$-cyclically monotone]\label{thm:opt_ccm}
Assume that $c:\Xx\times\Yy\to[0,+\infty)$ is continuous and that the Kantorovich value is finite. For any optimal plan $\pi$ solving~\eqref{eq-mk-generic}, $\supp(\pi)$ is $c$-cyclically monotone.
\end{thm}
\begin{proof}
Suppose that $\supp(\pi)$ is not $c$-cyclically monotone. Then there exist points $(x_i,y_i)_{i=1}^k$ in the support and a permutation $\sigma$ such that $\sum_i c(x_i,y_i)>\sum_i c(x_i,y_{\sigma(i)})$.
By continuity of $c$, after shrinking neighborhoods $U_i\ni x_i$ and $V_i\ni y_i$, the same strict inequality holds uniformly for every choice of points in these neighborhoods:
\(\sum_i c(u_i,v_i) > \sum_i c(u_i,\tilde v_{\sigma(i)}) \qquad (u_i\in U_i,\ v_i\in V_i,\ \tilde v_{\sigma(i)}\in V_{\sigma(i)}).\)
Choose the sets so that $m_i\eqdef\pi(U_i\times V_i)>0$. Taking
\(0<\lambda\leq\left(\sum_{i=1}^k\frac1{m_i}\right)^{-1}\)
ensures that the scaled restrictions
\(\pi_i=\lambda\frac{\pi|_{U_i\times V_i}}{m_i}\)
have common mass $\lambda$ and satisfy $\sum_i\pi_i\leq\pi$, even if some rectangles overlap. Let $\al_i=(P_\Xx)_\sharp\pi_i$ and $\be_i=(P_\Yy)_\sharp\pi_i$, and define $\tilde\pi=\pi-\sum_i\pi_i+\sum_i(\al_i\otimes\be_{\sigma(i)})/\lambda$.
The removed and reinserted first marginals are both $\sum_i\al_i$, and the removed and reinserted second marginals are both $\sum_i\be_i$ because $\sigma$ is a permutation. Hence $\tilde\pi\in\Couplings(\al,\be)$. Integrating the uniform strict inequality against the product probability $\otimes_i(\pi_i/\lambda)$ shows that the reinserted crossed terms have strictly smaller cost than the removed diagonal terms. This contradicts the optimality of $\pi$.
\end{proof}
\paragraph{Monotonicity.}
Assume the optimal plan is induced by a measurable map $T:\Xx\to\Yy$, meaning that $\pi=(\Id,T)_\sharp\al$. There is a set $G\subset\Xx$ of full $\al$-measure such that $(x,T(x))\in\supp(\pi)$ for every $x\in G$. For any $x_1,\ldots,x_k\in G$, cyclical monotonicity reads
\(\sum_{i=1}^k c(x_i,T(x_i)) \leq \sum_{i=1}^k c(x_i,T(x_{i+1})), \qquad x_{k+1}=x_1.\)
\paragraph{One dimension.}
In one space dimension, with cost $c(x,y)=|x-y|^p$, the two-point inequality becomes
\(|x-T(x)|^p+|y-T(y)|^p \leq |x-T(y)|^p+|y-T(x)|^p,\)
For $p>1$, strict convexity implies that a reversed pair $x<y$ and $T(x)>T(y)$ violates this inequality. Hence every optimal map is nondecreasing on the full-measure transport set, recovering the classical monotone rearrangement.

\section{Wasserstein Space}
\label{sec-wasserstein-space}
\subsection{Wasserstein Distances}
\label{sec-wasserstein-distances}
\paragraph{OT defines a distance.}
\begin{lem}[Discrete gluing lemma]\label{lem-gluing-discr}
Given $(\a,\b,\VectMode{c}) \in \simplex_n \times \simplex_p \times \simplex_m$,
let $\P \in \CouplingsD(\a,\b)$ and $\Q \in \CouplingsD(\b,\VectMode{c})$. Then there exists a nonnegative three-way array $\S \in \RR_+^{n \times p \times m}$
such that its two-way marginals satisfy
\(\sum_{k} \S_{i,j,k} = \P_{i,j} \qandq \sum_{i} \S_{i,j,k} = \Q_{j,k}.\)
$\R_{i,k}\eqdef\sum_j\S_{i,j,k}$, belongs to
$\CouplingsD(\a,\VectMode{c})$. For the canonical construction below, this
glued coupling is the twisted matrix product
$\R=\P\diag(1/\b)\Q$, with
$\R_{i,k}=\sum_{j:\b_j>0}\P_{i,j}\Q_{j,k}/\b_j$.
In the matrix notation, $1/\b_j$ is understood as $0$ when $\b_j=0$.
\end{lem}
\begin{proof}
\eqllead{One verifies that}{eq-glued-discr}{
\S_{i,j,k} = \choice{
\frac{\P_{i,j} \Q_{j,k}}{\b_j} \qifq \b_j \neq 0 \\
0 \text{ otherwise}
}
}
has the required marginals. Indeed, if $\b_j \neq 0$, \(\sum_{k} \S_{i,j,k} = \sum_{k} \frac{\P_{i,j} \Q_{j,k}}{\b_j} = \frac{\P_{i,j}}{\b_j}(\Q\ones_m)_j = \frac{\P_{i,j}}{\b_j}\b_j.\)
If $\b_j = 0$, then necessarily $\P_{i,j}=0$ and $\sum_{k} \S_{i,j,k} = 0 = \P_{i,j}$.
The same computation gives the other prescribed marginal:
\[
\sum_i \S_{i,j,k}
=
\choice{
\frac{\Q_{j,k}}{\b_j}\sum_i\P_{i,j}=\Q_{j,k} \qifq \b_j>0\\
0=\Q_{j,k} \qifq \b_j=0.
}
\]
Summing over $j$ then gives the displayed formula for $\R$. Its row and
column sums are $\sum_k\R_{i,k}=\sum_j\P_{i,j}=\a_i$ and
$\sum_i\R_{i,k}=\sum_j\Q_{j,k}=\VectMode{c}_k$, so
$\R\in\CouplingsD(\a,\VectMode{c})$.
\end{proof}
When the cost matrix is the $p$th power of a distance matrix, the discrete Kantorovich value becomes a metric on histograms.

\begin{defn}[Discrete Wasserstein distance]\label{def-discrete-wasserstein-distance}
Let $\distD\in\RR_+^{n\times n}$ be a distance matrix on $\range{n}$ and let $p\geq1$. The discrete $p$-Wasserstein distance between histograms $\a,\b\in\simplex_n$ is
\eql{\label{eq-wass-p-disc}
\WassD_p(\a,\b) \eqdef \MKD_{\distD^p}(\a,\b)^{1/p}.
}
It depends on the chosen ground distance $\distD$.
\end{defn}
\begin{prop}[Metric property of discrete Wasserstein distance]\label{prop-metric-histo}
For every distance matrix $\distD$ on $\range{n}$, Definition~\ref{def-discrete-wasserstein-distance} defines a distance on $\simplex_n$: $\WassD_p$ is symmetric, positive, $\WassD_p(\a,\b)=0$ if and only if $\a = \b$, and it satisfies the triangle inequality
\(\foralls \a,\b,\VectMode{c} \in \simplex_n, \quad \WassD_p(\a,\VectMode{c}) \leq \WassD_p(\a,\b) + \WassD_p(\b,\VectMode{c}).\)
\end{prop}
\begin{proof}
Transposition maps $\CouplingsD(\a,\b)$ bijectively onto $\CouplingsD(\b,\a)$ and preserves the cost because $\distD$ is symmetric. This proves symmetry. If $\a=\b$, the diagonal coupling $\diag(\a)$ has zero cost. Conversely, if an admissible coupling has zero cost, every off-diagonal entry must vanish because $\distD_{i,j}>0$ for $i\neq j$. The coupling is therefore diagonal, and its two marginals coincide. This proves definiteness.

For the triangle inequality, let $\P\in\CouplingsD(\a,\b)$ and $\Q\in\CouplingsD(\b,\VectMode{c})$ be optimal. Lemma~\ref{lem-gluing-discr} provides $\S\in\RR_+^{n\times n\times n}$ with these two marginals. Its first-third marginal
\(\R_{i,k}=\sum_j\S_{i,j,k}\)
belongs to $\CouplingsD(\a,\VectMode{c})$. The ground-metric triangle inequality and Minkowski's inequality give
\[
\begin{aligned}
\WassD_p(\a,\VectMode{c})
&\leq \left(\sum_{i,j,k}\distD_{i,k}^p\S_{i,j,k}\right)^{1/p}\\
&\leq \left(\sum_{i,j,k}(\distD_{i,j}+\distD_{j,k})^p\S_{i,j,k}\right)^{1/p}
\leq \left(\sum_{i,j,k}\distD_{i,j}^p\S_{i,j,k}\right)^{1/p}
+\left(\sum_{i,j,k}\distD_{j,k}^p\S_{i,j,k}\right)^{1/p}\\
&=\WassD_p(\a,\b)+\WassD_p(\b,\VectMode{c}).
\end{aligned}
\]
\end{proof}
\paragraph{Continuous gluing.}
\begin{lem}[Gluing lemma]\label{lem-gluing-general}
Let $(\al,\be,\ga) \in \Mm_+^1(\Xx) \times \Mm_+^1(\Yy) \times \Mm_+^1(\Zz)$, where $\Xx$, $\Yy$ and $\Zz$ are Polish spaces in the sense of Definition~\ref{def-polish-metric-space}.
Given $\pi \in \Couplings(\al,\be)$ and $\xi \in \Couplings(\be,\ga)$, there exists a probability measure
$\sigma \in \Mm_+^1(\Xx \times \Yy \times \Zz)$ such that
\((P_{\Xx,\Yy})_\sharp \sigma = \pi \qandq (P_{\Yy,\Zz})_\sharp \sigma = \xi\)
where $P_{\Xx,\Yy}(x,y,z)=(x,y)$ and $P_{\Yy,\Zz}(x,y,z)=(y,z)$ are the coordinate projections.
\end{lem}
\begin{proof}
Disintegrate the two couplings with respect to their common marginal $\be$ as
$\pi(\d x,\d y)=\pi_y(\d x)\be(\d y)$ and
$\xi(\d y,\d z)=\be(\d y)\xi_y(\d z)$.
The Polish hypotheses guarantee the existence of the probability kernels $y\mapsto\pi_y$ and $y\mapsto\xi_y$. Define their conditional product by
\(\int g\,\d\sigma \eqdef \int_\Yy\int_\Xx\int_\Zz g(x,y,z)\,\xi_y(\d z)\pi_y(\d x)\be(\d y)\)
for every bounded Borel function $g$. Taking $g(x,y,z)=h(x,y)$ recovers $\pi$, while taking $g(x,y,z)=h(y,z)$ recovers $\xi$, proving the two marginal identities.

In the discrete setting, $\pi_{y_j}=\sum_i(\P_{i,j}/\b_j)\de_{x_i}$ and $\xi_{y_j}=\sum_k(\Q_{j,k}/\b_j)\de_{z_k}$ whenever $\b_j>0$. Hence
$\sigma=\sum_{i,j,k}\S_{i,j,k}\de_{(x_i,y_j,z_k)}$, where
$\S_{i,j,k}=\P_{i,j}\Q_{j,k}/\b_j$, with the value zero when $\b_j=0$,
exactly as in~\eqref{eq-glued-discr}.
\end{proof}
\begin{defn}[$p$-Wasserstein space]\label{def-p-wasserstein-space}
Let $(\Xx,d)$ be Polish, let $p\geq1$, and fix $x_0\in\Xx$. The $p$-Wasserstein space is \(\Pp_p(\Xx) \eqdef \enscond{\al\in\Mm_+^1(\Xx)}{\int d(x,x_0)^p\d\al(x)<+\infty},\)
which is independent of the choice of $x_0$.
\end{defn}
If $\al,\be\in\Pp_p(\Xx)$, then the product coupling has finite $p$-cost by the triangle inequality, so the Kantorovich value is finite.

\begin{defn}[Wasserstein distance]\label{def-wasserstein-distance}
Let $(\X,\dist)$ be a Polish metric space and $p\geq1$. On the space $\Pp_p(\X)$ of Definition~\ref{def-p-wasserstein-space}, the $p$-Wasserstein distance is
\eql{\label{eq-defn-wass-dist}
\Wass_p(\al,\be) \eqdef \MK_{\dist^p}(\al,\be)^{1/p}
=
\left(\inf_{\pi\in\Couplings(\al,\be)}
\int_{\X\times\X}\dist(x,y)^p\d\pi(x,y)\right)^{1/p}.
}
It depends on the ground distance $\dist$.
\end{defn}
\begin{prop}[Metric property of the Wasserstein distance]\label{prop-metric-measure}
Definition~\ref{def-wasserstein-distance} defines a distance on $\Pp_p(\X)$: $\Wass_p$ is symmetric and nonnegative, $\Wass_p(\al,\be)=0$ if and only if $\al=\be$, and
\(\foralls (\al,\be,\ga)\in\Pp_p(\X)^3, \qquad \Wass_p(\al,\ga)\leq\Wass_p(\al,\be)+\Wass_p(\be,\ga).\)
\end{prop}
\begin{proof}
The value is finite on $\Pp_p(\X)$ because the product coupling has finite $p$-cost. If $\tau(x,y)=(y,x)$ is the flip map, then $\tau_\sharp\pi\in\Couplings(\be,\al)$ whenever $\pi\in\Couplings(\al,\be)$, and the transport cost is unchanged. This proves symmetry. The diagonal coupling $(\Id,\Id)_\sharp\al$ proves $\Wass_p(\al,\al)=0$.

Conversely, Proposition~\ref{prop-kantorovich-existence-compact} supplies an optimal coupling $\pi$. If $\Wass_p(\al,\be)=0$, then $\dist(x,y)=0$ for $\pi$-almost every $(x,y)$, so $x=y$ almost surely under $\pi$. For every bounded Borel function $f$,
\(\int f\,\d\al = \int f(x)\,\d\pi(x,y) = \int f(y)\,\d\pi(x,y) = \int f\,\d\be,\)
and therefore $\al=\be$.

For the triangle inequality, consider optimal couplings $\pi\in\Couplings(\al,\be)$ and $\xi\in\Couplings(\be,\ga)$,
and glue them into $\sigma$ by Lemma~\ref{lem-gluing-general}. Let $P_{1,3}(x,y,z)=(x,z)$ and define $\eta\eqdef(P_{1,3})_\sharp\sigma$. Then $\eta\in\Couplings(\al,\ga)$. If the two original couplings are induced by maps $T_1$ and $T_2$, this construction is induced by $T_2\circ T_1$; gluing therefore extends composition of deterministic maps.
The triangle inequality and Minkowski's inequality give
\[
\begin{aligned}
\Wass_p(\al,\ga)
&\leq\left(\int\dist(x,z)^p\d\sigma(x,y,z)\right)^{1/p}\\
&\leq\left(\int(\dist(x,y)+\dist(y,z))^p\d\sigma(x,y,z)\right)^{1/p}\\
&\leq\left(\int\dist(x,y)^p\d\sigma(x,y,z)\right)^{1/p}
+\left(\int\dist(y,z)^p\d\sigma(x,y,z)\right)^{1/p}\\
&=\left(\int\dist(x,y)^p\d\pi\right)^{1/p}
+\left(\int\dist(y,z)^p\d\xi\right)^{1/p}
=\Wass_p(\al,\be)+\Wass_p(\be,\ga).
\end{aligned}
\]
\end{proof}
\begin{prop}[Euclidean mean bound and quadratic decomposition]\label{prop-wasserstein-mean-decomposition}
Let $p\geq1$ and $\al,\be\in\Pp_p(\RR^d)$, and denote their means by
$m_\al\eqdef\int_{\RR^d}x\,\d\al(x)$ and
$m_\be\eqdef\int_{\RR^d}x\,\d\be(x)$. Then
$\norm{m_\al-m_\be}\leq\Wass_p(\al,\be)$. For $p=2$, define the centered
translates $\al^0\eqdef(x\mapsto x-m_\al)_\sharp\al$ and
$\be^0\eqdef(y\mapsto y-m_\be)_\sharp\be$.
Then the following exact decomposition holds: \(\Wass_2^2(\al,\be) = \norm{m_\al-m_\be}^2+\Wass_2^2(\al^0,\be^0).\)
\end{prop}
\begin{proof}
For every $\pi\in\Couplings(\al,\be)$,
$m_\al-m_\be=\int_{\RR^d\times\RR^d}(x-y)\,\d\pi(x,y)$.
Jensen's inequality therefore gives
$\norm{m_\al-m_\be}\leq
(\int\norm{x-y}^p\,\d\pi(x,y))^{1/p}$.
Taking the infimum over $\pi$ proves the first claim.

For $p=2$, set
$\pi^0\eqdef((x,y)\mapsto(x-m_\al,y-m_\be))_\sharp\pi$.
This translation defines a bijection from
$\Couplings(\al,\be)$ onto $\Couplings(\al^0,\be^0)$. Writing
$u=x-m_\al$ and $v=y-m_\be$, one has \(\int\norm{x-y}^2\,\d\pi = \norm{m_\al-m_\be}^2+\int\norm{u-v}^2\,\d\pi^0(u,v),\)
because $\int(u-v)\,\d\pi^0(u,v)=0$. Taking the infimum over the corresponding couplings proves the decomposition.
\end{proof}
\paragraph{Interpolation induced by an optimal plan.}
\label{sec-kantorovich-plan-interpolation}
\begin{defn}[$\Wass_2$ geodesic induced by an optimal plan]\label{def-w2-geodesic-induced-by-plan}
Let $\al_0,\al_1\in\Pp_2(\RR^d)$, and let $\pi^\star\in\Couplings(\al_0,\al_1)$ be optimal for $\Wass_2^2(\al_0,\al_1)$. For $t\in[0,1]$, define
\(e_t(x,y)\eqdef(1-t)x+t y, \qquad \al_t\eqdef(e_t)_\sharp\pi^\star.\)
The curve $(\al_t)_{t\in[0,1]}$ is the displacement, or McCann, $\Wass_2$ geodesic induced by $\pi^\star$.
\end{defn}
\begin{prop}[Optimal-plan interpolation is a $\Wass_2$ geodesic]\label{prop-plan-interpolation-w2-geodesic}
Let $(\al_t)_{t\in[0,1]}$ be defined by Definition~\ref{def-w2-geodesic-induced-by-plan}. Then, for every $0\leq s\leq t\leq1$, \(\Wass_2(\al_s,\al_t) = (t-s)\Wass_2(\al_0,\al_1).\)
Thus $t\mapsto\al_t$ is a constant-speed geodesic for the metric $\Wass_2$.
\end{prop}
\begin{proof}
Push the optimal plan $\pi^\star$ forward by $(e_s,e_t)$.
\(\int \norm{z-z'}^2\d\gamma_{s,t}(z,z') = \int \norm{e_t(x,y)-e_s(x,y)}^2\d\pi^\star(x,y) = (t-s)^2\Wass_2^2(\al_0,\al_1).\)
Hence $\Wass_2(\al_s,\al_t)\leq(t-s)\Wass_2(\al_0,\al_1)$. Applying this upper bound to the three pairs $(0,s)$, $(s,t)$ and $(t,1)$, and using the triangle inequality of Proposition~\ref{prop-metric-measure}, gives
\(\Wass_2(\al_0,\al_1) \leq \Wass_2(\al_0,\al_s)+\Wass_2(\al_s,\al_t)+\Wass_2(\al_t,\al_1) \leq \Wass_2(\al_0,\al_1).\)
All inequalities are therefore equalities, in particular the middle segment has the claimed length.
\end{proof}
\paragraph{Comparison with Monge.}
\begin{prop}[Kantorovich as the plan-space relaxation of Monge]\label{prop-kantorovich-relaxation-monge}
Let $(\Xx,\dist)$ be a compact metric space, let $p\geq1$, and let $\al,\be\in\Pp(\Xx)$ with $\al$ atomless. Then
\(\tilde\Wass_p(\al,\be)=\Wass_p(\al,\be).\)
\end{prop}
\begin{proof}
Let $\pi\in\Couplings(\al,\be)$. Choose finite Borel partitions $(A_i)_i$ and $(B_j)_j$ of $\Xx$ whose cells have diameter at most $\varepsilon$. Set $m_{i,j}=\pi(A_i\times B_j)$. Since $\al$ is atomless and $\sum_j m_{i,j}=\al(A_i)$, each $A_i$ can be partitioned into Borel sets $A_{i,j}$ with $\al(A_{i,j})=m_{i,j}$. If $m_{i,j}>0$, then $\be(B_j)>0$, and Proposition~\ref{prop-existence-transport-map-atomless} gives a measurable map
$T_{i,j}:A_{i,j}\to B_j$ satisfying
$(T_{i,j})_\sharp(\al|_{A_{i,j}}/m_{i,j})=\be|_{B_j}/\be(B_j)$.
Define $T$ by $T=T_{i,j}$ on $A_{i,j}$, ignoring null cells. Then $T_\sharp\al=\be$, because the contributions landing in $B_j$ sum to $\be|_{B_j}$, and the graph coupling $(\Id,T)_\sharp\al$ has the same masses $m_{i,j}$ as $\pi$ on all rectangles $A_i\times B_j$.

Let $h\in\Cc(\Xx\times\Xx)$. By uniform continuity on the compact product, the oscillation of $h$ on every rectangle $A_i\times B_j$ goes to zero as $\varepsilon\to0$. Since $\pi$ and $(\Id,T)_\sharp\al$ have the same mass on each rectangle, their integrals against $h$ differ by at most this maximal oscillation. Taking partitions with mesh tending to zero therefore gives graph couplings $(\Id,T_k)_\sharp\al$ converging weakly to $\pi$. Applying this to the continuous function $(x,y)\mapsto\dist(x,y)^p$ gives the convergence of costs.

Define
\[
\mathcal G(\al,\be)
\eqdef
\enscond{(\Id,T)_\sharp\al}{T_\sharp\al=\be}
\subset\Couplings(\al,\be),
\qquad
F_p(\pi)\eqdef\int_{\Xx\times\Xx}\dist(x,y)^p\d\pi(x,y),
\]
and let $G$ be the Monge graph functional which equals $F_p$ on $\mathcal G(\al,\be)$ and $+\infty$ otherwise. The function $F_p$ is weakly continuous on $\Couplings(\al,\be)$ and satisfies $F_p\leq G$. If $H$ is any weakly lower-semicontinuous function with $H\leq G$, then for the graph approximation above,
\(H(\pi)\leq\liminf_k H((\Id,T_k)_\sharp\al) \leq \lim_k F_p((\Id,T_k)_\sharp\al)=F_p(\pi).\)
Thus $F_p$ is the largest weakly lower-semicontinuous minorant of $G$, that is, the lower-semicontinuous envelope. Minimizing this envelope over all couplings gives the Kantorovich value, and the density argument gives the same infimum when restricting to graph couplings.
\end{proof}
\begin{cor}[Lower-semicontinuous envelope of the Monge \(p\)-cost]\label{cor-wasserstein-lsc-envelope-monge-distance}
Let $(\Xx,\dist)$ be a compact metric space and assume that atomless probability measures are dense in $\Pp(\Xx)$ for the $\Wass_p$ topology. Then $(\al,\be)\mapsto\Wass_p(\al,\be)^p$ is the lower-semicontinuous envelope, on $\Pp(\Xx)\times\Pp(\Xx)$ for the product $\Wass_p$ topology, of the extended Monge cost
\[
\tilde\Wass_p(\al,\be)^p
=
\inf_{\substack{T:\Xx\to\Xx\\ T_\sharp\al=\be}}
\int_\Xx \dist(x,T(x))^p\d\al(x),
\]
with the convention $\tilde\Wass_p(\al,\be)^p=+\infty$ if no admissible map exists.
\end{cor}
\begin{proof}
For every admissible map $T$, the graph plan $(\Id,T)_\sharp\al$ is a coupling, hence $\Wass_p(\al,\be)^p\leq \tilde\Wass_p(\al,\be)^p$. Since $\Wass_p^p$ is continuous for the product $\Wass_p$ topology, it is a lower-semicontinuous minorant of the extended Monge cost.

Conversely, let $H$ be any lower-semicontinuous minorant of the extended Monge cost. Fix $(\al,\be)$ and choose atomless $\al_k\to\al$ in $\Wass_p$. By Proposition~\ref{prop-kantorovich-relaxation-monge}, $\tilde\Wass_p(\al_k,\be)^p=\Wass_p(\al_k,\be)^p$. Therefore
\(H(\al,\be) \leq \liminf_k H(\al_k,\be) \leq \lim_k \Wass_p(\al_k,\be)^p = \Wass_p(\al,\be)^p.\)
Thus no larger lower-semicontinuous minorant exists.
\end{proof}
\subsection{Topology and Applications}
\label{sec-wasserstein-topology-applications}
\paragraph{Convergence in law topology.}
On a bounded metric space, all $\Wass_p$ distances define the same topology, although their comparison is generally H\"older rather than bi-Lipschitz.

\begin{prop}[Comparison of Wasserstein distances on bounded spaces]\label{prop-comp-wass-p}
Let $(\Xx,\dist)$ be a bounded Polish metric space, let $1\leq p\leq q<\infty$, and let $\al,\be\in\Pp_q(\Xx)$. If $p<q$, then
\eq{
\Wass_p(\al,\be) \leq \Wass_q(\al,\be) \leq \operatorname{diam}(\Xx)^{\frac{q-p}{q}} \Wass_p(\al,\be)^{\frac{p}{q}},
\qquad
\operatorname{diam}(\Xx)\eqdef\sup_{x,y\in\Xx}\dist(x,y).
}
If $p=q$, the comparison is simply the identity $\Wass_p=\Wass_q$.
\end{prop}
\begin{proof}
The case $p=q$ is immediate, so assume $p<q$.
The left inequality follows from Jensen's inequality, applied to any coupling $\pi$ and to the convex function $r\mapsto r^{q/p}$: \(\pa{\int \dist(x,y)^{p} \d\pi(x,y)}^{\frac{q}{p}} \leq \int \dist(x,y)^{q} \d\pi(x,y).\)
Taking the infimum gives $\Wass_p\leq\Wass_q$. The right inequality follows from
\(\dist(x,y)^q \leq \operatorname{diam}(\Xx)^{q-p} \dist(x,y)^p\)
and taking the infimum over couplings once more.
\end{proof}
\begin{defn}[Weak, or narrow, topology]\label{dfn-weak-conv}
A sequence $(\al_k)_k$ converges weakly, or narrowly, to $\al$ in $\Mm_+^1(\Xx)$, denoted $\al_k \rightharpoonup \al$, if and only if
\(\int_\Xx f\d\al_k\longrightarrow\int_\Xx f\d\al \qquad\text{for every }f\in\Cc_b(\Xx).\)
On a compact space, $\Cc_b(\Xx)=\Cc(\Xx)$ and this is precisely weak$^*$ convergence in the duality between measures and continuous functions. The term weak$^*$ is therefore also used informally in the book when no confusion can arise.
\end{defn}
\paragraph{Strong versus weak topology.}
In the present section we only use the recall that, for a signed difference $\al-\be$, \(\norm{\al-\be}_{\TV}=|\al-\be|(\Xx),\)
\begin{prop}[Total variation as a $0/1$-cost transport value]\label{prop-rel-wass-tv}
Let $\Xx$ be a standard Borel space, let $\al$ and $\be$ be probability measures on $\Xx$, and define the measurable cost $d_0(x,x)=0$ and $d_0(x,y)=1$ for $x\neq y$. Then
\(\MK_{d_0}(\al,\be) = \inf_{\pi\in\Couplings(\al,\be)}\int d_0(x,y)\d\pi(x,y) = \frac{1}{2}\norm{\al-\be}_{\TV}.\)
If $\Xx$ is countable and carries the discrete topology, then $d_0$ is a Polish metric and, for every $p\geq1$, the left-hand side is $\Wass_p(\al,\be)^p$ because $d_0^p=d_0$.
\end{prop}
\begin{proof}
Since $\Xx$ is standard Borel, its diagonal is measurable in $\Xx\times\Xx$; hence $d_0$ and the diagonal submeasure used below are well-defined.
Let $\lambda=\al+\be$, write $a=\d\al/\d\lambda$ and $b=\d\be/\d\lambda$, and define the common part $\eta$ by
\(\frac{\d\eta}{\d\lambda}=\min(a,b)\).
Its residuals $\al'=\al-\eta$ and $\be'=\be-\eta$ are mutually singular and have the same mass
\(r = 1-\eta(\Xx) = \frac12\int|a-b|\d\lambda = \frac12\norm{\al-\be}_{\TV}.\)
For any $\pi\in\Couplings(\al,\be)$, define its diagonal submeasure by
\(\zeta(A)=\pi\bigl(\{(x,x):x\in A\}\bigr)\).
Then $\zeta\leq\al$ and $\zeta\leq\be$, so its density with respect to $\lambda$ is bounded by $\min(a,b)$.
$\int d_0(x,y)\d\pi(x,y)=1-\pi(\{x=y\})\geq r$.
If $r=0$, the diagonal coupling attains this bound. If $r>0$, let $\Delta(x)=(x,x)$ and set
$\pi^\star=\Delta_\sharp\eta+r^{-1}\al^{\prime}\otimes\be^{\prime}$.
This is a coupling of $\al$ and $\be$. Since $\al'\perp\be'$, the product term is concentrated off the diagonal, so its $0/1$ cost is exactly its total mass $r$. Thus $\pi^\star$ attains the lower bound.
\end{proof}
\paragraph{Probabilistic interpretation.}
Probability theory gives weak convergence its most direct interpretation: it compares distributions rather than samplewise realizations. In terms of random vectors, if $X_n \sim \al_n$ and $X \sim \al$ (not necessarily defined on the same probability space), then $\al_n\rightharpoonup\al$ means precisely that $X_n$ converges in law to $X$.

In that setting, $X_n\to X$ almost surely means pointwise convergence outside a null set, while convergence in probability means
\(\foralls \epsilon>0,\qquad \PP(\norm{X_n-X}>\epsilon)\to0.\)
\paragraph{Central limit theorem and OT.}
\begin{defn}[Convolution of measures]\label{def-measure-convolution}
For probability measures $\al,\be$ on $\RR^d$, their convolution is \(\al*\be\eqdef\operatorname{add}_\sharp(\al\otimes\be), \qquad \operatorname{add}(x,y)=x+y.\)
Equivalently, for every bounded continuous $\varphi$, \(\int\varphi\,\d(\al*\be) = \iint\varphi(x+y)\,\d\al(x)\d\be(y).\)
\end{defn}
Thus $\al*\be$ is the law of $X+Y$ when $X$ and $Y$ are independent with laws $\al$ and $\be$. When $\al=\density{\al}\d x$ and $\be=\density{\be}\d x$, their convolution has density
\(\density{\al*\be}(z) = (\density{\al}*\density{\be})(z) = \int_{\RR^d}\density{\al}(x)\density{\be}(z-x)\d x.\)
\paragraph{Wasserstein metrizes weak convergence.}
One can then contrast the strong topology with the Wasserstein distance if $x_n \neq x$, \(\norm{\de_{x_n}-\de_x}_{\TV}=2 \qandq \Wass_p(\de_{x_n},\de_x) = d(x_n,x).\)
\begin{prop}[Wasserstein convergence on Polish spaces]\label{prop-wass-topology-polish}
Let $(\Xx,d)$ be a Polish metric space, let $1\leq p<+\infty$, and let $\al_k,\al\in\Pp_p(\Xx)$. For any reference point $x_0\in\Xx$, the following are equivalent:
\begin{enumerate}
\item $\Wass_p(\al_k,\al)\to0$;
\item $\al_k\rightharpoonup\al$ and the $p$th moments converge,
\(\int d(x,x_0)^p\d\al_k(x) \longrightarrow \int d(x,x_0)^p\d\al(x);\)
\item $\al_k\rightharpoonup\al$ and the $p$th moments are uniformly integrable,
\(\lim_{R\to+\infty}\sup_k \int_{\{d(x,x_0)>R\}}d(x,x_0)^p\d\al_k(x)=0.\)
\end{enumerate}
These conditions do not depend on the choice of $x_0$.
\end{prop}
\begin{proof}
Assume first that $\Wass_p(\al_k,\al)\to0$. Since $\Wass_1\leq\Wass_p$, integration against every bounded $1$-Lipschitz function converges; such functions determine weak convergence on Polish spaces. For any coupling $\pi\in\Couplings(\al_k,\al)$, the reverse triangle inequality in $L^p(\pi)$ gives
\[
\left|
\left(\int d(x,x_0)^p\d\al_k(x)\right)^{1/p}
-
\left(\int d(y,x_0)^p\d\al(y)\right)^{1/p}
\right|
\leq
\left(\int d(x,y)^p\d\pi(x,y)\right)^{1/p}.
\]
Taking the infimum over $\pi$ proves convergence of the $p$th moments, hence (ii).

Conversely, assume (ii). By the Skorokhod representation theorem, there exist random variables $X_k\sim\al_k$ and $X\sim\al$ on a common probability space such that $X_k\to X$ almost surely. In particular, the family $(d(X_k,x_0)^p)_k$ is uniformly integrable. Since
\(d(X_k,X)^p \leq 2^{p-1}\bigl(d(X_k,x_0)^p+d(X,x_0)^p\bigr),\)
the variables $d(X_k,X)^p$ are uniformly integrable and converge almost surely to zero. Vitali's theorem therefore gives $\EE[d(X_k,X)^p]\to0$. The law of $(X_k,X)$ is an admissible coupling, so
\(\Wass_p(\al_k,\al)^p\leq\EE[d(X_k,X)^p]\longrightarrow0.\)

Conversely, assume (iii) and define the bounded continuous truncation $h_R(x)=\min\{d(x,x_0)^p,R^p\}$. Weak convergence gives $\int h_R\d\al_k\to\int h_R\d\al$ for every $R$. The uniform tail bound in (iii), together with the integrability of $d(\cdot,x_0)^p$ under $\al$, then lets $R\to+\infty$ and proves convergence of the untruncated moments. This establishes all three equivalences. Since (i) is independent of the reference point, so are (ii) and (iii).
\end{proof}
\begin{prop}[Comparison with total variation on finite spaces]
Let $(\Xx,d)$ be a finite metric space with at least two points and define
\(d_{\min} \eqdef \min_{x\neq y}d(x,y)>0, \qquad d_{\max} \eqdef \max_{x,y}d(x,y)<+\infty.\)
Then, for every $\al,\be\in\Pp(\Xx)$, \(\frac{d_{\min}}{2}\norm{\al-\be}_{\TV} \leq \Wass_1(\al,\be) \leq \frac{d_{\max}}{2}\norm{\al-\be}_{\TV}.\)
\end{prop}
\begin{proof}
We denote $d_0(x,y)$ the distance such that $d_0(x,x)=0$ and $d_0(x,y)=1$ for $x \neq y$. One has
\(d_{\min} d_0(x,y) \leq d(x,y) \leq d_{\max} d_0(x,y)\),
so that integrating this against any $\pi \in \Couplings(\al,\be)$ and taking the minimum among those $\pi$ gives the result using Proposition~\ref{prop-rel-wass-tv}.
\end{proof}
The constants are optimal: for $\al=\de_x$ and $\be=\de_y$ with $x\neq y$, the comparison reduces to
\(d_{\min} \leq d(x,y) \leq d_{\max}.\)
\subsection{Distributional Robustness and \texorpdfstring{$\Wass_\infty$}{W-infinity}}
\label{sec-dro-wasserstein-infinity}
\paragraph{DRO ambiguity sets.}
Wasserstein distances are also used to define ambiguity sets around an empirical law. Given samples $z_i$ and $\widehat{\al}_n=\frac1n\sum_i\delta_{z_i}$, a distributionally robust optimization (DRO) problem replaces the empirical risk $\frac1n\sum_i \ell_\theta(z_i)$ by
\(\sup_{\be:\,\Wass_p(\be,\widehat{\al}_n)\leq \rho} \int \ell_\theta(z)\d\be(z),\)
\begin{prop}[Duality for empirical Wasserstein DRO]\label{prop-empirical-wasserstein-dro-duality}
Let $(\Zz,d)$ be a Polish metric space, let $1\leq p<\infty$ and $\rho>0$, and let $\widehat{\al}_n=n^{-1}\sum_{i=1}^n\delta_{z_i}$. Let $\ell_\theta:\Zz\to\RR$ be Borel measurable and integrable with respect to $\widehat{\al}_n$. Assume that it has finite upper $p$-growth: for some $z_0\in\Zz$ and $A,B\geq0$,
\(\ell_\theta(z)-\ell_\theta(z_0) \leq A d(z,z_0)^p+B \qquad\text{for every }z\in\Zz.\)
Then the robust value is finite and
\begin{equation}\label{eq-dro-dual-envelope}
\sup_{\substack{\be\in\Pp_p(\Zz)\\\Wass_p(\be,\widehat{\al}_n)^p\leq \rho^p}}
\int \ell_\theta\d\be
=
\min_{\lambda\geq0}
\left\{
\lambda\rho^p
+
\frac1n\sum_{i=1}^n
\sup_{z\in\Zz}\bigl[\ell_\theta(z)-\lambda d(z,z_i)^p\bigr]
\right\}.
\end{equation}
\end{prop}
\begin{proof}
Orient a coupling so that its first marginal is $\widehat{\al}_n$. It then disintegrates as $\pi=n^{-1}\sum_i\delta_{z_i}\otimes\nu_i$, its second marginal is $n^{-1}\sum_i\nu_i$, and its transport cost is $n^{-1}\sum_i\int d(z,z_i)^p\d\nu_i(z)$. Lagranging this single budget constraint with $\lambda\geq0$ gives weak duality and separates the maximization over the $\nu_i$ into the pointwise suprema in~\eqref{eq-dro-dual-envelope}. The reverse inequality is the strong-duality theorem for Wasserstein ambiguity sets under finite upper $p$-growth.

The growth bound and
$d(z,z_0)^p\leq2^{p-1}\bigl(d(z,z_i)^p+d(z_i,z_0)^p\bigr)$ show that every pointwise supremum is finite for all sufficiently large $\lambda$. The dual objective is lower semicontinuous and convex because it is a sum of suprema of affine functions of $\lambda$. Choosing $z=z_i$ in each supremum bounds the full objective from below by
$\lambda\rho^p+n^{-1}\sum_i\ell_\theta(z_i)$, so it is coercive because $\rho>0$. It therefore attains a finite minimum.
\end{proof}
The positive-radius hypothesis is deliberate. At $\rho=0$, the primal value is $n^{-1}\sum_i\ell_\theta(z_i)$. For $p=1$ and an $L_\theta$-Lipschitz loss, the Kantorovich--Rubinstein dual gives the transparent upper bound
\(\sup_{\be:\,\Wass_1(\be,\widehat{\al}_n)\leq \rho} \int \ell_\theta\d\be \leq \frac1n\sum_i\ell_\theta(z_i)+\rho L_\theta,\)
Applied to $c=d^p$, Proposition~\ref{prop-kantorovich-value-curvature} shows that $(\al,\be)\mapsto\Wass_p(\al,\be)^p$ is jointly convex. Thus $\Wass_1$ itself is jointly convex, but $\Wass_p$ need not be convex for $p>1$: on the real line,
\(\Wass_p\big((1-t)\delta_0+t\delta_1,\delta_0\big)=t^{1/p},\)
Therefore, if $z\mapsto\ell_\theta(z)$ is measurable and $\theta\mapsto\ell_\theta(z)$ is convex for every $z$, then
\(\theta\mapsto \sup_{\be:\,\Wass_p(\be,\widehat{\al}_n)\leq\rho} \int \ell_\theta\,\d\be\)
\paragraph{\texorpdfstring{$\Wass_\infty$}{W-infinity} robustness..}
\begin{defn}[$\Wass_\infty$ extended distance]\label{def-wasserstein-infinity}
Let $(\Xx,d)$ be a Polish metric space. For Borel probability measures $\al,\be\in\Pp(\Xx)$, define
\begin{equation}\label{eq-wass-infty}
\Wass_\infty(\al,\be)
\eqdef
\inf_{\pi\in\Couplings(\al,\be)}
\operatorname*{ess\,sup}_{(x,y)\sim\pi} d(x,y)
\end{equation}
with value $+\infty$ if no finite-displacement coupling exists.
\end{defn}
\begin{prop}[Metric and endpoint properties of $\Wass_\infty$]\label{prop-wasserstein-infinity-metric}
For every $\al,\be\in\Pp(\Xx)$ and every $r\geq0$, \(\Wass_\infty(\al,\be)\leq r \quad\Longleftrightarrow\quad \exists\,\pi\in\Couplings(\al,\be) \text{ supported on }\{(x,y):d(x,y)\leq r\}.\)
In particular, the infimum in~\eqref{eq-wass-infty} is attained whenever it is finite, and $\Wass_\infty$ is an extended metric on $\Pp(\Xx)$, hence a finite-valued metric on each finite-distance component. If $d$ is bounded, then
\(\lim_{q\to\infty}\Wass_q(\al,\be)=\Wass_\infty(\al,\be).\)
\end{prop}
\begin{proof}
If $\Wass_\infty(\al,\be)\leq r$, choose $\pi_k\in\Couplings(\al,\be)$ with essential displacement at most $r+1/k$. The coupling set is tight and weakly closed, hence weakly compact by Prokhorov's theorem. Along a subsequence, $\pi_k\rightharpoonup\pi$. For every $m\geq1$, the closed set $F_m=\{d(x,y)\leq r+1/m\}$ has $\pi_k(F_m)=1$ eventually, so Portmanteau gives $\pi(F_m)=1$. Thus $\pi(\cap_mF_m)=1$, proving the nontrivial implication and attainment. The converse is immediate.

Symmetry is inherited from transposing couplings, while a zero-cost coupling is supported on the diagonal and therefore has equal marginals. The triangle inequality is immediate if either term on its right-hand side is infinite. Otherwise, glue the attaining couplings of $(\al,\be)$ and $(\be,\ga)$ using Lemma~\ref{lem-gluing-general}; their $(x,z)$ marginal is supported where $d(x,z)\leq\Wass_\infty(\al,\be)+\Wass_\infty(\be,\ga)$. This proves the extended metric claim.

Assume now that $d$ is bounded. Monotonicity of $L^q$ norms gives $\Wass_q\leq\Wass_\infty$, so $L\eqdef\lim_{q\to\infty}\Wass_q$ exists and satisfies $L\leq\Wass_\infty$. Choose $q_k\to\infty$ and optimal $\pi_k$ for $\Wass_{q_k}$, whose existence follows from Proposition~\ref{prop-kantorovich-existence-compact}, and extract a weak limit $\pi$. For every $r>L$, Markov's inequality yields
\[
\pi_k\{d>r\}
\leq
\left(\frac{\Wass_{q_k}(\al,\be)}{r}\right)^{q_k}
\longrightarrow0.
\]
Since $\{d>r\}$ is open, Portmanteau gives $\pi\{d>r\}=0$. Letting $r\downarrow L$ shows that $\pi$ is supported on $\{d\leq L\}$, and hence $\Wass_\infty\leq L$.
\end{proof}
This limit concerns the distances, not the linear objectives defining $\Wass_p^p$. The endpoint minimizes the worst displacement rather than an average displacement.

\begin{prop}[$\Wass_\infty$ robust envelope around an empirical law]\label{prop-wasserstein-infty-dro}
Let $(\Zz,d)$ be a Polish metric space and let $\rho\geq0$. Let $\widehat{\al}=\sum_{i=1}^n a_i\delta_{z_i}$ with distinct $z_i$, positive $a_i$, and $\sum_i a_i=1$, and assume that the closed balls $\overline B(z_i,\rho)$ are compact. Repeated atoms can always be merged to meet this convention. For any real-valued upper-semicontinuous loss $\ell$,
\(\sup_{\be:\,\Wass_\infty(\be,\widehat{\al})\leq\rho} \int \ell(z)\d\be(z) = \sum_{i=1}^n a_i\sup_{z\in\overline B(z_i,\rho)}\ell(z).\)
\end{prop}
\begin{proof}
If $\Wass_\infty(\be,\widehat{\al})\leq\rho$, Proposition~\ref{prop-wasserstein-infinity-metric} provides a coupling $\pi\in\Couplings(\widehat{\al},\be)$ supported on $\{(x,z):d(x,z)\leq\rho\}$.
Disintegrating $\pi$ with respect to the first marginal gives $\pi=\sum_i a_i\delta_{z_i}\otimes\nu_i$, where each $\nu_i$ is supported in the closed ball $\overline B(z_i,\rho)$ and $\be=\sum_i a_i\nu_i$. Hence
\(\int\ell\,\d\be = \sum_i a_i\int\ell\,\d\nu_i \leq \sum_i a_i\sup_{\overline B(z_i,\rho)}\ell.\)
The reverse inequality follows by choosing, for each $i$, a maximizer $z_i^\star\in\overline B(z_i,\rho)$ and setting $\be=\sum_i a_i\delta_{z_i^\star}$. The coupling $\sum_i a_i\delta_{(z_i,z_i^\star)}$ has essential displacement at most $\rho$, so this $\be$ is feasible and attains the displayed value.
\end{proof}
\subsection{Measure-to-Vector Maps on Wasserstein Space}\label{sec-measure-to-vector-maps}
\paragraph{Functionals and Wasserstein--H\"older regularity.}
A scalar functional is a map
\(f:\Pp_p(\Xx)\longrightarrow\RR.\)
\begin{defn}[Uniform norm on functionals]\label{def-uniform-norm-functionals}
For a bounded functional \(f:\Pp_p(\Xx)\to\RR\), define
\(\norm{f}_\infty \eqdef \sup_{\al\in\Pp_p(\Xx)}|f(\al)|.\)
\end{defn}
There is no loss of principle in starting with scalar values. If \(F=(f_1,\ldots,f_s):\Pp_p(\Xx)\to\RR^s\), the definitions and constructions below apply coordinatewise, and the resulting bounds combine in any norm on \(\RR^s\).

\begin{defn}[Wasserstein--H\"older functional]\label{def-wasserstein-holder-functional}
Let \(p\geq1\) and \(0<\eta\leq1\). A functional \(f:\Pp(\Xx)\to\RR\) is \(\eta\)-H\"older for \(\Wass_p\), with constant \(L\), if
\begin{equation}\label{eq-wasserstein-holder-functional}
|f(\al)-f(\be)|
\leq
L\Wass_p(\al,\be)^\eta
\qquad
\text{for all }\al,\be\in\Pp(\Xx).
\end{equation}
\end{defn}
\paragraph{Examples and non-examples.}
If \(V:\Xx\to\RR\) is \(L\)-Lipschitz, then Kantorovich--Rubinstein duality~\eqref{eq-w1-metric} and the comparison \(\Wass_1\leq\Wass_p\) from Proposition~\ref{prop-comp-wass-p} give
\(f(\al)=\int_{\Xx} V\,\d\al, \qquad |f(\al)-f(\be)|\leq L\Wass_p(\al,\be).\)
\paragraph{Interaction polynomials.}
\begin{defn}[Wasserstein polynomial]\label{def-wasserstein-polynomial}
For an integer \(m\geq1\), a functional on \(\Pp(\Xx)\) is a \emph{Wasserstein polynomial of order at most \(m\)} if it can be written, for some bounded continuous symmetric kernel \(\varphi_m:\Xx^m\to\RR\), as
\begin{equation}\label{eq-wasserstein-interaction-polynomial}
\mathsf P_{\varphi_m}(\al)
\eqdef
\int_{\Xx^m}\varphi_m(x_1,\ldots,x_m)
\,\d\al(x_1)\cdots\d\al(x_m).
\end{equation}
\end{defn}
\paragraph{Empirical particle approximation.}
\begin{prop}[Density of Wasserstein polynomials]\label{prop-density-wasserstein-polynomials}
Let \(\Xx\) be a compact metric space and \(p\geq1\). For every \(\Wass_p\)-continuous functional \(f:\Pp(\Xx)\to\RR\) and every \(\varepsilon>0\), there exists a Wasserstein polynomial \(\mathsf P\) such that
\(\norm{f-\mathsf P}_{\infty} = \sup_{\al\in\Pp(\Xx)}|f(\al)-\mathsf P(\al)| <\varepsilon.\)
\end{prop}
\begin{proof}
Since \(\Xx\) is compact, \(\Pp(\Xx)\) is compact for the weak topology, which coincides with the topology induced by \(\Wass_p\). Wasserstein polynomials are continuous on this space. They contain the constants and are stable under linear combinations, after lifting lower-order kernels by adding dummy variables. They are also stable under products: if \(\mathsf P_\varphi\) and \(\mathsf P_\psi\) have orders \(m\) and \(r\), then
\(\mathsf P_\varphi(\al)\mathsf P_\psi(\al) = \int_{\Xx^{m+r}} \varphi(x_1,\ldots,x_m)\psi(x_{m+1},\ldots,x_{m+r}) \,\d\al^{\otimes(m+r)},\)
which is a Wasserstein polynomial after symmetrizing its kernel. This algebra separates points: if \(\al\neq\be\), there exists \(V\in\mathcal C(\Xx)\) such that \(\int V\,\d\al\neq\int V\,\d\be\), and \(\al\mapsto\int V\,\d\al\) is an order-one polynomial. The real Stone--Weierstrass theorem therefore gives the claimed uniform density in \(\mathcal C(\Pp(\Xx))\).
\end{proof}
\begin{defn}[Particle polynomial]\label{def-particle-polynomial}
Let \(f:\Pp(\Xx)\to\RR\) be bounded and continuous and let \(n\geq1\). Its symmetric \(n\)-particle kernel is
\[
\varphi_n^f(x_1,\ldots,x_n)
\eqdef
f\left(\frac1n\sum_{i=1}^n\de_{x_i}\right),
\]
and its associated \emph{particle polynomial} is
\begin{equation}\label{eq-empirical-particle-polynomial}
\mathsf B_n f(\al)
\eqdef
\int_{\Xx^n}\varphi_n^f(x_1,\ldots,x_n)
\,\d\al^{\otimes n}(x_1,\ldots,x_n)
=
\EE f(\widehat\al_n),
\end{equation}
where \(\widehat\al_n=n^{-1}\sum_{i=1}^n\de_{X_i}\) for i.i.d.\ random variables \(X_i\sim\al\).
\end{defn}
\begin{prop}[Particle-polynomial approximation of H\"older functionals]\label{prop-holder-particle-polynomial}
Let \(\Xx\subset\RR^d\) be compact and let \(f:\Pp(\Xx)\to\RR\) be \(\eta\)-H\"older for \(\Wass_p\), with \(p\geq1\), \(0<\eta\leq1\) and constant \(L\). Then, for every $n\geq1$ and every \(\al\in\Pp(\Xx)\),
\begin{equation}\label{eq-holder-particle-polynomial-general}
\left|\mathsf B_n f(\al)-f(\al)\right|
\leq
L\,\EE\!\left[\Wass_p(\widehat\al_n,\al)^\eta\right].
\end{equation}
Moreover, for every $n\geq2$,
\begin{equation}\label{eq-holder-particle-polynomial-rate}
\norm{\mathsf B_n f-f}_{\infty}
\leq
C_{\Xx,p,d,\eta}L\,r_{n,p,d}^{\,\eta},
\end{equation}
where the empirical Wasserstein scale \(r_{n,p,d}\) is defined in~\eqref{eq-empirical-wasserstein-scale} of Section~\ref{sec-sample-complexity}.
\end{prop}
\begin{proof}
The empirical representation~\eqref{eq-empirical-particle-polynomial} and H\"older regularity give
\[
|\mathsf B_n f(\al)-f(\al)|
\leq
\EE|f(\widehat\al_n)-f(\al)|
\leq
L\,\EE\!\left[\Wass_p(\widehat\al_n,\al)^\eta\right],
\]
which proves~\eqref{eq-holder-particle-polynomial-general}. The empirical Wasserstein moment bound~\eqref{eq-empirical-wasserstein-moment-scale} gives
\(\sup_{\al\in\Pp(\Xx)}\EE[\Wass_p(\widehat\al_n,\al)^p] \leq C_{\Xx,p,d}r_{n,p,d}^{\,p}.\)
For any nonnegative random variable \(Z\) and exponents \(0<\eta\leq p\), Lyapunov's moment inequality states that
\(\EE[Z^\eta]^{1/\eta}\leq\EE[Z^p]^{1/p}\).
It is the monotonicity of \(L^q\) norms on a probability space and follows directly from Jensen's inequality applied to the concave function \(s\mapsto s^{\eta/p}\); see, for instance,. Applying it to \(Z=\Wass_p(\widehat\al_n,\al)\) gives
\[
\EE\!\left[\Wass_p(\widehat\al_n,\al)^\eta\right]
\leq
\EE\!\left[\Wass_p(\widehat\al_n,\al)^p\right]^{\eta/p},
\]
and taking the supremum over \(\al\) concludes the proof.
\end{proof}
\subsection{Measure-to-Measure Maps on Wasserstein Space}\label{sec-measure-to-measure-maps}
\paragraph{Maps on Wasserstein space.}
Once Wasserstein distances provide a topology on probability measures, it is natural to study transformations of probability measures as maps
\(\Phi:\Pp_p(\Xx)\longrightarrow \Pp_p(\Xx)\)
\paragraph{Particle-preserving transport representations.}
The deterministic case is obtained by pushing each input measure through a map that may itself depend on that input measure:
\begin{equation}\label{eq-measure-map-transport-representation}
\Phi(\al) = \Gamma[\al]_\sharp \al,
\qquad
\Gamma[\al]: \Xx \to \Xx.
\end{equation}
Then, for every discrete measure,
\begin{equation}\label{eq-measure-map-preserve-atoms}
\al=\sum_{i=1}^n a_i\de_{x_i}
\quad\Longrightarrow\quad
\Phi(\al)=\sum_{i=1}^n a_i\de_{\Gamma[\al](x_i)}.
\end{equation}
When \(\Xx\subset\RR^d\), the family \(\al\mapsto\Gamma[\al]\) can equivalently be written as the jointly defined vector-valued map
\begin{equation}\label{eq-curried-transport-representative}
\Gamma:\Pp_p(\Xx)\times\Xx\to\RR^d,
\qquad
\Gamma(\al,x)\eqdef\Gamma[\al](x).
\end{equation}
A useful quantitative hypothesis is the uniform estimate
\(\sup_{x\in\Xx} \norm{\Gamma(\al,x)-\Gamma(\be,x)} \leq L\Wass_p(\al,\be)^\eta.\)
\paragraph{Wasserstein stability.}
\begin{prop}[Wasserstein stability of transport representations]\label{prop-measure-map-wass-lipschitz}
Let $(\Xx,d)$ be Polish, let $p\geq1$, let $E\subset\Xx$, and assume that, for all $\al,\be\in\Pp_p(\Xx)$ supported in $E$, the maps $T[\al]:E\to\Xx$ satisfy
\(d(T[\al](x),T[\al](y)) \leq L_x d(x,y) \quad\text{for all }x,y\in E,\)
and
\(\sup_{y\in E} d(T[\al](y),T[\be](y)) \leq L_{\rm law} \Wass_p(\al,\be).\)
Then $\Phi(\al)=T[\al]_\sharp\al$ is $(L_x+L_{\rm law})$-Lipschitz on probability measures supported in $E$: \(\Wass_p(\Phi(\al),\Phi(\be)) \leq (L_x+L_{\rm law})\Wass_p(\al,\be).\)
\end{prop}
\begin{proof}
Let $\pi$ be an optimal coupling between $\al$ and $\be$. The measure $(T[\al],T[\be])_\sharp\pi$ is a coupling between $\Phi(\al)$ and $\Phi(\be)$, hence
\[
\Wass_p(\Phi(\al),\Phi(\be))
\leq
\left(
\int d(T[\al](x),T[\be](y))^p\,\d\pi(x,y)
\right)^{1/p}.
\]
By the triangle inequality, \(d(T[\al](x),T[\be](y)) \leq L_x d(x,y)+L_{\rm law}\Wass_p(\al,\be).\)
Minkowski's inequality gives the claim.
\end{proof}
If $T:\Xx\to\Xx$ is $L$-Lipschitz, then
\(\Wass_p(T_\sharp\al,T_\sharp\be)\leq L\Wass_p(\al,\be).\)
\paragraph{Mean-field attention.}\label{par-mean-field-attention}
In the non-causal, unmasked setting used here, its discrete output is
\begin{equation}\label{eq-discrete-self-attention}
\operatorname{Att}^{(n)}_\theta(x_1,\ldots,x_n)_i
\eqdef
\sum_{j=1}^n a_{i,j}(x_1,\ldots,x_n)Vx_j,
\qquad
a_{i,j}(x_1,\ldots,x_n)
\eqdef
\frac{\exp(\dotp{Qx_i}{Kx_j})}
{\sum_{\ell=1}^n\exp(\dotp{Qx_i}{Kx_\ell})},
\end{equation}
Indeed, if \(\sigma\) is a permutation and \(x_i^\sigma=x_{\sigma(i)}\), reindexing the denominator gives \(a_{i,j}(x_1^\sigma,\ldots,x_n^\sigma)=a_{\sigma(i),\sigma(j)}(x_1,\ldots,x_n)\), and therefore
\(\operatorname{Att}^{(n)}_\theta(x_1^\sigma,\ldots,x_n^\sigma)_i = \operatorname{Att}^{(n)}_\theta(x_1,\ldots,x_n)_{\sigma(i)}.\)
\begin{defn}[Mean-field attention map]\label{def-mean-field-attention-map}
Fix \(p\geq1\), let $Q,K\in\RR^{r\times d}$ and $V\in\RR^{d\times d}$, write \(\theta=(Q,K,V)\), and let \(\al\in\Pp_p(\RR^d)\). Assume that, for \(\al\)-almost every \(x\),
\(0<\int e^{\dotp{Qx}{Kz}}\,\d\al(z)<+\infty \qandq \int e^{\dotp{Qx}{Kz}}\norm{Vz}\,\d\al(z)<+\infty.\)
The \emph{mean-field attention field} is
\begin{equation}\label{eq-mean-field-attention}
\Gamma_\theta[\al](x)
\eqdef
\frac{\int e^{\dotp{Qx}{Kz}}Vz\,\d\al(z)}
{\int e^{\dotp{Qx}{Kz}}\,\d\al(z)}
\end{equation}
and, whenever \(\Gamma_\theta[\al]\in L^p(\al)\), the induced attention map on measures is \(\operatorname{Att}_\theta(\al) \eqdef (\Gamma_\theta[\al])_\sharp\al.\)
\end{defn}
For an empirical token measure, the construction is an exact rewriting of discrete attention:
\[
\Gamma_\theta[\al](x_i)
=
\operatorname{Att}^{(n)}_\theta(x_1,\ldots,x_n)_i,
\qquad
\operatorname{Att}_\theta(\al)
=
\frac1n\sum_{i=1}^n
\de_{\operatorname{Att}^{(n)}_\theta(x_1,\ldots,x_n)_i}.
\]
\begin{prop}[Compact-support attention stability]\label{prop-attention-wass-lipschitz}
Assume $\Xx=\RR^d$. Fix $Q,K\in\RR^{r\times d}$ and $V\in\RR^{d\times d}$, write $\theta=(Q,K,V)$, and fix $R>0$. Let $\mathcal P_R$ be the set of probability measures supported in the Euclidean ball $B(0,R)$. For every $p\geq1$, there is a constant $C_{\theta,p}(R)$ such that
\(\Wass_p(\operatorname{Att}_\theta(\al),\operatorname{Att}_\theta(\be)) \leq C_{\theta,p}(R)\Wass_p(\al,\be), \qquad \al,\be\in\mathcal P_R.\)
Writing $A_R=\norm{Q}_{\rm op}\norm{K}_{\rm op}R^2$, one can take a constant bounded, up to polynomial factors in $R$ and the operator norms of $Q,K,V$, by \(C_{\theta,p}(R)\lesssim e^{2A_R}.\)
This is the exponential-in-score-radius behavior, equivalently $e^{O(R^2)}$ for dot-product scores on $B(0,R)$, refined in.
\end{prop}
\begin{proof}
Write \(A_R=\norm{Q}_{\rm op}\norm{K}_{\rm op}R^2\), and introduce the normalizer and numerator
\(Z_\al(x)=\int e^{\dotp{Qx}{Kz}}\,\d\al(z), \qquad N_\al(x)=\int e^{\dotp{Qx}{Kz}}Vz\,\d\al(z).\)
For \(x,z\in B(0,R)\), one has \(e^{-A_R}\leq e^{\dotp{Qx}{Kz}}\leq e^{A_R}\), and hence \(Z_\al(x)\geq e^{-A_R}\) and \(\norm{\Gamma_\theta[\al](x)}\leq\norm{V}_{\rm op}R\). As functions of \(z\), the scalar kernel \(e^{\dotp{Qx}{Kz}}\) and the vector-valued kernel \(e^{\dotp{Qx}{Kz}}Vz\) have Lipschitz constants bounded by \(e^{A_R}\) times polynomial factors in \(R\) and the operator norms of \(Q,K,V\). Kantorovich--Rubinstein duality, applied to the vector kernel after pairing it with arbitrary unit vectors, therefore gives
\(|Z_\al(x)-Z_\be(x)| + \norm{N_\al(x)-N_\be(x)} \leq C_{\theta}(R)e^{A_R}\Wass_1(\al,\be).\)
Using
\(\Gamma_\theta[\al](x)-\Gamma_\theta[\be](x) = \frac{N_\al(x)-N_\be(x)}{Z_\al(x)} + \Gamma_\theta[\be](x) \frac{Z_\be(x)-Z_\al(x)}{Z_\al(x)}\)
and the preceding bounds yield
\(\sup_{x\in B(0,R)} \norm{\Gamma_\theta[\al](x)-\Gamma_\theta[\be](x)} \leq C_{\theta}(R)e^{2A_R}\Wass_1(\al,\be).\)
For spatial regularity, introduce the attention law
\(\d\omega_{\al,x}(z) \eqdef \frac{e^{\dotp{Qx}{Kz}}}{Z_\al(x)}\,\d\al(z).\)
Differentiating~\eqref{eq-mean-field-attention} gives \(D_x\Gamma_\theta[\al](x) = \int \bigl(Vz-\Gamma_\theta[\al](x)\bigr) (Q^\top Kz)^\top \,\d\omega_{\al,x}(z),\)
so \(\norm{D_x\Gamma_\theta[\al](x)}_{\rm op}\leq2\norm{V}_{\rm op}\norm{Q}_{\rm op}\norm{K}_{\rm op}R^2\), uniformly in \(\al\in\mathcal P_R\). Finally, \(\Wass_1\leq\Wass_p\) on probability measures and the output is supported in \(B(0,\norm{V}_{\rm op}R)\). Proposition~\ref{prop-measure-map-wass-lipschitz}, applied with \(E=B(0,R)\), proves the claim.
\end{proof}
\paragraph{Mass-splitting Markov maps.}
To obtain a map on $\Pp_p(\Xx)$, assume the finite-moment bound
\(\int_{\Xx} d(x,x_0)^p\,K(y,\d x) \leq C\bigl(1+d(y,x_0)^p\bigr)\)
for some $x_0\in\Xx$. The associated linear map is defined weakly by
\begin{equation}\label{eq-measure-map-markov-kernel}
\int_{\Xx} f(x)\,\d\Psi_K(\al)(x)
\eqdef
\int_{\Xx}\int_{\Xx} f(x)\,K(y,\d x)\,\d\al(y).
\end{equation}
\begin{prop}[Wasserstein stability of Markov maps]\label{prop-markov-map-wasserstein-stability}
Let $(\Xx,d)$ be Polish, let $p\geq1$, and let $K$ satisfy the preceding moment bound. If
\(\Wass_p\bigl(K(y,\cdot),K(y',\cdot)\bigr) \leq Ld(y,y') \qquad\text{for all }y,y'\in\Xx,\)
then the map~\eqref{eq-measure-map-markov-kernel} is $L$-Lipschitz: \(\Wass_p\bigl(\Psi_K(\al),\Psi_K(\be)\bigr) \leq L\Wass_p(\al,\be) \qquad\text{for all }\al,\be\in\Pp_p(\Xx).\)
\end{prop}
\begin{proof}
Fix $\pi\in\Couplings(\al,\be)$. On Polish spaces, the measurable-selection theorem for optimal transport provides a probability kernel $(y,y')\mapsto q_{y,y'}$ such that
\(q_{y,y'}\in\Couplings\bigl(K(y,\cdot),K(y',\cdot)\bigr) \qandq \int d(x,x')^p\,\d q_{y,y'}(x,x') = \Wass_p\bigl(K(y,\cdot),K(y',\cdot)\bigr)^p.\)
Integrating this kernel against $\pi$ defines a coupling
\(q(\d x,\d x') \eqdef \int_{\Xx\times\Xx} q_{y,y'}(\d x,\d x')\,\d\pi(y,y')\)
between $\Psi_K(\al)$ and $\Psi_K(\be)$. Therefore
\[
\Wass_p\bigl(\Psi_K(\al),\Psi_K(\be)\bigr)^p
\leq
\int\Wass_p\bigl(K(y,\cdot),K(y',\cdot)\bigr)^p\,\d\pi(y,y')
\leq
L^p\int d(y,y')^p\,\d\pi(y,y').
\]
Taking the infimum over $\pi$ proves the claim.
\end{proof}

\section{Dual Problem}
\label{sec-dual}
\subsection{Discrete Dual}
\label{sec-discrete-dual}
The Kantorovich problem~\eqref{eq-kanto-discr} is a linear program so that one can equivalently compute its value by solving a dual linear program.

\begin{defn}[Admissible potentials]\label{def-admissible-potentials}
For a discrete problem with marginal sizes $n,m$ and cost matrix $\C\in\RR^{n\times m}$, a pair $(\fD,\gD)\in\RR^n\times\RR^m$ is admissible if it lies below the cost:
\eql{\label{eq-feasible-potential}
\PotentialsD(\C) \eqdef \enscond{
(\fD,\gD) \in \RR^n \times \RR^m
}{ \foralls (i,j) \in \range{n} \times \range{m}, \fD_i+\gD_j \leq \C_{i,j} }.
}
Equivalently, $\fD\oplus\gD\leq\C$ entrywise. In particular, the feasible set depends on the cost matrix, not on the marginal weights.
\end{defn}
\begin{prop}[Discrete Kantorovich duality]\label{prop-duality-discr}
One has
\eql{\label{eq-dual}
\MKD_\C(\a,\b) =
\umax{(\fD,\gD) \in \PotentialsD(\C)} \dotp{\fD}{\a} + \dotp{\gD}{\b}.
}
\end{prop}
\begin{proof}
For completeness, introduce multipliers $(\fD,\gD)$ for the two marginal constraints. The Lagrangian is
\eql{\label{eq-mk-lagr}
\mathcal L(\P,\fD,\gD)
=\dotp{\C}{\P}+\dotp{\a-\P\ones_m}{\fD}
+\dotp{\b-\P^\top\ones_n}{\gD}.
}
For fixed potentials, the dual function is \(\inf_{\P\geq0}\mathcal L(\P,\fD,\gD) =\dotp{\a}{\fD}+\dotp{\b}{\gD} +\inf_{\P\geq0}\dotp{\C-(\fD\oplus\gD)}{\P}.\)
The last infimum is
\[
\inf_{\P\geq0}\dotp{\Q}{\P}
=
\begin{cases}
0 & \text{if }\Q\geq0,\\
-\infty & \text{otherwise},
\end{cases}
\]
and is therefore finite exactly when $(\fD,\gD)\in\PotentialsD(\C)$. The primal polytope is nonempty because it contains the product coupling $\a\otimes\b$, and it is compact; hence the primal value is finite and attained. Finite-dimensional linear-programming strong duality then identifies this value with the maximum of the dual function and also gives dual attainment.
\end{proof}
\begin{prop}[Discrete complementary slackness]\label{prop-discrete-complementary-slackness}
For every $\P\in\CouplingsD(\a,\b)$ and $(\fD,\gD)\in\PotentialsD(\C)$, \(\dotp{\C}{\P}-\dotp{\fD}{\a}-\dotp{\gD}{\b} = \sum_{i,j}\P_{i,j}\bigl(\C_{i,j}-\fD_i-\gD_j\bigr) \geq0.\)
The plan and potentials are respectively primal and dual optimal if and only if
\eql{\label{eq-mk-pd-rel}
\Supp(\P)
\subset
\enscond{(i,j)\in\range{n}\times\range{m}}{\fD_i+\gD_j=\C_{i,j}}.
}
\end{prop}
\begin{proof}
The marginal constraints give
$\dotp{\fD}{\a}+\dotp{\gD}{\b}
=\sum_{i,j}\P_{i,j}(\fD_i+\gD_j)$, which proves the identity and nonnegativity. If both feasible objects are optimal, their values agree by Proposition~\ref{prop-duality-discr}; every nonnegative summand with $\P_{i,j}>0$ must therefore vanish. Conversely, the support condition makes the gap zero, so weak duality forces both objects to be optimal.
\end{proof}
\subsection{General Formulation}
\label{sec-dual-general}
\paragraph{Continuous duality.}
\begin{defn}[Continuous admissible potentials]\label{def-continuous-admissible-potentials}
For $c\in\Cc(\Xx\times\Yy)$, define
\eql{\label{eq-dfn-pot-dual}
\Potentials(c)
\eqdef
\enscond{(f,g)\in\Cc(\Xx)\times\Cc(\Yy)}
{f(x)+g(y)\leq c(x,y)\ \text{for all }(x,y)}.
}
Its elements are admissible continuous potentials.
\end{defn}
\begin{prop}[Kantorovich duality]\label{prop-kantorovich-duality-general}
Assume that $\Xx$ and $\Yy$ are compact metric spaces and that $c\in\Cc(\Xx\times\Yy)$. Then
\eql{\label{eq-dual-generic}
\MK_\c(\al,\be) =
\umax{(\f,\g) \in \Potentials(\c)}
\int_\X \f(x) \d\al(x) + \int_\Y \g(y) \d\be(y),
}
Here, $(\f,\g)$ is a pair of continuous functions, often called ``Kantorovich potentials''.
More general Polish-space versions require explicit lower-semicontinuity and integrability hypotheses; value duality and dual attainment are then separate statements. Proposition~\ref{prop-quadratic-dual-attainment} below records the noncompact quadratic case used in Brenier's theorem.
\end{prop}
\begin{proof}
Weak duality is immediate: if $f(x)+g(y)\leq c(x,y)$ and $\pi\in\Couplings(\al,\be)$, then
\(\int f\d\al+\int g\d\be = \int (f(x)+g(y))\d\pi(x,y) \leq \int c\d\pi.\)
Taking the supremum over admissible potentials and the infimum over couplings gives the first inequality.

For the reverse inequality, write $V=\MK_\c(\al,\be)$ and consider the convex cone
\(\mathcal K \eqdef \bigl\{(\pi_1,\pi_2,\textstyle\int c\d\pi+r): \pi\in\Mm_+(\X\times\Y),\ r\geq0\bigr\}\)
in $\Mm(\X)\times\Mm(\Y)\times\RR$, equipped with the weak topologies on the measure factors. This cone is closed. Indeed, convergence of the first marginal implies a uniform bound on the masses of the underlying nonnegative measures. Compactness of $\X\times\Y$ then gives a convergent subnet $\pi_k\rightharpoonup\pi$; continuity of marginalization and of $\pi\mapsto\int c\d\pi$ identifies the limit with an element of $\mathcal K$, while the limiting epigraph remainder remains nonnegative.

For every $t<V$, the point $z_t=(\al,\be,t)$ does not belong to $\mathcal K$. Strong separation of this point from the closed convex cone gives a continuous linear functional. Because $\mathcal K$ is a cone containing the origin, the functional can be oriented so that, for some $F\in\Cc(\X)$, $G\in\Cc(\Y)$ and $\lambda\in\RR$,
\(\int F\d\al+\int G\d\be+\lambda t<0 \leq \int F\d\eta+\int G\d\zeta+\lambda s \qquad \text{for all }(\eta,\zeta,s)\in\mathcal K.\)
Testing $(0,0,r)\in\mathcal K$ gives $\lambda\geq0$, while testing the element generated by $\pi=\de_{(x,y)}$ and $r=0$ gives
$F(x)+G(y)+\lambda c(x,y)\geq0$. The multiplier cannot vanish: if $\lambda=0$, integrating this last inequality against $\al\otimes\be$ would contradict the strict inequality at $z_t$. Thus $\lambda>0$, and
\(f\eqdef-F/\lambda,\qquad g\eqdef-G/\lambda\)
satisfy $f\oplus g\leq c$ and
$\int f\d\al+\int g\d\be>t$. Letting $t\uparrow V$ proves the reverse inequality.

Take a maximizing sequence $(f_k,g_k)$. Successively replace it by \(\widetilde g_k(y)=\min_x\{c(x,y)-f_k(x)\}, \qquad \widetilde f_k(x)=\min_y\{c(x,y)-\widetilde g_k(y)\}.\)
Each replacement preserves feasibility and cannot decrease the objective. Fixing $x_0\in\X$, use the gauge freedom
$(\widetilde f_k,\widetilde g_k)\mapsto
(\widetilde f_k-a_k,\widetilde g_k+a_k)$
to impose $\widetilde f_k(x_0)=0$. Uniform continuity of $c$ shows that $\widetilde g_k$ inherits a common modulus in the second variable and $\widetilde f_k$ a common modulus in the first. The gauge gives a uniform bound on $\widetilde f_k$, and the first envelope then gives one on $\widetilde g_k$. Arzel\`a--Ascoli yields uniformly convergent subsequences. Their limit remains admissible and attains the dual supremum.
\end{proof}
\paragraph{Complementary slackness.}
\begin{prop}[Continuous complementary slackness]\label{prop-continuous-complementary-slackness}
For $\pi\in\Couplings(\al,\be)$ and $(f,g)\in\Potentials(\c)$, \(\int c\d\pi-\int f\d\al-\int g\d\be = \int\bigl(c(x,y)-f(x)-g(y)\bigr)\d\pi(x,y) \geq0.\)
The coupling and potentials are respectively primal and dual optimal if and only if this gap vanishes. In that case,
\eql{\label{eq-mk-pd-rel-cont}
\supp(\pi)
\subset
\enscond{(x,y)\in\X\times\Y}{f(x)+g(y)=c(x,y)}.
}
\end{prop}
\begin{proof}
The marginal identities give the equality, and admissibility gives nonnegativity. A zero gap is equivalent to equality of the feasible primal and dual values, hence to joint optimality by Proposition~\ref{prop-kantorovich-duality-general}. Finally, the slack is continuous and nonnegative. If it were positive at a point of $\supp(\pi)$, it would be bounded below by a positive constant on a neighborhood of positive $\pi$-mass, contradicting its zero integral.
\end{proof}
The discrete case~\eqref{eq-dual} corresponds to the dual vectors being samples of the continuous potentials, namely $(\fD_i,\gD_j)=(\f(x_i),\g(y_j))$. Proposition~\ref{prop-continuous-complementary-slackness} is the exact analogue of Proposition~\ref{prop-discrete-complementary-slackness}.

The same optimality condition is even more explicit if one plots the \emph{slack}
\(s(x,y) \eqdef c(x,y)-f(x)-g(y).\)
Dual feasibility is the statement $s\geq0$, while complementary slackness says that an optimal coupling can only live where this slack vanishes.

\subsection{\texorpdfstring{$c$-transforms}{c-transforms}}
\label{sec-c-transfo}
\paragraph{\texorpdfstring{Best-response potentials and the $c$-transform.}{Best-response potentials and the c-transform.}}
Keeping a dual potential $\f$ fixed, one can maximize in closed form over the second potential in the dual problem~\eqref{eq-dual-generic}, which leads one to consider
\(\usup{\g \in \Cc(\Yy)} \enscond{ \int g \d \be }{ \foralls (x,y), \g(y) \leq c(x,y) - \f(x) }.\)
For each \(y\), feasibility imposes the upper bound
\(\g(y)\leq\inf_x\{c(x,y)-\f(x)\}\). Since \(\be\) is nonnegative, an optimal \(\g\) must saturate this bound \(\be\)-almost everywhere. This pointwise best response leads to the following definition.
\begin{defn}[$c$-transform]\label{def-c-transform}
For a function $f:\Xx\to\RR\cup\{-\infty\}$ that is not identically $-\infty$, its $c$-transform is
\eql{\label{eq-c-transform}
\foralls y \in \Y, \quad
\f^\c(y) \eqdef \uinf{x \in \X} \c(x,y) - \f(x).
}
Here $c(x,y)-(-\infty)=+\infty$. Because $c$ is finite and $f$ is finite somewhere, $f^c$ takes values in $\RR\cup\{-\infty\}$. For a function $g:\Yy\to\RR\cup\{-\infty\}$ that is not identically $-\infty$, the $\bar c$-transform associated with $\bar c(y,x)=c(x,y)$ is
\(\foralls x \in \X, \quad \g^{\bar\c}(x) \eqdef \uinf{y \in \Y} \c(x,y) - \g(y).\)
It likewise takes values in $\RR\cup\{-\infty\}$. A transform can be identically $-\infty$; compositions such as $f^{c\bar c}$ are used only when each intermediate transform is not identically $-\infty$. Under this qualification no indeterminate extended-real subtraction occurs. When $\Xx,\Yy$ are compact, $c$ is continuous and the input potential is real-valued and continuous, these infima are finite minima and the transformed potential is continuous.
\end{defn}
\begin{prop}[$c$-transforms solve the semi-relaxed problems]\label{prop-c-transform-semi-relaxed}
Assume the compactness and continuity hypotheses of Proposition~\ref{prop-kantorovich-duality-general}. For fixed $f\in\Cc(\X)$, the maximizers of the dual objective over all $g\in\Cc(\Y)$ such that $f\oplus g\leq c$ are exactly the functions satisfying $g=f^c$ $\be$-almost everywhere. Symmetrically, for fixed $g$, the maximizers over $f$ are exactly the functions satisfying $f=g^{\bar c}$ $\al$-almost everywhere.
\end{prop}
\begin{proof}
The constraint $f(x)+g(y)\leq c(x,y)$ for every $x$ is equivalent to
\(g(y)\leq\inf_x\{c(x,y)-f(x)\}=f^c(y)\).
Thus $f^c$ is feasible and pointwise dominates every feasible $g$. Integrating against $\be$ proves optimality, and equality holds exactly when $g=f^c$ $\be$-almost everywhere. The other assertion follows by exchanging the two marginals.
\end{proof}
\begin{prop}[Modulus stability of $c$-transforms]\label{prop-c-transform-lipschitz}
Suppose that $\omega$ is a modulus such that
\(|c(x,y)-c(x,y')| \leq\omega\bigl(d_\Y(y,y')\bigr) \qquad\text{for all }x,y,y'.\)
Then every finite $f^c$ has modulus $\omega$. Symmetrically, $g^{\bar c}$ inherits any uniform modulus of $c$ in its first variable. In particular, a uniform $L$-Lipschitz bound on the cost yields an $L$-Lipschitz bound on the corresponding transform.
\end{prop}
\begin{proof}
For each $x$, set $F_x(y)=c(x,y)-f(x)$ and $F(y)=f^c(y)=\inf_x F_x(y)$. Since all the functions $F_x$ share the modulus $\omega$,
\[
|F(y)-F(y')|
=
\left|\inf_x F_x(y)-\inf_x F_x(y')\right|
\leq
\sup_x |F_x(y)-F_x(y')|
\leq
\omega\bigl(d_\Y(y,y')\bigr).
\]
The proof in the first variable is identical.
\end{proof}
\paragraph{\texorpdfstring{Discrete $c$-transform.}{Discrete c-transform.}}
If $\al=\sum_{i=1}^n a_i\de_{x_i}$ and $\be=\sum_{j=1}^m b_j\de_{y_j}$, the same definition applies to the finite spaces $\X=\{x_i\}_i$ and $\Y=\{y_j\}_j$. The dual potentials are then vectors $(\fD,\gD)\in\RR^n\times\RR^m$, as in Section~\ref{sec-discrete-dual}, and the cost is the matrix $\C\in\RR^{n\times m}$ with entries
\(\C_{i,j}=c(x_i,y_j).\)
With the present sign convention $f\oplus g\leq c$, the discrete transform of $\fD$ is the vector $\fD^\C\in\RR^m$ defined by \((\fD^\C)_j = \min_{1\leq i\leq n} \C_{i,j}-\fD_i.\)
This is a minimum because feasibility imposes the upper bound $\gD_j\leq \C_{i,j}-\fD_i$ for every $i$. Thus the largest feasible best response is the minimum of these upper bounds; formulas with a maximum correspond to a different sign convention, not to the dual constraint~\eqref{eq-dual}. With $\bar\C\eqdef\C^\top$, the opposite transform of $\gD$ is
\(\gD^{\bar\C}\eqdef \gD^{\C^\top}\in\RR^n, \qquad (\gD^{\bar\C})_i = \min_{1\leq j\leq m} \C_{i,j}-\gD_j.\)
Thus the continuous best-response operation reduces exactly to taking column minima for $\fD^\C$, or row minima for $\gD^{\bar\C}$, in shifted cost matrices.

\paragraph{Euclidean case.}
The Euclidean quadratic cost is the model case where $c$-transforms become ordinary convex conjugates after removing the quadratic terms. This is the algebraic bridge between Kantorovich duality and Brenier maps.

For any $\pi\in\Couplings(\al,\be)$, \(\int c(x,y)\d\pi(x,y) = \frac12\int\norm{x}^2\d\al(x) +\frac12\int\norm{y}^2\d\be(y) -\int \dotp{x}{y}\d\pi(x,y).\)
\begin{defn}[Legendre--Fenchel transform and biconjugate]\label{def-legendre-fenchel-transform}
For $h:\RR^d\to\RR\cup\{+\infty\}$, define
\(h^*(y)\eqdef\sup_{x\in\RR^d}\{\dotp{x}{y}-h(x)\}, \qquad h^{**}\eqdef(h^*)^*.\)
\end{defn}
More precisely, if $\phi(x)=\frac12\norm{x}^2-f(x)$, then
\(f^c(y) = \frac12\norm{y}^2-\phi^*(y), \qquad f^{c\bar c}(x) = \frac12\norm{x}^2-\phi^{**}(x).\)
On the full Euclidean space, a nontrivial extended-valued function $f$ is $c$-closed precisely when $x\mapsto\frac12\norm{x}^2-f(x)$ is proper, lower-semicontinuous and convex. When $f$ is finite-valued, this is the usual statement that $f$ is $1$-semiconcave.

\begin{prop}[Quadratic dual attainment on $\RR^d$]\label{prop-quadratic-dual-attainment}
Let $\al,\be\in\Pp_2(\RR^d)$ and $c(x,y)=\norm{x-y}^2/2$. The dual supremum over pointwise feasible pairs in $L^1(\al)\times L^1(\be)$ is attained by Borel potentials $f,g:\RR^d\to\RR\cup\{-\infty\}$ that form a closed pair in the pointwise sense
\(f=g^{\bar c}, \qquad g=f^c, \qquad \MK_c(\al,\be)=\int f\d\al+\int g\d\be.\)
They can be represented by a proper lower-semicontinuous convex function $\phi:\RR^d\to\RR\cup\{+\infty\}$ through
\(f(x)=\frac12\norm{x}^2-\phi(x), \qquad g(y)=\frac12\norm{y}^2-\phi^*(y),\)
with the convention that a finite number minus $+\infty$ is $-\infty$.
For every optimal coupling $\pi$, the Fenchel contact relation
\(\phi(x)+\phi^*(y)=\dotp{x}{y}, \qquad\text{equivalently}\qquad y\in\partial\phi(x),\)
holds $\pi$-almost everywhere.
\end{prop}
\begin{proof}
The quadratic cost is finite and continuous and satisfies $0\leq c(x,y)\leq\norm{x}^2+\norm{y}^2$. In particular, its integral under $\al\otimes\be$ is finite, so the integrable dual-attainment theorem gives pointwise feasible Borel representatives $(f_0,g_0)$ whose classes lie in $L^1(\al)\times L^1(\be)$ and attain the dual value.

Set $\phi_0(x)=\norm{x}^2/2-f_0(x)$. This function is proper because $f_0$ is finite $\al$-almost everywhere. Choosing a point $y_0$ where $g_0(y_0)$ is finite, dual feasibility gives the affine minorant
\(\phi_0(x)\geq \dotp{x}{y_0}-\frac12\norm{y_0}^2+g_0(y_0).\)
Let $\phi=\phi_0^{**}$ and define
\(f(x)=\frac12\norm{x}^2-\phi(x), \qquad g(y)=\frac12\norm{y}^2-\phi^*(y).\)
The Fenchel--Moreau theorem shows that $\phi$ is proper, lower-semicontinuous and convex, with $\phi\leq\phi_0$ and $\phi^*=\phi_0^*$. Thus $f_0\leq f$. Moreover, feasibility of $(f_0,g_0)$ is equivalent to
\(\phi_0(x)+\frac12\norm{y}^2-g_0(y)\geq\dotp{x}{y},\)
so conjugating in $x$ gives $\phi_0^*(y)\leq\norm{y}^2/2-g_0(y)$ and hence $g_0\leq g$. Fenchel's inequality shows that $(f,g)$ is feasible. Choose $x_1,y_1$ such that $f(x_1)$ and $g(y_1)$ are finite. Feasibility then gives
$g(y)\leq c(x_1,y)-f(x_1)$ and $f(x)\leq c(x,y_1)-g(y_1)$. These quadratic upper bounds are integrable under the opposite marginal; together with $f_0\leq f$ and $g_0\leq g$, they show that $(f,g)\in L^1(\al)\times L^1(\be)$. Its objective is no smaller than the already optimal value and no larger by weak duality, so it is optimal. The conjugacy calculation above and $\phi=\phi^{**}$ give $g=f^c$ and $f=g^{\bar c}$ pointwise.

Finally, for any optimal $\pi$, the nonnegative slack $c-f-g$ is integrable and has integral zero. It therefore vanishes $\pi$-almost everywhere, which is exactly the displayed Fenchel equality and hence the subdifferential relation.
\end{proof}
\paragraph{Why hard alternating optimization stops.}
Starting from an admissible pair, the updates
\((f,g)\longmapsto(f,f^c)\longmapsto(f^{c\bar c},f^c)\)
For a proper function $f$ with an affine minorant, the Legendre transform satisfies $f^{***}=f^*$, while $f^{**}$ is the greatest lower-semicontinuous convex minorant of $f$. Analogous identities explain why exact alternating $c$-transform updates stop after one full cycle.

\begin{prop}[Algebra of $c$-transforms]\label{prop-c-transform-algebra}
Writing $\f^{c\bar c}\eqdef(\f^c)^{\bar c}$, the following pointwise identities hold whenever the transforms are well defined:
\[
\textnormal{(i) } f \leq f' \Rightarrow f^{c}\geq f'^{c},
\qquad
\textnormal{(ii) } f^{\c\bar{\c}} \geq f,
\qquad
\textnormal{(iii) } g^{\bar{\c}\c} \geq g,
\qquad
\textnormal{(iv) } \f^{\c\bar{\c}\c}=\f^{\c}.
\]
\end{prop}
\begin{proof}
The order reversal in (i) follows directly from the minus sign in the definition. For every $x$ and $y$, \(f^c(y)=\inf_{x'}\{c(x',y)-f(x')\} \leq c(x,y)-f(x),\)
so $c(x,y)-f^c(y)\geq f(x)$. Taking the infimum over $y$ proves (ii), and exchanging the two spaces proves (iii). Finally, (ii) and the order reversal imply
$f^c\geq f^{c\bar c c}$, whereas (iii) applied to $g=f^c$ gives
$f^{c\bar c c}\geq f^c$. This proves (iv).
\end{proof}
Identity~(iv) shows that the next best response leaves the pair obtained after one full cycle unchanged:
\eql{\label{eq-iter-c-trans}
(f^{c\bar c},f^c)
\longmapsto
(f^{c\bar c},f^{c\bar c c})
=
(f^{c\bar c},f^c).
}
Thus exact block maximization reaches a coordinatewise fixed point after one full cycle. The resulting pair is dual feasible, but it need not maximize the joint dual objective because its value still depends on the arbitrary initialization $f$.

\paragraph{\texorpdfstring{$c$-concave functions.}{c-concave functions.}}
The ranges of the two $c$-transform operators define the natural classes of dual potentials.

\begin{defn}[$c$-concave functions]\label{def-c-concave-functions}
A function on $\Yy$ is $c$-concave if it equals $f^c$ for some function $f$ on $\Xx$. A function on $\Xx$ is $\bar c$-concave if it equals $g^{\bar c}$ for some function $g$ on $\Yy$.
\end{defn}
\begin{defn}[Finite-valued $c$-concave class]\label{def-finite-c-concave-class}
The finite-valued $c$-concave class on $\Yy$ is \(\operatorname{Conc}_c(\Y) \eqdef \enscond{u:\Y\to\RR}{u=f^c\text{ for some }f:\X\to\RR\cup\{-\infty\}}.\)
\end{defn}
\begin{prop}[Stability under tropical barycenters]\label{prop-c-concave-tropical-barycenter}
If $u_r=f_r^c\in\operatorname{Conc}_c(\Y)$ and $a_r\in\RR$ for $1\leq r\leq R$, then
\[
\min_{1\leq r\leq R}\{u_r+a_r\}
=
\left(\max_{1\leq r\leq R}\{f_r-a_r\}\right)^c
\in\operatorname{Conc}_c(\Y).
\]
\end{prop}
\begin{proof}
Set $f=\max_r(f_r-a_r)$. For every $y\in\Yy$, \(f^c(y) = \inf_x\min_r\{c(x,y)-f_r(x)+a_r\} = \min_r\inf_x\{c(x,y)-f_r(x)+a_r\} = \min_r\{u_r(y)+a_r\}.\)
The middle equality holds because the minimum ranges over finitely many indices: both sides equal the infimum over $(x,r)\in\Xx\times\{1,\ldots,R\}$.
\end{proof}
Thus $\operatorname{Conc}_c(\Y)$ is a min-plus, or tropical, convex cone: tropical addition is pointwise minimum, while tropical scalar multiplication adds a constant. This universal closure does not imply stability under ordinary averages.

For $c_\tau(x,y)=\norm{x-y}^2/(2\tau)$, one has \(u\in\operatorname{Conc}_{c_\tau}(\RR^d) \quad\Longleftrightarrow\quad y\mapsto u(y)-\frac{\norm{y}^2}{2\tau}\text{ is concave},\)
Finally, on any metric space $(\X,d)$, \(\operatorname{Conc}_d(\X)=\enscond{u:\X\to\RR}{\Lip(u)\leq1},\)

\section{\texorpdfstring{Semi-discrete and $\Wass_1$}{Semi-discrete and W1}}
\label{sec-semidiscr-w1}
\subsection{Semi-dual}\label{sec-semi-dual}
\paragraph{General measure semi-dual.}
\begin{defn}[Full dual and semi-dual functionals]\label{def-full-dual-functional}
Under the compactness and continuity assumptions of Chapter~\ref{sec-dual}, define
\eql{\label{eq-full-dual-functional}
\Dd_0(f,g)
\eqdef
\begin{cases}
\displaystyle \int_\Xx f\d\al+\int_\Yy g\d\be,
& (f,g)\in\Potentials(\c),\\
-\infty, & \text{otherwise}.
\end{cases}
}
\eql{\label{eq-semi-dual}
\Ee_0(g)
\eqdef
\Dd_0(g^{\bar c},g)
=
\int_\Xx g^{\bar c}\d\al+\int_\Yy g\d\be.
}
\end{defn}
Thus the dual problem~\eqref{eq-dual-generic} is $\MK_\c(\al,\be)=\max_{f,g}\Dd_0(f,g)$. For fixed $g$, feasibility is equivalent to $f\leq g^{\bar c}$. Consequently,
\(\MK_\c(\al,\be) = \sup_{g\in\Cc(\Yy)}\Ee_0(g), \qquad \Ee_0(g)=\sup_{f\in\Cc(\Xx)}\Dd_0(f,g).\)
Partial maximization of a jointly concave objective preserves concavity, so $\Ee_0$ is concave. The semi-dual is invariant under $g\mapsto g+s$: indeed, $(g+s)^{\bar c}=g^{\bar c}-s$ and both measures have unit mass.

\paragraph{Discrete semi-dual.}
When both measures are discrete, the same elimination produces a finite-dimensional concave objective whose active affine pieces are selected by the discrete $c$-transform.

Let $\al=\sum_{i=1}^n\a_i\de_{x_i}$, $\be=\sum_{j=1}^m\b_j\de_{y_j}$, and $\C_{i,j}=c(x_i,y_j)$, where $\a\in\simplex_n$ and $\b\in\simplex_m$. We use the same notation for functions and vectors: in the discrete setting,
\[
\Dd_0(\fD,\gD)
\eqdef
\begin{cases}
\dotp{\fD}{\a}+\dotp{\gD}{\b},
& \fD\oplus\gD\leq\C,\\
-\infty, & \text{otherwise}.
\end{cases}
\]
Eliminating the source potential through the discrete $\bar C$-transform described in Section~\ref{sec-c-transfo} gives
\eql{\label{eq-discrete-semi-dual}
\MKD_\C(\a,\b)
=
\umax{\gD\in\RR^m}\Ee_0(\gD),
\qquad
\Ee_0(\gD)
\eqdef
\Dd_0(\gD^{\bar\C},\gD)
=
\sum_{i=1}^n \a_i(\gD^{\bar\C})_i
+
\sum_{j=1}^m \b_j\gD_j,
}
with \((\gD^{\bar\C})_i = \min_{1\leq j\leq m}(\C_{i,j}-\gD_j).\)
\begin{prop}[Supergradients of the discrete semi-dual]\label{prop-discrete-semidual-supergradient}
For any selection $\sigma_\gD(i)\in\argmin_j(\C_{i,j}-\gD_j)$, define $\widehat\b_j(\gD)\eqdef\sum_{i:\,\sigma_\gD(i)=j}\a_i$. Then $\b-\widehat\b(\gD)$ is a supergradient of $\Ee_0$ at $\gD$. If every row has a unique minimizer, $\Ee_0$ is differentiable at $\gD$ and $\nabla\Ee_0(\gD)=\b-\widehat\b(\gD)$.
\end{prop}
\begin{proof}
Let $\hD\in\RR^m$. Since $\sigma_\gD(i)$ is active at $\gD$,
\[
\min_j(\C_{i,j}-\hD_j)
\leq
\C_{i,\sigma_\gD(i)}-\hD_{\sigma_\gD(i)}
=
\min_j(\C_{i,j}-\gD_j)
-\bigl(\hD_{\sigma_\gD(i)}-\gD_{\sigma_\gD(i)}\bigr).
\]
Multiplying by $\a_i$, summing over $i$, and adding $\dotp{\b}{\hD-\gD}$ gives $\Ee_0(\hD)\leq\Ee_0(\gD)+\dotp{\b-\widehat\b(\gD)}{\hD-\gD}$, which is the supergradient inequality. Uniqueness of all active indices makes the same affine branch active in a neighborhood of $\gD$, yielding the gradient formula.
\end{proof}
\subsection{Auction Algorithm}
\label{sec-auction-dual-ascent}
The auction algorithm is derived from coordinate maximization of the semi-dual of the linear assignment problem. Its practical form uses bidder-specific dual-weight updates that cross the selected row's next indifference threshold by $\epsilon$; this controlled relaxation prevents jamming at nonsmooth ties.

\paragraph{Coordinate ascent and discrete Laguerre cells.}
Specializing~\eqref{eq-discrete-semi-dual} to uniform weights gives
\eql{\label{eq-auction-semidual}
\Ee_0(\gD)
=
\frac1n\sum_{i=1}^n (\gD^{\bar\C})_i
+
\frac1n\sum_{j=1}^n \gD_j,
\qquad
(\gD^{\bar\C})_i
=
\min_{1\leq j\leq n}(\C_{i,j}-\gD_j).
}
Thus $(\gD^{\bar\C},\gD)$ is dual feasible. This is precisely the discrete $\bar C$-transform described in Section~\ref{sec-c-transfo}, with the same sign convention as in the general and semi-discrete formulations.

The discrete Laguerre cell of target $j$ is
\eql{\label{eq-auction-discrete-laguerre}
\Laguerre_j^{\rm D}(\gD)
\eqdef
\enscond{i\in\range{n}}{
\C_{i,j}-\gD_j
\leq
\C_{i,k}-\gD_k
\text{ for every }k
}.
}
When every row has a unique minimizer, differentiation of~\eqref{eq-auction-semidual} gives
\eql{\label{eq-auction-coordinate-derivative}
\frac{\partial}{\partial \gD_j}\Ee_0(\gD)
=
\frac{1-\abs{\Laguerre_j^{\rm D}(\gD)}}{n}.
}
Hence an overfull cell calls for a decrease of its dual weight. More generally, complementary slackness (Proposition~\ref{prop-discrete-complementary-slackness}) shows that a perfect matching in the contact graph $i\in\Laguerre_j^{\rm D}(\gD)$ certifies optimality of $\gD$ in~\eqref{eq-auction-semidual}; conversely, assignment duality and complementary slackness produce such a matching at every maximizer.

For $i\in\Laguerre_j^{\rm D}(\gD)$, define the bid needed to make row $i$ indifferent between $j$ and its best alternative,
\eql{\label{eq-auction-bid}
\operatorname{bid}_j(\gD,i)
\eqdef
\min_{k\neq j}(\C_{i,k}-\gD_k)
-
(\C_{i,j}-\gD_j).
}
When $\Laguerre_j^{\rm D}(\gD)$ is nonempty, the largest maximizing decrement along the negative $j$th coordinate is given by
\(
\max_{i\in\Laguerre_j^{\rm D}(\gD)}\operatorname{bid}_j(\gD,i)
\)
\paragraph{Bids and relaxed contacts.}
Exact coordinate maximization can stall at a nonsmooth tie, so auction deliberately moves a bidder $\epsilon$ beyond its next indifference point. Suppose row $i$ is currently unassigned. Let $j_0$ and $j_1$ be its best and second-best targets for the current dual weights,
\eql{\label{eq-auction-reduced-costs}
j_0\in\argmin_j(\C_{i,j}-\gD_j),
\qquad
j_1\in\argmin_{j\neq j_0}(\C_{i,j}-\gD_j),
}
The bid decreases only the winning dual weight by
\eql{\label{eq-auction-dual-update}
\Delta_i
\eqdef
\operatorname{bid}_{j_0}(\gD,i)+\epsilon
=
(\C_{i,j_1}-\gD_{j_1})
-
(\C_{i,j_0}-\gD_{j_0})
+
\epsilon,
\qquad
\gD_{j_0}\leftarrow\gD_{j_0}-\Delta_i.
}
\begin{defn}[$\epsilon$-complementary slackness]\label{def-auction-eps-cs}
A partial permutation matrix $\P\in\{0,1\}^{n\times n}$ has at most one nonzero entry in each row and column. Such a matrix and target potential $\gD$ satisfy $\epsilon$-complementary slackness if
\eql{\label{eq-auction-epsilon-cs}
\P_{i,j}=1
\quad\Longrightarrow\quad
\C_{i,j}-\gD_j
\leq
(\gD^{\bar\C})_i+\epsilon.
}
\end{defn}
To expose the dual mechanism geometrically, define the auction reduced costs
\(\mathrm{r}_{i,j}(\gD) \eqdef \C_{i,j}-\gD_j-(\gD^{\bar\C})_i \geq 0.\)
At any intermediate state, write
\(
M\eqdef\{(i,j):\P_{i,j}=1\}
\)
\paragraph{Fixed tolerance.}
\begin{prop}[Fixed-$\epsilon$ auction convergence and complexity]\label{prop-auction-termination}
Let
\(R_\C\eqdef\max_{i,j}\C_{i,j}-\min_{i,j}\C_{i,j}.\)
Started from $\gD=0$, the fixed-tolerance auction method terminates after at most
\(
n\bigl(\lfloor R_\C/\epsilon\rfloor+1\bigr)
\)
bids. It returns a permutation matrix $\P$ satisfying $\epsilon$-complementary slackness and
\eql{\label{eq-auction-cost-certificate}
0
\leq
\frac1n\dotp{\C}{\P}
-
\umin{\P'\in\mathcal P_n^{\mathrm{perm}}}
\frac1n\dotp{\C}{\P'}
\leq
\epsilon.
}
With dense scans for the two smallest reduced costs, the algorithm requires
\(
O\bigl(n^2(1+R_\C/\epsilon)\bigr)
\)
arithmetic and comparison operations and $O(n^2)$ storage, including $\C$. If $\C$ has integer entries and $\epsilon<1/n$, the returned assignment is exactly optimal.
\end{prop}
\begin{proof}
At every bid,~\eqref{eq-auction-dual-update} decreases one target potential by at least $\epsilon$. Once a target is assigned, it remains assigned: a later bid can change its owner but never leaves it empty. While the algorithm is incomplete, there is therefore an unassigned target whose potential is still zero. If the selected target $j_0$ is already assigned, comparison with such a zero-potential target gives $\gD_{j_0}\geq-R_\C$; if it is unassigned, its potential is itself zero. Thus a target can be selected at most $\lfloor R_\C/\epsilon\rfloor+1$ times. Summing over the $n$ targets proves finite termination and the bid bound. At termination, every row and column has exactly one nonzero entry, so $\P$ is a permutation matrix.

The $\epsilon$-complementary-slackness invariant was established above. Summing it over the returned permutation gives the first inequality below. The second follows by summing the feasible dual inequalities $(\gD^{\bar\C})_i+\gD_j\leq\C_{i,j}$ against any optimal assignment:
\[
\frac1n\dotp{\C}{\P}
\leq
\frac1n\sum_i(\gD^{\bar\C})_i
+
\frac1n\sum_j\gD_j
+
\epsilon
\leq
\umin{\P'\in\mathcal P_n^{\mathrm{perm}}}
\frac1n\dotp{\C}{\P'}
+
\epsilon,
\]
which proves~\eqref{eq-auction-cost-certificate}. Each bid scans one row of $\C$ and otherwise uses constant work if the current owner of every target is stored, proving the complexity bound. For integer costs, the unnormalized assignment gap is an integer bounded by $n\epsilon<1$, hence it vanishes.
\end{proof}
\paragraph{\texorpdfstring{$\epsilon$-scaling.}{epsilon-scaling.}}
\begin{prop}[Complexity of $\epsilon$-scaling]\label{prop-auction-epsilon-scaling}
For a final tolerance $\eta>0$, the $\epsilon$-scaling auction method returns a permutation matrix satisfying $\eta$-complementary slackness and~\eqref{eq-auction-cost-certificate} with $\epsilon$ replaced by $\eta$. If $R_\C>0$, it uses at most
$1+\lceil\log_2^+(R_\C/\eta)\rceil$ auction phases, where
$\log_2^+(s)\eqdef\max\{0,\log_2 s\}$. Its dense worst-case complexity is
\[
O\left(n^3\left(1+\log_+\frac{R_\C}{\eta}\right)\right),
\qquad
\log_+(s)\eqdef\max\{0,\log s\}.
\]
When $R_\C=0$, the algorithm terminates immediately. If $\C$ has integer entries and $\eta<1/n$, the returned assignment is exactly optimal.
\end{prop}
\begin{proof}
Assume $R_\C>0$. The first tolerance is $\epsilon_0=\max\{R_\C,\eta\}$. If $\epsilon_0=\eta$, Proposition~\ref{prop-auction-termination} directly gives the claim with one phase. Otherwise, the first phase starts from the zero potential and costs $O(n^2)$ operations because $R_\C/\epsilon_0=1$.

Consider a later phase with tolerance $\epsilon$, initialized from a target potential produced at tolerance $\lambda$. Expressed in the present dual-potential convention, the warm-start estimate of shows that the new phase decreases each component of $\gD$ by at most $n(\lambda+\epsilon)$. Since every bid decreases one component by at least $\epsilon$, the phase performs at most
\(n^2\left(1+\frac{\lambda}{\epsilon}\right)\)
bids. Consecutive tolerances in the $\epsilon$-scaling auction method satisfy $\lambda/\epsilon\leq2$, so a warm-started phase uses $O(n^2)$ bids and $O(n^3)$ dense operations.

Halving from $R_\C$ to $\eta$ gives the announced phase count. The final phase enforces $\eta$-complementary slackness, and Proposition~\ref{prop-auction-termination} gives the $\eta$-optimality certificate. Its integer-cost conclusion also applies with $\epsilon=\eta$.
\end{proof}
\subsection{Semi-discrete}
\paragraph{Discrete target and Laguerre cells.}
A case of particular interest is when
\(
\be = \sum_{j=1}^m \b_j \de_{y_j}
\)
One can adapt Definition~\ref{def-c-transform} of the $\bar c$-transform to this setting by restricting the minimization to the finite support $\{y_j\}_{j=1}^m$ of $\be$. Equivalently, one applies the same definition with the discrete target space $\Y=\{y_j\}_{j=1}^m$ and identifies a vector $\gD\in\RR^m$ with the function $g:\Y\to\RR$ defined by $g(y_j)=\gD_j$:
\eql{\label{eq-disc-c-transfo}
\foralls \gD \in \RR^m, \;
\foralls x \in \Xx, \quad
\gD^{\bar \c}(x) \eqdef \umin{j \in \range{m}} \c(x,y_j) - \gD_j.
}
Crucially, using the discrete $\bar c$-transform, when $\be$ is a discrete measure, yields a finite-dimensional optimization,
\eql{\label{eq-semi-dual-discr}
\MK_\c(\al,\be) =
\umax{\gD \in \RR^m}
\Ee_0(\gD) \eqdef \Dd_0(\gD^{\bar\c},\gD) =
\int_\Xx \gD^{\bar \c}(x) \d\al(x) + \sum_j \gD_j \b_j.
}
The objective satisfies $\Ee_0(\gD+s\ones)=\Ee_0(\gD)$ for every $s\in\RR$. One may therefore impose the gauge $\sum_j\gD_j=0$ without changing the problem.

\begin{defn}[Laguerre cells and power diagrams]\label{def-laguerre-power-cells}
For sites $(y_j)_{j=1}^m$ and weights $\gD\in\RR^m$, the Laguerre cell associated with $y_j$ is
\eql{\label{eq-laguerre-cells}
\Laguerre_{j}(\gD) \eqdef \enscond{x \in \Xx}{ \foralls j' \neq j, \c(x,y_j) - \gD_j \leq \c(x,y_{j'}) - \gD_{j'} }.
}
The cells cover $\Xx$; after arbitrary tie-breaking on common boundaries, they induce a disjoint partition. When $c(x,y)=\norm{x-y}^2$, this Laguerre decomposition is also called a power diagram. If $\gD$ is constant, it reduces to the ordinary Voronoi diagram.
\end{defn}
\paragraph{One-dimensional semi-discrete OT.}
In one dimension, the Laguerre geometry of Definition~\ref{def-laguerre-power-cells} becomes explicit. Consider the quadratic cost $c(x,y)=\abs{x-y}^2$, assume that $\al$ has a density and that $y_1<\cdots<y_m$, and set $B_0=0$ and $B_j=\sum_{k=1}^j\b_k$. Consequently, the optimal cells agree, up to $\al$-null sets, with the quantile intervals:
\eql{\label{eq-semidiscrete-1d-laguerre-cells}
\foralls j\in\range{m},\qquad
\al\!\left(
\Laguerre_j(\gD^\star)
\mathbin{\triangle}
\left(\cumul{\al}^{-1}(B_{j-1}),\cumul{\al}^{-1}(B_j)\right]
\right)=0.
}
Here the harmless endpoint conventions are $\cumul{\al}^{-1}(0)=-\infty$ and $\cumul{\al}^{-1}(1)=+\infty$. Thus, after tie-breaking on the null boundaries, the optimal semi-discrete map sends the $j$-th quantile interval to $y_j$; this is the Laguerre-cell realization of the monotone rearrangement in Theorem~\ref{prop-1d-quantile-map}.

\begin{prop}[Closed-form one-dimensional semi-discrete transport]\label{prop-semidiscrete-1d-quantile}
Let \(\al=\rho_\al\,\d x\in\Pp_2(\RR)\), and define its integrated quantile by
\(
M_\al(r) \eqdef \int_0^r \cumul{\al}^{-1}(s)\d s.
\)
For
\(
\be=m^{-1}\sum_{j=1}^m\de_{y_j}
\)
with \(y_1<\cdots<y_m\), one has
\eql{\label{eq-semidiscrete-1d-w2-uniform}
\Wass_2(\al,\be)^2
=
\int_\RR x^2\d\al(x)
+\frac{1}{m}\sum_{j=1}^m y_j^2
-2\sum_{j=1}^m y_j
\left[M_\al\left(\frac{j}{m}\right)-M_\al\left(\frac{j-1}{m}\right)\right].
}
\end{prop}
\begin{proof}
Proposition~\ref{prop-wass-quantile-1d} gives \(\Wass_2(\al,\be)^2=\int_0^1\abs{\cumul{\al}^{-1}(r)-\cumul{\be}^{-1}(r)}^2\d r\). The target quantile satisfies \(\cumul{\be}^{-1}(r)=y_j\) for almost every \(r\in((j-1)/m,j/m]\). Expanding the square, using \(\int_0^1\abs{\cumul{\al}^{-1}(r)}^2\d r=\int_\RR x^2\d\al(x)\), and integrating \(\cumul{\al}^{-1}\) over each interval gives the stated formula.
\end{proof}
\paragraph{Mass balance.}
Splitting the integral in the semi-dual objective~\eqref{eq-semi-dual-discr} over this partition gives
\eql{\label{eq-semi-disc-energy}
\Ee_0(\gD) = \sum_{j=1}^m \int_{\Laguerre_{j}(\gD)} \pa{ c(x,y_j) - \gD_j } \d\al(x) + \dotp{\gD}{\b}.
}
\begin{prop}[Gradient of the semi-discrete dual]\label{prop-semidiscrete-dual-gradient}
If the minimizing index in~\eqref{eq-disc-c-transfo} is unique for $\al$-almost every $x$, then $\Ee_0$ is differentiable at $\gD$ and
\(\foralls j \in \range{m}, \quad \nabla\Ee_0(\gD)_j = \b_j - \int_{\Laguerre_{j}(\gD)} \d\al.\)
\end{prop}
\begin{proof}
For $\al$-almost every $x$, the minimizing index in $\min_j c(x,y_j)-\gD_j$ is unique. If this index is $j(x)$, then the directional derivative in a direction $h\in\RR^m$ is
\[
\left.\frac{\d}{\d t}\right|_{t=0}
\min_j \bigl(c(x,y_j)-(\gD_j+t h_j)\bigr)
=
-h_{j(x)}.
\]
The corresponding difference quotients are bounded by $\norm{h}_\infty$. Dominated convergence therefore gives \(\d\Ee_0(\gD)[h] = -\sum_j h_j\int_{\Laguerre_j(\gD)}\d\al +\sum_j h_j\b_j,\)
which is the announced gradient formula.
\end{proof}
At a differentiable maximizer, first-order optimality gives
\(
\al(\Laguerre_j(\gD))=\b_j
\)
Its graph lies in the contact set
\(
\gD^{\bar c}(x)+\gD_j=c(x,y_j),
\)
Expanding the squared distance shows that the active cell at $x$ minimizes the affine function
\(
x\mapsto -2\dotp{x}{y_j}+\norm{y_j}^2-\gD_j.
\)
\paragraph{Stochastic optimization.}
The semi-discrete formulation~\eqref{eq-semi-disc-energy} is useful because the objective is an expectation with respect to $\al$,
\eql{\label{eq-semi-disc-energy-entropy}
\Ee_0(\gD)
=
\int_\Xx E(\gD,x)\d\al(x)
=
\EE_X(E(\gD,X)),
\qquad
E(\gD,x)\eqdef \gD^{\bar c}(x)+\dotp{\gD}{\b}.
}
Away from cell boundaries, the stochastic gradient of the integrand is
\[
\nabla_{\gD}E(\gD,x)
=
\left(\b_j-\ones_{\Laguerre_j(\gD)}(x)\right)_{j=1}^m,
\]
Starting from $\gD^{(0)}=0$, stochastic gradient ascent draws $x^{(\ell)}\sim\al$ and performs
\eql{\label{eq-sgd}
\gD^{(\ell+1)}
\eqdef
\gD^{(\ell)}+\tau^{(\ell)}\nabla_{\gD}E(\gD^{(\ell)},x^{(\ell)}).
}
Equivalently, if
\(j^{(\ell)}\in\argmin_j\bigl(c(x^{(\ell)},y_j)-\gD_j^{(\ell)}\bigr),\)
then the coordinate update is \(\gD_j^{(\ell+1)} = \gD_j^{(\ell)} + \tau^{(\ell)}\bigl(\b_j-\ones_{\{j=j^{(\ell)}\}}\bigr).\)
For almost-sure stochastic-approximation convergence, one typically chooses
\eql{\label{eq-step-size-sgd}
\sum_{\ell=0}^{+\infty}\tau^{(\ell)}=+\infty,
\qquad
\sum_{\ell=0}^{+\infty}(\tau^{(\ell)})^2<+\infty,
}
\begin{prop}[Averaged stochastic semi-dual rate]\label{prop-semidiscrete-sgd-rate}
Let $\gD^\star$ maximize $\Ee_0$, set $R=\norm{\gD^{(0)}-\gD^\star}_2$, and suppose the stochastic supergradients are conditionally unbiased and satisfy $\norm{\nabla_\gD E(\gD,x)}_2\leq G$. For a horizon $L\geq1$, use the constant step $\tau=R/(G\sqrt L)$ and the average $\bar\gD_L=L^{-1}\sum_{\ell=0}^{L-1}\gD^{(\ell)}$. If $R=0$, the conclusion is immediate; otherwise,
\(\Ee_0(\gD^\star)-\EE\bigl[\Ee_0(\bar\gD_L)\bigr] \leq \frac{RG}{\sqrt L}.\)
\end{prop}
\begin{proof}
Conditionally on $\gD^{(\ell)}$, the stochastic supergradient is unbiased. Concavity and the squared-distance recursion give
\[
2\tau\,\EE\!\left[\Ee_0(\gD^\star)-\Ee_0(\gD^{(\ell)})\right]
\leq
\EE\norm{\gD^{(\ell)}-\gD^\star}_2^2
-
\EE\norm{\gD^{(\ell+1)}-\gD^\star}_2^2
+
\tau^2G^2.
\]
Summing from $0$ to $L-1$, discarding the final nonnegative squared distance, and using concavity at $\bar\gD_L$ yields
\[
\Ee_0(\gD^\star)-\EE\!\left[\Ee_0(\bar\gD_L)\right]
\leq
\frac{R^2}{2\tau L}+\frac{\tau G^2}{2}.
\]
The stated choice of $\tau$ gives the result.
\end{proof}
\subsection{Optimal Quantization}
\label{sec-optimal-quantization}
\paragraph{Free masses and prescribed weights.}
\begin{defn}[Optimal quantization problem]\label{def-optimal-quantization}
For $\al\in\Pp_p(\RR^d)$ and $m\in\NN\setminus\{0\}$, define
\eql{\label{eq-optimal-quantization}
\Qq_m(\al)
\eqdef
\umin{Y=(y_j)_{j=1}^m,\ \b\in\simplex_m}
\Wass_p(\al,\sum_{j=1}^m \b_j\de_{y_j}).
}
Zero weights and coincident codepoints are allowed, so the approximating measure has at most $m$ atoms.
\end{defn}
\begin{prop}[Quantization rate and curse of dimensionality]\label{prop-quantization-rate}
Let $\al$ be supported in a bounded subset $\Omega\subset\RR^d$ and assume that $\al=\rho\,\d x$ with $\rho\leq\rho_+<+\infty$. Then, for fixed $p\geq1$, there exist constants $0<c\leq C<+\infty$ such that
\(c\,m^{-1/d} \leq \Qq_m(\al) \leq C\,m^{-1/d}.\)
\end{prop}
\begin{proof}
Enclose $\Omega$ in a cube. With $k=\lfloor m^{1/d}\rfloor$, subdivide this cube into $k^d\leq m$ congruent subcubes and place one codepoint in each subcube meeting $\Omega$. Every source point then lies within distance $C/k\leq C'm^{-1/d}$ of a codepoint. The nearest-codepoint coupling proves the upper bound; unused codepoints may be assigned zero mass.

For the lower bound, fix any set $Y$ of at most $m$ codepoints and write $d_Y(x)=\min_j\norm{x-y_j}$. If $\omega_d$ is the volume of the Euclidean unit ball, then
\(
\al(\{d_Y\leq t\})\leq \rho_+m\omega_dt^d.
\)
Set $t_0=(2\rho_+m\omega_d)^{-1/d}$. Then $\al(\{d_Y>t\})\geq1/2$ for $0<t<t_0$, and the layer-cake formula gives \(\int d_Y(x)^p\d\al(x) = \int_0^{+\infty} p t^{p-1}\al(\{d_Y>t\})\d t \geq \frac{t_0^p}{2} \simeq c m^{-p/d}.\)
Taking the $p$-th root and minimizing over $Y$ proves the lower bound.
\end{proof}
For fixed codepoints $Y$, the powered transport cost
\(
\b\mapsto\Wass_p^p(\al,\sum_j\b_j\de_{y_j})
\)
\paragraph{Lloyd algorithm.}
\begin{prop}[Free masses give Voronoi cells]\label{prop-free-masses-voronoi}
For the cost $c(x,y)=d(x,y)^p$, fix distinct codepoints $Y=(y_j)_{j=1}^m$. Duplicate codepoints can be merged beforehand. Minimizing over the weights $\b\in\simplex_m$ gives
\(\min_{\b\in\simplex_m} \Wass_p^p(\al,\sum_j \b_j\de_{y_j}) = \int_\X \min_{1\leq j\leq m} c(x,y_j)\d\al(x).\)
An optimal coupling is induced by sending each $x$ to a nearest codepoint. The corresponding cells are the Voronoi cells
\(\VV_j(Y) \eqdef \enscond{x}{\foralls j',\ c(x,y_j)\leq c(x,y_{j'})},\)
up to arbitrary tie-breaking on common boundaries.
\end{prop}
\begin{proof}
For any coupling between $\al$ and a measure supported on $Y$, the conditional destination of a point $x$ belongs to $Y$, hence its conditional cost is at least $\min_j c(x,y_j)$. Integrating gives the lower bound. Conversely, choose a measurable nearest-codepoint map $T_Y(x)\in\argmin_j c(x,y_j)$, breaking ties measurably, and set $\b_j=\al(T_Y^{-1}(y_j))$. Then $(T_Y)_\sharp\al=\sum_j \b_j\de_{y_j}$ and the induced transport reaches the displayed lower bound.
\end{proof}
Consequently, the quantization energy can be written in the nearest-centroid form
\(\Qq_m(\al)^p = \min_Y \Ff(Y), \qquad \Ff(Y) \eqdef \int_\X \min_{1\leq j\leq m} c(x,y_j)\d\al(x).\)
At a differentiability point of this energy, any local minimizer with nonempty cells satisfies the centroid condition
\(y_j \in \uargmin{y} \int_{\VV_j(Y)} c(x,y)\d\al(x).\)
\paragraph{Continuous Lloyd flow.}
For a configuration \(Y\), define, on nonempty cells, \(a_j(Y)=\al(\VV_j(Y)), \qquad b_j(Y)=\frac{1}{a_j(Y)}\int_{\VV_j(Y)}x\,\d\al(x)\)
The relaxed Lloyd step
\(y_j^{(\ell+1)}=y_j^{(\ell)}+\tau\big(b_j(Y^{(\ell)})-y_j^{(\ell)}\big), \qquad 0<\tau\leq 1,\)
is the explicit-Euler discretization of the cell-mass preconditioned gradient flow of the quantization energy \(\Ff\),
\eql{\label{eq-lloyd-preconditioned-flow}
\dot y_j(t)=b_j(Y_t)-y_j(t).
}
Indeed, at differentiability points of \(\Ff\), the envelope theorem gives \(\nabla_{y_j}\Ff(Y) = 2\int_{\VV_j(Y)}(y_j-x)\d\al(x) = 2a_j(Y)(y_j-b_j(Y)),\)
so that
\(\dot y_j(t) = -\frac{1}{2a_j(Y_t)}\nabla_{y_j}\Ff(Y_t).\)
Along smooth portions of the flow, \(\frac{\d}{\d t}\Ff(Y_t) = -2\sum_j a_j(Y_t)\norm{b_j(Y_t)-y_j(t)}^2 \leq 0.\)
If \(\eta_t=\sum_j w_j\de_{y_j(t)}\) carries fixed positive weights, independent of the Voronoi masses \(a_j(Y_t)\), Proposition~\ref{prop-lagrangian-flow-continuity} turns this labelled particle ODE into the weak continuity equation
\(\partial_t\eta_t+\diverg(v_t\eta_t)=0, \qquad v_t(y_j(t))=b_j(Y_t)-y_j(t),\)
in the sense of the measure evolutions of Section~\ref{sec-dynamic-optimal-transport}. The weights \(w_j\) in this transport equation are auxiliary weights for the moving labelled particles; they are not the Voronoi masses used to define the quantization energy. If one instead records the free-weight projection
\(\nu_{Y_t}=\sum_j a_j(Y_t)\de_{y_j(t)},\)
then, formally, \(\partial_t\nu_{Y_t}+\diverg(v_t\nu_{Y_t}) = \sum_j \dot a_j(Y_t)\de_{y_j(t)}, \qquad v_t(y_j(t))=\dot y_j(t).\)
\paragraph{Mean-field limit and ultrafast diffusion.}
When the number \(m\) of codepoints is large, their empirical distribution
\(
m^{-1}\sum_{j=1}^m\de_{y_j}
\)
If \(\al=\rho\d x\), \(r=p/d\), the empirical site measures converge to \(\sigma\d x\), and the cells are asymptotically locally optimal, the high-resolution ansatz reads
\[
\begin{aligned}
m^{p/d}\int_\Omega \min_{1\leq j\leq m}\norm{x-y_j}^{p}\rho(x)\d x
&\longrightarrow C_{p,d}\,\Gg_\rho(\sigma),\\
\Gg_\rho(\sigma)
&\eqdef \int_\Omega \rho(x) \sigma(x)^{-r}\d x,
\qquad \int_\Omega \sigma\,\d x=1.
\end{aligned}
\]
Formally, the \(\Wass_2\)-gradient flow of \(\Gg_\rho\) is
\[
\partial_t \sigma
=
-r\,\diverg\!\left(
\sigma\nabla\!\left(\frac{\rho}{\sigma^{r+1}}\right)
\right),
\]
\paragraph{Quantization with fixed equal weights.}
In the equal-weight case, set
\(\nu_Y\eqdef \frac1m\sum_{j=1}^m\de_{y_j}, \qquad \Ff_{\rm eq}(Y) \eqdef \frac12\Wass_2^2(\al,\nu_Y),\)
Let \(C_j(Y)\) be the Laguerre cell transported to \(y_j\), so that \(\al(C_j(Y))=1/m\), and define its centroid
\(\bar x_j(Y) \eqdef m\int_{C_j(Y)} x\d\al(x).\)
At differentiability points, the envelope theorem gives \(\nabla_{y_j}\Ff_{\rm eq}(Y) = \frac1m\bigl(y_j-\bar x_j(Y)\bigr).\)
As long as the optimal matching between two nearby labelled configurations preserves labels, the \(\Wass_2\) metric on equal-weight empirical measures induces the local particle metric \(g_Y(U,V)=m^{-1}\sum_j\dotp{u_j}{v_j}\). Hence the associated \(\Wass_2\) gradient flow is the coupled system
\(\dot y_j(t) = \bar x_j(Y_t)-y_j(t), \qquad j=1,\ldots,m.\)
Equivalently, \(\nu_{Y_t}\) satisfies a continuity equation with velocity \(v_t(y_j(t))=\bar x_j(Y_t)-y_j(t)\). This is the so-called finite-particle \(\Wass_2\) gradient-flow viewpoint developed more systematically in Chapter~\ref{sec-wasserstein-gradient-flows}; the fixed-weight cells are Laguerre cells rather than the free-mass Voronoi cells used by Lloyd's method.

\paragraph{Equal-weight quantization on the line.}
\begin{prop}[One-dimensional equal-weight quantization]\label{prop-1d-equal-weight-quantization}
Let $\al\in\Mm_+^1(\RR)$ have finite second moment and quantile function $Q=\cumul{\al}^{-1}$. For
\(\Qq^{\mathrm{eq}}_m(\al)^2 \eqdef \min_{y_1\leq\cdots\leq y_m} \Wass_2^2(\al,\frac1m\sum_{i=1}^m\de_{y_i}),\)
set $I_i=((i-1)/m,i/m]$. Then the sorted minimizer is unique and its $i$th atom is \(y_i^\star = m\int_{I_i} Q(u)\d u,\)
and
\(\Qq^{\mathrm{eq}}_m(\al)^2 = \sum_{i=1}^m \int_{I_i} \abs{Q(u)-y_i^\star}^2\d u.\)
If $Q\in C^1([0,1])$, then
\(m^2\Qq^{\mathrm{eq}}_m(\al)^2 \longrightarrow \frac1{12}\int_0^1\abs{Q'(u)}^2\d u.\)
\end{prop}
\begin{proof}
After sorting the atoms, the quantile formula for $\Wass_2$ gives \(\Wass_2^2(\al,\frac1m\sum_{i=1}^m\de_{y_i}) = \sum_{i=1}^m \int_{I_i}\abs{Q(u)-y_i}^2\d u.\)
The minimization therefore decouples over the intervals $I_i$, and the best constant approximation of $Q$ on $I_i$ in $L^2$ is its average on $I_i$. Since $Q$ is nondecreasing, these interval averages are nondecreasing, so they satisfy the sorting constraint. Strict convexity in each scalar $y_i$ proves uniqueness and gives the displayed energy.

Denote by $\bar Q_{I_i}=m\int_{I_i}Q(u)\d u$ the interval average. If $Q$ is $C^1$, set $h=1/m$ and write $I_i=(a_i,a_i+h]$. Uniform Taylor expansion gives, for $v\in[0,1]$,
\(Q(a_i+hv)=Q(a_i)+hvQ'(a_i)+h\,r_i(v), \qquad \max_i\sup_{v\in[0,1]}\abs{r_i(v)}\to0.\)
Subtracting the average over $v\in[0,1]$ and integrating gives \(\int_{I_i}\abs{Q(u)-\bar Q_{I_i}}^2\d u = \frac{h^3}{12}\abs{Q'(a_i)}^2+o(h^3),\)
uniformly in $i$. Summing over $i$ yields a Riemann sum for $\frac1{12}\int_0^1\abs{Q'(u)}^2\d u$.
\end{proof}
Thus the common deterministic rule $m^{-1}\sum_i\de_{Q((i-1/2)/m)}$ should be read as a midpoint approximation of the optimal bin-average formula. It has the same leading squared error under the same smoothness assumptions. Indeed, orthogonal projection onto constants on each $I_i$ gives
\(\int_{I_i}|Q(u)-z|^2\d u = \int_{I_i}|Q(u)-y_i^\star|^2\d u + \frac1m|z-y_i^\star|^2,\)
\begin{prop}[Quantile-process asymptotics for random placement]\label{prop-1d-random-quantile-process}
Let $\al\in\Mm_+^1(\RR)$ have quantile $Q\in C^1([0,1])$, and let $\widehat\al_m=m^{-1}\sum_{i=1}^m\de_{X_i}$ with $X_i$ i.i.d.\ with law $\al$. Then
\[
m\,\EE\!\left[\Wass_2^2(\al,\widehat\al_m)\right]
\longrightarrow
\int_0^1 u(1-u)\abs{Q'(u)}^2\d u.
\]
\end{prop}
\begin{proof}
Let $\widehat Q_m$ be the empirical quantile function. The one-dimensional formula gives \(\Wass_2^2(\al,\widehat\al_m) = \int_0^1\abs{\widehat Q_m(u)-Q(u)}^2\d u.\)
Write $X_i=Q(U_i)$ with $U_i$ i.i.d.\ uniform on $[0,1]$, and let $\widehat U_m^{-1}$ be the empirical quantile function of the $U_i$. Then $\widehat Q_m=Q\circ \widehat U_m^{-1}$. The classical uniform quantile-process theorem and the functional delta method show that $\sqrt m\,(\widehat Q_m-Q)$ converges in law in $L^2(0,1)$ to $Q'$ times a standard Brownian bridge, up to an immaterial sign. Standard fourth-moment bounds for the uniform quantile process, together with the boundedness of $Q'$, give uniform integrability of the squared $L^2$ norms and hence convergence of expectations. A standard Brownian bridge has pointwise second moment $u(1-u)$, so Fubini's theorem gives the displayed limit.
\end{proof}
\subsection{Wasserstein-1 norm}
\label{sec-W1}
\paragraph{\texorpdfstring{$c$-transform for $\Wass_1$.}{c-transform for W1.}}
Assume that $d$ is a distance on $\X=\Y$ and take the ground cost $c(x,y)=d(x,y)$. We denote the Lipschitz constant of $f\in\Cc(\X)$ by

\begin{defn}[Lipschitz constant]\label{def-lipschitz-constant}
For a function $f:\X\to\RR$ on a metric space $(\X,d)$, its Lipschitz constant is
\eql{\label{eq-lip-constant}
\Lip(f)
\eqdef
\sup\enscond{\frac{|f(x)-f(y)|}{d(x,y)}}{x\neq y}.
}
The function is $1$-Lipschitz when $\Lip(f)\leq1$.
\end{defn}
\begin{prop}[$c$-transforms and $1$-Lipschitz functions]\label{prop-w1-c-transform-lipschitz}
Suppose $\X=\Y$ and $c(x,y)=d(x,y)$. Then there exists $g$ such that $f=g^c$ if and only if $\Lip(f)\leq1$. Furthermore, if $\Lip(f)\leq1$, then $f^c=-f$.
\end{prop}
\begin{proof}
First suppose $f=g^c$ for some $g$. For $x,y\in\X$,
\[
|f(x)-f(y)|
=
\left|\inf_z[d(x,z)-g(z)]-\inf_z[d(y,z)-g(z)]\right|
\leq
\sup_z |d(x,z)-d(y,z)|
\leq d(x,y),
\]
where the last inequality is the reverse triangle inequality. Thus $\Lip(f)\leq1$.

If $\Lip(f)\leq1$, then $f(x)\leq f(y)+d(x,y)$, so $d(x,y)-f(x)\geq -f(y)$ for all $x$ and hence $f^c(y)\geq -f(y)$. Taking $x=y$ gives $f^c(y)\leq -f(y)$. Therefore $f^c=-f$. Applying the same property to $-f$ gives $(-f)^c=f$, so every $1$-Lipschitz function is $c$-concave.
\end{proof}
The Kantorovich dual therefore becomes the Kantorovich--Rubinstein formula
\eql{\label{eq-w1-metric}
\Wass_1(\al,\be)
=
\umax{f}
\enscond{\int_\X f\d(\al-\be)}{\Lip(f)\leq1}
=:
\norm{\al-\be}_{W_1}.
}
For a discrete signed measure $\al-\be=\sum_{k=1}^N r_k\de_{z_k}$ on distinct sites, with $\sum_k r_k=0$,~\eqref{eq-w1-metric} becomes the finite-dimensional linear program
\eql{\label{eq-w1-discr}
\Wass_1(\al,\be)
=
\umax{(f_k)_k}
\enscond{\sum_k f_k r_k}{\foralls k,\ell,\ |f_k-f_\ell|\leq d(z_k,z_\ell)}.
}
This linear program can be solved by generic interior-point or first-order methods. If $N$ support points are involved, however, it still contains $O(N^2)$ Lipschitz constraints, mirroring the $O(nm)$ coupling variables of the original discrete Kantorovich LP; the dual formulation alone does not remove the all-pairs structure. When $d(x,y)=|x-y|$ on $\RR$, ordering the support points $z_1\leq z_2\leq\cdots$ reduces the constraints to neighboring pairs,
\(\Wass_1(\al,\be) = \umax{(f_k)_k} \enscond{\sum_k f_k r_k}{\foralls k,\ |f_{k+1}-f_k|\leq z_{k+1}-z_k}.\)
\paragraph{\texorpdfstring{$\Wass_1$ on Euclidean spaces.}{W1 on Euclidean spaces.}}
Let $\al,\be\in\Mm_+^1(\RR^d)$ have finite first moments and set $\xi=\al-\be$. Rademacher's theorem identifies $1$-Lipschitz functions with the continuous representatives of locally Sobolev functions whose Euclidean gradient has norm at most one almost everywhere. Normalizing at the origin fixes the additive gauge and ensures $|f(x)|\leq\norm{x}_2$, so integration against $\xi$ is well-defined. Hence
\eql{\label{eq-w1-cont}
\Wass_1(\al,\be) =
\usup{f\in C(\RR^d)\cap W_{\mathrm{loc}}^{1,\infty}(\RR^d)}
\enscond{ \int_{\RR^d} f\d\xi }
{ f(0)=0,\ \norm{\nabla f}_{L^\infty(\RR^d)} \leq 1 }.
}
\begin{defn}[Total variation norm of a vector-valued measure]\label{def-vector-measure-tv}
Let $\Mm(\RR^d;\RR^d)$ denote the finite $\RR^d$-valued Radon measures. For $\flow\in\Mm(\RR^d;\RR^d)$, define
\(\norm{\flow}_{\TV} \eqdef \sup_{\zeta\in C_c(\RR^d;\RR^d)} \enscond{\int_{\RR^d}\dotp{\zeta}{\d\flow}} {\sup_{x\in\RR^d}\norm{\zeta(x)}_2\leq1}.\)
\end{defn}
For $\flow\in\Mm(\RR^d;\RR^d)$, its divergence is defined distributionally by \(\langle\diverg\flow,\varphi\rangle \eqdef -\int_{\RR^d}\dotp{\nabla\varphi}{\d\flow}, \qquad \varphi\in C_c^1(\RR^d).\)
\begin{defn}[Beckmann problem]\label{def-beckmann-problem}
For $\al,\be\in\Mm_+^1(\RR^d)$, define
\(\operatorname{Beck}(\al,\be) \eqdef \inf_{\flow\in\Mm(\RR^d;\RR^d)} \enscond{\norm{\flow}_{\TV}}{\diverg\flow=\al-\be}.\)
\end{defn}
\begin{prop}[Euclidean Beckmann formula]\label{prop-euclidean-beckmann}
For $\al,\be\in\Mm_+^1(\RR^d)$ with finite first moments,
\eql{\label{eq-w1-cont-div}
\Wass_1(\al,\be)=\operatorname{Beck}(\al,\be).
}
\end{prop}
\begin{proof}
Let $\flow$ be feasible and let $f$ be admissible in~\eqref{eq-w1-cont}. Write $T_K(f)=\max\{-K,\min\{f,K\}\}$ and mollify it to a smooth function $f_{K,\delta}$ satisfying $\norm{f_{K,\delta}}_\infty\leq K$ and $\norm{\nabla f_{K,\delta}}_\infty\leq1$. Choose $\chi_R\in C_c^1(\RR^d)$ equal to one on $B_R$, with $0\leq\chi_R\leq1$ and $\norm{\nabla\chi_R}_\infty\leq C/R$. The distributional divergence identity applies to $\chi_R f_{K,\delta}$ and gives
\[
\left|\int \chi_R f_{K,\delta}\d(\al-\be)\right|
=
\left|\int\dotp{\chi_R\nabla f_{K,\delta}+f_{K,\delta}\nabla\chi_R}{\d\flow}\right|
\leq
\left(1+\frac{CK}{R}\right)\norm{\flow}_{\TV}.
\]
For fixed $K$ and $\delta$, dominated convergence lets $R\to+\infty$. Next $f_{K,\delta}\to T_K(f)$ uniformly as $\delta\to0$. Finally $T_K(f)\to f$ and $|T_K(f)|\leq|f|\leq\norm{x}_2$, so the finite first moments permit $K\to+\infty$. Equation~\eqref{eq-w1-cont} therefore gives $\Wass_1(\al,\be)\leq\norm{\flow}_{\TV}$.

Conversely, let $\pi$ be an optimal coupling and define a vector measure by its action on $\zeta\in C_c(\RR^d;\RR^d)$:
\(\int\dotp{\zeta(z)}{\d\flow(z)} \eqdef \int_{\RR^d\times\RR^d}\int_0^1 \dotp{\zeta((1-t)x+ty)}{y-x}\d t\d\pi(x,y).\)
For $\varphi\in C_c^1(\RR^d)$, the fundamental theorem of calculus gives \(\langle\diverg\flow,\varphi\rangle = -\int\!\int_0^1 \dotp{\nabla\varphi((1-t)x+ty)}{y-x}\d t\d\pi(x,y) = \int\varphi\d(\al-\be).\)
Moreover, $\norm{\flow}_{\TV}\leq\int\norm{x-y}_2\d\pi(x,y)=\Wass_1(\al,\be)$. Together with the first inequality, this proves equality and attainment.
\end{proof}
\paragraph{Graph distances and Beckmann flows.}
\begin{defn}[Graph geodesic distance]\label{def-graph-geodesic-distance}
Let $G=(V,E)$ be a connected finite graph with positive edge lengths $(\ell_e)_{e\in E}$. The graph geodesic distance between two vertices is \(d_G(i,j)=\min_{\gamma:i\leadsto j}\sum_{e\in\gamma}\ell_e.\)
The minimum is over all paths $\gamma$ joining $i$ to $j$.
\end{defn}
\begin{prop}[$\Wass_1$ and Beckmann flow on a graph]\label{prop-graph-w1-beckmann}
Let $G=(V,E)$ be a connected finite graph with positive edge lengths $(\ell_e)_{e\in E}$ and graph geodesic distance $d_G$.
For two probability vectors $\a,\b$ on $V$, set $r=\a-\b$ and orient each edge $e=(i,j)$. If
\((\nabla_G f)_e=f_j-f_i,\qquad \operatorname{div}_G=-\nabla_G^*\)
are the finite-difference gradient and its negative adjoint, then a positive flow on the oriented edge $i\to j$ has positive divergence at $i$ and negative divergence at $j$. With this convention,
\[
\Wass_{1,G}(\a,\b)
=
\max_f \enscond{\sum_{i\in V} f_i r_i}{|f_i-f_j|\leq \ell_e\quad\forall e=(i,j)}
=
\min_m \enscond{\sum_{e\in E}\ell_e |m_e|}{\operatorname{div}_G m=r}.
\]
The vector $m_e$ is an oriented edge flow, and the constraint $\operatorname{div}_G m=r$ is conservation of mass at each vertex.
\end{prop}
\begin{proof}
The edge constraints imply, by summing along every path and then minimizing over paths, that $|f_i-f_j|\leq d_G(i,j)$ for all vertices. Conversely, any $1$-Lipschitz function for $d_G$ satisfies $|f_i-f_j|\leq d_G(i,j)\leq\ell_e$ on every edge. The first equality is therefore the Kantorovich--Rubinstein formula on the metric space $(V,d_G)$.

For the second equality, write the graph Beckmann problem and dualize its equality constraint with a potential $f$: \(\inf_m \sum_e\ell_e|m_e| +\sup_f \sum_i f_i(r_i-(\operatorname{div}_G m)_i).\)
Using $\operatorname{div}_G=-\nabla_G^*$, the coupling term is $\sum_e m_e(\nabla_G f)_e$. The minimization over each scalar flow $m_e$ is finite exactly when $|(\nabla_G f)_e|\leq\ell_e$, and is then equal to zero. The dual problem is therefore precisely the graph Lipschitz dual above. Strong duality holds because this is a finite-dimensional linear program with a nonempty feasible set: connectedness and $\sum_i r_i=0$ allow the signed surplus to be routed along paths. This proves the graph Beckmann formula.
\end{proof}

\section{Divergences and Dual Norms}
\label{sec-divergences-dual-norms}
\subsection{Dual Norms (Integral Probability Metrics)}
\label{sec-dual-norms}
\paragraph{Integral probability metrics.}
\begin{defn}[Dual seminorm and integral probability metric]\label{def-dual-norm-ipm}
Let $B$ be a nonempty symmetric convex class of measurable functions that are integrable against the measures under consideration. For a signed measure $\xi$, define the extended dual seminorm
\eql{\label{eq-dual-norm-cont}
\norm{\xi}_B \eqdef
\usup{\f} \enscond{ \int_{\Xx} \f(x) \d\xi(x) }{ \f \in B}.
}
It is a norm when it is finite and $B$ separates signed measures. When applied to $\al-\be$ for probability measures, it is called an integral probability metric (IPM); it is a genuine metric precisely when $B$ separates probability measures.
\end{defn}
\begin{prop}[Metrization by dual norms]\label{prop-dual-norm-metrization}
Assume that $\Xx$ is compact, that $B=-B$, and that the measures considered are probability measures.
\begin{enumerate}
\item If every function in $\Cc(\Xx)$ can be uniformly approximated by elements of $\Span(B)$, then $\norm{\al_n-\al}_B \rightarrow 0$ implies $\al_n \rightharpoonup \al$.
\item If $B \subset \Cc(\Xx)$ is compact for $\norm{\cdot}_\infty$, then $\al_n \rightharpoonup \al$ implies $\norm{\al_n-\al}_B \rightarrow 0$.
\end{enumerate}
\end{prop}
\begin{proof}
For the first implication, $\norm{\al_n-\al}_B\to0$ and the symmetry of $B$ imply
\(\abs{\dotp{f}{\al_n-\al}}\leq \norm{\al_n-\al}_B \qforq f\in B.\)
If $h=\sum_{j=1}^J c_j f_j\in\Span(B)$, then
$\abs{\dotp{h}{\al_n-\al}}
\leq(\sum_{j=1}^J|c_j|)\norm{\al_n-\al}_B$,
so integrals converge for every $h\in\Span(B)$. Let $u\in\Cc(\Xx)$ and choose such an $h$ with $\norm{u-h}_\infty\leq\eta$. Since $\al_n$ and $\al$ are probabilities,
$\abs{\dotp{u}{\al_n-\al}}
\leq\abs{\dotp{h}{\al_n-\al}}+2\eta$.
Taking the limsup as $n\to\infty$ and then letting $\eta\to0$ gives $\dotp{u}{\al_n}\to\dotp{u}{\al}$ for all $u\in\Cc(\Xx)$, which is weak convergence.

For the second implication, assume $\al_n\rightharpoonup\al$ and choose a subsequence $(\al_{n_k})_k$ such that $\norm{\al_{n_k}-\al}_B\to\limsup_n\norm{\al_n-\al}_B$. Since $B$ is compact and the map $f\mapsto\dotp{f}{\al_{n_k}-\al}$ is continuous on $B$, the supremum is attained by some $f_{n_k}\in B$. Extracting a further subsequence if needed, $f_{n_k}\to f$ uniformly for some $f\in B$. Then
\(\dotp{f_{n_k}}{\al_{n_k}-\al} = \dotp{f}{\al_{n_k}-\al} +\dotp{f_{n_k}-f}{\al_{n_k}} -\dotp{f_{n_k}-f}{\al}.\)
The first term tends to zero by weak convergence and the last two by uniform convergence. Hence every limsup subsequence has limit zero, proving $\norm{\al_n-\al}_B\to0$.
\end{proof}
\begin{cor}[Wasserstein metrizes weak convergence]\label{cor-topol-wass}
On a compact metric space, $\Wass_p$ metrizes weak convergence on probability measures for every $1\leq p<+\infty$.
\end{cor}
\begin{proof}
For $p=1$, take $B=\enscond{f}{\Lip(f)\leq1}$. The span of $B$ contains all Lipschitz functions, and Lipschitz functions are dense in $\Cc(\Xx)$ on compact metric spaces.

Conversely, constants do not change the pairing with $\al_n-\al$. Fix $x_0\in\Xx$ and normalize potentials by $f(x_0)=0$. The normalized unit Lipschitz ball is uniformly bounded by $\mathrm{diam}(\Xx)$ and equicontinuous, hence compact in $\norm{\cdot}_\infty$ by Arzel\`a--Ascoli. Proposition~\ref{prop-dual-norm-metrization} gives $\Wass_1(\al_n,\al)\to0$. Proposition~\ref{prop-comp-wass-p} shows that all finite-order $\Wass_p$ distances induce the same topology on a compact space, so the result follows for every $1\leq p<+\infty$.
\end{proof}
\subsection{Dual RKHS Norms and Maximum Mean Discrepancies}
\label{sec-rkhs-mmd}
\begin{defn}[Positive and conditionally definite kernels]\label{def-positive-kernels}
A symmetric Borel-measurable function $\Krkhs:\Xx\times\Xx\to\RR$ is positive definite if for every $n\geq1$, every $x_1,\ldots,x_n\in\Xx$ and every $r\in\RR^n$,
\eql{\label{eq-dual-kern}
\sum_{i,j=1}^n r_i r_j \Krkhs(x_i,x_j)\geq0.
}
It is conditionally positive definite of order one if the same inequality is required only for zero-sum vectors, $\dotp{r}{\ones_n}=0$. It is conditionally negative definite of order one if the reverse inequality holds for every such vector, or equivalently if $-\Krkhs$ is conditionally positive definite.
\end{defn}
The conditional notions are the right ones for probability distances, because one applies quadratic forms to signed measures $\xi=\al-\be$ of total mass zero. Adding a term of the form $a(x)+a(y)$ to a kernel does not change $\iint \Krkhs(x,y)\d\xi(x)\d\xi(y)$ on such measures.

For a conditionally positive-definite kernel, fix $x_0\in\Xx$ and set
\(\Krkhs^\circ(x,y) \eqdef \Krkhs(x,y)-\Krkhs(x,x_0)-\Krkhs(x_0,y)+\Krkhs(x_0,x_0).\)
\begin{defn}[Kernel seminorm and MMD]\label{def-kernel-mmd-norm}
Assume that $\Phi^\circ$ is strongly measurable. A finite signed measure $\xi$ is $\Krkhs$-admissible if
$\int\sqrt{\Krkhs^\circ(x,x)}\d|\xi|(x)<+\infty$ and, in the conditionally positive-definite case, $\xi(\Xx)=0$. Its kernel seminorm is
\eql{\label{eq-kernel-dual}
\norm{\xi}^2_{\Krkhs}
\eqdef
\iint_{\Xx\times\Xx}\Krkhs^\circ(x,y)\d\xi(x)\d\xi(y).
}
For two probability measures such that $\al-\be$ is $\Krkhs$-admissible, the maximum mean discrepancy associated with $\Krkhs$ is \(\MMD_{\Krkhs}(\al,\be)\eqdef\norm{\al-\be}_{\Krkhs}.\)
It is a genuine distance when the kernel is characteristic, meaning that $\norm{\al-\be}_{\Krkhs}=0$ implies $\al=\be$.
\end{defn}
\begin{prop}[Kernel seminorm as an RKHS dual norm]\label{prop-kernel-rkhs-dual}
Let $\xi$ be $\Krkhs$-admissible in the sense of Definition~\ref{def-kernel-mmd-norm}. Then
$m_\xi\eqdef\int\Phi^\circ(x)\d\xi(x)$ is a well-defined Bochner integral and
$\norm{\xi}_{\Krkhs}
=\sup_{\norm{h}_{\RKHS^\circ}\leq1}\int h(x)\d\xi(x)$,
so $\norm{\cdot}_{\Krkhs}$ is the dual seminorm in the sense of~\eqref{eq-dual-norm-cont} associated with the RKHS unit ball.
\end{prop}
\begin{proof}
The integrability condition in Definition~\ref{def-kernel-mmd-norm} is precisely $\int\norm{\Phi^\circ(x)}_{\RKHS^\circ}\d|\xi|(x)<+\infty$, so $m_\xi$ exists. By the reproducing property,
$\int h(x)\d\xi(x)
=\langle h,m_\xi\rangle_{\RKHS^\circ}$.
Cauchy--Schwarz gives
$\sup_{\norm{h}_{\RKHS^\circ}\leq1}\int h\d\xi=\norm{m_\xi}_{\RKHS^\circ}$.
Finally,
$\norm{m_\xi}_{\RKHS^\circ}^2
=\iint\Krkhs^\circ(x,y)\d\xi(x)\d\xi(y)$,
which is exactly~\eqref{eq-kernel-dual}.
\end{proof}
\begin{prop}[Universal kernels metrize weak convergence]\label{prop-mmd-metrization}
Assume that $\Xx$ is compact and that the RKHS generated by the continuous kernel $\Krkhs$ is dense in $\Cc(\Xx)$ for the uniform norm. Then
\[
\MMD_{\Krkhs}(\al_n,\al)\to0
\quad\Longleftrightarrow\quad
\al_n\rightharpoonup\al
\]
for probability measures on $\Xx$.
\end{prop}
\begin{proof}
If $\MMD_{\Krkhs}(\al_n,\al)\to0$, then for every $g\in\RKHS$, \(\abs{\int g\d(\al_n-\al)} \leq \norm{g}_{\RKHS}\MMD_{\Krkhs}(\al_n,\al)\to0.\)
For any $h\in\Cc(\Xx)$ and any $\eta>0$, choose $g\in\RKHS$ with $\norm{h-g}_\infty\leq\eta$. Since $\al_n$ and $\al$ are probabilities,
\[
\left|\int h\,\d(\al_n-\al)\right|
\leq
2\eta+\left|\int g\,\d(\al_n-\al)\right|,
\]
and the last term tends to zero. This proves weak convergence. Conversely, if $\al_n\rightharpoonup\al$, then $\al_n\otimes\al_n$, $\al_n\otimes\al$ and $\al\otimes\al$ converge weakly on the compact product space. Applying this to the continuous bounded function $\Krkhs$ in the identity
\(\MMD_{\Krkhs}(\al_n,\al)^2 = \iint \Krkhs\,\d\al_n\d\al_n -2\iint \Krkhs\,\d\al_n\d\al +\iint \Krkhs\,\d\al\d\al\)
gives convergence to zero.
\end{proof}
\begin{defn}[Universal kernel]\label{def-universal-kernel}
A continuous positive-definite kernel $\Krkhs$ on a compact space $\Xx$ is universal if finite sums of sections $\sum_{i=1}^n a_i\Krkhs(x_i,\cdot)$ are dense in $\Cc(\Xx)$ for the uniform norm.
\end{defn}
For a positive-definite kernel and a discrete measure $\al$, one has the simple expression
\eq{
\norm{\al}_{\Krkhs}^2 = \sum_{i=1}^n \sum_{i'=1}^n \a_i\a_{i'} \KrkhsD_{i,i'} = \dotp{\KrkhsD\a}{\a}
\qwhereq
\KrkhsD_{i,i'} \eqdef \Krkhs(x_i,x_{i'}).
}
To compute the discrepancy between two discrete measures, one can use
\eql{\label{eq-mmd-discr}
\norm{\al-\be}_{\Krkhs}^2 =
\sum_{i,i'} \a_i \a_{i'} \Krkhs(x_i,x_{i'})	+
\sum_{j,j'} \b_j \b_{j'} \Krkhs(y_j,y_{j'}) - 2
\sum_{i,j} \a_i \b_j \Krkhs(x_i,y_j).
}
\subsection{\texorpdfstring{$\phi$-divergences}{phi-divergences}}
\label{sec-phi-div}
\paragraph{Definition by density ratios.}
\begin{defn}[Entropy function]
\label{def_entropy}
A function $\phi: \RR \to \RR \cup \{+\infty\}$ is an entropy function if it is proper, lower semicontinuous and convex, with $\dom \phi\subset [0,+\infty)$ and $\dom \phi \cap (0,+\infty) \neq \emptyset$. Its recession slope is
\eq{
\phi'_\infty = \lim_{s\rightarrow +\infty} \frac{\phi(s)}{s} \in \RR \cup \{+\infty\}.
}
\end{defn}
If $\phi'_\infty = \infty$, then $\phi$ grows faster than any linear function and $\phi$ is said to be \emph{superlinear}.

\begin{defn}[$\phi$-Divergences]
\label{def_divergence}
Let $\phi$ be an entropy function and let $\al,\be\in\Mm_+(\Xx)$. Write the Lebesgue decomposition of $\al$ with respect to $\be$ as
$\al=(\d\al/\d\be)\be+\al^\perp$, where $\al^\perp\perp\be$. The associated divergence is
\eql{\label{eq-phi-div}
\Divergm_\phi (\al|\be) \eqdef \int_\Xx \phi\left(\frac{\d \al}{\d \be} \right) \d \be
+ \phi'_\infty \al^{\perp}(\Xx),
}
with the convention $0\cdot(+\infty)=0$. It is extended by $+\infty$ whenever either argument is not a nonnegative measure.
\end{defn}
In the discrete setting, assuming
\eql{\label{eq-div-disc-meas}
\al=\sum_i \a_i \de_{x_i}
\qandq \be=\sum_i \b_i \de_{x_i}
}
are supported on the same set of $n$ points $(x_i)_{i=1}^n \subset \X$,~\eqref{eq-phi-div} defines a divergence on $\simplex_n$
\eql{\label{eq-discr-diverg}
\DivergmD_\phi(\a|\b) = \sum_{i \in \Supp(\b)} \phi\pa{ \frac{\a_i}{\b_i} } \b_i + \phi'_\infty \sum_{i \notin \Supp(\b)} \a_i,
}
where $\Supp(\b) \eqdef \enscond{i \in \range{n}}{ \b_i \neq 0 }$.

\begin{proposition}[Basic properties of $\phi$-divergences]\label{prop-basic-phi-divergence-properties}
If $\phi$ is an entropy function, then $\Divergm_\phi$ is jointly $1$-homogeneous, convex and weak-* lower semicontinuous in $(\al,\be)$.
\end{proposition}
\begin{proof}
Introduce the lower-semicontinuous perspective
\eq{
\foralls (u,v) \in (\RR_+)^2, \quad
\psi(u,v) = \choice{
v\phi(u/v) \qifq v>0 \\
u \phi'_\infty \qifq v=0.
}
}
The joint $1$-homogeneity follows directly from this formula. To see convexity, first take $v_0,v_1>0$ and $\tau\in[0,1]$. With $\bar v=(1-\tau)v_0+\tau v_1$,
$\theta_0=(1-\tau)v_0/\bar v$, and $\theta_1=\tau v_1/\bar v$, one has
$\theta_0+\theta_1=1$ and
$((1-\tau)u_0+\tau u_1)/\bar v
=\theta_0u_0/v_0+\theta_1u_1/v_1$.
Jensen's inequality for $\phi$, followed by multiplication by $\bar v$, proves convexity of $\psi$ when both second arguments are positive. Its recession value at $v=0$ gives the lower-semicontinuous convex extension to the closed quadrant. In the discrete case,
\(\DivergmD_\phi(\a|\b) = \sum_i \psi(\a_i,\b_i)\),
so homogeneity and convexity follow by summation. For general finite measures, representing both arguments with respect to a common dominating measure yields the integral of $\psi$; the lower-semicontinuity theorem for convex perspective functionals on measures gives weak-* lower semicontinuity, with the recession value accounting for singular mass.
\end{proof}
\begin{proposition}[Non-negativity of $\phi$-divergences]\label{phi-div-positive}
Assume that $\phi$ is normalized by $\phi(1)=0$. For probability distributions $(\al,\be) \in \Mm_+^1(\Xx)$, one has $\Divergm_\phi(\al|\be) \geq 0$. If $\phi$ is strictly convex, then one has $\Divergm_\phi(\al|\be)=0$ if and only if $\al=\be$.
\end{proposition}
\begin{proof}
Let $m=\al+\be$ and write $a=\d\al/\d m$ and $b=\d\be/\d m$. Then $a+b=1$ $m$-almost everywhere and, using the perspective $\psi$ above,
\(\Divergm_\phi(\al|\be)=\int \psi(a,b)\d m\).
Since $m(\Xx)=2$, the measure $m/2$ is a probability. Jensen's inequality, the $1$-homogeneity of $\psi$, and $\phi(1)=0$ give
\(\frac12\Divergm_\phi(\al|\be) \geq \psi\pa{\frac12\int a\d m,\frac12\int b\d m} =\psi(1/2,1/2)=0.\)
If $\phi$ is strictly convex, the map $s\mapsto\psi(s,1-s)$ is strictly convex on its effective domain: distinct points on the line $u+v=1$ cannot lie on the same ray, which is the only degeneracy of a strict perspective. Equality in Jensen therefore forces $a=1/2$, hence $a=b$ and $\al=\be$. For arbitrary finite nonnegative measures, $\phi\geq0$ also implies $\phi'_\infty\geq0$, so every term in~\eqref{eq-phi-div} is nonnegative.
\end{proof}
\paragraph{Classical examples and topology.}
\begin{thm}[Pinsker inequality]\label{thm-pinsker}
Let $\al,\be\in\Mm_+^1(\Xx)$ be probability measures on a measurable space. With the full-variation convention, \(\TV(\al|\be)^2 = \norm{\al-\be}_{\TV}^2 \leq 2\KL(\al|\be).\)
In particular, on a finite space, for $\a,\b\in\simplex_n$, \(\norm{\a-\b}_{\ell^1}^2 \leq 2\KLD(\a|\b).\)
\end{thm}
\begin{proof}
The claim is immediate if $\KL(\al|\be)=+\infty$, so assume $\al\ll\be$. Let $A$ be a positive Hahn set for the signed measure $\al-\be$, and set $a=\al(A)$ and $b=\be(A)$. Since $(\al-\be)(\Xx)=0$,
\(a-b = \frac12\norm{\al-\be}_{\TV}.\)
Applying Jensen's inequality separately on $A$ and $A^c$, or equivalently coarse-graining onto the partition $(A,A^c)$, gives \(\KL(\al|\be) \geq a\log\frac{a}{b} +(1-a)\log\frac{1-a}{1-b}.\)
For fixed $b\in(0,1)$, subtract $2(a-b)^2$ from the binary relative entropy on the right. The resulting function of $a$ has value and first derivative zero at $a=b$, while its second derivative is \(\frac1a+\frac1{1-a}-4\geq0.\)
Hence the binary relative entropy is at least $2(a-b)^2=\frac12\norm{\al-\be}_{\TV}^2$. The boundary cases follow by lower semicontinuity. For a finite space, the identity $\norm{\al-\be}_{\TV}=\norm{\a-\b}_{\ell^1}$ gives the second formula.
\end{proof}
\paragraph{\texorpdfstring{Main families of $\phi$-divergences.}{Main families of phi-divergences.}}
For $s>0$, the power-divergence family
\(\phi_\gamma(s)=\frac{s^\gamma-\gamma s+\gamma-1}{\gamma(\gamma-1)} \qquad(\gamma\neq0,1)\)
has the lower-semicontinuous boundary value and recession slope
\[
\phi_\gamma(0)=
\begin{cases}1/\gamma,&\gamma>0,\\+\infty,&\gamma<0,\end{cases}
\qquad
\phi'_{\gamma,\infty}=
\begin{cases}+\infty,&\gamma>1,\\1/(1-\gamma),&\gamma<1.\end{cases}
\]
If $\al=\rho_\al\lambda$ and $\be=\rho_\be\lambda$ relative to any common dominating measure $\lambda$, then
\(\Hellinger(\al,\be) \eqdef\Divergm_{\phi_H}(\al|\be)^{1/2} =\norm{\sqrt{\rho_\al}-\sqrt{\rho_\be}}_{L^2(\lambda)}.\)
For probability measures, the Jensen--Shannon distance is the square root of the symmetrized KL-to-the-mixture divergence
\[
\JS(\alpha,\beta)^2
=
\frac12\KL\!\left(\alpha\middle|\frac{\alpha+\beta}{2}\right)
+
\frac12\KL\!\left(\beta\middle|\frac{\alpha+\beta}{2}\right),
\]
and satisfies $0\leq\JS(\al,\be)^2\leq\log 2$. The same formula extends by $1$-homogeneity to finite positive measures; on pairs of equal mass $M$, its upper bound is $M\log 2$. Its exact entropy generator is
\[
\phi_{\JS}(s)
=\frac12\left[s\log s-(s+1)\log\pa{\frac{s+1}{2}}\right],
\qquad
\Divergm_{\phi_{\JS}}(\al|\be)=\JS(\al,\be)^2.
\]
\paragraph{Variational dual formula.}
\begin{proposition}[Variational representation of a $\phi$-divergence]\label{prop-phi-div-dual}
Assume that $\Xx$ is compact and let $\al,\be\in\Mm_+(\Xx)$. Define the Legendre transform of the entropy, already extended by $+\infty$ on negative arguments, by \(\phi^*(s) \eqdef \sup_{r\geq0}\{sr-\phi(r)\}.\)
\eqllead{Then}{eq-dual-div}{
\Divergm_\phi(\al|\be)
=
\sup_{\substack{f\in\Cc(\Xx)\\ f\leq\phi'_\infty}}
\left\{\int_\Xx f\d\al-\int_\Xx\phi^*(f)\d\be\right\}.
}
When $\phi'_\infty=+\infty$, the pointwise constraint is vacuous.
Equivalently, the convex conjugate of $\Divergm_\phi(\cdot|\be)$ on $\Mm(\Xx)$ is
\eql{\label{eq-legendre}
\Divergm_\phi^*(f|\be)
=
\begin{cases}
\displaystyle\int_\Xx \phi^*(f(x)) \d\be(x),
& f\leq\phi'_\infty\ \text{on }\Xx,\\
+\infty,&\text{otherwise},
\end{cases}
\qquad f\in\Cc(\Xx).
}
\end{proposition}
\begin{proof}
Set $L=\phi'_\infty$. If $L=+\infty$, finite divergence forces $\al=\rho\be$, and pointwise convex duality gives the stated conjugate; the pointwise cap is then vacuous. Assume henceforth that $L<+\infty$ and write $\al=\rho\be+\eta$, where $\rho\geq0$, $\eta\geq0$ and $\eta\perp\be$. For $f\in\Cc(\Xx)$,
\(\int f\d\al-\Divergm_\phi(\al|\be) = \int\bigl(f\rho-\phi(\rho)\bigr)\d\be + \int(f-L)\d\eta.\)
If $f\leq L$ on $\Xx$, the supremum of the last term over nonnegative singular measures is zero, attained at $\eta=0$. Pointwise convex duality for the first term then gives $\int\phi^*(f)\d\be$. If $f(x_0)>L$ at some point, continuity gives a neighborhood $U$ on which $f>L$. When $\be(U)>0$, one has $\phi^*(f)=+\infty$ on $U$ because $L$ is the recession slope. When $\be(U)=0$, taking arbitrarily large singular mass supported in $U$ makes the last term unbounded. This proves~\eqref{eq-legendre}. Proposition~\ref{prop-basic-phi-divergence-properties} gives convexity and weak-* lower semicontinuity, so Fenchel--Moreau yields~\eqref{eq-dual-div}.
\end{proof}
\subsection{GANs via Duality}
\label{sec-gan-duality}
\paragraph{Divergence-based adversarial losses.}
Writing $L=\phi'_\infty$ gives
\eq{
\inf_{\th} \Divergm_\phi(\al_\th|\be)
= \inf_{\th} \sup_{\substack{f\in\Cc(\Xx)\\f\leq L}}
\left\{\int_\Zz f(g_\th(z)) \d\zeta(z) -
\frac{1}{m}\sum_{j=1}^m \phi^*(f(y_j))\right\}.
}
When $L=+\infty$, the pointwise cap is vacuous. Replacing $f$ by a continuous neural critic $f_\xi$ whose output parametrization enforces $f_\xi\leq L$ on all of $\Xx$ gives the restricted adversarial problem
\(\inf_\theta\sup_{\xi:\,f_\xi\leq L} \int_\Zz f_\xi(g_\theta(z))\d\zeta(z) - \frac1m\sum_{j=1}^m \phi^*(f_\xi(y_j)).\)
The original vanilla GAN corresponds, up to an additive constant and the usual reparametrization by a discriminator taking values in $(0,1)$, to the unscaled Jensen--Shannon generator $\widehat\phi_{\JS}=2\phi_{\JS}$,
\(\widehat\phi_{\JS}(s)=s\log s-(s+1)\log\frac{s+1}{2}, \qquad \widehat\phi_{\JS}^*(u)=-\log(2-e^u),\quad u<\log 2.\)
Thus $\Divergm_{\widehat\phi_{\JS}}=2\JS^2$, consistently with the normalization above. In practice the min--max problem is solved by alternating stochastic gradient descent/ascent on $(\theta,\xi)$.

\paragraph{Dual norms and integral probability metrics.}
Instead of a density-ratio divergence, one can minimize a dual norm~\eqref{eq-dual-norm-cont}, also called an integral probability metric,
\eq{
\inf_{\th} \norm{\al_\th - \be}_B
= \inf_{\th}
\usup{\f \in B} \int_{\X} \f(x) \d(\al_\th-\be)(x)
= \inf_{\th}
\usup{\f \in B} \int_{\Zz} \f( g_\th(z)) \d\zeta - \frac{1}{m} \sum_{j=1}^m f(y_j).
}

\section{Entropic Regularization: Sinkhorn Algorithm}
\label{sec-sinkhorn}
\subsection{Entropic Regularization for Discrete Measures}
\label{sec-entropic-discrete}
\paragraph{Entropy penalty.}
\begin{defn}[Discrete Shannon--Boltzmann entropy]\label{def-discrete-shannon-boltzmann-entropy}
For a nonnegative matrix $\P$, its Shannon--Boltzmann entropy is \(\HD(\P) \eqdef -\sum_{i,j} \P_{i,j} \log(\P_{i,j}),\)
with the convention $0\log(0)=0$.
\end{defn}
\begin{defn}[Discrete entropic optimal transport]\label{def-discrete-entropic-ot}
For $\epsilon>0$, define
\eql{\label{eq-regularized-discr}
\MKD_\C^\epsilon(\a,\b) \eqdef
\umin{\P \in \CouplingsD(\a,\b)}
\dotp{\P}{\C} - \epsilon \HD(\P).
}
\end{defn}
\begin{prop}[Existence and uniqueness of entropic OT]\label{prop-entropic-unique}
Assume that $\a,\b$ are probability histograms and that $\C$ is finite. For every $\epsilon>0$, problem~\eqref{eq-regularized-discr} admits a unique minimizer. If all entries of $\a$ and $\b$ are positive, then this minimizer is positive on every entry.
\end{prop}
\begin{proof}
The transport polytope $\CouplingsD(\a,\b)$ is non-empty and compact, and the objective is continuous on it with the convention $0\log0=0$, so a minimizer exists. The scalar function $r\mapsto r\log r$ is strictly convex on $[0,+\infty)$. On the relative interior, the Hessian
\(-\partial^2 \HD(\P)=\diag(1/\P_{i,j})\)
is positive definite. Hence $-\HD$ is strictly convex on every nontrivial feasible segment, and $\dotp{\P}{\C}-\epsilon\HD(\P)$ has at most one minimizer.

If $\a_i,\b_j>0$ and a minimizer had $\P_{i,j}=0$, then for small $t>0$ the perturbation $\P_t=(1-t)\P+t\,\a\otimes\b$ remains feasible. The directional derivative of the entropic part at $t=0$ is $-\infty$ because the derivative of $r\log r$ at $0$ is $-\infty$ along a positive direction. Thus the objective decreases for small $t$, contradicting optimality. Therefore the minimizer is strictly positive.
\end{proof}
\subsection{Sinkhorn's Algorithm}
\begin{prop}[Scaling form of entropic OT]\label{prop-regularized-primal}
$\P$ is the unique solution to~\eqref{eq-regularized-discr} if and only if there exists $(\uD,\vD) \in \RR_+^n \times \RR_+^m$ such that
\eql{\label{eq-scaling-form}
\foralls (i,j) \in \range{n} \times \range{m}, \quad \P_{i,j} = \uD_i \K_{i,j} \vD_j
\qwhereq \K_{i,j} \eqdef e^{-\frac{\C_{i,j}}{\epsilon}},
}
and $\P \in \CouplingsD(\a,\b)$.
\end{prop}
\begin{proof}
Without loss of generality, we assume $\a_i,\b_j>0$ (otherwise, rows or columns with zero mass are fixed to zero and can be removed). By Proposition~\ref{prop-entropic-unique}, the minimizer is strictly positive on the remaining support.

We can thus ignore the positivity constraint when introducing two dual variables $\fD\in\RR^n,\gD\in\RR^m$ for each marginal constraint so that the Lagrangian of~\eqref{eq-regularized-discr} reads
\eq{\label{eq-sinkhorn-lagrangian}
\Lag(\P,\fD,\gD)=
\dotp{\P}{\C}
+\epsilon\sum_{i,j}\P_{i,j}\log(\P_{i,j})
+ \dotp{\fD}{\a - \P\ones_m}
+ \dotp{\gD}{\b - \transp{\P}\ones_n}.
}
The first-order conditions are

\(\frac{\partial\Lag(\P,\fD,\gD)}{\partial \P_{i,j}}= \C_{i,j} + \epsilon (\log\pa{ \P_{i,j} }+1) - \fD_i -\gD_j = 0.\)
Thus
$\P_{i,j}=e^{(\fD_i+\gD_j-\C_{i,j})/\epsilon-1}$,
which has the stated form with the positive scalings
$\uD_i\eqdef e^{\fD_i/\epsilon-1}$ and $\vD_j\eqdef e^{\gD_j/\epsilon}$.
\end{proof}
$\uD,\vD$ must therefore satisfy the nonlinear mass-conservation equations
\eql{\label{eq-dualsinkhorn-constraints}
\diag(\uD)\K\diag(\vD)\ones_m=\a,
\qandq
\diag(\vD)\K^\top \diag(\uD)\ones_n=\b,
}
Using entrywise multiplication, these equations reduce to
\eql{\label{eq-dualsinkhorn-constraints2}
\uD \odot (\K \vD) = \a
\qandq
\vD \odot (\transp{\K}\uD) = \b
}
where $\odot$ denotes entrywise multiplication.

When $n=m$ and $\a,\b$ are uniform, this is diagonal scaling toward bistochasticity. Proposition~\ref{prop-entropic-unique} guarantees a unique scaled matrix $\P$, although the two scaling vectors retain one multiplicative gauge degree of freedom. If some entries of $\K$ vanish, equivalently if $\C$ can take the value $+\infty$, existence and uniqueness require additional support conditions; here we keep $\K>0$.

\eql{\label{eq-sinkhorn}
\itt{\uD} \eqdef \frac{\a}{\K \it{\vD}}
\qandq
\itt{\vD} \eqdef \frac{\b}{\transp{\K}\itt{\uD}},
}
\subsection{Reformulation Using Relative Entropy}
\paragraph{Relative entropy.}
\begin{defn}[Discrete relative entropy]\label{def-discrete-relative-entropy}
For nonnegative matrices $\P,\Q$ of the same size, the generalized relative entropy, or Kullback--Leibler divergence, is
\eql{\label{eq-kl-defn}
\KLD(\P|\Q) \eqdef \sum_{i,j} \P_{i,j} \log\pa{\frac{\P_{i,j}}{\Q_{i,j}}} - \P_{i,j} + \Q_{i,j}.
}
The convention is $0\log(0)=0$, and $\KLD(\P|\Q)=+\infty$ if there exists $(i,j)$ such that $\Q_{i,j}=0$ but $\P_{i,j} \neq 0$.
\end{defn}
For the specific case of comparing probability distributions, where $\P$ and $\Q$ have the same total mass, this further simplifies to
\(\KLD(\P|\Q) = \sum_{i,j} \P_{i,j} \log\pa{\frac{\P_{i,j}}{\Q_{i,j}}}.\)
For the reference matrix $\Q=\ones_{n \times m}$, one has \(-\KLD(\P|\ones_{n \times m}) = \HD(\P)+\sum_{i,j}\P_{i,j}-n\,m.\)
On fixed-mass couplings the last two terms are constant, so KL regularization with reference $\ones_{n\times m}$ is equivalent to subtracting Shannon--Boltzmann entropy. $\KLD$ is a particular instance of both a $\phi$-divergence (as defined in Section~\ref{sec-phi-div}) and a Bregman divergence; up to standard affine rescalings, it is the canonical overlap between these two families.

\begin{prop}[Non-negativity and definiteness of relative entropy]\label{prop-kl-distance-like}
For all $\P,\Q\in\RR_+^{n\times m}$, one has $\KLD(\P|\Q)\geq0$, with equality if and only if $\P=\Q$.
\end{prop}
\begin{proof}
Write $\phi(s)=s\log s-s+1$. Convexity gives $\phi(s)\geq \phi(1)+\phi'(1)(s-1)=0$, and strict convexity gives equality only at $s=1$. Hence
\(\KLD(\P|\Q)=\sum_{(i,j):\,\Q_{i,j}>0}\Q_{i,j} \phi(\P_{i,j}/\Q_{i,j})\geq0,\)
whenever $\P\ll\Q$ entrywise; otherwise the divergence is $+\infty$. Equality forces $\P_{i,j}=\Q_{i,j}$ wherever $\Q_{i,j}>0$, while entries with $\Q_{i,j}=0$ vanish in both matrices. Thus equality is equivalent to $\P=\Q$.
\end{proof}
When $\P$ and $\Q$ are probability matrices and $\P\ll\Q$, the same divergence can also be written
\(\KLD(\P|\Q) =\sum_{i,j}\psi(\P_{i,j}/\Q_{i,j})\Q_{i,j}, \qquad \psi(s)=s\log s.\)
More generally, for any convex $\psi$ such that $\psi(1)=0$, Jensen's inequality gives
\eq{
\sum_{i,j}\psi(\P_{i,j}/\Q_{i,j})\Q_{i,j}
\geq \psi\!\left(\sum_{i,j}\P_{i,j}\right)=\psi(1)=0.
}
\paragraph{KL reformulation of regularized OT.}
Choosing the tensor product $\a \otimes \b = (\a_i \b_j)_{i,j}$ as reference measure leads to the normalized problem
\eql{\label{eq-regularized-discr-rescaled}
\umin{\P \in \CouplingsD(\a,\b)}
\dotp{\P}{\C} + \epsilon \KLD(\P|\a \otimes \b).
}
For every $\P\in\CouplingsD(\a,\b)$, the two objectives satisfy
\(\dotp{\P}{\C}+\epsilon\KLD(\P|\a\otimes\b) = \dotp{\P}{\C}-\epsilon\HD(\P) +\epsilon\bigl(\HD(\a)+\HD(\b)\bigr).\)
Thus~\eqref{eq-regularized-discr-rescaled} has exactly the same minimizer as the original entropic problem~\eqref{eq-regularized-discr}, while its optimal value is
\(
\MKD_\C^\epsilon(\a,\b)+\epsilon\bigl(\HD(\a)+\HD(\b)\bigr)
\).
Indeed, for every $\P\in\CouplingsD(\a,\b)$,
\eql{\label{eq-sinkhorn-tilted-kl-identity}
\dotp{\P}{\C}+\epsilon\KLD(\P|\a\otimes\b)
=
\epsilon\KLD\!\left(\P\middle|(\a\otimes\b)\odot e^{-\C/\epsilon}\right)
+
\epsilon\left(1-\sum_{i,j}a_i b_j e^{-\C_{i,j}/\epsilon}\right).
}
\begin{prop}[Reference measure shift for KL]\label{prop-kl-shift}
After removing zero-mass rows and columns, assume that $\a,\a'\in\simplex_n$ and $\b,\b'\in\simplex_m$ have positive entries. For every $\P \in \CouplingsD(\a,\b)$, one has
\(\KLD(\P|\a \otimes \b) = \KLD(\P|\a' \otimes \b') - \KLD(\a|\a') - \KLD(\b|\b').\)
In particular,~\eqref{eq-regularized-discr-rescaled} and~\eqref{eq-regularized-discr} have the same unique minimizer.
\end{prop}
\begin{proof}
Expanding the logarithm and using the marginal constraints gives
\begin{align*}
\KLD(\P|\a\otimes\b)
&=
\KLD(\P|\a'\otimes\b')
+
\sum_i \a_i\log\frac{\a_i'}{\a_i}
+
\sum_j \b_j\log\frac{\b_j'}{\b_j} \\[.25em]
&=
\KLD(\P|\a'\otimes\b')
-\KLD(\a|\a')-\KLD(\b|\b').
\end{align*}
\end{proof}
The tensor-product reference is nevertheless useful when supports vary, because it makes explicit which entries are allowed to vanish. It is also the normalization that passes cleanly to the continuous formulation below, where the reference measure is $\al\otimes\be$ rather than an ambient Lebesgue measure.

\begin{prop}[Convergence with $\epsilon$]\label{prop-convergence-eps}
Assume, after removing zero-mass rows and columns, that $\a$ and $\b$ are positive and that $\C$ is finite.
The unique solution $\P_\epsilon$ of~\eqref{eq-regularized-discr} converges to the optimal solution with maximal entropy within the set of all optimal solutions of the Kantorovich problem, namely
\eql{\label{eq-entropy-conv-1}
\P_\epsilon \overset{\epsilon \rightarrow 0}{\longrightarrow}
\uargmin{\P} \enscond{ -\HD(\P) }{
\P \in \CouplingsD(\a,\b), \dotp{\P}{\C} = \MKD_\C(\a,\b)
}
}
so that in particular
\eq{
\MKD_\C^\epsilon(\a,\b) \overset{\epsilon \rightarrow 0}{\longrightarrow} \MKD_\C(\a,\b).
}
\eqllead{Moreover,}{eq-entropy-conv-2}{
\P_\epsilon \overset{\epsilon \rightarrow \infty}{\longrightarrow}
\a \otimes \b.
}
\end{prop}
\begin{proof}
\textbf{Case $\epsilon \rightarrow 0$.}
We denote $\P_\ell$ the solution of~\eqref{eq-regularized-discr} for $\epsilon=\epsilon_\ell$.

Since $\CouplingsD(\a,\b)$ is bounded, we can extract a sequence (that we do not relabel for the sake of simplicity) such that $\P_\ell \rightarrow \P^\star$. Since $\CouplingsD(\a,\b)$ is closed, $\P^\star \in \CouplingsD(\a,\b)$. Using the equivalent KL-normalized formulation~\eqref{eq-regularized-discr-rescaled}, optimality of $\P$ and $\P_\ell$ for their respective optimization problems (for $\epsilon=0$ and $\epsilon=\epsilon_\ell$) gives
\eql{\label{eq-proof-gamma-conv}
0 \leq \dotp{\C}{\P_\ell} - \dotp{\C}{\P} \leq \epsilon_\ell ( \KLD(\P|\a \otimes \b)-\KLD(\P_\ell|\a \otimes \b) ).
}
Since $\KLD$ is continuous, taking the limit $\ell \rightarrow +\infty$ in this expression shows that
$\dotp{\C}{\P^\star} = \dotp{\C}{\P}$ so that $\P^\star$ is a feasible point of~\eqref{eq-entropy-conv-1}. Furthermore, dividing by $\epsilon_\ell$ in~\eqref{eq-proof-gamma-conv} and taking the limit shows that
$\KLD(\P^\star|\a \otimes \b) \leq \KLD(\P|\a \otimes \b)$, which shows that $\P^\star$ is a solution of~\eqref{eq-entropy-conv-1}. Since the solution $\P_0^\star$ to this program is unique by strict convexity of $\KLD(\cdot|\a \otimes \b)$ on the optimal face, one has $\P^\star = \P_0^\star$, and the whole sequence is converging.

\textbf{Case $\epsilon \rightarrow +\infty$.}
Subtracting $\min_{i,j}\C_{i,j}$ from the cost changes every feasible objective by the same constant, so it does not change the minimizer. We can therefore assume $\C\geq0$. Evaluating the energy at $\a \otimes \b$ (which belongs to the constraint set $\CouplingsD(\a,\b)$), one has
\(\dotp{\C}{\P_\epsilon} + \epsilon \KLD(\P_\epsilon|\a \otimes \b) \leq \dotp{\C}{\a \otimes \b} + \epsilon \times 0\)
and since $\dotp{\C}{\P_\epsilon} \geq 0$, this leads to
\(\KLD(\P_\epsilon|\a \otimes \b) \leq \epsilon^{-1} \dotp{\C}{\a \otimes \b} \leq \frac{\norm{\C}_\infty}{\epsilon}\)
so that $\KLD(\P_\epsilon|\a \otimes \b) \rightarrow 0$ and thus $\P_\epsilon \rightarrow \a \otimes \b$ since $\KLD$ is a valid divergence.
\end{proof}
\subsection{General Formulation}
\paragraph{Measure formulation.}
\begin{defn}[Relative entropy of measures]\label{def-measure-relative-entropy}
For finite nonnegative measures $\pi$ and $\xi$ on $\Xx\times\Yy$, the relative entropy is
\eql{\label{eq-defn-rel-entropy}
\KL(\pi|\xi) \eqdef
\int_{\Xx\times\Yy}\log\!\left(\frac{\d\pi}{\d\xi}\right)\d\pi
+\xi(\Xx\times\Yy)-\pi(\Xx\times\Yy).
}
By convention, $\KL(\pi|\xi)=+\infty$ if $\pi$ is not absolutely continuous with respect to $\xi$.
\end{defn}
\begin{defn}[Continuous entropic optimal transport]\label{def-continuous-entropic-ot}
For $\epsilon>0$, define
\eql{\label{eq-entropic-generic}
\MK_\c^\epsilon(\al,\be) \eqdef
\umin{\pi \in \Couplings(\al,\be)}
\int_{\Xx \times \Yy} c(x,y) \d\pi(x,y) + \epsilon \KL(\pi|\al\otimes\be).
}
\end{defn}
For fixed balanced marginals, the specific product reference in the entropy only matters up to additive constants, exactly as in Proposition~\ref{prop-kl-shift}. More precisely, an alternative product reference $\eta\otimes\zeta$ is admissible when $\al\ll\eta$, $\be\ll\zeta$ and the two marginal relative entropies are finite.

\paragraph{Probabilistic interpretation.}
\begin{defn}[Mutual information]\label{def-mutual-information}
If $(X,Y)\sim\pi$ have marginals $X\sim\al$ and $Y\sim\be$, the mutual information of the pair is \(\Ii(X,Y)\eqdef \KL(\pi|\al\otimes\be).\)
It is nonnegative and vanishes if and only if $X$ and $Y$ are independent.
\end{defn}
With this terminology, the entropic problem~\eqref{eq-entropic-generic} is equivalent to
\(\inf_{X\sim\al,\;Y\sim\be} \EE\bigl(c(X,Y)\bigr)+\epsilon\Ii(X,Y).\)
\paragraph{Sinkhorn for general measures.}
Define the Gibbs kernel and its two integral operators by
\eql{\label{eq-continuous-sinkhorn-operators}
\begin{aligned}
k_\epsilon(x,y)&\eqdef \exp\!\left(-\frac{c(x,y)}{\epsilon}\right),\\
(\mathcal K_\epsilon v)(x)&\eqdef\int_\Yy k_\epsilon(x,y)v(y)\d\be(y),
&
(\mathcal K_\epsilon^*u)(y)&\eqdef\int_\Xx k_\epsilon(x,y)u(x)\d\al(x).
\end{aligned}
}
The star is justified by Fubini's identity
\(
\int_\Xx u\,\mathcal K_\epsilon v\d\al
=
\int_\Yy v\,\mathcal K_\epsilon^*u\d\be
\).
For compact marginal supports and continuous $c$, Propositions~\ref{prop-continuous-entropic-duality} and~\ref{prop-entropic-dual-potentials} provide optimal dual potentials $(f_\epsilon,g_\epsilon)$. Set $u_\epsilon=e^{f_\epsilon/\epsilon}$ and $v_\epsilon=e^{g_\epsilon/\epsilon}$. The continuous density law~\eqref{eq-continuous-entropic-density-law} then becomes
\eql{\label{eq-continuous-sinkhorn-scaling}
\frac{\d\pi_\epsilon}{\d(\al\otimes\be)}(x,y)
=
u_\epsilon(x)k_\epsilon(x,y)v_\epsilon(y),
}
The marginal constraints are therefore equivalent to
\(u_\epsilon\,\mathcal K_\epsilon v_\epsilon=1 \quad \al\text{-a.e.}, \qquad v_\epsilon\,\mathcal K_\epsilon^*u_\epsilon=1 \quad \be\text{-a.e.}\)
Starting, for instance, from $v^{(0)}=1$, continuous Sinkhorn alternately enforces these two identities:
\eql{\label{eq-continuous-sinkhorn-iteration}
u^{(\ell+1)}=\frac{1}{\mathcal K_\epsilon v^{(\ell)}},
\qquad
v^{(\ell+1)}=\frac{1}{\mathcal K_\epsilon^*u^{(\ell+1)}}.
}
Equivalently, the updates read pointwise
\eq{
u^{(\ell+1)}(x)=\frac{1}{\displaystyle\int_\Yy k_\epsilon(x,y)v^{(\ell)}(y)\d\be(y)},
\qquad
v^{(\ell+1)}(y)=\frac{1}{\displaystyle\int_\Xx k_\epsilon(x,y)u^{(\ell+1)}(x)\d\al(x)}.
}
Given $v^{(\ell)}$, the first update produces an intermediate coupling with $\Xx$-marginal $\al$; the second produces one with $\Yy$-marginal $\be$, while generally perturbing the first marginal again. As in the discrete case, $(u,v)$ is defined only up to the gauge $(u,v)\mapsto(\lambda u,v/\lambda)$.

\begin{prop}[Closed-form Gibbs coupling]\label{prop-sinkhorn-gibbs-pushforward}
Let $\be_0$ be a sigma-finite reference measure on $\Yy$ and suppose that
\(Z_\epsilon(x)\eqdef\int_\Yy k_\epsilon(x,y)\d\be_0(y) \in(0,+\infty) \qquad\text{for $\al$-a.e.\ }x.\)
Define the normalized Gibbs transition and its output density by
\eql{\label{eq-gibbs-pushforward-target}
p_\epsilon(x,y)\eqdef\frac{k_\epsilon(x,y)}{Z_\epsilon(x)},
\qquad
q_\epsilon(y)\eqdef\int_\Xx p_\epsilon(x,y)\d\al(x),
\qquad
\d\be_\epsilon=q_\epsilon\,\d\be_0.
}
Assume that $\log Z_\epsilon\in L^1(\al)$, $\log q_\epsilon\in L^1(\be_\epsilon)$, and the plan below has finite entropic objective. Then the unique solution of~\eqref{eq-entropic-generic} between $\al$ and $\be_\epsilon$ is
\eql{\label{eq-closed-form-gibbs-coupling}
\d\pi_\epsilon(x,y)
=
p_\epsilon(x,y)\d\al(x)\d\be_0(y).
}
Relative to $\al\otimes\be_\epsilon$, its Sinkhorn scalings are explicitly
\[
\frac{\d\pi_\epsilon}{\d(\al\otimes\be_\epsilon)}(x,y)
=
\frac{k_\epsilon(x,y)}{Z_\epsilon(x)q_\epsilon(y)}
=
u_\epsilon(x)k_\epsilon(x,y)v_\epsilon(y),
\qquad
u_\epsilon=\frac1{Z_\epsilon},
\quad
v_\epsilon=\frac1{q_\epsilon}
\quad\text{on }\{q_\epsilon>0\}.
\]
The density identity is understood $\al\otimes\be_\epsilon$-almost everywhere; $v_\epsilon$ may be assigned any finite value, for instance zero, on the $\be_\epsilon$-null set $\{q_\epsilon=0\}$.
\end{prop}
\begin{proof}
By construction, $p_\epsilon(x,\cdot)\d\be_0$ is a probability measure. Fubini's theorem therefore shows that~\eqref{eq-closed-form-gibbs-coupling} has first marginal $\al$ and second marginal $\be_\epsilon$.

Let $\gamma\in\Couplings(\al,\be_\epsilon)$ have finite objective. Since $-\epsilon\log k_\epsilon=c$, the displayed density of $\pi_\epsilon$ gives
\[
\epsilon\KL(\gamma|\pi_\epsilon)
=
\int c\,\d\gamma
+\epsilon\KL(\gamma|\al\otimes\be_\epsilon)
+\epsilon\int\log Z_\epsilon\,\d\al
+\epsilon\int\log q_\epsilon\,\d\be_\epsilon.
\]
The last two terms depend only on the prescribed marginals. Thus minimizing the entropic objective is equivalent to minimizing $\KL(\gamma|\pi_\epsilon)$, whose unique minimizer is $\gamma=\pi_\epsilon$.
\end{proof}
If $k_\epsilon(x,\cdot)$ is already normalized with respect to $\be_0$, then $Z_\epsilon=1$ and~\eqref{eq-gibbs-pushforward-target} reduces to the particularly simple formula
\(\frac{\d\be_\epsilon}{\d\be_0}(y) = \int_\Xx k_\epsilon(x,y)\d\al(x).\)
Then $Z_\epsilon=(\pi\epsilon)^{d/2}$ and
\[
p_\epsilon(x,y)
=
(\pi\epsilon)^{-d/2}e^{-\norm{x-y}^2/\epsilon},
\qquad
\be_\epsilon
=
\al*\Gaussian\!\left(0,\frac{\epsilon}{2}\Id\right).
\]
Equivalently,~\eqref{eq-closed-form-gibbs-coupling} is the law of
\(
(X,X+\sqrt{\epsilon/2}\,G)
\)
For discrete measures, setting $\uD_i=\a_i u(x_i)$ and $\vD_j=\b_j v(y_j)$ gives
\(
\P_{i,j}=\a_i\b_j u(x_i)\K_{i,j}v(y_j)=\uD_i\K_{i,j}\vD_j
\),
\paragraph{Convergence with \texorpdfstring{$\epsilon$}{epsilon}.}
\begin{prop}[Convergence with $\epsilon$ for measures]\label{prop-continuous-convergence-epsilon}
Let $\Xx,\Yy$ be compact metric spaces, let $c\in\Cc(\Xx\times\Yy)$, and let $\al\in\Pp(\Xx)$, $\be\in\Pp(\Yy)$. Assume that finite-entropy couplings are dense in $\Couplings(\al,\be)$ for weak convergence with convergence of the cost integral. For $\epsilon>0$, let $\pi_\epsilon$ be a minimizer of~\eqref{eq-entropic-generic}. Then
\(\MK_\c^\epsilon(\al,\be)\longrightarrow \MK_\c(\al,\be) \quad\text{as}\quad \epsilon\downarrow0,\)
and every weak cluster point of $(\pi_\epsilon)_\epsilon$ is an optimal plan for the unregularized Kantorovich problem. If the latter optimal plan is unique, then $\pi_\epsilon$ converges weakly to it. In particular, for $\Xx=\Yy\subset\RR^d$, $c(x,y)=\norm{x-y}^2$ and $\al$ absolutely continuous, $\pi_\epsilon\rightharpoonup(\Id,T)_\sharp\al$, where $T$ is the Brenier map.
As $\epsilon\to+\infty$,
\[
\pi_\epsilon\longrightarrow \al\otimes\be
\quad\text{in total variation},\qquad
\MK_\c^\epsilon(\al,\be)\longrightarrow \int_{\Xx\times\Yy} c\,\d\al\d\be.
\]
\end{prop}
\begin{proof}
Existence follows from compactness of $\Couplings(\al,\be)$ and lower semicontinuity of the relative entropy. For the limit $\epsilon\downarrow0$, take any sequence $\epsilon_k\downarrow0$ and a weak cluster point $\pi^\star$ of $\pi_{\epsilon_k}$. Since $\KL(\pi|\al\otimes\be)\geq0$ on probability couplings,
\(\liminf_k \MK_\c^{\epsilon_k}(\al,\be) \geq \liminf_k \int c\,\d\pi_{\epsilon_k} = \int c\,\d\pi^\star \geq \MK_\c(\al,\be).\)
Conversely, by the density assumption, for every $\delta>0$ there exists a coupling $\pi^\delta$ with finite relative entropy and $\int c\,\d\pi^\delta\leq \MK_\c(\al,\be)+\delta$. Testing~\eqref{eq-entropic-generic} with $\pi^\delta$ gives
\(\limsup_k \MK_\c^{\epsilon_k}(\al,\be) \leq \int c\,\d\pi^\delta.\)
Letting $\delta\downarrow0$ proves convergence of the values and forces $\int c\,\d\pi^\star=\MK_\c(\al,\be)$. Uniqueness of the unregularized optimizer then identifies all cluster points. Brenier's theorem gives the stated quadratic consequence when $\al$ is absolutely continuous.

For the limit $\epsilon\to+\infty$, subtract a constant from $c$ and assume $c\geq0$. Since $\al\otimes\be$ is admissible and has zero relative entropy with respect to itself,
\(\int c\,\d\pi_\epsilon+\epsilon\KL(\pi_\epsilon|\al\otimes\be) \leq \int c\,\d\al\d\be.\)
Hence $\KL(\pi_\epsilon|\al\otimes\be)\leq \epsilon^{-1}\int c\,\d\al\d\be$. Pinsker's inequality from Theorem~\ref{thm-pinsker} gives total-variation convergence to $\al\otimes\be$, and the convergence of the values follows by continuity of $c$.
\end{proof}
\paragraph{Large-temperature expansion.}
When $\epsilon$ is large, the entropy dominates and the optimal plan is a small perturbation of the product coupling. The useful quantity is the component of the cost that cannot be absorbed into row and column potentials.

\begin{prop}[Large-temperature expansion]\label{prop-large-epsilon-expansion}
Let $r=\al\otimes\be$ and assume $c\in L^\infty(r)$. Define the product averages
\(\bar c_\Xx(x)\eqdef\int_\Yy c(x,y)\,\d\be(y),\qquad \bar c_\Yy(y)\eqdef\int_\Xx c(x,y)\,\d\al(x),\qquad \bar c\eqdef\int_{\Xx\times\Yy}c\,\d r,\)
and the centered interaction
\(c_0(x,y)\eqdef c(x,y)-\bar c_\Xx(x)-\bar c_\Yy(y)+\bar c.\)
Assume that the large-temperature branch $p_\epsilon=\d\pi_\epsilon/\d r$, $f_\epsilon,g_\epsilon$ admits an expansion to third order in $\epsilon^{-1}$ near $0$ in $L^\infty(r)$, with gauge $\int g_\epsilon\,\d\be=0$. Then
\(p_\epsilon(x,y)=1-\frac{c_0(x,y)}{\epsilon}+O(\epsilon^{-2}) \quad\text{in }L^2(r),\)
and
\[
\MK_\c^\epsilon(\al,\be)
=
\bar c-\frac{1}{2\epsilon}\int_{\Xx\times\Yy} c_0(x,y)^2\,\d r(x,y)
+\frac{1}{6\epsilon^2}\int_{\Xx\times\Yy} c_0(x,y)^3\,\d r(x,y)
+O(\epsilon^{-3}).
\]
Moreover, with \(A(x)\eqdef\int_\Yy c_0(x,y)^2\,\d\be(y),\qquad B(y)\eqdef\int_\Xx c_0(x,y)^2\,\d\al(x),\qquad \sigma^2\eqdef\int_{\Xx\times\Yy}c_0^2\,\d r,\)
one has \(f_\epsilon(x)=\bar c_\Xx(x)-\frac{A(x)}{2\epsilon}+O(\epsilon^{-2}),\)
and
\(g_\epsilon(y)=\bar c_\Yy(y)-\bar c +\frac{\sigma^2-B(y)}{2\epsilon}+O(\epsilon^{-2})\)
in the same gauge.
\end{prop}
\begin{proof}
The optimality equation can be written
\(p_\epsilon(x,y)=\exp\pa{\frac{f_\epsilon(x)+g_\epsilon(y)-c(x,y)}{\epsilon}},\)
with the constraints $\int p_\epsilon(x,y)\,\d\be(y)=1$ and $\int p_\epsilon(x,y)\,\d\al(x)=1$. Write $p_\epsilon=1+\epsilon^{-1}q+O(\epsilon^{-2})$. The linearized marginal constraints impose
\(\int q(x,y)\,\d\be(y)=0, \qquad \int q(x,y)\,\d\al(x)=0.\)
Expanding the objective gives
\[
\int c\,p_\epsilon\,\d r+\epsilon\int p_\epsilon\log(p_\epsilon)\,\d r
=
\bar c+\frac1\epsilon
\pa{\int c\,q\,\d r+\frac12\int q^2\,\d r}
+O(\epsilon^{-2}).
\]
Thus $q$ is the minimizer of the quadratic expression over zero-row and zero-column perturbations. Orthogonal projection in $L^2(r)$ gives $q=-c_0$, because $c_0$ has zero conditional means. This proves the first-order density and value expansions.

To identify the next displayed terms, write $p_\epsilon=1+\epsilon^{-1}q+\epsilon^{-2}s+O(\epsilon^{-3})$ and
$f_\epsilon=f_0+\epsilon^{-1}f_1+O(\epsilon^{-2})$, $g_\epsilon=g_0+\epsilon^{-1}g_1+O(\epsilon^{-2})$, where $f_0=\bar c_\Xx$ and $g_0=\bar c_\Yy-\bar c$. Since
\(p_\epsilon = \exp\pa{-\frac{c_0}{\epsilon}+\frac{f_1(x)+g_1(y)}{\epsilon^2}+O(\epsilon^{-3})},\)
one has $s=f_1+g_1+c_0^2/2$. The second-order marginal constraints and the gauge $\int g_1\,\d\be=0$ give $f_1=-A/2$ and $g_1=(\sigma^2-B)/2$. Finally,
\(p_\epsilon\log p_\epsilon = \epsilon^{-1}q+\epsilon^{-2}\pa{s+\frac12 q^2} +\epsilon^{-3}\pa{q s-\frac16 q^3}+O(\epsilon^{-4}),\)
and the zero-conditional-mean property of $c_0$ cancels all additive terms in $s$. The order-$\epsilon^{-2}$ contribution to the value is therefore
$\frac12\int c_0^3\,\d r-\frac13\int c_0^3\,\d r=\frac16\int c_0^3\,\d r$.
\end{proof}
\paragraph{Small-temperature expansion for smooth densities.}
\begin{prop}[Small-temperature quadratic expansion]\label{prop-small-epsilon-expansion}
Let $\al=\density{\al}\,\d x$ and $\be=\density{\be}\,\d x$ be probability measures on $\RR^d$ with bounded compactly supported densities. Let $\al_t=\rho_t\,\d x$ be their quadratic displacement interpolation and assume
\(\mathcal I_{\rm geo}(\al,\be)\eqdef \int_0^1\int_{\RR^d} \norm{\nabla\log\rho_t(x)}^2\rho_t(x)\,\d x\,\d t<+\infty.\)
For $c(x,y)=\norm{x-y}^2$, define
\[
\mathrm H(\al)\eqdef\int_{\RR^d}\density{\al}\log\density{\al}\,\d x,
\qquad
\mathrm H(\be)\eqdef\int_{\RR^d}\density{\be}\log\density{\be}\,\d x.
\]
Then, as $\epsilon\downarrow0$,
\[
\MK_{\norm{\cdot}^2}^{\epsilon}(\al,\be)
=
\Wass_2^2(\al,\be)
-\frac{d\epsilon}{2}\log(\pi\epsilon)
-\frac{\epsilon}{2}\pa{\mathrm H(\al)+\mathrm H(\be)}
+\frac{\epsilon^2}{16}\mathcal I_{\rm geo}(\al,\be)
+o(\epsilon^2).
\]
If, in addition, the normalized Kantorovich pair $(f_0,g_0)$ is unique, let $(f_\epsilon,g_\epsilon)$ be optimal Sinkhorn potentials normalized by $\int g_\epsilon\,\d\be=\int g_0\,\d\be=0$. Then
\(f_\epsilon\longrightarrow f_0\quad\text{in }L^1(\al), \qquad g_\epsilon\longrightarrow g_0\quad\text{in }L^1(\be)\)
as $\epsilon\downarrow0$. If the potentials are chosen as their continuous soft $c$-transform representatives, this convergence is uniform on $\Supp(\al)$ and $\Supp(\be)$. Moreover, the scalar part of the first potential is fixed by
\[
\int f_\epsilon\,\d\al
=
\Wass_2^2(\al,\be)
-\frac{d\epsilon}{2}\log(\pi\epsilon)
-\frac{\epsilon}{2}\bigl(\mathrm H(\al)+\mathrm H(\be)\bigr)
+\frac{\epsilon^2}{16}\mathcal I_{\rm geo}(\al,\be)
+o(\epsilon^2).
\]
\end{prop}
\begin{proof}
For the Brownian reference with heat kernel variance $T$, the dynamic Schr\"odinger formulation gives
\[
T\mathscr C_T(\al,\be)
=
\frac12\Wass_2^2(\al,\be)
+\frac{T}{2}\pa{\mathrm H(\al)+\mathrm H(\be)}
+\frac{T^2}{8}\mathcal I_{\rm geo}(\al,\be)+o(T^2),
\]
under the stated regularity assumptions. Here $\mathscr C_T$ uses entropy relative to the sigma-finite Brownian endpoint measure:
\[
\mathscr C_T(\al,\be)
=
\inf_{\pi\in\Couplings(\al,\be)}
\mathscr H(\pi|p_T(x,y)\,\d x\d y),
\qquad
p_T(x,y)=(2\pi T)^{-d/2}e^{-\norm{x-y}^2/(2T)},
\]
where $\mathscr H(\pi|\xi)\eqdef\int\log(\d\pi/\d\xi)\d\pi$ for sigma-finite $\xi$. When $\xi$ is a probability measure, this agrees with Definition~\ref{def-measure-relative-entropy}; unlike the generalized KL of that definition, it remains meaningful when the Brownian reference has infinite total mass. Expanding the heat kernel, multiplying by $2T$, and setting $\epsilon=2T$ gives
\(2T\mathscr C_T(\al,\be) = \MK_{\norm{\cdot}^2}^{\epsilon}(\al,\be) +2T\pa{\mathrm H(\al)+\mathrm H(\be)} +dT\log(2\pi T).\)
Equivalently, \(\MK_{\norm{\cdot}^2}^{\epsilon}(\al,\be) = 2T\mathscr C_T(\al,\be) -2T\pa{\mathrm H(\al)+\mathrm H(\be)} -dT\log(2\pi T),\)
which is exactly the displayed expansion because $2T=\epsilon$. At a dual optimum the exponential penalty integrates to zero, so the gauge $\int g_\epsilon\d\be=0$ gives $\int f_\epsilon\d\al=\MK_{\norm{\cdot}^2}^\epsilon(\al,\be)$. This proves the displayed scalar expansion. The cited compactness result gives subsequential $L^1$ convergence of the normalized potentials toward Kantorovich potentials; uniqueness identifies every cluster point and hence gives convergence of the full family.

On the two compact supports, the quadratic cost has a common modulus of continuity in either variable. The soft $c$-transform formula therefore gives, uniformly in $\epsilon$,
\(\abs{f_\epsilon(x)-f_\epsilon(x')} \leq\sup_{y\in\Supp(\be)}\abs{c(x,y)-c(x',y)},\)
and the same estimate holds for $g_\epsilon$ and for the limiting hard $c$-transforms. Thus both families are equicontinuous. If uniform convergence failed, equicontinuity would produce a fixed neighborhood on which the discrepancy stays bounded away from zero. Every neighborhood of a point in the support has positive marginal mass, contradicting $L^1$ convergence. Hence the convergence is uniform on both supports.
\end{proof}
\subsection{Dual of Sinkhorn}\label{sec-dual-sinkhorn}
\paragraph{Discrete dual.}
\begin{prop}[Dual of entropic OT]
The optimal value of~\eqref{eq-regularized-discr-rescaled} is
\eql{\label{eq-dual-formulation}
\umin{\P \in \CouplingsD(\a,\b)}
\dotp{\P}{\C}+\epsilon\KLD(\P|\a\otimes\b)
=
\umax{\fD \in \RR^n,\gD \in \RR^m}
\dotp{\fD}{\a} + \dotp{\gD}{\b}
- \epsilon \sum_{i,j} \exp\pa{
\frac{\fD_i+\gD_j-\C_{i,j}}{\epsilon}
} \a_i \b_j + \epsilon.
}
The optimal $(\fD,\gD)$ are linked to scalings $(\uD,\vD)$ appearing in~\eqref{eq-scaling-form} through
\eql{\label{eq-entropy-pd}
\uD_i=\a_i e^{\fD_i/\epsilon}
\qandq
\vD_j=\b_j e^{\gD_j/\epsilon}.
}
\end{prop}
\begin{proof}
We introduce Lagrange multipliers and consider
\(\umin{\P \geq 0} \umax{\fD,\gD} \dotp{\C}{\P} + \epsilon \KLD(\P|\a \otimes \b) + \dotp{\a-\P\ones}{\fD} + \dotp{\b-\P^\top\ones}{\gD}.\)
Finite-dimensional convex duality allows us to exchange the minimum over $\P$ with the maximum over $(\fD,\gD)$, giving
\eq{
\umax{\fD,\gD}
\dotp{\fD}{\a} + \dotp{\gD}{\b}
+
\epsilon \umin{\P \geq 0}
\pa{
\KLD(\P|\a \otimes \b)
-
\dotp{\frac{\fD\oplus\gD-\C}{\epsilon}}{\P}
}
=
\dotp{\fD}{\a} + \dotp{\gD}{\b} - \epsilon \KLD^*\pa{ \frac{\fD\oplus\gD-\C}{\epsilon}|\a \otimes \b }.
}
One concludes by using~\eqref{eq-legendre} for $\phi(r)=r \log(r)-r+1$
\(\KLD^*(H|\a \otimes \b) = \sum_{i,j} \phi^*(H_{i,j}) \a_i \b_j.\)
Indeed, the scalar maximization
\(\phi^*(s)=\sup_{r\geq0}\{rs-r\log r+r-1\}\)
has first-order condition $s-\log r=0$, hence $r=e^s$ and $\phi^*(s)=e^s-1$.
\end{proof}
\paragraph{\texorpdfstring{Discrete soft $c$-transforms.}{Discrete soft c-transforms.}}
For a fixed $\gD$, maximizing with respect to $\fD$ leads to the following equation after zeroing the derivative with respect to $\fD$:
\eq{
\a_i - e^{\frac{\fD_i}{\epsilon}} \a_i \sum_j \exp\pa{
\frac{\gD_j-\C_{i,j}}{\epsilon}
} \b_j = 0
}
which leads to the explicit solution
\eq{
\fD_i = -\epsilon \log \sum_j \exp\pa{ \frac{\gD_j-\C_{i,j}}{\epsilon} } \b_j.
}
\begin{defn}[Soft-min and discrete soft $c$-transform]\label{def-discrete-soft-c-transform}
For $h\in\RR^m$ and weights $\b\in\simplex_m$, the weighted soft-min at temperature $\epsilon>0$ is \({\min}_{\b}^\epsilon(h) \eqdef -\epsilon \log \sum_j e^{-h_j/\epsilon} \b_j.\)
It converges to $\min_{j:\,b_j>0}h_j$ as $\epsilon\to0$; for positive weights this is $\min_jh_j$. Given a cost matrix $\C$, the discrete soft $c$-transforms are
\eql{\label{eq-discrete-soft-c-transforms}
\fD_i = {\min}_{\b}^\epsilon( \C_{i,\cdot} - \gD ),
\qquad
\gD_j = {\min}_{\a}^\epsilon( \C_{\cdot,j} - \fD ).
}
\end{defn}
This follows from noticing that, similarly to the minimum operator, one has \({\min}_{\b}^\epsilon(h-\text{cst})={\min}_{\b}^\epsilon(h)-\text{cst}\)
\paragraph{Continuous dual and soft-transforms.}
\begin{prop}[Continuous entropic duality]\label{prop-continuous-entropic-duality}
Let $\Xx$ and $\Yy$ be compact metric spaces, let $\al\in\Pp(\Xx)$ and $\be\in\Pp(\Yy)$, and let $c\in\Cc(\Xx\times\Yy)$. For every $\epsilon>0$, the KL-regularized problem~\eqref{eq-entropic-generic} satisfies
\eql{\label{eq-dual-sinkh-cont}
\MK_\c^\epsilon(\al,\be)
=
\usup{f\in\Cc(\Xx),\,g\in\Cc(\Yy)}
\Dd_\epsilon(f,g),
}
where the concave dual objective is
\eql{\label{eq-dual-sinkhorn-objective}
\Dd_\epsilon(f,g)
\eqdef
\int_{\Xx} f(x)\d\al(x)
+
\int_{\Yy} g(y)\d\be(y)
-
\epsilon
\int_{\Xx\times\Yy}
\pa{
e^{\frac{f(x)+g(y)-c(x,y)}{\epsilon}}-1
}
\d\al(x)\d\be(y).
}
If $\pi^\star$ and $(f^\star,g^\star)$ are primal and dual optimizers, then
\eql{\label{eq-continuous-entropic-density-law}
\frac{\d\pi^\star}{\d(\al\otimes\be)}(x,y)
=
\exp\!\left(\frac{f^\star(x)+g^\star(y)-c(x,y)}{\epsilon}\right)
\qquad (\al\otimes\be)\text{-a.e.}
}
\end{prop}
\begin{proof}
Write $\xi=\al\otimes\be$ and $\phi(r)=r\log r-r+1$ for $r\geq0$. Since the primal objective is infinite unless $\pi\ll\xi$, put $r=\d\pi/\d\xi$. The marginal constraints are equivalent to
\(\int_{\Yy}r(x,y)\d\be(y)=1\quad\al\text{-a.e.}, \qquad \int_{\Xx}r(x,y)\d\al(x)=1\quad\be\text{-a.e.}\)
Introduce potentials $f$ and $g$ for these two constraints. The Lagrangian is
\[
\int_{\Xx}f\d\al+\int_{\Yy}g\d\be
+
\int_{\Xx\times\Yy}
\left[
\epsilon\phi(r(x,y))
+\bigl(c(x,y)-f(x)-g(y)\bigr)r(x,y)
\right]\d\xi(x,y).
\]
The scalar conjugate of $\phi$ is $\phi^*(s)=e^s-1$. Therefore pointwise minimization over $r\geq0$ gives
\[
\inf_{r\geq0}
\left\{
\epsilon\phi(r)+(c-f-g)r
\right\}
=
-\epsilon\phi^*\!\left(\frac{f+g-c}{\epsilon}\right)
=
-\epsilon\left(e^{(f+g-c)/\epsilon}-1\right).
\]
Integrating this identity yields $\Dd_\epsilon(f,g)$ and proves weak duality. Strong duality follows from Fenchel--Rockafellar entropy duality in the exponential Orlicz pair: the feasible density $r\equiv1$, corresponding to the product coupling $\xi$, is strictly positive and supplies the qualification point, while the boundedness of $c$ makes its objective finite. Because $c$ is continuous on the compact product, the dual supremum can be restricted to continuous potentials; equivalently, alternating soft $c$-transforms regularize bounded potentials into continuous ones without decreasing the dual value. This proves~\eqref{eq-dual-sinkh-cont}.

Finally, equality in the pointwise Fenchel inequality forces $r^\star(x,y)$ to be the unique minimizer of the scalar problem above. Its first-order condition is \(\epsilon\log r^\star(x,y) {}+c(x,y)-f^\star(x)-g^\star(y)=0.\)
Exponentiation gives~\eqref{eq-continuous-entropic-density-law}.
\end{proof}
\begin{defn}[Continuous soft $c$-transforms]\label{def-continuous-soft-c-transform}
For $f\in\Cc(\Xx)$ and $g\in\Cc(\Yy)$, define
\begin{align}
f^{c,\epsilon}(y)
&\eqdef
-\epsilon\log\!\left(
\int_{\Xx}
e^{\frac{f(x)-c(x,y)}{\epsilon}}\d\al(x)
\right),
\qquad y\in\Yy, \label{eq-soft-c-cont-f}\\
g^{\bar c,\epsilon}(x)
&\eqdef
-\epsilon\log\!\left(
\int_{\Yy}
e^{\frac{g(y)-c(x,y)}{\epsilon}}\d\be(y)
\right),
\qquad x\in\Xx. \label{eq-soft-c-cont-g}
\end{align}
\end{defn}
\begin{prop}[Existence and uniqueness of entropic dual potentials]\label{prop-entropic-dual-potentials}
Assume that $\Xx=\supp(\al)$ and $\Yy=\supp(\be)$ are compact and that $c$ is continuous. The dual problem~\eqref{eq-dual-sinkh-cont} has solutions, and its set of solutions is
\((f^\star+\lambda,g^\star-\lambda), \qquad \lambda\in\RR.\)
\end{prop}
\begin{proof}
Normalize potentials by imposing $\int f\d\al=0$. Replacing any pair $(f,g)$ by the corresponding soft $c$-transforms does not decrease the dual objective, because each soft transform is the exact maximizer in one block variable. For transformed potentials, the oscillations are bounded by the oscillation of the cost:
\(\norm{f}_V+\norm{g}_V \leq 2(\sup c-\inf c)\).
Moreover the modulus of continuity of the soft transforms is controlled by that of $c$; for instance
\(\abs{g^{\bar c,\epsilon}(x)-g^{\bar c,\epsilon}(x')} \leq \sup_y\abs{c(x,y)-c(x',y)}.\)
After normalization, maximizing sequences are therefore uniformly bounded and equicontinuous. Arzel\`a--Ascoli gives a uniformly converging subsequence, and continuity of~\eqref{eq-dual-sinkhorn-objective} gives a maximizer.

For uniqueness, use strict convexity of
\(H\mapsto \int e^{H/\epsilon}\d(\al\otimes\be)\)
on the image of $(f,g)\mapsto H=f\oplus g-c$, modulo constants. If two optimal pairs exist, their midpoint is also optimal; strict convexity forces the two functions $f\oplus g$ to agree $\al\otimes\be$-almost everywhere. Full support and continuity then imply equality everywhere on $\Xx\times\Yy$, so $f-f'$ is constant and $g-g'$ is the opposite constant. Without replacing the ambient spaces by the marginal supports, uniqueness has this meaning only on those supports; values outside them do not enter the dual objective.
\end{proof}
\paragraph{Dual Sinkhorn for general measures.}\label{par-continuous-dual-sinkhorn}
Indeed, for fixed $g$ and a perturbation $h\in\Cc(\Xx)$,
\eq{
\left.\frac{\d}{\d s}\Dd_\epsilon(f+s h,g)\right|_{s=0}
=
\int_\Xx h(x)\left[
1-e^{f(x)/\epsilon}
\int_\Yy e^{(g(y)-c(x,y))/\epsilon}\d\be(y)
\right]\d\al(x).
}
Consequently, maximizing first over $f$ and then over $g$ gives the alternating updates
\eql{\label{eq-continuous-dual-sinkhorn-iteration}
f^{(\ell+1)}=(g^{(\ell)})^{\bar c,\epsilon},
\qquad
g^{(\ell+1)}=(f^{(\ell+1)})^{c,\epsilon}.
}
More precisely, setting
\(u^{(\ell)}=e^{f^{(\ell)}/\epsilon}, \qquad v^{(\ell)}=e^{g^{(\ell)}/\epsilon}\)
and using the kernel operators~\eqref{eq-continuous-sinkhorn-operators}, exponentiation of~\eqref{eq-continuous-dual-sinkhorn-iteration} yields
\(u^{(\ell+1)}=\frac{1}{\mathcal K_\epsilon v^{(\ell)}}, \qquad v^{(\ell+1)}=\frac{1}{\mathcal K_\epsilon^*u^{(\ell+1)}}.\)
Thus the coupling density reconstructed from the current potentials is
\(
u^{(\ell)}(x)k_\epsilon(x,y)v^{(\ell)}(y)
\)
with respect to $\al\otimes\be$, exactly as in the scaling law~\eqref{eq-continuous-sinkhorn-scaling}. At a fixed point, its two marginals are $\al$ and $\be$; equivalently, the potentials jointly maximize the continuous dual.

\paragraph{Neural dual solvers.}
For the bilinear cost \(c(x,y)=-\dotp{x}{y}\), the signs of the dual potentials are convex: writing \(\Phi=-f\) and \(\Psi=-g\), the zero-temperature constraint \(f(x)+g(y)\leq-\dotp{x}{y}\) becomes the Fenchel-type inequality
\(\Phi(x)+\Psi(y)\geq \dotp{x}{y}.\)
For the quadratic cost this is the same statement after subtracting the quadratic terms, using the convexity properties of soft transforms. One can therefore maximize the continuous dual~\eqref{eq-dual-sinkhorn-objective} over parameterized convex potentials, estimating the integrals by stochastic samples.

\subsection{Marginal-Dependent Problems}
\label{sec-sinkhorn-marginal-dependent}
Let $\mathcal F$ and $\mathcal G$ be proper convex lower semicontinuous functionals on finite nonnegative measures on $\Xx$ and $\Yy$. The unregularized marginal-dependent transport problem is
\begin{equation}\label{eq-marginal-dependent-unregularized-cont}
\inf_{\pi\in\Mm_+(\Xx\times\Yy)}
\int c(x,y)\d\pi(x,y)
+
\mathcal F(\pi_1)
+
\mathcal G(\pi_2),
\end{equation}
where $\pi_1=(\mathrm{p}_{\Xx})_\sharp\pi$ and $\pi_2=(\mathrm{p}_{\Yy})_\sharp\pi$ are the two marginals of $\pi$. Entropic regularization turns this problem into the scaling-friendly form below.

\begin{defn}[Marginal-dependent entropic transport]\label{def-marginal-dependent-entropic-transport}
For reference measures \(\al,\be\), proper convex marginal functionals \(\mathcal F,\mathcal G\), and \(\epsilon>0\), define
\begin{equation}\label{eq-marginal-dependent-cont}
\inf_{\pi\in\Mm_+(\Xx\times\Yy)}
\int c(x,y)\d\pi(x,y)
+
\mathcal F(\pi_1)
+
\mathcal G(\pi_2)
+
\epsilon\KL(\pi|\al\otimes\be).
\end{equation}
For discrete reference weights \(\a,\b\) and convex functions
\(\mathsf F:\RR_+^n\to\RR\cup\{+\infty\}\),
\(\mathsf G:\RR_+^m\to\RR\cup\{+\infty\}\), its matrix form is
\begin{equation}\label{eq-marginal-dependent-discrete}
\inf_{\P\in\RR_+^{n\times m}}
\dotp{\C}{\P}
+
\mathsf F(\P\ones_m)
+
\mathsf G(\transp{\P}\ones_n)
+
\epsilon\KLD(\P|\a\otimes\b).
\end{equation}
\end{defn}
In this subsection, when the total mass of $\pi$ is not fixed, the KL term is understood in the generalized sense associated with $\varphi(s)=s\log s-s+1$. Thus, if $\lambda=\al\otimes\be$ and $\pi=\rho\lambda+\pi^\perp$ is the Lebesgue decomposition of $\pi$ with respect to $\lambda$, then
\(\KL(\pi|\lambda) = \int \bigl(\rho\log\rho-\rho+1\bigr)\d\lambda\)
Equivalently, if $\K_{i,j}=a_i b_j e^{-\C_{i,j}/\epsilon}$, the terms involving $\C$ can be absorbed into the Gibbs reference and the problem is, up to the additive constant $\epsilon\sum_{i,j}(a_i b_j-\K_{i,j})$,
\(\inf_{\P\geq0} \mathsf F(\P\ones_m) + \mathsf G(\transp{\P}\ones_n) + \epsilon\KLD(\P|\K).\)
\begin{prop}[Dual and scaling for marginal penalties]\label{prop-marginal-dependent-dual-scaling}
Assume $a_i,b_j>0$, $\epsilon>0$, that $\C$ has finite entries, and that $\mathsf F,\mathsf G$ are proper lower-semicontinuous convex functions. Assume also a Fenchel qualification condition, for instance the existence of a matrix $\P>0$ such that $\P\ones_m$ and $\transp{\P}\ones_n$ belong to the relative interiors of $\dom(\mathsf F)$ and $\dom(\mathsf G)$. The Fenchel dual of~\eqref{eq-marginal-dependent-discrete} is
\begin{equation}\label{eq-marginal-dependent-dual}
\sup_{\fD\in\RR^n,\gD\in\RR^m}
-
\mathsf F^*(-\fD)
-
\mathsf G^*(-\gD)
-
\epsilon\sum_{i,j}a_i b_j
\left[
\exp\left(\frac{\fD_i+\gD_j-\C_{i,j}}{\epsilon}\right)-1
\right].
\end{equation}
Equivalently, up to the additive constant $\epsilon\sum_{i,j}a_i b_j$, the last term is
\[
-
\epsilon\sum_{i,j}\K_{i,j}
\exp\left(\frac{\fD_i+\gD_j}{\epsilon}\right).
\]
If $\uD=e^{\fD/\epsilon}$ and $\vD=e^{\gD/\epsilon}$, exact block ascent in the two dual variables is the generalized Sinkhorn cycle
\begin{equation}\label{eq-generalized-sinkhorn-kl-prox}
\begin{aligned}
r &\leftarrow \prox_{\mathsf F/\epsilon}^{\KLD}(\K\vD),
&\qquad
\uD &\leftarrow r\oslash(\K\vD),\\
s &\leftarrow \prox_{\mathsf G/\epsilon}^{\KLD}(\transp{\K}\uD),
&\qquad
\vD &\leftarrow s\oslash(\transp{\K}\uD),\\
\P&=\diag(\uD)\K\diag(\vD),
\end{aligned}
\end{equation}
where divisions are entrywise and
\begin{equation}\label{eq-kl-prox-marginal}
\prox_{\mathsf h}^{\KLD}(z)
\eqdef
\uargmin{r\in\RR_+^d}
\mathsf h(r)+\KLD(r|z).
\end{equation}
In particular, $\prox_{\mathsf F/\epsilon}^{\KLD}(z)$ is the minimizer of $\mathsf F(r)+\epsilon\KLD(r|z)$. Each of the two proximal problems has a unique minimizer. Under the strict-feasibility assumption above, these minimizers are positive, so every division in~\eqref{eq-generalized-sinkhorn-kl-prox} is well-defined.
\end{prop}
\begin{proof}
Introduce independent variables $r$ and $s$ for the two marginals and use dual variables $\fD,\gD$ for the constraints $r=\P\ones_m$ and $s=\transp{\P}\ones_n$. The Lagrangian contains
\(\mathsf F(r)+\dotp{\fD}{r} + \mathsf G(s)+\dotp{\gD}{s} + \dotp{\C}{\P} + \epsilon\KLD(\P|\a\otimes\b) - \dotp{\fD\oplus\gD}{\P}.\)
Minimizing over $r$ and $s$ gives $-\mathsf F^*(-\fD)$ and $-\mathsf G^*(-\gD)$. For each scalar entry, the convex conjugate of $p\mapsto \C_{i,j}p+\epsilon(p\log(p/(a_i b_j))-p+a_i b_j)$ is $q\mapsto \epsilon a_i b_j(e^{(q-\C_{i,j})/\epsilon}-1)$. Minimizing over $\P$ with $q=\fD_i+\gD_j$ gives~\eqref{eq-marginal-dependent-dual}.

For the scaling form, fix $\vD>0$, set $\widetilde{\K}=\K\diag(\vD)$ and $z=\widetilde{\K}\ones_m=\K\vD$. Updating the row scaling means looking for $\P=\diag(\uD)\widetilde{\K}$, so its row marginal is $r=\P\ones_m=\uD\odot z$. The row-wise chain rule gives
\(\KLD(\diag(\uD)\widetilde{\K}|\widetilde{\K}) = \sum_i \bigl(r_i\log(r_i/z_i)-r_i+z_i\bigr) = \KLD(r|z).\)
Thus exact optimization of this block is equivalent to minimizing $\mathsf F(r)+\epsilon\KLD(r|z)$ over $r\geq0$, hence $r=\prox_{\mathsf F/\epsilon}^{\KLD}(z)$ and $\uD=r\oslash z$. The column update is identical with $w=\transp{\K}\uD$.

The KL term is strictly convex and superlinear for $z>0$. Together with proper lower semicontinuity and convexity of $\mathsf F$, this gives existence and uniqueness of the row proximal minimizer. Since $\C$ is finite and the current scalings are positive, $z>0$. Let $\bar r>0$ belong to $\dom(\mathsf F)$, as supplied by strict feasibility. If a proximal minimizer $r$ had a zero coordinate, then along $r_t=(1-t)r+t\bar r$ convexity bounds the change of $\mathsf F$ by $O(t)$, whereas the corresponding KL term contains $t\bar r_i\log t$. This negative $t\log(1/t)$ variation dominates all $O(t)$ terms and contradicts minimality. Thus $r>0$; the same argument applies to the column block.
\end{proof}
When $\mathsf F=\iota_{\{\a\}}$ and $\mathsf G=\iota_{\{\b\}}$, the first two dual terms are $\dotp{\fD}{\a}+\dotp{\gD}{\b}$, and one recovers the usual entropic dual~\eqref{eq-dual-sinkhorn-objective}. The classical Sinkhorn update is the special case in which the KL proximal maps return the prescribed marginals $\a$ and $\b$.

\begin{itemize}
\item \textbf{Hard marginal constraint.} If $\mathsf F=\iota_{\{\a\}}$, then $\prox_{\mathsf F/\epsilon}^{\KLD}(z)=\a$.
\item \textbf{KL marginal relaxation.} If $\mathsf F(r)=\tau\KLD(r|\a)$ with $\tau>0$, then, coordinatewise,
\(\prox_{\mathsf F/\epsilon}^{\KLD}(z) = z^{\epsilon/(\tau+\epsilon)}\odot \a^{\tau/(\tau+\epsilon)},\)
and the scaling update becomes $\uD\leftarrow(\a\oslash z)^{\tau/(\tau+\epsilon)}$, the damped scaling of unbalanced Sinkhorn.
\item \textbf{Pointwise bounds.} If $\mathsf F=\iota_{\{\ell\leq r\leq u\}}$, then the proximal map is the coordinatewise clipping $\prox_{\mathsf F/\epsilon}^{\KLD}(z)_i=\min\{u_i,\max\{\ell_i,z_i\}\}$.
\item \textbf{Total-variation marginal relaxation.} If $\mathsf F(r)=\tau\norm{r-\a}_1$ and $\lambda=\tau/\epsilon$, then the proximal map is coordinatewise
\[
\prox_{\mathsf F/\epsilon}^{\KLD}(z)_i
=
\begin{cases}
z_i e^{\lambda}, & z_i<a_i e^{-\lambda},\\
a_i, & a_i e^{-\lambda}\leq z_i\leq a_i e^{\lambda},\\
z_i e^{-\lambda}, & z_i>a_i e^{\lambda}.
\end{cases}
\]
\item \textbf{Fixed total mass.} If $\mathsf F=\iota_{\{\dotp{r}{\ones}=m\}}$, then $\prox_{\mathsf F/\epsilon}^{\KLD}(z)=m z/\dotp{z}{\ones}$.
\end{itemize}
\subsection{Heat Kernels and Hopf--Cole Transforms}
\label{sec-sinkhorn-heat-hopf-cole}
\paragraph{Geodesics in heat.}
On $\RR^d$, the heat kernel for $\partial_t u=\Delta u$ is
\[
h_t(x,y)=(4\pi t)^{-d/2}\exp\!\left(-\frac{\norm{x-y}^2}{4t}\right).
\]
For the quadratic cost $c(x,y)=\norm{x-y}^2$, the Gibbs kernel used by Sinkhorn satisfies
\eql{\label{eq-sinkhorn-heat-kernel}
\K_\epsilon(x,y)
\eqdef e^{-\norm{x-y}^2/\epsilon}
=
(\pi\epsilon)^{d/2}h_{\epsilon/4}(x,y).
}
On a Riemannian manifold or surface $M$, let $L=-\Delta_M$ be the positive Laplace--Beltrami operator and replace the dense Gibbs kernel by the intrinsic heat operator
\(\mathsf H_\epsilon\eqdef e^{-(\epsilon/4)L},\)
For histograms on the same discretized domain, the scaling steps then read
\eql{\label{eq-convolutional-sinkhorn}
\uD^{(\ell+1)}
=
\a\oslash\bigl(\mathsf H_\epsilon\vD^{(\ell)}\bigr),
\qquad
\vD^{(\ell+1)}
=
\b\oslash\bigl(\mathsf H_\epsilon^\top\uD^{(\ell+1)}\bigr),
}
where $\mathsf H_\epsilon$ denotes the discretized integral operator, including quadrature weights, and $\mathsf H_\epsilon^\top$ its discrete adjoint. In a symmetric mass-normalized discretization, the two operators coincide. Equivalently, it uses the effective cost
\(c_\epsilon(x,y)\eqdef-\epsilon\log h_{\epsilon/4}(x,y).\)
Varadhan's formula gives \(c_\epsilon(x,y)\longrightarrow d_M(x,y)^2 \qquad (\epsilon\downarrow0),\)
The heat operator can itself be applied through shifted-Laplacian solves:
\[
\mathsf H_\epsilon
=
\lim_{q\to\infty}
\left(\Id+\frac{\epsilon}{4q}L\right)^{-q}.
\]
\paragraph{Soft Hopf--Lax and Hopf--Cole.}
We use the normalized quadratic cost $c(x,y)=\frac12\norm{x-y}^2$, so that the Hopf--Lax operator applied to an initial datum $h$ is precisely
\[
(-h)^{\bar c}(x)=\inf_y\left\{h(y)+\frac12\norm{x-y}^2\right\},
\]
where the sign only reflects the convention of Definition~\ref{def-c-transform}. In the present entropic setting, the parameter of interest is instead the temperature $\epsilon$.

In the notation of Definition~\ref{def-continuous-soft-c-transform}, and using Lebesgue measure on $\RR^d$,
\begin{equation}\label{eq-soft-hopf-lax-heat}
(-h)^{\bar c,\epsilon}(x)
=
-\epsilon\log
\int
\exp\!\left(
-\frac{h(y)+\norm{x-y}^2/2}{\epsilon}
\right)\d y.
\end{equation}
If
\[
G_\epsilon(z)\eqdef (2\pi\epsilon)^{-d/2}\exp\!\left(-\frac{\norm{z}^2}{2\epsilon}\right)
\]
is the normalized Gaussian kernel with covariance $\epsilon\Id$, then~\eqref{eq-soft-hopf-lax-heat} is exactly
\begin{equation}\label{eq-soft-c-transform-gaussian-convolution}
(-h)^{\bar c,\epsilon}(x)
=
-\epsilon\log\Big(G_\epsilon\ast e^{-h/\epsilon}\Big)(x)
-\frac{\epsilon d}{2}\log(2\pi\epsilon).
\end{equation}
Thus a soft quadratic $c$-transform is a Gaussian convolution followed by a logarithm, up to an explicit additive constant independent of $x$. This is the bridge between soft-minimum operators, heat kernels and entropic transport potentials.

\begin{prop}[Soft quadratic $c$-transform and Legendre approximation]\label{prop-soft-legendre-convolution}
Let $f:\RR^d\to\RR\cup\{+\infty\}$ be such that the integrals below are positive and finite, and introduce the quadratic shift $\mathsf S f(y)\eqdef f(y)-\frac12\norm{y}^2$.
For $\epsilon>0$, define the soft conjugate by applying the soft $\bar c$-transform to the shifted function:
\begin{equation}\label{eq-soft-legendre-definition}
f^{*,\epsilon}(p)
\eqdef
\frac12\norm{p}^2-\big(-\mathsf S f\big)^{\bar c,\epsilon}(p).
\end{equation}
\eqllead{Then}{eq-soft-legendre-logsumexp}{
f^{*,\epsilon}(p)
=
\epsilon\log\int_{\RR^d}
\exp\!\left(\frac{\dotp{p}{y}-f(y)}{\epsilon}\right)\d y,
}
and equivalently
\begin{equation}\label{eq-soft-legendre-convolution}
f^{*,\epsilon}(p)
=
\frac12\norm{p}^2
+\epsilon\log\Big(G_\epsilon\ast e^{-\mathsf S f/\epsilon}\Big)(p)
+\frac{\epsilon d}{2}\log(2\pi\epsilon).
\end{equation}
If, for instance, $f$ is finite, continuous and superlinear, then $f^{*,\epsilon}(p)\to f^*(p)$ for every $p$ as $\epsilon\to0$.
\end{prop}
\begin{proof}
The hard identity follows by completing the square:
\[
\inf_y\left\{\mathsf S f(y)+\frac12\norm{p-y}^2\right\}
=
\frac12\norm{p}^2-
\sup_y\{\dotp{p}{y}-f(y)\}
=
\frac12\norm{p}^2-f^*(p).
\]
This explains the shift: it turns the Legendre--Fenchel transform into a quadratic $\bar c$-transform. Replacing the hard transform by its soft version gives~\eqref{eq-soft-legendre-definition}. Expanding the square in~\eqref{eq-soft-hopf-lax-heat} with the function $\mathsf S f$ gives
\[
\exp\!\left(-\frac{\mathsf S f(y)+\norm{p-y}^2/2}{\epsilon}\right)
=
\exp\!\left(-\frac{\norm{p}^2}{2\epsilon}\right)
\exp\!\left(\frac{\dotp{p}{y}-f(y)}{\epsilon}\right),
\]
which proves~\eqref{eq-soft-legendre-logsumexp}. Formula~\eqref{eq-soft-legendre-convolution} is the Gaussian-convolution identity~\eqref{eq-soft-c-transform-gaussian-convolution} applied to $\mathsf S f$. If $f$ is finite, continuous and superlinear, then $h_p(y)=\dotp{p}{y}-f(y)$ is continuous, attains its maximum, and has compact upper level sets. Every neighborhood of a maximizer has positive Lebesgue measure, while superlinearity controls the tails. The Laplace principle applied to~\eqref{eq-soft-legendre-logsumexp} therefore gives the pointwise supremum $\sup_y h_p(y)=f^*(p)$. Lower semicontinuity alone would not suffice: a downward spike of $f$ on a Lebesgue-null set changes $f^*$ but is invisible to the integral.
\end{proof}
With the same temperature normalization, the change of variables $u_s=e^{-\phi_s/\epsilon}$ converts the viscous Hamilton--Jacobi equation
\(\partial_s\phi_s+\frac12\norm{\nabla\phi_s}^2=\frac{\epsilon}{2}\Delta\phi_s,\)
Conversely, starting from a positive heat solution $u_s$ and setting $\phi_s=-\epsilon\log u_s$ gives a Hamilton--Jacobi solution, while $v_s=\nabla\phi_s=-\epsilon\nabla\log u_s$ solves the gradient viscous Burgers equation
\(\partial_s v_s+(v_s\cdot\nabla)v_s=\frac{\epsilon}{2}\Delta v_s.\)
\subsection{Other Convex Regularizers}
\label{sec-sinkhorn-other-regularizers}
\paragraph{\texorpdfstring{$\phi$-divergence regularization.}{phi-divergence regularization.}}
Let $\phi$ be an entropy function in the sense of Definition~\ref{def_entropy}, and recall the $\phi$-divergence $\Divergm_\phi$ from Definition~\ref{def_divergence}. For $\epsilon>0$, define the $\phi$-regularized transport value.

\begin{defn}[$\phi$-regularized optimal transport]\label{def-phi-regularized-ot}
\eql{\label{eq-phi-regularized-ot}
\MK_{\c,\phi}^{\epsilon}(\al,\be)
\eqdef
\umin{\pi\in\Couplings(\al,\be)}
\int_{\Xx\times\Yy} c(x,y)\d\pi(x,y)
+
\epsilon\Divergm_\phi(\pi|\al\otimes\be).
}
\end{defn}
\begin{prop}[Dual and primal--dual law for $\phi$-regularized OT]\label{prop-phi-regularized-ot-dual}
Set $L\eqdef\phi'_\infty$ and assume exact Fenchel--Rockafellar duality for the marginalization constraint. This holds, in particular, when $\Xx,\Yy$ are compact metric spaces, $\al,\be$ have full support, $c$ is Lipschitz, and $\phi$ is of Legendre type. Then
\eql{\label{eq-phi-regularized-ot-dual}
\MK_{\c,\phi}^{\epsilon}(\al,\be)
=
\usup{\substack{f\in\Cc(\Xx),\,g\in\Cc(\Yy)\\
f\oplus g-c\leq\epsilon L}}
\int f\d\al+\int g\d\be
-
\epsilon
\int
\phi^*\!\left(\frac{f(x)+g(y)-c(x,y)}{\epsilon}\right)
\d\al(x)\d\be(y).
}
The cap is vacuous when $L=+\infty$. If primal and dual optimizers exist, decompose
$\pi^\star=r^\star(\al\otimes\be)+(\pi^\star)^\perp$. Then
\eql{\label{eq-phi-regularized-ot-density-law}
\frac{f^\star(x)+g^\star(y)-c(x,y)}{\epsilon}
\in
\partial\phi(r^\star(x,y))
\qquad
\al\otimes\be\text{-a.e.}
}
If $L<+\infty$, the singular part $(\pi^\star)^\perp$ is concentrated on the recession contact set
\(\{(x,y):f^\star(x)+g^\star(y)-c(x,y)=\epsilon L\}.\)
If $L=+\infty$ and the objective is finite, then $(\pi^\star)^\perp=0$. In the smooth interior, the absolutely continuous density reads
\[
r^\star(x,y)
=
(\phi')^{-1}
\!\left(\frac{f^\star(x)+g^\star(y)-c(x,y)}{\epsilon}\right).
\]
\end{prop}
\begin{proof}
Introduce dual variables $(f,g)$ for the two marginal constraints. For fixed $(f,g)$, the minimization over $\pi$ gives
\[
\int f\d\al+\int g\d\be
+
\inf_{\pi\in\Mm_+(\Xx\times\Yy)}
\left\{
\int\big(c-(f\oplus g)\big)\d\pi
+
\epsilon\Divergm_\phi(\pi|\al\otimes\be)
\right\}.
\]
Using the Legendre formula~\eqref{eq-legendre} for the convex functional $\Divergm_\phi(\cdot|\al\otimes\be)$, the infimum equals
\[
-\epsilon
\int
\phi^*\!\left(\frac{f(x)+g(y)-c(x,y)}{\epsilon}\right)
\d\al(x)\d\be(y),
\]
provided $f\oplus g-c\leq\epsilon L$, and it is $-\infty$ otherwise. Equality in the Fenchel inequality yields both the subgradient inclusion
\(\frac{f^\star\oplus g^\star-c}{\epsilon} \in \partial\phi(r^\star),\)
for the absolutely continuous part and the recession equality
$\int[\epsilon L-(f^\star\oplus g^\star-c)]\d(\pi^\star)^\perp=0$ for the singular part. The latter proves the contact-set statement. If $L=+\infty$, finite divergence excludes singular mass by Definition~\ref{def_divergence}. Inversion of $\phi'$ gives the density law when $\phi$ is differentiable and the density lies in the interior of its domain. See for this finite-recession duality and the associated generalized Sinkhorn scheme.
\end{proof}
For the KL entropy $\phi(r)=r\log r-r+1$, one has $\phi^*(s)=e^s-1$, and~\eqref{eq-phi-regularized-ot-dual} reduces exactly to the continuous Sinkhorn dual~\eqref{eq-dual-sinkhorn-objective}. Other choices replace the exponential density law by another scalar transfer function:
\[
\begin{array}{lll}
\phi(r)=r\log r-r+1 &\Rightarrow& r^\star=e^s,\\[.25em]
\phi(r)=r-\log r-1 &\Rightarrow& r^\star=(1-s)^{-1}\quad (s<1),\\[.25em]
\phi(r)=\frac12(r-1)^2 &\Rightarrow& r^\star=(1+s)_+,
\end{array}
\qquad
s\eqdef\frac{f^\star\oplus g^\star-c}{\epsilon}.
\]
For Burg entropy the displayed formula describes only the absolutely continuous part. Since $\phi'_\infty=1$, an optimizer may additionally carry singular mass on the contact set $s=1$.

\paragraph{Bregman-divergence regularization.}
\begin{defn}[Measure Bregman divergence]\label{def-measure-bregman-divergence}
If $\Phi$ is a differentiable convex functional on a convex class of nonnegative measures and $\xi$ is a reference measure in its domain, the measure Bregman divergence generated by $\Phi$ is
\begin{equation}\label{eq-measure-bregman-divergence}
B_\Phi(\pi|\xi)
\eqdef
\Phi(\pi)-\Phi(\xi)
-
\int \delta\Phi(\xi)\d(\pi-\xi),
\end{equation}
where $\delta\Phi(\xi)$ is the first variation in the sense of Definition~\ref{def-first-variation}, and the formula is understood whenever the right-hand side is well-defined.
\end{defn}
\begin{defn}[Bregman-regularized optimal transport]\label{def-bregman-regularized-ot}
Fix $\xi=\al\otimes\be$. For $\epsilon>0$, define
\[
\MK_{\c,\Phi}^{\epsilon}(\al,\be)
\eqdef
\inf_{\pi\in\Couplings(\al,\be)}
\left\{\int c\,\d\pi+\epsilon B_\Phi(\pi|\xi)\right\}.
\]
\end{defn}
To type the function-space dual precisely, let $(\mathsf E,\mathsf E^*)$ be a separated locally convex dual pair in which $\mathsf E$ contains the signed measures on $\Xx\times\Yy$ under consideration and $\mathsf E^*$ contains the admissible test functions.

\begin{prop}[Dual and primal--dual law for Bregman-regularized OT]\label{prop-bregman-regularized-ot-dual}
Fix the marginals $\al,\be$ and set $\xi\eqdef\al\otimes\be$. Extend $\Phi$ by $+\infty$ outside the nonnegative cone of $\mathsf E$, and assume that it is proper, lower semicontinuous, convex, and G\^ateaux differentiable at $\xi$, with $q_\xi\eqdef\delta\Phi(\xi)\in\mathsf E^*$. For $u\in\mathsf E^*$, define
\[
\Phi^*(u)
\eqdef
\sup_{\pi\in\mathsf E}
\left\{
\langle u,\pi\rangle-\Phi(\pi)
\right\}.
\]
Assume $c\in\mathsf E^*$ and that Fenchel--Rockafellar duality is exact for the affine constraint $\mathsf A\pi=(\al,\be)$ in the chosen dual pair. In finite dimension, the concrete condition that a feasible $\pi_0$ belongs to $\operatorname{ri}(\dom\Phi)$ is sufficient. Then
\eql{\label{eq-bregman-regularized-ot-dual}
\MK_{\c,\Phi}^{\epsilon}(\al,\be)
=
\sup_{f\in\mathsf U_\Xx,\,g\in\mathsf U_\Yy}
\int f\d\al+\int g\d\be
-
\epsilon
\left[
\Phi^*\!\left(
q_\xi+\frac{f\oplus g-c}{\epsilon}
\right)
-
\Phi^*(q_\xi)
\right].
}
If $(f^\star,g^\star)$ and $\pi^\star$ are optimal and $\Phi$ is differentiable at $\pi^\star$, then
\eql{\label{eq-bregman-regularized-ot-density-law}
\delta\Phi(\pi^\star)
=
q_\xi+\frac{f^\star\oplus g^\star-c}{\epsilon}.
}
\end{prop}
\begin{proof}
Using~\eqref{eq-measure-bregman-divergence}, the Bregman-regularized objective can be written, up to a constant independent of $\pi$, as
\[
\epsilon\Phi(\pi)
+
\left\langle c-\epsilon q_\xi,\pi\right\rangle
-\epsilon\Phi(\xi)+\epsilon\langle q_\xi,\xi\rangle.
\]
Introduce dual potentials $(f,g)$ for the two marginal constraints. The inner minimization over nonnegative measures $\pi$ gives
\[
\inf_{\pi\in\mathsf E}
\left\{
\epsilon\Phi(\pi)
+
\left\langle c-(f\oplus g)-\epsilon q_\xi,\pi\right\rangle
\right\}
=
-\epsilon
\Phi^*\!\left(
q_\xi+\frac{f\oplus g-c}{\epsilon}
\right).
\]
Fenchel equality at $\xi$ gives \(-\Phi(\xi)+\langle q_\xi,\xi\rangle = \Phi^*(q_\xi),\)
which yields~\eqref{eq-bregman-regularized-ot-dual}. Equality in Fenchel's inequality gives the primal--dual relation~\eqref{eq-bregman-regularized-ot-density-law}.
\end{proof}
When $\Phi$ is separable with respect to a fixed dominating measure $\xi_0$, say $\Phi(\pi)=\int h(\d\pi/\d\xi_0)\d\xi_0$ with $\xi\ll\xi_0$, the Bregman optimality condition becomes
\[
h'\!\left(\frac{\d\pi^\star}{\d\xi_0}\right)
=
h'\!\left(\frac{\d\xi}{\d\xi_0}\right)
+
\frac{f^\star\oplus g^\star-c}{\epsilon}.
\]
This is an additive update in entropy coordinates. The density-ratio formulation instead uses the scalar law associated with $\phi$ relative to $\al\otimes\be$.

\begin{prop}[KL is the common Bregman and $\phi$ case]\label{prop-kl-only-bregman-phi}
Let $\omega$ be a finite reference measure on $\Xx$ with a measurable set $A$ satisfying $0<\omega(A)<\omega(\Xx)$. Work on probability measures $\al=p\omega$ and $\be=q\omega$ whose densities are bounded above and below away from $0$. Assume that $\Phi$ is twice G\^ateaux differentiable along bounded zero-mass density perturbations, and that $\phi\in C^2(0,+\infty)$ is convex with $\phi(1)=0$. If
\(B_\Phi(\al|\be)=\Divergm_\phi(\al|\be)\)
for all such $\al,\be$, then there exist $\kappa\geq0$ and $a\in\RR$ such that
\(\phi(t)=\kappa\,t\log t+a(t-1).\)
Hence the common divergence is $\kappa\KL(\al|\be)$, and $\Phi(p\omega)$ differs from $\kappa\int p\log p\,\d\omega$ by an affine functional on the positive probability simplex.
\end{prop}
\begin{proof}
Fix $\be=q\omega$ and perturb it by $\al_t=(q+t h)\omega$, where $h$ is bounded, $\int h\d\omega=0$, and $t$ is small enough that $q+t h>0$. Differentiating twice at $t=0$ gives
\(D^2\Phi(q)[h,h] = \phi''(1)\int \frac{h^2}{q}\d\omega.\)
Indeed the left-hand side is the second variation of the Bregman error, while
\[
\Divergm_\phi(\al_t|\be)
=
\int q\,\phi\left(1+t h/q\right)\d\omega
\]
has second derivative $\phi''(1)\int h^2/q\,\d\omega$. Setting $\kappa=\phi''(1)\geq0$, the functional $\Psi(p\omega)=\Phi(p\omega)-\kappa\int p\log p\d\omega$ has zero second variation along every zero-mass line segment in the positive simplex. It is therefore affine there, and $B_\Psi=0$. Thus $B_\Phi=\kappa\KL$.

Since $\Divergm_\phi=\kappa\KL$, the function $r(t)=\phi(t)-\kappa t\log t$ generates the zero $\phi$-divergence. Testing on densities for which the ratio $p/q$ takes two values $x<1<y$, with weights chosen so that the mean ratio is $1$, gives $(y-1)r(x)+(1-x)r(y)=0$. Hence $r(t)/(t-1)$ is constant on $(0,+\infty)\setminus\{1\}$, so $r(t)=a(t-1)$. This proves the claim.
\end{proof}
Thus, except for KL, the two generalizations lead to different duals and different algorithms. Bregman regularization by $B_\Phi(\pi|\xi)$ keeps the projection geometry of Section~\ref{sec-convergence-init}: linear costs tilt the reference in dual coordinates and alternating marginal updates are Bregman projections.

\paragraph{\texorpdfstring{Generalized soft $c$-transforms and alternate dual maximization method.}{Generalized soft c-transforms and alternate dual maximization method.}}
The two constrained $\phi$-soft $c$-transforms are
\eql{\label{eq-phi-soft-c-transform}
\begin{aligned}
g^{\bar c,\epsilon,\phi}(x)
&\in
\uargmin{\substack{u\in\RR\\u+g(y)-c(x,y)\leq\epsilon L\ \forall y}}
\left\{
\epsilon\int_{\Yy}
\phi^*\!\left(\frac{u+g(y)-c(x,y)}{\epsilon}\right)\d\be(y)-u
\right\},\\
f^{c,\epsilon,\phi}(y)
&\in
\uargmin{\substack{v\in\RR\\f(x)+v-c(x,y)\leq\epsilon L\ \forall x}}
\left\{
\epsilon\int_{\Xx}
\phi^*\!\left(\frac{f(x)+v-c(x,y)}{\epsilon}\right)\d\al(x)-v
\right\}.
\end{aligned}
}
When $\phi^*$ is differentiable and the first scalar minimizer lies in the interior of the cap, it is characterized by
\[
\int_{\Yy}
(\phi^*)'\!\left(\frac{u+g(y)-c(x,y)}{\epsilon}\right)\d\be(y)=1,
\]
and the second satisfies the symmetric equation. Thus an interior update normalizes one absolutely continuous conditional density.

The Bregman dual~\eqref{eq-bregman-regularized-ot-dual} has analogous block transforms, but the global conjugate $\Phi^*$ prevents a pointwise definition in general. Under the dual-pair and attainment assumptions of Proposition~\ref{prop-bregman-regularized-ot-dual}, with $\xi=\al\otimes\be$ and $q_\xi=\delta\Phi(\xi)$, define
\eql{\label{eq-bregman-soft-c-transform}
\begin{aligned}
g^{\bar c,\epsilon,\Phi}
&\in\uargmin{u\in\mathsf U_\Xx}
\left\{
\epsilon\Phi^*\!\left(q_\xi+\frac{u\oplus g-c}{\epsilon}\right)
-\int u\d\al
\right\},\\
f^{c,\epsilon,\Phi}
&\in\uargmin{v\in\mathsf U_\Yy}
\left\{
\epsilon\Phi^*\!\left(q_\xi+\frac{f\oplus v-c}{\epsilon}\right)
-\int v\d\be
\right\}.
\end{aligned}
}
Alternate dual maximization therefore reads
\eql{\label{eq-generalized-soft-c-alternate-maximization}
\begin{aligned}
f^{(\ell+1)}&=\big(g^{(\ell)}\big)^{\bar c,\epsilon,\phi},
&g^{(\ell+1)}&=\big(f^{(\ell+1)}\big)^{c,\epsilon,\phi},\\
f^{(\ell+1)}&=\big(g^{(\ell)}\big)^{\bar c,\epsilon,\Phi},
&g^{(\ell+1)}&=\big(f^{(\ell+1)}\big)^{c,\epsilon,\Phi},
\end{aligned}
}
In the discrete setting, the quadratic choices $\phi(r)=\frac12(r-1)^2$ and $\Phi(\P)=\frac12\norm{\P}_{\mathrm F}^2$ give, respectively,
\eql{\label{eq-quadratic-regularized-density-laws}
\P^\star_{i,j}
=
a_i b_j\left(1+\frac{\fD_i^\star+\gD_j^\star-\C_{i,j}}{\epsilon}\right)_+,
\qquad
\P^\star_{i,j}
=
\left(a_i b_j+\frac{\fD_i^\star+\gD_j^\star-\C_{i,j}}{\epsilon}\right)_+.
}
\[
\sum_j b_j\left(1+\frac{u+\gD_j-\C_{i,j}}{\epsilon}\right)_+=1.
\]
This left-hand side is continuous, nondecreasing, and piecewise affine. Sorting its breakpoints computes the update in $O(m\log m)$ operations for a row with $m$ entries; equivalently, it is a weighted simplex projection.

\subsection{Sinkhorn Divergences}
\label{sec-sinkhorn-div}
\paragraph{Entropic bias.}
Consequently, a minimizer
\(\be_\epsilon\in\uargmin{\be}\MK_\c^\epsilon(\al,\be)\)
\begin{prop}[Large-temperature entropic bias]
Let $\Xx,\Yy$ be compact and assume that $c$ is continuous. Then $\MK_\c^\epsilon(\al,\be) \rightarrow \iint c(x,y)\d\al(x)\d\be(y)$ as $\epsilon \rightarrow +\infty$.
\end{prop}
\begin{proof}
Let $(f_\epsilon,g_\epsilon)$ be optimal dual potentials, normalized by $\int g_\epsilon\d\be=0$. The soft $c$-transform equation gives
\[
f_\epsilon(x)
=
-\epsilon\log\int
\exp\!\left(\frac{g_\epsilon(y)-c(x,y)}{\epsilon}\right)\d\be(y).
\]
For bounded $c$, the oscillations of normalized entropic potentials are bounded uniformly in $\epsilon$ by the oscillation of $c$. Hence the log-sum-exp expansion is uniform:
\(f_\epsilon(x) = -\int (g_\epsilon(y)-c(x,y))\d\be(y)+O(\epsilon^{-1}) = \int c(x,y)\d\be(y)+O(\epsilon^{-1}).\)
At optimality the exponential penalty in the dual integrates to zero, so
\(\MK_\c^\epsilon(\al,\be) = \int f_\epsilon\d\al+\int g_\epsilon\d\be = \iint c(x,y)\d\al(x)\d\be(y)+O(\epsilon^{-1}).\)
This proves the limit.
\end{proof}
Thus, at large $\epsilon$, the raw entropic cost approaches the bilinear cross-energy $(\al,\be)\mapsto\iint c\,\d\al\d\be$, not a distance.

\paragraph{Sinkhorn divergences.}
\begin{defn}[Sinkhorn divergence]\label{def-sinkhorn-divergence}
Let $c$ be a symmetric cost on a common state space. For $\epsilon>0$, the debiased Sinkhorn divergence associated with the entropic OT value $\MK_\c^\epsilon$ is
\eql{\label{eq-sinkhorn-divergence}
\bar\MK_\c^\epsilon(\al,\be) \eqdef
\MK_\c^\epsilon(\al,\be) - \frac{1}{2} \MK_\c^\epsilon(\al,\al) - \frac{1}{2}\MK_\c^\epsilon(\be,\be).
}
\end{defn}
\begin{lem}[Entropic dual cost at optimum]
Let $(f_{\al,\be},g_{\al,\be})$ be optimal dual potentials, normalized arbitrarily. Then
\eql{\label{eq-formula-cost-dual}
\MK_\c^\epsilon(\al,\be) = \dotp{f_{\al,\be}}{\al} + \dotp{g_{\al,\be}}{\be}.
}
\end{lem}
\begin{proof}
We first notice that, at optimality, $f_{\al,\be} = -\epsilon\log\int_\Yy e^{\frac{g_{\al,\be}(y)-c(x,y)}{\epsilon}} \d\be(y)$. After exponentiation, this equivalently reads
\eq{
1 = \int_\Yy e^{\frac{f_{\al,\be}(x) + g_{\al,\be}(y)-c(x,y)}{\epsilon}} \d\be(y)
\qarrq
\int_{\Xx\times\Yy} \pa{ e^{\frac{f_{\al,\be} \oplus g_{\al,\be}-c}{\epsilon}} -1 } \d(\al \otimes \be) = 0.
}
Substituting this identity in~\eqref{eq-dual-sinkh-cont} gives the result.
\end{proof}
\begin{prop}[Asymptotics of Sinkhorn divergences]\label{prop-sinkhorn-divergence-asymptotics}
Let $\Xx$ be a compact metric space, let $\al,\be\in\Pp(\Xx)$, and assume that $c:\Xx\times\Xx\to\RR_+$ is continuous and symmetric, with $c(x,x)=0$.
Then $\bar\MK_\c^\epsilon(\al,\be) \rightarrow \MK_\c(\al,\be)$ when $\epsilon\rightarrow 0$ and
\eq{
\bar\MK_\c^\epsilon(\al,\be) \rightarrow
\frac{1}{2}\int -c \d(\al-\be) \otimes \d(\al-\be)
\qwhenq
\epsilon \rightarrow +\infty.
}
\end{prop}
\begin{proof}
The discrete convergence result above gives the finite-dimensional analogue; we now pass to continuous measures.
\textbf{Case $\epsilon \rightarrow 0$.}
The first limit is the $\Gamma$-convergence of entropic OT to the Kantorovich problem. Non-negativity of the relative entropy and continuity of $c$ give the liminf inequality. For the matching limsup, one approximates an optimal coupling by finite-entropy couplings and chooses the approximation scale so that its entropy multiplied by $\epsilon$ tends to zero. Since $c(x,x)=0$, the two self-costs in the debiased expression converge to zero, while the cross term converges to $\MK_\c(\al,\be)$.

\textbf{Case $\epsilon \rightarrow +\infty$.} We denote by $(f_\epsilon,g_\epsilon)$ optimal dual potentials. After normalizing them and using boundedness of $c$, their oscillations stay uniformly bounded, so the following expansion is uniform. The optimality condition on $f_\epsilon$ (equivalently the Sinkhorn fixed point on $f_\epsilon$) reads
\begin{align*}
f_\epsilon &= -\epsilon \log \int \exp\pa{ \frac{g_\epsilon(y) - c(\cdot,y)}{\epsilon} } \d \be(y)
=	-\epsilon \log \int \pa{ 1 + \frac{g_\epsilon(y) - c(\cdot,y)}{\epsilon} + o(1/\epsilon) } \d \be(y) \\
&=	-\epsilon \log\pa{1+\frac{1}{\epsilon}\int (g_\epsilon(y)-c(\cdot,y))\d\be(y)+o(1/\epsilon)}
= - \int g_\epsilon \d\be + \int c(\cdot,y) \d \be(y) + o(1).
\end{align*}
Plugging this relation in the dual expression~\eqref{eq-formula-cost-dual}
\(\MK_\c^\epsilon(\al,\be) = \int f_\epsilon \d\al + \int g_\epsilon \d \be = \iint c(x,y) \d\al(x)\d\be(y) + o(1).\)
Applying this expansion to $(\al,\be)$, $(\al,\al)$ and $(\be,\be)$ gives
\[
\bar\MK_\c^\epsilon(\al,\be)
\to
\int c\,\d\al\otimes\d\be
-\frac12\int c\,\d\al\otimes\d\al
-\frac12\int c\,\d\be\otimes\d\be
=
-\frac12\int c\,\d(\al-\be)\otimes\d(\al-\be).
\]
\end{proof}
\begin{prop}[Positivity of Sinkhorn divergences]\label{prop-sinkhorn-positive}
Assume that the entropic dual admits optimal potentials and that the symmetric kernel $k_\epsilon(x,y)=e^{-c(x,y)/\epsilon}$ is positive semidefinite in the sense of Definition~\ref{def-positive-kernels}. Then $\bar\MK_\c^\epsilon(\al,\be)\geq0$.
If, in addition, the common state space is compact, $c$ is continuous, and $k_\epsilon$ is universal in the sense of Definition~\ref{def-universal-kernel}, then
\(\bar\MK_\c^\epsilon(\al,\be)=0 \quad\Longleftrightarrow\quad \al=\be,\)
and $\bar\MK_\c^\epsilon(\al_n,\al)\to0$ is equivalent to $\al_n\rightharpoonup\al$.
\end{prop}
\begin{proof}
In the following, we denote by $(f_{\al,\be},g_{\al,\be})$ optimal dual potentials for the
dual Schr\"odinger problem between $\al$ and $\be$.
We denote by $f_{\al,\al}=g_{\al,\al}$ (one can assume they are equal by symmetry) the solution for the problem between $\al$ and itself.

Using the suboptimal function $(f_{\al,\al},g_{\be,\be})$ in the dual maximization problem, and using relation~\eqref{eq-formula-cost-dual} for the simplified expression of the dual cost, one obtains
\eq{
\MK_\c^\epsilon(\al,\be) \geq \dotp{f_{\al,\al}}{\al} + \dotp{g_{\be,\be}}{\be}
- \epsilon \dotp{ e^{\frac{f_{\al,\al} \oplus g_{\be,\be}-c}{\epsilon}} -1 }{\al \otimes \be}
}
Moreover $\dotp{f_{\al,\al}}{\al} = \frac{1}{2}\MK_\c^\epsilon(\al,\al)$, and similarly for $\be$, so the previous inequality equivalently reads
\eq{
\frac{1}{\epsilon} \bar \MK_\c^\epsilon(\al,\be)
\geq 1 - \dotp{ e^{\frac{f_{\al,\al} \oplus g_{\be,\be}-c}{\epsilon}} }{\al \otimes \be}
= 1 - \dotp{\tilde\al}{\tilde\be}_{k_\epsilon}
}
where $\tilde\al=e^{f_{\al,\al}/\epsilon} \al$, $\tilde\be=e^{f_{\be,\be}/\epsilon} \be$
and we introduced the positive-semidefinite kernel pairing
$\dotp{\tilde\al}{\tilde\be}_{k_\epsilon} \eqdef \int k_\epsilon(x,y) \d\tilde\al(x) \d\tilde\be(y)$.
The self Sinkhorn fixed point equation, once exponentiated, reads pointwise
$e^{f_{\al,\al}(x)/\epsilon}\int k_\epsilon(x,y)\d\tilde\al(y)=1$ for $\al$-a.e.\ $x$,
and hence
\eq{
\norm{\tilde \al}_{k_\epsilon}^2
= \dotp{\tilde\al}{\tilde\al}_{k_\epsilon}
= \int e^{f_{\al,\al}(x)/\epsilon}
\left(\int k_\epsilon(x,y)\d\tilde\al(y)\right)\d\al(x) = 1
}
and similarly $\norm{\tilde \be}_{k_\epsilon}^2=1$. Therefore, by Cauchy--Schwarz, one has
$1 - \dotp{\tilde\al}{\tilde\be}_{k_\epsilon} \geq 0$.

Suppose now that $k_\epsilon$ is universal and $\bar\MK_\c^\epsilon(\al,\be)=0$. Both inequalities above are then equalities. Equality in Cauchy--Schwarz, together with the unit norms, implies that $\tilde\al$ and $\tilde\be$ have the same kernel mean embedding; universality gives $\tilde\al=\tilde\be\eqdef\tilde\mu$. The two self-consistency equations yield
\(\d\al=(k_\epsilon\tilde\mu)\,\d\tilde\mu=\d\be, \qquad (k_\epsilon\tilde\mu)(x)\eqdef\int k_\epsilon(x,y)\d\tilde\mu(y),\)
so $\al=\be$. The converse follows directly from the definition. Finally, on a compact space, continuity of the entropic value under weak convergence and compactness of $\Pp(\Xx)$ turn this definiteness into the stated equivalence of convergences; see.
\end{proof}
\subsection{Complex \texorpdfstring{$\epsilon$}{epsilon}}
\label{sec-complex-epsilon}
\paragraph{Measure fixed point.}
Let $\al\in\Pp(\Xx)$ and $\be\in\Pp(\Yy)$. For $\epsilon\in\CC\setminus\{0\}$, reuse the Gibbs kernel and integral operators of~\eqref{eq-continuous-sinkhorn-operators} whenever the defining integrals exist.

\paragraph{Discrete histograms.}
For discrete histograms, evaluate the Gibbs kernel entrywise at $(x_i,y_j)$. The scaling factorization~\eqref{eq-scaling-form} and Sinkhorn updates~\eqref{eq-sinkhorn} are then used verbatim over $\CC$, with the same marginal constraints and multiplicative gauge, whenever $\epsilon\neq0$ and all divisions are defined.

\begin{thm}[Local holomorphic continuation of Sinkhorn scalings]\label{thm-carlier-complex-sinkhorn}
Fix positive histograms $\a^0\in\simplex_n$, $\b^0\in\simplex_m$, a finite real cost matrix $\C^0$, and a real temperature $\epsilon_0>0$. Choose positive scalings $(\uD^0,\vD^0)$ of the corresponding Sinkhorn coupling. Then there are complex neighborhoods of $(\epsilon_0,\a^0,\b^0,\C^0)$, inside the affine constraint $\sum_i \a_i=\sum_j \b_j$, and a unique holomorphic map
\((\epsilon,\a,\b,\C)\mapsto(\uD_\epsilon,\vD_\epsilon)\in\CC^n\times\CC^m\)
satisfying the linear gauge $\sum_i\a_i^0u_{\epsilon,i}=\sum_i\a_i^0u_i^0$, such that
$\P_\epsilon=\diag(\uD_\epsilon)\K_\epsilon(\C)\diag(\vD_\epsilon)$ has marginals $(\a,\b)$. In particular, the gauge-fixed scalings and $\P_\epsilon$ are holomorphic in all four arguments near the base point.
\end{thm}
\begin{proof}
Apply the holomorphic implicit-function theorem to the $n$ row residuals, the first $m-1$ column residuals, and the gauge residual $G(\uD)=\sum_i\a_i^0(u_i-u_i^0)$. There are $n+m$ equations for the $n+m$ scaling coordinates.

Let $(\delta\uD,\delta\vD)$ belong to its kernel and write
\(r_i\eqdef\frac{\delta u_i}{u_i^0}, \qquad s_j\eqdef\frac{\delta v_j}{v_j^0}.\)
Since $\P^0=\diag(\uD^0)\K_{\epsilon_0}(\C^0)\diag(\vD^0)$, the induced perturbation is $\delta\P_{i,j}=\P^0_{i,j}(r_i+s_j)$.
The linearized row sums vanish. The first $m-1$ linearized column sums vanish by definition of the residuals, and the last one vanishes because the total row sum equals the total column sum. Therefore
\(0 = \sum_i \overline{r_i}\sum_j\delta \P_{i,j} + \sum_j \overline{s_j}\sum_i\delta \P_{i,j} = \sum_{i,j}\P^0_{i,j}\abs{r_i+s_j}^2.\)
Strict positivity of $\P^0$ gives $r_i=\theta$ and $s_j=-\theta$ for some $\theta\in\CC$. The linearized gauge is $0=\sum_i\a_i^0\delta u_i=\theta\sum_i\a_i^0u_i^0$; its last factor is positive, hence $\theta=0$. The derivative is injective and therefore invertible. The implicit-function theorem now gives the asserted holomorphic scalings, and their matrix product gives the holomorphic coupling.
\end{proof}

\section{Entropic Regularization: Convergence}
\label{sec-sinkhorn-advanced}
\label{sec-entropic-convergence}
\label{sec-convergence-dual}
The chapter revisits Sinkhorn convergence through several complementary lenses. Bregman projections explain the alternating-projection geometry, while Fortet's order argument gives qualitative fixed-point convergence.

\subsection{Sinkhorn Convergence: Bregman Point of View}
\label{sec-convergence-init}
\paragraph{\texorpdfstring{Alternating $\KL$ projections.}{Alternating KL projections.}}
\begin{defn}[Bregman divergence]\label{def-bregman-divergence}
Let $\Omega\subset\RR^{n\times m}$ be convex with nonempty interior, and let $\Phi:\RR^{n\times m}\to\RR\cup\{+\infty\}$ be a proper lower-semicontinuous Legendre function with $\operatorname{dom}(\Phi)=\Omega$. Thus $\Phi=+\infty$ outside $\Omega$, while $\Phi$ is differentiable and strictly convex on $\operatorname{int}(\Omega)$. For $\P\in\Omega$ and $\Q\in\operatorname{int}(\Omega)$, the Bregman divergence generated by $\Phi$ is
\(B_\Phi(\P|\Q)\eqdef \Phi(\P)-\Phi(\Q)-\dotp{\nabla\Phi(\Q)}{\P-\Q}.\)
\end{defn}
\begin{defn}[Bregman projection]\label{def-bregman-projection}
For $\Q\in\operatorname{int}(\Omega)$ and a nonempty convex set $\Cc$, define
\(\Proj_{\Cc}^{B_\Phi}(\Q) \eqdef \argmin_{\P\in\Cc\cap\Omega} B_\Phi(\P|\Q).\)
This set is empty if $\Cc\cap\Omega$ is empty or if the infimum is not attained, and it is set-valued when the minimizer is not unique. We identify it with its unique element only when existence and uniqueness are ensured.
\end{defn}
\paragraph{Linear tilts and Gibbs references.}
\begin{prop}[Linear tilts of Bregman penalties]\label{prop-bregman-linear-tilt}
Let $\Phi$ be a Legendre function, and let $B_\Phi$ be its Bregman divergence. Fix $\Q,\Q^\C\in\operatorname{int}(\operatorname{dom}\Phi)$ such that
\(\nabla\Phi(\Q^\C)=\nabla\Phi(\Q)-\C/\epsilon.\)
Then, for all $\P\in\operatorname{dom}\Phi$, \(\dotp{\P}{\C}+\epsilon B_\Phi(\P|\Q) = \epsilon B_\Phi(\P|\Q^\C)+\text{\upshape cst},\)
where the constant does not depend on $\P$.
\end{prop}
\begin{proof}
Subtract the two Bregman divergences: \(B_\Phi(\P|\Q^\C)-B_\Phi(\P|\Q) = \dotp{\nabla\Phi(\Q)-\nabla\Phi(\Q^\C)}{\P} +\text{cst}.\)
Using $\nabla\Phi(\Q)-\nabla\Phi(\Q^\C)=\C/\epsilon$ and multiplying by $\epsilon$ gives the claim.
\end{proof}
For the negative entropy $\Phi(\P)=\sum_{i,j}\P_{i,j}\log\P_{i,j}$, one has $B_\Phi=\KLD$. After removing zero-mass rows and columns, as in the proof of Proposition~\ref{prop-regularized-primal}, the histograms are positive and $\Q=\a\otimes\b$ lies in the interior of the entropy domain. Proposition~\ref{prop-bregman-linear-tilt} then gives the tilted reference
\(\K_{\a,\b}^\epsilon \eqdef (\a\otimes\b)\odot e^{-\C/\epsilon}.\)
Thus
\(\dotp{\P}{\C}+\epsilon\KLD(\P|\a\otimes\b) = \epsilon\KLD(\P|\K_{\a,\b}^\epsilon)+\text{\upshape cst}.\)
Thus the unique solution $\P_\epsilon$ of~\eqref{eq-regularized-discr} is the KL projection of the tilted Gibbs reference onto $\CouplingsD(\a,\b)$:
\eql{\label{eq-kl-proj}
\P_\epsilon = \Proj_{\CouplingsD(\a,\b)}^\KLD(\K_{\a,\b}^\epsilon) \eqdef \uargmin{\P \in \CouplingsD(\a,\b)} \KLD(\P|\K_{\a,\b}^\epsilon).
}
\paragraph{Cyclic projection convergence.}
Given two closed convex constraint sets $\Cc_1$ and $\Cc_2$, the cyclic method projects first onto $\Cc_1$, then onto $\Cc_2$, and repeats. Full iterates satisfy the second constraint and half-iterates satisfy the first; both constraints are attained together only in the limit.

\begin{prop}[Convergence of cyclic Bregman projections]\label{prop-cyclic-kl-affine}
Let $\Phi$ be a Legendre generator on a finite-dimensional convex domain, and let $\Cc_1,\Cc_2$ be closed convex sets whose intersection meets the interior of the domain. Generate $(\P^{(\ell-1/2)},\P^{(\ell)})$ by the cyclic Bregman projection iteration from an interior point $\P^{(0)}$. Assume that the projections are well-defined and that the full half-step sequence remains in a compact subset of the interior of the domain. Then there exists $\bar\P\in\Cc_1\cap\Cc_2$ such that $\P^{(\ell)}\to\bar\P$ and $\P^{(\ell-1/2)}\to\bar\P$.
If, in addition, $\Cc_1$ and $\Cc_2$ are affine, then
\(\bar\P = \Proj_{\Cc_1\cap\Cc_2}^{B_\Phi}(\P^{(0)}).\)
\end{prop}
\begin{proof}
For $\Q\in\operatorname{dom}\Phi$ and $\P,\P^+\in\operatorname{int}(\operatorname{dom}\Phi)$, the Bregman three-point formula reads
\(B_\Phi(\Q|\P) = B_\Phi(\Q|\P^+) + B_\Phi(\P^+|\P) + \dotp{\nabla\Phi(\P^+)-\nabla\Phi(\P)}{\Q-\P^+}.\)
If $\P^+=\Proj_\Cc^{B_\Phi}(\P)$ and $\Cc$ is closed and convex, first-order optimality gives \(\dotp{\nabla\Phi(\P^+)-\nabla\Phi(\P)}{\Q-\P^+}\geq0 \qquad \forall \Q\in\Cc\cap\operatorname{dom}\Phi.\)
\(B_\Phi(\Q|\P) \geq B_\Phi(\Q|\P^+)+B_\Phi(\P^+|\P) \qquad\forall \Q\in\Cc\cap\operatorname{dom}\Phi.\)

Let $(Z^{(r)})_r$ be the half-step sequence, so that $Z^{(2\ell)}=\P^{(\ell)}$ and $Z^{(2\ell+1)}=\P^{(\ell+1/2)}$. Fix $\Q\in\Cc_1\cap\Cc_2\cap\operatorname{int}(\operatorname{dom}\Phi)$, which exists by assumption. Applying the inequality at each half-step gives
\(B_\Phi(\Q|Z^{(r)})-B_\Phi(\Q|Z^{(r+1)}) \geq B_\Phi(Z^{(r+1)}|Z^{(r)})\geq0.\)
Thus $B_\Phi(\Q|Z^{(r)})$ decreases and $\sum_r B_\Phi(Z^{(r+1)}|Z^{(r)})<+\infty$. The projection drops tend to zero; compactness and strict convexity imply $\norm{Z^{(r+1)}-Z^{(r)}}\to0$. Every cluster point of either subsequence consequently belongs to both closed sets. Let $\bar\P\in\Cc_1\cap\Cc_2$ be one such cluster point. The sequence $B_\Phi(\bar\P|Z^{(r)})$ is nonincreasing and tends to zero along a convergent subsequence, hence tends to zero altogether. Compactness and strict convexity then give $Z^{(r)}\to\bar\P$, proving the first claim.

Suppose now that $\Cc_1$ and $\Cc_2$ are affine. The first-order inequality above is then an equality because both signs of every feasible direction are admissible. Moreover, each dual displacement
\(\nabla\Phi(Z^{(r+1)})-\nabla\Phi(Z^{(r)})\)
belongs to the normal space of the affine set used at that half-step. Telescoping and passing to the limit gives \(\nabla\Phi(\bar \P)-\nabla\Phi(\P^{(0)}) \in N_{\Cc_1}+N_{\Cc_2} = N_{\Cc_1\cap\Cc_2},\)
where the last equality uses affinity and the nonempty intersection. This is precisely the first-order optimality condition for minimizing $\Q\mapsto B_\Phi(\Q|\P^{(0)})$ over $\Cc_1\cap\Cc_2$, proving the displayed projection identity.
\end{proof}
\paragraph{Row and column scalings.}
Denote the two affine marginal sets by
\eq{\label{eq-affine-marginal-sets}
\Cc^1_\a \eqdef \enscond{\P\in\RR^{n\times m}}{\P\ones_m=\a}
\qandq
\Cc^2_\b \eqdef \enscond{\P\in\RR^{n\times m}}{\transp{\P}\ones_n=\b}.
}
Positivity is a separate constraint, so
\(\CouplingsD(\a,\b) = \Cc^1_\a\cap\Cc^2_\b\cap\RR_+^{n\times m}.\)
For negative entropy, however, the effective domain of the Legendre generator is $\Omega=\RR_+^{n\times m}$ and $\Phi=+\infty$ outside $\Omega$. Thus the two half-steps of the cyclic Bregman projection iteration specialize to
\eql{\label{eq-kl-sinkh-proj}
\itt{\P} \eqdef \Proj_{\Cc^1_\a}^{\KLD}(\it{\P})
\qandq
\ittt{\P} \eqdef \Proj_{\Cc^2_\b}^{\KLD}(\itt{\P}).
}
\begin{prop}[KL projections are scalings]
Let $\P\geq0$, set $r_i\eqdef(\P\ones_m)_i$ and $s_j\eqdef(\P^\top\ones_n)_j$, and assume the support-compatibility conditions
\(
a_i>0\Rightarrow r_i>0
\)
and
\(
b_j>0\Rightarrow s_j>0
\).
Then the finite KL projections are unique and satisfy
\[
\left[\Proj_{\Cc^1_\a}^{\KLD}(\P)\right]_{i,j}
=
\begin{cases}
a_i\P_{i,j}/r_i,&r_i>0,\\
0,&r_i=0,
\end{cases}
\qquad
\left[\Proj_{\Cc^2_\b}^{\KLD}(\P)\right]_{i,j}
=
\begin{cases}
b_j\P_{i,j}/s_j,&s_j>0,\\
0,&s_j=0.
\end{cases}
\]
If a compatibility condition fails, the corresponding affine set contains no matrix at finite KL divergence from $\P$. In particular, when $\P>0$ and $\a,\b>0$, these formulas reduce to
\(\Proj_{\Cc^1_\a}^{\KLD}(\P) = \diag\pa{\frac{\a}{\P\ones_m}}\P, \qquad \Proj_{\Cc^2_\b}^{\KLD}(\P) = \P\diag\pa{\frac{\b}{\P^\top\ones_n}}.\)
\end{prop}
\begin{proof}
Consider one row or column, with reference vector $q\geq0$, reference mass $Q\eqdef\sum_iq_i$, and prescribed mass $m\geq0$. If $Q>0$ and $m>0$, finite KL forces $p_i=0$ wherever $q_i=0$. On the positive support of $q$, strict convexity and the Lagrange equation give $p=u q$; the mass constraint yields $u=m/Q$. If $Q>0$ and $m=0$, nonnegativity forces the unique solution $p=0$, which is the same formula with $u=0$. If $Q=m=0$, again $p=0$ is the unique feasible vector. Finally, if $Q=0<m$, every feasible $p$ has infinite KL divergence from $q$, so no finite projection exists. Applying these four cases independently to every row or column proves the claim.
\end{proof}
Since only $\vD^{(0)}$ is initialized, the physical plan sequence starts after the first row update and is
\eql{\label{eq-sink-matrix}
\P^{(2\ell+1)} \eqdef \diag(\itt{\uD})\K\diag(\it{\vD}),
\qquad
\P^{(2\ell+2)} \eqdef \diag(\itt{\uD})\K\diag(\itt{\vD}),
\qquad \ell\geq0.
}
Thus odd iterates have source marginal $\a$, whereas even iterates have target marginal $\b$. In practice, however, one should prefer using~\eqref{eq-sinkhorn}, which only requires manipulating scaling vectors and multiplying by a Gibbs kernel, and can often be accelerated when the kernel has separable, sparse, low-rank or geometric structure.

\paragraph{Other divergences.}
Positivity must then be imposed explicitly in each marginal constraint by defining
\eql{\label{eq-positive-marginal-sets}
\Cc^1_{\a,+}
\eqdef
\Cc^1_\a\cap\RR_+^{n\times m}
\qandq
\Cc^2_{\b,+}
\eqdef
\Cc^2_\b\cap\RR_+^{n\times m},
}
Projection onto $\Cc^1_{\a,+}$ is rowwise projection onto a simplex: \(\big[\Proj_{\Cc^1_{\a,+}}^{B_\Phi}(\Q)\big]_{i,j} = (\Q_{i,j}-\tau_i)_+, \qquad \sum_j(\Q_{i,j}-\tau_i)_+=a_i,\)
For the quadratic generator and the product reference $\xi=\a\otimes\b$, the Bregman penalty is \(B_\Phi(\P|\xi) = \frac12\sum_{i,j}(\P_{i,j}-a_i b_j)^2.\)
This is closely related, but not identical, to the quadratic $\phi$-divergence generated by $\phi(r)=\frac12(r-1)^2$, which for positive $a_i,b_j$ is
\(\Divergm_\phi(\P|\a\otimes\b) = \frac12\sum_{i,j}\frac{(\P_{i,j}-a_i b_j)^2}{a_i b_j}.\)
\subsection{Sinkhorn Convergence: Monotone Point of View}\label{sec-sinkhorn-monotone}
\paragraph{Variation seminorm and topical maps.}
\begin{defn}[Variation seminorm]\label{def-variation-seminorm}
For a bounded real-valued function $h$, its variation seminorm is \(\norm{h}_V\eqdef \sup h-\inf h.\)
It vanishes exactly on constant functions and therefore induces a norm after quotienting by additive constants.
\end{defn}
\begin{defn}[Topical map]\label{def-topical-map}
Let $E$ be an ordered vector space of functions containing the constants. A map $\mathcal T:E\to E$ is \emph{topical} if
\(f\leq g\Longrightarrow\mathcal T(f)\leq\mathcal T(g), \qquad \mathcal T(f+s)=\mathcal T(f)+s \quad(s\in\RR).\)
\end{defn}
\begin{prop}[Topical maps are variation-nonexpansive]\label{prop-topical-variation-nonexpansive}
Let $E$ be a vector space of bounded real-valued functions, ordered pointwise. Every topical map $\mathcal T:E\to E$ satisfies \(\norm{\mathcal T(f)-\mathcal T(g)}_V \leq \norm{f-g}_V.\)
The same conclusion holds if $\mathcal T$ is order reversing and additively anti-homogeneous, meaning $\mathcal T(f+s)=\mathcal T(f)-s$.
\end{prop}
\begin{proof}
Set $a=\inf(f-g)$ and $b=\sup(f-g)$, so that $g+a\leq f\leq g+b$. If $\mathcal T$ is topical, then
\(\mathcal T(g)+a\leq\mathcal T(f)\leq\mathcal T(g)+b.\)
Hence the oscillation of $\mathcal T(f)-\mathcal T(g)$ is at most $b-a=\norm{f-g}_V$. In the order-reversing, additively anti-homogeneous case, the corresponding bounds are $\mathcal T(g)-b\leq\mathcal T(f)\leq\mathcal T(g)-a$ and give the same conclusion.
\end{proof}
\paragraph{Generalized Sinkhorn maps.}
Consider the $\phi$-divergence regularized transport problem of Section~\ref{sec-sinkhorn-other-regularizers}. Its one-sided block updates, expressed using the Legendre transform $\phi^*$ defined in~\eqref{eq-legendre}, are the generalized soft $c$-transforms of~\eqref{eq-phi-soft-c-transform}; their alternating dual-maximization iteration is~\eqref{eq-generalized-soft-c-alternate-maximization}. After choosing a consistent extremal minimizer whenever a scalar update is not unique, one full cycle acting on the first potential is
\eql{\label{eq-generalized-sinkhorn-map}
\mathcal A_\phi(f)
\eqdef
\big(f^{c,\epsilon,\phi}\big)^{\bar c,\epsilon,\phi}.
}
\begin{prop}[Generalized Sinkhorn maps are topical]\label{prop-phi-double-soft-transform-monotone}
Assume that the scalar minimizer sets in~\eqref{eq-phi-soft-c-transform} are nonempty compact intervals, and select either their smallest elements everywhere or their largest elements everywhere. Each selected one-sided transform is order reversing and additively anti-homogeneous:
\[
g\leq g'
\Longrightarrow
g^{\bar c,\epsilon,\phi}\geq {g'}^{\bar c,\epsilon,\phi},
\qquad
(g+s)^{\bar c,\epsilon,\phi}
=
g^{\bar c,\epsilon,\phi}-s,
\]
and likewise for $f\mapsto f^{c,\epsilon,\phi}$.
\(\norm{\mathcal A_\phi(f)-\mathcal A_\phi(h)}_V \leq \norm{f-h}_V.\)
No selection convention is needed when the scalar minimizers are unique.
\end{prop}
\begin{proof}
Since $\phi$ is extended by $+\infty$ on $(-\infty,0)$, its Legendre transform $\phi^*$ from~\eqref{eq-legendre} is convex and nondecreasing. For fixed $x$, let $H_g^x(u)$ denote the scalar objective in the first line of~\eqref{eq-phi-soft-c-transform}. If $g\leq g'$, then $u\mapsto H_{g'}^x(u)-H_g^x(u)$ is nondecreasing: this follows by integrating the fact that $s\mapsto\phi^*(s+\delta)-\phi^*(s)$ is nondecreasing for every $\delta\geq0$.

Let $I=\argmin H_g^x$ and $I'=\argmin H_{g'}^x$. If $a=\min I$ and $a'=\min I'$ satisfied $a'>a$, the optimality inequalities $H_g^x(a)\leq H_g^x(a')$ and $H_{g'}^x(a')\leq H_{g'}^x(a)$ would contradict the preceding monotonicity; in the equality case, $a$ would also belong to $I'$, contradicting the definition of $a'$. Hence $\min I'\leq\min I$. The same argument with the maximal minimizers gives $\max I'\leq\max I$. Finally, $H_{g+s}^x(u)=H_g^x(u+s)+s$, so either extremal selection is additively anti-homogeneous. The second one-sided transform is treated symmetrically. Their composition is topical, and Proposition~\ref{prop-topical-variation-nonexpansive} gives the final estimate.
\end{proof}
\paragraph{Monotone convergence for generalized Sinkhorn.}
\begin{prop}[Monotone convergence of generalized Sinkhorn cycles]\label{prop-fortet-monotone}
Let $\Xx$ and $\Yy$ be compact metric spaces and $c\in\Cc(\Xx\times\Yy)$. Suppose that the assumptions of Proposition~\ref{prop-phi-double-soft-transform-monotone} hold, that the selected transforms are finite valued, and that $\mathcal A_\phi$ has a fixed point $f^\star\in\Cc(\Xx)$.
If $f^{(0)}\in\Cc(\Xx)$ is a subsolution, $f^{(0)}\leq\mathcal A_\phi(f^{(0)})$, then the iterates $f^{(\ell+1)}=\mathcal A_\phi(f^{(\ell)})$ converge uniformly to a fixed point of $\mathcal A_\phi$. The analogous conclusion holds for a supersolution, with a decreasing sequence. If the fixed point is unique modulo additive constants, the corresponding potential classes converge to this unique class.
\end{prop}
\begin{proof}
Because $\mathcal A_\phi$ is additively homogeneous, shifting a subsolution preserves the subsolution inequality. After such a shift, compactness of $\Xx$ allows us to assume $f^{(0)}\leq f^\star$. Topicality then gives
\(f^{(0)}\leq f^{(1)}\leq\cdots\leq f^\star,\)
so the orbit converges pointwise.

The one-sided transforms inherit the modulus of continuity of the cost. Indeed, if
\(\delta(x,x')\eqdef\sup_{y\in\Yy}|c(x,y)-c(x',y)|\),
order reversal and additive anti-homogeneity give
\[
\left|g^{\bar c,\epsilon,\phi}(x)-g^{\bar c,\epsilon,\phi}(x')\right|
\leq \delta(x,x').
\]
Thus $(f^{(\ell)})_{\ell\geq1}$ is uniformly bounded and equicontinuous. Arzel\`a--Ascoli makes the orbit relatively compact in the uniform norm, and its monotonicity forces every uniformly convergent subsequence to have the same pointwise limit. Hence the full sequence converges uniformly.

Moreover, every selected one-sided transform is nonexpansive for the uniform norm. If $\norm{g-h}_\infty\leq r$, then $g-r\leq h\leq g+r$; order reversal and additive anti-homogeneity imply $\norm{g^{\bar c,\epsilon,\phi}-h^{\bar c,\epsilon,\phi}}_\infty\leq r$, and similarly for the other block. Therefore $\mathcal A_\phi$ is uniformly continuous, and the limit is a fixed point. The supersolution argument reverses all inequalities. Uniqueness modulo constants gives the last claim.
\end{proof}
\subsection{Sinkhorn Convergence: Sublinear Robust Rate}
\begin{prop}[Robust $O(1/\ell)$ dual rate for discrete Sinkhorn]\label{prop-sinkhorn-dual-rate}
Let $\a\in\simplex_n$ and $\b\in\simplex_m$ be positive histograms, and let $\C\in\RR_+^{n\times m}$ be finite.
Initialize $\gD^{(0)}=0$, perform complete Sinkhorn dual cycles, and denote the resulting potentials by $(\fD^{(\ell)},\gD^{(\ell)})$ for $\ell\geq1$. Let
\(\Delta^{(\ell)} \eqdef \Dd_\epsilon(\fD^\star,\gD^\star) - \Dd_\epsilon(\fD^{(\ell)},\gD^{(\ell)})\)
be the dual suboptimality gap for the entropic dual objective~\eqref{eq-dual-sinkhorn-objective}. Then
\(0\leq\Delta^{(\ell)}\leq \frac{2\norm{\C}_\infty^2}{\epsilon \ell}, \qquad \ell\geq1.\)
\end{prop}
\begin{proof}
Write $M\eqdef\norm{\C}_\infty$; the case $M=0$ is immediate because one complete cycle produces $\a\otimes\b$. We first prove the oscillation estimate used below. If $\fD$ is the soft transform of $\gD$, as in~\eqref{eq-discrete-soft-c-transforms}, set
\[
Z_i\eqdef\sum_j b_j\exp\!\left(\frac{\gD_j-\C_{i,j}}{\epsilon}\right),
\qquad \fD_i=-\epsilon\log Z_i.
\]
Since $0\leq \C_{i,j}\leq M$, one has $e^{-M/\epsilon}Z_{i'}\leq Z_i\leq e^{M/\epsilon}Z_{i'}$ for every $i,i'$. Therefore $\abs{\fD_i-\fD_{i'}}\leq M$ and $\norm{\fD}_V\leq M$. The same argument applies to the other soft transform. Every Sinkhorn potential is produced by such a transform, and the block optimality equations imply the same property for $(\fD^\star,\gD^\star)$.
\(\norm{\fD^{(\ell)}}_V,\ \norm{\gD^{(\ell)}}_V, \ \norm{\fD^\star}_V,\ \norm{\gD^\star}_V \leq M.\)

Associate with any potentials the nonnegative matrix
\[
\P(\fD,\gD)
\eqdef
(\a\otimes\b)
\odot
\exp\!\left(\frac{\fD\oplus\gD-\C}{\epsilon}\right).
\]
At the end of cycle $\ell$, the matrix $\P^{(\ell)}\eqdef\P(\fD^{(\ell)},\gD^{(\ell)})$ has column marginal $\b$ and hence total mass one. Denote its row marginal by $r^{(\ell)}\eqdef\P^{(\ell)}\ones_m$ and set $\delta_\ell\eqdef\norm{\a-r^{(\ell)}}_1$. The dual gradient at this point is $(\a-r^{(\ell)},0)$. Concavity of $\Dd_\epsilon$ therefore gives
\(\Delta^{(\ell)} \leq \dotp{\fD^\star-\fD^{(\ell)}}{\a-r^{(\ell)}}.\)
For every vector $h$ and every zero-sum vector $z$, subtracting $(\max h+\min h)/2$ from $h$ proves
\(\abs{\dotp{h}{z}} \leq \frac12\norm{h}_V\norm{z}_1.\)
The oscillation estimate and the triangle inequality for $\norm{\cdot}_V$ thus yield
\begin{equation}\label{eq-sinkhorn-gap-residual}
\Delta^{(\ell)}\leq M\delta_\ell.
\end{equation}

Its row update satisfies
\[
\fD_i^{(\ell+1)}-\fD_i^{(\ell)}
=
\epsilon\log\!\left(\frac{a_i}{r_i^{(\ell)}}\right).
\]
The matrices before and after this row update both have total mass one. Substitution in~\eqref{eq-dual-sinkhorn-objective} therefore gives the exact gain $\epsilon\KLD(\a|r^{(\ell)})$. The subsequent column update can only increase the dual objective, so Pinsker's inequality, Theorem~\ref{thm-pinsker}, and~\eqref{eq-sinkhorn-gap-residual} give
\[
\Delta^{(\ell)}-\Delta^{(\ell+1)}
\geq
\epsilon\KLD(\a|r^{(\ell)})
\geq
\frac{\epsilon}{2}\delta_\ell^2
\geq
\frac{\epsilon}{2M^2}(\Delta^{(\ell)})^2.
\]

Set $x=M/\epsilon$, and let $\widehat\P^{(1)}$ be the matrix after the first row update from $\gD^{(0)}=0$. Directly from the update, \(e^{-x} \leq \frac{\widehat{\P}^{(1)}_{i,j}}{a_i b_j} \leq e^x.\)
Its column marginal $\widehat s$ consequently satisfies $e^{-x}\leq \widehat s_j/b_j\leq e^x$. The first column update gives $\P^{(1)}_{i,j}=\widehat{\P}^{(1)}_{i,j}b_j/\widehat s_j$, and hence
\(\frac{r_i^{(1)}}{a_i} = \sum_j\frac{\widehat{\P}^{(1)}_{i,j}}{a_i}\frac{b_j}{\widehat s_j} \in[e^{-x},e^x],\)
because $(\widehat{\P}^{(1)}_{i,j}/a_i)_j$ is a probability vector. Thus $\delta_1\leq\min\{2,e^x-1\}\leq2\min\{1,x\}$, where the last inequality uses $e^x-1\leq2x$ for $0\leq x\leq1$. Equation~\eqref{eq-sinkhorn-gap-residual} gives $\Delta^{(1)}\leq2M^2/\epsilon$.

Finally, whenever $\Delta^{(\ell+1)}>0$,
\[
\frac1{\Delta^{(\ell+1)}}-\frac1{\Delta^{(\ell)}}
=
\frac{\Delta^{(\ell)}-\Delta^{(\ell+1)}}{\Delta^{(\ell)}\Delta^{(\ell+1)}}
\geq
\frac{\epsilon}{2M^2}.
\]
Together with the bound on $\Delta^{(1)}$, telescoping proves $\Delta^{(\ell)}\leq2M^2/(\epsilon\ell)$. If a gap vanishes, all subsequent gaps vanish and the conclusion is immediate.
\end{proof}
\begin{cor}[Approximating unregularized OT by regularized dual costs]\label{cor-sinkhorn-dual-complexity}
Under the assumptions of Proposition~\ref{prop-sinkhorn-dual-rate}, suppose $nm>1$ and define the unregularized cost $L_0\eqdef\MKD_\C(\a,\b)$. For any $\delta>0$, choose $\epsilon=\delta/\log(nm)$. Then
\(\ell\geq \frac{2\norm{\C}_\infty^2\log(nm)}{\delta^2} \quad\Longrightarrow\quad \abs{L_0-\Dd_\epsilon(\fD^{(\ell)},\gD^{(\ell)})}\leq\delta.\)
\end{cor}
\begin{proof}
Let $L_\epsilon\eqdef\max_{\fD,\gD}\Dd_\epsilon(\fD,\gD)$ and let $\P^0$ solve the unregularized problem. The KL-normalized primal formulation~\eqref{eq-regularized-discr-rescaled} gives
\(0\leq L_\epsilon-L_0 \leq \epsilon\KLD(\P^0|\a\otimes\b) \leq \epsilon\log(nm).\)
The last inequality follows from $\KLD(\P^0|\a\otimes\b)=\HD(\a)+\HD(\b)-\HD(\P^0)\leq\log(nm)$.
Moreover, Proposition~\ref{prop-sinkhorn-dual-rate} gives $0\leq L_\epsilon-\Dd_\epsilon(\fD^{(\ell)},\gD^{(\ell)})\leq2\norm{\C}_\infty^2/(\epsilon\ell)$. Since the difference between the computed dual value and $L_0$ is the difference of these two nonnegative errors, its absolute value is bounded by their maximum. The prescribed $\epsilon$ and $\ell$ make both errors at most $\delta$.
\end{proof}
For dense $n\times n$ problems, one Sinkhorn cycle costs $O(n^2)$ operations. Corollary~\ref{cor-sinkhorn-dual-complexity} therefore yields the arithmetic complexity
\[
O\!\left(\frac{n^2\norm{\C}_\infty^2\log n}{\delta^2}\right)
\]
\subsection{Sinkhorn Convergence: Linear Hilbert Metric Rate}
\label{sec-sinkhorn-hilbert}
\paragraph{Projective contraction.}
\begin{defn}[Hilbert metric]\label{def-hilbert-metric}
On $\RR_{+,*}^n$, Hilbert's projective metric is
\eql{\label{eq-hilbert-metric}
\foralls (\uD,\uD') \in (\RR_{+,*}^n)^2, \quad
\Hilbert(\uD,\uD') \eqdef
\norm{\log(\uD)-\log(\uD')}_V.
}
Here, $\norm{\cdot}_V$ is the variation seminorm of Definition~\ref{def-variation-seminorm}, specialized to functions on the finite space $\{1,\ldots,n\}$: $\norm{z}_V=\max_i z_i-\min_i z_i$.
\end{defn}
\begin{prop}[Hilbert metric on the projective cone]
The function $\Hilbert$ defines a complete distance on the projective cone $\RR_{+,*}^n/\sim$, where $\uD \sim \uD'$ means that $\uD=s\uD'$ for some $s>0$.
\end{prop}
\begin{proof}
The map $\uD \mapsto \log \uD$ identifies $\RR_{+,*}^n/\sim$ with the quotient vector space $\RR^n/\Span(\ones_n)$, because multiplying $\uD$ by $s>0$ adds the constant vector $\log(s)\ones_n$. The variation seminorm $\norm{z}_V=\max_i z_i-\min_i z_i$ vanishes exactly on constant vectors, so it induces a norm on this quotient. Symmetry, the triangle inequality and separation therefore follow from the corresponding norm properties. Completeness follows because $\RR^n/\Span(\ones_n)$ is finite-dimensional and all finite-dimensional normed spaces are complete.
\end{proof}
\begin{thm}[Birkhoff contraction theorem]\label{thm-birkhoff}
Let $\K \in \RR_{+,*}^{n \times m}$. Then, for $(\vD,\vD') \in (\RR_{+,*}^m)^2$,
\[
\Hilbert(\K \vD,\K \vD') \leq \la(\K) \Hilbert(\vD,\vD'),
\qquad
\la(\K) \eqdef \frac{\sqrt{\eta(\K)}-1}{\sqrt{\eta(\K)}+1}<1,
\qquad
\eta(\K) \eqdef \umax{i,j,k,\ell}\frac{\K_{i,k}\K_{j,\ell}}{\K_{j,k}\K_{i,\ell}}.
\]
\end{thm}
\begin{proof}
Work in logarithmic coordinates and define
\(F(z)\eqdef\log(\K e^z), \qquad \mathsf P(z)_{i,k} \eqdef \frac{\K_{i,k}e^{z_k}}{(\K e^z)_i}.\)
Every row of $\mathsf P(z)$ is a probability vector and $\d F(z)[h]=\mathsf P(z)h$. We first bound its contraction in variation seminorm. If $p$ and $q$ are two probability rows, then
\(\abs{\dotp{p-q}{h}} \leq \frac12\norm{p-q}_1\norm{h}_V.\)
Indeed, subtracting $\min_k h_k$ and dividing by $\norm{h}_V$ reduces this inequality to the variational characterization of total variation by functions taking values in $[0,1]$.
\begin{equation}\label{eq-birkhoff-proof-dobrushin}
\norm{\mathsf P(z)h}_V
\leq
\delta(\mathsf P(z))\norm{h}_V,
\qquad
\delta(\mathsf P)
\eqdef
\frac12\max_{i,j}\sum_k\abs{\mathsf P_{i,k}-\mathsf P_{j,k}}.
\end{equation}

Fix rows $p,q$ of $\mathsf P(z)$, with indices $i,j$, and put $r_k=p_k/q_k$. Then $\sum_kq_kr_k=1$, and the definition of $\eta(\K)$ gives
\(\frac{\max_k r_k}{\min_k r_k} = \max_{k,\ell} \frac{\K_{i,k}\K_{j,\ell}}{\K_{j,k}\K_{i,\ell}} \leq\eta(\K).\)
Write $a=\min_k r_k\leq1\leq b=\max_k r_k$. If $a=b=1$, then $p=q$ and the desired bound is immediate. Otherwise, $(r-1)_+\leq (b-1)(r-a)/(b-a)$ on $[a,b]$, so
\(\frac12\norm{p-q}_1 = \sum_kq_k(r_k-1)_+ \leq \frac{(1-a)(b-1)}{b-a}.\)
Set $a=e^{-u}$ and $b=e^v$. The last quotient equals
\[
\frac{2\sinh(u/2)\sinh(v/2)}{\sinh((u+v)/2)}
\leq
\tanh\!\left(\frac{u+v}{4}\right)
\leq
\tanh\!\left(\frac{\log\eta(\K)}{4}\right)
=
\la(\K),
\]
where the first inequality follows from $\sinh(u/2)\sinh(v/2)\leq\sinh^2((u+v)/4)$. Thus~\eqref{eq-birkhoff-proof-dobrushin} holds with $\delta(\mathsf P(z))\leq\la(\K)$, uniformly in $z$.

For $z_t=(1-t)z'+tz$, the fundamental theorem of calculus and the triangle inequality for $\norm{\cdot}_V$ now give
\(\norm{F(z)-F(z')}_V \leq \int_0^1\norm{\mathsf P(z_t)(z-z')}_V\d t \leq \la(\K)\norm{z-z'}_V.\)
Taking $z=\log\vD$ and $z'=\log\vD'$ proves the result because $F(z)=\log(\K\vD)$.
\end{proof}
Theorem~\ref{thm-birkhoff} and completeness of that cone give a unique fixed ray $u^\star>0$, hence a unique positive leading eigenvector up to scale, and
\(\Hilbert(\K^\ell u^{(0)},u^\star) \leq \la(\K)^\ell\Hilbert(u^{(0)},u^\star).\)
Thus Hilbert contraction supplies both the Perron--Frobenius eigenvector and the global linear convergence of power iteration. The contraction can be visualized for every initialization at once on the three-state probability simplex. With the column-vector convention, every positive column-stochastic matrix $\K$ maps $\simplex_3$ into its interior, and
\(\K^{\ell+1}\simplex_3 = \K^\ell(\K\simplex_3) \subseteq \K^\ell\simplex_3.\)
To compare isotropic contraction, rotation, and anisotropic contraction, set
\[
J_3\eqdef\frac{1}{3}\ones_3\ones_3^\top,
\qquad
A\eqdef
\begin{pmatrix}
0&-1&1\\
1&0&-1\\
-1&1&0
\end{pmatrix},
\qquad
\K_{\rho,\delta}\eqdef \rho\Id_3+(1-\rho)J_3+\delta A.
\]
The three kernels used below are
\[
\K_1=\K_{.90,0},
\qquad
\K_2=\K_{.90,.014},
\qquad
\K_3=
\begin{pmatrix}
.950&.018&.032\\
.042&.880&.078\\
.008&.102&.890
\end{pmatrix}.
\]
\paragraph{Sinkhorn contraction.}
\begin{thm}[Projective linear convergence of Sinkhorn]\label{thm-sinkhorn-hilbert-linear}
Starting from $\vD^{(0)}>0$, use the scaling iterates of~\eqref{eq-sinkhorn} and the physical plan sequence $\P^{(r)}$ defined in~\eqref{eq-sink-matrix}. Fix any optimal scaling pair $(\uD^\star,\vD^\star)$.
Writing $\la=\la(\K)$, one has
\eql{\label{eq-convlin-sinkh}
\Hilbert(\vD^{(\ell)},\vD^\star)
\leq
\la^{2\ell}\Hilbert(\vD^{(0)},\vD^\star),
\qquad
\Hilbert(\uD^{(\ell+1)},\uD^\star)
\leq
\la^{2\ell+1}\Hilbert(\vD^{(0)},\vD^\star).
}
Thus the scaling rays converge linearly. After fixing any continuous gauge, their representatives converge, and $\P^{(2\ell)}\to\P^\star$ entrywise.
\eql{\label{eq-convsinkh-control}
\Hilbert(\uD^{(\ell)},\uD^\star)
\leq
\frac{\Hilbert(\P^{(2\ell)}\ones_m,\a)}{1-\la^2}
\qandq
\Hilbert(\vD^{(\ell)},\vD^\star)
\leq
\frac{\Hilbert((\P^{(2\ell+1)})^\top\ones_n,\b)}{1-\la^2}.
}
\eqllead{Lastly, for $\ell\geq1$,}{eq-convlin-sinkh-prim}{
\norm{\log(\P^{(2\ell)})-\log(\P^\star)}_\infty
\leq
\Hilbert(\uD^{(\ell)},\uD^\star)
+
\Hilbert(\vD^{(\ell)},\vD^\star),
}
where $\P^\star$ is the unique solution of~\eqref{eq-regularized-discr}.
\end{thm}
\begin{proof}
Entrywise multiplication by a fixed positive vector and entrywise inversion are isometries of Hilbert's metric. Therefore the row map $R(\vD)=\a/(\K\vD)$ and the column map $C(\uD)=\b/(\K^\top\uD)$ are both $\la$-Lipschitz by Theorem~\ref{thm-birkhoff}. Their full-cycle compositions $C\circ R$ and $R\circ C$ are $\la^2$-contractions. Since $\vD^\star=C(R(\vD^\star))$ and $\uD^\star=R(C(\uD^\star))$, iteration gives~\eqref{eq-convlin-sinkh}.

For any contraction $F$ with factor $q<1$ and fixed point $x^\star$, the triangle inequality gives \(d(x,x^\star) \leq d(x,Fx)+d(Fx,Fx^\star) \leq d(x,Fx)+q\,d(x,x^\star),\)
hence $d(x,x^\star)\leq d(x,Fx)/(1-q)$. Apply this with $q=\la^2$ to the two full-cycle maps. The scaling identities
\(
\Hilbert(\uD^{(\ell+1)},\uD^{(\ell)})
=\Hilbert(\a,\P^{(2\ell)}\ones_m)
\)
and
\(
\Hilbert(\vD^{(\ell+1)},\vD^{(\ell)})
=\Hilbert(\b,(\P^{(2\ell+1)})^\top\ones_n)
\)
then prove~\eqref{eq-convsinkh-control}.

Finally, write $\xi_i=\log(u_i^{(\ell)}/u_i^\star)$ and $\zeta_j=\log(v_j^{(\ell)}/v_j^\star)$. Then $\log(\P_{i,j}^{(2\ell)}/\P_{i,j}^\star)=\xi_i+\zeta_j$ and $\operatorname{osc}(\xi\oplus\zeta)=\operatorname{osc}(\xi)+\operatorname{osc}(\zeta)$.
Both plans have total mass one, so $\sum_{i,j}\P_{i,j}^\star e^{\xi_i+\zeta_j}=1$. This proves~\eqref{eq-convlin-sinkh-prim}.
\end{proof}
\paragraph{Nonlinear Sinkhorn images of the simplex.}
Using the row and column maps from the proof, define the full-cycle map on the left scaling by
\[
F_{\uD}(\uD)
\eqdef
R(C(\uD))
=
\a\oslash\left[\K\left(\b\oslash(\K^\top\uD)\right)\right].
\]
It is positively homogeneous, $F_{\uD}(s\uD)=sF_{\uD}(\uD)$ for $s>0$, and therefore induces the projective self-map
\(\widehat F_{\uD}(p) \eqdef \frac{F_{\uD}(p)}{\dotp{F_{\uD}(p)}{\ones_3}}, \qquad p\in\simplex_3.\)
Thus $\widehat\uD^{(\ell)}\eqdef\uD^{(\ell)}/\dotp{\uD^{(\ell)}}{\ones_3}$ follows $\widehat\uD^{(\ell+1)}=\widehat F_{\uD}(\widehat\uD^{(\ell)})$ when $\ell$ indexes complete Sinkhorn cycles. Since $\widehat F_{\uD}$ maps $\simplex_3$ into itself, its set-valued iterates are nested:
\[
\widehat F_{\uD}^{\,\ell+1}(\simplex_3)
=
\widehat F_{\uD}^{\,\ell}\!\left(\widehat F_{\uD}(\simplex_3)\right)
\subseteq
\widehat F_{\uD}^{\,\ell}(\simplex_3).
\]
Theorem~\ref{thm-sinkhorn-hilbert-linear} quantifies their projective shrinkage: since the first image is strictly positive,
\[
\operatorname{diam}_{\Hilbert}\!\left(\widehat F_{\uD}^{\,\ell}(\simplex_3)\right)
\leq
\la(\K)^{2(\ell-1)}
\operatorname{diam}_{\Hilbert}\!\left(\widehat F_{\uD}(\simplex_3)\right),
\qquad \ell\geq1.
\]
\paragraph{Dual-potential form of the contraction.}
By~\eqref{eq-entropy-pd}, the KL-normalized potentials satisfy
\(
\fD^{(\ell)}=\epsilon\log(\uD^{(\ell)}\oslash\a)
\)
and
\(
\gD^{(\ell)}=\epsilon\log(\vD^{(\ell)}\oslash\b)
\).
The Hilbert contraction can therefore be read directly on potential classes because
\[
\norm{\fD^{(\ell)}-\fD^\star}_V
=
\epsilon\Hilbert(\uD^{(\ell)},\uD^\star),
\qquad
\norm{\gD^{(\ell)}-\gD^\star}_V
=
\epsilon\Hilbert(\vD^{(\ell)},\vD^\star).
\]
Thus a complete cycle contracts both potential classes at the rate $\la(\K)^2$. Moreover, the temperature factor must be retained when passing from potentials to densities:
\eq{
\norm{\log\pa{\P^{(2\ell)}\oslash\P^\star}}_\infty
=
\frac1\epsilon
\norm{ (\fD^{(\ell)}-\fD^\star) \oplus (\gD^{(\ell)}-\gD^\star) }_\infty
\leq
\frac{\norm{\fD^{(\ell)}-\fD^\star}_V + \norm{\gD^{(\ell)}-\gD^\star}_V}{\epsilon}.
}
With $\eta=\eta(\K_\epsilon)$ as in Theorem~\ref{thm-birkhoff}, \(\la=\frac{\sqrt{\eta}-1}{\sqrt{\eta}+1} \leq \tanh(R/(2\epsilon))<1.\)
The posterior estimates~\eqref{eq-convsinkh-control} become
\(\norm{\fD^{(\ell)}-\fD^\star}_V \leq \frac{\epsilon}{1-\la^2} \norm{\log((\P^{(2\ell)}\ones_m)\oslash\a)}_V,\)
and
\(\norm{\gD^{(\ell)}-\gD^\star}_V \leq \frac{\epsilon}{1-\la^2} \norm{\log(((\P^{(2\ell+1)})^\top\ones_n)\oslash\b)}_V.\)
\subsection{Local Convergence Analysis and Acceleration}
\label{sec-sinkhorn-local-acceleration}
\paragraph{Conditional operator and maximal correlation.}
Let
\(
L_0^2(\al)\eqdef\enscond{h\in L^2(\al)}{\int h\d\al=0}
\)
\begin{defn}[Conditional coupling operator]\label{def-sinkhorn-conditional-operator}
For a coupling $\pi\in\Couplings(\al,\be)$, let
\(
\pi(\d x,\d y)=\al(\d x)\pi_x(\d y)
=\be(\d y)\pi^y(\d x)
\)
be its two disintegrations. Define
\(
T_\pi:L^2(\be)\to L^2(\al)
\)
and
\(
T_\pi^*:L^2(\al)\to L^2(\be)
\)
by
\eql{\label{eq-sinkhorn-conditional-operator}
(T_\pi k)(x)\eqdef\int_\Yy k(y)\d\pi_x(y),
\qquad
(T_\pi^*h)(y)\eqdef\int_\Xx h(x)\d\pi^y(x).
}
The centered maximal correlation of $\pi$ is the restricted operator norm
\eql{\label{eq-sinkhorn-maximal-correlation}
\sigma(\pi)
\eqdef
\norm{T_\pi:L_0^2(\be)\to L_0^2(\al)}_{\rm op},
}
with the convention that this norm is zero when the domain is $\{0\}$. If both centered spaces are nonzero, equivalently,
\[
\sigma(\pi)
=
\sup_{\substack{k\in L_0^2(\be)\\\norm{k}_{L^2(\be)}=1}}
\norm{T_\pi k}_{L^2(\al)}
=
\sup_{\substack{h\in L_0^2(\al),\,k\in L_0^2(\be)\\
\norm{h}_{L^2(\al)}=\norm{k}_{L^2(\be)}=1}}
\int h(x)k(y)\d\pi(x,y).
\]
For the optimal entropic coupling $\pi_\epsilon$ of~\eqref{eq-entropic-generic}, set
\(
T_\epsilon\eqdef T_{\pi_\epsilon}
\)
and
\(
\sigma_\epsilon\eqdef\sigma(\pi_\epsilon).
\)
\end{defn}
In the discrete setting of~\eqref{eq-regularized-discr}, assuming the weights are positive, the optimizer $\P_\epsilon$ gives
\eql{\label{eq-sinkhorn-conditional-operator-discrete}
T_\epsilon=\diag(\a)^{-1}\P_\epsilon,
\qquad
T_\epsilon^*=\diag(\b)^{-1}\P_\epsilon^\top,
}
where adjointness is understood in the $\a$- and $\b$-weighted Euclidean inner products. If $\min(n,m)\geq2$, then $\sigma_\epsilon$ is the second singular value of
\(
\diag(\a)^{-1/2}\P_\epsilon\diag(\b)^{-1/2}
\),
\paragraph{Dual Hessian and exact local rate.}
\begin{prop}[Dual Hessian and local Sinkhorn rate]\label{prop-sinkhorn-local-rate}
Assume the setting of Proposition~\ref{prop-continuous-entropic-duality}, and that differentiation under the Gibbs integrals and the stated $L^2$ Taylor expansions are valid. These requirements are automatic in finite spaces and form the local regularity hypothesis in continuous settings. For bounded directions $(h,k)$ and $(h',k')$,
\eql{\label{eq-sinkhorn-dual-hessian-form}
-\epsilon D^2\Dd_\epsilon(f_\epsilon,g_\epsilon)
[(h,k),(h',k')]
=
\int_{\Xx\times\Yy}
(h(x)+k(y))(h'(x)+k'(y))\d\pi_\epsilon(x,y).
}
Hence the Hessian form extends continuously to
\(L^2(\al)\oplus L^2(\be)\), where the positive dual curvature operator is
\eql{\label{eq-sinkhorn-dual-hessian-operator}
-\nabla^2\Dd_\epsilon(f_\epsilon,g_\epsilon)
=
\frac1\epsilon
\begin{pmatrix}
\Id&T_\epsilon\\
T_\epsilon^*&\Id
\end{pmatrix}.
}
After quotienting by the gauge direction $(1,-1)$ and using the product-$L^2$ quotient norm, its spectrum is contained in $[(1-\sigma_\epsilon)/\epsilon,2/\epsilon]$. If at least one centered space is nonzero, both endpoints belong to the spectrum, although the lower one need not be an eigenvalue. In particular, when $\sigma_\epsilon<1$, the spectral condition number of the full dual is \(2/(1-\sigma_\epsilon)\). If both centered spaces are $\{0\}$, the quotient is one-dimensional, its sole eigenvalue is $2/\epsilon$, and its condition number is one.
The derivatives of the two soft-transform updates~\eqref{eq-soft-c-cont-f}--\eqref{eq-soft-c-cont-g} at the optimum are
\eql{\label{eq-sinkhorn-soft-transform-jacobians}
D[g\mapsto g^{\bar c,\epsilon}](g_\epsilon)=-T_\epsilon,
\qquad
D[f\mapsto f^{c,\epsilon}](f_\epsilon)=-T_\epsilon^*.
}
\(g\mapsto(g^{\bar c,\epsilon})^{c,\epsilon}\)
is $T_\epsilon^*T_\epsilon$. If the gauges are fixed so that
\(e^{(\ell)}\eqdef g^{(\ell)}-g_\epsilon\in L_0^2(\be)\), then
\eql{\label{eq-sinkhorn-local-error}
e^{(\ell+1)}
=
T_\epsilon^*T_\epsilon e^{(\ell)}
+
o\!\left(\norm{e^{(\ell)}}_{L^2(\be)}\right).
}
Unless the centered error vanishes in finitely many cycles, this gives
\[
\limsup_{\substack{\ell\to\infty\\e^{(\ell)}\neq0}}
\frac{\norm{e^{(\ell+1)}}_{L^2(\be)}}
{\norm{e^{(\ell)}}_{L^2(\be)}}
\leq\sigma_\epsilon^2.
\]
If $\sigma_\epsilon^2$ is an isolated eigenvalue and the asymptotic error has a nonzero projection onto its eigenspace, this ratio converges to $\sigma_\epsilon^2$. When $L_0^2(\be)=\{0\}$, there is no centered target-potential degree of freedom and the local cycle statement is vacuous.
\end{prop}
\begin{proof}
Twice differentiating the exponential term in~\eqref{eq-dual-sinkhorn-objective} gives~\eqref{eq-sinkhorn-dual-hessian-form}; the density law~\eqref{eq-continuous-entropic-density-law} identifies its reference measure with $\pi_\epsilon$. Its two diagonal terms are the $L^2$ inner products of the marginals, and its cross term is represented by $T_\epsilon$, which proves~\eqref{eq-sinkhorn-dual-hessian-operator}.

Each quotient class has a unique representative orthogonal to the gauge direction, hence of the form
\(
(h_0+m,k_0+m)
\)
with $h_0\in L_0^2(\al)$ and $k_0\in L_0^2(\be)$. On the centered component, Cauchy--Schwarz and~\eqref{eq-sinkhorn-maximal-correlation} give
\[
(1-\sigma_\epsilon)
\bigl(\norm{h_0}_{L^2(\al)}^2+\norm{k_0}_{L^2(\be)}^2\bigr)
\leq
\norm{h_0}_{L^2(\al)}^2+\norm{k_0}_{L^2(\be)}^2
+2\langle h_0,T_\epsilon k_0\rangle_{L^2(\al)}
\leq
(1+\sigma_\epsilon)
\bigl(\norm{h_0}_{L^2(\al)}^2+\norm{k_0}_{L^2(\be)}^2\bigr).
\]
The common-constant direction $(1,1)$ has eigenvalue $2/\epsilon$, whereas $(1,-1)$ is the gauge kernel. If both centered spaces are nonzero, approximate maximizers in~\eqref{eq-sinkhorn-maximal-correlation} show that the lower spectral endpoint is $(1-\sigma_\epsilon)/\epsilon$. If exactly one centered space is nonzero, every vector in that centered factor has eigenvalue $1/\epsilon$, which is the same endpoint because $\sigma_\epsilon=0$. If both centered spaces vanish, only the common constant direction remains after quotienting. This proves all stated condition-number formulas without requiring a singular-vector basis.

Differentiating either log-partition formula in Definition~\ref{def-continuous-soft-c-transform} produces minus the corresponding conditional expectation under $\pi_\epsilon$, proving~\eqref{eq-sinkhorn-soft-transform-jacobians}. The chain rule and a Taylor expansion then give~\eqref{eq-sinkhorn-local-error}. Since
\(\norm{T_\epsilon^*T_\epsilon}_{L_0^2(\be)\to L_0^2(\be)}
=\sigma_\epsilon^2\),
the rate follows.
\end{proof}
When $L_0^2(\be)\neq\{0\}$, the same spectral gap has a direct probabilistic interpretation. The law of total variance gives
\eql{\label{eq-sinkhorn-conditional-variance-gap}
1-\sigma_\epsilon^2
=
\inf_{\substack{k\in L_0^2(\be)\\\norm{k}_{L^2(\be)}=1}}
\int_\Xx
\operatorname{Var}_{\pi_{\epsilon,x}}\!\bigl(k(Y)\bigr)
\d\al(x).
}
Thus Sinkhorn slows down when the optimal conditional laws $Y|X=x$ become nearly deterministic. This is exactly what happens as $\epsilon\to0$ in a smooth Monge regime.

\begin{prop}[The global Hilbert factor controls the local rate]\label{prop-sinkhorn-hilbert-controls-local}
In the discrete positive-kernel setting of Theorem~\ref{thm-sinkhorn-hilbert-linear}, let $\la(\K)$ be the Birkhoff contraction factor defined in Theorem~\ref{thm-birkhoff}. Then
\eql{\label{eq-sinkhorn-hilbert-controls-local}
\sigma_\epsilon\leq\la(\K)<1,
\qquad
\sigma_\epsilon^2\leq\la(\K)^2.
}
\end{prop}
\begin{proof}
Let $R(\vD)=\a\oslash(\K\vD)$ and $C(\uD)=\b\oslash(\K^\top\uD)$ be the row and column scaling maps used in the proof of Theorem~\ref{thm-sinkhorn-hilbert-linear}. Entrywise inversion and multiplication by a fixed positive vector are isometries of Hilbert's metric. Theorem~\ref{thm-birkhoff}, applied successively to $\K$ and $\K^\top$, therefore gives the pairwise estimate
\(\Hilbert\bigl(C(R(\vD)),C(R(\widetilde\vD))\bigr) \leq \la(\K)^2\Hilbert(\vD,\widetilde\vD),\)
where $\la(\K^\top)=\la(\K)$ follows directly from its cross-ratio definition. Under the logarithmic change of variables $g=\epsilon\log\vD$, Hilbert's metric becomes $\epsilon^{-1}\norm{\cdot}_V$ by~\eqref{eq-hilbert-metric}. Hence the complete dual update
\(
\mathcal S(g)\eqdef(g^{\bar c,\epsilon})^{c,\epsilon}
\)
is $\la(\K)^2$-Lipschitz for the variation norm on the quotient by constants: \(\norm{\mathcal S(g)-\mathcal S(\widetilde g)}_V \leq \la(\K)^2\norm{g-\widetilde g}_V.\)
Differentiate this inequality at the fixed point $g_\epsilon$. Proposition~\ref{prop-sinkhorn-local-rate} gives
\(D\mathcal S(g_\epsilon)=T_\epsilon^*T_\epsilon\), and therefore
\(\norm{T_\epsilon^*T_\epsilon h}_V \leq \la(\K)^2\norm{h}_V \qquad \text{for every }h\in L_0^2(\be).\)
In the present finite-dimensional setting, $T_\epsilon^*T_\epsilon$ is positive and self-adjoint for the $\be$-weighted inner product, and its largest eigenvalue on $L_0^2(\be)$ is $\sigma_\epsilon^2$. If $\sigma_\epsilon>0$, take a corresponding eigenvector $h_\star$. It is nonconstant, so $\norm{h_\star}_V>0$, and the preceding bound gives
\(\sigma_\epsilon^2\norm{h_\star}_V = \norm{T_\epsilon^*T_\epsilon h_\star}_V \leq \la(\K)^2\norm{h_\star}_V.\)
The conclusion is immediate; the case $\sigma_\epsilon=0$ is trivial.
\end{proof}
If a uniform Laplace expansion turns the conditional-variance form in~\eqref{eq-sinkhorn-conditional-variance-gap} into $\epsilon$ times a coercive weighted Dirichlet form, then
\(
1-\sigma_\epsilon^2\asymp\epsilon.
\)
\paragraph{Blockwise over-relaxation.}
A minimal acceleration extrapolates each block toward its soft-transform maximizer: for $\omega>0$,
\eql{\label{eq-sinkhorn-block-overrelaxation}
f^{(\ell+1)}=(1-\omega)f^{(\ell)}+\omega(g^{(\ell)})^{\bar c,\epsilon},
\qquad
g^{(\ell+1)}=(1-\omega)g^{(\ell)}+\omega(f^{(\ell+1)})^{c,\epsilon}.
}
\begin{prop}[Optimal local block relaxation]\label{prop-sinkhorn-optimal-overrelaxation}
Assume the hypotheses of Proposition~\ref{prop-sinkhorn-local-rate} and $\sigma_\epsilon<1$. For every $0<\omega<2$, the iteration~\eqref{eq-sinkhorn-block-overrelaxation} is locally linearly convergent modulo the additive gauge. Its optimal relaxation parameter and corresponding asymptotic factor are
\eql{\label{eq-sinkhorn-optimal-overrelaxation}
\omega_\star
=
\frac{2}{1+\sqrt{1-\sigma_\epsilon^2}},
\qquad
r_\star
=
\omega_\star-1
=
\frac{1-\sqrt{1-\sigma_\epsilon^2}}
{1+\sqrt{1-\sigma_\epsilon^2}}.
}
For $\omega=1$, the factor is $\sigma_\epsilon^2$, as in Proposition~\ref{prop-sinkhorn-local-rate}.
\end{prop}
\begin{proof}
Use~\eqref{eq-sinkhorn-soft-transform-jacobians}. The polar decomposition of $T_\epsilon$ and the spectral theorem for $(T_\epsilon^*T_\epsilon)^{1/2}$ reduce the centered linearized iteration to scalar spectral values $s\in[0,\sigma_\epsilon]$; in finite dimensions these are the usual singular values. On the corresponding two-dimensional mode, one relaxed cycle has matrix
\[
\mathsf M_\omega(s)
=
\begin{pmatrix}
1-\omega&-\omega s\\
-\omega s(1-\omega)&1-\omega+\omega^2s^2
\end{pmatrix},
\]
whose characteristic polynomial is $z^2-(2(1-\omega)+\omega^2s^2)z+(1-\omega)^2$.
The supremum of the root moduli is reached, or approached, as $s\to\sigma_\epsilon$. When $\sigma_\epsilon>0$, the discriminant of this limiting mode vanishes when
\(\omega^2\sigma_\epsilon^2=4(\omega-1)\):
the roots coalesce there and form an equal-modulus complex-conjugate pair beyond it. The smaller solution is $\omega_\star$ in~\eqref{eq-sinkhorn-optimal-overrelaxation}; at this value both root moduli equal $\omega_\star-1$. Below this threshold the dominant real root decreases with $\omega$, while above it the conjugate-root modulus $\omega-1$ increases. If $\sigma_\epsilon=0$, the polynomial is $(z-(1-\omega))^2$ and the optimum $\omega_\star=1$ follows directly. This proves optimality and local convergence for $0<\omega<2$.
The nongauge constant mode has factor $(1-\omega)^2$ and never dominates these centered modes.
\end{proof}
Set $\delta_\epsilon=1-\sigma_\epsilon^2$. When $\delta_\epsilon$ is small, the exact expression in~\eqref{eq-sinkhorn-optimal-overrelaxation} gives the complete-cycle factor
\(r_\star = \frac{1-\sqrt{\delta_\epsilon}}{1+\sqrt{\delta_\epsilon}} = 1-2\sqrt{\delta_\epsilon} +2\delta_\epsilon +O(\delta_\epsilon^{3/2}),\)
Per half-step, the corresponding factor is
\(
\sqrt{r_\star}=1-\sqrt{\delta_\epsilon}+O(\delta_\epsilon)
\).
In the finite positive-kernel setting, the Hilbert estimate used in Theorem~\ref{thm-sinkhorn-hilbert-linear} yields global convergence of the relaxed iteration for the conservative range
\(
0<\omega<2/(1+\la(\K))
\);
\paragraph{Entropic semi-dual and variable projection.}
The hard semi-dual \(\Ee_0\) in~\eqref{eq-semi-dual} eliminates one potential from the full dual \(\Dd_0\). At positive temperature, the full dual is \(\Dd_\epsilon\) from~\eqref{eq-dual-sinkhorn-objective}, and the soft \(\bar c\)-transform \(g^{\bar c,\epsilon}\) is defined by~\eqref{eq-soft-c-cont-g}. Applying the same partial maximization gives the entropic semi-dual
\eql{\label{eq-entropic-semidual-local}
\Ee_\epsilon(g)
\eqdef
\Dd_\epsilon(g^{\bar c,\epsilon},g)
=
\int_\Xx g^{\bar c,\epsilon}\d\al
+
\int_\Yy g\d\be.
}
For a potential $g$, let $\pi_g$ be the Gibbs measure obtained from the density law~\eqref{eq-continuous-entropic-density-law} by replacing $(f_\epsilon,g_\epsilon)$ with $(g^{\bar c,\epsilon},g)$. Write
\(
\pi_g(\d x,\d y)=\al(\d x)\pi_{g,x}(\d y)
\)
and denote its second marginal by $\be_g$. The first marginal is exactly $\al$ by construction.

A representative application is training a two-layer network by \(\min_{U,V}\frac12\norm{U\sigma(VX)-Y}_{\rm F}^2,\)
\begin{prop}[Semi-dual curvature and conditioning]\label{prop-sinkhorn-semidual-curvature}
For bounded directions $k,k'$,
\eql{\label{eq-entropic-semidual-derivatives}
D\Ee_\epsilon(g)[k]
=
\int_\Yy k\d(\be-\be_g),
\qquad
D^2\Ee_\epsilon(g)[k,k']
=
-\frac1\epsilon
\int_\Xx
\operatorname{Cov}_{\pi_{g,x}}\!\bigl(k(Y),k'(Y)\bigr)
\d\al(x).
}
At $g=g_\epsilon$, its positive curvature on $L_0^2(\be)$ is
\eql{\label{eq-entropic-semidual-hessian}
-\nabla^2\Ee_\epsilon(g_\epsilon)
=
\frac1\epsilon(\Id-T_\epsilon^*T_\epsilon),
\qquad
\frac{1-\sigma_\epsilon^2}{\epsilon}\Id
\preceq
-\nabla^2\Ee_\epsilon(g_\epsilon)
\preceq
\frac1\epsilon\Id.
}
If $L_0^2(\be)\neq\{0\}$, its spectral condition number is at most
\(
1/(1-\sigma_\epsilon^2)
\). If $L_0^2(\be)=\{0\}$, the semi-dual has no centered degree of freedom and no spectral condition number is assigned.
\end{prop}
\begin{proof}
The envelope theorem gives the first derivative in~\eqref{eq-entropic-semidual-derivatives}. Differentiating the normalized conditional Gibbs law gives the covariance formula. At the optimum, total covariance yields
\(\int_\Xx \operatorname{Cov}_{\pi_{\epsilon,x}}(k(Y),k'(Y)) \d\al(x) = \int_\Yy kk'\d\be - \int_\Xx(T_\epsilon k)(T_\epsilon k')\d\al,\)
which is~\eqref{eq-entropic-semidual-hessian}. Equivalently, this operator is the Schur complement of the first identity block in~\eqref{eq-sinkhorn-dual-hessian-operator}. Its spectrum on centered functions lies in
\(
[(1-\sigma_\epsilon^2)/\epsilon,1/\epsilon]
\).
\end{proof}
\subsection{Entropic Optimal Transport between Gaussians}
\label{sec-gaussian-sinkhorn}
\begin{prop}[Quadratic closure of Sinkhorn iterates]\label{prop-gaussian-sinkhorn-closure}
Let $\be=\Gaussian(\mean_\be,\cov_\be)$ on $\RR^d$ and take $c(x,y)=\norm{x-y}^2$. If $g(y)$ is a quadratic polynomial such that the Gaussian integral below is finite, then the soft transform
\[
f(x)
=
-\epsilon\log
\int
\exp\!\left(\frac{g(y)-\norm{x-y}^2}{\epsilon}\right)
\d\be(y)
\]
is a quadratic polynomial in $x$.
\end{prop}
\begin{proof}
The exponent is the sum of a quadratic polynomial in $y$ and the logarithm of the Gaussian density of $\be$. Completing the square in $y$ evaluates the integral as a positive constant times the exponential of a quadratic polynomial in $x$. Taking $-\epsilon\log$ therefore gives a quadratic polynomial.
\end{proof}
\begin{prop}[Balanced entropic OT between Gaussians]\label{prop-gaussian-sinkhorn-closed-form}
Let $\al=\Gaussian(\mean_\al,\cov_\al)$ and $\be=\Gaussian(\mean_\be,\cov_\be)$ have positive-definite covariances, and let $\cov_\al^{1/2}\cov_\be^{1/2}=U\diag(\sigma_i)V^\top$ be a singular-value decomposition. For the balanced entropic problem~\eqref{eq-entropic-generic} with $c(x,y)=\norm{x-y}^2$,
the optimizer is Gaussian with cross-covariance
\(K_\epsilon = \cov_\al^{1/2} U\diag(s_i)V^\top \cov_\be^{1/2}, \qquad s_i = \frac{\sqrt{\epsilon^2+16\sigma_i^2}-\epsilon}{4\sigma_i}.\)
The optimal value is
\[
\norm{\mean_\al-\mean_\be}^2
+
\tr(\cov_\al)+\tr(\cov_\be)
+
\sum_i
\left(
-2\sigma_i s_i
-\frac{\epsilon}{2}\log(1-s_i^2)
\right).
\]
As $\epsilon\downarrow0$, $s_i\to1$ and the full covariance contribution, including the two trace terms, converges to $\Bb(\cov_\al,\cov_\be)^2$.
\end{prop}
\begin{proof}
Let $(X,Y)$ be any coupling with finite second moments and cross-covariance
\(K=\EE\big[(X-\mean_\al)(Y-\mean_\be)^\top\big]\).
Replacing $(X,Y)$ by the Gaussian vector with the same mean and covariance leaves the quadratic cost unchanged. Since the marginals are fixed, \(\KL(\pi|\al\otimes\be) = -h(X,Y)+h(\al)+h(\be),\)
where $h$ denotes differential entropy when it is finite. Among laws with a fixed covariance, the Gaussian maximizes entropy; if the entropy is not finite, the relative entropy is already $+\infty$. Thus the Gaussian replacement cannot increase the objective, and it is enough to optimize over Gaussian couplings.

Any such coupling has covariance
\[
\begin{pmatrix}
\cov_\al & K\\
K^\top & \cov_\be
\end{pmatrix}.
\]
Write $K=\cov_\al^{1/2}S\cov_\be^{1/2}$. The block covariance constraint is equivalent to the singular values of $S$ being at most one, and finite entropy forces them to be strictly smaller than one. The cost depends on $K$ through
\(\norm{\mean_\al-\mean_\be}^2+\tr(\cov_\al)+\tr(\cov_\be)-2\tr(K),\)
while $\KL(\pi|\al\otimes\be)=-\frac12\log\det(\Id-SS^\top)$.
By von Neumann's trace inequality, the minimizer aligns $S$ with the singular vectors of $\cov_\al^{1/2}\cov_\be^{1/2}$, so $S=U\diag(s_i)V^\top$. The problem separates into scalar minimizations
\(\min_{0\leq s<1} -2\sigma_i s-\frac{\epsilon}{2}\log(1-s^2).\)
The first-order condition is $2\sigma_i=\epsilon s/(1-s^2)$, whose positive solution is the displayed $s_i$. Substitution gives the value formula. Since $s_i\to1$ and $\epsilon\log(1-s_i^2)\to0$ as $\epsilon\downarrow0$, the spectral sum converges to $-2\sum_i\sigma_i$. The identity $\sum_i\sigma_i=\tr((\cov_\al^{1/2}\cov_\be\cov_\al^{1/2})^{1/2})$ then gives the Bures--Wasserstein covariance contribution.
\end{proof}
\begin{cor}[Gaussian Sinkhorn divergence and smoothed Bures term]\label{cor-gaussian-sinkhorn-divergence}
Let $\al=\Gaussian(\mean_\al,\cov_\al)$ and $\be=\Gaussian(\mean_\be,\cov_\be)$ have positive-definite covariances. For $r>0$, define
\[
\tau_\epsilon(r)
\eqdef
\frac{\sqrt{\epsilon^2+16r^2}-\epsilon}{4r},
\qquad
\psi_\epsilon(r)
\eqdef
-2r\,\tau_\epsilon(r)
-\frac{\epsilon}{2}\log\bigl(1-\tau_\epsilon(r)^2\bigr).
\]
If $\sigma_i(\Sigma,\Lambda)$ denotes the singular values of $\Sigma^{1/2}\Lambda^{1/2}$ and $\lambda_i(\Sigma)$ the eigenvalues of $\Sigma$, then the debiased Sinkhorn divergence~\eqref{eq-sinkhorn-divergence} is
\(\bar\MK_{\norm{\cdot-\cdot}^2}^{\epsilon}(\al,\be) = \norm{\mean_\al-\mean_\be}^2 + \Bb_\epsilon(\cov_\al,\cov_\be)^2,\)
where the Gaussian covariance contribution is the debiased smoothed Bures term
\[
\Bb_\epsilon(\Sigma,\Lambda)^2
\eqdef
\sum_i \psi_\epsilon\bigl(\sigma_i(\Sigma,\Lambda)\bigr)
-\frac12\sum_i\psi_\epsilon\bigl(\lambda_i(\Sigma)\bigr)
-\frac12\sum_i\psi_\epsilon\bigl(\lambda_i(\Lambda)\bigr).
\]
Moreover $\Bb_\epsilon(\Sigma,\Lambda)^2\to\Bb(\Sigma,\Lambda)^2$ as $\epsilon\downarrow0$, where $\Bb$ is the Bures--Wasserstein metric of Proposition~\ref{prop-gaussian-w2-bures}.
\end{cor}
\begin{proof}
Proposition~\ref{prop-gaussian-sinkhorn-closed-form} writes the raw entropic value as the squared mean displacement plus trace terms and a spectral sum. With the notation above, the spectral part is exactly $\sum_i\psi_\epsilon(\sigma_i(\cov_\al,\cov_\be))$. Applying the same formula to the self-costs $(\al,\al)$ and $(\be,\be)$ replaces these singular values by the eigenvalues of $\cov_\al$ and $\cov_\be$. In the polarization formula~\eqref{eq-sinkhorn-divergence}, the trace terms cancel because
$\tr\cov_\al+\tr\cov_\be-\frac12(2\tr\cov_\al)-\frac12(2\tr\cov_\be)=0$, while the polarization of the squared mean terms leaves exactly $\norm{\mean_\al-\mean_\be}^2$. Since $\tau_\epsilon(r)\to1$ and $\epsilon\log(1-\tau_\epsilon(r)^2)\to0$, one has $\psi_\epsilon(r)\to-2r$. The limit is therefore
\(\tr\Sigma+\tr\Lambda -2\sum_i\sigma_i(\Sigma,\Lambda) = \Bb(\Sigma,\Lambda)^2,\)
which is the Bures formula~\eqref{eq-bures-defn}.
\end{proof}
\begin{prop}[One-dimensional Gaussian Sinkhorn rate]\label{prop-gaussian-sinkhorn-1d-rate}
Consider $\al=\be=\Gaussian(0,1)$ on $\RR$ with $c(x,y)=(x-y)^2$. Initialize $g^{(0)}=0$ and, after $\ell$ complete Sinkhorn cycles, write
\(g^{(\ell)}(y)=q^{(\ell)}y^2+\mathrm{cst}.\)
One soft transform maps a quadratic coefficient $q$ to
\(\mathsf Q_\epsilon(q) = 1-\frac{1}{1-q+\epsilon/2}, \qquad q<1+\epsilon/2,\)
so that $q^{(\ell+1)}=\mathsf Q_\epsilon(\mathsf Q_\epsilon(q^{(\ell)}))$. Set
\(A_\ell\eqdef 1-q^{(\ell)}+\frac{\epsilon}{2}, \qquad A_\star\eqdef 1-q_\star+\frac{\epsilon}{2} = \frac{\epsilon+\sqrt{\epsilon^2+16}}{4}.\)
Then $q^{(\ell)}\to q_\star=1+\epsilon/2-A_\star$, and the convergence obeys the exact cross-ratio identity
\(\frac{A_\ell-A_\star}{A_\ell+A_\star^{-1}} = A_\star^{-4\ell} \frac{A_0-A_\star}{A_0+A_\star^{-1}}.\)
In particular,
\[
0\leq q_\star-q^{(\ell)}
\leq
q_\star
\left(\frac{4}{\epsilon+\sqrt{\epsilon^2+16}}\right)^{4\ell}.
\]
\end{prop}
\begin{proof}
Completing the square in $\int\exp((q y^2-(x-y)^2)/\epsilon)\d\Gaussian(0,1)(y)$ gives $\mathsf Q_\epsilon(q)$. Write $a=\epsilon/2$. Applying this map twice shows that $A_\ell=1-q^{(\ell)}+a$ satisfies
\(A_{\ell+1}=M(A_\ell), \qquad M(A)\eqdef a+\frac{A}{1+aA}.\)
The fixed points of $M$ solve $A^2-aA-1=0$ and are $A_\star>0$ and $-A_\star^{-1}$. Direct algebra for this M\"obius map gives \(\frac{M(A)-A_\star}{M(A)+A_\star^{-1}} = A_\star^{-4} \frac{A-A_\star}{A+A_\star^{-1}},\)
which proves the exact identity by iteration. Moreover, $A_0=1+a>A_\star$ and $M$ leaves $[A_\star,A_0]$ invariant. On this interval, \(0<M'(A)=\frac{1}{(1+aA)^2} \leq \frac{1}{(1+aA_\star)^2} =A_\star^{-4},\)
where $1+aA_\star=A_\star^2$. Hence $0\leq A_\ell-A_\star\leq A_\star^{-4\ell}(A_0-A_\star)$. Since $A_0-A_\star=q_\star$ and $A_\ell-A_\star=q_\star-q^{(\ell)}$, this is the announced bound.
\end{proof}
\subsection{Continuous \texorpdfstring{$\varepsilon$}{epsilon}-Sinkhorn Flow}\label{sec-continuous-epsilon-sinkhorn}
\paragraph{Parabolic Monge--Amp\`ere limit.}
Take the quadratic torus cost $c(x,y)=d_{\TT^d}(x,y)^2/2$. In the zero-temperature limit, the soft transforms concentrate near the map $x\mapsto x+\nabla u(x)$, while the first Laplace correction records its Jacobian determinant. Rescaling the accumulated correction in iteration time leads to the following flow.

\begin{defn}[Continuous \texorpdfstring{$\varepsilon$}{epsilon}-Sinkhorn flow]\label{def-continuous-epsilon-sinkhorn}
Let $\al$ and $\be$ be probability measures on the unit-volume flat torus $\TT^d$ with smooth positive densities $\density{\al}=e^{-F}$ and $\density{\be}=e^{-G}$. A gauge-fixed potential $u_t:\TT^d\to\RR$, with $\int_{\TT^d}u_t\d x=0$, follows the continuous $\epsilon$-Sinkhorn flow if
\begin{equation}\label{eq-continuous-epsilon-sinkhorn-pde}
\partial_t u_t(x)
=
\log\det\bigl(\Id+\nabla^2u_t(x)\bigr)
-G\bigl(x+\nabla u_t(x)\bigr)
+F(x)-\bar r_t,
\end{equation}
where $x+\nabla u_t(x)$ is understood modulo $\ZZ^d$, $\Id+\nabla^2u_t\succ0$, and $\bar r_t$ is the unique scalar preserving the mean-zero gauge.
\end{defn}
\begin{prop}[Stationary continuous Sinkhorn potentials]\label{prop-continuous-sinkhorn-stationary}
Assume that a smooth stationary solution of~\eqref{eq-continuous-epsilon-sinkhorn-pde} satisfies $\Id+\nabla^2u\succ0$, and define $T(x)=x+\nabla u(x)$ modulo $\ZZ^d$. Then $T_\sharp\al=\be$. Conversely, any smooth map of this form satisfying $\Id+\nabla^2u\succ0$ and $T_\sharp\al=\be$ solves the stationary equation, up to the additive gauge of $u$.
\end{prop}
\begin{proof}
At stationarity, the right-hand side before subtracting $\bar r_t$ is a constant $r$, hence
\(e^{-G(T(x))}\det\nabla T(x)=e^{-F(x)}e^r.\)
The positive-Jacobian condition $\Id+\nabla^2u\succ0$ makes $T$ an orientation-preserving local diffeomorphism of the torus, homotopic to the identity and therefore of degree one. It is thus a diffeomorphism. Integrating the displayed identity and using that both measures have unit mass gives $e^r=1$. The remaining equality is precisely the change-of-variables formula for $T_\sharp(e^{-F}\d x)=e^{-G}\d x$. Reversing the argument proves the converse.
\end{proof}
\paragraph{Rescaled continuous Sinkhorn iterates.}
Let $\mathsf S_\epsilon$ denote one complete continuous dual Sinkhorn cycle~\eqref{eq-continuous-dual-sinkhorn-iteration}, written in the sign convention $u=-f$. Using the soft transforms of Definition~\ref{def-continuous-soft-c-transform},
\(\mathsf S_\epsilon u \eqdef -\Big((-u)^{c,\epsilon}\Big)^{\bar c,\epsilon}.\)
For a smooth potential with $\Id+\nabla^2u\succ0$, set
\(\mathcal R(u)(x) \eqdef \log\det(\Id+\nabla^2u(x)) -G(x+\nabla u(x))+F(x).\)
\begin{prop}[Conditional rescaled continuous Sinkhorn limit]\label{prop-scaled-log-sinkhorn-limit}
Let $u_\epsilon^{(0)}=u_0$ be smooth and mean zero, and generate
\(u_\epsilon^{(\ell+1)} = \mathsf S_\epsilon u_\epsilon^{(\ell)} - \int_{\TT^d}\mathsf S_\epsilon u_\epsilon^{(\ell)}\d x.\)
For $t\in[\ell\epsilon,(\ell+1)\epsilon]$, let
\(\widetilde u_\epsilon(t) \eqdef (1-\theta)u_\epsilon^{(\ell)} +\theta u_\epsilon^{(\ell+1)}, \qquad \theta\eqdef t/\epsilon-\ell.\)
Assume that, for every $T>0$, there are $\kappa_T>0$ and $\eta_\epsilon(T)\to0$ such that, whenever $\ell\epsilon\leq T$,
\[
\Id+\nabla^2u_\epsilon^{(\ell)}\succeq\kappa_T\Id,
\qquad
\left\|
u_\epsilon^{(\ell+1)}-u_\epsilon^{(\ell)}
-\epsilon\left(
\mathcal R(u_\epsilon^{(\ell)})
-\int_{\TT^d}\mathcal R(u_\epsilon^{(\ell)})\d x
\right)
\right\|_{C^0}
\leq
\epsilon\eta_\epsilon(T).
\]
If $\widetilde u_\epsilon\to u$ in $C([0,T];C^2(\TT^d))$ for every $T$, then $u\in C^1([0,T];C^0(\TT^d))$ and solves~\eqref{eq-continuous-epsilon-sinkhorn-pde}, with
\(
\bar r_t=\int_{\TT^d}\mathcal R(u_t)\d x
\).
\end{prop}
\begin{proof}
At a grid time $t_\epsilon=\epsilon\lfloor t/\epsilon\rfloor$, summing the normalized consistency estimate gives
\[
u_\epsilon^{(\lfloor t/\epsilon\rfloor)}-u_0
=
\epsilon\sum_{r<\lfloor t/\epsilon\rfloor}
\left[
\mathcal R(u_\epsilon^{(r)})
-\int_{\TT^d}\mathcal R(u_\epsilon^{(r)})\d x
\right]
+e_\epsilon(t),
\qquad
\norm{e_\epsilon(t)}_{C^0}\leq T\eta_\epsilon(T).
\]
Define the left-endpoint interpolation $\underline u_\epsilon(s)\eqdef u_\epsilon^{(\lfloor s/\epsilon\rfloor)}$. Uniform convergence of the affine interpolants to the continuous limit implies that $\underline u_\epsilon\to u$ uniformly in $C^2$. Together with the uniform Jacobian positivity and smoothness of $F,G$, this gives uniform convergence of the summand to
\(
\mathcal R(u_s)-\int_{\TT^d}\mathcal R(u_s)\d x
\).
Indeed, the sum in the preceding display is exactly
\[
\int_0^{t_\epsilon}
\left[
\mathcal R(\underline u_\epsilon(s))
-\int_{\TT^d}\mathcal R(\underline u_\epsilon(s))\d x
\right]\d s,
\]
and therefore converges to the corresponding integral along $u$. The affine remainder between $t_\epsilon$ and $t$ tends to zero, so
\[
u_t-u_0
=
\int_0^t
\left[
\mathcal R(u_s)-\int_{\TT^d}\mathcal R(u_s)\d x
\right]\d s.
\]
The integrand is continuous in $C^0$, hence $u$ is $C^1$ in time with values in $C^0$. Differentiating this identity gives~\eqref{eq-continuous-epsilon-sinkhorn-pde}.
\end{proof}
\paragraph{One-dimensional Gaussian closure.}
In one dimension, the flow reads \(\partial_tu_t(x) = \log\bigl(1+u_t''(x)\bigr) -G\bigl(x+u_t'(x)\bigr)+F(x)-\bar r_t.\)
Although it was introduced on the torus, the same formal equation applies on $\RR$ for confining Gaussian densities. Let $\al=\Gaussian(m_\al,\sigma_\al^2)$ and $\be=\Gaussian(m_\be,\sigma_\be^2)$. The equation is closed on quadratic potentials: up to the irrelevant additive constant, write
\(x+u_t'(x)=q_t+a_t(x-m_\al), \qquad a_t>0.\)
Substituting this ansatz into~\eqref{eq-continuous-epsilon-sinkhorn-pde} and matching the quadratic and linear coefficients gives the closed two-dimensional system
\eql{\label{eq-continuous-sinkhorn-gaussian-1d}
\dot a_t
=
\frac{1}{\sigma_\al^2}-\frac{a_t^2}{\sigma_\be^2},
\qquad
\dot q_t
=
-\frac{a_t}{\sigma_\be^2}(q_t-m_\be).
}

\section{Statistical Optimal Transport}
\label{sec-statistical-ot}
\subsection{Law of Large Numbers and Central Limit Theorem}
\label{sec-law-large-numbers-clt}
\label{sec-quantitative-clt}
\paragraph{Empirical laws.}
If \(X_1,\ldots,X_n\) are i.i.d.\ samples with common law \(\alpha\), the associated empirical measure is the random probability measure
\begin{equation}\label{eq-empirical-law-alpha-n}
\hat\alpha_n \eqdef \frac1n\sum_{i=1}^n\delta_{X_i}.
\end{equation}
The ordinary law of large numbers says that empirical averages converge to expectations. In measure language this means that \(\hat\alpha_n\) converges weakly toward \(\alpha\), because testing \(\hat\alpha_n\) against a bounded continuous function \(\varphi\) gives the sample average \(n^{-1}\sum_i\varphi(X_i)\).

\begin{prop}[Empirical law of large numbers in \texorpdfstring{$\Wass_p$}{Wp}]\label{prop-empirical-lln-wasserstein}
Let \((\Xx,d)\) be a Polish metric space, let \(p\geq1\), and let \(\alpha\in\Pp_p(\Xx)\). Let \((X_i)_{i\geq1}\) be i.i.d.\ random variables with law \(\alpha\), and define \(\hat\alpha_n\) by~\eqref{eq-empirical-law-alpha-n}. Then
\[
\hat\alpha_n\rightharpoonup\alpha
\qandq
\Wass_p(\hat\alpha_n,\alpha)\longrightarrow0
\]
almost surely. Moreover,
$\EE\,\Wass_p(\hat\alpha_n,\alpha)^p\longrightarrow0$.
\end{prop}
\begin{proof}
Fix a reference point \(x_0\in\Xx\), and write \(r(x)=d(x,x_0)\). Since \(\Xx\) is Polish, the weak topology on \(\Pp(\Xx)\) admits a countable convergence-determining class \((\varphi_k)_{k\geq1}\subset C_b(\Xx)\). For each fixed \(k\), the strong law of large numbers gives
$\int \varphi_k\d\hat\alpha_n = \frac1n\sum_{i=1}^n\varphi_k(X_i) \longrightarrow \EE\varphi_k(X_1) = \int\varphi_k\d\alpha$
almost surely. Intersecting the corresponding probability-one events over the countable set of indices gives, on a single event of probability one, convergence against every \(\varphi_k\). By the convergence-determining property, this is exactly \(\hat\alpha_n\rightharpoonup\alpha\).

The moment condition \(\alpha\in\Pp_p(\Xx)\) means \(\int r^p\d\alpha<+\infty\). Applying the strong law again to the integrable random variable \(r(X_1)^p\) gives
\(\int r^p\d\hat\alpha_n = \frac1n\sum_{i=1}^n r(X_i)^p \longrightarrow \int r^p\d\alpha\)
almost surely. Proposition~\ref{prop-wass-topology-polish} states that weak convergence plus convergence of the \(p\)-th moments is equivalent to \(\Wass_p\)-convergence on \(\Pp_p(\Xx)\). Hence \(\Wass_p(\hat\alpha_n,\alpha)\to0\) almost surely.

Set \(A_n=\int r^p\d\hat\alpha_n\) and \(M=\int r^p\d\alpha\). The triangle inequality through the Dirac mass \(\delta_{x_0}\), followed by \((a+b)^p\leq2^{p-1}(a^p+b^p)\), gives the deterministic bound
$\Wass_p(\hat\alpha_n,\alpha)^p \leq 2^{p-1}\big(A_n+M\big)$.
The family \((A_n)_n\) is uniformly integrable. Indeed, by the de la Vall\'ee--Poussin criterion, choose a convex superlinear function \(\Psi\) such that \(\EE\Psi(r(X_1)^p)<+\infty\); Jensen's inequality gives
$\EE\Psi(A_n) \leq \frac1n\sum_{i=1}^n \EE\Psi(r(X_i)^p) = \EE\Psi(r(X_1)^p)$.
Thus \((A_n+M)_n\), and hence \((\Wass_p(\hat\alpha_n,\alpha)^p)_n\), is uniformly integrable. Together with almost-sure convergence to zero, this implies \(\EE\Wass_p(\hat\alpha_n,\alpha)^p\to0\).
\end{proof}
\paragraph{Central-limit fluctuations.}
\begin{prop}[Berry--Esseen bound in $\Wass_1$]\label{prop-berry-esseen-w1}
Let $(X_i)_{i=1}^n$ be i.i.d.\ real random variables such that $\EE X_i=0$, $\EE X_i^2=1$ and $\EE |X_i|^3<+\infty$. If $\alpha_n$ is the law of $n^{-1/2}\sum_i X_i$ and $\gamma$ is the standard Gaussian law, then $\Wass_1(\alpha_n,\gamma)\leq C\,\EE |X_1|^3/\sqrt n$, where $C$ is a universal constant.
\end{prop}
\begin{proof}
By Kantorovich--Rubinstein duality,
$\Wass_1(\alpha_n,\gamma)=\sup_{\Lip(h)\leq1}\left|\EE h(S_n)-\EE h(G)\right|$, where $S_n=n^{-1/2}\sum_i X_i$ and $G\sim\gamma$.
For each such $h$, solve Stein's equation
\(f_h'(x)-x f_h(x)=h(x)-\EE h(G)\).
The solution satisfies \(\norm{f_h'}_\infty+\norm{f_h''}_\infty\leq C\) for a universal constant. Hence
\(\EE h(S_n)-\EE h(G)=\EE[f_h'(S_n)-S_nf_h(S_n)].\)
Write \(S_n^{(i)}=S_n-X_i/\sqrt n\). Independence and \(\EE X_i=0\) give
\[
\EE[S_nf_h(S_n)]
=
\frac1{\sqrt n}\sum_{i=1}^n
\EE\!\left[X_i\big(f_h(S_n^{(i)}+X_i/\sqrt n)-f_h(S_n^{(i)})\big)\right].
\]
Taylor's formula with integral remainder and \(\EE X_i^2=1\) control the difference between \(\EE[S_nf_h(S_n)]\) and \(n^{-1}\sum_i\EE f_h'(S_n^{(i)})\) by \(C\EE|X_1|^3/\sqrt n\). The Lipschitz bound on \(f_h'\) controls the difference between \(\EE f_h'(S_n)\) and the same average by \(C\EE|X_1|/\sqrt n\). Since \(\EE X_1^2=1\), the triangle inequality gives \(\left|\EE f_h'(S_n)-\EE[S_nf_h(S_n)]\right|\leq C\EE|X_1|^3/\sqrt n\). Taking the supremum over \(h\) proves the result. Sharper constants and higher-order transport-distance refinements are available, but are omitted from the compact handout.
\end{proof}
\begin{prop}[Sharp symmetric lattice asymptotic]\label{prop-sharp-lattice-w1-clt}
Let \(X\) be centered, symmetric, of unit variance, and supported on a finite subset of \(a+h\ZZ\), where \(h>0\) is the maximal lattice span. If \(\alpha_n\) is the law of \(n^{-1/2}\sum_{i=1}^nX_i\) and \(\gamma=\mathcal N(0,1)\), then
\(\sqrt n\,\Wass_1(\alpha_n,\gamma)\longrightarrow \frac h4.\)
\end{prop}
\begin{proof}
Denote by \(F_n\) the CDF of \(\alpha_n\), and by \(\Phi\) and \(\varphi\) the standard Gaussian CDF and density. With the centered periodic sawtooth \(\psi(u)=\frac12-\{u\}\), the integrated lattice Edgeworth expansion gives
\[
F_n(x)-\Phi(x)
=
\frac h{\sqrt n}\,
\psi\!\left(\frac{\sqrt n\,x-na}{h}\right)\varphi(x)+r_n(x),
\qquad
\norm{r_n}_{L^1(\RR)}=o(n^{-1/2});
\]
see,, and. Vallender's one-dimensional identity states that \(\Wass_1(\alpha_n,\gamma)=\int_\RR|F_n-\Phi|\). Therefore \(\big||u+v|-|u|\big|\leq|v|\) reduces the claim to
\[
\int_\RR
\left|\psi\!\left(\frac{\sqrt n\,x-na}{h}\right)\right|
\varphi(x)\d x
\longrightarrow
\int_0^1|\psi(u)|\d u
=
\frac14.
\]
The convergence is periodic averaging: prove it first for compactly supported step functions, where it is a Riemann sum over periods, and conclude by \(L^1\)-approximation of \(\varphi\).
\end{proof}
\begin{prop}[Sharp symmetric density asymptotic]\label{prop-sharp-density-w1-clt}
Let \(X\) be centered, symmetric, of unit variance, with a density and moments of every order. Set \(\kappa_4=\EE[X^4]-3\) and \(H_3(x)=x^3-3x\). If \(\alpha_n\) is the law of \(n^{-1/2}\sum_{i=1}^nX_i\) and \(\gamma=\mathcal N(0,1)\), then
\(\Wass_1(\alpha_n,\gamma) = \frac{|\kappa_4|}{24n} \int_\RR |H_3(x)|\varphi(x)\d x +o(n^{-1}).\)
\end{prop}
\begin{proof}
The non-lattice Edgeworth expansion through order \(n^{-1}\) has an \(n^{-1/2}\) term proportional to the third cumulant, and \(n^{-1}\) terms proportional to the fourth cumulant and to the square of the third one. Symmetry makes the third cumulant vanish and gives
\(F_n(x)-\Phi(x) =-\frac{\kappa_4}{24n}H_3(x)\varphi(x)+r_n(x), \qquad \norm{r_n}_{L^1(\RR)}=o(n^{-1});\)
see and. Inserting this expansion into Vallender's identity and using \(\big||u+v|-|u|\big|\leq|v|\) proves the formula.
\end{proof}
Since
\(\int_\RR |H_3(x)|\varphi(x)\d x =2\varphi(0)+8\varphi(\sqrt3) =\frac{2+8e^{-3/2}}{\sqrt{2\pi}},\)
one obtains the two sharp equivalents
\begin{equation}\label{eq-bernoulli-uniform-sharp-w1-clt}
\Wass_1(\alpha_n,\gamma)
\sim
\begin{cases}
\dfrac1{2\sqrt n},
& X\sim\frac12(\delta_{-1}+\delta_1),\\[2mm]
\dfrac{1+4e^{-3/2}}{10\sqrt{2\pi}}\dfrac1n,
& X\sim\operatorname{Unif}[-\sqrt3,\sqrt3].
\end{cases}
\end{equation}
\subsection{Sample Complexity}
\label{sec-sample-complexity}
\paragraph{Unregularized OT.}
For \(n\geq2\), \(p\geq1\), and \(d\geq1\), define
\begin{equation}\label{eq-empirical-wasserstein-scale}
r_{n,p,d}
\eqdef
\begin{cases}
n^{-1/(2p)}, & d<2p,\\
n^{-1/(2p)}(\log(1+n))^{1/p}, & d=2p,\\
n^{-1/d}, & d>2p.
\end{cases}
\end{equation}
\begin{prop}[Empirical OT has intrinsic-dimension value rates]\label{prop-empirical-ot-rate}
Let \(p\geq1\), let \(\Xx\subset\RR^d\) be compact, and, for each \(\alpha\in\Pp(\Xx)\), let \(\widehat\alpha_n=n^{-1}\sum_{i=1}^n\de_{X_i}\) for i.i.d.\ samples \(X_i\sim\alpha\). Then
\begin{equation}\label{eq-empirical-wasserstein-moment-scale}
\sup_{\alpha\in\Pp(\Xx)}
\EE\!\left[\Wass_p(\widehat\alpha_n,\alpha)^p\right]^{1/p}
\leq
C_{\Xx,p,d}\,r_{n,p,d}.
\end{equation}
\[
\EE\!\left[
\abs{\Wass_p(\hat\alpha_n,\hat\beta_m)-\Wass_p(\alpha,\beta)}^p
\right]^{1/p}
\lesssim_{p,d}
r_{n,p,d}+r_{m,p,d}.
\]
\end{prop}
\begin{prop}[Dyadic partition bound for \texorpdfstring{$\Wass_1$}{W1}]\label{prop-dyadic-partition-w1}
Let \(\mathcal Q_j\) be nested dyadic partitions of \([0,1]^d\) into half-open cubes of side length \(2^{-j}\), with cubes on the outer face closed along the boundary of \([0,1]^d\). For all \(\alpha,\beta\in\Pp([0,1]^d)\) and every integer \(J\geq0\),
\(\Wass_1(\alpha,\beta) \leq \sqrt d \sum_{j=0}^{J-1}2^{-j} \sum_{Q\in\mathcal Q_{j+1}}\abs{\alpha(Q)-\beta(Q)} + \sqrt d\,2^{-J}.\)
\end{prop}
\begin{proof}
The proof is constructive: one decomposes the two measures into pieces which can be coupled at different spatial scales. Write \(\Omega=[0,1]^d\). We build positive measures \((\alpha_j,\beta_j)_{j=0}^J\) such that \(\sum_{j=0}^J\alpha_j=\alpha\), \(\sum_{j=0}^J\beta_j=\beta\), and \(\alpha_j(Q)=\beta_j(Q)\) for every \(Q\in\mathcal Q_j\).
Assuming such a decomposition for a moment, each pair \((\alpha_j,\beta_j)\) can be coupled inside the cells of \(\mathcal Q_j\). Indeed, on every cell \(Q\in\mathcal Q_j\) with positive mass, take the product coupling between the normalized restrictions of \(\alpha_j\) and \(\beta_j\) to \(Q\), multiplied by the common mass \(\alpha_j(Q)=\beta_j(Q)\). Summing over cells and over \(j\) gives a coupling of \(\alpha\) and \(\beta\). Since the diameter of a cell in \(\mathcal Q_j\) is at most \(\sqrt d\,2^{-j}\), this coupling has cost bounded by
$\Wass_1(\alpha,\beta) \leq \sqrt d\sum_{j=0}^J2^{-j}\alpha_j(\Omega)$.

Start at the finest scale. For \(Q\in\mathcal Q_J\), keep the common part of the two masses by setting \(\alpha_J(A\cap Q)=\frac{\alpha(Q)\wedge\beta(Q)}{\alpha(Q)}\alpha(A\cap Q)\) and \(\beta_J(A\cap Q)=\frac{\alpha(Q)\wedge\beta(Q)}{\beta(Q)}\beta(A\cap Q)\), with the convention that the corresponding term is zero when the denominator vanishes. For \(j=J-1,\ldots,1\), assume that the pieces at all finer levels \(k>j\) have been constructed and define the residuals \(\alpha'_j=\alpha-\sum_{k=j+1}^J\alpha_k\) and \(\beta'_j=\beta-\sum_{k=j+1}^J\beta_k\).
On each \(Q\in\mathcal Q_j\), set \(\alpha_j\) and \(\beta_j\) equal to the common part of \(\alpha'_j\) and \(\beta'_j\) inside \(Q\), exactly as above. Finally set \(\alpha_0=\alpha-\sum_{k=1}^J\alpha_k\) and \(\beta_0=\beta-\sum_{k=1}^J\beta_k\). Each step removes only common residual mass, so all measures are positive, and by construction \(\alpha_j(Q)=\beta_j(Q)\) on all cells \(Q\in\mathcal Q_j\).

For notational consistency, set \(\alpha'_0=\alpha_0\), \(\beta'_0=\beta_0\), \(\alpha'_J=\alpha\), and \(\beta'_J=\beta\). Let \(0\leq j<J\) and \(R\in\mathcal Q_j\). Since all pieces constructed at levels finer than \(j+1\) have equal mass on every child \(Q\in\mathcal Q_{j+1}\), the residual imbalance on \(Q\) just before the matching at level \(j+1\) is the original imbalance, \(\alpha'_{j+1}(Q)-\beta'_{j+1}(Q)=\alpha(Q)-\beta(Q)\).
After removing the common part at level \(j+1\), the remaining source residual inside \(Q\) has mass
$\alpha'_j(Q) = \alpha'_{j+1}(Q)-\alpha_{j+1}(Q) = \big(\alpha'_{j+1}(Q)-\beta'_{j+1}(Q)\big)_+$.
Because \(\alpha_j\) is itself a common submeasure of the residual \(\alpha'_j\), the mass matched at level \(j\) inside \(R\) satisfies
\[
\alpha_j(R)=\beta_j(R)
\leq \alpha'_j(R)
=
\sum_{Q\subset R}
\big(\alpha'_{j+1}(Q)-\beta'_{j+1}(Q)\big)_+
\leq
\sum_{Q\subset R}
\abs{\alpha(Q)-\beta(Q)},
\]
where the sums are over \(Q\in\mathcal Q_{j+1}\). Summing over \(R\in\mathcal Q_j\) gives
$\alpha_j(\Omega) \leq \sum_{Q\in\mathcal Q_{j+1}}\abs{\alpha(Q)-\beta(Q)}$.
The last level satisfies \(\alpha_J(\Omega)\leq1\). Substituting these bounds into the cost estimate proves the claim.
\end{proof}
\begin{proof}[Proof of Proposition~\ref{prop-empirical-ot-rate}]
General empirical Wasserstein moment estimates prove the one-sample statement for arbitrary \(p\); the two-sample estimate then follows from the reverse triangle inequality and Minkowski's inequality.
By the triangle inequality, \(\abs{\Wass_1(\hat\alpha_n,\hat\beta_m)-\Wass_1(\alpha,\beta)} \leq \Wass_1(\hat\alpha_n,\alpha)+\Wass_1(\hat\beta_m,\beta).\)
It is therefore enough to bound \(\EE\Wass_1(\hat\alpha_n,\alpha)\). Apply Proposition~\ref{prop-dyadic-partition-w1} to \((\hat\alpha_n,\alpha)\). If \(Q\in\mathcal Q_{j+1}\), then \(n\hat\alpha_n(Q)\) is a binomial random variable with parameter \(\alpha(Q)\), and hence
\(\EE\abs{\hat\alpha_n(Q)-\alpha(Q)} \leq \sqrt{\EE\abs{\hat\alpha_n(Q)-\alpha(Q)}^2} \leq \sqrt{\frac{\alpha(Q)}{n}}.\)
Summing over the \(2^{(j+1)d}\) cells and using the Cauchy--Schwarz inequality gives
\[
\sum_{Q\in\mathcal Q_{j+1}}
\EE\abs{\hat\alpha_n(Q)-\alpha(Q)}
\leq
\frac{1}{\sqrt n}
\sum_{Q\in\mathcal Q_{j+1}}\sqrt{\alpha(Q)}
\leq
\frac{2^{(j+1)d/2}}{\sqrt n}.
\]
Therefore, for every \(J\geq0\),
$\EE\Wass_1(\hat\alpha_n,\alpha) \lesssim_d 2^{-J} + \frac1{\sqrt n} \sum_{j=0}^{J-1}2^{j(d/2-1)}$.
If \(d=1\), the sum is uniformly bounded and we choose \(J\) such that \(2^{-J}\lesssim n^{-1/2}\). If \(d=2\), the sum is \(J\), and choosing \(2^J\simeq n^{1/2}\) gives \((\log n)n^{-1/2}\). If \(d\geq3\), the sum is dominated by its last term; choosing \(2^J\simeq n^{1/d}\) gives \(2^{-J}+n^{-1/2}2^{J(d/2-1)}\simeq n^{-1/d}\).
This proves the displayed \(\Wass_1\) rates. The elementary dyadic estimate is not sharp for every law in the critical case \(d=2\). The proposition deliberately retains the distribution-free \((\log n)n^{-1/2}\) bound supplied by the dyadic proof.
\end{proof}
\paragraph{Lower bounds and minimax optimality.}
\begin{prop}[Minimax lower bound for bounded densities]\label{prop-minimax-lower-w1-density}
Let \(d\geq3\). Let \(\mathcal A\) be the class of probability measures on \([0,1]^d\) admitting a density \(\rho\) with respect to Lebesgue measure such that
$\frac12\leq \rho(x)\leq \frac32$ for a.e.\ $x\in[0,1]^d$.
For each \(n\geq1\), let \(\mathfrak E_n\) be the class of measurable estimators \(\widetilde\alpha_n:([0,1]^d)^n\to\Pp([0,1]^d)\).
There exists \(c_d>0\), depending only on \(d\), such that
\[
\inf_{\widetilde\alpha_n\in\mathfrak E_n}
\sup_{\alpha\in\mathcal A}
\EE_{(X_1,\ldots,X_n)\sim\alpha^{\otimes n}}
\left[
\Wass_1\big(\widetilde\alpha_n(X_1,\ldots,X_n),\alpha\big)
\right]
\geq
c_d\,n^{-1/d}.
\]
The expectation is over i.i.d.\ samples \(X_i\sim\alpha\). In particular, the high-dimensional rate of Proposition~\ref{prop-empirical-ot-rate} is minimax sharp for \(\Wass_1\), even when one allows arbitrary estimators rather than the empirical measure.
\end{prop}
\begin{proof}
We use a standard hypercube construction. Choose a smooth function \(\psi\) compactly supported in \((0,1)^d\), with \(\int\psi=0\), \(\norm{\psi}_\infty\leq1\), and whose positive and negative parts are separated in the first coordinate. There exists a compactly supported Lipschitz function \(\varphi\) such that \(\int\varphi\psi\d x>0\). We extend both functions by zero outside \((0,1)^d\), and absorb \(\Lip(\varphi)\) into the constants below.

Let \(m\) be an integer, \(h=1/m\), and partition \([0,1]^d\) into \(M=m^d\) cubes \(Q_k=z_k+h[0,1)^d\). Define the rescaled bumps
$\psi_k(x)=\psi\!\left(\frac{x-z_k}{h}\right)\mathbf 1_{Q_k}(x)$.
For \(\theta=(\theta_1,\ldots,\theta_M)\in\{-1,1\}^M\) and a fixed small \(\eta>0\), set
$\rho_\theta(x)=1+\eta\sum_{k=1}^M\theta_k\psi_k(x)$ and $\alpha_\theta=\rho_\theta\d x$.
For \(\eta\) small enough, all these densities belong to \(\mathcal A\), and they integrate to one because every \(\psi_k\) has zero mean.

If \(\theta\) and \(\theta'\) differ on a set \(S\) of coordinates, the dual formula for \(\Wass_1\) gives an additive lower bound. Indeed, rescale \(\varphi\) as
\(\varphi_k(x)=h\,\varphi\!\left(\frac{x-z_k}{h}\right)\).
Each \(\varphi_k\) is supported in a fixed strict subcube of \(Q_k\). After dividing by this constant, Kantorovich--Rubinstein duality and the choice \(s_k=\operatorname{sign}(\theta_k-\theta'_k)\) on \(S\) give \(\Wass_1(\alpha_\theta,\alpha_{\theta'})\geq c\eta h^{d+1}\mathrm{Ham}(\theta,\theta')\), where \(\mathrm{Ham}\) denotes Hamming distance and \(c>0\) depends only on the chosen bump.

Neighboring vertices of the hypercube are statistically close. If \(\theta^{[k]}\) is obtained from \(\theta\) by flipping only the \(k\)-th sign, then the one-sample Kullback--Leibler divergence is bounded by
$\KL(\alpha_\theta|\alpha_{\theta^{[k]}}) \leq C\eta^2 h^d$,
because the two densities differ only on \(Q_k\), the perturbation has size \(O(\eta)\), and the densities are uniformly bounded away from zero. For \(n\) independent samples the divergence is multiplied by \(n\). Choose \(m\) so that \(m^d\simeq n\). Then \(n h^d\) is bounded, and, by taking \(\eta\) sufficiently small, Pinsker's inequality (Theorem~\ref{thm-pinsker}) gives a uniform bound
$\TV\big(\alpha_\theta^{\otimes n},\alpha_{\theta^{[k]}}^{\otimes n}\big) \leq q<1$.
Assouad's lemma converts this \(M\)-dimensional hypercube into \(M\) binary testing problems. More precisely, if the loss between two vertices is bounded below by \(\delta\) times their Hamming distance, while every pair of one-coordinate neighboring experiments has total variation at most \(q<1\), then every estimator has worst-case expected loss bounded below by a universal constant times \(M\delta(1-q)\). Here \(\delta=c\eta h^{d+1}\), so the lemma gives
\[
\inf_{\widetilde\alpha_n\in\mathfrak E_n}
\sup_{\theta\in\{-1,1\}^M}
\EE_{\alpha_\theta^{\otimes n}}
\left[
\Wass_1\big(\widetilde\alpha_n(X_1,\ldots,X_n),\alpha_\theta\big)
\right]
\geq
c'\eta\,M h^{d+1}(1-q)
=
c'\eta(1-q)h
\simeq
n^{-1/d}.
\]
Decreasing the constant handles the finitely many small values of \(n\), and the claim follows.
\end{proof}
\paragraph{Leveraging smoothness.}
\begin{prop}[Smoothed plug-in rates]\label{prop-smooth-plugin-w1-rate}
Let \(d\geq3\), \(s>0\), and let \(\alpha=\density{\alpha}\d x\), \(\beta=\density{\beta}\d x\) be probability measures on the unit-volume flat torus \(\TT^d\), with \(\max\{\norm{\density{\alpha}}_{H^s},\norm{\density{\beta}}_{H^s}\}\leq L\). Choose an even function \(\chi\in\Cc^\infty(\RR^d;[0,1])\) such that \(\chi(\xi)=1\) for \(\norm{\xi}\leq1\) and \(\chi(\xi)=0\) for \(\norm{\xi}\geq2\). For \(0<h\leq1\), let \(\kappa_h\) be the periodic kernel with Fourier coefficients \(\widehat\kappa_h(k)=\chi(hk)\), and define \(\tilde\alpha_{n,h}:=\hat\alpha_n\ast\kappa_h\) and \(\tilde\beta_{m,h}:=\hat\beta_m\ast\kappa_h\). These are signed measures of total mass one. Then
\[
\EE\norm{\tilde\alpha_{n,h}-\alpha}_{W_1}
\leq
C\left(h^{s+1}+n^{-1/2}h^{1-d/2}\right).
\]
The analogous bound holds for \(\tilde\beta_{m,h}-\beta\). Define the signed plug-in estimator
\(\widehat{\Wass}_{1,n,m} \eqdef \norm{\tilde\alpha_{n,h_n}-\tilde\beta_{m,h_m}}_{W_1}.\)
Choosing \(h_n\simeq n^{-1/(d+2s)}\) and \(h_m\simeq m^{-1/(d+2s)}\) gives
\[
\EE\abs{
\widehat{\Wass}_{1,n,m}
-
\Wass_1(\alpha,\beta)}
\leq
C\left(n^{-\frac{s+1}{d+2s}}+m^{-\frac{s+1}{d+2s}}\right).
\]
The constant depends only on \(d,s,L\) and \(\chi\).
\end{prop}
\begin{proof}
It suffices to prove the one-sample bound. For a zero-mean function \(w\) on the torus, set
\(\norm{w}_{\dot H^{-1}(\TT^d)}^2 \eqdef \sum_{k\in\ZZ^d\setminus\{0\}} \frac{\abs{\widehat w(k)}^2}{4\pi^2\norm{k}^2}.\)
Kantorovich--Rubinstein duality and Cauchy--Schwarz give \(\norm{w\d x}_{W_1}\lesssim\norm{w}_{\dot H^{-1}(\TT^d)}\), because a \(1\)-Lipschitz test function has \(L^2\)-gradient norm at most one on the unit-volume torus. Moreover, \(\EE\tilde\alpha_{n,h}=(\density{\alpha}\ast\kappa_h)\d x\), so the error splits into a deterministic smoothing bias and a centered empirical fluctuation.

Since \(\chi\) equals one near the origin, for every \(s>0\) one has \(\abs{1-\chi(hk)}\lesssim_s(h\norm{k})^{s+1}\). Fourier estimates and Parseval's identity therefore give
\[
\begin{aligned}
\norm{\density{\alpha}\ast\kappa_h-\density{\alpha}}_{\dot H^{-1}}
&\lesssim h^{s+1}\norm{\density{\alpha}}_{H^s},\\
\EE\norm{\density{\tilde\alpha_{n,h}}-\density{\alpha}\ast\kappa_h}_{\dot H^{-1}}^2
&\lesssim \frac1n\sum_{k\neq0}\frac{\abs{\chi(hk)}^2}{\norm{k}^2}
\lesssim \frac{h^{2-d}}n.
\end{aligned}
\]
In the second line, the support of \(\chi\) restricts the sum to \(0<\norm{k}\leq2/h\), and \(\sum_{0<\norm{k}\leq2/h}\norm{k}^{-2}\lesssim h^{2-d}\) because \(d\geq3\). The triangle and Jensen inequalities give the one-sample bound. Equating its two terms yields \(h\simeq n^{-1/(d+2s)}\). Finally, the reverse triangle inequality for the norm \(\norm{\cdot}_{W_1}\) gives the two-sample claim.
\end{proof}
\paragraph{MMD.}
MMD contains no transport optimization: its square is a combination of expectations of kernel evaluations, estimated empirically by sums of kernel values. Ordinary Monte-Carlo averaging therefore gives the dimension-free parametric rate below for bounded kernels.

\begin{prop}[MMD has a parametric value rate]\label{prop-mmd-sample-rate}
Let $k$ be a bounded positive definite kernel with $k(x,x)\leq\kappa^2$, and let $\MMD_k$ be the discrepancy of Definition~\ref{def-kernel-mmd-norm}. If $\hat\alpha_n$ and $\hat\beta_m$ are independent empirical measures, then
\[
\EE\abs{\MMD_k(\hat\alpha_n,\hat\beta_m)-\MMD_k(\alpha,\beta)}
\leq
\kappa\left(\frac1{\sqrt n}+\frac1{\sqrt m}\right).
\]
\end{prop}
\begin{proof}
Let $\Phi(x)=k(x,\cdot)$ be the feature map and $m_\alpha=\EE\Phi(X)$. The reverse triangle inequality for the RKHS norm gives
\(\abs{\MMD_k(\hat\alpha_n,\hat\beta_m)-\MMD_k(\alpha,\beta)} \leq \MMD_k(\hat\alpha_n,\alpha)+\MMD_k(\hat\beta_m,\beta).\)
Independence cancels the cross terms after taking the squared norm and expectation, giving
\(\EE\MMD_k(\hat\alpha_n,\alpha)^2 =\frac1n\EE\norm{\Phi(X)-m_\alpha}_{\RKHS_k}^2 =\frac{1}{n}\Big(\EE k(X,X)-\EE k(X,X')\Big).\)
The same estimate applies to $\hat\beta_m$, and Jensen's inequality together with $k(x,x)\leq\kappa^2$ gives the displayed bound.
\end{proof}
\paragraph{Entropic OT.}
\begin{prop}[Fixed-temperature Sinkhorn divergence has a parametric rate]\label{prop-sinkhorn-sample-rate}
Let \(c(x,y)=\norm{x-y}^2/2\), and assume that \(\alpha\) and \(\beta\) are \(\sigma^2\)-subgaussian probability measures on \(\RR^d\), in the sense that
$\int e^{\norm{x}^2/(2d\sigma^2)}\d\alpha(x)\leq2$ and $\int e^{\norm{y}^2/(2d\sigma^2)}\d\beta(y)\leq2$.
Set
$q_d\eqdef\left\lceil\frac{5d}{2}\right\rceil+6$ and $\Lambda_{d,\sigma}(\epsilon) \eqdef \epsilon\left[1+\left(\frac{\sigma}{\sqrt\epsilon}\right)^{q_d}\right]$.
Let \(\hat\alpha_n\) and \(\hat\beta_m\) be independent empirical measures. Then, for every \(\epsilon>0\),
\[
\EE\abs{\bar\MK_\c^\epsilon(\hat\alpha_n,\hat\beta_m)-\bar\MK_\c^\epsilon(\alpha,\beta)}
\leq
C_d\,\Lambda_{d,\sigma}(\epsilon)
\left(\frac1{\sqrt n}+\frac1{\sqrt m}\right).
\]
\end{prop}
\begin{proof}
Adding a separable term \(a(x)+b(y)\) to the cost adds
$\int a\d\mu+\int b\d\nu$ to $\MK_c^\epsilon(\mu,\nu)$, because every feasible coupling has marginals $(\mu,\nu)$. When the same function is added on both sides, $b=a$, these marginal terms cancel in the Sinkhorn divergence. Since
$c(x,y)=\norm{x-y}^2/2=-\dotp{x}{y}+\norm{x}^2/2+\norm{y}^2/2$, we may therefore work throughout the proof with the bilinear cost $c(x,y)=-\dotp{x}{y}$. We first estimate its raw entropic value and then pass to the debiased expression.

For the measures \(\alpha,\beta\) in the statement, write the dual objective, following~\eqref{eq-dual-sinkhorn-objective}, as
\[
\Dd_\epsilon^{\alpha,\beta}(u,v)
\eqdef
\int u\d\alpha+\int v\d\beta
-\epsilon\iint
\left(
\exp\left(\frac{u(x)+v(y)-c(x,y)}{\epsilon}\right)-1
\right)
\d\alpha(x)\d\beta(y).
\]
One can choose normalized optimal potentials for which the Sinkhorn equations hold pointwise on \(\RR^d\). Thus, at an optimal pair \((u,v)\),
$\int \exp\left(\frac{u(x)+v(y)-c(x,y)}{\epsilon}\right) \d\beta(y) =1$ for all $x$.
Equivalently, the potentials are fixed points of the soft \(c\)-transform, but only the dual normalization is needed below.
Let \((u,v)\) and \((u_n,v_n)\) be optimal dual potentials for \((\alpha,\beta)\) and \((\hat\alpha_n,\beta)\). By optimality,
$\Dd_\epsilon^{\alpha,\beta}(u_n,v_n) \leq \Dd_\epsilon^{\alpha,\beta}(u,v)$ and $\Dd_\epsilon^{\hat\alpha_n,\beta}(u,v) \leq \Dd_\epsilon^{\hat\alpha_n,\beta}(u_n,v_n)$.
Hence the difference of the two optimal values is sandwiched by evaluating the two objectives at these two pairs of potentials. For the pair \((u,v)\) associated with \((\alpha,\beta)\), the preceding Sinkhorn equation cancels the exponential term in the change of the first marginal:
\(\Dd_\epsilon^{\alpha,\beta}(u,v) - \Dd_\epsilon^{\hat\alpha_n,\beta}(u,v) = \int u\d(\alpha-\hat\alpha_n).\)
The same identity holds for \((u_n,v_n)\), using its normalization with respect to \(\beta\).
\[
\abs{\MK_\c^\epsilon(\hat\alpha_n,\beta)-\MK_\c^\epsilon(\alpha,\beta)}
\leq
\sup_{w\in\mathcal F_\epsilon(\alpha,\hat\alpha_n;\beta)}
\abs{\int w\d(\hat\alpha_n-\alpha)},
\]
where \(\mathcal F_\epsilon(\alpha,\hat\alpha_n;\beta)\) is the union of the normalized first-marginal potentials for the pairs \((\alpha,\beta)\) and \((\hat\alpha_n,\beta)\). This pathwise class is random; it will be controlled through the subgaussian envelopes below.

Set $s_0=\lceil d/2\rceil+1$, and let $\widetilde\sigma_n$ be the smallest common subgaussian proxy for $\alpha$, $\beta$, and $\hat\alpha_n$. This random variable is finite almost surely and satisfies $\EE\widetilde\sigma_n^{2k}\leq C_k\sigma^{2k}$ for every positive integer $k$. For $\tau\geq0$, let $\mathcal F_{1,\tau}$ denote the class of normalized first-marginal potentials associated with pairs of $\tau^2$-subgaussian measures at unit temperature. The regularity estimates of Mena--Niles-Weed show that every $w\in\mathcal F_{1,\tau}$ satisfies $w/(1+\tau^{3s_0})\in\mathcal F^{s_0}$, where $\mathcal F^{s_0}$ is a fixed polynomial-growth $C^{s_0}$ class whose elements $h$ obey
\[
\abs{h(x)}\leq C_{s_0,d}\bigl(1+\norm{x}^2\bigr),
\qquad
\abs{D^\kappa h(x)}\leq C_{s_0,d}\bigl(1+\norm{x}^{s_0}\bigr),
\quad 1\leq|\kappa|\leq s_0.
\]
In particular, the preceding random pathwise class is contained in $\mathcal F_{1,\widetilde\sigma_n}$. For a random envelope $L$ controlled by the empirical subgaussian moment, the empirical $L_2$-covering numbers of $\mathcal F^{s_0}$ satisfy
\(\log N(\delta,\mathcal F^{s_0},L_2(\hat\alpha_n)) \leq C_d\,L^{d/(2s_0)}\delta^{-d/s_0}(1+\sigma^{2d}),\)
with $\EE L\leq2$. Since $d/(2s_0)<1$, Dudley's entropy integral is finite. Rademacher symmetrization, the entropy-integral chaining bound, Cauchy--Schwarz, and the moment bounds for $\widetilde\sigma_n$ therefore yield
\(\EE\sup_{w\in\mathcal F_{1,\widetilde\sigma_n}} \abs{\int w\d(\hat\alpha_n-\alpha)} \leq C_d(1+\sigma^{q_d})\,n^{-1/2}.\)

The general \(\epsilon\) case follows by scaling. Set \(\alpha^\epsilon=(D_{1/\sqrt\epsilon})_\sharp\alpha\) and \(\beta^\epsilon=(D_{1/\sqrt\epsilon})_\sharp\beta\). Then
$\MK_\c^\epsilon(\alpha,\beta) = \epsilon\,\MK_\c^1(\alpha^\epsilon,\beta^\epsilon)$.
The push-forward of a \(\sigma^2\)-subgaussian law by \(D_{1/\sqrt\epsilon}\) is \((\sigma^2/\epsilon)\)-subgaussian, so the standard-deviation proxy \(\sigma\) is replaced by \(\sigma/\sqrt\epsilon\). Keeping this half-integer power avoids an unjustified rounding when $q_d$ is odd and $\epsilon>1$; for $0<\epsilon\leq1$ it can of course be bounded by an integer power of $1/\epsilon$.
Changing both marginals is handled by the triangle inequality,
\[
\abs{\MK_\c^\epsilon(\hat\alpha_n,\hat\beta_m)-\MK_\c^\epsilon(\alpha,\beta)}
\leq
\abs{\MK_\c^\epsilon(\hat\alpha_n,\hat\beta_m)-\MK_\c^\epsilon(\alpha,\hat\beta_m)}
+
\abs{\MK_\c^\epsilon(\alpha,\hat\beta_m)-\MK_\c^\epsilon(\alpha,\beta)}.
\]
Empirical measures drawn from a subgaussian law have such a random proxy, with moments controlled by the population proxy. Integrating over this proxy gives the two-sample extension. The argument does not require the two empirical marginals to have the same sample size, so the two contributions scale as \(n^{-1/2}\) and \(m^{-1/2}\).

Finally, use the definition
$\bar\MK_\c^\epsilon(\alpha,\beta) = \MK_\c^\epsilon(\alpha,\beta) -\frac12\MK_\c^\epsilon(\alpha,\alpha) -\frac12\MK_\c^\epsilon(\beta,\beta)$.
The cross term is controlled by the two-marginal estimate. For the self term, use the telescoping inequality
\[
\abs{\MK_\c^\epsilon(\hat\alpha_n,\hat\alpha_n)-\MK_\c^\epsilon(\alpha,\alpha)}
\leq
\abs{\MK_\c^\epsilon(\hat\alpha_n,\hat\alpha_n)-\MK_\c^\epsilon(\alpha,\hat\alpha_n)}
+
\abs{\MK_\c^\epsilon(\alpha,\hat\alpha_n)-\MK_\c^\epsilon(\alpha,\alpha)}.
\]
Each term is a one-marginal perturbation in the same pathwise sense: the other marginal may itself be empirical, but it is part of the common random subgaussian class. Thus the self term
\(\MK_\c^\epsilon(\hat\alpha_n,\hat\alpha_n)-\MK_\c^\epsilon(\alpha,\alpha)\)
is bounded at the same order as the cross term involving \(\alpha\), and similarly for \(\beta\). Summing these three contributions gives the displayed estimate, up to changing \(C_d\).
\end{proof}
\paragraph{Sample complexity of estimating OT maps.}
If \(\gD_{n,m}^\epsilon\) is the target-side dual potential, and if \(\bar c(y,x)\eqdef c_2(x,y)\), define
\[
u_{n,m}^\epsilon(x)
\eqdef
(\gD_{n,m}^\epsilon)^{\bar c,\epsilon}(x)
=
-\epsilon\log\sum_j b_j
\exp\!\left(\frac{\gD_{n,m,j}^\epsilon-c_2(x,Y_j)}{\epsilon}\right),
\qquad
T_{n,m}^\epsilon(x)\eqdef x-\nabla u_{n,m}^\epsilon(x).
\]
With the sign convention used here, the associated soft convex potential is
\[
\varphi_{n,m}^\epsilon(x)
\eqdef
\frac12\norm{x}^2-u_{n,m}^\epsilon(x)
=
\epsilon\log\sum_j b_j
\exp\!\left(
\frac{\dotp{x}{Y_j}+\gD_{n,m,j}^\epsilon-\frac12\norm{Y_j}^2}{\epsilon}
\right),
\qquad
T_{n,m}^\epsilon=\nabla\varphi_{n,m}^\epsilon.
\]
Differentiating the log-sum-exp gives the explicit half-squared formula
\begin{equation}\label{eq-empirical-entropic-map-barycentric}
T_{n,m}^\epsilon(x)
=
\sum_j \omega_j^\epsilon(x)Y_j,
\qquad
\omega_j^\epsilon(x)
=
\frac{
b_j\exp\!\left((\gD_{n,m,j}^\epsilon-c_2(x,Y_j))/\epsilon\right)
}{
\sum_k b_k\exp\!\left((\gD_{n,m,k}^\epsilon-c_2(x,Y_k))/\epsilon\right)
}.
\end{equation}
At a sampled source point \(X_i\), the Sinkhorn normalization gives \(\omega_j^\epsilon(X_i)=\P_{i,j}^\epsilon/a_i\) and \(T_{n,m}^\epsilon(X_i)=a_i^{-1}\sum_j \P_{i,j}^\epsilon Y_j\). Thus the extrapolated soft-transform map coincides on the empirical support with the barycentric projection~\eqref{eq-barycentric-projection} of the optimal entropic plan.

\begin{prop}[Consistency of empirical barycentric maps]\label{prop-empirical-barycentric-map-consistency}
Let \(\Omega\subset\RR^d\) be compact, let \(\alpha,\beta\in\Pp_2(\Omega)\), and assume that \(\alpha\) is absolutely continuous. Let \(\phi\) be a Brenier potential from \(\alpha\) to \(\beta\) for the cost \(c_2(x,y)=\norm{x-y}^2/2\), and choose a bounded Borel selection \(T(x)\in\partial\phi(x)\) on \(\Omega\). At every differentiability point of \(\phi\), this selection necessarily equals \(\nabla\phi(x)\), so it is a representative of the unique Brenier map. Let \(X_i,Y_j\in\Omega\), and assume that \(\hat\alpha_n=n^{-1}\sum_{i=1}^n\delta_{X_i}\to\alpha\) and \(\hat\beta_m=m^{-1}\sum_{j=1}^m\delta_{Y_j}\to\beta\) in \(\Wass_2\). Let \(\epsilon_{n,m}\geq0\) satisfy \(\epsilon_{n,m}\to0\) and \(\epsilon_{n,m}\log(\min\{n,m\})\to0\). For each \(n,m\), let \(\P_{n,m}^{\epsilon}\) be an optimal coupling matrix between \(\hat\alpha_n\) and \(\hat\beta_m\), with \(\epsilon=\epsilon_{n,m}\): unregularized if \(\epsilon=0\), entropic if \(\epsilon>0\). Form the aggregated plan \(\pi_{n,m}^\epsilon=\sum_{i,j}\P_{i,j}^\epsilon\delta_{(X_i,Y_j)}\), disintegrate it as \(\pi_{n,m}^\epsilon(\d x,\d y)=\hat\alpha_n(\d x)\pi_{n,m,x}^\epsilon(\d y)\), and define
$\bar T_{n,m}^{\epsilon}(x) \eqdef \int y\,\d\pi_{n,m,x}^\epsilon(y)$.
Equivalently, if \(I_x\eqdef\{i:X_i=x\}\), then
\(\bar T_{n,m}^{\epsilon}(x) = \frac{n}{|I_x|}\sum_{i\in I_x}\sum_j\P_{i,j}^{\epsilon}Y_j.\)
When the source sites are distinct this reduces to \(\bar T_{n,m}^{\epsilon}(X_i)=n\sum_j\P_{i,j}^{\epsilon}Y_j\). The aggregated definition remains single-valued when sites are repeated and satisfies
\(\int \norm{\bar T_{n,m}^{\epsilon}(x)-T(x)}^2 \d\hat\alpha_n(x) = \frac1n\sum_i \norm{\bar T_{n,m}^{\epsilon}(X_i)-T(X_i)}^2 \longrightarrow0\)
as \(n,m\to+\infty\).
If \(\beta\) is also absolutely continuous, the analogous backward barycentric projections converge in \(L^2(\hat\beta_m)\) to the Brenier map \(S:\beta\to\alpha\).
\end{prop}
\begin{proof}
The assumptions imply \(\MK_{c_2}(\hat\alpha_n,\hat\beta_m)\to\MK_{c_2}(\alpha,\beta)\). If \(\epsilon=0\), \(\P_{n,m}^{\epsilon}\) is exactly optimal. If \(\epsilon>0\), compare its entropic objective with an unregularized optimal empirical plan \(Q_{n,m}\). Regarding \(Q_{n,m}\) as the joint law of two uniform discrete indices, its relative entropy with respect to their product law is their mutual information. Hence
$\KLD(Q_{n,m}|\hat\alpha_n\otimes\hat\beta_m) \leq \min\{\log n,\log m\}$.
It follows that
$\int c_2\,\d \pi_{n,m}^{\epsilon} \leq \MK_{c_2}(\hat\alpha_n,\hat\beta_m) +\epsilon\log(\min\{n,m\})$.
Hence the transport costs of \(\pi_{n,m}^{\epsilon}\) converge to the population optimum. By compactness, every weak limit point of these aggregated plans has marginals \(\alpha,\beta\) and optimal quadratic cost. Brenier uniqueness therefore forces the whole sequence to converge weakly to \((\Id,T)_\sharp\alpha\).

Disintegrate the aggregated plan \(\pi_{n,m}^{\epsilon}\) with respect to its first marginal and denote its barycentric projection by \(\bar T_{n,m}^{\epsilon}\). Jensen's inequality gives
\(\int \norm{\bar T_{n,m}^{\epsilon}(x)-T(x)}^2\d\hat\alpha_n(x) \leq \iint \norm{y-T(x)}^2\d \pi_{n,m}^{\epsilon}(x,y).\)
At every differentiability point of \(\phi\), the subdifferential is the singleton \(\{\nabla\phi(x)\}\). Since the target is compact, the relevant subgradients are bounded, and the closed-graph property of \(\partial\phi\) implies continuity of the selected gradient at such a point. Convex functions are differentiable Lebesgue-almost everywhere; because \(\alpha\) is absolutely continuous, the bounded function \((x,y)\mapsto\norm{y-T(x)}^2\) is continuous \((\Id,T)_\sharp\alpha\)-almost everywhere. The portmanteau theorem for bounded functions whose discontinuity set has zero limiting mass gives
\(\iint \norm{y-T(x)}^2\d \pi_{n,m}^{\epsilon}(x,y) \longrightarrow \iint \norm{y-T(x)}^2\d[(\Id,T)_\sharp\alpha](x,y)=0.\)
The proof for the backward map is identical when the optimal plan is also induced by the Brenier map from \(\beta\) to \(\alpha\).
\end{proof}
\begin{prop}[Finite-sample rate for the entropic map]\label{prop-entropic-map-rate}
Assume that \(\alpha\) and \(\beta\) have densities on a common compact set \(\Omega\), both bounded above and with \(\density{\beta}\) bounded below away from zero. Let \(T=\nabla\phi\) be the quadratic Brenier map from \(\alpha\) to \(\beta\), and assume that \(\phi\in C^2(\Omega)\), its Legendre conjugate satisfies \(\phi^*\in C^{s+1}(\Omega)\) for some \(s>1\), and
\(\mu\Id_d\preceq\nabla^2\phi(x)\preceq L\Id_d \qquad (x\in\Omega)\)
for constants \(0<\mu\leq L\). Set \(\bar s\eqdef s\wedge3\), assume that \(\mathcal I_{\rm geo}(\alpha,\beta)<+\infty\) as in Proposition~\ref{prop-small-epsilon-expansion}, and draw \(n\) independent samples from each measure. If \(\epsilon_n\simeq n^{-1/(d+\bar s+1)}\), then the estimator \(T_{n,n}^{\epsilon_n}\) defined in~\eqref{eq-empirical-entropic-map-barycentric} satisfies
\[
\EE\norm{T_{n,n}^{\epsilon_n}-T}_{L^2(\alpha)}^2
\lesssim
\bigl(1+\mathcal I_{\rm geo}(\alpha,\beta)\bigr)
n^{-\frac{\bar s+1}{2(d+\bar s+1)}}\log n.
\]
\end{prop}
\paragraph{Sliced Wasserstein.}
For a direction \(\theta\in\Sphere^{d-1}\), let \(P_\theta(x)=\dotp{\theta}{x}\) be the one-dimensional projection and let \(\sigma\) denote the normalized surface measure on the sphere. For \(\alpha,\beta\in\Pp_p(\RR^d)\), the sliced Wasserstein distance is
\(\SW_p(\alpha,\beta)^p \eqdef \int_{\Sphere^{d-1}} \Wass_p\big((P_\theta)_\sharp\alpha,(P_\theta)_\sharp\beta\big)^p \d\sigma(\theta).\)
\begin{thm}[Dimension-free empirical rate for \texorpdfstring{$\SW_1$}{sliced Wasserstein-1}]\label{thm-sliced-sample-complexity}
Let \(\alpha,\beta\in\Pp(\RR^d)\) satisfy
\(\supp(\alpha)\cup\supp(\beta)\subset B(0,R)\), and let
\(\hat\alpha_n\) and \(\hat\beta_m\) be their empirical measures from \(n\) and \(m\) i.i.d.\ samples, respectively. Then
\[
\EE\abs{\SW_1(\hat\alpha_n,\hat\beta_m)-\SW_1(\alpha,\beta)}
\leq
R\left(\frac{1}{\sqrt n}+\frac{1}{\sqrt m}\right).
\]
\end{thm}
\begin{proof}
The triangle inequality for \(\SW_1\) gives \(\abs{\SW_1(\hat\alpha_n,\hat\beta_m)-\SW_1(\alpha,\beta)} \leq \SW_1(\hat\alpha_n,\alpha)+\SW_1(\hat\beta_m,\beta).\)
Fix \(\theta\in\Sphere^{d-1}\), set \(\alpha_\theta=(P_\theta)_\sharp\alpha\), and denote the cumulative functions of \(\alpha_\theta\) and its empirical measure by \(F_\theta\) and \(\hat F_{\theta,n}\). Both measures are supported on \([-R,R]\). The one-dimensional identity~\eqref{eq-w1-1d} and the Cauchy--Schwarz inequality give
\[
\EE\,\Wass_1\big((P_\theta)_\sharp\hat\alpha_n,\alpha_\theta\big)
=
\int_{-R}^R \EE\abs{\hat F_{\theta,n}(s)-F_\theta(s)}\d s
\leq
\int_{-R}^R \sqrt{\frac{F_\theta(s)(1-F_\theta(s))}{n}}\d s
\leq
\frac{R}{\sqrt n}.
\]
Indeed, \(n\hat F_{\theta,n}(s)\) is binomial with success probability \(F_\theta(s)\). Integrating this bound over \(\theta\) yields \(\EE\,\SW_1(\hat\alpha_n,\alpha)\leq R/\sqrt n\). The same argument for \(\beta\), followed by the preceding triangle inequality, proves the result.
\end{proof}
\subsection{Bias and Variance of OT}
\label{sec-bias-variance-ot}
\paragraph{Bias, variance and approximation errors.}
We now ask for a finer statistical description of the KL-normalized entropic OT value \(\MK_c^\epsilon\) from Definition~\ref{def-continuous-entropic-ot}, with the convention
\(\MK_c^0\eqdef\MK_c\) for unregularized OT\@. Fixing \(c\), \(\alpha\) and \(\beta\), define the statistical bias
\(B_{n,m}^\epsilon \eqdef \EE \MK_c^\epsilon(\hat\alpha_n,\hat\beta_m) -\MK_c^\epsilon(\alpha,\beta)\)
and the centered fluctuation
\(Z_{n,m}^\epsilon \eqdef \MK_c^\epsilon(\hat\alpha_n,\hat\beta_m) -\EE \MK_c^\epsilon(\hat\alpha_n,\hat\beta_m).\)
If the target is the exact cost but the statistic is entropically regularized, these quantities isolate the three sources of error:
\[
\MK_c^\epsilon(\hat\alpha_n,\hat\beta_m)-\MK_c^0(\alpha,\beta)
=
\underbrace{\MK_c^\epsilon(\alpha,\beta)-\MK_c^0(\alpha,\beta)}_{\text{regularization bias}}
+
\underbrace{B_{n,m}^\epsilon}_{\text{statistical bias}}
+
\underbrace{Z_{n,m}^\epsilon}_{\text{centered fluctuation}}.
\]
For the debiased Sinkhorn divergence, define \(\bar B_{n,m}^\epsilon\) and \(\bar Z_{n,m}^\epsilon\) by replacing \(\MK_c^\epsilon\) with \(\bar\MK_c^\epsilon\) above; the same decomposition then holds with bars. It is usually read after choosing a temperature \(\epsilon=\epsilon_n\).

\begin{prop}[Finite-space bias and CLT for OT]\label{prop-finite-ot-clt}
Let \(\alpha=\sum_{i=1}^{N}a_i\delta_{x_i}\) and \(\beta=\sum_{j=1}^{M}b_j\delta_{y_j}\), where all weights are positive and \(c\) is finite on the product support. Let \(\hat\alpha_n\) and \(\hat\beta_n\) be the independent empirical laws of \(n\) i.i.d.\ samples from \(\alpha\) and \(\beta\), respectively. Fix \(\epsilon\geq0\), with \(\MK_c^0=\MK_c\), and let \((f_\epsilon^\star,g_\epsilon^\star)\) be optimal dual potentials for \(\MK_c^\epsilon(\alpha,\beta)\). When \(\epsilon=0\), assume that this pair is unique up to the additive gauge; for \(\epsilon>0\), this uniqueness follows after fixing the gauge. Define
\[
\mathsf{v}_\epsilon
\eqdef
\sum_{i=1}^{N}a_i\bigl(f_\epsilon^\star(x_i)-\bar f_\epsilon\bigr)^2
+
\sum_{j=1}^{M}b_j\bigl(g_\epsilon^\star(y_j)-\bar g_\epsilon\bigr)^2,
\]
where \(\bar f_\epsilon=\sum_i a_i f_\epsilon^\star(x_i)\) and \(\bar g_\epsilon=\sum_j b_j g_\epsilon^\star(y_j)\). Then, as \(n\to+\infty\) with \(N,M\) fixed,
\[
\sqrt n\,B_{n,n}^\epsilon\longrightarrow0,
\qquad
\sqrt n\,Z_{n,n}^\epsilon\overset{\mathrm{law}}{\longrightarrow}\Gaussian(0,\mathsf{v}_\epsilon),
\qquad
n\,\operatorname{Var}\!\left[\MK_c^\epsilon(\hat\alpha_n,\hat\beta_n)\right]
\longrightarrow\mathsf{v}_\epsilon.
\]
Equivalently,
\[
\sqrt n\left(
\MK_c^\epsilon(\hat\alpha_n,\hat\beta_n)
-
\MK_c^\epsilon(\alpha,\beta)
\right)
\overset{\mathrm{law}}{\longrightarrow}\Gaussian(0,\mathsf{v}_\epsilon).
\]
For every fixed \(\epsilon>0\), one moreover has \(B_{n,n}^\epsilon=O(n^{-1})\).
\end{prop}
\begin{proof}
Writing \(\hat\a_n\) and \(\hat\b_n\) for the two empirical histograms, the independent multinomial central-limit theorems give \(\sqrt n(\hat\a_n-\a,\hat\b_n-\b)\overset{\mathrm{law}}{\longrightarrow}(G_\a,G_\b)\), where \(G_\a\) and \(G_\b\) are independent centered Gaussian vectors with covariance matrices \(\diag(\a)-\a\a^\top\) and \(\diag(\b)-\b\b^\top\). Proposition~\ref{prop-ot-first-variations-unregularized} gives the directional derivative for every \(\epsilon\geq0\); it is linear under the assumed uniqueness at \(\epsilon=0\), and automatically at positive temperature. Thus,
\(D\MK_c^\epsilon(\alpha,\beta)[h,k] = \sum_i f_\epsilon^\star(x_i)h_i + \sum_j g_\epsilon^\star(y_j)k_j.\)
The finite-dimensional delta method therefore gives the centered Gaussian limit stated in the proposition. Independence and the two multinomial covariance formulas give its variance \(\mathsf{v}_\epsilon\).

Set \(R\eqdef\max_{i,j}c(x_i,y_j)-\min_{i,j}c(x_i,y_j)\). Exact dual potentials may be replaced by their \(c\)-transforms, and entropic potentials by their soft \(c\)-transforms; in either case their oscillations are at most \(R\). Since differences of histograms have zero sum, integrating the corresponding subgradient bounds along segments in the two marginal simplices gives \(\abs{\MK_c^\epsilon(\alpha',\beta')-\MK_c^\epsilon(\alpha,\beta)}\leq\frac{R}{2}(\norm{\a'-\a}_1+\norm{\b'-\b}_1)\) for any histograms \((\a',\b')\) and their associated measures \((\alpha',\beta')\). Because \(N,M\) are fixed, the fourth moments of \(\sqrt n\norm{\hat\a_n-\a}_1\) and \(\sqrt n\norm{\hat\b_n-\b}_1\) are uniformly bounded. The squared rescaled OT values are therefore uniformly integrable. The Gaussian limit then yields convergence of the means and second moments, which proves the assertions for \(B_{n,n}^\epsilon\), \(Z_{n,n}^\epsilon\), and the variance.

For fixed \(\epsilon>0\), the entropic value is twice continuously differentiable in a neighborhood of the positive pair \((\a,\b)\). On that neighborhood, a second-order Taylor expansion has a uniformly bounded remainder; its linear term has zero expectation and \(\EE(\norm{\hat\a_n-\a}^2+\norm{\hat\b_n-\b}^2)=O(n^{-1})\). The probability of leaving the neighborhood is exponentially small, and the preceding global Lipschitz bound controls its contribution. Hence \(B_{n,n}^\epsilon=O(n^{-1})\).
\end{proof}
\paragraph{Fixed support versus a continuum.}
More explicitly, its limiting variable is
\[
\sum_{i=1}^{N}f_\epsilon^\star(x_i)(G_\a)_i
+
\sum_{j=1}^{M}g_\epsilon^\star(y_j)(G_\b)_j,
\qquad
\begin{cases}
\operatorname{Cov}(G_\a)=\diag(\a)-\a\a^\top,\\
\operatorname{Cov}(G_\b)=\diag(\b)-\b\b^\top.
\end{cases}
\]
If \(a_{\min}=\min_i a_i\) and \(b_{\min}=\min_j b_j\), then
\(\PP[\text{some source or target atom is unobserved}] \leq N e^{-n a_{\min}}+M e^{-n b_{\min}}.\)
For uniform weights, the source contribution is at most $\delta/2$ as soon as $n\geq N[\log N+\log(2/\delta)]$, and the analogous target condition uses $M$. Thus a vanishing failure probability requires $n/N-\log N\to+\infty$ and $n/M-\log M\to+\infty$; nearly uniform weights have the same scale up to constants.

\[
\alpha_t=\left(\frac{1+t}{2},\frac{1-t}{2}\right),
\qquad
\beta_t=\left(\frac{1-t}{2},\frac{1+t}{2}\right).
\]
Then $\Wass_p(\alpha_t,\beta_t)^p=t$ and $\Wass_p(\alpha_t,\beta_t)=t^{1/p}$. Thus the right directional derivative of $\Wass_p^p$ at $t=0$ is nonzero, while that of $\Wass_p$ is nonzero for $p=1$ and infinite for $p>1$.

\subsection{Sketching Sinkhorn in Linear Time}
\label{sec-sketching-sinkhorn}
\paragraph{PSD kernels and ordinary feature sketches.}
The RKHS/MMD section, Section~\ref{sec-rkhs-mmd}, already used positive semidefinite kernels to define Hilbertian discrepancies between measures. Recall that a symmetric kernel \(k:\Xx\times\Xx\to\RR\) is positive semidefinite if, for every finite family \((x_i)_{i=1}^n\), the Gram matrix \((k(x_i,x_j))_{i,j}\) is positive semidefinite. Many scalable kernel methods start from an integral feature representation
\(k(x,y)=\int_Z \dotp{\phi(x,z)}{\phi(y,z)}_{\RR^p}\,d\xi(z), \qquad \phi:\Xx\times Z\to\RR^p,\)
Drawing \(z_1,\ldots,z_r\sim\xi\) gives the Monte-Carlo sketch
\(\widetilde k_r(x,y) = \frac1r\sum_{\ell=1}^r \dotp{\phi(x,z_\ell)}{\phi(y,z_\ell)}_{\RR^p}.\)
Set \(R=rp\), flatten \((\ell,q)\) into a single feature index, and define
\eql{\label{eq-rectangular-feature-sketch}
(\Phi_X)_{i,(\ell,q)}=r^{-1/2}\phi_q(x_i,z_\ell),
\qquad
(\Phi_Y)_{j,(\ell,q)}=r^{-1/2}\phi_q(y_j,z_\ell),
\qquad
\widetilde K=\Phi_X\Phi_Y^\top.
}
Here \(\Phi_X\in\RR^{n\times R}\), \(\Phi_Y\in\RR^{m\times R}\), and
\(\widetilde K_{i,j} = \frac1r\sum_{\ell=1}^r \dotp{\phi(x_i,z_\ell)}{\phi(y_j,z_\ell)}_{\RR^p} = \widetilde k_r(x_i,y_j).\)
\begin{prop}[Entrywise accuracy of rectangular feature sketches]
\label{prop-sinkhorn-sketch-positive-guarantee}
Let \(z_1,\ldots,z_r\) be independent with law \(\xi\), and let \(K_{i,j}=k(x_i,y_j)\) and \(\widetilde K\) be the rectangular kernel matrix and sketch defined above, with \(\widetilde K\) given by~\eqref{eq-rectangular-feature-sketch}. Assume that, for some \(M<+\infty\),
$\abs{\dotp{\phi(x,z)}{\phi(y,z)}_{\RR^p}}\leq M$ for all $x,y$ and $\xi$-a.e.\ $z$.
If \(K_{\min}\eqdef\min_{i,j}K_{i,j}>0\), then, for every \(0<\delta\leq1/2\),
\eql{\label{eq-rectangular-sketch-relative-concentration}
\PP\pa{
\max_{i,j}\abs{\frac{\widetilde K_{i,j}}{K_{i,j}}-1}>\delta
}
\leq
2nm\exp\pa{-\frac{rK_{\min}^2\delta^2}{2M^2}}.
}
On the complementary event, \(\widetilde K\) is entrywise positive and, for every \(\epsilon>0\), \(\max_{i,j} \abs{-\epsilon\log\widetilde K_{i,j}+\epsilon\log K_{i,j}} \leq 2\epsilon\delta.\)
\end{prop}
\begin{proof}
For each fixed \((i,j)\), the summands in the definition of \(\widetilde K_{i,j}\) have expectation \(K_{i,j}\) and lie in \([-M,M]\). Hoeffding's inequality therefore gives
\[
\PP\pa{
\abs{\widetilde K_{i,j}-K_{i,j}}>\delta K_{i,j}
}
\leq
2\exp\pa{-\frac{r\delta^2K_{i,j}^2}{2M^2}}
\leq
2\exp\pa{-\frac{r\delta^2K_{\min}^2}{2M^2}}.
\]
A union bound over the \(nm\) entries proves~\eqref{eq-rectangular-sketch-relative-concentration}. On its complementary event, write \(\widetilde K_{i,j}=K_{i,j}(1+s_{i,j})\), where \(\abs{s_{i,j}}\leq\delta\leq1/2\). Hence \(\widetilde K_{i,j}>0\), and \(\abs{\log(1+s_{i,j})}\leq2\abs{s_{i,j}}\) gives the logarithmic estimate.
\end{proof}
\paragraph{Application to Sinkhorn kernels.}
Let $\alpha=\sum_{i=1}^n a_i\delta_{x_i}$ and $\beta=\sum_{j=1}^m b_j\delta_{y_j}$, and write $\a=(a_i)_i$ and $\b=(b_j)_j$ for their weight vectors. Specialize the preceding construction to the Gibbs kernel \(k_\epsilon(x,y)=e^{-c(x,y)/\epsilon}\), assumed to be PSD and to admit the displayed feature representation. Replacing \(K\) by its rank-\(R\) sketch~\eqref{eq-rectangular-feature-sketch}, the two matrix-vector products are evaluated as
\(\widetilde K v=\Phi_X(\Phi_Y^\top v), \qquad \widetilde K^\top u=\Phi_Y(\Phi_X^\top u),\)
The Nystr\"om--Sinkhorn method first sets
\[
\zeta\eqdef\min\left\{1,
\frac{\tau/\epsilon}{50\bigl(4D^2/\epsilon+\log(N\epsilon/\tau)\bigr)}\right\},
\qquad
\rho\eqdef\frac{\zeta}{2}e^{-4D^2/\epsilon}.
\]
If \(I\subset\{1,\ldots,N\}\) is the resulting landmark set, it forms \(V=K_{:,I}\), \(A=K_{I,I}\), and stores only the rank-at-most-\(R\) factorization
\(\widetilde K=VA^\dagger V^\top, \qquad R=\abs{I},\)
\begin{prop}[Accuracy and complexity of Gaussian Nystr\"om--Sinkhorn]
\label{prop-gaussian-nystrom-sinkhorn-complexity}
Let \(z_1,\ldots,z_N\in B(0,D)\subset\RR^d\), where \(N\geq2\) and \(D>1\), and let \(\alpha=\sum_{i=1}^N a_i\delta_{z_i}\) and \(\beta=\sum_{i=1}^N b_i\delta_{z_i}\) have positive weights. Write \(\a=(a_i)_i\), \(\b=(b_i)_i\), and \(\C_{i,j}=\norm{z_i-z_j}^2\). Fix \(\epsilon\in[1/N,1]\) and \(\tau,\delta\in(0,1)\), and apply the preceding Nystr\"om--Sinkhorn procedure. Its kernel approximation \(\widetilde K\) is entrywise positive, and it returns a feasible coupling \(\widehat{\P}\in\CouplingsD(\a,\b)\) and a value \(\widehat W\) satisfying
\[
0\leq
\dotp{\C}{\widehat{\P}}
+\epsilon\KL(\widehat{\P}\,|\,\a\otimes\b)
-\MK_c^\epsilon(\alpha,\beta)
\leq\tau,
\qquad
\abs{\widehat W-\MK_c^\epsilon(\alpha,\beta)}\leq\tau.
\]
With probability at least \(1-\delta\), the selected sketch rank satisfies
\eql{\label{eq-gaussian-nystrom-rank}
R
\lesssim_d
\pa{
1+\frac{D^2}{\epsilon}
+\log\frac{N(1+D^2)}{\tau}
}^d
\log\frac{N}{\delta},
}
and the complete computation uses
\(\widetilde O\pa{ NR\pa{R+\frac{D^4}{\epsilon\tau}} } \quad\text{operations and}\quad O\pa{N(R+d)} \quad\text{memory}.\)
Here the constant hidden in \(\lesssim_d\) depends only on \(d\).
\end{prop}
\begin{proof}
Set the inverse temperature \(\eta=1/\epsilon\). The assumptions above are precisely the normalized range \(\eta\in[1,N]\), \(D>1\), and \(\tau\leq1\) used in the cited Nystr\"om--Sinkhorn theorem. Since \(\norm{z_i-z_j}\leq2D\), the smallest kernel entry \(K_{\min}:=\min_{i,j}K_{i,j}\) satisfies \(K_{\min}\geq e^{-4D^2/\epsilon}\). The adaptive Nystr\"om guarantee gives \(\max_{i,j}\abs{K_{i,j}-\widetilde K_{i,j}}\leq\rho\). Hence \(\rho\leq\zeta K_{\min}/2\), so \(\widetilde K\) is entrywise positive and \(\max_{i,j}\abs{\log K_{i,j}-\log\widetilde K_{i,j}}\leq\zeta\). The Sinkhorn stability and rounding estimates then give the two deterministic error bounds. The Gaussian effective-dimension estimate yields~\eqref{eq-gaussian-nystrom-rank} with probability at least \(1-\delta\). Constructing the factorization costs \(O(NR^2)\), each kernel product costs \(O(NR)\), and the quantitative scaling bound contributes \(\widetilde O(NRD^4/(\epsilon\tau))\), which yields the stated time and memory estimates. The entropy convention in the cited theorem differs from the KL-normalized value \(\MK_c^\epsilon\) only by a marginal-dependent constant, so the same additive guarantees hold after applying that known correction to \(\widehat W\).
\end{proof}
\paragraph{Positive features and complete positivity.}
\begin{defn}[Doubly nonnegative and completely positive matrices and kernels]\label{def-dnn-cp-kernels}
In finite dimension, set
\[
\mathrm{DNN}_n=\enscond{A\in\RR^{n\times n}}{A\succeq0,\ A_{i,j}\geq0},
\qquad
\mathrm{CP}_n=\enscond{BB^\top}{B\in\RR_+^{n\times q}\hbox{ for some }q\geq1}.
\]
Matrices in these sets are called doubly nonnegative and completely positive, respectively. A kernel is doubly positive, respectively completely positive, if every finite Gram matrix belongs to \(\mathrm{DNN}_n\), respectively \(\mathrm{CP}_n\).
It is a positive-feature kernel if there are a probability space \((Z,\xi)\), an integer \(p\geq1\), and measurable features \(\phi(x,\cdot):Z\to\RR_+^p\) such that
\(k(x,y)=\int_Z\dotp{\phi(x,z)}{\phi(y,z)}_{\RR^p}\d\xi(z) \qquad (x,y\in\Xx).\)
\end{defn}
\begin{prop}[Complete positivity and positive features]
\label{prop-cp-kernel-positive-features}
Let \(\Xx\) be a separable metric space and let \(k:\Xx\times\Xx\to\RR\) be a finite-valued continuous PSD kernel. Then \(k\) is completely positive if and only if it is a positive-feature kernel. Moreover, scalar features \(p=1\) suffice. For a finite set \(\Xx=\{x_1,\ldots,x_n\}\), this reduces to \(K\in\mathrm{CP}_n\) if and only if \(K=BB^\top\) for some \(B\geq0\).
\end{prop}
\begin{proof}
Suppose first that \(k\) has a positive-feature representation. For \(x_1,\ldots,x_n\), define the nonnegative vector \(v_q(z)\eqdef(\phi_q(x_i,z))_{i=1}^n\geq0\). Its Gram matrix is \(K=\sum_{q=1}^p\int_Z v_q(z)v_q(z)^\top\d\xi(z)\).
The integrand belongs to \(\mathrm{CP}_n\), and the closed convex cone \(\mathrm{CP}_n\) is stable under such integrals; hence \(K\in\mathrm{CP}_n\).

Conversely, choose a dense sequence \((x_i)_{i\geq1}\) in \(\Xx\) and positive weights \((\omega_i)_{i\geq1}\) such that \(S:=\sum_i\omega_i k(x_i,x_i)<+\infty\). The case \(S=0\) is trivial. For each sufficiently large \(r\), complete positivity gives a factorization \(K^{(r)}=B^{(r)}(B^{(r)})^\top\) of the Gram matrix on \(x_1,\ldots,x_r\). Write \(b_\ell\) for its nonzero columns and set \(t_\ell:=\sum_{i=1}^r\omega_i(b_\ell)_i^2\) and \(S_r:=\sum_{i=1}^r\omega_i k(x_i,x_i)=\sum_\ell t_\ell\). On the compact metrizable product \(Z:=\prod_{i\geq1}[0,\sqrt{S/\omega_i}]\), define \(z^{(r,\ell)}_i=\sqrt{S_r/t_\ell}(b_\ell)_i\) for \(i\leq r\) and \(z^{(r,\ell)}_i=0\) otherwise. Then \(\xi_r:=\sum_\ell(t_\ell/S_r)\delta_{z^{(r,\ell)}}\) is a probability measure on \(Z\), and \(\int_Zz_i z_j\d\xi_r(z)=k(x_i,x_j)\) whenever \(i,j\leq r\). Compactness yields a weakly convergent subsequence with limit \(\xi\). Since every coordinate product \(z_i z_j\) is continuous and bounded on \(Z\), the scalar features \(\phi(x_i,z):=z_i\) satisfy \(\int_Z\phi(x_i,z)\phi(x_j,z)\d\xi(z)=k(x_i,x_j)\) for \(i,j\geq1\).
For indices \(i,j\), \(\norm{\phi(x_i,\cdot)-\phi(x_j,\cdot)}_{L^2(\xi)}^2 =k(x_i,x_i)+k(x_j,x_j)-2k(x_i,x_j).\)
Continuity of \(k\) therefore makes \(\phi(x_i,\cdot)\) Cauchy whenever \(x_i\) converges in \(\Xx\). Density extends \(\phi\) continuously to all of \(\Xx\), independently of the approximating sequence; its values remain nonnegative because \(L^2_+(\xi)\) is closed. Taking limits in the inner-product identity on the dense sequence gives \(\langle\phi(x,\cdot),\phi(y,\cdot)\rangle_{L^2(\xi)}=k(x,y)\), proving the positive-feature representation.
\end{proof}
For a rectangular source--target matrix, Proposition~\ref{prop-cp-kernel-positive-features} is applied to the Gram matrix on the union of the two sampled point sets, after which one extracts the rectangular cross block.

Despite this certification difficulty, completely positive kernels, equivalently positive-sketchable kernels under the assumptions of Proposition~\ref{prop-cp-kernel-positive-features}, enjoy a useful closure algebra. It provides a systematic way to construct new kernels whose low-rank features remain nonnegative.

\begin{prop}[Basic closure of completely positive kernels]
\label{prop-completely-positive-kernel-closure}
Let \(k_1,\ldots,k_J\) be completely positive kernels, \(a_1,\ldots,a_J\geq0\), and \((k_\theta)_{\theta\in\Theta}\) a measurable family of completely positive kernels. Whenever they are finite, the sum \(k_{\mathrm{sum}}:=\sum_{j=1}^J a_jk_j\), the product \(k_{\mathrm{prod}}:=\prod_{j=1}^J k_j\), and the mixture \(k_{\mathrm{mix}}:=\int_\Theta k_\theta\,\d\eta(\theta)\), with \(\eta\geq0\), are completely positive. In particular, so is every positive autocorrelation
$k(x,y)=\int g(u-x)g(u-y)\,du$ with $g\geq0$,
and every nonnegative mixture of such kernels.
\end{prop}
\begin{proof}
Restrict the kernels to arbitrary points \(x_1,\ldots,x_n\). If \(K_1=B_1B_1^\top\) and \(K_2=B_2B_2^\top\) with \(B_1,B_2\geq0\), then a nonnegative weighted sum is factored by concatenating the correspondingly rescaled columns. For \(w_{\ell,s}:=(B_1)_{\cdot,\ell}\odot(B_2)_{\cdot,s}\geq0\), one has \(K_1\odot K_2=\sum_{\ell,s}w_{\ell,s}w_{\ell,s}^\top\), so the pointwise product is completely positive. A measurable nonnegative mixture is an integral of matrices in the closed convex cone \(\mathrm{CP}_n\), and therefore remains in that cone. Finally, the autocorrelation formula is itself a scalar positive-feature representation, so its finite Gram matrices are completely positive by the same integral argument.
\end{proof}
\begin{prop}[Generalized Gaussian positive features]
\label{prop-gaussian-positive-features}
Let \(0<p\leq2\), \(\epsilon>0\), and
$k_{p,\epsilon}(x,y) = \exp\pa{-\frac{\norm{x-y}^p}{\epsilon}}$ for $x,y\in\RR^d$.
Set \(a=p/2\), let \(S_a\) be the positive \(a\)-stable random variable normalized by
$\EE\pa{e^{-tS_a}}=e^{-t^a}$ for $t\geq0$,
and draw independently \(\Lambda=\epsilon^{-1/a}S_a\) and \(Z\sim\Gaussian(0,\Id_d)\).
For any \(x_0\in\RR^d\), define the positive feature
\(\phi_{p,\epsilon}(x;\Lambda,Z) := \exp\pa{ \sqrt{2\Lambda}\dotp{Z}{x-x_0} -2\Lambda\norm{x-x_0}^2 }.\)
Then \(k_{p,\epsilon}(x,y)=\EE[\phi_{p,\epsilon}(x;\Lambda,Z)\phi_{p,\epsilon}(y;\Lambda,Z)]\), so \(k_{p,\epsilon}\) is completely positive. For i.i.d.\ copies \((\Lambda_\ell,Z_\ell)_{\ell=1}^r\), define the positive sketching features \(\varphi_\ell(x)=r^{-1/2}\phi_{p,\epsilon}(x;\Lambda_\ell,Z_\ell)\), for \(1\leq\ell\leq r\).
They satisfy \(\varphi_\ell\geq0\) and \(\EE[\sum_{\ell=1}^r\varphi_\ell(x)\varphi_\ell(y)]=k_{p,\epsilon}(x,y)\). For \(p=2\), one has \(S_1=1\) almost surely and \(\Lambda=1/\epsilon\), hence
\(\phi_{2,\epsilon}(x;Z) = \exp\pa{ \sqrt{\frac{2}{\epsilon}}\dotp{Z}{x-x_0} -\frac{2}{\epsilon}\norm{x-x_0}^2 }.\)
\end{prop}
\begin{proof}
The stable-law normalization gives
$\EE e^{-t\Lambda} = \exp\pa{-\frac{t^a}{\epsilon}}$,
while the Gaussian moment-generating function gives \(\EE_Z[\phi_{p,\epsilon}(x;\Lambda,Z)\phi_{p,\epsilon}(y;\Lambda,Z)]=e^{-\Lambda\norm{x-y}^2}\) conditionally on \(\Lambda\).
Taking \(t=\norm{x-y}^2\) in the first identity proves the feature formula. The Monte Carlo statement follows by averaging independent copies.
\end{proof}
Indeed, writing $u=x-x_0$, $v=y-x_0$, and $F_{x,y}=\phi_{p,\epsilon}(x;\Lambda,Z)\phi_{p,\epsilon}(y;\Lambda,Z)$, a second Gaussian-moment calculation gives \(\EE_Z[F_{x,y}^2\mid\Lambda] = \exp(8\Lambda\dotp{u}{v}).\)
At $x=y\neq x_0$, this is $\exp(8\Lambda\norm{u}^2)$. For $0<p<2$, the positive stable mixing variable has no positive exponential moment, so the second moment is infinite; for $p=2$, it grows as $\exp(8\norm{u}^2/\epsilon)$.

\paragraph{Positive sketches for Sinkhorn.}
For the \(p\)-power cost \(c_p(x,y)=\norm{x-y}^p\), with \(0<p\leq2\), let
\(K_{i,j}=k_{p,\epsilon}(x_i,y_j) =\exp\pa{-\frac{\norm{x_i-y_j}^p}{\epsilon}}.\)
For \(p\geq1\), this is the usual \(p\)-Wasserstein power cost; for \(0<p<1\), it is simply a concave power transport cost. The factor \(1/p\), often used in \(\Wass_p^p\), only rescales \(\epsilon\). Sinkhorn applied to \(K_r\) uses only products by the two feature matrices, and its plan is represented as
\(P_r = \diag(u)\Phi_X\Phi_Y^\top\diag(v) = LR^\top, \qquad L=\diag(u)\Phi_X, \quad R=\diag(v)\Phi_Y.\)
\paragraph{Connection with linear time attention.}
Feature-based linear attention replaces this positive kernel by \(\phi(q_i)^\top\phi(k_j)\) and, whenever the denominator is positive, computes the normalized output as
\(y_i = \frac{\phi(q_i)^\top\sum_j\phi(k_j)v_j} {\phi(q_i)^\top\sum_j\phi(k_j)},\)

\section{Generalized Wasserstein Distances}
\label{sec-extensions}
\label{sec-generalized-wasserstein-distances}
\subsection{Unbalanced OT}
\label{sec-unbalanced}
\paragraph{Relaxed formulation.}
\begin{defn}[Unbalanced optimal transport]\label{def-unbalanced-optimal-transport}
For nonnegative measures $(\al,\be) \in \Mm_+(\Xx) \times \Mm_+(\Yy)$, a nonnegative measurable cost $c$, and proper lower-semicontinuous convex entropy functions $\psi_s:[0,+\infty)\to[0,+\infty]$ satisfying $\psi_s(1)=0$, define
\eql{\label{eq-unbalanced-primal}
\UW_c(\al,\be)\eqdef \uinf{\pi \in \Mm_+(\Xx \times \Yy)} \int_{\Xx \times \Yy} c(x,y) \d\pi(x,y)
+ \Dd_{\psi_1}(\pi_1|\al) + \Dd_{\psi_2}(\pi_2|\be),
}
where $\Dd_{\psi_s}$ is the $\phi$-divergence of Definition~\ref{def_divergence}.
\end{defn}
Writing $\psi_s=\tau\bar\psi_s$ exposes the relaxation scale: \(\UW_{c,\tau}(\alpha,\beta) = \inf_{\pi\geq0} \int c\d\pi + \tau\Dd_{\bar\psi_1}(\pi_1|\alpha) + \tau\Dd_{\bar\psi_2}(\pi_2|\beta).\)
To state this precisely, extend the Kantorovich value~\eqref{eq-mk-generic} by positive homogeneity: for finite measures $(\widetilde\alpha,\widetilde\beta)$ with common mass $M$,
\begin{equation}\label{eq-finite-mass-kantorovich-extension}
\MK_c(\widetilde\alpha,\widetilde\beta)
\eqdef
\begin{cases}
M\MK_c(\widetilde\alpha/M,\widetilde\beta/M),&M>0,\\
0,&M=0.
\end{cases}
\end{equation}
\begin{prop}[Balanced transport between relaxed marginals]\label{prop-uot-balanced-relaxed-marginals}
Under the hypotheses of Definition~\ref{def-unbalanced-optimal-transport},
\begin{equation}\label{eq-uot-relaxed-marginal-reformulation}
\UW_c(\alpha,\beta)
=
\inf_{\substack{\widetilde\alpha\in\Mm_+(\Xx),\ \,\widetilde\beta\in\Mm_+(\Yy)\\
\widetilde\alpha(\Xx)=\widetilde\beta(\Yy)}}
\left\{
\MK_c(\widetilde\alpha,\widetilde\beta)
+\Dd_{\psi_1}(\widetilde\alpha|\alpha)
+\Dd_{\psi_2}(\widetilde\beta|\beta)
\right\}.
\end{equation}
A plan $\pi^*$ solves the unbalanced problem if and only if its marginals $(\pi_1^*,\pi_2^*)$ solve~\eqref{eq-uot-relaxed-marginal-reformulation} and $\pi^*$ is an optimal balanced coupling between them.
\end{prop}
\begin{proof}
Every plan $\pi$ has equal-mass marginals $(\widetilde\alpha,\widetilde\beta)=(\pi_1,\pi_2)$. For fixed marginals, the divergence terms are constant and minimizing the remaining transport term gives~\eqref{eq-finite-mass-kantorovich-extension}. Grouping the plans in~\eqref{eq-unbalanced-primal} by their marginals proves~\eqref{eq-uot-relaxed-marginal-reformulation}. The same argument shows that an optimizer must be transport-optimal conditionally on its marginals and that these marginals minimize the outer problem. Conversely, an optimal balanced coupling between any minimizing pair in~\eqref{eq-uot-relaxed-marginal-reformulation} attains the unbalanced value.
\end{proof}
\paragraph{KL mass--shape separation.}
Let $A=\alpha(\Xx)>0$, $B=\beta(\Yy)>0$, and write the normalized input shapes as $\widehat\alpha=\alpha/A$ and $\widehat\beta=\beta/B$.

\begin{prop}[Exact mass--shape separation for KL relaxation]\label{prop-uot-kl-mass-shape}
Assume $\psi_1=\tau_0\phi_{\KL}$ and $\psi_2=\tau_1\phi_{\KL}$, where $\phi_{\KL}(s)=s\log s-s+1$ for $s\geq0$, with $0\log0=0$ and $\tau_0,\tau_1>0$. Set $T=\tau_0+\tau_1$, $\theta_i=\tau_i/T$, and
\begin{equation}\label{eq-uot-kl-shape-envelope}
\mathcal E_{c,\tau_0,\tau_1}(\widehat\alpha,\widehat\beta)
\eqdef
\inf_{\substack{\widetilde\alpha\in\Pp(\Xx)\\
\widetilde\beta\in\Pp(\Yy)}}
\left\{
\MK_c(\widetilde\alpha,\widetilde\beta)
+\tau_0\KL(\widetilde\alpha|\widehat\alpha)
+\tau_1\KL(\widetilde\beta|\widehat\beta)
\right\}.
\end{equation}
Then, with $e^{-\infty}=0$,
\begin{equation}\label{eq-uot-kl-mass-shape}
\UW_c(\alpha,\beta)
=
\tau_0A+\tau_1B
-T A^{\theta_0}B^{\theta_1}
\exp\!\left(
-\frac{\mathcal E_{c,\tau_0,\tau_1}(\widehat\alpha,\widehat\beta)}{T}
\right).
\end{equation}
If $(\widetilde\alpha^*,\widetilde\beta^*)$ attains the finite shape envelope and $\bar\pi^*$ is an optimal probability coupling between these marginals, then $\pi^*=M^*\bar\pi^*$ is optimal, where
\begin{equation}\label{eq-uot-kl-optimal-mass}
M^*
=
A^{\theta_0}B^{\theta_1}
\exp\!\left(
-\frac{\mathcal E_{c,\tau_0,\tau_1}(\widehat\alpha,\widehat\beta)}{T}
\right)
\end{equation}
\end{prop}
\begin{proof}
Write every nonzero plan as $\pi=M\bar\pi$, where $M=\pi(\Xx\times\Yy)>0$ and $\bar\pi$ is a probability coupling with marginals $(\widetilde\alpha,\widetilde\beta)$. The KL homogeneity identity
\begin{equation}\label{eq-kl-mass-shape-identity}
\KL(M\widetilde\alpha|A\widehat\alpha)
=
M\KL(\widetilde\alpha|\widehat\alpha)
+M\log(M/A)-M+A
\end{equation}
holds in the extended-value sense, and similarly for the second marginal. Minimizing first over $\bar\pi$ and its normalized marginals therefore reduces the objective to
\(M\mathcal E_{c,\tau_0,\tau_1}(\widehat\alpha,\widehat\beta) +\tau_0\big(M\log(M/A)-M+A\big) +\tau_1\big(M\log(M/B)-M+B\big).\)
When the shape value is finite, this strictly convex function of $M$ has the unique minimizer~\eqref{eq-uot-kl-optimal-mass}; substitution gives~\eqref{eq-uot-kl-mass-shape}. If the shape value is infinite, every nonzero plan has infinite objective and the same formula returns the value $\tau_0A+\tau_1B$ of the zero plan.
\end{proof}
These figures should be read as pictures of a single relaxed plan $\pi$: the same nonnegative measure determines both the transported coupling and its two relaxed marginals. The small-$\tau$ result below formalizes the opposite regime, where transport becomes negligible compared with local mass variation.

\begin{prop}[Small-transport-scale limit for marginal penalties]\label{prop-unbalanced-small-scale-limit}
Assume that $\alpha,\beta$ are finite measures on a compact metric space $\Xx$, that $c$ is continuous, $c\geq0$, and $c(x,y)=0$ if and only if $x=y$. Assume also that the marginal divergences are nonnegative, weak-* lower semicontinuous, and have weak-* compact sublevel sets on $\Mm_+(\Xx)$. Then
\[
\lim_{\tau\downarrow0}\frac1\tau\UW_{c,\tau}(\alpha,\beta)
=
\inf_{\chi\in\Mm_+(\Xx)}
\Dd_{\bar\psi_1}(\chi|\alpha)
+
\Dd_{\bar\psi_2}(\chi|\beta).
\]
The right-hand side is the infimal gluing divergence obtained by matching the two measures through a common zero-transport marginal $\chi$. In the dominated case, if $\alpha=a\lambda$, $\beta=b\lambda$, and the minimizing common marginal is absolutely continuous, $\chi=r\lambda$, this divergence decouples pointwise as
\[
\int \mathfrak m_{\bar\psi_1,\bar\psi_2}(a(x),b(x))\d\lambda(x),
\qquad
\mathfrak m_{\bar\psi_1,\bar\psi_2}(a,b)
\eqdef
\inf_{r\geq0} a\,\bar\psi_1(r/a)+b\,\bar\psi_2(r/b),
\]
with the usual recession conventions when $a=0$ or $b=0$. For superlinear entropies, and in particular for KL, finite energy forces this dominated form. Thus, when $\bar\psi_1=\bar\psi_2$ is the KL entropy,
\(\inf_{\chi\in\Mm_+(\Xx)} \KL(\chi|\alpha)+\KL(\chi|\beta) = \int (\sqrt{a}-\sqrt{b})^2\d\lambda.\)
Thus the KL marginal relaxation contains the squared Hellinger distance as its local mass-variation limit.
\end{prop}
\begin{proof}
For the upper bound, restrict to diagonal plans $\pi=(\Id,\Id)_\sharp\chi$, whose transport cost is zero and whose two marginals are both $\chi$.

For the lower bound, let $\tau_n\downarrow0$ and let $\pi_n$ be almost minimizing plans with bounded scaled values $\tau_n^{-1}\UW_{c,\tau_n}(\alpha,\beta)$. Since the divergences are nonnegative, $\int c\d\pi_n=O(\tau_n)$, hence $\int c\d\pi_n\to0$. The bounded scaled values also put the two marginals in compact divergence sublevel sets. Since a coupling has the same total mass as each marginal, the couplings are tight on $\Xx\times\Xx$. Up to subsequences, $\pi_n\rightharpoonup\pi_0$. Lower semicontinuity of the transport cost yields $\int c\d\pi_0=0$, so $\pi_0$ is concentrated on the diagonal. Its two marginals are therefore equal to a common measure $\chi$. Lower semicontinuity of the marginal divergences gives
\(\liminf_n\frac1{\tau_n}\UW_{c,\tau_n}(\alpha,\beta) \geq \Dd_{\bar\psi_1}(\chi|\alpha)+\Dd_{\bar\psi_2}(\chi|\beta),\)
and optimizing over $\chi$ gives the lower bound.

In the dominated case, writing $\chi=r\lambda$ gives \(\Dd_{\bar\psi_1}(\chi|\alpha)+\Dd_{\bar\psi_2}(\chi|\beta) = \int a\,\bar\psi_1(r/a)+b\,\bar\psi_2(r/b) \d\lambda,\)
so the minimization over $\chi$ decouples into the scalar envelope $\mathfrak m_{\bar\psi_1,\bar\psi_2}$. For KL, no singular part is admissible when $\alpha$ and $\beta$ are dominated by $\lambda$. The pointwise objective is $r\log(r/a)-r+a+r\log(r/b)-r+b$. Its optimality condition is $\log(r/a)+\log(r/b)=0$, hence $r=\sqrt{ab}$, and the minimum is $a+b-2\sqrt{ab}=(\sqrt a-\sqrt b)^2$.
\end{proof}
\begin{prop}[Dual of unbalanced optimal transport]\label{prop-dual-unbalanced-ot}
Assume that $\Xx,\Yy$ are compact metric spaces, that $c$ is finite and continuous, and that $\alpha,\beta$ are finite Radon measures. Assume also that $\psi_1,\psi_2$ are entropy functions in the sense of Definition~\ref{def_entropy} and that either both are superlinear, or
$\psi'_{1,\infty}+\psi'_{2,\infty}+\min c>0$.
Assume finally that at least one plan has finite primal objective in~\eqref{eq-unbalanced-primal}.
Then one has equality between~\eqref{eq-unbalanced-primal} and
\[
\UW_c(\al,\be)
=
\sup_{\substack{f\in C(\Xx),\ g\in C(\Yy)\\ f\oplus g\leq c}}
-
\Dd_{\psi_1}^*(-f|\al)
-
\Dd_{\psi_2}^*(-g|\be).
\]
\end{prop}
\begin{proof}
Use the variational formula~\eqref{eq-legendre} for the dual of a divergence and introduce the marginal variables through continuous potentials:
\[
\inf_{\pi\geq0}\sup_{\substack{f\in C(\Xx)\\g\in C(\Yy)}}
\int c\d\pi + \int -f\d\pi_1 + \int -g\d\pi_2
-
\Dd_{\psi_1}^*(-f|\al)
-
\Dd_{\psi_2}^*(-g|\be).
\]
Under the compactness and coercivity assumptions in the statement, the
entropy--transport duality theorem of Liero, Mielke and Savar\'e
justifies exchanging the infimum and the supremum.
\[
\sup_{\substack{f\in C(\Xx)\\g\in C(\Yy)}}
-
\Dd_{\psi_1}^*(-f|\al)
-
\Dd_{\psi_2}^*(-g|\be)
+
\inf_{\pi\geq0}\int\big(c-(f\oplus g)\big)\d\pi.
\]
The last infimum is $0$ when $f\oplus g\leq c$ and $-\infty$ otherwise, which gives the displayed dual.
\end{proof}
\paragraph{Reverse and homogeneous formulations.}
Assuming first that the reference measures and transported marginals have mutually absolutely continuous parts, one can factor the objective as
\begin{align*}
&\int c(x,y)\d\pi(x,y)
+\Dd_{\psi_1}(\pi_1|\al)+\Dd_{\psi_2}(\pi_2|\be) \\
&\qquad =
\int
\left(
c(x,y)
+
\psi_1\!\left(\frac{\d\pi_1}{\d\al}(x)\right)\frac{\d\al}{\d\pi_1}(x)
+
\psi_2\!\left(\frac{\d\pi_2}{\d\be}(y)\right)\frac{\d\be}{\d\pi_2}(y)
\right)
\d\pi(x,y).
\end{align*}
\begin{defn}[Local reverse cost and homogeneous perspective]\label{def-unbalanced-local-costs}
For $i\in\{1,2\}$ and $r\geq0$, let $R_i(r)=r\psi_i(1/r)$ when $r>0$ and $R_i(0)=\psi'_{i,\infty}$. For $c\in[0,+\infty]$ and $r,s\geq0$, define the reverse cost
\eql{\label{eq-unbalanced-reverse-local-cost}
L_c(r,s) \eqdef c+R_1(r)+R_2(s).
}
For finite $c$, first set
\[
\widetilde H_c(r,s)
\eqdef
\inf_{\theta>0}\theta L_c(r/\theta,s/\theta)
=
\inf_{\theta>0}
\bigl\{r\psi_1(\theta/r)+s\psi_2(\theta/s)+\theta c\bigr\},
\]
where the second expression is initially read for $r,s>0$. Its lower-semicontinuous envelope on $[0,+\infty)^2$ is denoted by
\eql{\label{eq-unbalanced-homogeneous-local-cost}
H_c \eqdef \operatorname{lsc}\widetilde H_c.
}
For $c=+\infty$, set $H_{+\infty}(r,s)=\psi_1(0)r+\psi_2(0)s$, which is also the monotone limit of $H_c$ as $c\to+\infty$.
\end{defn}
For the unit KL entropy $\psi(t)=t\log t-t+1$, direct minimization in the positive quadrant followed by this closure gives, with $e^{-\infty}=0$, \(H_c(r,s)=r+s-2\sqrt{rs}\,e^{-c/2}, \qquad c\in[0,+\infty],\quad r,s\geq0.\)
In particular, $H_c(r,0)=r$, $H_c(0,s)=s$ and $H_c(0,0)=0$. If $\al=F\pi_1+\al^\perp$ and $\be=G\pi_2+\be^\perp$ are the Lebesgue decompositions of the reference marginals with respect to the transported marginals, then the reverse formulation reads
\(\UW_c(\al,\be) = \inf_{\pi\geq0} \int L_{c(x,y)}(F(x),G(y))\d\pi(x,y) + \psi_1(0)\al^\perp(\Xx) + \psi_2(0)\be^\perp(\Yy).\)
Here one uses a base measure $\lambda$ on $\Xx\times\Yy$ and two nonnegative weights $u,v$ attached to each transported pair:
\eql{\label{eq-homogeneous}
\HW_c(\al,\be)
=
\inf_{\substack{\lambda\in\Mm_+(\Xx\times\Yy),\ u,v\geq0\\(\mathrm p_1)_\sharp(u\lambda)=\al,\ (\mathrm p_2)_\sharp(v\lambda)=\be}}
\int H_{c(x,y)}\big(u(x,y),v(x,y)\big)\d\lambda(x,y).
}
\begin{prop}[Homogenization does not change the unbalanced cost]\label{prop-homogeneous-unbalanced}
Assume that $\Xx,\Yy$ are Polish spaces, that $c:\Xx\times\Yy\to[0,+\infty]$ is lower semicontinuous, and that the entropy functions satisfy Definition~\ref{def_entropy}. Assume that either both entropies are superlinear, or
$\psi'_{1,\infty}+\psi'_{2,\infty}+\inf c>0$ and $c$ has compact sublevels. Suppose moreover that the primal problem has a finite competitor. Then
$\HW_c(\al,\be)=\UW_c(\al,\be)$.
\end{prop}
\begin{proof}
This is the homogeneous-reformulation theorem of Liero, Mielke and Savar\'e. Its interior mechanism is worth recording. The reverse representation and $H_c\leq\widetilde H_c\leq L_c$ show that every unbalanced competitor induces a semi-coupling of no larger cost, including singular mass through the closed axis values; hence $\HW_c\leq\UW_c$.

Conversely, first take a semi-coupling $(\lambda,u,v)$ with $u,v>0$ and finite $c$. Measurable $\eta$-minimizers $\theta>0$ of the raw perspective exist because the integrand is normal. With $\widetilde\pi=\theta\lambda$, disintegration and conditional Jensen for the perspective $(a,m)\mapsto a\psi_s(m/a)$ give
\[
\int c\d\widetilde\pi
+\Dd_{\psi_1}(\widetilde\pi_1|\alpha)
+\Dd_{\psi_2}(\widetilde\pi_2|\beta)
\leq
\int\bigl[c\theta+u\psi_1(\theta/u)+v\psi_2(\theta/v)\bigr]\d\lambda.
\]
Letting $\eta\downarrow0$ gives the converse inequality for interior semi-couplings and the raw perspective $\widetilde H_c$. For arbitrary semi-couplings, the relaxation theorem cited above passes to the lower-semicontinuous envelope $H_c$ while preserving the homogeneous marginal constraints; this is precisely the step that handles zero weights, singular mass and $c=+\infty$. The coercivity assumptions provide tightness. Minimizing over semi-couplings yields $\UW_c\leq\HW_c$.
\end{proof}
\paragraph{Conic lifting.}
Assume now that $\Xx=\Yy$ and $\psi_1=\psi_2=\psi$. The homogeneous formulation naturally lifts mass to a radial coordinate.

\begin{defn}[Cone lift]\label{def-unbalanced-cone-lift}
The cone over $\Xx$ is $\mathfrak C[\Xx]\eqdef(\Xx\times\RR_+)/\sim$, where all $(x,0)$ form one apex. For $p\geq1$, its ground cost is \(\De((x,r),(y,s))\eqdef H_{c(x,y)}(r^p,s^p)^{1/p}.\)
\end{defn}
\begin{itemize}
\item $\Dd_\psi=\KL$, $p=2$, and $c(x,y)=-\log\cos^2(d(x,y)\wedge\pi/2)$ give the Hellinger--Kantorovich or Wasserstein--Fisher--Rao cone metric
\(\De((x,r),(y,s))^2=r^2+s^2-2rs\cos(d(x,y)\wedge\pi/2).\)
\item $\Dd_\psi=\KL$, $p=2$, and $c(x,y)=d(x,y)^2$ give the Gaussian Hellinger formula
\(\De((x,r),(y,s))^2=r^2+s^2-2rs e^{-d(x,y)^2/2}.\)
It is a cone metric whenever the Gaussian kernel $k(x,y)=e^{-d(x,y)^2/2}$ is positive definite. This holds, for example, when $(\Xx,d)$ is a subset of a Hilbert space: if $k(x,y)=\langle\Phi(x),\Phi(y)\rangle$ with $\|\Phi(x)\|=1$, then the displayed quantity is $\|r\Phi(x)-s\Phi(y)\|^2$. For a general metric space, positive definiteness of this kernel is an additional hypothesis.
\item $\Dd_\psi=\TV$, $p=1$, and $c(x,y)=d(x,y)$ give the partial-transport cone cost
\(\De((x,r),(y,s))=r+s-(r\wedge s)(2-d(x,y))_+.\)
\end{itemize}
For a finite measure $\eta$ on the cone, define its weighted base projection $\mathsf P_p\eta\in\Mm_+(\Xx)$ by \(\int_{\Xx}\varphi(x)\d(\mathsf P_p\eta)(x) = \int_{\mathfrak C[\Xx]}\varphi(x)r^p\d\eta(x,r)\)
for every bounded Borel $\varphi$. This is well defined at the apex because the weight $r^p$ vanishes there. The corresponding cone value is
\[
\CW(\al,\be)
=
\inf_{\gamma\in\Mm_+(\mathfrak C[\Xx]^2)}
\left\{
\int \De((x,r),(y,s))^p\d\gamma
\;; \;
\mathsf P_p\gamma_1=\al,\quad
\mathsf P_p\gamma_2=\be
\right\}.
\]
\begin{defn}[Wasserstein--Fisher--Rao distance]\label{def-wfr-distance}
For $p=2$ and
\(\De((x,r),(y,s))^2=r^2+s^2-2rs\cos(d(x,y)\wedge\pi/2),\)
the Wasserstein--Fisher--Rao, or Hellinger--Kantorovich, distance is
$\WFR(\al,\be)\eqdef\CW(\al,\be)^{1/2}$.
\end{defn}
\begin{thm}[Cone formulation of unbalanced OT]\label{thm-cone-unbalanced-ot}
Under the hypotheses of Proposition~\ref{prop-homogeneous-unbalanced}, assume in addition that $\Xx$ is a compact metric space, that $(r,s)\mapsto H_{c(x,y)}(r,s)$ is convex for every $(x,y)$, and that $(x,y,r,s)\mapsto H_{c(x,y)}(r,s)$ is lower semicontinuous. Then
\(\UW_c(\alpha,\beta)=\HW_c(\alpha,\beta)=\CW(\alpha,\beta).\)
If, in addition, $\De$ is a continuous distance on $\mathfrak C[\Xx]$, then $\CW^{1/p}$ is a distance between nonnegative finite measures.
\end{thm}
\begin{proof}
The equality $\UW=\HW$ is Proposition~\ref{prop-homogeneous-unbalanced}.

Let $(\lambda,u,v)$ be feasible in~\eqref{eq-homogeneous}. Define the cone coupling
\(\gamma = \big((x,y)\mapsto ((x,u(x,y)^{1/p}),(y,v(x,y)^{1/p}))\big)_\sharp\lambda.\)
Then $\mathsf P_p\gamma_1=\al$ and $\mathsf P_p\gamma_2=\be$ by the
weighted marginal constraints. Moreover
\(\int \De((x,r),(y,s))^p\d\gamma = \int H_{c(x,y)}(u(x,y),v(x,y))\d\lambda(x,y).\)
Taking the infimum gives $\CW\leq\HW$.

Conversely, let $\gamma$ be feasible for $\CW$. Choose arbitrary
representatives of cone points at the apex. This does not affect the
weighted constraints, and it does not change the cost because, by the
recession convention, $H_{c(x,y)}(0,s)$ and $H_{c(x,y)}(r,0)$ do not depend
on the arbitrary spatial representative attached to the zero radius. Let
$\lambda$ be the projection of $\gamma$ onto the spatial variables
$(x,y)$, and disintegrate
$\gamma=\gamma_{x,y}\d\lambda(x,y)$ with respect to this projection. Set
$u(x,y)=\int r^p\d\gamma_{x,y}(r,s)$ and $v(x,y)=\int s^p\d\gamma_{x,y}(r,s)$.
The cone marginal constraints imply
$(\mathrm{p}_1)_\sharp(u\lambda)=\al$ and
$(\mathrm{p}_2)_\sharp(v\lambda)=\be$. Since $H_c$ is convex in its two
arguments, as a perspective of the convex reverse cost $L_c$, Jensen's
inequality gives, for $\lambda$-a.e.\ $(x,y)$, \(H_{c(x,y)}(u(x,y),v(x,y)) \leq \int H_{c(x,y)}(r^p,s^p)\d\gamma_{x,y}(r,s).\)
After integration and using
$\De((x,r),(y,s))^p=H_{c(x,y)}(r^p,s^p)$, this yields
$\HW\leq\CW$. Hence $\UW=\HW=\CW$.

Assume now that $\De$ is a distance. The nontrivial point in the triangle inequality is that two cone plans sharing the same weighted base projection need not have the same ordinary cone marginal. The radial homogeneity removes this obstruction. Indeed, common radial rescaling
\(((x,r),(y,s))\mapsto((x,r/\vartheta),(y,s/\vartheta)), \qquad \gamma\mapsto\vartheta^p\gamma,\)
preserves both weighted projections and the $p$-action. After adding mass at the cone apex, the normalization lemma for homogeneous lifts therefore lets one rescale two almost optimal plans for $(\alpha,\beta)$ and $(\beta,\zeta)$ so that their ordinary middle cone marginal is the same lift
$\eta_\beta = M\delta_{\mathfrak o} + (x\mapsto(x,1))_\sharp\beta$.
Here $M<+\infty$ is chosen large enough for both plans; additional mass at $(\mathfrak o,\mathfrak o)$ changes neither the weighted projections nor the action. The ordinary gluing lemma now produces a joint cone plan $(Z_0,Z_1,Z_2)$. Since $\De$ is a metric, Minkowski's inequality gives
\[
\left(\int\De(Z_0,Z_2)^p\right)^{1/p}
\leq
\left(\int\De(Z_0,Z_1)^p\right)^{1/p}
+
\left(\int\De(Z_1,Z_2)^p\right)^{1/p}.
\]
Taking infima proves the triangle inequality. This normalization-and-gluing argument is the metric engine of the homogeneous cone construction.

For completeness, normalize each non-apex pair by taking $\vartheta=(r^p+s^p)^{1/p}$. The transformed radii then satisfy $r^p+s^p=1$, while the transformed plan has total mass $\alpha(\Xx)+\beta(\Xx)$; apex mass can again be fixed without changing the problem. Thus minimizing plans may be restricted to a weakly compact family with bounded radii. Compactness of $\Xx$ and lower semicontinuity of the cone cost give an optimal plan. If $\CW(\alpha,\beta)=0$, such a plan is concentrated on the diagonal of the cone. There $x=y$ whenever the common radius is positive and $r=s$, so its two weighted base projections coincide and $\alpha=\beta$. Nonnegativity and symmetry are immediate, completing the metric proof.
\end{proof}
\paragraph{Entropic KL relaxation.}
On compact metric spaces, assume the continuous-cost setting of Proposition~\ref{prop-dual-unbalanced-ot}, with $\psi_1,\psi_2$ superlinear, and let $\phi$ be a superlinear entropy finite on $(0,+\infty)$. A generic entropic regularization of unbalanced OT then reads
\begin{equation}\label{eq-unbalanced-entropic-primal}
\inf_{\pi\in\Mm_+(\Xx\times\Yy)}
\int c\d\pi
+
\Dd_{\psi_1}(\pi_1|\al)
+
\Dd_{\psi_2}(\pi_2|\be)
+
\epsilon\Dd_\phi(\pi|\al\otimes\be).
\end{equation}
Fenchel--Rockafellar duality on finite Radon measures gives
\[
\sup_{\substack{f\in C(\Xx)\\g\in C(\Yy)}}
-
\Dd_{\psi_1}^*(-f|\al)
-
\Dd_{\psi_2}^*(-g|\be)
-
\epsilon\Dd_\phi^*\left(\frac{f\oplus g-c}{\epsilon}\middle|\al\otimes\be\right).
\]
For $\Dd_\phi=\KL$, the primal-dual relation is $\d\pi=e^{(f\oplus g-c)/\epsilon}\d\al\d\be$. If, in addition, $\Dd_{\psi_1}=\Dd_{\psi_2}=\tau\KL$, the dual specializes to
\[
\sup_{\substack{f\in C(\Xx)\\g\in C(\Yy)}}
-\tau\int(e^{-f/\tau}-1)\d\al
-\tau\int(e^{-g/\tau}-1)\d\be
-\epsilon\iint
\left(e^{(f\oplus g-c)/\epsilon}-1\right)
\d\al\d\be,
\]
and coordinate maximization gives the damped soft-transform updates
\[
f \leftarrow \omega\, g^{\bar c,\epsilon},
\qquad
g \leftarrow \omega\, f^{c,\epsilon},
\qquad
\omega\eqdef\frac{\tau}{\tau+\epsilon},
\]
where the soft $c$-transforms are those of Definition~\ref{def-continuous-soft-c-transform}. Equivalently,
\begin{align*}
f(x)
&=
-
\frac{\tau\epsilon}{\tau+\epsilon}
\log\int_{\Yy}
\exp\left(\frac{g(y)-c(x,y)}{\epsilon}\right)\d\be(y),\\
g(y)
&=
-
\frac{\tau\epsilon}{\tau+\epsilon}
\log\int_{\Xx}
\exp\left(\frac{f(x)-c(x,y)}{\epsilon}\right)\d\al(x).
\end{align*}
In the discrete case, with $K_{i,j}=e^{-\C_{i,j}/\epsilon}\a_i\b_j$, this gives the generalized Sinkhorn scaling
\[
u_i\leftarrow \left(\frac{\a_i}{(K v)_i}\right)^\omega,
\qquad
v_j\leftarrow \left(\frac{\b_j}{(\transp{K}u)_j}\right)^\omega,
\qquad
\P=\diag(u)K\diag(v).
\]
\paragraph{Metric contraction of the damped updates.}
For balanced entropic OT, potentials are only defined up to the gauge $(f,g)\sim(f+\lambda,g-\lambda)$, so Hilbert's projective metric is natural. The KL marginal penalties break this invariance: under the usual finite-dual assumptions, the dual potentials are unique up to $\al$- and $\be$-null sets, and the natural distance is the ordinary $L^\infty$ distance on potentials. In scaling variables $u=e^{f/\epsilon}$ and $v=e^{g/\epsilon}$, this is the Thompson metric on the product positive cone,
\(d_T\big((u,v),(u',v')\big) \eqdef \max\{\norm{\log u-\log u'}_\infty,\norm{\log v-\log v'}_\infty\},\)
\begin{prop}[Linear contraction of unbalanced Sinkhorn]\label{prop-unbalanced-sinkhorn-contraction}
Assume that the damped soft-transform map sends bounded potentials to bounded potentials; for instance, this holds when \(c\) is bounded. Let
$\omega=\tau/(\tau+\epsilon)<1$, and define
\(d_\infty((f,g),(f',g')) \eqdef \max\{\norm{f-f'}_\infty,\norm{g-g'}_\infty\}.\)
The damped soft-transform map
\[
\mathcal T_{\rm GS}(f,g)
\eqdef
\big(\tilde f,\tilde g\big),
\qquad
\tilde f=\omega\,g^{\bar c,\epsilon},
\qquad
\tilde g=\omega\,\tilde f^{c,\epsilon},
\]
is a contraction: \(d_\infty\big(\mathcal T_{\rm GS}(f,g),\mathcal T_{\rm GS}(f',g')\big) \leq \omega\, d_\infty((f,g),(f',g')).\)
$(f^\star,g^\star)$ and the iterates satisfy
\(d_\infty((f^{(\ell)},g^{(\ell)}),(f^\star,g^\star)) \leq \omega^\ell d_\infty((f^{(0)},g^{(0)}),(f^\star,g^\star)).\)
Equivalently, for $u^{(\ell)}=e^{f^{(\ell)}/\epsilon}$ and
$v^{(\ell)}=e^{g^{(\ell)}/\epsilon}$, \(d_T((u^{(\ell)},v^{(\ell)}),(u^\star,v^\star)) \leq \omega^\ell d_T((u^{(0)},v^{(0)}),(u^\star,v^\star)).\)
\end{prop}
\begin{proof}
The soft $c$-transform is $1$-Lipschitz in $L^\infty$. Indeed, if
$\norm{g-g'}_\infty\leq \delta$, then $g-\delta\leq g'\leq g+\delta$.
Since $g\mapsto g^{\bar c,\epsilon}$ is order reversing and satisfies
$(g+\lambda)^{\bar c,\epsilon}=g^{\bar c,\epsilon}-\lambda$, this implies
$g^{\bar c,\epsilon}-\delta \leq (g')^{\bar c,\epsilon} \leq g^{\bar c,\epsilon}+\delta$.
Hence
$\norm{g^{\bar c,\epsilon}-(g')^{\bar c,\epsilon}}_\infty\leq
\norm{g-g'}_\infty$, and the same argument applies to
$f\mapsto f^{c,\epsilon}$.

Let $(\tilde f,\tilde g)=\mathcal T_{\rm GS}(f,g)$ and
$(\tilde f',\tilde g')=\mathcal T_{\rm GS}(f',g')$. With
$D=d_\infty((f,g),(f',g'))$, the previous Lipschitz estimate gives
$\norm{\tilde f-\tilde f'}_\infty \leq \omega\norm{g-g'}_\infty \leq \omega D$,
and then
$\norm{\tilde g-\tilde g'}_\infty \leq \omega\norm{\tilde f-\tilde f'}_\infty \leq \omega^2D \leq \omega D$.
This proves the contraction. The Banach fixed-point theorem gives
existence, uniqueness and the displayed linear rate.
\end{proof}
\paragraph{Gaussian Hellinger--Kantorovich transport.}
\begin{defn}[Gaussian Hellinger--Kantorovich distance]\label{def-gaussian-ghk}
For $\alpha,\beta\in\Mm_+(\RR^d)$ and $\tau>0$, define
\begin{equation}\label{eq-gaussian-ghk-endpoint}
\operatorname{GHK}_\tau^2(\alpha,\beta)
\eqdef
\inf_{\pi\in\Mm_+(\RR^d\times\RR^d)}
\left\{
\int\norm{x-y}^2\d\pi
+\tau\KL(\pi_1|\alpha)
+\tau\KL(\pi_2|\beta)
\right\}.
\end{equation}
\end{defn}
We now specialize~\eqref{eq-gaussian-ghk-endpoint} to scaled Gaussian inputs, \(\alpha=r_0\Gaussian(m_0,\Sigma_\alpha), \qquad \beta=r_1\Gaussian(m_1,\Sigma_\beta), \qquad r_0,r_1>0,\quad \Sigma_\alpha,\Sigma_\beta\succ0,\)
where $r_i$ is the total mass and $m_i\in\RR^d$ is the mean, while $\Sigma_\alpha$ and $\Sigma_\beta$ are the respective covariances. These inputs provide one of the few continuous unbalanced problems with an explicit endpoint cost.

Write $\bar\alpha=\Gaussian(m_0,\Sigma_\alpha)$ and $\bar\beta=\Gaussian(m_1,\Sigma_\beta)$. The covariance relaxation associated with these normalized inputs is governed by a matrix Bregman divergence.

\begin{defn}[LogDet divergence and KL-softened Bures cost]\label{def-gaussian-soft-bures}
For $P,\Sigma\in\mathbb S_{++}^d$, define
\begin{equation}\label{eq-logdet-divergence-gaussian-uot}
D_{\mathrm{LD}}(P|\Sigma)
\eqdef
B_{\Phi_{\mathrm{LD}}}(P|\Sigma)
=
\tr(\Sigma^{-1}P)-\log\det(\Sigma^{-1}P)-d.
\end{equation}
For $\tau>0$, define the KL-softened Bures covariance cost
\begin{equation}\label{eq-gaussian-ghk-covariance-envelope}
\Bb_\tau^2(\Sigma_\alpha,\Sigma_\beta)
\eqdef
\min_{P,Q\in\mathbb S_{++}^d}
\left\{
\Bb^2(P,Q)
+\frac{\tau}{2}D_{\mathrm{LD}}(P|\Sigma_\alpha)
+\frac{\tau}{2}D_{\mathrm{LD}}(Q|\Sigma_\beta)
\right\},
\qquad
\Bb_\infty^2(\Sigma_\alpha,\Sigma_\beta)
\eqdef
\Bb^2(\Sigma_\alpha,\Sigma_\beta).
\end{equation}
\end{defn}
\begin{prop}[Closed form of the KL-softened Bures cost]\label{prop-gaussian-ghk-adjusted-covariances}
Set $\eta=2/\tau$ and $R_\gamma=\Sigma_\gamma^{-1}+\eta\Id$ for $\gamma\in\{\alpha,\beta\}$. Define
\begin{equation}\label{eq-gaussian-ghk-adjusted-map}
L_\tau
=
R_\beta^{-1/2}
\left(R_\beta^{1/2}R_\alpha R_\beta^{1/2}\right)^{1/2}
R_\beta^{-1/2}.
\end{equation}
Then $L_\tau R_\beta L_\tau=R_\alpha$, and the unique minimizer of~\eqref{eq-gaussian-ghk-covariance-envelope} is
\begin{equation}\label{eq-gaussian-ghk-adjusted-covariances}
P_\tau=\left[\Sigma_\alpha^{-1}+\eta(\Id-L_\tau)\right]^{-1},
\qquad
Q_\tau=L_\tau P_\tau L_\tau.
\end{equation}
Its value is
\begin{equation}\label{eq-gaussian-ghk-covariance-value}
\Bb_\tau^2(\Sigma_\alpha,\Sigma_\beta)
=
\frac{\tau}{2}
\log\frac{\det\Sigma_\alpha\det\Sigma_\beta}
{\det P_\tau\det Q_\tau}.
\end{equation}
\end{prop}
\begin{proof}
For $P,Q\succ0$, let $L$ be the Gaussian optimal map from $P$ to $Q$, so $Q=LPL$. The Frobenius gradients of the squared Bures term are $\Id-L$ and $\Id-L^{-1}$; hence stationarity of~\eqref{eq-gaussian-ghk-covariance-envelope} reads
\begin{equation}\label{eq-gaussian-ghk-covariance-stationarity}
P^{-1}=\Sigma_\alpha^{-1}+\eta(\Id-L),
\qquad
Q^{-1}=\Sigma_\beta^{-1}+\eta(\Id-L^{-1}).
\end{equation}
Substituting $Q^{-1}=L^{-1}P^{-1}L^{-1}$ gives $LR_\beta L=R_\alpha$, whose unique positive-definite solution is~\eqref{eq-gaussian-ghk-adjusted-map}. Back-substitution gives~\eqref{eq-gaussian-ghk-adjusted-covariances}. Positivity follows directly from
\(2\eta^{-1}P_\tau^{-1} =(\Id-L_\tau)^2 +\eta^{-1}\Sigma_\alpha^{-1} +\eta^{-1}L_\tau\Sigma_\beta^{-1}L_\tau \succ0.\)
The squared Bures term is jointly convex, while each LogDet term is strictly convex and coercive on $\mathbb S_{++}^d$; the stationary pair is therefore the unique global minimizer. Multiplying the two equations in~\eqref{eq-gaussian-ghk-covariance-stationarity} by $P$ and $Q$, respectively, and taking traces gives
\[
\Bb^2(P,Q)
+\frac{\tau}{2}
\left[
\tr(\Sigma_\alpha^{-1}P)+\tr(\Sigma_\beta^{-1}Q)-2d
\right]
=0.
\]
At $(P,Q)=(P_\tau,Q_\tau)$, this identity cancels the Bures term against the trace parts of the two LogDet divergences. The remaining logarithms give~\eqref{eq-gaussian-ghk-covariance-value}.
\end{proof}
\begin{prop}[Gaussian Hellinger--Kantorovich endpoint]\label{prop-gaussian-ghk-endpoint}
Set $G_\tau=\Sigma_\alpha+\Sigma_\beta+\frac{\tau}{2}\Id$. Then
\begin{equation}\label{eq-gaussian-ghk-closed-form}
\operatorname{GHK}_\tau^2(\alpha,\beta)
=
\tau\left\{
r_0+r_1
-2\sqrt{r_0r_1}\,
\exp\!\left(
-\frac14\norm{m_0-m_1}_{G_\tau^{-1}}^2
-\frac{\Bb_\tau^2(\Sigma_\alpha,\Sigma_\beta)}{2\tau}
\right)
\right\}.
\end{equation}
\end{prop}
\begin{proof}
Apply the equal-penalty KL mass--shape formula~\eqref{eq-uot-kl-mass-shape} to the normalized inputs $\bar\alpha$ and $\bar\beta$. Gelbrich's inequality, Theorem~\ref{thm-gelbrich-projection}, and the information-projection identity for KL show that replacing each competing normalized marginal by the Gaussian with the same mean and covariance cannot increase the objective. The Gaussian formulas for $\Wass_2^2$ and KL then separate the mean and covariance variables. Direct minimization over the two means gives $\frac{\tau}{2}\norm{m_0-m_1}_{G_\tau^{-1}}^2$, while Definition~\ref{def-gaussian-soft-bures} identifies the covariance minimum with $\Bb_\tau^2(\Sigma_\alpha,\Sigma_\beta)$. Substitution into~\eqref{eq-uot-kl-mass-shape} gives~\eqref{eq-gaussian-ghk-closed-form}.
\end{proof}
\paragraph{Partial optimal transport.}
\begin{defn}[Fixed-mass, penalized and TV-unbalanced transport]\label{def-partial-optimal-transport}
Let $A=\al(\Xx)$, $B=\be(\Yy)$ and $M=\min\{A,B\}$. For $0\leq m\leq M$, the fixed-mass partial transport value is
\begin{equation}\label{eq-partial-ot-fixed-mass}
\operatorname{POT}_m(\al,\be)
\eqdef
\inf_{\substack{\pi\in\Mm_+(\Xx\times\Yy)\\
\pi_1\leq\al,\ \pi_2\leq\be\\
\pi(\Xx\times\Yy)=m}}
\int c\,\d\pi.
\end{equation}
For an unmatched-mass price $\lambda\geq0$, its penalized form is
\begin{equation}\label{eq-partial-ot-tv-penalized}
\operatorname{POT}^{\TV}_\lambda(\al,\be)
\eqdef
\inf_{\substack{\pi\in\Mm_+(\Xx\times\Yy)\\
\pi_1\leq\al,\ \pi_2\leq\be}}
\left\{
\int c\,\d\pi
+\lambda\bigl(A+B-2\pi(\Xx\times\Yy)\bigr)
\right\}.
\end{equation}
Finally, the TV-unbalanced transport value is
\begin{equation}\label{eq-partial-ot-tv-unbalanced}
\UW_{c,\lambda}^{\TV}(\al,\be)
\eqdef
\inf_{\pi\in\Mm_+(\Xx\times\Yy)}
\left\{
\int c\,\d\pi
+\lambda\norm{\al-\pi_1}_{\TV}
+\lambda\norm{\be-\pi_2}_{\TV}
\right\}.
\end{equation}
\end{defn}
The last problem is precisely the unbalanced formulation~\eqref{eq-unbalanced-primal} with $\psi_1=\psi_2=\lambda\phi_{\TV}$, where $\phi_{\TV}(s)=|s-1|$ for $s\geq0$. The constraints in~\eqref{eq-partial-ot-fixed-mass} and~\eqref{eq-partial-ot-tv-penalized} mean that the transported marginals are submeasures of the available ones.

\begin{prop}[Equivalence of partial and TV-unbalanced transport]\label{prop-tv-partial-ot-lagrange}
Assume that $\Xx$ and $\Yy$ are compact metric spaces and that $c$ is continuous and nonnegative. Then, for every $\lambda\geq0$,
\begin{equation}\label{eq-partial-tv-uot-equivalence}
\UW_{c,\lambda}^{\TV}(\al,\be)
=
\operatorname{POT}^{\TV}_\lambda(\al,\be)
=
\lambda(A+B)
+
\min_{0\leq m\leq M}
\left\{
\operatorname{POT}_m(\al,\be)-2\lambda m
\right\}.
\end{equation}
The value function $m\mapsto\operatorname{POT}_m(\al,\be)$ is convex on $[0,M]$. In particular, a minimizer of the TV-unbalanced problem can be chosen with $\pi_1\leq\al$ and $\pi_2\leq\be$. If such a minimizer has mass $m_\lambda$, then it solves $\operatorname{POT}_{m_\lambda}(\al,\be)$ and $m_\lambda$ minimizes the scalar problem in~\eqref{eq-partial-tv-uot-equivalence}. Conversely, a minimizer of $\operatorname{POT}_m(\al,\be)$ solves both penalized problems whenever
\(2\lambda\in\partial\bigl[m'\mapsto\operatorname{POT}_{m'}(\al,\be)\bigr](m),\)
where the value function is extended by $+\infty$ outside $[0,M]$.
\end{prop}
\begin{proof}
Let $J_\lambda(\pi)$ denote the objective in~\eqref{eq-partial-ot-tv-unbalanced}. Starting from any $\pi$, trim its source marginal to the greatest common submeasure $\pi_1\wedge\al$, obtained by taking the pointwise minimum of the two densities with respect to $\pi_1+\al$. This is realized at the plan level by multiplying the disintegration of $\pi$ with respect to $\pi_1$ by the corresponding Radon--Nikodym factor. If $\delta$ units are removed, the source TV penalty decreases by $\lambda\delta$, while the target TV penalty increases by at most $\lambda\delta$ by the triangle inequality. Since $c\geq0$, the transport term cannot increase. Trimming the resulting target marginal to a submeasure of $\be$ gives, by the same argument, a plan $\widetilde\pi\leq\pi$ such that $\widetilde\pi_1\leq\al$, $\widetilde\pi_2\leq\be$ and $J_\lambda(\widetilde\pi)\leq J_\lambda(\pi)$.

For a subcoupling $\pi$ of mass $m$, one has
\(\norm{\al-\pi_1}_{\TV}=A-m\) and \(\norm{\be-\pi_2}_{\TV}=B-m\).
This proves the first equality in~\eqref{eq-partial-tv-uot-equivalence}. Partitioning the admissible subcouplings according to their mass then gives the second equality. Convex combinations of admissible plans show that $m\mapsto\operatorname{POT}_m$ is convex. Compactness and continuity ensure attainment of all three minima. Equality throughout the preceding bounds yields the optimizer correspondence. Finally, a mass $m$ minimizes $m'\mapsto\operatorname{POT}_{m'}-2\lambda m'$ exactly when the displayed subgradient condition holds.
\end{proof}
\begin{prop}[Fixed total mass reduces partial OT to balanced OT]\label{prop-partial-ot-metric-slice}
Specialize Definition~\ref{def-partial-optimal-transport} to $\Yy=\Xx$ and $c(x,y)=d(x,y)^p$, where $(\Xx,d)$ is a metric space and $p\geq1$. Fix $m>0$. If $\al$ and $\be$ have total mass $m$ and finite $p$th moments, then
\begin{equation}\label{eq-partial-ot-balanced-slice}
\operatorname{POT}_m(\al,\be)
=
m\Wass_p\!\left(\frac{\al}{m},\frac{\be}{m}\right)^p.
\end{equation}
\end{prop}
\begin{proof}
For an admissible plan in~\eqref{eq-partial-ot-fixed-mass}, the inequalities $\pi_1\leq\al$ and $\pi_2\leq\be$ are equalities because all four measures have mass $m$. Hence the admissible set is the usual set of couplings between $\al$ and $\be$. Dividing a coupling by $m$ identifies it with a probability coupling between $\al/m$ and $\be/m$, which proves~\eqref{eq-partial-ot-balanced-slice}. The distance claim follows from Proposition~\ref{prop-metric-measure}.
\end{proof}
\subsection{Sliced Wasserstein Distances}
\label{sec-sliced-wasserstein}
\paragraph{Spherical averaging.}
\begin{defn}[Sliced Wasserstein distance]\label{def-sliced-wasserstein}
Let $p\geq1$, let $\alpha,\beta\in\Pp_p(\RR^d)$, and let $\sigma$ be the uniform probability measure on the sphere $\Sphere^{d-1}$. The sliced $p$-Wasserstein distance is
\eql{\label{eq-sliced-wasserstein}
\SW_p(\al,\be)^p
\eqdef
\int_{\Sphere^{d-1}}
\Wass_p\big((P_\theta)_\sharp\al,(P_\theta)_\sharp\be\big)^p
\d\sigma(\theta).
}
\end{defn}
\begin{prop}[Metric properties of sliced Wasserstein]\label{prop-sliced-wasserstein-metric}
For $p\geq1$, $\SW_p$ is a distance on $\Pp_p(\RR^d)$. Moreover, $\SW_p$ metrizes weak convergence together with convergence of the $p$th moment. Finally,
\[
\SW_p(\alpha,\beta)^p
\leq
\kappa_{d,p}\Wass_p(\alpha,\beta)^p,
\qquad
\kappa_{d,p}
\eqdef
\int_{\Sphere^{d-1}}|\theta_1|^p\d\sigma(\theta)
=
\frac{\Gamma(d/2)\Gamma((p+1)/2)}{\sqrt{\pi}\,\Gamma((d+p)/2)}.
\]
In particular, $\kappa_{d,2}=1/d$. Conversely, if $d\geq2$, $R>0$, and $\alpha,\beta\in\Pp_p(\RR^d)$ are supported in
\(B_R\eqdef\enscond{x\in\RR^d}{\norm{x}\leq R}\), then there is a constant
$C_{d,p}$ depending only on $d$ and $p$ such that Bonnotte's historical estimate reads \(\Wass_p(\alpha,\beta)^p \leq C_{d,p}R^{p-\frac{1}{d+1}} \SW_p(\alpha,\beta)^{\frac{1}{d+1}}.\)
The sharp reverse estimate for $p=1$ also gives, for every $p\geq1$, constants $\widetilde C_{d,p},\widetilde c_{d,p}>0$ such that
\[
\Wass_p(\alpha,\beta)
\leq
\widetilde C_{d,p}
R^{1-\frac{1}{pd}}
\SW_p(\alpha,\beta)^{\frac{1}{pd}},
\qquad\text{equivalently}\qquad
\SW_p(\alpha,\beta)
\geq
\widetilde c_{d,p}R
\left(\frac{\Wass_p(\alpha,\beta)}{R}\right)^{pd}.
\]
For $p=1$, the exponent $1/d$ in this last comparison is sharp.
\end{prop}
\begin{proof}
Non-negativity and symmetry follow from the one-dimensional Wasserstein distance. For the triangle inequality, apply the triangle inequality of $\Wass_p$ for each direction $\theta$ and then Minkowski's inequality in $L^p(\Sphere^{d-1})$.

If $\SW_p(\alpha,\beta)=0$, then $(P_\theta)_\sharp\alpha=(P_\theta)_\sharp\beta$ for almost every direction. By continuity of characteristic functions this holds for all directions, and the Cram\'er--Wold theorem implies $\alpha=\beta$. This proves separation.

Let $\pi$ be any coupling between $\alpha$ and $\beta$. Projecting $\pi$ gives an admissible coupling between the corresponding one-dimensional projections, hence
\(\Wass_p\big((P_\theta)_\sharp\alpha,(P_\theta)_\sharp\beta\big)^p \leq \int_{\RR^d\times\RR^d}|\dotp{\theta}{x-y}|^p\d\pi(x,y).\)
Rotational invariance of $\sigma$ gives
$\int_{\Sphere^{d-1}}|\dotp{\theta}{z}|^p\d\sigma(\theta) = \kappa_{d,p}\norm{z}^p$.
Integrating the first inequality in $\theta$ and optimizing over $\pi$ proves the stated direct comparison. The beta-integral formula for the first coordinate of a uniform point on the sphere gives the displayed value of $\kappa_{d,p}$; in particular, $\kappa_{d,2}=1/d$.

To prove the topological statement, first suppose that $\Wass_p(\alpha_n,\alpha)\to0$. The inequality $\SW_p\leq\Wass_p$ immediately gives $\SW_p(\alpha_n,\alpha)\to0$. Conversely, suppose that $\SW_p(\alpha_n,\alpha)\to0$. The spherical identity
\[
\int_{\Sphere^{d-1}}\int_{\RR^d}|\langle\theta,x\rangle|^p\d\alpha_n(x)\d\sigma(\theta)
=
\kappa_{d,p}\int_{\RR^d}\|x\|^p\d\alpha_n(x),
\qquad
\kappa_{d,p}>0,
\]
and the triangle inequality in $L^p(\Sphere^{d-1}\times(0,1))$ show that the $p$th moments of $(\alpha_n)$ are uniformly bounded. Hence the sequence is tight. From every subsequence one can extract a further subsequence for which the projected $\Wass_p$ distances converge to zero for $\sigma$-a.e.\ $\theta$. Every weak limit therefore has the same one-dimensional projections as $\alpha$, so Cram\'er--Wold identifies it with $\alpha$. Thus $\alpha_n\rightharpoonup\alpha$. Finally, let $Q_{n,\theta}$ and $Q_\theta$ be the quantile functions of the projected measures. The one-dimensional quantile formula gives
\(\SW_p(\alpha_n,\alpha) = \|Q_{n,\theta}-Q_\theta\|_{L^p(\Sphere^{d-1}\times(0,1))}.\)
Hence the corresponding $L^p$ norms converge. By the spherical identity, their $p$th powers are respectively $\kappa_{d,p}\int\|x\|^p\d\alpha_n(x)$ and $\kappa_{d,p}\int\|x\|^p\d\alpha(x)$. Thus the $p$th moments converge. Weak convergence plus convergence of these moments is equivalent to convergence in $\Wass_p$.

The reverse estimates on compact sets are deeper and use the Radon structure of slices. Bonnotte's Lemma~5.1.4 proves
$\Wass_1(\alpha,\beta) \leq C_dR^{\frac{d}{d+1}}\SW_1(\alpha,\beta)^{\frac{1}{d+1}}$
by mollifying a Kantorovich--Rubinstein test function, representing the regularized test through one-dimensional projections, bounding the resulting action by $\SW_1$, and optimizing the smoothing scale. On $B_R$,
$\Wass_p^p\leq(2R)^{p-1}\Wass_1$ and $\SW_1\leq\SW_p$,
where the first inequality follows by evaluating the $p$-cost on a $\Wass_1$-optimal coupling, and the second follows from the monotonicity of one-dimensional Wasserstein distances and H\"older's inequality in $\theta$. These inequalities give Bonnotte's stated general-$p$ estimate. The sharper $p=1$ comparison theorem of Carlier, Figalli, M\'erigot and Wang replaces $1/(d+1)$ by the optimal exponent $1/d$.
$\Wass_p^p \leq C_{d,p}R^{p-\frac1d}\SW_p^{\frac1d}$,
which is equivalent to the final two-sided formulation in the statement. This supplies a lower bound on $\SW_p$ in terms of $\Wass_p$ for every $p$, but the asserted sharpness concerns only $p=1$.
\end{proof}
\paragraph{Radon point of view.}\label{rem-sliced-radon-viewpoint}
\begin{defn}[Measure-valued Radon transform]\label{def-measure-radon-transform}
For $\al\in\Pp(\RR^d)$, define
\(\mathfrak R\al(\theta)\eqdef(P_\theta)_\sharp\al, \qquad \theta\in\Sphere^{d-1}.\)
\end{defn}
Slicing can also be understood as applying Definition~\ref{def-measure-radon-transform} before measuring discrepancy. Thus $\mathfrak R\al$ is a collection of one-dimensional projected measures indexed by directions. If $\al=\rho(x)\d x$ has a density, then $\mathfrak R\al(\theta)$ has density given by the classical Radon transform
\(R\rho(\theta,s) = \int_{\{x:\dotp{\theta}{x}=s\}}\rho(x)\,\d\mathcal H^{d-1}(x).\)
\[
d_{\mathrm{Rad},p}(h_0,h_1)
\eqdef
\left(
\int_{\Sphere^{d-1}}
\Wass_p\big(h_0(\theta),h_1(\theta)\big)^p
\d\sigma(\theta)
\right)^{1/p}.
\]
Indeed, associate with each Radon field $h$ the joint law
\(\widehat h(\d\theta,\d s) \eqdef h(\theta)(\d s)\sigma(\d\theta)\)
Its first marginal is $\sigma$, so $\widehat h\in\Pp_{p,\sigma}(\Sphere^{d-1}\times\RR)$, and Definition~\ref{def-conditional-wasserstein-distance}, with condition space $S=\Sphere^{d-1}$ and fiber $\Omega=\RR$, gives
\(d_{\mathrm{Rad},p}(h_0,h_1) = \Wass_{p,\sigma}(\widehat h_0,\widehat h_1).\)
For any map $F:E\to(Y,d_Y)$, the pull-back of the metric $d_Y$ through $F$ is the function on $E\times E$ defined by \((F^*d_Y)(x,x')\eqdef d_Y\big(F(x),F(x')\big).\)
Consequently, \eqref{eq-sliced-wasserstein} reads \(\SW_p(\alpha,\beta) = d_{\mathrm{Rad},p}(\mathfrak R\alpha,\mathfrak R\beta) = (\mathfrak R^*d_{\mathrm{Rad},p})(\alpha,\beta).\)
\paragraph{Infinitesimal SW geometry.}
\begin{prop}[First-order sliced comparison along Brenier perturbations]\label{prop-sliced-first-order-tangent}
Let $\alpha\in\Pp_2(\RR^d)$ be absolutely continuous, and let $\varphi\in C^2(\RR^d)$ be such that $\nabla\varphi\in L^2(\alpha;\RR^d)$ and $\nabla^2\varphi$ is bounded. For $\abs{t}$ small enough, set
\(T_t(x)\eqdef x+t\nabla\varphi(x), \qquad \alpha_t\eqdef(T_t)_\sharp\alpha.\)
Then $T_t$ is the quadratic Brenier map from $\alpha$ to $\alpha_t$, and
\(\Wass_2(\alpha,\alpha_t)^2 = t^2\int_{\RR^d}\norm{\nabla\varphi(x)}^2\d\alpha(x).\)
Moreover, \(\SW_2(\alpha,\alpha_t)^2 \leq \frac{t^2}{d}\int_{\RR^d}\norm{\nabla\varphi(x)}^2\d\alpha(x).\)
\end{prop}
\begin{proof}
The map $T_t$ is the gradient of
\(\psi_t(x)\eqdef \frac12\norm{x}^2+t\varphi(x).\)
Since $\nabla^2\varphi$ is bounded, $\nabla^2\psi_t=\Id+t\nabla^2\varphi$ is positive semidefinite for $\abs{t}$ small enough. Hence $\psi_t$ is convex and $T_t=\nabla\psi_t$ is a Brenier map. Its graph is cyclically monotone, so
\(\Wass_2(\alpha,\alpha_t)^2 = \int\norm{x-T_t(x)}^2\d\alpha(x) = t^2\int\norm{\nabla\varphi(x)}^2\d\alpha(x).\)
For a fixed direction $\theta$, the coupling
\(\big(P_\theta,P_\theta\circ T_t\big)_\sharp\alpha\)
is admissible between $(P_\theta)_\sharp\alpha$ and $(P_\theta)_\sharp\alpha_t$. Therefore
\(\Wass_2\big((P_\theta)_\sharp\alpha,(P_\theta)_\sharp\alpha_t\big)^2 \leq t^2\int\abs{\dotp{\theta}{\nabla\varphi(x)}}^2\d\alpha(x).\)
Integrating over the sphere and using
\(\int_{\Sphere^{d-1}}\abs{\dotp{\theta}{w}}^2\d\sigma(\theta) = \frac1d\norm{w}^2\)
gives the sliced estimate.
\end{proof}
\paragraph{Intrinsic sliced length.}\label{par-sliced-intrinsic-length}
Its intrinsic, or path, metric is
\begin{equation}\label{eq-intrinsic-sliced-length}
\ell_{\SW_2}(\alpha,\beta)
\eqdef
\inf_{\substack{\gamma_0=\alpha,\ \gamma_1=\beta\\
\gamma\ \text{$\SW_2$-absolutely continuous}}}
\int_0^1 \abs{\dot\gamma_t}_{\SW_2}\,\d t,
\end{equation}
where \(\abs{\dot\gamma_t}_{\SW_2} \eqdef \lim_{h\to0}\frac{\SW_2(\gamma_{t+h},\gamma_t)}{\abs{h}}\)
Consequently, \(\SW_2\leq\ell_{\SW_2}\leq d^{-1/2}\Wass_2.\)
Integrating along that $\Wass_2$-geodesic gives, for every sufficiently small $t\neq0$, \(\ell_{\SW_2}(\alpha,\alpha_t) <d^{-1/2}\Wass_2(\alpha,\alpha_t).\)
\paragraph{Sliced Wasserstein between Gaussians.}
\begin{prop}[Sliced Wasserstein between Gaussians]\label{prop-sliced-gaussian}
Let $\alpha=\Gaussian(m_\alpha,\Sigma_\alpha)$ and
$\beta=\Gaussian(m_\beta,\Sigma_\beta)$ on $\RR^d$, with positive-semidefinite
covariance matrices. Then
\begin{equation}\label{eq-sliced-gaussian}
\SW_2(\alpha,\beta)^2
=
\frac{\norm{m_\alpha-m_\beta}^2}{d}
+
\int_{\Sphere^{d-1}}
\pa{
\sqrt{\transp{\theta}\Sigma_\alpha\theta}
-
\sqrt{\transp{\theta}\Sigma_\beta\theta}
}^2
\d\sigma(\theta).
\end{equation}
\end{prop}
\begin{proof}
For every $\theta\in\Sphere^{d-1}$, \((P_\theta)_\sharp\alpha = \Gaussian\bigl(\transp{\theta}m_\alpha, \transp{\theta}\Sigma_\alpha\theta\bigr),\)
and similarly for $\beta$. The one-dimensional Gaussian formula gives
\[
\Wass_2\bigl((P_\theta)_\sharp\alpha,
(P_\theta)_\sharp\beta\bigr)^2
=
\abs{\transp{\theta}(m_\alpha-m_\beta)}^2
+
\pa{
\sqrt{\transp{\theta}\Sigma_\alpha\theta}
-
\sqrt{\transp{\theta}\Sigma_\beta\theta}
}^2.
\]
Integrating over $\theta$ and using
$\int_{\Sphere^{d-1}}\theta\transp{\theta}\,\d\sigma(\theta)=\Id/d$
gives~\eqref{eq-sliced-gaussian}.
\end{proof}
\paragraph{\texorpdfstring{$L^q$-sliced, max-sliced and subspace variants.}{L^q-sliced, max-sliced and subspace variants.}}
\begin{defn}[$L^q$-sliced and subspace-sliced Wasserstein]\label{def-sliced-variants}
Let $p\in[1,\infty)$, $q\in[1,\infty]$, and $1\leq k\leq d$. Denote by \(\mathrm{St}(d,k) \eqdef \enscond{U\in\RR^{d\times k}}{\transp{U}U=\Id_k}\)
the Stiefel manifold, equipped with its normalized invariant probability
measure $\sigma_{d,k}$. For $q<\infty$, define
\begin{equation}\label{eq-lq-subspace-sliced}
\SW_{p,q,k}(\alpha,\beta)
\eqdef
\left(
\int_{\mathrm{St}(d,k)}
\Wass_p\big((\transp{U})_\sharp\alpha,
(\transp{U})_\sharp\beta\big)^q
\d\sigma_{d,k}(U)
\right)^{1/q},
\end{equation}
and extend the definition to $q=\infty$ by \(\SW_{p,\infty,k}(\alpha,\beta) \eqdef \sup_{U\in\mathrm{St}(d,k)} \Wass_p\big((\transp{U})_\sharp\alpha, (\transp{U})_\sharp\beta\big).\)
We use the abbreviations
\(\SW_{p,q}\eqdef\SW_{p,q,1}, \qquad \SW_p=\SW_{p,p}, \qquad \MaxSW_{p,k}\eqdef\SW_{p,\infty,k}, \qquad \MaxSW_p=\SW_{p,\infty}.\)
Thus $k=1$ gives line-based slicing, $q=\infty$ gives max-slicing,
and $k=d$ gives $\SW_{p,q,d}=\Wass_p$ for every $q$. Quadratic max-slicing is useful
when only a small set of directions carries the discrepancy, for instance in
generative modeling. For $p=2$, the $k$-dimensional
max-sliced case is the projection-robust Wasserstein distance
$\MaxSW_{2,k}=\PRW_{2,k}$, which is related to the
subspace-robust distance of
Definition~\ref{def-projection-subspace-robust-wasserstein}, studied later in
Section~\ref{sec-spectral-subspace-wasserstein}.
\end{defn}
The supremum in the max-sliced case is attained because the Stiefel manifold is compact and the projected Wasserstein distance depends continuously on $U$. Moreover, the integrand is unchanged under $U\mapsto UQ$ for $Q\in\mathrm O(k)$, because $Q^\top$ is an isometry of $\RR^k$.

\begin{prop}[Basic bounds for sliced variants]\label{prop-sliced-variant-bounds}
Let $p\in[1,\infty)$, $1\leq q\leq r\leq\infty$, and
$1\leq k\leq \ell\leq d$. Each $\SW_{p,q,k}$ is a finite distance on
$\Pp_p(\RR^d)$ and metrizes the same topology as $\Wass_p$. For every
$\alpha,\beta\in\Pp_p(\RR^d)$,
\begin{align}
\SW_{p,q,k}(\alpha,\beta)
&\leq
\SW_{p,r,k}(\alpha,\beta)
\leq
\Wass_p(\alpha,\beta),
\label{eq-lq-sliced-monotonicity}\\
\SW_{p,q,k}(\alpha,\beta)
&\leq
\SW_{p,q,\ell}(\alpha,\beta)
\leq
\Wass_p(\alpha,\beta).
\label{eq-subspace-dimension-monotonicity}
\end{align}
For a unit vector $e\in\Sphere^{d-1}$, set
\[
\kappa_{d,k,p}
\eqdef
\int_{\mathrm{St}(d,k)}\norm{\transp{U}e}^p\d\sigma_{d,k}(U)
=
\frac{\Gamma(d/2)\Gamma((k+p)/2)}
{\Gamma(k/2)\Gamma((d+p)/2)}.
\]
This constant is independent of $e$, satisfies $\kappa_{d,d,p}=1$ and
$\kappa_{d,1,p}=\kappa_{d,p}$ from
Proposition~\ref{prop-sliced-wasserstein-metric}, and gives the sharper bounds
\begin{equation}\label{eq-lq-subspace-sliced-upper-bound}
\SW_{p,q,k}(\alpha,\beta)
\leq
\begin{cases}
\kappa_{d,k,p}^{1/p}\Wass_p(\alpha,\beta), & 1\leq q\leq p,\\
\kappa_{d,k,p}^{1/q}\Wass_p(\alpha,\beta), & p\leq q<\infty,\\
\Wass_p(\alpha,\beta), & q=\infty.
\end{cases}
\end{equation}
\eqllead{Moreover,}{eq-subspace-sliced-lower-bound}{
\SW_p(\alpha,\beta)
\leq
\SW_{p,p,k}(\alpha,\beta)
\leq
\kappa_{d,k,p}^{1/p}\Wass_p(\alpha,\beta).
}
In particular, $\kappa_{d,k,2}=k/d$ and
$\SW_{2,2,k}^2\leq(k/d)\Wass_2^2$. If $q\geq p$, the compact-support reverse
estimates in Proposition~\ref{prop-sliced-wasserstein-metric} remain valid
with $\SW_{p,q,k}$ in place of $\SW_p$.
\end{prop}
\begin{proof}
Write
$D_U(\alpha,\beta)\eqdef
\Wass_p((\transp{U})_\sharp\alpha,(\transp{U})_\sharp\beta)$.
Non-negativity and symmetry are inherited from $\Wass_p$.
Minkowski's inequality in $L^q(\sigma_{d,k})$, or the supremum triangle
inequality when $q=\infty$, proves the triangle inequality. The reverse
triangle inequality for $\Wass_p$ shows that $U\mapsto D_U(\alpha,\beta)$ is
continuous. Indeed, for $U_n\to U$,
\[
\Wass_p\big((\transp{U_n})_\sharp\alpha,(\transp{U})_\sharp\alpha\big)
\leq
\norm{U_n-U}_{\mathrm{op}}
\left(\int\norm{x}^p\d\alpha(x)\right)^{1/p},
\]
and likewise for $\beta$. If $\SW_{p,q,k}(\alpha,\beta)=0$, then $D_U=0$ for
all $U$: this is immediate for $q=\infty$, while for finite $q$ it follows
from continuity and the full support of $\sigma_{d,k}$. For any
$\theta\in\Sphere^{d-1}$, choose a frame $U$ whose first column is $\theta$.
Equality of the two $k$-dimensional projected measures then implies equality
of their first-coordinate marginals, hence
$(P_\theta)_\sharp\alpha=(P_\theta)_\sharp\beta$. The Cram\'er--Wold theorem
gives $\alpha=\beta$, proving separation.

Every $U^\top$ is $1$-Lipschitz, so $D_U\leq\Wass_p(\alpha,\beta)$.
Monotonicity of $L^q$ norms on a probability space proves
\eqref{eq-lq-sliced-monotonicity}. To compare projection dimensions, let
$V\in\mathrm{St}(d,\ell)$ and $R\in\mathrm{St}(\ell,k)$. Since
$(VR)^\top=R^\top V^\top$ and $R^\top$ is $1$-Lipschitz,
$D_{VR}(\alpha,\beta)\leq D_V(\alpha,\beta)$. Under the product of the
invariant measures, $VR$ is uniformly distributed on $\mathrm{St}(d,k)$.
Integration proves~\eqref{eq-subspace-dimension-monotonicity} for finite $q$;
for $q=\infty$, extend each $k$-frame to an $\ell$-frame and use the same
contraction.

For the sharper constants, let $\pi\in\Gamma(\alpha,\beta)$. Then
$D_U(\alpha,\beta)^p \leq \int\norm{\transp{U}(x-y)}^p\d\pi(x,y)$.
Integration in $U$, followed by minimization in $\pi$, gives
$\norm{D_\cdot}_{L^p(\sigma_{d,k})}^p
\leq\kappa_{d,k,p}\Wass_p(\alpha,\beta)^p$. If $q\leq p$, use
$\norm{D_\cdot}_{L^q}\leq\norm{D_\cdot}_{L^p}$. If $p\leq q<\infty$, use
$D_U\leq\Wass_p$ to obtain
\(\int D_U^q\d\sigma_{d,k}(U) \leq \Wass_p(\alpha,\beta)^{q-p} \int D_U^p\d\sigma_{d,k}(U).\)
This proves~\eqref{eq-lq-subspace-sliced-upper-bound}. The beta-integral for
the $p$th power of the length of a random $k$-dimensional projection gives
the displayed formula for $\kappa_{d,k,p}$.

Finally, let $\omega\in\Sphere^{k-1}$. Since
$P_{U\omega}=P_\omega\circ U^\top$, projection contraction gives \(\Wass_p\big((P_{U\omega})_\sharp\alpha, (P_{U\omega})_\sharp\beta\big) \leq D_U(\alpha,\beta).\)
Averaging its $p$th power first over $\omega$ and then over $U$ makes
$U\omega$ uniform on $\Sphere^{d-1}$ and proves the first inequality
in~\eqref{eq-subspace-sliced-lower-bound}. For $q\geq p$,
\eqref{eq-lq-sliced-monotonicity} gives
$\SW_p\leq\SW_{p,p,k}\leq\SW_{p,q,k}$, which transfers the stated reverse
estimates.

For
$\eta\in\Pp_p(\RR^d)$, set
$M_p(\eta)\eqdef(\int\norm{x}^p\d\eta(x))^{1/p}$. Since the only coupling
between $(U^\top)_\sharp\eta$ and $\delta_0$ sends every point to the origin, \(D_U(\eta,\delta_0)^p = \int\norm{U^\top x}^p\d\eta(x), \qquad \int D_U(\eta,\delta_0)^p\d\sigma_{d,k}(U) = \kappa_{d,k,p}M_p(\eta)^p.\)
Also $D_U(\eta,\delta_0)\leq M_p(\eta)$.
\begin{equation}\label{eq-lq-sliced-radial-moment}
\SW_{p,q,k}(\eta,\delta_0)
\geq
\begin{cases}
\kappa_{d,k,p}^{1/q}M_p(\eta), & 1\leq q\leq p,\\
\kappa_{d,k,p}^{1/p}M_p(\eta), & p\leq q\leq\infty.
\end{cases}
\end{equation}
Indeed, the second line follows from monotonicity of $L^q$ norms. For
$q\leq p$ and $M_p(\eta)>0$, the pointwise bound
$D_U\leq M_p(\eta)$ gives
$D_U^q\geq M_p(\eta)^{q-p}D_U^p$; the zero-moment case is immediate.

Suppose now that $\SW_{p,q,k}(\alpha_n,\alpha)\to0$. The triangle inequality
and~\eqref{eq-lq-sliced-radial-moment} imply
$\sup_n M_p(\alpha_n)<\infty$. Write
$D_{n,U}\eqdef D_U(\alpha_n,\alpha)$. The reverse triangle inequality and
the coupling of each point with itself give
\(\abs{D_{n,U}-D_{n,V}} \leq \norm{U-V}_{\mathrm{op}} \big(M_p(\alpha_n)+M_p(\alpha)\big),\)
so $(D_{n,\cdot})_n$ is equicontinuous on the compact Stiefel manifold. For
$q<\infty$, equicontinuity and the full support of $\sigma_{d,k}$ upgrade
$\norm{D_{n,\cdot}}_{L^q}\to0$ to $\norm{D_{n,\cdot}}_{L^\infty}\to0$;
for $q=\infty$ this is immediate. Given $\theta\in\Sphere^{d-1}$, choose a
frame $U$ whose first column is $\theta$. Projection contraction then yields
\(\Wass_p\big((P_\theta)_\sharp\alpha_n,(P_\theta)_\sharp\alpha\big) \leq D_{n,U},\)
uniformly in $\theta$. Hence $\SW_p(\alpha_n,\alpha)\to0$, and
Proposition~\ref{prop-sliced-wasserstein-metric} gives
$\Wass_p(\alpha_n,\alpha)\to0$. The converse follows directly from
$\SW_{p,q,k}\leq\Wass_p$.
\end{proof}
\paragraph{Min-SW lifted transport plans.}
For a direction $\theta$ along which both projected point families have distinct coordinates, let $\sigma_\theta,\tau_\theta\in\mathfrak S_n$ be their sorting permutations,
\[
\dotp{x_{\sigma_\theta(1)}}{\theta}<\cdots<
\dotp{x_{\sigma_\theta(n)}}{\theta},
\qquad
\dotp{y_{\tau_\theta(1)}}{\theta}<\cdots<
\dotp{y_{\tau_\theta(n)}}{\theta}.
\]
Lifting the one-dimensional monotone matching gives
\begin{equation}\label{eq-min-sw-empirical-plan}
\pi_\theta
=
\frac1n\sum_{i=1}^n
\delta_{(x_{\sigma_\theta(i)},y_{\tau_\theta(i)})}
\in\Couplings(\alpha,\beta).
\end{equation}
Set
\(\alpha_\theta=(P_\theta)_\sharp\alpha, \qquad \beta_\theta=(P_\theta)_\sharp\beta,\)
and let
\(\varpi_\theta = (q_{\alpha_\theta},q_{\beta_\theta})_\sharp \mathrm{Leb}_{[0,1]}\)
\begin{defn}[Min-SW discrepancy and lifted plans]\label{def-min-sw}
\begin{equation}\label{eq-min-sw-constrained-lifts}
\mathcal C_\theta(\alpha,\beta)
\eqdef
\left\{
\pi\in\Couplings(\alpha,\beta)
\;:\;
(P_\theta,P_\theta)_\sharp\pi=\varpi_\theta
\right\}.
\end{equation}
The Min-SW discrepancy is
\begin{equation}\label{eq-min-sw-definition}
\MinSW_2(\alpha,\beta)^2
\eqdef
\min_{\theta\in\Sphere^{d-1}}
\min_{\pi\in\mathcal C_\theta(\alpha,\beta)}
\int_{\RR^d\times\RR^d}\norm{x-y}^2\d\pi(x,y).
\end{equation}
\end{defn}
Indeed, disintegrate
\(\alpha(\d x)=\int_\RR \alpha^s(\d x)\d\alpha_\theta(s), \qquad \beta(\d y)=\int_\RR \beta^t(\d y)\d\beta_\theta(t)\)
The conditionally independent lift
\begin{equation}\label{eq-min-sw-fiber-lift}
\pi_\theta^0(\d x,\d y)
=
\int_{\RR^2}
\alpha^s(\d x)\beta^t(\d y)
\d\varpi_\theta(s,t)
\end{equation}
\begin{prop}[Min-SW lower bound, topology and metric status]\label{prop-min-sw-comparison}
For all $\alpha,\beta\in\Pp_2(\RR^d)$,
\begin{equation}\label{eq-min-sw-comparison}
\Wass_2(\alpha,\beta)
\leq
\MinSW_2(\alpha,\beta).
\end{equation}
The converse fails in every dimension $d\geq2$: there are measures $\alpha_n$ such that
\begin{equation}\label{eq-min-sw-topological-counterexample}
\Wass_2(\alpha_n,\gamma_d)\longrightarrow0
\qquad\text{but}\qquad
\inf_n\MinSW_2(\alpha_n,\gamma_d)>0,
\qquad
\gamma_d=\Gaussian(0,\Id_d).
\end{equation}
Thus $\MinSW_2$ does not metrize weak convergence, or even $\Wass_2$ convergence, on $\Pp_2(\RR^d)$, and no estimate
$\MinSW_2(\alpha,\beta)\leq\omega(\Wass_2(\alpha,\beta))$ can hold for all measures with a modulus $\omega(r)\to0$ as $r\to0$. In particular, there is no global H\"older bound $\MinSW_2\leq C\Wass_2^a$ with $a>0$.
The function $\MinSW_2$ is non-negative, symmetric and separates measures, but it does not satisfy the triangle inequality in general. It is therefore a separating symmetric discrepancy, not a distance. In dimension one, by contrast, $\MinSW_2=\Wass_2$.
\end{prop}
\begin{proof}
Every $\pi\in\mathcal C_\theta(\alpha,\beta)$ is an admissible coupling, which proves~\eqref{eq-min-sw-comparison}. Symmetry follows by transposing the constrained plans, and separation follows from the same lower bound.

To see that the triangle inequality can fail, let $\alpha_X,\alpha_Y,\alpha_Z$ be the uniform laws on the rows of
\[
X=\begin{pmatrix}-4&-4\\0&3\\1&-2\\4&-4\end{pmatrix},
\qquad
Z=\begin{pmatrix}-3&2\\2&2\\0&0\\2&-4\end{pmatrix},
\qquad
Y=\frac{X+Z}{2}.
\]
A direct enumeration of the $4!$ assignments shows that the identity assignment is quadratically optimal for both $(X,Y)$ and $(Y,Z)$. It is induced by sorting along directions proportional to $(2,-1)$ and $(3,1)$, respectively. Hence
\(\MinSW_2(\alpha_X,\alpha_Y)^2 = \MinSW_2(\alpha_Y,\alpha_Z)^2 =\frac{51}{16}.\)
For $(X,Z)$, the projected orders change only across the ten lines orthogonal to pairwise differences of rows of $X$ or $Z$. Since the cost is linear in the plan, it suffices to enumerate the resulting angular cells and, on their boundaries, the compatible extreme permutations.
\(\frac14\{63,67,73,79,91,95,105,131,135\}\).
Thus $\MinSW_2(\alpha_X,\alpha_Z)^2=63/4$, and consequently
\(\MinSW_2(\alpha_X,\alpha_Z)=\frac{\sqrt{63}}{2} > \frac{\sqrt{51}}{2} = \MinSW_2(\alpha_X,\alpha_Y) + \MinSW_2(\alpha_Y,\alpha_Z).\)

Write $\gamma_k=\Gaussian(0,\Id_k)$ and let
\begin{equation}\label{eq-min-sw-quantization-obstruction}
e_d^2
\eqdef
\inf_{z_1,\ldots,z_d\in\RR^{d-1}}
\int_{\RR^{d-1}}
\min_{1\leq j\leq d}\norm{u-z_j}^2
\d\gamma_{d-1}(u)>0.
\end{equation}
This is the optimal quadratic quantization error of a $(d-1)$-dimensional standard Gaussian by at most $d$ points. It is positive because the set of probability measures supported on at most $d$ points is closed for $\Wass_2$, whereas $\gamma_{d-1}$ is not finitely supported.

Choose a sequence $\alpha_n=n^{-1}\sum_{i=1}^n\delta_{x_i}$ converging to $\gamma_d$ in $\Wass_2$ and in affine general position, meaning that no affine hyperplane contains more than $d$ of the $x_i$. Such a deterministic sequence exists: an i.i.d.\ Gaussian sequence has both properties almost surely. Fix $\theta$ and $\pi\in\mathcal C_\theta(\alpha_n,\gamma_d)$, and let $(X,Y)\sim\pi$. Write
$S=\dotp{\theta}{X}$, $T=\dotp{\theta}{Y}$ and $Z=P_{\theta^\perp}Y$, where $P_{\theta^\perp}=\Id_d-\theta\theta^\top$ is identified with an orthogonal projection onto $\RR^{d-1}$. The monotone coupling between the discrete law of $S$ and the atomless Gaussian law of $T$ makes $S$ a function of $T$. Since $Y\sim\gamma_d$, the variables $T$ and $Z\sim\gamma_{d-1}$ are independent; hence the conditional law of $Z$ given $S$ is still $\gamma_{d-1}$. For each value $s$ of $S$, affine general position gives
\(\#\{i:\dotp{\theta}{x_i}=s\}\leq d.\)
Conditionally on $S=s$, the vector $P_{\theta^\perp}X$ therefore takes at most $d$ values. By~\eqref{eq-min-sw-quantization-obstruction},
\[
\EE\!\left[\norm{P_{\theta^\perp}(X-Y)}^2\mid S=s\right]\geq e_d^2.
\]
Averaging and discarding the non-negative displacement along $\theta$ gives
$\int\norm{x-y}^2\d\pi(x,y)\geq e_d^2$. This holds for every admissible $\pi$ and every $\theta$, so
$\MinSW_2(\alpha_n,\gamma_d)\geq e_d$, which proves~\eqref{eq-min-sw-topological-counterexample} and rules out every vanishing modulus $\omega$.
\end{proof}
\begin{prop}[Min-SW between nondegenerate Gaussians]\label{prop-min-sw-gaussians}
Let $\alpha=\Gaussian(m_\alpha,\Sigma_\alpha)$ and $\beta=\Gaussian(m_\beta,\Sigma_\beta)$, where $\Sigma_\alpha,\Sigma_\beta\succ0$. Then
\begin{equation}\label{eq-min-sw-gaussians}
\MinSW_2(\alpha,\beta)^2
=
\Wass_2(\alpha,\beta)^2
=
\norm{m_\alpha-m_\beta}^2
+
\tr\pa{
\Sigma_\alpha+\Sigma_\beta
-2\pa{\Sigma_\alpha^{1/2}\Sigma_\beta\Sigma_\alpha^{1/2}}^{1/2}
}.
\end{equation}
\end{prop}
\begin{proof}
Let $T(x)=m_\beta+A(x-m_\alpha)$ be the Gaussian Brenier map of Proposition~\ref{prop-gaussian-w2-bures}, where the symmetric positive-definite matrix $A$ is given by~\eqref{eq-bures-map}. Choose a unit eigenvector $\theta$ of $A$, with $A\theta=\lambda\theta$ and $\lambda>0$. If $X\sim\alpha$ and $Y=T(X)$, then
\(\dotp{\theta}{Y-m_\beta} = \dotp{A\theta}{X-m_\alpha} = \lambda\dotp{\theta}{X-m_\alpha}.\)
Hence the optimal Gaussian coupling $(\Id,T)_\sharp\alpha$ belongs to $\mathcal C_\theta(\alpha,\beta)$ in Definition~\ref{def-min-sw}. It is feasible for~\eqref{eq-min-sw-definition}, which gives $\MinSW_2\leq\Wass_2$. The reverse inequality is~\eqref{eq-min-sw-comparison}; formula~\eqref{eq-dist-gauss} concludes the proof.
\end{proof}
\subsection{Quotient Wasserstein and Wasserstein--Procrustes}
\label{sec-quotient-wasserstein-procrustes}
\paragraph{Quotient distances.}
\begin{defn}[Quotient Wasserstein distance]\label{def-quotient-wasserstein}
Let a group $\mathcal G$ act on a metric space $(\Xx,d)$ by isometries, and assume that the action preserves finite $p$th moments. Write $g_\sharp\alpha$ for the push-forward of a measure $\alpha$ by $g\in\mathcal G$, and $[\alpha]\eqdef\{g_\sharp\alpha\;:\;g\in\mathcal G\}$ for its orbit. The quotient $p$-Wasserstein distance between the orbits of $\alpha,\beta\in\Pp_p(\Xx)$ is
\[
\Wass_{p,\mathcal G}([\alpha],[\beta])
\eqdef
\inf_{g,h\in\mathcal G}\Wass_p(g_\sharp\alpha,h_\sharp\beta)
=
\inf_{g\in\mathcal G}\Wass_p(\alpha,g_\sharp\beta).
\]
\end{defn}
\begin{prop}[Metric property on the quotient]\label{prop-quotient-wasserstein-metric}
If $\mathcal G$ acts by isometries and preserves $\Pp_p(\Xx)$, then $\Wass_{p,\mathcal G}$ is a pseudo-distance on $\Pp_p(\Xx)$. If the infimum is attained between every pair of orbits, then it is a distance on the orbit space $\Pp_p(\Xx)/\mathcal G$. This applies in particular to continuous actions of compact groups.
\end{prop}
\begin{proof}
Isometry of the action gives $\Wass_p(g_\sharp\alpha,g_\sharp\beta)=\Wass_p(\alpha,\beta)$, which proves the second formula in Definition~\ref{def-quotient-wasserstein}. Non-negativity and symmetry are inherited from $\Wass_p$. For the triangle inequality, choose $g_1,g_2\in\mathcal G$. Then
\[
\Wass_p(\alpha,(g_1)_\sharp\beta)
+
\Wass_p(\beta,(g_2)_\sharp\gamma)
=
\Wass_p(\alpha,(g_1)_\sharp\beta)
+
\Wass_p((g_1)_\sharp\beta,(g_1g_2)_\sharp\gamma).
\]
Applying the triangle inequality for $\Wass_p$ and then taking infima over $g_1,g_2$ gives the result. If the infimum is attained and the quotient distance is zero, there exist $g,h\in\mathcal G$ such that $\Wass_p(g_\sharp\alpha,h_\sharp\beta)=0$, hence $g_\sharp\alpha=h_\sharp\beta$ and therefore $[\alpha]=[\beta]$. For a continuous action of a compact group, the objective is continuous on a compact parameter set, so attainment follows. Without attainment, the construction is naturally a metric only after identifying distinct orbits at zero orbit distance.
\end{proof}
\paragraph{Rigid motions and Wasserstein--Procrustes.}
\begin{defn}[Wasserstein--Procrustes distance]\label{def-wasserstein-procrustes}
For $\alpha,\beta\in\Pp_2(\RR^d)$, define
\[
\operatorname{WP}_2(\alpha,\beta)^2
\eqdef
\inf_{\substack{R\in\mathrm O(d),\,t\in\RR^d\\ \pi\in\Couplings(\alpha,\beta)}}
\int\norm{Rx+t-y}^2\d\pi(x,y).
\]
Restricting \(R\) to \(\mathrm{SO}(d)\) defines the orientation-preserving variant.
\end{defn}
This is an extrinsic counterpart of the Gromov--Wasserstein viewpoint developed in Section~\ref{sec-gromov-wasserstein}. Procrustes alignment searches only over ambient rigid motions, whereas GW compares the metric-measure structures themselves.

Given a current rigid motion $(R^{(\ell)},t^{(\ell)})$, first compute the OT coupling between the registered source cloud and the target cloud,
\begin{equation}\label{eq-wasserstein-procrustes-coupling-update}
\P^{(\ell)}
\in
\argmin_{\P\in\CouplingsD(\a,\b)}
\sum_{i,j}\P_{i,j}\norm{R^{(\ell)}x_i+t^{(\ell)}-y_j}^2.
\end{equation}
Then freeze this coupling and update the rigid motion by
\begin{equation}\label{eq-wasserstein-procrustes-rigid-update}
(R^{(\ell+1)},t^{(\ell+1)})
\in
\argmin_{R\in\mathrm O(d),\,t\in\RR^d}
\sum_{i,j}\P^{(\ell)}_{i,j}\norm{Rx_i+t-y_j}^2,
\end{equation}
\begin{prop}[Rigid update for a fixed coupling]\label{prop-wasserstein-procrustes-rigid-update}
Let $\alpha=\sum_i a_i\delta_{x_i}$ and $\beta=\sum_j b_j\delta_{y_j}$, and fix $\P\in\CouplingsD(\a,\b)$. Define
\(\bar x=\sum_i a_i x_i,\qquad \bar y=\sum_j b_j y_j,\qquad M_\P=\sum_{i,j} \P_{i,j}(y_j-\bar y)(x_i-\bar x)^\top.\)
If $M_\P=U\Sigma V^\top$ is a singular value decomposition, then the minimizers over the full orthogonal group of
\(\min_{R\in\mathrm O(d),\,t\in\RR^d} \sum_{i,j}\P_{i,j}\norm{Rx_i+t-y_j}^2\)
are given by $R^\star=UV^\top$ and $t^\star=\bar y-R^\star\bar x$, up to the usual non-uniqueness when $M_\P$ is rank deficient. If one imposes $R\in\mathrm{SO}(d)$, set
\(D=\diag(1,\ldots,1,\det(UV^\top)), \qquad R^\star=UDV^\top, \qquad t^\star=\bar y-R^\star\bar x.\)
\end{prop}
\begin{proof}
For fixed $R$, differentiating with respect to $t$ gives $t=\bar y-R\bar x$. Substituting this value centers the variables and leaves
\(\sum_{i,j}\P_{i,j}\norm{R(x_i-\bar x)-(y_j-\bar y)}^2.\)
The two quadratic terms independent of $R$ are fixed, so the problem is equivalent to maximizing
\(\sum_{i,j}\P_{i,j}\dotp{R(x_i-\bar x)}{y_j-\bar y} = \tr(R^\top M_\P).\)
Von Neumann's trace inequality gives $\tr(R^\top U\Sigma V^\top)\leq\tr(\Sigma)$ for $R\in\mathrm O(d)$, with equality for $R=UV^\top$. Under the constraint $R\in\mathrm{SO}(d)$, the same argument applies with the additional determinant constraint on $U^\top R V$, which gives the displayed diagonal correction.
\end{proof}
\subsection{Vector Quantiles and Linear Optimal Transport}
\label{sec-linear-ot}
\label{sec-vector-quantiles-linearized-transport}
\paragraph{Vector quantiles.}
Assume that $\gamma\in\Pp_2(\RR^d)$ is absolutely continuous. For a target law $\al\in\Pp_2(\RR^d)$, its vector quantile relative to $\gamma$ is the Brenier map
\(\T_\al=\nabla\phi_\al, \qquad (\T_\al)_\sharp\gamma=\al,\)
or equivalently the solution of
\(\min_{\T_\sharp\gamma=\al} \int \norm{x-\T(x)}^2\d\gamma(x).\)
\paragraph{Linearized Wasserstein coordinates.}
\begin{defn}[Linear optimal-transport embedding]\label{def-lot-embedding}
Fix an absolutely continuous reference measure $\gamma\in\Pp_2(\RR^d)$, and let $T_\alpha$ denote the Brenier map from $\gamma$ to $\alpha\in\Pp_2(\RR^d)$. Then
\begin{equation}\label{eq-lot-embedding}
\alpha \mapsto T_\alpha-\Id\in L^2(\gamma;\RR^d),
\qquad
\LOT_\gamma(\alpha,\beta)=\norm{T_\alpha-T_\beta}_{L^2(\gamma)}.
\end{equation}
\end{defn}
\begin{prop}[Closed convex range of linear OT]\label{prop-lot-closed-convex-range}
Fix an absolutely continuous reference $\gamma\in\Pp_2(\RR^d)$. The set
\(\enscond{T_\alpha\in L^2(\gamma;\RR^d)}{\alpha\in\Pp_2(\RR^d)}\)
is a closed convex cone. The map $\alpha\mapsto T_\alpha$ is a bijection onto this cone, with inverse $T\mapsto T_\sharp\gamma$. Equivalently, the range of the displacement embedding $\alpha\mapsto T_\alpha-\Id$ in~\eqref{eq-lot-embedding} is closed and convex.
\end{prop}
\begin{proof}
By Brenier's theorem, $T_\alpha=\nabla\phi_\alpha$ for a convex function $\phi_\alpha$. If $a,b\geq0$, then $aT_\alpha+bT_\beta$ is the gradient of the convex function $a\phi_\alpha+b\phi_\beta$. Setting $\eta=(aT_\alpha+bT_\beta)_\sharp\gamma$, cyclical monotonicity and uniqueness in Brenier's theorem show that $aT_\alpha+bT_\beta=T_\eta$ $\gamma$-almost everywhere. Thus the range is a convex cone, and the inverse formula is immediate.

For closedness, suppose that $T_{\alpha_n}\to T$ in $L^2(\gamma;\RR^d)$ and set $\alpha=T_\sharp\gamma$. The common-source coupling gives \(\Wass_2(\alpha_n,\alpha)^2 \leq \int\norm{T_{\alpha_n}-T}^2\d\gamma \longrightarrow0.\)
Moreover,
$\int\norm{x-T_{\alpha_n}(x)}^2\d\gamma(x)=\Wass_2(\gamma,\alpha_n)^2$.
Both sides converge, respectively, to
$\int\norm{x-T(x)}^2\d\gamma(x)$ and $\Wass_2(\gamma,\alpha)^2$.
Hence $(\Id,T)_\sharp\gamma$ is optimal between $\gamma$ and $\alpha$; uniqueness of the Brenier map gives $T=T_\alpha$ $\gamma$-almost everywhere. The range is therefore closed. Translating it by $-\Id$ preserves closedness and convexity, and~\eqref{eq-lot-embedding} identifies $\LOT_\gamma$ with the ambient Hilbert distance on that translated range.
\end{proof}
For instance, \(\LOT_\gamma(\gamma,\alpha) = \norm{T_\alpha-\Id}_{L^2(\gamma)} = \Wass_2(\gamma,\alpha).\)
For two arbitrary targets, the coupling $(T_\alpha,T_\beta)_\sharp\gamma$ is admissible but not generally optimal, so $\LOT_\gamma$ is a tangent-space approximation of the Wasserstein geometry.

Proposition~\ref{prop-lot-closed-convex-range} makes linear interpolation intrinsic to the LOT geometry: the unique constant-speed $\LOT_\gamma$-geodesic between $\alpha_0$ and $\alpha_1$ is
\(\alpha_t=\big((1-t)T_{\alpha_0}+tT_{\alpha_1}\big)_\sharp\gamma, \qquad t\in[0,1].\)
Likewise, for a family $(\alpha_s)_s$ with weights $\lambda_s\geq0$ satisfying $\sum_s\lambda_s=1$, the unique $\LOT_\gamma$-barycenter is obtained by averaging maps,
\(\bar T=\sum_s\lambda_s T_{\alpha_s}, \qquad \bar\alpha_{\LOT}=\bar T_\sharp\gamma.\)
This is exact in one dimension, where quantile functions linearize $\Wass_2$, and it is especially useful when many barycenters with changing weights must be evaluated quickly.

\begin{prop}[Quantitative stability of linear OT]\label{prop-linear-ot-stability}
Let $X,Y\subset\RR^d$ be compact convex sets, and let $\gamma\in\Pp(X)$ satisfy $\d\gamma=\rho_\gamma\d x$ with $\rho_\gamma\in L^\infty(X)$. Then there exists a constant $C=C(d,X,Y,\norm{\rho_\gamma}_{L^\infty})$ such that, for all $\alpha,\beta\in\Pp(Y)$,
\(\Wass_2(\alpha,\beta)\leq \LOT_\gamma(\alpha,\beta), \qquad \LOT_\gamma(\alpha,\beta)^4 \leq C\Wass_1(\alpha,\beta).\)
\end{prop}
\begin{proof}
The first inequality is immediate: $(T_\alpha,T_\beta)_\sharp\gamma$ is a feasible coupling between $\alpha$ and $\beta$. The second is the sharp stability estimate of M\'erigot. Its proof bounds a symmetrized divergence of the Kantorovich functional by $\Wass_1(\alpha,\beta)$ and controls from below the same divergence by the fourth power of $\norm{T_\alpha-T_\beta}_{L^2(\gamma)}$.
In one dimension, quantile functions make the embedding isometric to $\Wass_2$; since the targets lie in $Y$, one then has the sharper estimate $\LOT_\gamma(\alpha,\beta)^2\leq\operatorname{diam}(Y)\Wass_1(\alpha,\beta)$. In dimensions $d\geq2$, the uniform exponent $1/4$ is optimal over broad classes of compactly supported reference measures with bounded densities, although better exponents can hold near regular targets. It should not be read as a global Lipschitz estimate in $\Wass_2$.
\end{proof}
\paragraph{Principal components in linear OT coordinates.}
Given training measures \((\alpha_i)_{i=1}^N\), set
\(\xi_i \eqdef T_{\alpha_i}-\Id, \qquad \bar\xi \eqdef \frac1N\sum_{i=1}^N \xi_i.\)
The empirical covariance operator on \(L^2(\gamma;\RR^d)\) is the finite-rank operator
\(\mathcal C_{\LOT} h \eqdef \frac1N\sum_{i=1}^N \dotp{\xi_i-\bar\xi}{h}_{L^2(\gamma)} (\xi_i-\bar\xi),\)
If \(Gv_k=\lambda_k v_k\) with \(\lambda_k>0\) and \(\norm{v_k}=1\), then
\(e_k=\frac{1}{\sqrt{N\lambda_k}}\sum_{i=1}^N (v_k)_i(\xi_i-\bar\xi)\)
The score of \(\alpha_i\) along mode \(e_k\) is \(a_{i,k} \eqdef \dotp{\xi_i-\bar\xi}{e_k}_{L^2(\gamma)}.\)
A low-rank reconstruction, or a synthetic excursion with prescribed coefficients \(a=(a_1,\ldots,a_m)\), is then obtained by pushing the reference through
\(T_a(x) \eqdef x+\bar\xi(x)+\sum_{k=1}^m a_k e_k(x), \qquad \alpha_a \eqdef (T_a)_\sharp\gamma.\)
\subsection{Spectral and Robust Wasserstein Distances}
\label{sec-spectral-subspace-wasserstein}
\begin{defn}[Monotone spectral gauge]\label{def-monotone-spectral-gauge}
A monotone spectral gauge on positive semidefinite matrices is a convex, positively $1$-homogeneous map $\gamma:\mathbb S_+^d\to\RR_+$ such that $\gamma(M)=0$ only for $M=0$, $\gamma(QM\transp{Q})=\gamma(M)$ for every orthogonal matrix $Q$, and
\(0\preceq M\preceq N \quad\Longrightarrow\quad \gamma(M)\leq\gamma(N).\)
\end{defn}
\begin{defn}[Schatten gauge]\label{def-schatten-gauge}
For $1\leq q\leq+\infty$, define the Schatten gauge by
\[
\gamma_q(M)\eqdef\norm{M}_{S_q}
=
\begin{cases}
\pa{\sum_{i=1}^d\lambda_i(M)^q}^{1/q}, & 1\leq q<+\infty,\\
\lambda_{\max}(M), & q=+\infty,
\end{cases}.
\]
It is a monotone spectral gauge. The cases $q=1$, $q=2$ and $q=+\infty$ are respectively the trace, Frobenius and spectral gauges.
\end{defn}
\begin{defn}[Spectral Wasserstein distance]\label{def-spectral-wasserstein}
Let $\gamma$ be a monotone spectral gauge. For a coupling $\pi\in\Couplings(\alpha,\beta)$, define its displacement second-moment matrix
\(M_\pi\eqdef\int (x-y)(x-y)^\top\d\pi(x,y).\)
The spectral Wasserstein distance associated with $\gamma$ is
\eql{\label{eq-spectral-wasserstein}
\Wass_\gamma(\alpha,\beta)^2
\eqdef
\inf_{\pi\in\Couplings(\alpha,\beta)}\gamma(M_\pi).
}
\end{defn}
The special case $\gamma(M)=\tr(M)$ gives the usual quadratic Wasserstein distance $\Wass_2$. The spectral gauge $\gamma(M)=\lambda_{\max}(M)$ instead measures the largest mean squared projected displacement. For $A\succeq0$, define the quadratic projected transport cost
\eql{\label{eq-quadratic-projected-cost}
\Wass_{2,A}(\alpha,\beta)^2
\eqdef
\inf_{\pi\in\Couplings(\alpha,\beta)}
\int (x-y)^\top A(x-y)\d\pi(x,y)
=
\Wass_2((A^{1/2})_\sharp\alpha,(A^{1/2})_\sharp\beta)^2.
}
The polar set of the gauge is
\eql{\label{eq-spectral-polar-set}
\mathcal B_\gamma
\eqdef
\enscond{A\succeq0}{\tr(AM)\leq\gamma(M)\quad\text{for all }M\succeq0},
}
For the Schatten gauge $\gamma_q$, Schatten H\"older duality gives \(\mathcal B_{\gamma_q} = \enscond{A\succeq0}{\norm{A}_{S_{q^\ast}}\leq1}, \qquad \frac1q+\frac1{q^\ast}=1,\)
with the usual endpoint conventions. Thus the trace gauge has polar set $\{A:0\preceq A\preceq\Id\}$, the Frobenius gauge is self-polar on $\mathbb S_+^d$, and the spectral gauge has polar set $\{A\succeq0:\tr(A)\leq1\}$.

\begin{prop}[Robust representation and metric equivalence]\label{prop-spectral-wasserstein-robust}
Assume, for simplicity, that the measures are compactly supported and that $\gamma$ is closed and finite on the positive semidefinite cone. Then
\(\Wass_\gamma(\alpha,\beta)^2 = \sup_{A\in\mathcal B_\gamma} \Wass_{2,A}(\alpha,\beta)^2.\)
If there exist constants $0<a\leq b<+\infty$ such that $a\Id\in\mathcal B_\gamma$ and $\mathcal B_\gamma\subset\enscond{A}{0\preceq A\preceq b\Id}$, equivalently
$a\tr(M)\leq\gamma(M)\leq b\tr(M)$ for $M\succeq0$,
then
\(\sqrt a\,\Wass_2(\alpha,\beta) \leq \Wass_\gamma(\alpha,\beta) \leq \sqrt b\,\Wass_2(\alpha,\beta).\)
In particular, $\Wass_\gamma$ is a distance. When $\gamma$ is the restriction of a norm to the positive semidefinite cone, these bounds hold automatically in finite dimension, so $\Wass_\gamma$ is equivalent to $\Wass_2$ on measures with finite second moments.
\end{prop}
\begin{proof}
Using the polar representation of $\gamma$,
$\Wass_\gamma(\alpha,\beta)^2 = \inf_{\pi\in\Couplings(\alpha,\beta)} \sup_{A\in\mathcal B_\gamma}\tr(AM_\pi)$.
The coupling set is convex and compact for weak convergence under compact support. Since $\gamma$ is a finite gauge on a finite-dimensional cone and vanishes only at the origin, it is equivalent to the trace norm on the slice $\tr(M)=1$, so $\mathcal B_\gamma$ is convex and compact. The map $(\pi,A)\mapsto\tr(AM_\pi)$ is affine in each variable and continuous. Sion's minimax theorem gives
\[
\inf_\pi\sup_{A\in\mathcal B_\gamma}\tr(AM_\pi)
=
\sup_{A\in\mathcal B_\gamma}\inf_\pi\tr(AM_\pi)
=
\sup_{A\in\mathcal B_\gamma}\Wass_{2,A}(\alpha,\beta)^2.
\]
For fixed $A\succeq0$, $\Wass_{2,A}$ is the Wasserstein pseudodistance associated with the seminorm $x\mapsto\norm{A^{1/2}x}$. Since all terms are nonnegative, the robust identity also gives
$\Wass_\gamma(\alpha,\beta) = \sup_{A\in\mathcal B_\gamma}\Wass_{2,A}(\alpha,\beta)$.
A supremum of pseudodistances is symmetric and satisfies the triangle inequality.

If $a\Id\in\mathcal B_\gamma$ and $A\preceq b\Id$ for all $A\in\mathcal B_\gamma$, then
\(a\Wass_2(\alpha,\beta)^2 = \Wass_{2,a\Id}(\alpha,\beta)^2 \leq \Wass_\gamma(\alpha,\beta)^2 \leq b\Wass_2(\alpha,\beta)^2,\)
which proves definiteness and equivalence with $\Wass_2$. The equivalence between these operator bounds and $a\tr(M)\leq\gamma(M)\leq b\tr(M)$ follows directly from the polar formula. In finite dimension, any norm restricted to the positive semidefinite cone is equivalent to the trace norm on that cone. The finite-second-moment case follows by truncation when these norm-equivalence bounds hold.
\end{proof}
\begin{defn}[Projection- and subspace-robust Wasserstein]\label{def-projection-subspace-robust-wasserstein}
For $1\leq k\leq d$, let
$\mathcal P_k=\enscond{P}{P^2=P=\transp{P},\ \tr(P)=k}$.
The Paty--Cuturi projection-robust and subspace-robust Wasserstein distances are respectively the max--min and min--max quantities
\[
\PRW_{2,k}(\alpha,\beta)
\eqdef
\sup_{\transp{U}U=\Id_k}
\Wass_2((\transp{U})_\sharp\alpha,(\transp{U})_\sharp\beta)
=
\sup_{P\in\mathcal P_k}\Wass_{2,P}(\alpha,\beta),
\qquad
\SRW_{2,k}(\alpha,\beta)^2
\eqdef
\inf_{\pi\in\Couplings(\alpha,\beta)}
\sup_{P\in\mathcal P_k}\tr(PM_\pi).
\]
\end{defn}
\begin{prop}[Projection/subspace robust comparison]\label{prop-ky-fan-srw-comparison}
Let $\lambda_1(M)\geq\cdots\geq\lambda_d(M)\geq0$ and define the Ky Fan gauge $\gamma_k(M)=\sum_{\ell=1}^k\lambda_\ell(M)$. Then
\[
\max\left\{\PRW_{2,k}(\alpha,\beta),\sqrt{\frac{k}{d}}\,\Wass_2(\alpha,\beta)\right\}
\leq
\SRW_{2,k}(\alpha,\beta)
=\Wass_{\gamma_k}(\alpha,\beta)
\leq
\Wass_2(\alpha,\beta).
\]
For $k=d$, all displayed bounds are equalities and $\PRW_{2,d}=\SRW_{2,d}=\Wass_2$.
\end{prop}
\begin{proof}
Diagonalizing $M$ gives Ky Fan's variational formula
\(\gamma_k(M) = \max_{P\in\mathcal P_k}\tr(PM) = \max_{0\preceq A\preceq\Id,\ \tr(A)\leq k}\tr(AM).\)
Hence minimax weak duality yields
\[
\PRW_{2,k}(\alpha,\beta)^2
=
\sup_{P\in\mathcal P_k}\inf_{\pi\in\Couplings(\alpha,\beta)}\tr(PM_\pi)
\leq
\inf_{\pi\in\Couplings(\alpha,\beta)}\sup_{P\in\mathcal P_k}\tr(PM_\pi)
=
\Wass_{\gamma_k}(\alpha,\beta)^2.
\]
Finally, $(k/d)\tr(M)\leq\gamma_k(M)\leq\tr(M)$ for every $M\succeq0$. Taking infima over couplings and square roots proves the remaining bounds.
\end{proof}
For $k=1$, $\gamma_1(M)=\lambda_{\max}(M)$. The associated spectral gradient flows and their connection with the Muon algorithm are studied in Section~\ref{sec-normalized-spectral-wasserstein-dynamics}.

\subsection{Conditional Wasserstein Distances}
\label{sec-conditional-wasserstein-distances}
\begin{defn}[Conditional couplings and conditional OT]\label{def-conditional-ot}
Let $(S,\lambda)$ be a standard Borel probability space of conditions and let $(\Omega,\dist)$ be a Polish metric space. For $p\geq1$, denote by $\Pp_{p,\lambda}(S\times\Omega)$ the set of probability measures on $S\times\Omega$ whose first marginal is $\lambda$ and whose disintegration
\(\alpha(\d s,\d x)=\alpha_s(\d x)\lambda(\d s)\)
satisfies $\int_S\int_\Omega \dist(x,x_0)^p\,\d\alpha_s(x)\d\lambda(s)<+\infty$ for some, hence every, $x_0\in\Omega$.
For $\alpha,\beta\in\Pp_{p,\lambda}(S\times\Omega)$, a conditional coupling is a measure
\(\Pi(\d s,\d x,\d y) = \pi_s(\d x,\d y)\lambda(\d s)\)
such that $\pi_s\in\Couplings(\alpha_s,\beta_s)$ for $\lambda$-a.e.\ $s$. The set of such couplings is denoted $\Couplings_\lambda(\alpha,\beta)$.
Let $(s,x,y)\mapsto c_s(x,y)$ be jointly Borel, nonnegative, and lower semicontinuous in $(x,y)$ for every $s$. The conditional OT value is
\begin{equation}\label{eq-conditional-ot-general}
\begin{aligned}
\MK_c^\lambda(\alpha,\beta)
&\eqdef
\inf_{\Pi\in\Couplings_\lambda(\alpha,\beta)}
\int_S\int_{\Omega\times\Omega} c_s(x,y)\d\pi_s(x,y)\d\lambda(s) \\
&=
\int_S
\inf_{\pi_s\in\Couplings(\alpha_s,\beta_s)}
\int_{\Omega\times\Omega} c_s(x,y)\d\pi_s(x,y)
\d\lambda(s).
\end{aligned}
\end{equation}
\end{defn}
\begin{defn}[Conditional Wasserstein distance]\label{def-conditional-wasserstein-distance}
For the fiber-independent cost $c_s(x,y)=c_p(x,y)\eqdef\dist(x,y)^p$, define
\begin{equation}\label{eq-conditional-wasserstein-distance}
\Wass_{p,\lambda}(\alpha,\beta)
\eqdef
\left(\MK_{c_p}^\lambda(\alpha,\beta)\right)^{1/p}
=
\left(
\int_S \Wass_p^p(\alpha_s,\beta_s)\d\lambda(s)
\right)^{1/p}.
\end{equation}
\end{defn}
Thus $\MK_c^\lambda$ is the general conditional Kantorovich value, whereas $\Wass_{p,\lambda}$ is its metric specialization to the constant family $c_s=c_p=\dist^p$, after taking the $p$th root. Equivalently, $\Wass_{p,\lambda}(\alpha,\beta)^p = \MK_{c_p}^\lambda(\alpha,\beta)$; the two definitions use exactly the same conditional coupling problem.

\begin{prop}[Metric property]\label{prop-conditional-wasserstein-distance}
For $p\geq1$, $\Wass_{p,\lambda}$ is a distance on $\Pp_{p,\lambda}(S\times\Omega)$. Moreover, this metric space is Polish.
\end{prop}
\begin{proof}
Non-negativity and symmetry follow from the corresponding properties of $\Wass_p$ on each fiber. If $\Wass_{p,\lambda}(\alpha,\beta)=0$, then $\Wass_p(\alpha_s,\beta_s)=0$ for $\lambda$-a.e.\ $s$, hence $\alpha_s=\beta_s$ for $\lambda$-a.e.\ $s$ and the disintegrated measures $\alpha$ and $\beta$ coincide. For the triangle inequality, let $\gamma\in\Pp_{p,\lambda}(S\times\Omega)$. Since $\Wass_p(\alpha_s,\gamma_s)\leq\Wass_p(\alpha_s,\beta_s)+\Wass_p(\beta_s,\gamma_s)$ for a.e.\ $s$, Minkowski's inequality in $L^p(S,\lambda)$ gives
\(\Wass_{p,\lambda}(\alpha,\gamma) \leq \Wass_{p,\lambda}(\alpha,\beta)+\Wass_{p,\lambda}(\beta,\gamma).\)
To see completeness and separability, identify a measure with the $\lambda$-a.e.\ equivalence class of its measurable disintegration map $s\mapsto\alpha_s\in\Pp_p(\Omega)$. Formula~\eqref{eq-conditional-wasserstein-distance} is the corresponding $L^p(S,\lambda;\Pp_p(\Omega))$ metric. Since $(\Pp_p(\Omega),\Wass_p)$ is Polish, the standard $L^p$ space of measurable metric-valued maps is Polish as well.
\end{proof}
\begin{prop}[Conditional geodesics]\label{prop-conditional-wasserstein-geodesics}
Assume that $(\Omega,\dist)$ is a geodesic Polish space and let $\alpha_0,\alpha_1\in\Pp_{p,\lambda}(S\times\Omega)$. For $\lambda$-a.e.\ $s$, choose a constant-speed $\Wass_p$ geodesic $t\mapsto\alpha_{t,s}$ between $\alpha_{0,s}$ and $\alpha_{1,s}$, measurably in $s$, and set
$\alpha_t(\d s,\d x)=\alpha_{t,s}(\d x)\lambda(\d s)$.
Then $t\mapsto\alpha_t$ is a constant-speed geodesic for $\Wass_{p,\lambda}$.
\end{prop}
\begin{proof}
For $0\leq r\leq t\leq1$, the fiberwise geodesic property gives \(\Wass_p(\alpha_{r,s},\alpha_{t,s}) = (t-r)\Wass_p(\alpha_{0,s},\alpha_{1,s}) \quad\text{for }\lambda\text{-a.e.\ }s.\)
Integrating the $p$th power over $S$ gives
$\Wass_{p,\lambda}(\alpha_r,\alpha_t) = (t-r)\Wass_{p,\lambda}(\alpha_0,\alpha_1)$.
Hence the conditional curve has constant speed and realizes the distance between its endpoints. In Euclidean space one may take, for each $s$, an optimal plan $\pi_s\in\Couplings(\alpha_{0,s},\alpha_{1,s})$ and set $\alpha_{t,s}=((1-t)\mathrm p_1+t\mathrm p_2)_\sharp\pi_s$, where $\mathrm p_1,\mathrm p_2$ are the two coordinate projections. On a general geodesic space, the same construction uses a measurable family of optimal dynamical plans. Such a family can be selected measurably in the Polish setting; when geodesics are non-unique, different selections can produce different conditional geodesics.
\end{proof}

\section{Generalized OT Problems}
\label{sec-generalized-ot-problems}
\subsection{OT Barycenters}
\label{sec-barycenters}
\paragraph{Fr\'echet means.}
The natural formulation is a Fr\'echet-mean problem on the space of probability measures: the unknown is the barycenter measure itself, and its support is not prescribed. It uses the continuous Kantorovich value $\MK_c$ defined in~\eqref{eq-mk-generic}.

\begin{defn}[Optimal-transport barycenter]\label{def-ot-barycenter}
Given input measures $(\be_s)_{s=1}^S$ on a space $\Xx$ and weights $\la\in\simplex_S$, discard any zero-weight inputs. An \emph{optimal-transport barycenter} of this weighted family is any solution of
\eql{\label{eq-barycenter-generic}
\umin{\al \in \Mm_+^1(\Xx)} \sum_{s=1}^S \la_s \MK_{\c}(\al,\be_s).
}
\end{defn}
\paragraph{Fixed-support restriction.}
Assume the inputs are discrete, \(\be_s=\sum_{j=1}^{n_s} b_{s,j}\de_{x_{s,j}}, \qquad \b_s=(b_{s,j})_{j=1}^{n_s}\in\simplex_{n_s}.\)
For each input $s$, the cost $\MK_\c(\al,\be_s)$ then becomes a finite Kantorovich problem with cost matrix
\((\C_s)_{i,j}=c(y_i,x_{s,j}) \in \RR^{n\times n_s}.\)
Thus its value is $\MKD_{\C_s}(\a,\b_s)$, in the notation of the discrete Kantorovich problem~\eqref{eq-kanto-discr}.

\begin{defn}[Fixed-support discrete OT barycenter]\label{def-fixed-support-discrete-barycenter}
With the candidate sites and cost matrices above, a \emph{fixed-support discrete OT barycenter} is a measure $\al^\star=\sum_i a_i^\star\de_{y_i}$ whose weight vector $\a^\star=(a_i^\star)_{i=1}^n$ solves
\eql{\label{eq-wass-discr}
\umin{\a \in \simplex_n} \sum_{s=1}^S \la_s \MKD_{\C_s}(\a,\b_s),
}
\end{defn}
For quadratic costs on $\RR^d$, this gives
\[
\umin{\substack{\a\in\simplex_n\\Y=(y_i)_{i=1}^n\in(\RR^d)^n}}
\sum_{s=1}^S\la_s\MKD_{\C_s(Y)}(\a,\b_s),
\qquad
[\C_s(Y)]_{i,j}=\norm{y_i-x_{s,j}}^2.
\]
For fixed $Y$, the objective is the convex problem~\eqref{eq-wass-discr} in $\a$; its dependence on $Y$ is generally nonconvex, so the joint problem is nonconvex. Thus choosing the support is itself a significant computational part of the barycenter problem, rather than a harmless preprocessing step.

\begin{prop}[Two-measure barycenters are Wasserstein geodesics]\label{prop-two-measure-barycenter-geodesic}
Let $\be_0,\be_1\in\Pp_2(\RR^d)$ and $t\in[0,1]$. For $S=2$, $c(x,y)=\norm{x-y}^2$ and weights $(1-t,t)$, the minimizers of the barycenter problem~\eqref{eq-barycenter-generic} are exactly the time-$t$ points of constant-speed $\Wass_2$ geodesics from $\be_0$ to $\be_1$. The minimum value is
\(t(1-t)\Wass_2^2(\be_0,\be_1).\)
In particular, if $\pi$ is any optimal coupling between $\be_0$ and $\be_1$, then
\(\al_t \eqdef \bigl((x,y)\mapsto(1-t)x+ty\bigr)_\sharp\pi\)
is a barycenter. If $\be_0$ has a density, the barycenter is unique and
\(\al_t=((1-t)\Id+tT)_\sharp\be_0,\)
where $T$ is the Brenier map from $\be_0$ to $\be_1$.
\end{prop}
\begin{proof}
The endpoint cases $t\in\{0,1\}$ are immediate, so assume $0<t<1$. For any candidate $\al\in\Pp_2(\RR^d)$, set
\(A=\Wass_2(\be_0,\al), \qquad B=\Wass_2(\al,\be_1), \qquad D=\Wass_2(\be_0,\be_1).\)
The identity
\((1-t)A^2+tB^2 = t(1-t)(A+B)^2+\bigl((1-t)A-tB\bigr)^2\)
and the triangle inequality $A+B\geq D$ give
\((1-t)\Wass_2^2(\al,\be_0) +t\Wass_2^2(\al,\be_1) \geq t(1-t)D^2.\)
For the interpolation $\al_t$ induced by an optimal coupling, Proposition~\ref{prop-plan-interpolation-w2-geodesic} gives $A=tD$ and $B=(1-t)D$, so the lower bound is attained.

Conversely, every minimizer must attain equality in both inequalities above. Hence $A+B=D$ and $(1-t)A=tB$, which imply $A=tD$ and $B=(1-t)D$. Parameterize a constant-speed geodesic from $\be_0$ to $\al$ on $[0,t]$ and one from $\al$ to $\be_1$ on $[t,1]$; both pieces then have speed $D$. Their concatenation is globally geodesic: any shorter path between points on opposite pieces, combined with the remaining subsegments, would contradict $D=\Wass_2(\be_0,\be_1)$. Thus the concatenated geodesic passes through $\al$ at time $t$. Finally, if $\be_0$ has a density, Brenier's theorem (Theorem~\ref{thm-brenier}) gives the unique optimal coupling $(\Id,T)_\sharp\be_0$, and the displayed McCann interpolation follows from Definition~\ref{def-monge-mccann-interpolation}.
\end{proof}
\begin{prop}[Mean and support of quadratic Wasserstein barycenters]\label{prop-w2-barycenter-mean-support}
Let $\be_1,\ldots,\be_S\in\Pp_2(\RR^d)$, let $\la\in\simplex_S$, and let $\al^\star$ be a barycenter for the squared Euclidean cost. Define
\[
K
\eqdef
\overline{\operatorname{conv}}
\left(
\bigcup_{\{s:\,\la_s>0\}}\supp(\be_s)
\right).
\]
Then
\(\supp(\al^\star)\subset K, \qquad \int_{\RR^d}x\,\d\al^\star(x) = \sum_{s=1}^S\la_s\int_{\RR^d}x\,\d\be_s(x).\)
\end{prop}
\begin{proof}
Write $m_\al=\int x\,\d\al(x)$ and let $\al^0$ be its centered translate, as in Proposition~\ref{prop-wasserstein-mean-decomposition}. That proposition gives, for every candidate $\al$,
\(\sum_{s=1}^S\la_s\Wass_2^2(\al,\be_s) = \sum_{s=1}^S\la_s\norm{m_\al-m_{\be_s}}^2 +\sum_{s=1}^S\la_s\Wass_2^2(\al^0,\be_s^0).\)
For a fixed centered law $\al^0$, the second sum is independent of $m_\al$, whereas the first is uniquely minimized at
$m_\al=\sum_s\la_s m_{\be_s}$. Translating $\al^\star$ without changing its centered law therefore proves the mean identity.

For the support statement, let $P_K$ be the Euclidean metric projection onto the closed convex set $K$. For every $y\in K$, the projection inequality gives \(\norm{x-y}^2 \geq \norm{P_K(x)-y}^2+\norm{x-P_K(x)}^2.\)
For each $s$ with $\la_s>0$, let $\pi_s$ be an optimal coupling between $\al^\star$ and $\be_s$. Since $\be_s(K)=1$, pushing $\pi_s$ forward by $(P_K,\Id)$ and integrating the preceding inequality yields
\(\Wass_2^2\bigl((P_K)_\sharp\al^\star,\be_s\bigr) \leq \Wass_2^2(\al^\star,\be_s) - \int_{\RR^d}\operatorname{dist}(x,K)^2\,\d\al^\star(x).\)
After multiplication by $\la_s$ and summation, barycenter optimality forces the last integral to vanish. Hence $\al^\star(K)=1$, and the closedness of $K$ gives $\supp(\al^\star)\subset K$.
\end{proof}
If all input supports are compact, their finite union is compact and so is its convex hull; in that case the closure in the definition of $K$ is unnecessary.

\paragraph{One-dimensional case.}
\begin{prop}[Quantile barycenters on the line]\label{prop-quantile-barycenters}
Let $\be_1,\ldots,\be_S\in\Pp_2(\RR)$. For $\Xx=\RR$ and $c(x,y)=|x-y|^2$, the quantile function of their Wasserstein barycenter is the weighted average of the input quantile functions:
\(\cumul{\al^\star}^{-1}(r) = \sum_{s=1}^S\la_s\cumul{\be_s}^{-1}(r), \qquad r\in(0,1).\)
\end{prop}
\begin{proof}
The one-dimensional formula~\eqref{eq-wass-cumul} gives \(\sum_s\la_s\Wass_2^2(\al,\be_s) = \int_0^1 \sum_s\la_s \abs{\cumul{\al}^{-1}(r)-\cumul{\be_s}^{-1}(r)}^2 \d r.\)
The minimization decouples pointwise in $r$. For each fixed $r$, the minimizer of $z\mapsto\sum_s\la_s|z-\cumul{\be_s}^{-1}(r)|^2$ is the weighted average $\sum_s\la_s\cumul{\be_s}^{-1}(r)$. This function is nondecreasing because it is a positive weighted sum of nondecreasing quantile functions. It also belongs to $L^2(0,1)$ because every input quantile does, so it is the quantile function of a measure in $\Pp_2(\RR)$.
\end{proof}
\paragraph{Gaussian case.}
\begin{prop}[Nondegenerate Gaussian inputs remain Gaussian]\label{prop-gaussian-barycenter}
Let the positive-weight inputs be $\be_s=\Gaussian(\mean_s,\cov_s)$, with $\cov_s\succeq0$, and assume that at least one $\cov_s$ is positive definite. The quadratic Wasserstein barycenter is unique and has the form $\al^\star=\Gaussian(\mean,\cov)$, where
\(\mean=\sum_s\lambda_s\mean_s.\)
Its positive-definite covariance $\cov$ is the unique minimizer of the Bures objective
\eql{\label{eq-gaussian-barycenter-energy}
\mathcal E(X)
\eqdef
\sum_s \la_s \Bb(X,\cov_s)^2.
}
For $X\in\mathbb S_{++}^d$, define
\begin{align}
R_s(X)
&\eqdef
\pa{X^{1/2}\cov_sX^{1/2}}^{1/2},
&
T_s(X)
&\eqdef
X^{-1/2}R_s(X)X^{-1/2},
\label{eq-gaussian-barycenter-maps}
\\
\overline T(X)
&\eqdef
\sum_s\la_sT_s(X),
&
\Psi_{\lambda}(X)
&\eqdef
\sum_s\la_sR_s(X).
\label{eq-gaussian-barycenter-average-map}
\end{align}
Here $T_s(X)$ is the positive-semidefinite linear optimal map from $\Gaussian(0,X)$ to $\Gaussian(0,\cov_s)$. The covariance $\cov$ is equivalently characterized by either of the stationarity equations
\eql{\label{eq-gaussian-barycenter-fixed-point}
\overline T(\cov)=\Id
\qquad\Longleftrightarrow\qquad
\Psi_{\lambda}(\cov)=\cov.
}
In dimension one, if $\sigma_s=\sqrt{\cov_s}$ denotes the input standard deviation, then the barycenter standard deviation is $\sigma=\sum_s\lambda_s\sigma_s$.
\end{prop}
\begin{proof}
Let $\GaussProj(\al)$ denote the Gaussian measure with the same mean and covariance as $\al$. For every competitor $\al\in\Pp_2(\RR^d)$ and every Gaussian input $\be_s$, the Gelbrich contraction of Theorem~\ref{thm-gelbrich-projection} gives
\(\Wass_2^2\bigl(\GaussProj(\al),\be_s\bigr) = \Wass_2^2\bigl(\GaussProj(\al),\GaussProj(\be_s)\bigr) \leq \Wass_2^2(\al,\be_s).\)
Summing with weights $\lambda_s$ shows that moment-matched Gaussian projection cannot increase the barycenter objective. Since a barycenter exists, projecting any minimizer therefore produces a Gaussian barycenter. The input with positive-definite covariance is absolutely continuous; the uniqueness criterion stated after Definition~\ref{def-ot-barycenter} then implies that the barycenter itself is this Gaussian measure.

For a Gaussian candidate, formula~\eqref{eq-dist-gauss} separates the objective as
\[
\sum_s\lambda_s\Wass_2^2\bigl(\Gaussian(\mean,\cov),\Gaussian(\mean_s,\cov_s)\bigr)
=
\sum_s\lambda_s\norm{\mean-\mean_s}^2
+
\sum_s\lambda_s\Bb(\cov,\cov_s)^2.
\]
The first term is uniquely minimized at $\mean=\sum_s\lambda_s\mean_s$. The second is the Bures barycenter problem. Uniqueness of the Wasserstein barycenter makes its covariance minimizer unique, and the presence of a positive-definite input makes this minimizer positive definite.

First suppose that $\cov_s\succ0$. Cyclicity of the trace gives \(\tr R_s(X) = \tr\pa{\cov_s^{1/2}X\cov_s^{1/2}}^{1/2}.\)
For a symmetric perturbation $Z$, differentiation of the trace of the matrix square root yields
\[
D\!\left[\tr R_s(X)\right][Z]
=
\frac12\tr\bigl(T_s(X)Z\bigr).
\]
\eqllead{Consequently,}{eq-gaussian-barycenter-first-variation}{
D\mathcal E(X)[Z]
=
\tr\bigl((\Id-\overline T(X))Z\bigr).
}
The same identity for singular $\cov_s$ follows by replacing $\cov_s$ with $\cov_s+\delta\Id$ and letting $\delta\downarrow0$. Since the minimizer lies in $\mathbb S_{++}^d$, first-order optimality is $\overline T(\cov)=\Id$. Finally, definitions~\eqref{eq-gaussian-barycenter-maps}--\eqref{eq-gaussian-barycenter-average-map} give
\(\Psi_\lambda(X) = X^{1/2}\overline T(X)X^{1/2},\)
so $\overline T(\cov)=\Id$ is equivalent to $\Psi_\lambda(\cov)=\cov$. Conversely, a positive-definite solution of this equation is stationary for the convex objective~\eqref{eq-gaussian-barycenter-energy}, hence is its unique minimizer. In dimension one, $\Bb(\sigma^2,\sigma_s^2)^2=(\sigma-\sigma_s)^2$, whose minimizer is $\sigma=\sum_s\lambda_s\sigma_s$.
\end{proof}
Once the mean has been set to $\mean=\sum_s\lambda_s\mean_s$, the optimal affine map from $\Gaussian(\mean,X)$ to $\Gaussian(\mean_s,\cov_s)$ is \(x\longmapsto \mean_s+T_s(X)(x-\mean).\)
Pushing the current Gaussian forward by this averaged map gives the \emph{transport fixed-point update} of \'Alvarez-Esteban et al., whose covariance is
\eql{\label{eq-gaussian-barycenter-transport-update}
\mathcal K_1(X)
\eqdef
\overline T(X)X\overline T(X)
=
X^{-1/2}\Psi_\lambda(X)^2X^{-1/2}.
}
More generally, damping the displacement by $\eta\in(0,1]$ replaces the averaged map by \(B_\eta(X) \eqdef (1-\eta)\Id+\eta\overline T(X)\)
\eqllead{and gives}{eq-gaussian-barycenter-damped-update}{
\mathcal K_\eta(X)
\eqdef
B_\eta(X)X B_\eta(X).
}
\begin{prop}[Convergence of the Gaussian transport fixed-point iteration]\label{prop-gaussian-barycenter-fixed-point-convergence}
Assume that $\lambda_s>0$, $\cov_s\succeq0$, and that at least one $\cov_s$ is positive definite. Fix $\eta\in(0,1]$ and $X^{(0)}\succ0$, and iterate
\(X^{(\ell+1)}=\mathcal K_\eta(X^{(\ell)}).\)
Then every $X^{(\ell)}$ is positive definite and $X^{(\ell)}\to\cov$, where $\cov$ is the covariance of the Gaussian barycenter in Proposition~\ref{prop-gaussian-barycenter}. More precisely, with
\(\mathcal R(X) \eqdef \tr\bigl((\Id-\overline T(X))X(\Id-\overline T(X))\bigr),\)
\eqllead{one has}{eq-gaussian-barycenter-descent}{
\mathcal E(X^{(\ell+1)})
\leq
\mathcal E(X^{(\ell)})
-
\eta(2-\eta)\mathcal R(X^{(\ell)}),
}
and therefore $\sum_{\ell\geq0}\mathcal R(X^{(\ell)})<+\infty$. Equivalently, \(\Bb\bigl(X^{(\ell)},X^{(\ell+1)}\bigr)^2 = \eta^2\mathcal R(X^{(\ell)}).\)
\end{prop}
\begin{proof}
Fix $X\succ0$, abbreviate $T_s=T_s(X)$, $\overline T=\overline T(X)$ and $B=B_\eta(X)$, and let $Z\sim\Gaussian(0,X)$. Since $T_sZ\sim\Gaussian(0,\cov_s)$, the pair $(BZ,T_sZ)$ is an admissible coupling between $\Gaussian(0,\mathcal K_\eta(X))$ and $\Gaussian(0,\cov_s)$. For every $z\in\RR^d$, the Euclidean variance decomposition around $\overline Tz=\sum_s\lambda_sT_sz$ gives
\begin{align*}
\sum_s\lambda_s\norm{Bz-T_sz}^2
&=
\sum_s\lambda_s\norm{\overline Tz-T_sz}^2
+(1-\eta)^2\norm{z-\overline Tz}^2,
\\
\sum_s\lambda_s\norm{z-T_sz}^2
&=
\sum_s\lambda_s\norm{\overline Tz-T_sz}^2
+\norm{z-\overline Tz}^2.
\end{align*}
Integrating and bounding each Wasserstein cost by the admissible coupling $(BZ,T_sZ)$ proves~\eqref{eq-gaussian-barycenter-descent}. Moreover, $B\succ0$: this is immediate when $\eta<1$, and for $\eta=1$ it follows because the weighted sum $\overline T$ contains at least one positive-definite term. Hence $z\mapsto Bz$ is the optimal map from $\Gaussian(0,X)$ to $\Gaussian(0,\mathcal K_\eta(X))$, which proves the displayed identity for the Bures step length.

Set
\(\delta(X)\eqdef\det(X)^{1/(2d)}, \qquad c_0\eqdef\sum_s\lambda_s\det(\cov_s)^{1/(2d)}>0.\)
The Minkowski determinant inequality, applied to the positive-semidefinite matrices $\Id$ and $T_s(X)$ with weights $1-\eta$ and $\eta\lambda_s$, gives
\begin{align*}
\delta\bigl(\mathcal K_\eta(X)\bigr)
&=
\delta(X)\det(B_\eta(X))^{1/d}
\\
&\geq
(1-\eta)\delta(X)
+
\eta\sum_s\lambda_s
\delta(X)\det(T_s(X))^{1/d}
\\
&=
(1-\eta)\delta(X)+\eta c_0,
\end{align*}
where the last identity uses $T_s(X)XT_s(X)=\cov_s$. It follows that
$\delta(X^{(\ell)})\geq\min\{\delta(X^{(0)}),c_0\}>0$ for every $\ell$.

The decreasing energy also bounds the traces. Indeed, the Schatten Cauchy--Schwarz inequality in the Bures formula gives \(\Bb(X,\cov_s)^2 \geq \pa{\sqrt{\tr X}-\sqrt{\tr\cov_s}}^2,\)
so~\eqref{eq-gaussian-barycenter-descent} bounds $\tr X^{(\ell)}$ uniformly. A uniform trace bound together with the determinant lower bound places all iterates in a compact subset of $\mathbb S_{++}^d$.

Summing~\eqref{eq-gaussian-barycenter-descent} gives $\mathcal R(X^{(\ell)})\to0$. Every subsequence therefore has a further subsequence converging to some $X_\infty\succ0$. Continuity of the matrix square root and inverse on the preceding compact set yields $\mathcal R(X_\infty)=0$, hence $\overline T(X_\infty)=\Id$. Proposition~\ref{prop-gaussian-barycenter} identifies $X_\infty$ with the unique covariance $\cov$. Since every cluster point is $\cov$, the full sequence converges to $\cov$. The undamped case $\eta=1$ is the Gaussian specialization of Theorem~4.2 of \'Alvarez-Esteban et al.; the argument above also proves the damped extension.
\end{proof}
\paragraph{Sinkhorn for barycenters.}
One can then use the entropy-only convention of~\eqref{eq-regularized-discr} and approximate~\eqref{eq-wass-discr} by
\eql{\label{eq-entropic-bary}
\umin{\a\in\simplex_n}
\sum_{s=1}^S\la_s\MKD_{\C_s}^\epsilon(\a,\b_s)
}
Problem~\eqref{eq-entropic-bary} has the same minimizers as the weighted KL projection problem
\eql{\label{eq-bary-entropy-couplings}
\umin{(\P_s)_s}
\enscond{
\epsilon\sum_s\la_s\KLD(\P_s|\K_s)
}{
\foralls s,\ \transp{\P_s}\ones_n=\b_s,
\quad
\P_1\ones_{n_1}=\ldots=\P_S\ones_{n_S}
},
}
where $\K_s\eqdef e^{-\C_s/\epsilon}$. The barycenter $\a$ is implicitly encoded in the common row marginal
\(\a=\P_1\ones_{n_1}=\ldots=\P_S\ones_{n_S}.\)
The optimal couplings solving~\eqref{eq-bary-entropy-couplings} then have scaling form
\eql{\label{eq-bary-opt}
\P_s=\diag(\uD_s)\K_s\diag(\vD_s),
}
and the generalized Sinkhorn iterations are
\begin{align}
\foralls s\in\range{1,S},\qquad
\itt{\vD}_s
&\eqdef
\frac{\b_s}{\transp{\K_s}\it{\uD}_s},
\label{eq-sinkhorn-bary}
\\
\foralls s\in\range{1,S},\qquad
\itt{\uD}_s
&\eqdef
\frac{\itt{\a}}{\K_s\itt{\vD}_s},
\label{eq-sinkhorn-bary-2}
\\
\qwhereq
\itt{\a}
&\eqdef
\prod_s(\K_s\itt{\vD}_s)^{\la_s}.
\label{eq-sinkhorn-bary-3}
\end{align}
Write
\(\uD_s=e^{\fD_s/\epsilon}\) and \(\vD_s=e^{\gD_s/\epsilon}\), with all operations understood componentwise. Proposition~\ref{prop-dual-entropic-barycenters} below shows that the potentials maximize the concave dual~\eqref{eq-dual-bary-entropy}. With $(\fD_s)_s$ fixed, this objective separates over $s$, and its exact maximizer in each $\gD_s$ is
\(\gD_s^+ = \epsilon\log\pa{\frac{\b_s}{\transp{\K_s}e^{\fD_s/\epsilon}}},\)
Conversely, with $(\gD_s^+)_s$ fixed, exact maximization over the coupled block $(\fD_s)_s$ under $\sum_s\lambda_s\fD_s=0$ gives
\(\q_s \eqdef \K_s e^{\gD_s^+/\epsilon}, \qquad \a^+ = \prod_s\q_s^{\lambda_s}, \qquad \fD_s^+ = \epsilon\log\pa{\frac{\a^+}{\q_s}}.\)
Thus a complete generalized Sinkhorn cycle is exact two-block coordinate ascent on the dual, or equivalently alternating minimization of its negative. Indeed, the constraint follows from $\log\a^+=\sum_s\lambda_s\log\q_s$, while the first-order conditions require $e^{\fD_s^+/\epsilon}\q_s=\a^+$ for every $s$.

\begin{prop}[Dual of entropic barycenters]\label{prop-dual-entropic-barycenters}
Under the preceding positivity assumptions, the optimal scalings in~\eqref{eq-bary-opt} can be written as $(\uD_s,\vD_s)=(e^{\fD_s/\epsilon},e^{\gD_s/\epsilon})$, where $(\fD_s,\gD_s)_s$ solve the dual problem
\eql{\label{eq-dual-bary-entropy}
\umax{(\fD_s,\gD_s)_s}
\enscond{
\sum_s\la_s
\left(
\dotp{\gD_s}{\b_s}
-\epsilon\dotp{\K_s e^{\gD_s/\epsilon}}{e^{\fD_s/\epsilon}}
+\epsilon\sum_{i,j}(\K_s)_{i,j}
\right)
}{
\sum_s\la_s\fD_s=0
}.
}
\end{prop}
\begin{proof}
Introduce Lagrange multipliers in~\eqref{eq-bary-entropy-couplings}:
\begin{align*}
\umin{(\P_s)_s,\a}
\umax{(\fD_s,\gD_s)_s}
\sum_s\la_s\Big(
\epsilon\KLD(\P_s|\K_s)
+\dotp{\a-\P_s\ones_{n_s}}{\fD_s}
+\dotp{\b_s-\transp{\P_s}\ones_n}{\gD_s}
\Big).
\end{align*}
The explicit constraint $\a\in\simplex_n$ may be dropped here: nonnegativity of the couplings, together with $\P_s\ones_{n_s}=\a$ and $\transp{\P_s}\ones_n=\b_s$, already forces $\a\in\simplex_n$. We may therefore minimize the Lagrangian over $\a\in\RR^n$. Strong duality holds, so one can exchange the minimum and maximum. Finiteness of the minimum with respect to $\a$ gives the vector constraint $\sum_s\la_s\fD_s=0$, while minimization with respect to $\P_s$ gives the Legendre transform of $\KLD(\cdot|\K_s)$:
\[
\umax{(\fD_s,\gD_s)_s}
\sum_s\la_s
\left[
\dotp{\gD_s}{\b_s}
-\epsilon
\KLD^*\left(\frac{\fD_s\oplus\gD_s}{\epsilon}\middle|\K_s\right)
\right],
\qquad
\sum_s\la_s\fD_s=0.
\]
\eqllead{The separable conjugate is}{eq-legendre-kl-bary}{
\KLD^*(\VectMode{U}|\K)
=
\sum_{i,j}\K_{i,j}\big(e^{\VectMode{U}_{i,j}}-1\big),
}
because for $k>0$, \(\sup_{r\geq0} ur-\big(r\log(r/k)-r+k\big) = k(e^u-1),\)
and the case $k=0$ follows by lower semicontinuity. Substituting this conjugate gives~\eqref{eq-dual-bary-entropy}, including its additive constant. The coordinate maximization in $\gD_s$ gives~\eqref{eq-sinkhorn-bary}; the block maximization in all $(\fD_s)_s$ under $\sum_s\lambda_s\fD_s=0$ gives the weighted geometric mean~\eqref{eq-sinkhorn-bary-3} and then~\eqref{eq-sinkhorn-bary-2}.
\end{proof}
\paragraph{Doubly regularized entropic barycenters.}
\begin{defn}[Doubly regularized entropic barycenter]\label{def-doubly-regularized-entropic-barycenter}
For $\epsilon,\tau\geq0$, define
\eql{\label{eq-doubly-regularized-entropic-barycenter}
\mathcal E_{\epsilon,\tau}(\al)
\eqdef
\sum_{s=1}^S\la_s\MK_\c^\epsilon(\al,\be_s)
+\tau\KL(\al|\ga),
\qquad
\al_{\epsilon,\tau}^\star
\in
\argmin_{\al\in\Mm_+^1(\Xx)}\mathcal E_{\epsilon,\tau}(\al),
}
with the convention $\MK_\c^0=\MK_\c$.
\end{defn}
On a compact space with continuous cost, $\mathcal E_{\epsilon,\tau}$ is proper, lower semicontinuous and convex, so it admits a minimizer. If $\tau>0$, strict convexity of $\KL(\cdot|\ga)$ makes this minimizer unique; finite outer entropy also imposes $\al_{\epsilon,\tau}^\star\ll\ga$.

The KL reformulation~\eqref{eq-regularized-discr-rescaled} gives
\begin{equation}\label{eq-doubly-regularized-barycenter-cancellation}
\mathcal E_{\epsilon,\tau}(\a)=\sum_{s=1}^S\la_s\MKD_{\C_s}^\epsilon(\a,\b_s)+(\epsilon-\tau)\HD(\a)-\tau\sum_{i=1}^n a_i\log g_i+\epsilon\sum_{s=1}^S\la_s\HD(\b_s).
\end{equation}
Indeed, $\KLD(\a|\gD)=-\HD(\a)-\sum_i a_i\log g_i$, while each normalized entropic OT value contributes $\epsilon\HD(\a)+\epsilon\HD(\b_s)$. If $\ga$ is uniform on the grid and $\tau=\epsilon$, every term outside the first sum in~\eqref{eq-doubly-regularized-barycenter-cancellation} is constant.

\paragraph{Wasserstein-over-Wasserstein and barycenters.}
Assume, for instance, that $c$ is lower semicontinuous and that there exists at least one $\beta_0\in\Pp(\Omega)$ with \(\int_{\Pp(\Omega)}\MK_c(\beta_0,\alpha)\d\mathfrak A(\alpha)<+\infty,\)
The barycenter correspondence is then
\(\mathcal B_c(\mathfrak A) \eqdef \operatorname*{Argmin}_{\beta\in\Pp(\Omega)} \int_{\Pp(\Omega)}\MK_c(\beta,\alpha)\d\mathfrak A(\alpha).\)
When this set is a singleton, we denote its element by \(\widetilde\alpha_{\mathfrak A} \eqdef \uargmin{\beta\in\Pp(\Omega)} \int_{\Pp(\Omega)}\MK_c(\beta,\alpha)\d\mathfrak A(\alpha),\)
\begin{prop}[Law of large numbers for barycenters over measures]\label{prop-wow-barycenter-lln}
Let $\Omega$ be a compact metric space, let $c:\Omega\times\Omega\to\RR$ be continuous, and let $\mathfrak A\in\Pp(\Pp(\Omega))$. Let $\alpha_1,\alpha_2,\ldots$ be independent random probability measures with common law $\mathfrak A$, and set
\(\widehat{\mathfrak A}_p \eqdef \frac1p\sum_{i=1}^p\delta_{\alpha_i} \in\Pp(\Pp(\Omega)).\)
Then, almost surely,
\begin{equation}\label{eq-wow-empirical-law-lln}
\widehat{\mathfrak A}_p \rightharpoonup \mathfrak A\ \text{ in }\Pp(\Pp(\Omega)),
\qquad
\bar\alpha_{\widehat{\mathfrak A}_p}\rightharpoonup\bar\alpha_{\mathfrak A}\ \text{ in }\Pp(\Omega).
\end{equation}
Moreover, if $\mathcal B_c(\mathfrak A)=\{\widetilde\alpha_{\mathfrak A}\}$ is a singleton and if $\widetilde\alpha_{\widehat{\mathfrak A}_p}\in\mathcal B_c(\widehat{\mathfrak A}_p)$ is any empirical barycenter, then, almost surely,
\eql{\label{eq-wow-population-barycenter-lln}
\widetilde\alpha_{\widehat{\mathfrak A}_p}
\rightharpoonup
\widetilde\alpha_{\mathfrak A}
\quad\text{in }\Pp(\Omega).
}
\end{prop}
\begin{proof}
Set $K=\Pp(\Omega)$. Since $\Omega$ is compact metric, $K$ is compact metric for weak convergence, and so is $\Pp(K)$. The space $\Cc(K)$ is separable for the uniform norm. Applying the scalar strong law to a countable dense family of test functions and then using uniform approximation gives, almost surely, for every $\Phi\in\Cc(K)$,
\(\int \Phi(\alpha)\d\widehat{\mathfrak A}_p(\alpha) = \frac1p\sum_{i=1}^p\Phi(\alpha_i) \longrightarrow \int \Phi(\alpha)\d\mathfrak A(\alpha)\)
This is exactly~\eqref{eq-wow-empirical-law-lln}.

For the collapsed mixtures, take $f\in\Cc(\Omega)$ and define $\Phi_f(\alpha)=\int_\Omega f\d\alpha$. This function is continuous on $\Pp(\Omega)$. Hence
\[
\int_\Omega f\d\bar\alpha_{\widehat{\mathfrak A}_p}
=
\int_{\Pp(\Omega)}\Phi_f(\alpha)\d\widehat{\mathfrak A}_p(\alpha)
\longrightarrow
\int_{\Pp(\Omega)}\Phi_f(\alpha)\d\mathfrak A(\alpha)
=
\int_\Omega f\d\bar\alpha_{\mathfrak A},
\]
which proves the second convergence in~\eqref{eq-wow-empirical-law-lln}.

Since $c$ is continuous on the compact product $\Omega\times\Omega$, the value $(\beta,\alpha)\mapsto\MK_c(\beta,\alpha)$ is continuous on $K^2$. Therefore the map $\beta\mapsto h_\beta$, where $h_\beta(\alpha)=\MK_c(\beta,\alpha)$, is continuous from the compact space $K$ to $\Cc(K)$. Its image
\(\mathcal H \eqdef \{h_\beta:\beta\in K\}\)
is compact in $\Cc(K)$. The convergence of $\widehat{\mathfrak A}_p$ to $\mathfrak A$ is uniform over $\mathcal H$: given $\eta>0$, cover $\mathcal H$ by finitely many $\eta$-balls in $\norm{\cdot}_\infty$, use weak convergence for the centers, and bound the error on the balls by the total variation of $\widehat{\mathfrak A}_p-\mathfrak A$. Hence the empirical objectives
$F_p(\beta)\eqdef\int\MK_c(\beta,\alpha)\d\widehat{\mathfrak A}_p(\alpha)$
converge uniformly on $K$ to
$F(\beta)\eqdef\int\MK_c(\beta,\alpha)\d\mathfrak A(\alpha)$.
Let $\beta_p\in\mathcal B_c(\widehat{\mathfrak A}_p)$. Compactness gives a subsequence $\beta_{p_k}\rightharpoonup\beta$. Uniform convergence, continuity of $F$, and optimality of $\beta_{p_k}$ give, for any $\gamma\in K$,
\(F(\beta) = \lim_k F_{p_k}(\beta_{p_k}) \leq \lim_k F_{p_k}(\gamma) = F(\gamma),\)
so $\beta\in\mathcal B_c(\mathfrak A)$. If this set is the singleton $\{\widetilde\alpha_{\mathfrak A}\}$, every converging subsequence has the same limit, and therefore the whole sequence $\widetilde\alpha_{\widehat{\mathfrak A}_p}$ converges to $\widetilde\alpha_{\mathfrak A}$, proving~\eqref{eq-wow-population-barycenter-lln}.
\end{proof}
\paragraph{Toward central limit theorems on Wasserstein space.}
There are nevertheless important settings where such results can be proved. In one dimension, the quantile representation linearizes $\Wass_2$, so barycenter fluctuations can be studied through empirical averages of quantile functions.

\paragraph{Sliced and Radon barycenters.}
Slicing gives a scalable surrogate for high-dimensional barycenters by applying the one-dimensional quantile formula of Proposition~\ref{prop-quantile-barycenters} in every projection direction. For inputs $\be_s\in\Pp_2(\RR^d)$, define the sliced barycenter by replacing $\Wass_2$ with $\SW_2$, introduced in Definition~\ref{def-sliced-wasserstein} and interpreted through the Radon transform in Section~\ref{sec-sliced-wasserstein}:
\[
\umin{\al\in\Pp_2(\RR^d)}
\sum_{s=1}^S\la_s\SW_2(\al,\be_s)^2
=
\umin{\al\in\Pp_2(\RR^d)}
\int_{\Sphere^{d-1}}
\sum_{s=1}^S\la_s
\Wass_2\big((P_\theta)_\sharp\al,(P_\theta)_\sharp\be_s\big)^2
\d\sigma(\theta).
\]
A cheaper Radon-domain approximation drops this consistency constraint and minimizes directly over measurable collections $(\gamma_\theta)_{\theta\in\Sphere^{d-1}}$ with $\gamma_\theta\in\Pp_2(\RR)$,
\[
\umin{(\gamma_\theta)_\theta}
\int_{\Sphere^{d-1}}
\sum_{s=1}^S\la_s
\Wass_2(\gamma_\theta,\be_{s,\theta})^2
\d\sigma(\theta),
\qquad
\be_{s,\theta}\eqdef(P_\theta)_\sharp\be_s.
\]
For each $\theta$, the minimizer is the one-dimensional barycenter $\gamma_\theta$ with quantile
\(\cumul{\gamma_\theta}^{-1}(r) = \sum_{s=1}^S\la_s\cumul{\be_{s,\theta}}^{-1}(r), \qquad r\in(0,1).\)
For two inputs $\be_0,\be_1$ and interpolation weight $t\in[0,1]$, it is useful to make this directionwise construction explicit. For $r\in(0,1)$, define
\begin{equation}\label{eq-radon-barycenter-quantile-field}
Q_i(\theta,r)
\eqdef
\cumul{\be_{i,\theta}}^{-1}(r),
\qquad
Q_t(\theta,r)
\eqdef
(1-t)Q_0(\theta,r)+tQ_1(\theta,r),
\qquad
\gamma_{t,\theta}
\eqdef
\bigl(Q_t(\theta,\cdot)\bigr)_\sharp\mathrm{Leb}_{[0,1]}.
\end{equation}
The values assigned to $Q_t(\theta,\cdot)$ at $r=0$ and $r=1$ are immaterial to this push-forward. Thus $Q_t$ is the quantile field of the independently interpolated projected laws $(\gamma_{t,\theta})_\theta$.

Let $h(\theta,t)$ denote a density of $\gamma_\theta$. We use the one-dimensional Fourier transform in the variable $t$,
\begin{equation}\label{eq-radon-sinogram-fourier}
\widehat h(\theta,\omega)
\eqdef
\int_{\RR}e^{-\imath\omega t}h(\theta,t)\d t,
\qquad
h(\theta,t)
=
\frac{1}{2\pi}\int_{\RR}e^{\imath\omega t}\widehat h(\theta,\omega)\d\omega,
\end{equation}
\begin{prop}[Radon least-squares pseudoinverse]\label{prop-radon-pseudoinverse}
Let $d\geq2$, let $\sigma$ be the uniform probability measure on $\Sphere^{d-1}$, and let $R$ be the density Radon transform introduced in the \hyperref[rem-sliced-radon-viewpoint]{Radon point of view} paragraph. Write
\(\mathcal D(R) \eqdef \enscond{\rho\in L^2(\RR^d)}{R\rho\in L^2(\Sphere^{d-1}\times\RR,\d\sigma\,\d t)}.\)
Let $h\in L^2(\Sphere^{d-1}\times\RR,\d\sigma\,\d t)$ and assume that the Fourier expressions below define an element of $\mathcal D(R)$. Then the unique solution of
\begin{equation}\label{eq-radon-least-squares}
\umin{\rho\in\mathcal D(R)}
\int_{\Sphere^{d-1}}\int_{\RR}
\abs{R\rho(\theta,t)-h(\theta,t)}^2
\d t\,\d\sigma(\theta)
\end{equation}
is $\rho^\dagger=R^\dagger h$, where
\begin{equation}\label{eq-radon-pseudoinverse}
R^\dagger h(x)
=
\frac{\abs{\Sphere^{d-1}}}{2(2\pi)^d}
\int_{\Sphere^{d-1}}\int_{\RR}
e^{\imath\omega\dotp{\theta}{x}}
|\omega|^{d-1}
\widehat h(\theta,\omega)
\d\omega\,\d\sigma(\theta).
\end{equation}
Thus $R^\dagger h$ is an $L^2$ pseudoinverse reconstruction. Without an additional $L^1$ assumption, it need not be the density of a finite signed measure. If $h=R\rho$ is Radon-consistent and $\rho$ is sufficiently regular, then $R^\dagger R\rho=\rho$.
\end{prop}
\begin{proof}
For the $d$-dimensional Fourier transform, use the compatible convention
\[
\widehat\rho(\xi)
=
\int_{\RR^d}e^{-\imath\dotp{\xi}{x}}\rho(x)\d x,
\qquad
\rho(x)
=
\frac{1}{(2\pi)^d}\int_{\RR^d}e^{\imath\dotp{\xi}{x}}\widehat\rho(\xi)\d\xi.
\]
\eqllead{The Fourier-slice theorem gives}{eq-radon-fourier-slice}{\widehat{R\rho}(\theta,\omega)=\widehat\rho(\omega\theta).}
Applying the one-dimensional Plancherel identity in $t$ rewrites the objective in~\eqref{eq-radon-least-squares}, up to the positive factor $1/(2\pi)$, as
$\int_{\Sphere^{d-1}}\int_{\RR} \abs{\widehat\rho(\omega\theta)-\widehat h(\theta,\omega)}^2 \d\omega\,\d\sigma(\theta)$.
Every $\xi\neq0$ has the two signed-polar representations
$(\theta,\omega)=(\xi/\norm{\xi},\norm{\xi})$ and
$(-\xi/\norm{\xi},-\norm{\xi})$. The pointwise least-squares minimizer is therefore
\begin{equation}\label{eq-radon-pseudoinverse-fourier}
\widehat{\rho^\dagger}(\xi)=\frac12\left(\widehat h\left(\frac{\xi}{\norm{\xi}},\norm{\xi}\right)+\widehat h\left(-\frac{\xi}{\norm{\xi}},-\norm{\xi}\right)\right).
\end{equation}
This also proves uniqueness under the stated existence assumptions. Applying the inverse $d$-dimensional Fourier transform to~\eqref{eq-radon-pseudoinverse-fourier} and using signed polar coordinates yields~\eqref{eq-radon-pseudoinverse}; the factor $\abs{\Sphere^{d-1}}/2$ accounts for the normalization of $\sigma$ and the two signed representations of each nonzero frequency. Finally, if $h=R\rho$, then~\eqref{eq-radon-fourier-slice} makes both terms in~\eqref{eq-radon-pseudoinverse-fourier} equal to $\widehat\rho(\xi)$, proving $R^\dagger R\rho=\rho$.
\end{proof}
Because this multiplier amplifies high frequencies, one typically chooses a bandwidth $\Omega>0$ and replaces it by
\begin{equation}\label{eq-radon-windowed-ramp}
m_\Omega(\omega)
\eqdef
|\omega|^{d-1}\chi(\omega/\Omega),
\qquad
R_\Omega^\dagger h(x)
\eqdef
\frac{\abs{\Sphere^{d-1}}}{2(2\pi)^d}
\int_{\Sphere^{d-1}}\int_{\RR}
e^{\imath\omega\dotp{\theta}{x}}
m_\Omega(\omega)\widehat h(\theta,\omega)
\d\omega\,\d\sigma(\theta),
\end{equation}
where $\chi$ is an even low-pass window with $\chi(0)=1$. For the two-input path~\eqref{eq-radon-barycenter-quantile-field}, choose $\eta_t\geq0$ so that $((R_\Omega^\dagger h_t)-\eta_t)_+\in L^1(\RR^d)$ has finite nonzero mass, and define the nonnegative, unit-mass reconstruction
\begin{equation}\label{eq-radon-display-reconstruction}
A_t(x)
\eqdef
\frac{\bigl((R_\Omega^\dagger h_t)(x)-\eta_t\bigr)_+}
{\displaystyle\int_{\RR^d}\bigl((R_\Omega^\dagger h_t)(z)-\eta_t\bigr)_+\d z}.
\end{equation}
\subsection{Multimarginal OT}
\label{sec-multimarginal-ot}
\label{subsec:multi-marginal}
\paragraph{Definition and basic structure.}
\begin{defn}[Multimarginal optimal transport]\label{def-multimarginal-ot}
Let $\al_s\in\Pp(\Xx_s)$ for $s=1,\ldots,S$, and denote by
$\Couplings(\al_1,\ldots,\al_S)$ the probability measures whose
$s$-th marginal is $\al_s$. For a measurable cost $c$ bounded from below,
define
\begin{equation}\label{eq-multimarginal-ot}
\operatorname{MMOT}_c(\al_1,\ldots,\al_S)
\eqdef
\inf_{\pi\in\Couplings(\al_1,\ldots,\al_S)}
\int_{\Xx_1\times\cdots\times\Xx_S}
c(x_1,\ldots,x_S)\d\pi(x_1,\ldots,x_S).
\end{equation}
\end{defn}
\paragraph{Monge structure and splitting-set twist.}
\begin{defn}[Twist on splitting sets]\label{def-twist-splitting-sets}
Fix $x_1\in\Xx_1$. A set $M\subset\Xx_2\times\cdots\times\Xx_S$ is a $c$-splitting set at $x_1$ if there exist functions $u_s:\Xx_s\to\RR\cup\{-\infty\}$, for $s=2,\ldots,S$, such that
\(\sum_{s=2}^S u_s(x_s)\leq c(x_1,x_2,\ldots,x_S)\)
for all $(x_2,\ldots,x_S)$, with equality on $M$. Assume $c$ is differentiable in $x_1$. The cost is twisted on splitting sets if, for every $x_1$ and every $c$-splitting set $M$ at $x_1$, the map
\((x_2,\ldots,x_S) \longmapsto \nabla_{x_1}c(x_1,x_2,\ldots,x_S)\)
is injective on $M$.
\end{defn}
\begin{prop}[Multi-marginal Monge structure]\label{prop-multimarginal-monge-structure}
Assume that $\Xx_s\subset\RR^d$, that $c$ is continuous and differentiable with respect to $x_1$, and that $c$ is twisted on splitting sets. Suppose that Kantorovich dual maximizers $(\varphi_s)_{s=1}^S$ exist, that $\al_1$ is absolutely continuous, and that $\varphi_1$ is differentiable $\al_1$-a.e.\ Then every optimal plan $\pi^\star\in\Couplings(\al_1,\ldots,\al_S)$ is concentrated on the graph of maps
\(\pi^\star=(\Id,\T_2,\ldots,\T_S)_\sharp\al_1, \qquad (\T_s)_\sharp\al_1=\al_s.\)
In particular, under these hypotheses the optimizer is unique.
\end{prop}
\begin{proof}
Let $(\varphi_s)_{s=1}^S$ be optimal dual potentials. Complementary slackness gives a Borel contact set $\Gamma$ of full $\pi^\star$-measure on which $\sum_s\varphi_s(x_s)=c(x_1,\ldots,x_S)$. After disintegrating with respect to the first marginal, fix a point $x_1$ where $\varphi_1$ is differentiable and where the conditional plan is concentrated on the fiber
\(M(x_1)=\{(x_2,\ldots,x_S):(x_1,x_2,\ldots,x_S)\in\Gamma\}.\)
Indeed, set $u_2=\varphi_2+\varphi_1(x_1)$ and $u_s=\varphi_s$ for $s\geq3$. Dual feasibility gives $\sum_{s=2}^S u_s(x_s)\leq c(x_1,x_2,\ldots,x_S)$, and equality holds on $M(x_1)$. If $(x_2,\ldots,x_S)\in M(x_1)$, the function
\(z\longmapsto c(z,x_2,\ldots,x_S)-\sum_{s=2}^S\varphi_s(x_s)\)
touches $\varphi_1$ from above at $z=x_1$. Differentiating at this contact point gives \(\nabla\varphi_1(x_1)=\nabla_{x_1}c(x_1,x_2,\ldots,x_S).\)
All points in the fiber therefore have the same value of $\nabla_{x_1}c$. Twist on splitting sets makes the fiber a singleton for $\al_1$-a.e.\ $x_1$. Disintegrating $\pi^\star$ with respect to its first marginal gives Dirac conditional measures, hence measurable maps $(\T_2,\ldots,\T_S)$. If $\pi^1$ and $\pi^2$ are two optimal plans, their average is also optimal. The conditional measure of this average over $x_1$ is the average of the two Dirac conditionals, and it must again be a Dirac mass by the preceding argument. Hence the two Dirac masses coincide for $\al_1$-a.e.\ $x_1$, proving uniqueness.
\end{proof}
\paragraph{Coulomb cost and density-functional theory.}
For $N$ electrons in $\RR^3$, the repulsive Coulomb interaction is the multi-body cost
\(c_{\mathrm{Coul}}(x_1,\ldots,x_N) \eqdef \sum_{1\leq i<j\leq N}\frac{1}{\norm{x_i-x_j}},\)
Away from collisions, \(\nabla_{x_1}c_{\mathrm{Coul}}(x_1,\ldots,x_N) = -\sum_{j=2}^N\frac{x_1-x_j}{\norm{x_1-x_j}^3}.\)
For $N=2$, the map from $x_2$ to this vector is injective, so the ordinary two-marginal twist argument can be recovered under the required existence, duality and differentiability hypotheses. For $N\geq3$, however, the displayed total force does not by itself determine the entire tuple $(x_2,\ldots,x_N)$; twist on splitting sets, and hence a Monge representation, is not automatic.

If $\rho$ is an electron density with $\int_{\RR^3}\rho(x)\d x=N$ and $\al=\rho/N$ is the associated probability density, the strictly-correlated-electrons relaxation of density-functional theory is the equal-marginal problem
\[
V_{\mathrm{ee}}^{\mathrm{SCE}}[\rho]
\eqdef
\inf_{\pi\in\Couplings(\al,\ldots,\al)}
\int_{(\RR^3)^N}
c_{\mathrm{Coul}}(x_1,\ldots,x_N)
\d\pi(x_1,\ldots,x_N).
\]
The deterministic ansatz writes a plan through co-motion maps
\(\pi=(\Id,\T_2,\ldots,\T_N)_\sharp\al, \qquad (\T_i)_\sharp\al=\al,\)
\begin{prop}[Cyclic co-motion plans]\label{prop-cyclic-comotion-plans}
Let $\al\in\Pp(\RR^3)$ and let $\T:\RR^3\to\RR^3$ be measurable with $\T_\sharp\al=\al$ and $\T^N=\Id$ $\al$-a.e.\ Set $\T^0=\Id$ and
\(\pi_\T=(\T^0,\T^1,\ldots,\T^{N-1})_\sharp\al.\)
Then $\pi_\T\in\Couplings(\al,\ldots,\al)$. If $R_\sigma(x_1,\ldots,x_N)=(x_{\sigma(1)},\ldots,x_{\sigma(N)})$ for $\sigma\in\Perm(N)$, then
\(\bar\pi_\T\eqdef\frac1{N!}\sum_{\sigma\in\Perm(N)}(R_\sigma)_\sharp\pi_\T\)
is a symmetric admissible plan and
\[
\int c_{\mathrm{Coul}}\d\bar\pi_\T
=
\int c_{\mathrm{Coul}}\d\pi_\T
=
\int_{\RR^3}
\sum_{0\leq i<j\leq N-1}
\frac{1}{\norm{\T^i(x)-\T^j(x)}}\d\al(x).
\]
\end{prop}
\begin{proof}
Since $\T_\sharp\al=\al$, all iterates $\T^i$ preserve $\al$, so every marginal of $\pi_\T$ is $\al$. The plan $\bar\pi_\T$ is an average of coordinate permutations of $\pi_\T$, hence has the same marginals and is invariant under coordinate permutations. The Coulomb cost is symmetric in its arguments, so its integral is unchanged by each $R_\sigma$. The last identity follows by evaluating $c_{\mathrm{Coul}}$ on the graph $(x,\T(x),\ldots,\T^{N-1}(x))$.
\end{proof}
Proposition~\ref{prop-multimarginal-monge-structure} explains the general mechanism that can force graph solutions, while Proposition~\ref{prop-cyclic-comotion-plans} records the additional equal-marginal symmetry used by co-motion maps. Thus the DFT problem is both a central application and a warning that multi-marginal OT is richer than a naive deterministic matching problem.

\paragraph{Multi-marginal formulation of barycenters.}
For the squared Euclidean cost, one can introduce a latent barycenter point and eliminate it explicitly, leading to the multi-marginal cost
\(c_{\mathrm{bar}}(x_1,\ldots,x_S) = \min_{x\in\RR^d} \sum_{s=1}^S\la_s\norm{x-x_s}^2.\)
\begin{prop}[Multi-marginal formula for quadratic barycenters]\label{prop-multimarginal-barycenter}
Let $\be_1,\ldots,\be_S\in\Pp_2(\RR^d)$ and $\la\in\simplex_S$. Define
\(B(x_1,\ldots,x_S)=\sum_{s=1}^S\la_s x_s, \qquad c_{\mathrm{bar}}(x_1,\ldots,x_S)=\min_x\sum_s\la_s\norm{x-x_s}^2.\)
If $\pi^\star$ solves the multi-marginal OT problem with marginals $(\be_s)_s$ and cost $c_{\mathrm{bar}}$, then $\al^\star=B_\sharp\pi^\star$ is a Wasserstein barycenter. Conversely, every barycenter is obtained this way from an optimal multi-marginal plan.
\end{prop}
\begin{proof}
For any candidate barycenter $\al$ and couplings $\pi_s\in\Couplings(\al,\be_s)$, glue the couplings along their common $\al$ marginal to obtain a joint law of $(X,Y_1,\ldots,Y_S)$. Conditioning on $(Y_s)_s$ and minimizing over $X$ gives
\(\sum_s\la_s\EE\norm{X-Y_s}^2 \geq \EE\min_x\sum_s\la_s\norm{x-Y_s}^2 = \EE c_{\mathrm{bar}}(Y_1,\ldots,Y_S).\)
Taking the infimum over the couplings gives that the barycenter value is at least the multi-marginal value. Conversely, from an optimal multi-marginal plan $\pi^\star$, set $X=B(Y_1,\ldots,Y_S)$. The couplings between $X$ and each $Y_s$ are feasible for the barycenter problem and attain exactly the multi-marginal cost, proving equality and the formula.
If $\al^\star$ is any barycenter, choose optimal couplings between $\al^\star$ and each $\be_s$ and glue them along the common $\al^\star$ marginal. Since the barycenter and multi-marginal values are equal, the conditional minimization inequality above must be an equality. Thus $X=B(Y_1,\ldots,Y_S)$ almost surely for the induced optimal multi-marginal plan, and $\al^\star=B_\sharp\pi^\star$.
\end{proof}
\begin{cor}[Sparse discrete barycenters]\label{cor-discrete-barycenters}
Let
\(\be_s=\sum_{i_s=1}^{n_s}b_{s,i_s}\de_{x_{s,i_s}}, \qquad s=1,\ldots,S,\)
be discrete input measures. There exists a quadratic Wasserstein barycenter $\al^\star$ supported on weighted averages $\sum_s\la_s x_{s,i_s}$ of one support point from each input and satisfying
\(\#\supp(\al^\star) \leq \sum_{s=1}^S n_s-S+1.\)
\end{cor}
\begin{proof}
Write a discrete multi-marginal plan as a nonnegative tensor
$\P=(\P_{i_1,\ldots,i_S})$ and let $\mathcal A_{\mathrm{marg}}$ collect its $S$ marginals:
$\bigl(\mathcal A_{\mathrm{marg}}\P\bigr)_{s,i_s} = \sum_{(i_r)_{r\neq s}}\P_{i_1,\ldots,i_S}$.
After vectorizing $\P$, Proposition~\ref{prop-lp-rank-sparsity} applies to the multi-marginal linear program. Its constraint operator has rank
$\rank(\mathcal A_{\mathrm{marg}}) = \sum_{s=1}^S n_s-S+1$.
Indeed, a family of vectors $u_s\in\RR^{n_s}$ belongs to the annihilator of its image precisely when
$\sum_{s=1}^S u_{s,i_s}=0$ for every $(i_1,\ldots,i_S)$.
Varying one index at a time shows that every $u_s$ is constant, say $u_s=c_s\ones_{n_s}$, and the remaining condition is $\sum_s c_s=0$. The annihilator therefore has dimension $S-1$, which proves the rank formula.

The rank-controlled sparsity result therefore supplies an optimal multi-marginal tensor $\P^\star$ with at most $\sum_s n_s-S+1$ positive entries. Proposition~\ref{prop-multimarginal-barycenter} gives the barycenter $\al^\star=B_\sharp\P^\star$. Each positive tensor entry produces at most one atom under $B$, and collisions can only reduce the support size. This proves the claim.
\end{proof}
\paragraph{Entropic regularization of multi-marginal OT.}
As in the two-marginal case, one can add an entropic penalty with respect to the product measure $\al_1\otimes\cdots\otimes\al_S$:
\(\inf_{\pi\in\Couplings(\al_1,\ldots,\al_S)} \int c\d\pi+\epsilon\KL(\pi|\al_1\otimes\cdots\otimes\al_S).\)
Whenever such potentials $(f_s)_s$ attain the entropic dual, the optimizer has the generalized Gibbs form
\[
\d\pi^\star(x_1,\ldots,x_S)
=
\exp\!\left(\frac{\sum_s f_s(x_s)-c(x_1,\ldots,x_S)}{\epsilon}\right)
\prod_s\d\al_s(x_s),
\]
\paragraph{Treewidth and graphical structure.}
The important exception is when the cost factors over a sparse interaction graph. Let $G=(V,E)$ be a finite undirected graph with $V=\{1,\ldots,S\}$ and suppose, for simplicity, that
\begin{equation}\label{eq-multimarginal-graph-cost}
c(x_1,\ldots,x_S)
=
\sum_{(r,s)\in E}c_{r,s}(x_r,x_s).
\end{equation}
\begin{defn}[Tree decomposition and treewidth]\label{def-treewidth}
A \emph{tree decomposition} of a finite graph $G=(V,E)$ is a tree $\mathcal T=(Q,F)$ together with bags $B_q\subseteq V$, $q\in Q$, such that
\begin{enumerate}
\item $\bigcup_{q\in Q}B_q=V$;
\item for every edge $(r,s)\in E$, some bag contains both $r$ and $s$;
\item for every vertex $s\in V$, the nodes $\enscond{q\in Q}{s\in B_q}$ form a connected subtree of $\mathcal T$.
\end{enumerate}
The width of this decomposition is $\max_{q\in Q}|B_q|-1$, and the \emph{treewidth} $\operatorname{tw}(G)$ is the minimum width over all tree decompositions.
\end{defn}
\paragraph{Junction-tree contractions inside Sinkhorn.}
In the discrete setting, define the edge Gibbs matrices and the current unary factors
\[
\K^{r,s}_{i_r,i_s}
=
\exp\!\left(-\frac{\C^{r,s}_{i_r,i_s}}{\epsilon}\right),
\qquad
h_s(i_s)
=
(\a_s)_{i_s}(u_s)_{i_s}.
\]
Under~\eqref{eq-multimarginal-graph-cost}, the current scaled coupling is represented implicitly as
\begin{equation}\label{eq-multimarginal-graph-factorization}
\P_{i_1,\ldots,i_S}
=
\prod_{s\in V}h_s(i_s)
\prod_{(r,s)\in E}\K^{r,s}_{i_r,i_s}.
\end{equation}
When $G$ is a tree, let $N_G(r)$ denote the neighbors of $r$. The directed sum-product messages satisfy
\begin{equation}\label{eq-multimarginal-tree-message}
m_{r\to s}(i_s)
=
\sum_{i_r=1}^{n_r}
\K^{r,s}_{i_r,i_s}h_r(i_r)
\prod_{q\in N_G(r)\setminus\{s\}}
m_{q\to r}(i_r).
\end{equation}
A leaf-to-root pass followed by a root-to-leaf pass computes every directed message and hence every current marginal,
\begin{equation}\label{eq-multimarginal-tree-marginal}
(\widehat{\a}_s)_{i_s}
=
h_s(i_s)
\prod_{r\in N_G(s)}m_{r\to s}(i_s).
\end{equation}
For the block $s$ currently selected by the cyclic scaling algorithm, the exact coordinate update is then $u_s \leftarrow u_s\odot\frac{\a_s}{\widehat{\a}_s}$, with componentwise products and quotients. This changes only the unary factor $h_s$.

For a general tree decomposition, choose one host bag for each unary factor, assign each edge factor in~\eqref{eq-multimarginal-graph-factorization} to a bag containing both endpoints, and denote the product assigned to bag $B_q$ by $\psi_q(i_{B_q})$. For adjacent bags $q,q'\in Q$, set $\mathcal S_{q,q'}=B_q\cap B_{q'}$. Junction-tree messages take the form
\begin{equation}\label{eq-multimarginal-junction-message}
M_{q\to q'}(i_{\mathcal S_{q,q'}})
=
\sum_{i_{B_q\setminus\mathcal S_{q,q'}}}
\psi_q(i_{B_q})
\prod_{\ell\in N_{\mathcal T}(q)\setminus\{q'\}}
M_{\ell\to q}(i_{\mathcal S_{\ell,q}}).
\end{equation}
After a collect-distribute pass, the calibrated bag belief is
\begin{equation}\label{eq-multimarginal-junction-belief}
\mathfrak b_q(i_{B_q})
=
\psi_q(i_{B_q})
\prod_{\ell\in N_{\mathcal T}(q)}
M_{\ell\to q}(i_{\mathcal S_{\ell,q}}).
\end{equation}
For any $s\in B_q$, summing $\mathfrak b_q$ over $B_q\setminus\{s\}$ gives $\widehat{\a}_s$; the running-intersection property ensures that every bag containing $s$ gives the same result.

Redundant adjacent bags may be contracted, so one can assume that the decomposition is reduced. If its width is $w$, every bag contains at most $w+1$ variables and every separator contains at most $w$.

\subsection{Low-Rank Optimal Transport}
\label{sec-low-rank-ot}
\begin{defn}[Low-rank factored couplings]\label{def-low-rank-couplings}
Let $\a\in\simplex_n$, $\b\in\simplex_m$ and let $r\geq1$. A rank-$r$ factored coupling is a triple $(\Q,\R,g)$ such that
$\Q\in\RR_+^{n\times r}$, $\R\in\RR_+^{m\times r}$, and $g\in\simplex_r$,
with
$\Q\ones_r=\a$, $\R\ones_r=\b$, and $\transp \Q\ones_n=\transp \R\ones_m=g$.
It induces the coupling
\begin{equation}\label{eq-low-rank-coupling-factor}
\P(\Q,\R,g)
=
\Q\diag(g)^{-1}\transp \R,
\qquad
\P_{i,j}=\sum_{k=1}^r \frac{\Q_{i,k}\R_{j,k}}{g_k},
\end{equation}
where columns with $g_k=0$ are discarded before applying the formula.
\end{defn}
\begin{prop}[Factored couplings and nonnegative rank]\label{prop-low-rank-factorization}
For every feasible triple $(\Q,\R,g)$ in Definition~\ref{def-low-rank-couplings}, the matrix $\P(\Q,\R,g)$ belongs to $\CouplingsD(\a,\b)$ and has nonnegative rank at most $r$. Conversely, every coupling $\P\in\CouplingsD(\a,\b)$ with nonnegative rank at most $r$ admits a representation of the form~\eqref{eq-low-rank-coupling-factor}, after possibly deleting zero latent components.
\end{prop}
\begin{proof}
The marginal constraints follow by direct summation: \(\P\ones_m=\Q\diag(g)^{-1}\transp \R\ones_m =\Q\diag(g)^{-1}g=\Q\ones_r=\a,\)
and similarly $\transp \P\ones_n=\b$. Since $\P$ is a sum of $r$ nonnegative rank-one matrices $\Q_{:,k}\R_{:,k}^\top/g_k$, its nonnegative rank is at most $r$.

Conversely, suppose $\P=UV^\top$ with $U\in\RR_+^{n\times r}$ and $V\in\RR_+^{m\times r}$. Set $q_k=\sum_i U_{i,k}$ and $s_k=\sum_j V_{j,k}$. Components with $q_ks_k=0$ do not contribute and can be removed. Define $g_k=q_ks_k$, $\Q_{i,k}=U_{i,k}s_k$ and $\R_{j,k}=V_{j,k}q_k$. Since $\P$ has total mass one, $g\in\simplex_r$. Moreover $\Q\ones_r=\P\ones_m=\a$, $\R\ones_r=\transp \P\ones_n=\b$, and both column marginals are equal to $g$. Finally,
$\sum_k\frac{\Q_{i,k}\R_{j,k}}{g_k} = \sum_k\frac{U_{i,k}s_kV_{j,k}q_k}{q_ks_k} =\P_{i,j}$.
\end{proof}
For a cost matrix $\C\in\RR^{n\times m}$, the low-rank constrained OT value is therefore
\(\min_{(\Q,\R,g)} \dotp{\C}{\Q\diag(g)^{-1}\transp \R},\)
The resulting fixed-latent-mass model is
\begin{equation}\label{eq-low-rank-entropic-ot}
\min_{\substack{\Q\ones_r=\a,\ \Q^\top\ones_n=g\\
\R\ones_r=\b,\ \R^\top\ones_m=g}}
\dotp{\C}{\Q\diag(g)^{-1}\transp \R}
+
\epsilon\KLD(\Q|\a\otimes g)
+
\epsilon\KLD(\R|\b\otimes g),
\end{equation}
where $\Q,\R$ are nonnegative. For fixed $g$, this differs only by constants from adding $-\epsilon(\HD(\Q)+\HD(\R))$ to the factorized transport cost.

\subsection{Capacity-Constrained Optimal Transport}
\label{sec-capacity-constrained-ot}
\begin{defn}[Capacity-constrained optimal transport]\label{def-capacity-constrained-ot}
Let $\alpha\in\Mm_+^1(\Xx)$, $\beta\in\Mm_+^1(\Yy)$, let $c:\Xx\times\Yy\to[0,+\infty]$ be measurable, and let $u:\Xx\times\Yy\to[0,+\infty]$ be a measurable capacity. Define
\begin{equation}\label{eq-capacity-constrained-ot}
\MK_c^u(\alpha,\beta)
\eqdef
\inf_{\pi\in\Couplings(\alpha,\beta)}
\enscond{
\int_{\Xx\times\Yy} c(x,y)\d\pi(x,y)
}{
\pi\ll\alpha\otimes\beta,\quad
\frac{\d\pi}{\d(\alpha\otimes\beta)}(x,y)\leq u(x,y)
\quad (\alpha\otimes\beta)\text{-a.e.}
}.
\end{equation}
By convention, $u\equiv+\infty$ removes both the density bound and the absolute-continuity requirement, so that $\MK_c^{+\infty}(\alpha,\beta)=\MK_c(\alpha,\beta)$.
\end{defn}
The product coupling $\alpha\otimes\beta$ is feasible whenever $u\geq1$. Thus the constraint is not meant to promote the product plan; rather, it prevents the optimizer from using any pair more than the prescribed density ratio.

For discrete measures $\alpha=\sum_i a_i\de_{x_i}$ and $\beta=\sum_j b_j\de_{y_j}$, a capacity is an upper matrix $\overline{\P}\in\RR_+^{n\times m}$. Given a cost matrix $\C$, the density-ratio discretization of Definition~\ref{def-capacity-constrained-ot}, with $\overline{\P}_{i,j}=u_{i,j}a_i b_j$, is the capped linear program
\begin{equation}\label{eq-discrete-capacity-constrained-ot}
\min_{\P\in\CouplingsD(\a,\b)}
\dotp{\C}{\P}
\quad\text{subject to}\quad
0\leq \P_{i,j}\leq \overline{\P}_{i,j}\quad\forall(i,j).
\end{equation}
\begin{prop}[Feasibility of a capped transport polytope]\label{prop-capacity-feasibility}
The constraints in~\eqref{eq-discrete-capacity-constrained-ot} are feasible if and only if
\begin{equation}\label{eq-capacity-cut-condition}
a(I)+b(J)-1\leq \overline{\P}(I,J)
\qquad
\text{for every }I\subset\{1,\ldots,n\},\ J\subset\{1,\ldots,m\}.
\end{equation}
\end{prop}
\begin{proof}
If $\P$ is feasible, then
\(\P(I,J) = a(I)-\P(I,J^c) \geq a(I)-b(J^c) = a(I)+b(J)-1.\)
Since $\P(I,J)\leq \overline{\P}(I,J)$, condition~\eqref{eq-capacity-cut-condition} is necessary.

For sufficiency, build a flow network with capacity $a_i$ from the source to row vertex $i$, capacity $\overline{\P}_{i,j}$ from row $i$ to column $j$, and capacity $b_j$ from column $j$ to the sink. A cut containing row set $I$ and column set $J^c$ has capacity
\(1-a(I)+\overline{\P}(I,J)+1-b(J)\).
Condition~\eqref{eq-capacity-cut-condition} says that every such cut has capacity at least one. The max-flow/min-cut theorem therefore gives a unit flow, and its row-to-column edge values form the required matrix $\P$.
\end{proof}
With $\K_{i,j}=a_i b_j e^{-\C_{i,j}/\epsilon}$, the regularized problem is
\begin{equation}\label{eq-entropic-capacity-constrained-ot}
\min_{\P\in\CouplingsD(\a,\b),\,0\leq \P\leq \overline{\P}}
\dotp{\C}{\P}
+
\epsilon\KLD(\P|\a\otimes\b).
\end{equation}
On the active support, all iterates and correction factors are positive, so every division in the KL--Dykstra iteration is well defined. KL-Dykstra convergence on finite-dimensional affine and box constraints then shows that, when the capped polytope is nonempty, the iterates converge to the unique minimizer of~\eqref{eq-entropic-capacity-constrained-ot}.

\subsection{Metric Learning and Inverse OT}
\label{sec-metric-learning-inverse-ot}
\subsubsection{Differentiating OT losses}
\paragraph{First variations of OT values.}
\begin{prop}[First variations of OT values]\label{prop-ot-first-variations-unregularized}\label{prop-ot-first-variations-entropic}
Let $\alpha$ and $\beta$ be probability measures with compact supports $\Xx=\supp(\alpha)$ and $\Yy=\supp(\beta)$, let $c\in\Cc(\Xx\times\Yy)$, and let $\epsilon\geq0$. Use the convention $\MK_c^0=\MK_c$, where the unregularized and positive-temperature values are defined in Definitions~\ref{def-continuous-kantorovich-problem} and~\ref{def-continuous-entropic-ot}. Denote by $\mathcal O_c^\epsilon(\alpha,\beta)$ and $\mathcal D_c^\epsilon(\alpha,\beta)$ the sets of primal and dual optimizers of the corresponding problem.
If $\chi$ is a signed measure with $\chi(\Xx)=0$ and $\alpha_t=\alpha+t\chi$ is a probability measure for $0\leq t\leq t_0$, then
\[
\left.\frac{\d}{\d t}\right|_{t=0^+}
\MK_c^\epsilon(\alpha_t,\beta)
=
\sup_{(f,g)\in\mathcal D_c^\epsilon(\alpha,\beta)}
\int_{\Xx} f(x)\d\chi(x).
\]
If $h\in\Cc(\Xx\times\Yy)$ and $c_t=c+th$, then
\[
\left.\frac{\d}{\d t}\right|_{t=0^+}
\MK_{c_t}^\epsilon(\alpha,\beta)
=
\inf_{\pi\in\mathcal O_c^\epsilon(\alpha,\beta)}
\int_{\Xx\times\Yy} h(x,y)\d\pi(x,y).
\]
For fixed $(\alpha,\beta)$, a finite signed Radon measure $\eta$ on $\Xx\times\Yy$ represents the first variation of the functional $c\mapsto\MK_c^\epsilon(\alpha,\beta)$ if, for every $h\in\Cc(\Xx\times\Yy)$,
\(\MK_{c+t h}^\epsilon(\alpha,\beta) = \MK_c^\epsilon(\alpha,\beta) +t\int_{\Xx\times\Yy}h\d\eta+o(t) \qquad(t\to0).\)
In particular, if the normalized optimal potential $f_\epsilon^\star$ and the optimal plan $\pi_\epsilon^\star$ are unique (so in particular when $\epsilon>0$), then, with the marginal first variation understood as in Definition~\ref{def-first-variation},
\[
\frac{\delta \MK_c^\epsilon}{\delta\alpha}(\alpha,\beta)=f_\epsilon^\star,
\qquad
\frac{\delta \MK_c^\epsilon}{\delta c}(\alpha,\beta)=\pi_\epsilon^\star.
\]
\end{prop}
\begin{proof}
For $\epsilon=0$, the Kantorovich duality theorem is stated in Proposition~\ref{prop-kantorovich-duality-general}. Its dual value is a supremum of affine functions of $\alpha$, so Danskin's theorem gives the first formula as the supremum over active dual potentials. The condition $\chi(\Xx)=0$ makes this expression independent of the additive gauge $(f,g)\mapsto(f+\lambda,g-\lambda)$.

For $\epsilon>0$, use the continuous entropic duality and primal--dual density law of Proposition~\ref{prop-continuous-entropic-duality}. At an optimal pair, the marginal constraint makes the density in~\eqref{eq-continuous-entropic-density-law} integrate to one with respect to $\beta$ for every $x$ on the source support. Therefore the derivative of the dual objective~\eqref{eq-dual-sinkhorn-objective} with respect to $\alpha$ at fixed optimal potentials is $\int f_\epsilon^\star\d\chi$: the derivative of its exponential term vanishes after the row normalization. Danskin's theorem then gives the same first formula, now with a singleton set of normalized active potentials.

For every $\epsilon\geq0$, the primal objective depends on $c$ only through the affine term $\int c\d\pi$. The minimum form of Danskin's theorem therefore gives the second formula as the infimum of $\int h\d\pi$ over active primal optimizers. When the normalized potential or plan is unique, the corresponding directional derivative is linear in the perturbation, which gives the two first variations.
\end{proof}
In the discrete unregularized case, the primal problem~\eqref{eq-kanto-discr} and its dual~\eqref{eq-dual} show that any optimal dual vector $f^\star$ is a subgradient with respect to the source weights $\a$, while any optimal plan $\P^\star$ is a supergradient with respect to the cost matrix $\C$, because the value is concave in $\C$:
\(f^\star\in\partial_{\a}\MKD_{\C}(\a,\b), \qquad \P^\star\in\partial_{\C}^{\mathrm{sup}}\MKD_{\C}(\a,\b).\)
For $\epsilon>0$, these are ordinary derivatives on the relative interiors of finite-dimensional simplices. For a finite-dimensional parametrization $c_\theta$, the entropic formula gives the backpropagation rule
\(\partial_{\theta}\MK_{c_\theta}^\epsilon(\alpha,\beta) = \int \partial_\theta c_\theta(x,y)\d\pi_{\theta,\epsilon}^\star(x,y),\)
where $\pi_{\theta,\epsilon}^\star$ is the entropic optimizer between $(\alpha,\beta)$ for the cost $c_\theta$. Here we extend $\MK_c^\epsilon$ to this signed cost by the same variational formula. It is proper because every $\pi\in\Couplings(\alpha,\beta)$ satisfies
\[
\int\abs{\dotp{Ax}{y}}\d\pi(x,y)
\leq
\norm{A}_{\mathrm{op}}
\left(\int\norm{x}^2\d\alpha(x)\right)^{1/2}
\left(\int\norm{y}^2\d\beta(y)\right)^{1/2}.
\]
Indeed, after splitting $\{\norm{x}\norm{y}>R\}$ according to whether $\norm{x}>\sqrt R$ or $\norm{y}>\sqrt R$, Cauchy--Schwarz bounds the two tails, up to the factor $\norm{A}_{\mathrm{op}}$, by
\[
\left(\int_{\{\norm{x}>\sqrt R\}}\norm{x}^2\d\alpha(x)\right)^{1/2}
\left(\int\norm{y}^2\d\beta(y)\right)^{1/2}
\qandq
\left(\int\norm{x}^2\d\alpha(x)\right)^{1/2}
\left(\int_{\{\norm{y}>\sqrt R\}}\norm{y}^2\d\beta(y)\right)^{1/2},
\]
which vanish as $R\to+\infty$. Thus the bilinear term is continuous on the weakly compact coupling set, the forward problem is attained, and for $\epsilon>0$ its optimizer is unique. For every perturbation $H\in\RR^{d\times d}$,
\[
D_A\MK_{c_A}^\epsilon(\alpha,\beta)[H]
=
\int \dotp{Hx}{y}\d\pi_{A,\epsilon}^\star(x,y)
=
\left\langle H,\int yx^\top\d\pi_{A,\epsilon}^\star(x,y)\right\rangle_{\mathrm F}.
\]
Consequently,
\[
\nabla_A\MK_{c_A}^\epsilon(\alpha,\beta)
=
\int yx^\top\d\pi_{A,\epsilon}^\star(x,y)
=
\operatorname{Cov}_{\pi_{A,\epsilon}^\star}(Y,X)+m_\beta m_\alpha^\top,
\]
where $(X,Y)\sim\pi_{A,\epsilon}^\star$, $m_\alpha=\int x\d\alpha(x)$ and $m_\beta=\int y\d\beta(y)$. Thus the gradient is the raw cross-moment of the optimal coupling, and it is its cross-covariance when both marginals are centered.

\subsubsection{Inverse Optimal Transport}
\begin{defn}[Regularized inverse-OT loss]\label{def-inverse-ot-loss}
Let $\epsilon\geq0$ and $\widehat\pi\in\Couplings(\alpha,\beta)$, with $\KL(\widehat\pi|\alpha\otimes\beta)<+\infty$ when $\epsilon>0$. Let $\mathcal C_{\mathrm{adm}}$ be a convex class of real-valued measurable costs such that, for every $c\in\mathcal C_{\mathrm{adm}}$, one has $\int|c|\d\widehat\pi<+\infty$, the forward value $\MK_c^\epsilon(\alpha,\beta)$ is finite, and its infimum is attained. Extending the notation of~\eqref{eq-entropic-generic} to these proper signed costs, define
\eql{\label{eq-inverse-ot-loss}
\mathcal F_\epsilon(c\mid\widehat\pi)
\eqdef
\int c\d\widehat\pi
+\epsilon\KL(\widehat\pi|\alpha\otimes\beta)
-\MK_c^\epsilon(\alpha,\beta),
}
for $c\in\mathcal C_{\mathrm{adm}}$, where the entropy term is set to zero when $\epsilon=0$.
\end{defn}
\begin{prop}[Convexity and calibration of the inverse-OT loss]\label{prop-inverse-ot-convex}
For fixed $(\alpha,\beta,\widehat\pi)$, the map $c\mapsto\mathcal F_\epsilon(c\mid\widehat\pi)$ is convex and nonnegative on $\mathcal C_{\mathrm{adm}}$. Moreover,
\(\mathcal F_\epsilon(c\mid\widehat\pi)=0 \quad\Longleftrightarrow\quad \widehat\pi\text{ solves~\eqref{eq-entropic-generic} for the cost }c.\)
When $\epsilon>0$, the regularized optimizer is unique whenever the value is finite, so the zero set identifies the forward coupling for a fixed cost.
\end{prop}
\begin{proof}
The first two terms in~\eqref{eq-inverse-ot-loss} are respectively linear and constant in $c$, whereas $c\mapsto\MK_c^\epsilon(\alpha,\beta)$ is concave because it is the infimum of affine functions of $c$, as in Proposition~\ref{prop-kantorovich-value-curvature}. Testing~\eqref{eq-entropic-generic} with $\widehat\pi$ proves nonnegativity, and equality holds precisely when this test plan attains the infimum. Uniqueness for $\epsilon>0$ follows from the strict convexity of relative entropy on finite-entropy couplings.
\end{proof}
\paragraph{Bilinear cost learning.}
For concreteness, assume $\alpha,\beta\in\Pp_2(\RR^d)$ and consider the bilinear model
\(c_A(x,y)=\dotp{Ax}{y}, \qquad A\in\RR^{d\times d}.\)
Let $\mathcal A\subset\RR^{d\times d}$ be a convex admissible set and, when $\epsilon>0$, assume $\KL(\widehat\pi|\alpha\otimes\beta)<+\infty$. The moment bound above applies to $\widehat\pi$ and to every feasible coupling, so the bilinear family lies in a proper finite-value domain. After fixing the cost invariances through $\mathcal A$, one estimates
\eql{\label{eq-inverse-ot-bilinear-estimator}
A^\star\in\uargmin{A\in\mathcal A}
\mathcal F_\epsilon(c_A\mid\widehat\pi).
}
\paragraph{Finite-sample polyhedrality and population curvature.}
Let $A_0=-\Id$ and let $\pi_0$ be the associated population optimal plan, equivalently the quadratic $\Wass_2$ plan. From i.i.d.\ pairs $(X_i,Y_i)\sim\pi_0$, form
\(\widehat\pi_n\eqdef\frac1n\sum_{i=1}^n\de_{(X_i,Y_i)}.\)
\paragraph{Statistical estimation.}
Under their nondegenerate-certificate assumptions and with a regularization parameter of order $n^{-1/2}$, up to logarithmic, confidence and $\epsilon$-dependent factors, they prove
\(\norm{A_n^\star-A_0}_{\mathrm F}=O_\PP(n^{-1/2})\)
\subsection{Weak Optimal Transport}
\label{sec-weak-ot}
\paragraph{Barycentric projection of a coupling.}
\begin{defn}[Barycentric projection of a coupling]\label{def-barycentric-projection}
Let $\alpha,\beta\in\Pp_1(\RR^d)$ and let $\pi\in\Couplings(\alpha,\beta)$. Disintegrate $\pi$ with respect to its first marginal as
$\pi(\d x,\d y)=\pi_x(\d y)\alpha(\d x)$. Since $\beta$ has finite first moment, the conditional mean is finite for $\alpha$-a.e.\ $x$, and the barycentric projection of $\pi$ is the map
\begin{equation}\label{eq-barycentric-projection}
\bar T_\pi(x)
\eqdef
\int_{\RR^d}y\d\pi_x(y),
\qquad
\bar\beta_\pi
\eqdef
(\bar T_\pi)_\sharp\alpha.
\end{equation}
\end{defn}
Then
\(\pi_{x_i}=\sum_{j=1}^m\frac{\P_{i,j}}{a_i}\de_{y_j} \qandq \bar T_\pi(x_i)=\frac{1}{a_i}\sum_{j=1}^m\P_{i,j}y_j.\)
\begin{prop}[Barycentric projection of a quadratic optimal plan]\label{prop-barycentric-projection-optimal}
Let $\pi\in\Couplings(\alpha,\beta)$ be optimal for the quadratic cost $\norm{x-y}^2$ between $\alpha,\beta\in\Pp_2(\RR^d)$, and define $\bar T_\pi$ and $\bar\beta_\pi$ by~\eqref{eq-barycentric-projection}. Then $(\Id,\bar T_\pi)_\sharp\alpha$ is an optimal coupling between $\alpha$ and $\bar\beta_\pi$. Equivalently, $\bar T_\pi$ is a quadratic optimal transport map from $\alpha$ to the projected target $\bar\beta_\pi$.
\end{prop}
\begin{proof}
By Theorem~\ref{thm:opt_ccm}, $\pi$ is concentrated on a $c$-cyclically monotone set $\Gamma$ for $c(x,y)=\norm{x-y}^2$. For the quadratic cost, and since it is enough to check cyclic permutations, this means that every finite cycle $(x_i,y_i)_{i=1}^m\subset\Gamma$ satisfies
$\sum_{i=1}^m \dotp{x_i}{y_i} \geq \sum_{i=1}^m \dotp{x_i}{y_{i+1}}$, where $y_{m+1}=y_1$.
After changing the disintegration on an $\alpha$-negligible set, $\pi_x$ is supported on the section
\(\Gamma_x=\enscond{y}{(x,y)\in\Gamma}\)
for $\alpha$-a.e.\ $x$.
Choose $x_1,\ldots,x_m$ in this full-measure set and independently sample $Y_i\sim\pi_{x_i}$. Applying the cyclic inequality to $(x_i,Y_i)$ and taking expectations gives
$\sum_{i=1}^m \dotp{x_i}{\bar T_\pi(x_i)} \geq \sum_{i=1}^m \dotp{x_i}{\bar T_\pi(x_{i+1})}$.
Thus $(\Id,\bar T_\pi)_\sharp\alpha$ is concentrated on a cyclically monotone graph. By the cyclic-monotonicity characterization of quadratic optimality, this plan is optimal between its two marginals, namely $\alpha$ and $\bar\beta_\pi$.
\end{proof}
\paragraph{Weak transport costs.}
\begin{defn}[Weak optimal transport]\label{def-weak-optimal-transport}
Let $\Xx$ and $\Yy$ be Polish spaces, let $\alpha\in\Pp(\Xx)$ and $\beta\in\Pp(\Yy)$, and let $C:\Xx\times\Pp(\Yy)\to\RR\cup\{+\infty\}$ be measurable. For each $\pi\in\Couplings(\alpha,\beta)$, let
$\pi(\d x,\d y)=\alpha(\d x)\pi_x(\d y)$ be a disintegration. Then
\begin{equation}\label{eq-weak-ot}
\WOT_C(\alpha,\beta)
\eqdef
\inf_{\pi\in\Couplings(\alpha,\beta)}
\int C(x,\pi_x)\d\alpha(x).
\end{equation}
\end{defn}
\begin{prop}[Weak Kantorovich duality]\label{prop-weak-ot-duality}
Let $\Xx$ and $\Yy$ be compact metric spaces, and let $C:\Xx\times\Pp(\Yy)\to\RR\cup\{+\infty\}$ be proper, jointly lower semicontinuous, bounded from below and convex in its second argument. Fix $\alpha\in\Pp(\Xx)$ and $\beta\in\Pp(\Yy)$, and assume that $\WOT_C(\alpha,\beta)<+\infty$. For $g\in\Cc(\Yy)$, define the weak $C$-transform
\[
g^C(x)
\eqdef
\inf_{\nu\in\Pp(\Yy)}
\left\{
C(x,\nu)-\int g(y)\d\nu(y)
\right\}.
\]
Then
\[
\WOT_C(\alpha,\beta)
=
\sup_{g\in\Cc(\Yy)}
\left\{
\int g^C(x)\d\alpha(x)+\int g(y)\d\beta(y)
\right\}.
\]
When $C(x,\nu)=\int c(x,y)\d\nu(y)$, this reduces to the usual Kantorovich dual with $g^C(x)=\inf_y(c(x,y)-g(y))$.
\end{prop}
\begin{proof}
For any coupling $\pi$ and any $g\in\Cc(\Yy)$, the definition of $g^C$ gives
\(C(x,\pi_x)\geq g^C(x)+\int g(y)\d\pi_x(y)\).
After integration with respect to $\alpha$, the second term becomes $\int g\d\beta$ because the second marginal of $\pi$ is $\beta$. This proves weak duality.

For the converse, lift a kernel $x\mapsto\pi_x$ to the measure $\alpha(\d x)\delta_{\pi_x}(\d\nu)$ on $\Xx\times\Pp(\Yy)$. Relaxing this graph constraint gives probability measures $P$ whose first marginal is $\alpha$ and whose intensity satisfies
\(\int_{\Xx\times\Pp(\Yy)}\nu\,\d P(x,\nu)=\beta\).
This relaxation does not change the value: disintegrating $P(\d x,\d\nu)=P_x(\d\nu)\alpha(\d x)$ and replacing $P_x$ by its barycenter $\bar\nu_x=\int\nu\,\d P_x(\nu)$ preserves the intensity constraint and cannot increase the objective, by convexity of $C(x,\cdot)$. The relaxed feasible set is compact, and the integral functional is lower semicontinuous. Fenchel--Rockafellar duality for the affine intensity constraint therefore has no gap; its continuous multiplier is $g\in\Cc(\Yy)$. Minimization in $\nu$ gives $g^C(x)$, while the constraint contributes $\int g\d\beta$. This is the weak-cost analogue of Proposition~\ref{prop-duality-discr}; see for the Polish-space formulation.
\end{proof}
\begin{prop}[Barycentric weak transport is weaker than $\Wass_2$]\label{prop-barycentric-weak-ot}
Let $\alpha,\beta\in\Pp_2(\RR^d)$ and define
\(C_{\mathrm{bar}}(x,\nu) = \norm{x-\int y\d\nu(y)}^2.\)
Equivalently,
\begin{equation}\label{eq-barycentric-weak-primal}
\WOT_{C_{\mathrm{bar}}}(\alpha,\beta)
=
\inf_{\pi\in\Couplings(\alpha,\beta)}
\int_{\RR^d}\norm{x-\bar T_\pi(x)}^2\d\alpha(x).
\end{equation}
Then
$\WOT_{C_{\mathrm{bar}}}(\alpha,\beta) \leq \Wass_2^2(\alpha,\beta)$.
\end{prop}
\begin{proof}
Let $\pi$ be any coupling and disintegrate it as $\pi_x\alpha$. By Jensen's inequality,
$\norm{x-\bar T_\pi(x)}^2 \leq \int\norm{x-y}^2\d\pi_x(y)$.
Integrating in $x$ gives $\int C_{\mathrm{bar}}(x,\pi_x)\d\alpha(x)\leq\int\norm{x-y}^2\d\pi(x,y)$. Taking the infimum over $\pi$ proves the claim.
\end{proof}
For discrete measures $\alpha=\sum_{i=1}^n a_i\de_{x_i}$ and $\beta=\sum_{j=1}^m b_j\de_{y_j}$, delete zero-mass source atoms so that $a_i>0$. Formula~\eqref{eq-barycentric-weak-primal} becomes
\begin{equation}\label{eq-barycentric-weak-discrete}
\min_{\P\in\CouplingsD(\a,\b)}
\sum_{i=1}^n a_i
\norm{x_i-\frac{1}{a_i}\sum_{j=1}^m\P_{i,j}y_j}^2.
\end{equation}
Writing its unregularized objective as $F_{\mathrm{bar}}(\pi)$, a linearized ground cost at $\pi$ is
\begin{equation}\label{eq-barycentric-weak-linearized-cost}
c_\pi^{\mathrm{bar}}(x,y)
\eqdef
2\dotp{x-\bar T_\pi(x)}{x-y}.
\end{equation}
Indeed, for every $\widehat\pi\in\Couplings(\alpha,\beta)$,
\[
\left.\frac{\d}{\d t}\right|_{t=0}
F_{\mathrm{bar}}\bigl((1-t)\pi+t\widehat\pi\bigr)
=
\int c_\pi^{\mathrm{bar}}\d(\widehat\pi-\pi).
\]
The term $2\dotp{x-\bar T_\pi(x)}{x}$ in~\eqref{eq-barycentric-weak-linearized-cost} depends only on $x$ and is therefore a harmless source-marginal gauge; equivalently, one may use $2\dotp{\bar T_\pi(x)-x}{y}$. The first-order condition for the KL-regularized convex problem is consequently the self-consistent entropic OT problem
\begin{equation}\label{eq-barycentric-weak-entropic-fixed-point}
\pi^\star
\in
\uargmin{\widehat\pi\in\Couplings(\alpha,\beta)}
\left\{
\int c_{\pi^\star}^{\mathrm{bar}}\d\widehat\pi
+\epsilon\KL(\widehat\pi|\alpha\otimes\beta)
\right\}.
\end{equation}
\subsubsection{Martingale Optimal Transport}
\label{sec-martingale-ot}
\begin{defn}[Martingale couplings and martingale OT]\label{def-martingale-coupling}
Let $\alpha,\beta\in\Pp_1(\RR^d)$. A coupling $\pi\in\Couplings(\alpha,\beta)$ is a martingale coupling if, for the disintegration $\pi(\d x,\d y)=\pi_x(\d y)\alpha(\d x)$,
\begin{equation}\label{eq-martingale-coupling}
\int_{\RR^d}y\d\pi_x(y)=x
\qquad\text{for $\alpha$-a.e.\ }x.
\end{equation}
Equivalently, $\bar T_\pi=\Id$ in~\eqref{eq-barycentric-projection}. The set of such couplings is denoted
\(\Couplings_{\mathrm{mart}}(\alpha,\beta) \eqdef \enscond{\pi\in\Couplings(\alpha,\beta)}{\bar T_\pi=\Id\quad\alpha\text{-a.e.}}.\)
\end{defn}
For a cost $c$, the martingale OT problem minimizes $\int c\d\pi$ over $\Couplings_{\mathrm{mart}}(\alpha,\beta)$, with value $+\infty$ when the admissible set is empty.

The terminology comes from probability: if $(X,Y)\sim\pi$, then~\eqref{eq-martingale-coupling} is exactly $\EE[Y|X]=X$. Hence martingale OT is a Kantorovich problem with the usual two marginal constraints plus a barycentric constraint on the conditional laws. Equivalently, the barycentric projected coupling

\paragraph{Stochastic orders.}
\begin{defn}[Classical stochastic order]\label{def-classical-stochastic-order}
For $\alpha,\beta\in\Pp(\RR)$,
\[
\alpha\preceq_{\mathrm{st}}\beta
\quad\Longleftrightarrow\quad
\int\varphi\,\d\alpha\leq\int\varphi\,\d\beta
\quad\text{for every bounded increasing }\varphi.
\]
\end{defn}
\begin{prop}[Strassen's theorem for stochastic order]\label{prop-strassen-stochastic-order}
Let $\alpha,\beta\in\Pp(\RR)$. Then $\alpha\preceq_{\mathrm{st}}\beta$ if and only if there exists $\pi\in\Couplings(\alpha,\beta)$ concentrated on $\enscond{(x,y)}{x\leq y}$.
\end{prop}
\begin{proof}
A coupling supported on $x\leq y$ immediately gives the integral inequality for every increasing $\varphi$. Conversely, testing approximate indicators of intervals $(t,+\infty)$ gives $F_\alpha(t)\geq F_\beta(t)$ for all continuity points of the distribution functions. If $U$ is uniform on $(0,1)$ and $F_\alpha^{-1},F_\beta^{-1}$ are the generalized quantiles, then $X=F_\alpha^{-1}(U)$ and $Y=F_\beta^{-1}(U)$ have laws $\alpha$ and $\beta$, and the inequality between distribution functions gives $X\leq Y$ almost surely.
\end{proof}
\paragraph{Convex order and martingale feasibility.}
\begin{defn}[Convex order]\label{def-convex-order}
For $\alpha,\beta\in\Pp_1(\RR^d)$, write $\alpha\preceq_{\mathrm{cx}}\beta$ when
$\int \varphi\,\d\alpha\leq\int \varphi\,\d\beta$
for every continuous convex $\varphi:\RR^d\to\RR$ with at most linear growth.
\end{defn}
\begin{thm}[Strassen's martingale theorem]\label{thm-strassen-martingale}
Let $\alpha,\beta\in\Pp_1(\RR^d)$. Then
$\alpha\preceq_{\mathrm{cx}}\beta$ if and only if $\Couplings_{\mathrm{mart}}(\alpha,\beta)\neq\emptyset$.
Equivalently, $\beta$ is above $\alpha$ in convex order if and only if there exists a coupling $\pi(\d x,\d y)=\pi_x(\d y)\alpha(\d x)$ such that $\int y\d\pi_x(y)=x$ for $\alpha$-almost every $x$.
\end{thm}
\begin{proof}
If $\pi$ is a martingale coupling, then Jensen's inequality gives, for every convex $\varphi$,
$\varphi(x) = \varphi\!\left(\int y\d\pi_x(y)\right) \leq \int \varphi(y)\d\pi_x(y)$,
and integration in $x$ gives $\int\varphi\d\alpha\leq\int\varphi\d\beta$.

Conversely, assume $\alpha\preceq_{\mathrm{cx}}\beta$. First suppose that both measures are supported in a compact convex set $K$. Separate $\beta$ from the closed convex set of probability measures on $K$ that are reachable from $\alpha$ by martingale kernels supported in $K$. If $\beta$ were not in this set, a continuous separating function $\psi:K\to\RR$ would satisfy
\[
\int\psi\d\beta
>
\sup
\left\{
\int\psi\d\eta
\;:\;
\eta\in\Pp(K),\
\Couplings_{\mathrm{mart}}(\alpha,\eta)\neq\emptyset
\right\}.
\]
For a fixed $x\in K$, optimizing over probability measures on $K$ with barycenter $x$ gives the concave envelope $\operatorname{conc}_K\psi(x)$; hence the right-hand side is $\int\operatorname{conc}_K\psi\d\alpha$. Since convex order is equivalently the reverse inequality for concave functions,
\(\int\operatorname{conc}_K\psi\d\alpha \geq \int\operatorname{conc}_K\psi\d\beta \geq \int\psi\d\beta,\)
a contradiction. For general measures in $\Pp_1(\RR^d)$, the same separation argument is carried out in the $\Wass_1$ topology, whose continuous test functions have at most linear growth. Compactness is replaced by tightness and uniform integrability of first moments; these properties also make the set of attainable second marginals closed and preserve the martingale constraint under limits. This is the standard extension in Strassen's theorem.
\end{proof}
\begin{prop}[Brenier--Strassen projection formula]\label{prop-brenier-strassen-projection}
Let $\alpha,\beta\in\Pp_2(\RR^d)$ and let $C_{\mathrm{bar}}$ be the barycentric quadratic cost of Proposition~\ref{prop-barycentric-weak-ot}. Then
\begin{equation}\label{eq-brenier-strassen-projection}
\WOT_{C_{\mathrm{bar}}}(\alpha,\beta)
=
\inf_{\eta\in\Pp_2(\RR^d):\,\eta\preceq_{\mathrm{cx}}\beta}
\Wass_2^2(\alpha,\eta).
\end{equation}
\end{prop}
\begin{proof}
Let $(X,Y)$ have law $\pi\in\Couplings(\alpha,\beta)$, set $Z=\EE[Y\mid X]$, and denote its law by $\eta$. Conditional Jensen gives
\(\EE[\varphi(Z)] \leq \EE[\varphi(Y)]\)
for every convex $\varphi$, so $\eta\preceq_{\mathrm{cx}}\beta$. Since $(X,Z)$ couples $\alpha$ and $\eta$, \(\Wass_2^2(\alpha,\eta) \leq \EE\norm{X-Z}^2 = \int C_{\mathrm{bar}}(x,\pi_x)\d\alpha(x).\)
Taking the infimum over $\pi$ proves that the left-hand side of~\eqref{eq-brenier-strassen-projection} is no smaller than the right-hand side.

Conversely, fix $\eta\preceq_{\mathrm{cx}}\beta$. By Theorem~\ref{thm-strassen-martingale}, there is a martingale coupling from $\eta$ to $\beta$. Glue it to an optimal coupling $(X,Z)$ between $\alpha$ and $\eta$, conditionally independently given $Z$. The resulting variables satisfy $\EE[Y\mid Z]=Z$, hence
\(\EE[Y\mid X]=\EE[Z\mid X]\).
Conditional Jensen therefore gives \(\EE\norm{X-\EE[Y\mid X]}^2 = \EE\norm{X-\EE[Z\mid X]}^2 \leq \EE\norm{X-Z}^2 = \Wass_2^2(\alpha,\eta).\)
Taking the infimum over $\eta$ proves the reverse inequality.
\end{proof}
For Gaussian measures with the same mean, convex order reduces to the Loewner order on covariance matrices:
\(\Gaussian(m,\Sigma_0)\preceq_{\mathrm{cx}}\Gaussian(m,\Sigma_1) \quad\Longleftrightarrow\quad \Sigma_1-\Sigma_0\succeq0.\)
Conversely, fix $u\in\RR^d$ and use the convex, at-most-linear functions
\[
\varphi_R(x)
\eqdef
\begin{cases}
\dotp{u}{x-m}^2, & |\dotp{u}{x-m}|\leq R,\\
2R|\dotp{u}{x-m}|-R^2, & |\dotp{u}{x-m}|>R.
\end{cases}
\]

\section{Beyond Comparing Measures}
\label{sec-beyond-comparing-measures}
\subsection{Vector and Matrix-Valued Measures}
\label{sec-vector-matrix-valued-measures}
\paragraph{Positive vector-valued measures.}
\begin{defn}[Positive vector-valued measure]\label{def-positive-vector-valued-measure}
A positive $\RR_+^m$-valued measure on $\X$ is a tuple
\(\al=(\al^1,\ldots,\al^m)\in\Mm_+(\X;\RR_+^m),\)
where each component $\al^k$ is a nonnegative finite measure.
\end{defn}
For the dynamic formulas below, assume that $\X$ is either $\RR^d$, the flat torus, or a bounded convex domain in $\RR^d$ equipped with no-flux boundary conditions. First suppose that the endpoints and the curve have densities.

\begin{defn}[Vector-valued dynamic transport]\label{def-vector-valued-dynamic-transport}
For a nonnegative extended-valued action density $\Phi:\RR_+^m\times(\RR^m)^d\to[0,+\infty]$, define
\begin{equation}\label{eq-vector-valued-bb}
\mathcal W_{\Phi}^2(\al_0,\al_1)
\eqdef
\inf_{u,V}
\int_0^1\!\int_\X
\Phi(u_t(x),V_t(x))\,\d x\,\d t
\end{equation}
subject to the endpoint constraints $u_0\d x=\al_0$, $u_1\d x=\al_1$ and the componentwise continuity equation
\begin{equation}\label{eq-vector-valued-continuity}
\partial_t u_t+\nabla_x\cdot V_t=0,
\qquad
(\nabla_x\cdot V_t)^k=\sum_{\ell=1}^d \partial_{x_\ell} V_{t,\ell}^k.
\end{equation}
\end{defn}
For the linear mobility actions $\Phi_{\mathsf M}$ above, these requirements hold: the action is even and quadratic in $V$, while the range convention and finite-dimensional coercivity make it separating on its effective domain. Their lower-semicontinuous measure extensions therefore define extended distances on the finite-action components with prescribed channel masses.

\paragraph{Positive matrix-valued measures.}
\begin{defn}[Positive matrix-valued measure]\label{def-positive-matrix-valued-measure}
Write $\mathbb S^m$ for real symmetric matrices and $\mathbb S_+^m$ for the positive semidefinite cone. A positive $\mathbb S_+^m$-valued measure is a countably additive map from Borel sets of $\X$ to $\mathbb S_+^m$; we denote the cone of such measures by
\(\mathcal A\in\Mm_+(\X;\mathbb S_+^m).\)
\end{defn}
\begin{defn}[Matrix-valued dynamic transport]\label{def-matrix-valued-dynamic-transport}
\begin{equation}\label{eq-matrix-valued-bb}
\begin{aligned}
\mathcal W_{\mathrm{mat}}^2(\mathcal A_0,\mathcal A_1)
\eqdef
\inf_{A,P}\int_0^1\!\int_\X
\sum_{\ell=1}^d \tr\big(P_{t,\ell}^{\top} A_t^\dagger P_{t,\ell}\big)\d x\d t
\end{aligned}
\end{equation}
subject to $A_0\d x=\mathcal A_0$, $A_1\d x=\mathcal A_1$ and to the matrix-valued continuity equation
\begin{equation}\label{eq-matrix-valued-continuity}
\partial_t A_t+\nabla_x\cdot P_t=0,
\qquad
\nabla_x\cdot P_t=\sum_{\ell=1}^d \partial_{x_\ell}P_{t,\ell}.
\end{equation}
Here $A_t^\dagger$ is the Moore--Penrose inverse, and the integrand is $+\infty$ unless
$\operatorname{Ran}(P_{t,\ell})\subseteq\operatorname{Ran}(A_t)$ for every $\ell$.
\end{defn}
Here $\mathcal W_{\mathrm{mat}}$ is an extended distance. Indeed, the action is invariant under $P\mapsto-P$, quadratic under time rescaling and zero only for $P=0$ on its effective domain.

\begin{prop}[Diagonal matrix subproblem]\label{prop-matrix-diagonal-reduction}
Assume that the endpoints are diagonal in a fixed orthonormal basis, \(\mathcal A_i=\diag(\al_i^1,\ldots,\al_i^m), \qquad i=0,1,\)
and that $\al_0^k(\X)=\al_1^k(\X)=m_k$ for every $k$. If one restricts the admissible curves in~\eqref{eq-matrix-valued-bb} to remain diagonal in that basis,
\(A_t=\diag(u_t^1,\ldots,u_t^m), \qquad P_{t,\ell}=\diag(V_{t,\ell}^1,\ldots,V_{t,\ell}^m),\)
then the value of this restricted matrix problem is
\[
\sum_{k:m_k>0} m_k\,\Wass_2^2\!\left(\frac{\al_0^k}{m_k},\frac{\al_1^k}{m_k}\right),
\]
with zero contribution from zero-mass channels. Thus the commuting matrix submodel is exactly the positive vector-valued Benamou--Brenier model associated with $\mathsf M_{\mathrm{diag}}(u)=\diag(u_1,\ldots,u_m)$.
\end{prop}
\begin{proof}
The continuity equation~\eqref{eq-matrix-valued-continuity} is diagonal entry by diagonal entry and gives $\partial_t u^k+\nabla\cdot V^k=0$. Moreover,
\(\sum_{\ell=1}^d \tr\big(P_{t,\ell}^{\top}A_t^\dagger P_{t,\ell}\big) = \sum_{k=1}^m \frac{|V_t^k|^2}{u_t^k}\)
with the same scalar perspective convention as before. The admissible set and the action are therefore exactly those of the diagonal vector model.
\end{proof}
\subsection{Wasserstein over Wasserstein}
\label{sec-wasserstein-over-wasserstein}
\paragraph{Wasserstein spaces as ground spaces.}
\begin{prop}[Wasserstein spaces as ground spaces]\label{prop-wasserstein-space-polish}
Let $1\leq p<\infty$. If $(\X,\dist)$ is Polish, then $(\Pp_p(\X),\Wass_p)$ is Polish. If $\X$ is compact, then $\Pp_p(\X)=\Pp(\X)$ and this space is compact for $\Wass_p$.
\end{prop}
\begin{proof}
This is a standard structural theorem for Wasserstein spaces. For completeness, extract from a $\Wass_p$-Cauchy sequence a subsequence $(\alpha_k)_k$ such that $\sum_k\Wass_p(\alpha_k,\alpha_{k+1})<\infty$. Gluing optimal couplings of consecutive terms produces random variables $(X_k)_k$ with laws $\alpha_k$ and $\sum_k\norm{\dist(X_k,X_{k+1})}_{L^p}<\infty$. Hence $X_k$ converges almost surely and in $L^p$ to an $\X$-valued random variable; its law is the $\Wass_p$ limit of the subsequence, and the Cauchy property gives convergence of the original sequence. Separability follows by approximating measures with finitely supported measures on a countable dense subset and rational weights. If $\X$ is compact, Prokhorov's theorem gives compactness of $\Pp(\X)$ for weak convergence, and Proposition~\ref{prop-wass-topology-polish} identifies this topology with the $\Wass_p$ topology.
\end{proof}
\paragraph{Random measures and collapsed mixtures.}
Fix $1\leq p<\infty$. We denote elements of $\Pp_p(\X)$ by $\alpha,\beta,\ldots$. Elements of $\Pp_p(\Pp_p(\X))$ are denoted by fraktur letters, for instance $\mathfrak A,\mathfrak B$; they are probability laws over probability measures, or random probability measures. A basic parametric example is obtained from a measurable family $(\alpha_\zeta)_{\zeta\in Z}$ and a probability law $\gamma$ on the parameter space:
\begin{equation}\label{eq-wow-parametric-law}
\mathfrak A=(\zeta\mapsto\alpha_\zeta)_\sharp\gamma.
\end{equation}
\begin{defn}[Collapsed, or barycentric, mixture]\label{def-collapsed-barycentric-mixture}
For $\mathfrak A\in\Pp(\Pp(\X))$, the collapsed, or barycentric, mixture associated with $\mathfrak A$ is the measure $\bar\alpha_{\mathfrak A}$ defined by
\begin{equation}\label{eq-wow-barycentric-mixture}
\int_\X f(x)\d\bar\alpha_{\mathfrak A}(x)
=
\int_{\Pp(\X)}
\left(\int_\X f(x)\d\alpha(x)\right)
\d\mathfrak A(\alpha),
\end{equation}
for bounded continuous $f$.
\end{defn}
Moreover, if $\mathfrak A\in\Pp_p(\Pp_p(\X))$, then $\bar\alpha_{\mathfrak A}\in\Pp_p(\X)$ because
\(\int_\X \dist(x,x_0)^p\,\d\bar\alpha_{\mathfrak A}(x) = \int_{\Pp_p(\X)}\Wass_p(\alpha,\delta_{x_0})^p\,\d\mathfrak A(\alpha).\)
\paragraph{Transport between random measures.}
\begin{defn}[Wasserstein-over-Wasserstein distance]\label{def-wasserstein-over-wasserstein}
Let $1\leq p<+\infty$ and let $\X$ be a Polish metric space. For $\mathfrak A,\mathfrak B\in\Pp_p(\Pp_p(\X))$, define
\begin{equation}\label{eq-wow-distance}
\mathbb W_p^p(\mathfrak A,\mathfrak B)
\eqdef
\inf_{\Pi\in\Couplings(\mathfrak A,\mathfrak B)}
\int_{\Pp_p(\X)\times\Pp_p(\X)}
\Wass_p^p(\alpha,\beta)\d\Pi(\alpha,\beta).
\end{equation}
\end{defn}
For $p=2$, Gaussian mixtures provide a particularly explicit example with two levels of geometry. A mixture $\sum_i a_i\Gaussian(m_i,\Sigma_i)$ can either be viewed as the collapsed measure on $\X$, or as the component law
\(\mathfrak A=\sum_i a_i\delta_{\Gaussian(m_i,\Sigma_i)}\)
Given two component laws
\(\mathfrak A=\sum_i a_i\delta_{\Gaussian(m_i,\Sigma_i)}, \qquad \mathfrak B=\sum_j b_j\delta_{\Gaussian(n_j,\Lambda_j)},\)
the discrete problem induced by~\eqref{eq-wow-distance} uses the component cost
\(\C_{i,j}=\norm{m_i-n_j}^2+\Bb(\Sigma_i,\Lambda_j)^2.\)
Let $\P^\star$ be an optimal coupling between the weights $a$ and $b$. Assume first that all component covariances $\Sigma_i$ and $\Lambda_j$ are positive definite. If $A_{i,j}$ denotes the Brenier linear part~\eqref{eq-bures-map} from $\Sigma_i$ to $\Lambda_j$, then each active component pair is interpolated by
\(m_{ij,t}=(1-t)m_i+t n_j, \qquad \Sigma_{ij,t} = \big((1-t)\Id+tA_{i,j}\big)\Sigma_i \big((1-t)\Id+tA_{i,j}\big)^\top,\)
and collapsing these component geodesics gives \(\bar\alpha_t= \sum_{i,j}\P^\star_{i,j}\Gaussian(m_{ij,t},\Sigma_{ij,t}).\)
\begin{prop}[Collapsing is non-expansive]\label{prop-wow-collapsed-bound}
Let $(\X,\dist)$ be Polish, let $1\leq p<\infty$, and let $\mathfrak A,\mathfrak B\in\Pp_p(\Pp_p(\X))$. For the collapsed mixtures defined by~\eqref{eq-wow-barycentric-mixture},
\(\Wass_p(\bar\alpha_{\mathfrak A},\bar\beta_{\mathfrak B}) \leq \mathbb W_p(\mathfrak A,\mathfrak B).\)
\end{prop}
\begin{proof}
Fix $\Pi\in\Couplings(\mathfrak A,\mathfrak B)$ and $\eta>0$. A standard measurable-selection theorem for optimal transport on Polish spaces gives a measurable kernel $(\alpha,\beta)\mapsto\pi_{\alpha,\beta}\in\Couplings(\alpha,\beta)$ such that
\(\int_{\X\times\X}\dist(x,y)^p\,\d\pi_{\alpha,\beta}(x,y) \leq \Wass_p(\alpha,\beta)^p+\eta.\)
Integrating this kernel against $\Pi$ gives a coupling $\bar\pi$ between $\bar\alpha_{\mathfrak A}$ and $\bar\beta_{\mathfrak B}$. Its cost satisfies
\(\int_{\X\times\X}\dist(x,y)^p\d\bar\pi(x,y) \leq \int_{\Pp_p(\X)^2}\Wass_p^p(\alpha,\beta)\d\Pi(\alpha,\beta)+\eta.\)
Letting $\eta\downarrow0$, taking the infimum over $\Pi$, and then taking the $p$-th root proves the claim.
\end{proof}
\subsection{Gromov--Wasserstein}
\label{sec-gromov-wasserstein}
\paragraph{Discrete formulation.}
Optimal transport needs a ground cost $\C$ to compare histograms $(\a,\b)$, and thus cannot be used directly if the histograms are not defined on the same underlying space, or if one cannot pre-register these spaces to define a ground cost.

To address this issue, one can instead use a weaker requirement: two matrices $\distD \in \RR^{n \times n}$ and $\distD' \in \RR^{m \times m}$ are available and represent relationships between the points on which the histograms are defined. A typical scenario is when these matrices are powers of distance matrices.

Define the quadratic distortion energy
\begin{equation}\label{eq-gw-def}
\Ee_{\distD,\distD'}(\P)
\eqdef
\sum_{i,j,i',j'} \De(\distD_{i,i'},\distD'_{j,j'})^p \P_{i,j}\P_{i',j'},
\qquad
\GWD( (\a,\distD), (\b,\distD') )^p
\eqdef
\min_{\P \in \CouplingsD(\a,\b) }
\Ee_{\distD,\distD'}(\P).
\end{equation}
where $p \geq 1$ and $\De$ is a distance on $\RR$, typically $\De(u,v)=|u-v|$.

When $\distD$ and $\distD'$ satisfy the metric axioms on $\{1,\ldots,n\}$ and $\{1,\ldots,m\}$, respectively, these weighted sets define finite metric-measure spaces with measures $\sum_i\a_i\delta_i$ and $\sum_j\b_j\delta_j$. Theorem~\ref{thm-gw-metric} therefore applies directly to~\eqref{eq-gw-def}: $\GWD$ satisfies the triangle inequality and defines a distance modulo measure-preserving isometries.

\paragraph{General setting.}
\begin{defn}[Metric-measure space]\label{def-metric-measure-space}
A metric-measure space is a triple
$\XX=(\X,\dist_\X,\al)$,
where $(\X,\dist_\X)$ is Polish and $\al$ is a Borel probability measure on $\X$.
\end{defn}
The general setting corresponds to computing couplings between metric-measure spaces $\XX=(\X,\dist_\X,\al)$ and $\YY=(\Y,\dist_\Y,\be)$, where the distance and the measure are both part of the data. The required moment condition depends on how their distances are compared and is therefore included in the next definition.

\begin{defn}[Gromov--Wasserstein distance]\label{def-gromov-wasserstein}
Fix $p\geq1$ and a distance $\De$ on $[0,+\infty)$ such that $(r,s)\mapsto\De(r,s)$ is lower semicontinuous for the usual topology. Define the distortion-dependent size
\(\mathsf S_{\De,p}(\XX)^p \eqdef \int_{\X\times\X} \De\bigl(\dist_\X(x,x'),0\bigr)^p \d\al(x)\d\al(x'),\)
and analogously $\mathsf S_{\De,p}(\YY)$. If both sizes are finite, then
\begin{align}
\label{eq-gw-generic}
\GW( \XX,\YY )^p \eqdef
\umin{ \pi \in \Couplings(\al,\be) }
\int_{\X^2 \times \Y^2}
\De(\dist_\X(x,x'),\dist_\Y(y,y'))^p
\d\pi(x,y)\d\pi(x',y').
\end{align}
\end{defn}
\begin{defn}[Infinite-order Gromov--Wasserstein distance]\label{def-gromov-wasserstein-infinity}
For a lower-semicontinuous distance $\De$ on $[0,+\infty)$, the infinite-order Gromov--Wasserstein distance is
\begin{equation}\label{eq-gw-infinity}
\GW_\infty(\XX,\YY)
\eqdef
\inf_{\pi\in\Couplings(\al,\be)}
\operatorname*{ess\,sup}_{\pi\otimes\pi}
\De\bigl(\dist_\X(x,x'),\dist_\Y(y,y')\bigr).
\end{equation}
\end{defn}
\begin{prop}[Infinite-order limit of Gromov--Wasserstein]\label{prop-gw-infinite-order-limit}
Let $\XX$ and $\YY$ be compact metric-measure spaces, and assume that $\De$ is continuous. Then
\begin{equation}\label{eq-gw-infinity-support}
\GW_\infty(\XX,\YY)
=
\inf_{\pi\in\Couplings(\alpha,\beta)}
\sup_{\bigl((x,y),(x',y')\bigr)\in\supp(\pi)^2}
\De\bigl(\dist_\X(x,x'),\dist_\Y(y,y')\bigr).
\end{equation}
Moreover, $\GW_p(\XX,\YY)$ is nondecreasing in $p$ and
\begin{equation}\label{eq-gw-p-to-infinity}
\lim_{p\to\infty}\GW_p(\XX,\YY)
=
\GW_\infty(\XX,\YY).
\end{equation}
\end{prop}
\begin{proof}
For each coupling $\pi$, continuity of the distortion integrand identifies its essential supremum under $\pi\otimes\pi$ with its supremum over $\supp(\pi)^2$, which proves~\eqref{eq-gw-infinity-support}. Monotonicity follows from monotonicity of probability-space $L^p$ norms, and gives a limit $L\leq\GW_\infty(\XX,\YY)$.

Let $p_k\to\infty$ and let $\pi_k$ minimize the order-$p_k$ problem. Compactness of $\Couplings(\alpha,\beta)$ gives, after extraction, $\pi_k\rightharpoonup\pi_\star$. If the distortion vanishes on $\supp(\pi_\star)^2$, the result is immediate. Otherwise, fix $(x,y),(x',y')\in\supp(\pi_\star)$ with positive distortion and a number
$0<r<\De\bigl(\dist_\X(x,x'),\dist_\Y(y,y')\bigr)$.
By continuity, there are neighborhoods $U$ and $V$ of these points on whose product the distortion exceeds $r$. The Portmanteau theorem gives $\liminf_k\pi_k(U)>0$ and $\liminf_k\pi_k(V)>0$. Hence, for some $c>0$ and all sufficiently large $k$,
\(\GW_{p_k}(\XX,\YY) \geq r\bigl(\pi_k(U)\pi_k(V)\bigr)^{1/p_k} \geq rc^{1/p_k}.\)
Letting first $k\to\infty$ and then $r$ increase to the distortion of the selected pair shows that $L$ dominates every distortion on $\supp(\pi_\star)^2$. Therefore
\(L \geq \sup_{\bigl((x,y),(x',y')\bigr)\in\supp(\pi_\star)^2} \De\bigl(\dist_\X(x,x'),\dist_\Y(y,y')\bigr) \geq \GW_\infty(\XX,\YY),\)
which proves~\eqref{eq-gw-p-to-infinity}.
\end{proof}
Taking $\X=\{1,\ldots,n\}$, $\Y=\{1,\ldots,m\}$, $\al=\sum_i \a_i\delta_i$ and $\be=\sum_j \b_j\delta_j$ reduces~\eqref{eq-gw-generic} exactly to the discrete problem~\eqref{eq-gw-def}, since every coupling is represented by a matrix $\P\in\CouplingsD(\a,\b)$. Thus the construction above is precisely the finite-space, or equivalently finitely supported, specialization of the metric-measure formulation.

\paragraph{Comparison with Wasserstein and Wasserstein--Procrustes.}
\begin{prop}[GW, Wasserstein--Procrustes and Wasserstein]\label{prop-gw-wasserstein-procrustes-comparison}\label{prop-gw-controlled-by-wasserstein}\label{prop-gw-procrustes-upper-certificate}
Let $p\geq1$, let $\alpha,\beta\in\Pp_p(\X)$ be probability measures on the same metric space $(\X,\dist_\X)$, and take $\De(u,v)=|u-v|$ in~\eqref{eq-gw-generic}. Then
\begin{equation}\label{eq-gw-fixed-space-wasserstein}
\GW((\X,\dist_\X,\alpha),(\X,\dist_\X,\beta))
\leq
2\Wass_p(\alpha,\beta),
\end{equation}
where $\Wass_p$ uses the same ground metric $\dist_\X$. In particular, when $\X=\RR^d$ is equipped with the Euclidean distance, write $\XX_\alpha=(\RR^d,\norm{\cdot},\alpha)$ and $\XX_\beta=(\RR^d,\norm{\cdot},\beta)$. Then
\begin{equation}\label{eq-gw-wp-wasserstein-comparison}
\frac12\GW(\XX_\alpha,\XX_\beta)
\leq
\Wass_{p,\mathrm E(d)}([\alpha],[\beta])
\leq
\Wass_p(\alpha,\beta).
\end{equation}
For $p=2$, the middle term is the Wasserstein--Procrustes distance of Definition~\ref{def-wasserstein-procrustes}. The factor $1/2$ reflects the normalization in~\eqref{eq-gw-generic}; it disappears under the alternative convention $\De(u,v)=|u-v|/2$.
\end{prop}
\begin{proof}
Fix $\pi\in\Couplings(\alpha,\beta)$, and draw $(x,y)$ and $(x',y')$ independently from $\pi$. The reverse triangle inequality gives \(\abs{\dist_\X(x,x')-\dist_\X(y,y')} \leq \dist_\X(x,y)+\dist_\X(x',y').\)
Taking the $L^p(\pi\otimes\pi)$ norm, applying Minkowski's inequality, and optimizing over $\pi$ proves~\eqref{eq-gw-fixed-space-wasserstein}.

Now specialize to $\RR^d$. Choosing the identity in Definition~\ref{def-quotient-wasserstein} proves the right inequality in~\eqref{eq-gw-wp-wasserstein-comparison}. Rigid motions preserve pairwise Euclidean distances, so the GW value is unchanged when $\alpha$ and $\beta$ are pushed forward by arbitrary $g,h\in\mathrm E(d)$. Applying~\eqref{eq-gw-fixed-space-wasserstein} to these aligned measures gives
\(\GW(\XX_\alpha,\XX_\beta) \leq 2\Wass_p(g_\sharp\alpha,h_\sharp\beta).\)
Taking the infimum over $g,h$ proves the left inequality.
\end{proof}
\begin{defn}[Isometric metric-measure spaces]\label{def-isometric-mm-spaces}
Two metric-measure spaces $\XX=(\X,\dist_\X,\al)$ and $\YY=(\Y,\dist_\Y,\be)$ are isometric if there exists a bijection $\phi:\supp(\al)\to\supp(\be)$ such that $\phi_\sharp\al=\be$ and
\(\dist_\Y(\phi(x),\phi(x'))=\dist_\X(x,x')\)
for all $x,x'\in\supp(\al)$.
\end{defn}
\begin{thm}[Gromov--Wasserstein metric modulo isometries]\label{thm-gw-metric}
For compact metric-measure spaces, $p\geq1$ and $\De(u,v)=|u-v|$, $\GW$ defines a distance up to measure-preserving isometries.
\end{thm}
\begin{proof}
Compactness of the supports makes the coupling set compact and the distortion functional continuous, so an optimal coupling exists. If $\phi$ is a measure-preserving isometry, the graph coupling $(\Id,\phi)_\sharp\al$ has zero distortion; in particular, $\GW(\XX,\XX)=0$. Symmetry and non-negativity are immediate.

If $\GW(\XX,\YY)=0$ and $\pi$ is an optimal plan, then $\dist_\X(x,x')=\dist_\Y(y,y')$ holds $\pi\otimes\pi$-almost everywhere. By continuity, this equality holds on $\supp(\pi)^2$.
\(\dist_\pi((x,y),(x',y')) \eqdef \frac12\dist_\X(x,x')+\frac12\dist_\Y(y,y').\)
The first projection $\psi:(x,y)\mapsto x$ is measure-preserving. For $((x,y),(x',y'))\in\supp(\pi)^2$, \(\dist_\X(\psi(x,y),\psi(x',y')) = \dist_\X(x,x') = \dist_\Y(y,y') = \dist_\pi((x,y),(x',y')),\)
so $\psi$ is an isometry and therefore injective. To see surjectivity onto $\supp(\al)$, take $x\in\supp(\al)$. Since $\psi_\sharp\pi=\al$, there is a sequence $(x_k,y_k)\in\supp(\pi)$ with $x_k\to x$. The equality of distances on $\supp(\pi)$ makes $(y_k)_k$ Cauchy, and compactness gives a convergent subsequence with limit $(x,y)\in\supp(\pi)$. The same argument for the second projection shows that the support space is also isometric to $\YY$.

For the triangle inequality, let $\pi$ be an optimal coupling between $\XX$ and $\YY$, and $\xi$ an optimal coupling between $\YY$ and $\ZZ=(\Z,\dist_\Z,\ga)$. By the gluing lemma, take $\sigma$ on $\X\times\Y\times\Z$ whose $(\X,\Y)$ and $(\Y,\Z)$ marginals are $\pi$ and $\xi$. Let $\rho=(P_{\X,\Z})_\sharp\sigma$, and write $\bar\sigma=\sigma\otimes\sigma$ for the product law of two independent triples $(x,y,z)$ and $(x',y',z')$. Then $\rho$ is feasible between $\XX$ and $\ZZ$, and
\begin{align*}
\GW(\XX,\ZZ)
&\leq
\left(
\int
\abs{\dist_\X(x,x')-\dist_\Z(z,z')}^p
\d\bar\sigma
\right)^{1/p} \\
&\leq
\left(
\int
\abs{\dist_\X(x,x')-\dist_\Y(y,y')}^p
\d\bar\sigma
\right)^{1/p}
+
\left(
\int
\abs{\dist_\Y(y,y')-\dist_\Z(z,z')}^p
\d\bar\sigma
\right)^{1/p} \\
&=
\GW(\XX,\YY)+\GW(\YY,\ZZ),
\end{align*}
where the second inequality uses the pointwise triangle inequality followed by Minkowski's inequality. This proves the triangle inequality and completes the metric proof.
\end{proof}
\paragraph{Marginal stability and geodesics.}
\begin{prop}[Marginal stability and empirical GW rates]\label{prop-gw-empirical-stability}
Let $\XX=(\X,\dist_\X,\alpha)$ and $\YY=(\Y,\dist_\Y,\beta)$ be compact metric-measure spaces, with $\De(u,v)=|u-v|$ and exponent $p\geq1$. For any probability measures $\tilde\alpha$ on $\X$ and $\tilde\beta$ on $\Y$, set
$\widetilde\XX=(\X,\dist_\X,\tilde\alpha)$ and $\widetilde\YY=(\Y,\dist_\Y,\tilde\beta)$.
Then, deterministically, \(\abs{\GW(\widetilde\XX,\widetilde\YY)-\GW(\XX,\YY)} \leq 2\Wass_p^\X(\tilde\alpha,\alpha) + 2\Wass_p^\Y(\tilde\beta,\beta),\)
where $\Wass_p^\X$ and $\Wass_p^\Y$ denote Wasserstein distances computed with $\dist_\X$ and $\dist_\Y$. If $\hat\alpha_n$ and $\hat\beta_m$ are independent empirical measures built from iid samples with laws $\alpha$ and $\beta$, write $\widehat\XX_n=(\X,\dist_\X,\hat\alpha_n)$ and $\widehat\YY_m=(\Y,\dist_\Y,\hat\beta_m)$.
If, in addition, $\X$ and $\Y$ are regular compact subsets of Euclidean spaces of intrinsic dimensions $d_\X$ and $d_\Y$, and the high-dimensional Wasserstein regime $d_\X,d_\Y>2p$ applies, then
\(\EE\abs{\GW(\widehat\XX_n,\widehat\YY_m)-\GW(\XX,\YY)} = O(n^{-1/d_\X}+m^{-1/d_\Y}).\)
For $p=1$ on subsets of $[0,1]^d$, the rates, including the low-dimensional logarithmic cases, are those of Proposition~\ref{prop-empirical-ot-rate}.
\end{prop}
\begin{proof}
The triangle inequality of Theorem~\ref{thm-gw-metric} gives \(\GW(\widetilde\XX,\widetilde\YY) - \GW(\XX,\YY) \leq \GW(\widetilde\XX,\XX) + \GW(\YY,\widetilde\YY).\)
Exchanging the two pairs gives the same bound for the opposite difference, hence the absolute value. The two terms on the right compare metric-measure spaces with the same underlying metric and only different measures, so Proposition~\ref{prop-gw-controlled-by-wasserstein} bounds them by $2\Wass_p^\X(\tilde\alpha,\alpha)$ and $2\Wass_p^\Y(\tilde\beta,\beta)$. Substituting $\tilde\alpha=\hat\alpha_n$ and $\tilde\beta=\hat\beta_m$, taking expectations and applying the empirical Wasserstein rates recalled in Section~\ref{sec-sample-complexity} gives the stated Euclidean rate.
\end{proof}
\begin{prop}[Gromov--Wasserstein geodesics]\label{prop-gw-geodesics}
Let $\XX_0=(\X_0,\dist_{\X_0},\al_0)$ and $\XX_1=(\X_1,\dist_{\X_1},\al_1)$ be compact metric-measure spaces, take $\De(u,v)=|u-v|$ in~\eqref{eq-gw-generic}, and let $\pi^\star$ be an optimal coupling. Define, on $\Z=\X_0\times\X_1$,
\(\dist_t((x_0,x_1),(x'_0,x'_1)) \eqdef (1-t)\dist_{\X_0}(x_0,x'_0)+t\dist_{\X_1}(x_1,x'_1), \qquad \XX_t=(\Z,\dist_t,\pi^\star).\)
For $0<t<1$, $\dist_t$ is a metric. At $t=0$ and $t=1$, one quotients $\Z$ by the zero-distance relation associated with $\dist_t$. Then $t\mapsto\XX_t$ is a constant-speed geodesic:
\(\GW(\XX_s,\XX_t)=|t-s|\GW(\XX_0,\XX_1) \qquad \forall s,t\in[0,1].\)
\end{prop}
\begin{proof}
Write $D=\GW(\XX_0,\XX_1)$. For $s<t$, couple $\XX_s$ and $\XX_t$ by the diagonal coupling induced by the identity on $\Z$ and the measure $\pi^\star$. For two independent points $z=(x_0,x_1)$ and $z'=(x'_0,x'_1)$ sampled from $\pi^\star$,
\(\dist_t(z,z')-\dist_s(z,z') = (t-s)\big(\dist_{\X_1}(x_1,x'_1)-\dist_{\X_0}(x_0,x'_0)\big).\)
Using this feasible coupling gives $\GW(\XX_s,\XX_t)\leq(t-s)D$. The same construction with the projections from $\Z$ to $\X_0$ and $\X_1$ gives $\GW(\XX_0,\XX_t)\leq tD$ and $\GW(\XX_t,\XX_1)\leq(1-t)D$. The triangle inequality for $\GW$ then yields, for $0\leq s\leq t\leq1$,
\(D \leq \GW(\XX_0,\XX_s)+\GW(\XX_s,\XX_t)+\GW(\XX_t,\XX_1) \leq sD+\GW(\XX_s,\XX_t)+(1-t)D,\)
so $\GW(\XX_s,\XX_t)\geq(t-s)D$. This proves equality.
\end{proof}
Its generator uses squared Euclidean dissimilarities, normalizes each space by its own maximum squared distance $M_\X$ and $M_\Y$, and displays
\[
r_{i,i'}
=
\left|
\frac{\dist_\X(x_i,x_{i'})^2}{M_\X}
-
\frac{\dist_\Y(y_{\sigma(i)},y_{\sigma(i')})^2}{M_\Y}
\right|.
\]
Thus $r_{i,i'}^2$ is the local square-loss integrand optimized by GW, whereas the displayed value $r_{i,i'}$ is its absolute-residual scale.

\paragraph{Distance-profile lower bound.}
Local distance distributions provide an inexpensive intrinsic certificate and a useful initialization principle for the non-convex solver. The construction maps each point to its distance-profile distribution and then pushes forward the ambient measure, thereby producing a distribution over distributions.

\begin{prop}[M\'emoli profile lower bound]\label{prop-memoli-gw-profile-lower-bound}
Let $\XX=(\X,\dist_\X,\alpha)$ and $\YY=(\Y,\dist_\Y,\beta)$ be compact metric-measure spaces and take $\De(u,v)=|u-v|$ in~\eqref{eq-gw-generic}, with the same exponent $p\geq1$. For each $x\in\X$ and $y\in\Y$, define the distance-profile measures on $\RR_+$ by
\(\alpha_x\eqdef(\dist_\X(x,\cdot))_\sharp\alpha, \qquad \beta_y\eqdef(\dist_\Y(y,\cdot))_\sharp\beta.\)
Let $\mathfrak D_\XX=(x\mapsto\alpha_x)_\sharp\alpha$ and $\mathfrak D_\YY=(y\mapsto\beta_y)_\sharp\beta$, which are probability measures on $\Pp_p(\RR_+)$. Then
\(\mathbb W_p\bigl(\mathfrak D_\XX,\mathfrak D_\YY\bigr) \leq \GW(\XX,\YY).\)
Here $\mathbb W_p$ is the Wasserstein-over-Wasserstein distance~\eqref{eq-wow-distance}: its ground space is $(\Pp_p(\RR_+),\Wass_p)$.
\end{prop}
\begin{proof}
Fix any $\pi\in\Couplings(\alpha,\beta)$. It induces a coupling $(x,y)\mapsto(\alpha_x,\beta_y)$ between $\mathfrak D_\XX$ and $\mathfrak D_\YY$, hence
\(\mathbb W_p\bigl(\mathfrak D_\XX,\mathfrak D_\YY\bigr)^p \leq \int_{\X\times\Y}\Wass_p(\alpha_x,\beta_y)^p\,\d\pi(x,y).\)
For fixed $(x,y)$, the map $(x',y')\mapsto(\dist_\X(x,x'),\dist_\Y(y,y'))$ pushes the same coupling $\pi$ to a coupling between $\alpha_x$ and $\beta_y$. Therefore
\(\Wass_p(\alpha_x,\beta_y)^p \leq \int_{\X\times\Y} \abs{\dist_\X(x,x')-\dist_\Y(y,y')}^p \d\pi(x',y').\)
Integrating in $(x,y)$ gives \(\mathbb W_p\bigl(\mathfrak D_\XX,\mathfrak D_\YY\bigr)^p \leq \int_{\X^2\times\Y^2} \abs{\dist_\X(x,x')-\dist_\Y(y,y')}^p \d\pi(x,y)\d\pi(x',y').\)
Taking the infimum over $\pi$ and then the $p$-th root proves the claim.
\end{proof}
Combining the two propositions gives the extended Euclidean sandwich
\begin{equation}\label{eq-gw-profile-procrustes-sandwich}
\mathbb W_p(\mathfrak D_\XX,\mathfrak D_\YY)
\leq
\GW(\XX,\YY)
\leq
2\,\Wass_{p,\mathrm E(d)}([\alpha],[\beta])
\leq
2\,\Wass_p(\alpha,\beta),
\end{equation}
where $\XX=(\RR^d,\norm{\cdot},\alpha)$ and $\YY=(\RR^d,\norm{\cdot},\beta)$. The profile term is an inexpensive intrinsic certificate, while the Wasserstein--Procrustes term gives an extrinsic certificate after optimal rigid registration.

\paragraph{First-order optimality condition as self-consistent OT.}
To isolate the relevant curvature, set $\mathcal Z=\X\times\Y$ and
\(L\bigl((x,y),(x',y')\bigr) \eqdef \De\bigl(\dist_\X(x,x'),\dist_\Y(y,y')\bigr)^p.\)
For finite signed measures $\xi,\eta\in\Mm(\mathcal Z)$ for which the integrals are absolutely convergent, define the symmetric bilinear form and its associated quadratic form by
\begin{equation}\label{eq-gw-quadratic-form}
\mathcal B_L(\xi,\eta)
\eqdef
\iint_{\mathcal Z^2}L(z,z')\d\xi(z)\d\eta(z'),
\qquad
\mathcal Q_L(\xi)
\eqdef\mathcal B_L(\xi,\xi).
\end{equation}
Thus~\eqref{eq-gw-generic} minimizes $\mathcal Q_L(\pi)$ over $\pi\in\Couplings(\alpha,\beta)$. At a coupling $\pi$, define the linearized ground cost
\begin{equation}\label{eq-gw-linearized-cost}
c_\pi(z)
\eqdef
2\int_{\mathcal Z}L(z,z')\d\pi(z').
\end{equation}
For any $\widehat\pi\in\Couplings(\alpha,\beta)$ and $\xi=\widehat\pi-\pi$, the exact quadratic expansion is
\begin{equation}\label{eq-gw-exact-linearization}
\mathcal Q_L(\pi+t\xi)
=
\mathcal Q_L(\pi)
{}+t\int_{\mathcal Z}c_\pi\d\xi
{}+t^2\mathcal Q_L(\xi).
\end{equation}
\begin{defn}[First-order stationary GW coupling]\label{def-gw-stationary-coupling}
A coupling $\pi\in\Couplings(\alpha,\beta)$ is \emph{first-order stationary for GW} if it satisfies the variational inequality
\begin{equation}\label{eq-gw-stationary-variational-inequality}
\int_{\mathcal Z}c_\pi\d(\widehat\pi-\pi)\geq0
\qquad\text{for every }\widehat\pi\in\Couplings(\alpha,\beta).
\end{equation}
\end{defn}
\begin{prop}[Frozen-OT characterization of stationarity]\label{prop-gw-frozen-ot-stationarity}
A coupling $\pi$ is first-order stationary in the sense of Definition~\ref{def-gw-stationary-coupling} if and only if
\begin{equation}\label{eq-gw-frozen-ot-fixed-point}
\pi\in\argmin_{\widehat\pi\in\Couplings(\alpha,\beta)}
\int_{\mathcal Z}c_\pi\d\widehat\pi.
\end{equation}
In particular, every local minimizer for the narrow topology, and hence every global minimizer, of $\mathcal Q_L$ is an optimal Kantorovich coupling for its own linearized cost $c_\pi$.
\end{prop}
\begin{proof}
The variational inequality~\eqref{eq-gw-stationary-variational-inequality} is exactly the assertion that $\pi$ minimizes the linear functional $\widehat\pi\mapsto\int c_\pi\d\widehat\pi$, proving the equivalence. If $\pi$ is a local minimizer and $\widehat\pi$ is arbitrary, the segment $\pi_t=(1-t)\pi+t\widehat\pi$ remains feasible. The right derivative at $t=0$ in~\eqref{eq-gw-exact-linearization} must be nonnegative, which gives~\eqref{eq-gw-stationary-variational-inequality}.
\end{proof}
\paragraph{Conditional concavity of GW in Euclidean spaces.}
To make these statements precise, define the space of signed zero-marginal directions
\begin{equation}\label{eq-gw-zero-marginal-space}
\mathcal T_{\X,\Y}
\eqdef
\enscond{\xi\in\Mm(\X\times\Y)}{(P_\X)_\sharp\xi=0,\ (P_\Y)_\sharp\xi=0}.
\end{equation}
For finite spaces, identifying $\xi$ with an $n\times m$ matrix $\Xi$ gives \(\mathcal T_{\X,\Y} = \enscond{\Xi\in\RR^{n\times m}}{\Xi\ones_m=0,\ \transp{\Xi}\ones_n=0}, \qquad \dim\mathcal T_{\X,\Y}=(m-1)(n-1).\)
\begin{prop}[Conditional concavity of squared GW]\label{prop-gw-conditional-concavity}
Let $\X$ and $\Y$ be compact metric spaces, and let $\alpha\in\Mm_+^1(\X)$ and $\beta\in\Mm_+^1(\Y)$. In~\eqref{eq-gw-generic}, take $p=2$ and $\De(u,v)=|u-v|$. Assume that the two distance kernels $\dist_\X$ and $\dist_\Y$ are either both conditionally positive definite or both conditionally negative definite in the sense of Definition~\ref{def-positive-kernels}. Then
\begin{equation}\label{eq-gw-conditional-negative-form}
\mathcal Q_L(\xi)\leq0
\qquad\text{for every }\xi\in\mathcal T_{\X,\Y}.
\end{equation}
\end{prop}
\begin{proof}
The zero marginal conditions make the two pure terms in the squared loss vanish, and therefore
\(\mathcal Q_L(\xi) = -2\iint \dist_\X(x,x')\dist_\Y(y,y')\d\xi(x,y)\d\xi(x',y').\)
Choose $s\in\{-1,1\}$ so that $s\dist_\X$ and $s\dist_\Y$ are conditionally positive definite. Fix base points $x_0\in\X$ and $y_0\in\Y$, and define the centered kernels
\begin{align*}
\widetilde{\dist}_\X(x,x')
&=s\bigl(\dist_\X(x,x')-\dist_\X(x,x_0)-\dist_\X(x_0,x')+\dist_\X(x_0,x_0)\bigr),\\
\widetilde{\dist}_\Y(y,y')
&=s\bigl(\dist_\Y(y,y')-\dist_\Y(y,y_0)-\dist_\Y(y_0,y')+\dist_\Y(y_0,y_0)\bigr).
\end{align*}
These kernels are positive definite. Since $\xi$ has zero marginals and $s^2=1$, every centering term vanishes after integration. Let $\Phi_\X$ and $\Phi_\Y$ be Hilbert feature maps for $\widetilde{\dist}_\X$ and $\widetilde{\dist}_\Y$. Compactness and continuity make these maps bounded. Writing $z=(x,y)$ and $z'=(x',y')$, one obtains
\(\iint_{\mathcal Z^2}\dist_\X(x,x')\dist_\Y(y,y')\d\xi(z)\d\xi(z') = \norm{\int_{\mathcal Z}\Phi_\X(x)\otimes\Phi_\Y(y)\d\xi(z)}^2 \geq0.\)
This proves~\eqref{eq-gw-conditional-negative-form}; the second-derivative identity above then proves concavity on every feasible segment.
\end{proof}
\begin{prop}[Extreme-point and permutation optimizers in the concave regime]\label{prop-gw-extreme-optimal-coupling}
Assume the hypotheses of Proposition~\ref{prop-gw-conditional-concavity}. Then the squared GW problem has an optimal coupling that is an extreme point of the convex set $\Couplings(\alpha,\beta)$. If, in addition, $\alpha$ and $\beta$ are supported on respectively $n$ and $m$ points, such an optimizer can be chosen with at most $n+m-1$ nonzero entries. If $n=m$ and both measures are uniform, it can be chosen as a permutation coupling $\P_\sigma/n$.
\end{prop}
\begin{proof}
Set $K=\Couplings(\alpha,\beta)$. Because $\X$ and $\Y$ are compact metric spaces, $K$ is weakly compact, metrizable and convex in the locally convex space of finite signed measures on $\X\times\Y$. The kernel $L$ is continuous and bounded, so $\mathcal Q_L$ is weakly continuous and admits a minimizer $\pi_\star\in K$. The Choquet representation theorem for metrizable compact convex sets provides a probability measure $\nu$ concentrated on the extreme points of $K$ whose barycenter is $\pi_\star$. Concavity of $\mathcal Q_L$ and minimality of $\pi_\star$ give
\(\mathcal Q_L(\pi_\star) \geq \int_K\mathcal Q_L(\eta)\d\nu(\eta) \geq \mathcal Q_L(\pi_\star).\)
Both inequalities are therefore equalities. Since $\mathcal Q_L(\eta)\geq\mathcal Q_L(\pi_\star)$ throughout $K$, equality of the integral implies that $\mathcal Q_L(\eta)=\mathcal Q_L(\pi_\star)$ for $\nu$-almost every $\eta$. Because $\nu$ is concentrated on extreme couplings, at least one extreme coupling is optimal. This proves directly the instance of the Bauer maximum principle needed here.

For finitely supported measures, the coupling set is a transport polytope. Its marginal operator has rank $n+m-1$, as computed in the proof of Proposition~\ref{prop-sparse-optimal-plans}. If an extreme coupling had more than $n+m-1$ positive entries, the restriction of this operator to its support would have a nonzero kernel direction $H$; for sufficiently small $t>0$, both $\P+tH$ and $\P-tH$ would be distinct feasible couplings with midpoint $\P$, contradicting extremality. This proves the support bound. In the equal-size uniform case, Theorem~\ref{thm-birkhoff-von-neumann} identifies the extreme points of the scaled Birkhoff polytope as $\P_\sigma/n$.
\end{proof}
\begin{prop}[Tangent majorization and frozen-OT replacement]\label{prop-gw-frozen-ot-replacement}
Assume that $\mathcal Q_L$ is concave on $\Couplings(\alpha,\beta)$. Then, for every $\pi,\widehat\pi\in\Couplings(\alpha,\beta)$,
\begin{equation}\label{eq-gw-tangent-majorization}
\mathcal Q_L(\widehat\pi)
\leq
\mathcal Q_L(\pi)
+
\int_{\mathcal Z}c_\pi\d(\widehat\pi-\pi)
=
2\mathcal B_L(\pi,\widehat\pi)-\mathcal Q_L(\pi).
\end{equation}
\begin{enumerate}
\item if $\pi^+\in\argmin_{\eta\in\Couplings(\alpha,\beta)}\int c_\pi\d\eta$, then $\mathcal Q_L(\pi^+)\leq\mathcal Q_L(\pi)$;
\item if $\pi$ is a local minimizer of $\mathcal Q_L$ and $\pi^+$ minimizes its frozen OT problem, then
\(\mathcal Q_L\bigl((1-t)\pi+t\pi^+\bigr) = \mathcal Q_L(\pi) \qquad\text{for every }t\in[0,1];\)
\item if $\pi_\star$ is globally GW-optimal, then every minimizer of the frozen OT problem with cost $c_{\pi_\star}$ is also globally GW-optimal.
\end{enumerate}
\end{prop}
\begin{proof}
Apply~\eqref{eq-gw-exact-linearization} at $t=1$ with $\xi=\widehat\pi-\pi$. Concavity gives $\mathcal Q_L(\xi)\leq0$, which proves~\eqref{eq-gw-tangent-majorization}. For the first conclusion, frozen optimality gives $\int c_\pi\d(\pi^+-\pi)\leq0$.

Now assume that $\pi$ is a local minimizer. Proposition~\ref{prop-gw-frozen-ot-stationarity} shows that $\pi$ itself minimizes its frozen objective. Hence every other frozen minimizer $\pi^+$ satisfies
$\int c_\pi\d(\pi^+-\pi)=0$. With $\xi=\pi^+-\pi$, the exact expansion~\eqref{eq-gw-exact-linearization} becomes
$\mathcal Q_L(\pi+t\xi) = \mathcal Q_L(\pi)+t^2\mathcal Q_L(\xi)$.
Local minimality for sufficiently small $t>0$ gives $\mathcal Q_L(\xi)\geq0$, while concavity gives the reverse inequality. Thus $\mathcal Q_L(\xi)=0$, proving the second conclusion for every $t\in[0,1]$. If $\pi=\pi_\star$ is globally optimal, the common value is the global minimum, which proves the third conclusion.
\end{proof}
Thus freezing one occurrence of the coupling produces a global upper bound, not merely a formal linearization. The second conclusion is a flat-replacement theorem for a coupling already known to be locally minimizing; it is not a converse to Proposition~\ref{prop-gw-frozen-ot-stationarity}.

\paragraph{GW as cost-regularized robust Wasserstein transport.}
\begin{prop}[Concave coupling energies as cost-regularized OT]\label{prop-concave-coupling-cost-regularized-ot}
Let $\X$ and $\Y$ be compact, set $\mathcal Z=\X\times\Y$, and let $Q:\Couplings(\alpha,\beta)\to\RR\cup\{-\infty\}$ be proper, upper semicontinuous and concave. Let $\mathfrak F_Q$ be any proper convex weak-* lower-semicontinuous extension of $-Q$ from $\Couplings(\alpha,\beta)$ to finite signed measures on $\mathcal Z$. Such an extension always exists: one possible choice is the canonical constrained extension
\[
\mathfrak F_Q^{\mathrm{can}}(\xi)
\eqdef
\begin{cases}
-Q(\xi), & \xi\in\Couplings(\alpha,\beta),\\
+\infty, & \text{otherwise},
\end{cases}.
\]
Define the convex conjugate of the chosen extension by
\[
\mathfrak F_Q^*(h)
\eqdef
\sup_{\xi\in\Mm(\mathcal Z)}
\left\{\int_{\mathcal Z}h\d\xi-\mathfrak F_Q(\xi)\right\},
\qquad h\in\Cc(\mathcal Z).
\]
Then the convex cost penalty $\mathfrak R_Q(c)\eqdef\mathfrak F_Q^*(-c)$ satisfies
\begin{equation}\label{eq-concave-coupling-cost-regularized-ot}
\inf_{\pi\in\Couplings(\alpha,\beta)}Q(\pi)
=
\inf_{c\in\Cc(\X\times\Y)}
\left\{\MK_c(\alpha,\beta)+\mathfrak R_Q(c)\right\}.
\end{equation}
\end{prop}
\begin{proof}
Fenchel--Moreau gives, for every feasible $\pi$,
$Q(\pi) = \inf_{c\in\Cc(\mathcal Z)} \left\{\int_{\mathcal Z}c\d\pi+\mathfrak F_Q^*(-c)\right\}$.
Taking the infimum in $\pi$ and exchanging the two infima yields~\eqref{eq-concave-coupling-cost-regularized-ot}; the inner infimum in $\pi$ is exactly $\MK_c(\alpha,\beta)$.
\end{proof}
After collecting the terms affine in $\pi$ into a fixed cost $c_0$, one can write
\[
\mathcal Q_L(\pi)
=
C(\alpha,\beta)
+
\int_{\mathcal Z}c_0\d\pi
-
2\norm{m_\pi}_{\RKHS_{\mathcal Z}}^2,
\qquad
m_\pi\eqdef\int_{\mathcal Z}\Phi(z)\d\pi(z).
\]
Completing the square therefore gives the equivalent learned-cost representation
\begin{equation}\label{eq-gw-rkhs-cost-regularization}
\mathcal Q_L(\pi)
=
C(\alpha,\beta)
+
\inf_{c\in c_0+\RKHS_{\mathcal Z}}
\left\{
\int_{\mathcal Z}c\d\pi
+
\frac18\norm{c-c_0}_{\RKHS_{\mathcal Z}}^2
\right\}.
\end{equation}
Thus, after quotienting out separable costs that integrate to constants on $\Couplings(\alpha,\beta)$, the quadratic extension in Proposition~\ref{prop-concave-coupling-cost-regularized-ot} makes $\mathfrak R_Q$, up to the additive marginal constant, a squared RKHS norm, with value $+\infty$ outside the admissible affine cost space.

\begin{prop}[Finite-dimensional cost learning for Euclidean GW]\label{prop-gw-euclidean-cost-regularization}
Let $\X\subset\RR^{d_x}$ and $\Y\subset\RR^{d_y}$ be compact, and use $p=2$ and $\De(u,v)=|u-v|$ throughout. Retain the notation $\GW$ of~\eqref{eq-gw-generic} when the within-space functions are symmetric dissimilarities rather than metrics.
For each of the two choices below, set $\XX=(\X,\dist_\X,\alpha)$ and $\YY=(\Y,\dist_\Y,\beta)$. There is a constant $C(\alpha,\beta)$, depending only on the marginals and on the chosen dissimilarities, such that
\begin{equation}\label{eq-gw-euclidean-cost-learning}
\GW(\XX,\YY)^2
=
C(\alpha,\beta)
+
\inf_{M\in\RR^{d_y\times d_x}}
\left\{\MK_{c_M}(\alpha,\beta)+\frac12\norm{M}_{\mathrm F}^2\right\}.
\end{equation}
The constant $C(\alpha,\beta)$ is different in the two cases, while the learned cost is as follows.
\begin{enumerate}
\item For the inner-product dissimilarities
\begin{equation}\label{eq-gw-inner-product-dissimilarities}
\dist_\X(x,x')=-\dotp{x}{x'},
\qquad
\dist_\Y(y,y')=-\dotp{y}{y'},
\end{equation}
the learned costs form the linear family
$c_M(x,y)=-2\dotp{Mx}{y}$.
\item For the squared-distance dissimilarities
\begin{equation}\label{eq-gw-squared-distance-dissimilarities}
\dist_\X(x,x')=\norm{x-x'}^2,
\qquad
\dist_\Y(y,y')=\norm{y-y'}^2,
\end{equation}
translate $\alpha$ and $\beta$ independently so that both have zero mean, which leaves $\GW$ unchanged. The learned costs then form the affine family
$c_M(x,y)=-4\norm{x}^2\norm{y}^2-4\dotp{Mx}{y}$.
\end{enumerate}
\end{prop}
\begin{proof}
For $\pi\in\Couplings(\alpha,\beta)$, write $S_\pi\eqdef\int yx^\top\d\pi(x,y)$.
For two independent pairs $(x,y),(x',y')\sim\pi$, \(\iint\dotp{x}{x'}\dotp{y}{y'}\d\pi\d\pi = \norm{S_\pi}_{\mathrm F}^2.\)
For~\eqref{eq-gw-inner-product-dissimilarities}, expanding the square therefore gives $\mathcal Q_L(\pi)=C(\alpha,\beta)-2\norm{S_\pi}_{\mathrm F}^2$. For fixed $\pi$, completing the square gives
\[
\inf_M
\left\{-2\dotp{M}{S_\pi}_{\mathrm F}+\frac12\norm{M}_{\mathrm F}^2\right\}
=-2\norm{S_\pi}_{\mathrm F}^2,
\qquad M=2S_\pi,
\]
which proves~\eqref{eq-gw-euclidean-cost-learning} in the first case after minimizing over $\pi$.

For centered marginals, direct expansion yields
\[
\iint\norm{x-x'}^2\norm{y-y'}^2\d\pi\d\pi
=
2\int\norm{x}^2\norm{y}^2\d\pi
+2
\left(\int_\X\norm{x}^2\d\alpha(x)\right)
\left(\int_\Y\norm{y}^2\d\beta(y)\right)
+4\norm{S_\pi}_{\mathrm F}^2.
\]
Substitution into the squared-distance plan objective gives
\begin{equation}\label{eq-gw-squared-distance-cost-regularization}
\mathcal Q_L(\pi)
=
C(\alpha,\beta)
-4\int\norm{x}^2\norm{y}^2\d\pi(x,y)
-8\norm{S_\pi}_{\mathrm F}^2.
\end{equation}
Finally,
\[
\inf_M
\left\{-4\dotp{M}{S_\pi}_{\mathrm F}+\frac12\norm{M}_{\mathrm F}^2\right\}
=-8\norm{S_\pi}_{\mathrm F}^2,
\qquad M=4S_\pi,
\]
which proves~\eqref{eq-gw-euclidean-cost-learning} in the second case.
\end{proof}
\paragraph{GW between Gaussian measures.}
Let
\(\alpha=\Gaussian(m_\alpha,\Sigma_\alpha)\in\Pp_2(\RR^{d_x}), \qquad \beta=\Gaussian(m_\beta,\Sigma_\beta)\in\Pp_2(\RR^{d_y}),\)
\begin{prop}[Exact inner-product and Gaussian-constrained squared GW]\label{prop-gw-between-gaussians}
For the inner-product dissimilarities~\eqref{eq-gw-inner-product-dissimilarities}, the unrestricted problem admits a jointly Gaussian minimizer for arbitrary $m_\alpha$ and $m_\beta$. If both means vanish, then
\begin{equation}\label{eq-centered-gaussian-inner-product-gw}
\GW(\XX,\YY)^2
=
\sum_{i=1}^{D}(\lambda_i-\eta_i)^2.
\end{equation}
For the squared-distance dissimilarities~\eqref{eq-gw-squared-distance-dissimilarities}, define the Gaussian-constrained squared distortion
\begin{equation}\label{eq-gaussian-constrained-squared-gw-definition}
\GW^{\mathrm{Gau}}(\XX,\YY)^2
\eqdef
\inf\left\{\mathcal Q_L(\pi):
\pi\in\Couplings(\alpha,\beta),
\pi\text{ is jointly Gaussian}\right\}.
\end{equation}
Then
\begin{equation}\label{eq-gaussian-constrained-squared-gw}
\GW^{\mathrm{Gau}}(\XX,\YY)^2
=
4\left(\sum_{i=1}^{D}\lambda_i-\sum_{i=1}^{D}\eta_i\right)^2
+8\sum_{i=1}^{D}(\lambda_i-\eta_i)^2.
\end{equation}
If $d_x\geq d_y$ and $\Sigma_\alpha\succ0$, a minimizer in~\eqref{eq-gaussian-constrained-squared-gw-definition} is induced by a deterministic affine map that aligns the principal covariance directions, rescales the first $d_y$ coordinates and discards the remaining ones.
\end{prop}
\begin{proof}
For a coupling $\pi$, set
\(K_\pi \eqdef \int (y-m_\beta)(x-m_\alpha)^\top\d\pi(x,y).\)
Its block covariance is positive semidefinite,
$\left(\begin{smallmatrix}\Sigma_\alpha&K_\pi^\top\\K_\pi&\Sigma_\beta\end{smallmatrix}\right)\succeq0$.
Conversely, every matrix $K$ satisfying this constraint is the cross-covariance of a jointly Gaussian coupling. The expansion in the proof of Proposition~\ref{prop-gw-euclidean-cost-regularization} shows that the inner-product objective $\mathcal Q_L(\pi)$ depends on $\pi$ only through the cross-moment
$\int yx^\top\d\pi=m_\beta m_\alpha^\top+K_\pi$. Replacing any coupling by this Gaussian coupling therefore leaves the objective unchanged. This proves Gaussian optimality, including for nonzero means.

When the means vanish, the positive-semidefinite block constraint is equivalent to the exact factorization
$K=\Sigma_\beta^{1/2}R^\top\Sigma_\alpha^{1/2}$ with
$R\in\RR^{d_x\times d_y}$ and $\norm{R}_{\mathrm{op}}\leq1$; this remains valid for singular covariances. Von Neumann's trace inequality then gives \(\max_K\norm{K}_{\mathrm F}^2 = \sum_{i=1}^{\min(d_x,d_y)}\lambda_i\eta_i.\)
Substitution into the inner-product instance of~\eqref{eq-gw-euclidean-cost-learning} yields~\eqref{eq-centered-gaussian-inner-product-gw}.

For a centered jointly Gaussian pair $(X,Y)$ with cross-covariance $K$, Isserlis' formula gives
\[
\mathcal Q_L(\pi)
=
4\bigl(\tr\Sigma_\alpha-\tr\Sigma_\beta\bigr)^2
+8\left(\tr\Sigma_\alpha^2+\tr\Sigma_\beta^2-2\norm{K}_{\mathrm F}^2\right).
\]
Squared distances are translation invariant, so the same expression applies for arbitrary means. Maximizing $\norm{K}_{\mathrm F}^2$ as above proves~\eqref{eq-gaussian-constrained-squared-gw}. More explicitly, write
$\Sigma_\alpha=U\diag(\lambda_i)U^\top$ and
$\Sigma_\beta=V\diag(\eta_j)V^\top$. For arbitrary signs $s_j\in\{-1,1\}$, the map
\[
T(x)
=
m_\beta+V
\begin{bmatrix}
\diag\!\left(s_j\sqrt{\eta_j/\lambda_j}\right)_{j=1}^{d_y}&0
\end{bmatrix}
U^\top(x-m_\alpha)
\]
pushes $\alpha$ to $\beta$ and attains the maximizing cross-covariance. This is the claimed principal-component alignment.
\end{proof}
For the squared-distance dissimilarities~\eqref{eq-gw-squared-distance-dissimilarities}, the unrestricted problem is subtler because~\eqref{eq-gw-squared-distance-cost-regularization} also contains the mixed fourth moment $\int\norm{x}^2\norm{y}^2\d\pi$, which is not fixed by the cross-moment $\int yx^\top\d\pi$ outside the Gaussian class.

\begin{prop}[Failure of Gaussian closure for squared-distance GW]\label{prop-gw-gaussian-closure-fails}
For all sufficiently large integers $n$, consider
\(\alpha=\Gaussian(0,1)\quad\text{on }\RR, \qquad \beta_n=\Gaussian(0,\Id_n)\quad\text{on }\RR^n,\)
with the squared-distance dissimilarities~\eqref{eq-gw-squared-distance-dissimilarities}, and write $\XX_n=(\RR,\dist_\X,\alpha)$ and $\YY_n=(\RR^n,\dist_\Y,\beta_n)$. No jointly Gaussian coupling is optimal for the unrestricted GW problem between $\alpha$ and $\beta_n$. In particular,
\(\GW(\XX_n,\YY_n)^2<\GW^{\mathrm{Gau}}(\XX_n,\YY_n)^2.\)
\end{prop}
\begin{proof}
For $(X,Y)\sim\pi$, equation~\eqref{eq-gw-squared-distance-cost-regularization} shows that minimizing the GW objective is equivalent to maximizing
\(\mathcal G(\pi) \eqdef \EE[X^2\norm{Y}^2] + 2\norm{\EE[YX]}^2.\)
If $(X,Y)$ is jointly Gaussian, write $k=\EE[YX]\in\RR^n$. Positivity of its block covariance gives $\norm{k}^2\leq1$, while Isserlis' formula gives \(\EE[X^2\norm{Y}^2]=n+2\norm{k}^2.\)

We construct a better non-Gaussian coupling. Let $Y\sim\Gaussian(0,\Id_n)$, set $R=\norm{Y}^2$, and denote by $G_j$ the cumulative distribution function of a $\chi_j^2$ variable. Then $U=G_n(R)$ is uniform on $(0,1)$ and $A=G_1^{-1}(U)$ has law $\chi_1^2$. If $\varepsilon$ is an independent Rademacher variable and $X=\varepsilon\sqrt A$, then $X\sim\Gaussian(0,1)$ and $\EE[YX]=0$. This coupling therefore satisfies
\(\mathcal G(\pi) = \EE[AR] = \int_0^1G_1^{-1}(u)G_n^{-1}(u)\d u.\)
The central limit theorem gives weak convergence of the standardized $\chi_n^2$ variable to a standard Gaussian. Since their second moments also converge, this convergence holds in $\Wass_2$; the quantile formula of Proposition~\ref{prop-wass-quantile-1d} therefore gives
\(\frac{G_n^{-1}(u)-n}{\sqrt{2n}} \longrightarrow \Phi^{-1}(u) \qquad\text{in }L^2(0,1),\)
where $\Phi$ is the standard Gaussian cumulative distribution function. Since $G_1^{-1}\in L^2(0,1)$, \(\mathcal G(\pi) = n+\sqrt{2n}\bigl(c+o(1)\bigr), \qquad c\eqdef\int_0^1G_1^{-1}(u)\Phi^{-1}(u)\d u>0.\)
Here $c$ is the covariance of the two strictly increasing, nonconstant functions $G_1^{-1}$ and $\Phi^{-1}$ of a uniform random variable, so $c>0$. Thus $\mathcal G(\pi)>n+4$ for all sufficiently large $n$, which is strictly better than every jointly Gaussian coupling.
\end{proof}
\paragraph{Monge--GW maps.}
\begin{defn}[Gromov--Monge problem]\label{def-gromov-monge}
Fix $p\geq1$ and the discrepancy $\De$ used in Definition~\ref{def-gromov-wasserstein}. The directed Gromov--Monge value is
\begin{equation}\label{eq-gromov-monge-p}
\GW_p^{\mathrm M}(\XX,\YY)
\eqdef
\inf_{\substack{T:\X\to\Y\ \mathrm{measurable}\\T_\sharp\alpha=\beta}}
\left(
\iint_{\X^2}
\De\bigl(\dist_\X(x,x'),\dist_\Y(T(x),T(x'))\bigr)^p
\d\alpha(x)\d\alpha(x')
\right)^{1/p},
\end{equation}
with value $+\infty$ when no such measure-preserving map exists.
\end{defn}
Its infinite-order extension is
\begin{equation}\label{eq-gromov-monge-infinity}
\GW_\infty^{\mathrm M}(\XX,\YY)
\eqdef
\inf_{\substack{T:\X\to\Y\ \mathrm{measurable}\\T_\sharp\alpha=\beta}}
\operatorname*{ess\,sup}_{\alpha\otimes\alpha}
\De\bigl(\dist_\X(x,x'),\dist_\Y(T(x),T(x'))\bigr).
\end{equation}
For $1\leq p\leq q<\infty$, monotonicity of $L^p$ norms and the restriction from couplings to graph couplings give
\begin{equation}\label{eq-gromov-monge-basic-comparisons}
\begin{aligned}
\GW_p^{\mathrm M}(\XX,\YY)
&\leq
\GW_q^{\mathrm M}(\XX,\YY)
\leq
\GW_\infty^{\mathrm M}(\XX,\YY),\\
\GW_p(\XX,\YY)
&\leq
\GW_p^{\mathrm M}(\XX,\YY),
\qquad
\GW_\infty(\XX,\YY)\leq\GW_\infty^{\mathrm M}(\XX,\YY).
\end{aligned}
\end{equation}
\begin{prop}[Finite uniform Gromov--Monge limits]\label{prop-finite-uniform-gromov-monge-limit}
Let $\X=\{x_i\}_{i=1}^n$ and $\Y=\{y_j\}_{j=1}^n$ carry uniform measures. Then
\begin{align}
\GW_p^{\mathrm M}(\XX,\YY)
&=
\min_{\sigma\in\mathfrak S_n}
\left(
\frac1{n^2}\sum_{i,i'=1}^n
\De\bigl(\dist_\X(x_i,x_{i'}),\dist_\Y(y_{\sigma(i)},y_{\sigma(i')})\bigr)^p
\right)^{1/p},
\label{eq-finite-gromov-monge-p}\\
\lim_{p\to\infty}\GW_p^{\mathrm M}(\XX,\YY)
&=
\GW_\infty^{\mathrm M}(\XX,\YY)
=
\GW_\infty(\XX,\YY).
\label{eq-finite-gw-gm-infinity-limit}
\end{align}
\end{prop}
\begin{proof}
A measure-preserving map between equal-size uniform spaces is a permutation; this proves~\eqref{eq-finite-gromov-monge-p}. Since $\mathfrak S_n$ is finite, the minimum commutes with $p\to\infty$.

For the last equality, let $\P_\infty$ minimize $\GW_\infty$. The scaled matrix $n\P_\infty$ is doubly stochastic, so Theorem~\ref{thm-birkhoff-von-neumann} provides a permutation whose graph is contained in $\supp(\P_\infty)$. Its distortion cannot exceed the maximal distortion on $\supp(\P_\infty)^2$, and hence $\GW_\infty^{\mathrm M}\leq\GW_\infty$. The reverse inequality is~\eqref{eq-gromov-monge-basic-comparisons}.
\end{proof}
For the inner-product dissimilarities~\eqref{eq-gw-inner-product-dissimilarities}, the learned-cost representation yields a Monge optimizer; the stronger transformed-Brenier description obtained from the twist condition has a precise rank limitation.

\begin{prop}[Monge--GW map for inner-product dissimilarities]\label{prop-gw-inner-product-monge-map}
Let $\X\subset\RR^{d_x}$ and $\Y\subset\RR^{d_y}$ be compact, assume $d_x\geq d_y$, and let $\alpha$ be absolutely continuous with respect to Lebesgue measure. Then the GW problem with dissimilarities~\eqref{eq-gw-inner-product-dissimilarities} admits an optimal coupling of the form $(\Id,T)_\sharp\alpha$. In the inner-product instance of~\eqref{eq-gw-euclidean-cost-learning}, let $M^\star$ be an optimal learned matrix. If $\rank(M^\star)=d_y$, then the map $y\mapsto\nabla_x c_{M^\star}(x,y)$ is injective for every $x$; the optimal coupling for this fixed cost is unique and
\begin{equation}\label{eq-gw-inner-product-transformed-brenier-map}
T(x)=\nabla\phi(M^\star x)
\end{equation}
for the Brenier potential $\phi$ transporting $(M^\star)_\sharp\alpha$ to $\beta$.
\end{prop}
\begin{proof}
Compactness of the coupling set and coercivity of the quadratic penalty in $M$ give an optimal pair $(\pi^\star,M^\star)$ in the inner-product instance of~\eqref{eq-gw-euclidean-cost-learning}. With $M^\star$ frozen, $\pi^\star$ solves linear OT for $c_{M^\star}$. The linear-cost Monge-existence theorem of Dumont, Lacombe and Vialard, also used in the cost-regularized setting in, gives a graph-supported optimizer for every matrix $M^\star$ under the stated dimension and absolute-continuity assumptions. Its rank-deficient case combines transport on the image of $M^\star$ with measurable transport along its fibers. Replacing $\pi^\star$ by this optimizer leaves the joint objective unchanged.

If $M^\star$ has full row rank, then
\(y\longmapsto\nabla_x c_{M^\star}(x,y)=-2(M^\star)^\top y\)
is injective, so the cost is twisted. Moreover, $(M^\star)_\sharp\alpha$ is absolutely continuous on $\RR^{d_y}$, and minimizing $-2\dotp{z}{y}$ is equivalent, up to marginal-only terms, to quadratic transport between $z=M^\star x$ and $y$. Brenier's theorem, Theorem~\ref{thm-brenier}, therefore gives the unique map $z\mapsto\nabla\phi(z)$, proving~\eqref{eq-gw-inner-product-transformed-brenier-map}. This is precisely the full-rank construction of.
\end{proof}
If $M^\star$ is rank deficient, the elementary twist argument fails and uniqueness need not hold; the existence assertion nevertheless remains valid by the linear-cost analysis cited in the proof. This distinction is essential: a Monge optimizer always exists here under the hypotheses of Proposition~\ref{prop-gw-inner-product-monge-map}, but a literal full-rank Brenier characterization does not always apply.

\paragraph{Entropic regularization and biconvex alternating minimization.}
Assume in this paragraph that $\X$ and $\Y$ are compact metric spaces, that $L$ in~\eqref{eq-gw-quadratic-form} is continuous, and set $\mathfrak m=\alpha\otimes\beta$. With the relative entropy of Definition~\ref{def-measure-relative-entropy}, define
\[
\mathcal R_\epsilon(\pi)
\eqdef
\begin{cases}
0, & \epsilon=0,\\
\epsilon\KL(\pi|\mathfrak m), & \epsilon>0.
\end{cases}
\]
\begin{defn}[Entropic GW and its biconvex relaxation]\label{def-entropic-gw-biconvex}
For $\epsilon\geq0$, set
\begin{align}
\mathcal F_\epsilon(\pi)
&\eqdef \mathcal Q_L(\pi)+\mathcal R_\epsilon(\pi),
&
\GW_\epsilon(\XX,\YY)^p
&\eqdef \umin{\pi\in\Couplings(\alpha,\beta)}\mathcal F_\epsilon(\pi),
\label{eq-gw-entropic-general}\\
\mathcal F_\epsilon^{\mathrm{bi}}(\pi,\xi)
&\eqdef \mathcal B_L(\pi,\xi)
+\frac12\mathcal R_\epsilon(\pi)
+\frac12\mathcal R_\epsilon(\xi),
&
\GW_{\epsilon,\mathrm{bi}}(\XX,\YY)^p
&\eqdef \umin{\pi,\xi\in\Couplings(\alpha,\beta)}
\mathcal F_\epsilon^{\mathrm{bi}}(\pi,\xi).
\label{eq-gw-biconvex-relaxation}
\end{align}
\end{defn}
Diagonal consistency also gives the lower bound
\begin{equation}\label{eq-gw-biconvex-lower-bound}
\GW_{\epsilon,\mathrm{bi}}(\XX,\YY)^p
\leq
\GW_\epsilon(\XX,\YY)^p.
\end{equation}
Block optimality explains both the usefulness and the limitation of the relaxation. If $(\pi_\star,\xi_\star)$ minimizes~\eqref{eq-gw-biconvex-relaxation}, then each component minimizes an OT problem with the other component frozen. Indeed, for the linearized cost $c_\xi$ defined in~\eqref{eq-gw-linearized-cost}, one has $\mathcal B_L(\pi,\xi)=\frac12\int c_\xi\d\pi$ and therefore
\[
\pi_\star\in\uargmin{\pi\in\Couplings(\alpha,\beta)}
\left\{\int c_{\xi_\star}\d\pi+\mathcal R_\epsilon(\pi)\right\},
\]
and symmetrically for $\xi_\star$. Thus the relaxation replaces the quadratic problem by two OT problems, but their ground costs are themselves unknown.

In the unregularized discrete case, these block problems are linear programs. If $\alpha$ and $\beta$ are both uniform on $n$ points, at least one \emph{global minimizer of the lifted problem} has the form $(\P_\sigma/n,\P_\tau/n)$.

\begin{prop}[Tightness of the biconvex GW relaxation]\label{prop-gw-biconvex-tightness}
Assume that $\mathcal Q_L$ is concave on $\Couplings(\alpha,\beta)$. Since $\mathcal Q_L$ is quadratic, this is equivalent to
\begin{equation}\label{eq-gw-biconvex-negative-condition}
\mathcal Q_L(\pi-\xi)\leq0
\qquad
\text{for every }\pi,\xi\in\Couplings(\alpha,\beta).
\end{equation}
Then~\eqref{eq-gw-biconvex-lower-bound} is an equality for every $\epsilon\geq0$. Moreover, every minimizing pair $(\pi_\star,\xi_\star)$ satisfies
\(\mathcal F_\epsilon(\pi_\star) = \mathcal F_\epsilon(\xi_\star) = \GW_\epsilon(\XX,\YY)^p, \qquad \mathcal Q_L(\pi_\star-\xi_\star)=0.\)
If $\epsilon>0$, then $\pi_\star=\xi_\star$. The same conclusion holds for $\epsilon=0$ when the inequality in~\eqref{eq-gw-biconvex-negative-condition} is strict for every nonzero feasible difference.
\end{prop}
\begin{proof}
Symmetry and bilinearity give the polarization identity
\begin{equation}\label{eq-gw-biconvex-polarization}
\mathcal F_\epsilon^{\mathrm{bi}}(\pi,\xi)
=
\frac12\mathcal F_\epsilon(\pi)
+
\frac12\mathcal F_\epsilon(\xi)
-
\frac12\mathcal Q_L(\pi-\xi).
\end{equation}
Under~\eqref{eq-gw-biconvex-negative-condition}, the right-hand side is at least $\GW_\epsilon(\XX,\YY)^p$. Conversely, diagonal pairs recover $\mathcal F_\epsilon$, proving equality of the optimal values. Equality in the preceding bounds gives the two asserted identities. Strict negativity immediately forces $\pi_\star=\xi_\star$. Suppose now that $\epsilon>0$. The two diagonal objective values are finite, so both relative entropies are finite. If $\pi_\star\ne\xi_\star$, then $\mathcal Q_L(\pi_\star-\xi_\star)=0$ makes the quadratic term affine on their joining segment, while strict convexity of relative entropy makes its value at the midpoint strictly smaller than the average of its endpoint values. The midpoint would therefore have a smaller diagonal objective, a contradiction.
\end{proof}
Starting from $(\pi^{(0)},\xi^{(0)})\in\Couplings(\alpha,\beta)^2$, exact alternating minimization of the general-measure relaxation reads
\begin{align}\label{eq-gw-biconvex-alternating-general}
\pi^{(\ell+1)}
&\in
\uargmin{\pi\in\Couplings(\alpha,\beta)}
\left\{\int c_{\xi^{(\ell)}}\d\pi+\mathcal R_\epsilon(\pi)\right\},
\\[-.2em]
\xi^{(\ell+1)}
&\in
\uargmin{\xi\in\Couplings(\alpha,\beta)}
\left\{\int c_{\pi^{(\ell+1)}}\d\xi+\mathcal R_\epsilon(\xi)\right\}.
\end{align}
For discrete measures, let $\HD$ denote the discrete Shannon--Boltzmann entropy of Definition~\ref{def-discrete-shannon-boltzmann-entropy}. Given $\P\in\CouplingsD(\a,\b)$, define its contracted GW cost by
\begin{equation}\label{eq-gw-contracted-cost}
\C(\P)_{i,j}
\eqdef
\sum_{i'=1}^{n}\sum_{j'=1}^{m}
\De\bigl(\distD_{i,i'},\distD'_{j,j'}\bigr)^p
\P_{i',j'}.
\end{equation}
The KL identity on $\CouplingsD(\a,\b)$ shows that, after removing the fixed additive constant $\epsilon(\HD(\a)+\HD(\b))$, the discrete biconvex problem is
\begin{equation}\label{eq-gw-biconvex-discrete}
\umin{\P,\Q\in\CouplingsD(\a,\b)}
\dotp{\C(\P)}{\Q}
-
\frac{\epsilon}{2}\bigl(\HD(\P)+\HD(\Q)\bigr),
\end{equation}
and the original lifted optimal value is obtained by adding this constant back. There is no factor $1/2$ in front of the bilinear term under the convention~\eqref{eq-gw-contracted-cost}: $\C(\P)$ is half of the full linearized cost in~\eqref{eq-gw-linearized-cost}. Restricting~\eqref{eq-gw-biconvex-discrete} to the diagonal $\P=\Q$ recovers the usual entropic GW problem
\begin{equation}\label{eq-gw-entropy}
\umin{\P\in\CouplingsD(\a,\b)}
\Ee_{\distD,\distD'}(\P)-\epsilon\HD(\P).
\end{equation}
For an unstructured distortion kernel, evaluating one entry of~\eqref{eq-gw-contracted-cost} requires $nm$ terms. Forming all $nm$ entries therefore costs $O(n^2m^2)$ dense arithmetic operations; explicitly storing the kernel would likewise require $O(n^2m^2)$ memory. The squared loss admits a substantially cheaper contraction without forming this four-index tensor.

\begin{prop}[Fast contracted cost for squared GW]\label{prop-gw-squared-cost-update}
Assume that
\(\De\bigl(\distD_{i,i'},\distD'_{j,j'}\bigr)^p = \bigl(\distD_{i,i'}-\distD'_{j,j'}\bigr)^2.\)
Then, for every $\P\in\CouplingsD(\a,\b)$,
\begin{equation}\label{eq-gw-sinkh}
\C(\P)
=
\distD^{\odot2}\a\,\ones_m^\top
+
\ones_n(\distD'^{\odot2}\b)^\top
-
2\distD\,\P\,\transp{\distD'}.
\end{equation}
Its ambient Frobenius gradient satisfies
$\nabla\Ee_{\distD,\distD'}(\P)=2\C(\P)$.
\end{prop}
\begin{proof}
Expanding the square in~\eqref{eq-gw-contracted-cost} gives
\[
\C(\P)_{i,j}
=
\sum_{i',j'}\distD_{i,i'}^2\P_{i',j'}
+
\sum_{i',j'}(\distD'_{j,j'})^2\P_{i',j'}
-
2\sum_{i',j'}\distD_{i,i'}\distD'_{j,j'}\P_{i',j'}.
\]
The marginal identities $\P\ones_m=\a$ and $\transp{\P}\ones_n=\b$ turn the first two sums into $(\distD^{\odot2}\a)_i$ and $(\distD'^{\odot2}\b)_j$, while the last sum is $(\distD\P\transp{\distD'})_{i,j}$. This proves~\eqref{eq-gw-sinkh}. The underlying four-index kernel is symmetric, so differentiating
$\Ee_{\distD,\distD'}(\P)=\dotp{\C(\P)}{\P}$
gives the factor two in the gradient. Finally, computing $\distD\P$ and then $(\distD\P)\transp{\distD'}$ costs respectively $O(n^2m)$ and $O(nm^2)$ operations; the marginal terms are cheaper.
\end{proof}
Since the lifted functional is symmetric in its two blocks, interleaving
\(\Q^{(0)},\ \P^{(1)},\ \Q^{(1)},\ \P^{(2)},\ldots\)
\paragraph{Fused GW: adding features.}
In machine-learning applications, points in the two spaces often carry features in a common metric space $(\mathcal Z_{\mathrm f},\dist_{\mathcal Z_{\mathrm f}})$ through embeddings
\(\phi_\X:\X\to\mathcal Z_{\mathrm f}, \qquad \phi_\Y:\Y\to\mathcal Z_{\mathrm f},\)
\begin{defn}[Fused Gromov--Wasserstein problem]\label{def-fused-gromov-wasserstein}
Under the notation of~\eqref{eq-gw-quadratic-form}, assume that the two distortion-dependent sizes of Definition~\ref{def-gromov-wasserstein} are finite. Let
$c_{\X,\Y}:\X\times\Y\to\RR_+$ be a lower-semicontinuous cross-space cost satisfying
$\int c_{\X,\Y}(x,y)\d\alpha(x)\d\beta(y)<+\infty$, and let $0\leq\lambda\leq1$. The fused GW value is
\begin{equation}\label{eq-fused-gromov-wasserstein}
\operatorname{FGW}_{\lambda,p}(\XX,\YY;c_{\X,\Y})^p
\eqdef
\umin{\pi\in\Couplings(\alpha,\beta)}
(1-\lambda)\int_{\X\times\Y}c_{\X,\Y}\d\pi
+
\lambda\mathcal Q_L(\pi).
\end{equation}
At $\lambda=1$, the feature term is omitted, including for couplings on which its integral is infinite.
\end{defn}
The objective is lower semicontinuous on the weakly compact coupling set, while $\alpha\otimes\beta$ is a finite-valued competitor; hence the minimum is attained. The first term compares features by ordinary Kantorovich transport, while the second compares intrinsic geometry.

For discrete measures, setting $M_{i,j}=c_{\X,\Y}(x_i,y_j)$ reduces~\eqref{eq-fused-gromov-wasserstein} to
\(\umin{\P\in\CouplingsD(\a,\b)} (1-\lambda)\dotp{M}{\P} + \lambda\Ee_{\distD,\distD'}(\P),\)
The algorithms above extend by the same linearization principle. With the coupling regularizer $\mathcal R_\epsilon$ defined before~\eqref{eq-gw-entropic-general}, the regularized fused objective is
\[
\mathcal F_{\epsilon,\lambda}^{\mathrm{FGW}}(\pi)
\eqdef
(1-\lambda)\int_{\X\times\Y}c_{\X,\Y}\d\pi
+
\lambda\mathcal Q_L(\pi)
+
\mathcal R_\epsilon(\pi).
\]
Its first variation is an ordinary OT objective with ground cost
\begin{equation}\label{eq-fused-gw-linearized-cost}
c_{\pi,\lambda}^{\mathrm{FGW}}(x,y)
\eqdef
(1-\lambda)c_{\X,\Y}(x,y)
+
2\lambda
\int_{\X\times\Y}
L\bigl((x,y),(x',y')\bigr)\d\pi(x',y').
\end{equation}
Equivalently, diagonal consistency gives the symmetric biconvex lift
\[
\mathcal F_{\epsilon,\lambda}^{\mathrm{FGW,bi}}(\pi,\xi)
\eqdef
\frac{1-\lambda}{2}
\int_{\X\times\Y}c_{\X,\Y}\d(\pi+\xi)
+
\lambda\mathcal B_L(\pi,\xi)
+
\frac12\mathcal R_\epsilon(\pi)
+
\frac12\mathcal R_\epsilon(\xi).
\]
Multiplying a block objective by two shows that, with $\xi$ fixed, its $\pi$-update is precisely
\[
\uargmin{\pi\in\Couplings(\alpha,\beta)}
\left\{
\int_{\X\times\Y}c_{\xi,\lambda}^{\mathrm{FGW}}\d\pi
+
\mathcal R_\epsilon(\pi)
\right\}.
\]
\paragraph{Hausdorff and Gromov--Hausdorff viewpoints.}
\begin{defn}[Hausdorff distance]\label{def-hausdorff-distance}
If $A,B$ are nonempty compact subsets of a metric space $(\Z,\dist_\Z)$, their Hausdorff distance is
\[
d_{\mathrm H}^{\Z}(A,B)
\eqdef
\max\left\{
\sup_{a\in A}\inf_{b\in B}\dist_\Z(a,b),
\sup_{b\in B}\inf_{a\in A}\dist_\Z(a,b)
\right\}.
\]
\end{defn}
\begin{defn}[Gromov--Hausdorff distance]\label{def-gromov-hausdorff-distance}
For compact metric spaces $(\X,\dist_\X)$ and $(\Y,\dist_\Y)$, their Gromov--Hausdorff distance is \(d_{\mathrm{GH}}(\X,\Y) \eqdef \inf_{\Z,\phi,\psi} d_{\mathrm H}^{\Z}(\phi(\X),\psi(\Y)).\)
Here the infimum ranges over metric spaces $(\Z,\dist_\Z)$ and isometric embeddings $\phi:\X\hookrightarrow\Z$ and $\psi:\Y\hookrightarrow\Z$.
\end{defn}
One has $d_{\mathrm{GH}}(\X,\Y)=0$ exactly when $\X$ and $\Y$ are isometric. Hence $d_{\mathrm{GH}}$ is a genuine metric on isometry classes of compact metric spaces.

A correspondence is a set $R\subset\X\times\Y$ whose two coordinate projections are surjective, and its distortion is \(\operatorname{dis}(R) \eqdef \sup_{(x,y),(x',y')\in R} \abs{\dist_\X(x,x')-\dist_\Y(y,y')}.\)
Writing $\mathfrak R(\X,\Y)$ for the set of correspondences, compact metric spaces satisfy the classical identity
\begin{equation}\label{eq-gh-correspondence-distortion}
2d_{\mathrm{GH}}(\X,\Y)
=
\inf_{R\in\mathfrak R(\X,\Y)}\operatorname{dis}(R).
\end{equation}
\begin{prop}[Gromov--Hausdorff as an infinite-order intrinsic matching]\label{prop-gh-gm-gw-infinity}
Let $\X$ and $\Y$ be compact metric spaces and let $\alpha,\beta$ have full support. Then
\begin{equation}\label{eq-gh-gw-gm-infinity-chain}
2d_{\mathrm{GH}}(\X,\Y)
\leq
\GW_\infty(\XX,\YY).
\end{equation}
Moreover, removing the prescribed weights recovers Gromov--Hausdorff exactly:
\begin{equation}\label{eq-gh-as-weight-envelope-gw-infinity}
2d_{\mathrm{GH}}(\X,\Y)
=
\inf_{\substack{\alpha\in\Mm_+^1(\X),\ \supp(\alpha)=\X\\
\beta\in\Mm_+^1(\Y),\ \supp(\beta)=\Y}}
\GW_\infty\bigl((\X,\dist_\X,\alpha),(\Y,\dist_\Y,\beta)\bigr).
\end{equation}
\end{prop}
\begin{proof}
The support of any $\pi\in\Couplings(\alpha,\beta)$ is compact. The support of each marginal is the closure of the corresponding coordinate projection of $\supp(\pi)$. These projections are compact, hence already closed; because $\alpha$ and $\beta$ have full support, they are therefore equal to $\X$ and $\Y$. Thus $\supp(\pi)$ is a correspondence, and~\eqref{eq-gh-correspondence-distortion} gives~\eqref{eq-gh-gw-gm-infinity-chain}.

For~\eqref{eq-gh-as-weight-envelope-gw-infinity}, the preceding lower bound applies to every pair of full-support weights. Conversely, choose a correspondence $R$ with distortion arbitrarily close to the right-hand side of~\eqref{eq-gh-correspondence-distortion}. Replacing $R$ by its closure preserves both the correspondence property and its distortion, so $R$ may be assumed compact. It then carries a probability $\pi$ with support exactly $R$. Let $\alpha$ and $\beta$ be its marginals. They have full support, while~\eqref{eq-gw-infinity-support} gives $\GW_\infty(\XX,\YY)\leq\operatorname{dis}(R)$. Taking the infimum over $R$ proves the identity.
\end{proof}
Thus Gromov--Hausdorff is the mass-free, hard-support counterpart of GW: it minimizes worst-case distortion over correspondences between the two supports. By contrast, finite-$p$ GW fixes the masses and averages distortion under a coupling; optimizing these masses and taking $p=\infty$ recovers $2d_{\mathrm{GH}}$ through~\eqref{eq-gh-as-weight-envelope-gw-infinity}.

\subsection{Quantum Optimal Transport}
\label{sec-quantum-ot}
\paragraph{Finite-dimensional states and couplings.}
\begin{defn}[Hermitian and density matrices]\label{def-hermitian-density-matrices}
Let $\mathbb H_n$ be the real vector space of $n\times n$ Hermitian matrices, \(\mathbb H_n^+ = \enscond{A\in\mathbb H_n}{A\succeq0}, \qquad \mathbb H_n^{+,1} = \enscond{A\in\mathbb H_n^+}{\tr(A)=1}.\)
Here $A\succeq0$ means that $A$ is positive semidefinite, namely $z^*Az\geq0$ for every $z\in\CC^n$, or equivalently that all eigenvalues of $A$ are nonnegative.
Elements of $\mathbb H_n^{+,1}$ are density matrices.
\end{defn}
For $A\in\mathbb H_n$ and $B\in\mathbb H_m$, the tensor product $A\otimes B$ is the linear operator determined by
\begin{equation}\label{eq-qot-tensor-product}
(A\otimes B)(x\otimes y)=(Ax)\otimes(By),
\qquad
A\otimes B=[A_{i,k}B]_{i,k=1}^n,
\qquad
(A\otimes B)_{(i,j),(k,\ell)}=A_{i,k}B_{j,\ell}.
\end{equation}
\begin{defn}[Joint quantum states and partial traces]\label{def-joint-quantum-state}
A joint quantum state between $\CC^n$ and $\CC^m$ is a density matrix $T\in\mathbb H_{nm}^{+,1}$ acting on $\CC^n\otimes\CC^m$. Its marginals are the partial traces $\operatorname{Tr}_B(T)\in\mathbb H_n^{+,1}$ and $\operatorname{Tr}_A(T)\in\mathbb H_m^{+,1}$ defined, for all $F\in\mathbb H_n$ and $G\in\mathbb H_m$, by
\begin{equation}\label{eq-qot-partial-traces}
\tr(F\,\operatorname{Tr}_B T)=\tr((F\otimes\Id_m)T),
\qquad
\tr(G\,\operatorname{Tr}_A T)=\tr((\Id_n\otimes G)T),
\end{equation}
and they play the role of the two marginals of a classical coupling.
\end{defn}
In product bases $(e_i)_{i=1}^n$ and $(f_j)_{j=1}^m$, the operator $T$ can equivalently be viewed as the four-index tensor
\(T_{(i,j),(k,\ell)} = \langle e_i\otimes f_j,T(e_k\otimes f_\ell)\rangle.\)
The full and partial traces then read
\begin{equation}\label{eq-qot-partial-traces-indices}
\tr(T)=\sum_{i=1}^n\sum_{j=1}^m T_{(i,j),(i,j)},
\qquad
[\operatorname{Tr}_B T]_{i,k}=\sum_{j=1}^m T_{(i,j),(k,j)},
\qquad
[\operatorname{Tr}_A T]_{j,\ell}=\sum_{i=1}^n T_{(i,j),(i,\ell)}.
\end{equation}
\begin{defn}[Finite-dimensional quantum OT]\label{def-finite-dimensional-qot}
Let $A\in\mathbb H_n^{+,1}$, $B\in\mathbb H_m^{+,1}$ and let $C\in\mathbb H_{nm}$ be a Hermitian cost observable. The quantum OT value is the semidefinite program
\begin{equation}\label{eq-qot-primal}
\mathrm{QOT}_C(A,B)
\eqdef
\min_{T\in\mathbb H_{nm}^+}
\enscond{\tr(CT)}{\operatorname{Tr}_B(T)=A,\ \operatorname{Tr}_A(T)=B}.
\end{equation}
\end{defn}
\paragraph{Dual formulation.}
\begin{prop}[Quantum Kantorovich duality]\label{prop-qot-duality}
For $A\in\mathbb H_n^{+,1}$ and $B\in\mathbb H_m^{+,1}$, the dual of~\eqref{eq-qot-primal} is
\begin{equation}\label{eq-qot-dual}
\mathrm{QOT}_C(A,B)
=
\sup_{F\in\mathbb H_n,\ G\in\mathbb H_m}
\left\{
\tr(FA)+\tr(GB)
:\,
F\otimes\Id_m+\Id_n\otimes G\preceq C
\right\}.
\end{equation}
There is no duality gap. If $A$ and $B$ are positive definite, the supremum is attained. For singular marginals, the primal reduces to the tensor product of their supports; on this reduced space the corresponding dual maximum is attained, while the full-space formula above is always valid as a supremum.
\end{prop}
\begin{proof}
Introduce Hermitian Lagrange multipliers $F$ and $G$ for the two marginal constraints. The Lagrangian is
\[
\tr(CT)+\tr\!\left(F(A-\operatorname{Tr}_B T)\right)
+\tr\!\left(G(B-\operatorname{Tr}_A T)\right)
=
\tr(FA)+\tr(GB)
+\tr\!\left((C-F\otimes\Id_m-\Id_n\otimes G)T\right),
\]
where~\eqref{eq-qot-partial-traces} was used in the last equality. Minimizing over $T\succeq0$ gives a finite lower bound if and only if $C-F\otimes\Id_m-\Id_n\otimes G\succeq0$, in which case the infimum in $T$ is $0$. When $A,B\succ0$, the coupling $A\otimes B$ is strictly feasible, so Slater's theorem gives equality of primal and dual values and dual attainment.

For singular marginals, let $P_A$ and $P_B$ be their support projections. If $T$ is feasible, then
\(\tr\bigl(((\Id_n-P_A)\otimes\Id_m)T\bigr) = \tr((\Id_n-P_A)A)=0.\)
Positivity implies that $T$ vanishes on $(\supp A)^\perp\otimes\CC^m$; the analogous argument for $B$ yields
\(T=(P_A\otimes P_B)T(P_A\otimes P_B)\).
Thus the primal is exactly the problem obtained by compressing $A$, $B$ and $C$ to $\supp A\otimes\supp B$. The reduced marginals are positive definite, so Slater applies there. Equivalently, the unreduced dual values approach the reduced maximum after assigning sufficiently negative potentials on the orthogonal complements. This proves the stated supremum formula without a gap.
\end{proof}
\paragraph{Quantum-to-classical OT connection.}
\begin{defn}[Antisymmetric lift of a ground cost]\label{def-qot-ground-cost-lift}
Assume $m=n$, and let $\C_0=(\C_{0,i,j})\in\RR_+^{n\times n}$ be symmetric, with $\C_{0,i,i}=0$ and $\C_{0,i,j}>0$ for $i\neq j$. In the product basis $(e_i\otimes e_j)_{i,j=1}^n$, define
\begin{equation}\label{eq-qot-ground-cost-lift}
s_{i,j}^-\eqdef\frac{e_i\otimes e_j-e_j\otimes e_i}{\sqrt2},
\qquad
C^{\mathrm Q}[\C_0]
\eqdef
2\sum_{1\leq i<j\leq n}\C_{0,i,j}s_{i,j}^-(s_{i,j}^-)^*.
\end{equation}
For $A,B\in\mathbb H_n^{+,1}$ and $\a,\b\in\simplex_n$, define the associated quantum and diagonal classical costs by
\begin{equation}\label{eq-qot-ground-cost-values}
\mathcal W_{\C_0}^{\mathrm Q}(A,B)
\eqdef\mathrm{QOT}_{C^{\mathrm Q}[\C_0]}(A,B),
\qquad
\mathcal W_{\C_0}^{\mathrm{cl}}(\a,\b)
\eqdef\mathcal W_{\C_0}^{\mathrm Q}(\diag(\a),\diag(\b)).
\end{equation}
\end{defn}
Entrywise,
\begin{equation}\label{eq-qot-ground-cost-lift-entries}
[C^{\mathrm Q}[\C_0]]_{(i,j),(k,\ell)}
=
\C_{0,i,j}\bigl(\delta_{i,k}\delta_{j,\ell}-\delta_{i,\ell}\delta_{j,k}\bigr).
\end{equation}
Thus the $n^2\times n^2$ observable is sparse: it has only $2n(n-1)$ nonzero entries, rather than order $n^4$. For each unordered pair $\{i,j\}$ it contains the block $\C_{0,i,j}\left(\begin{smallmatrix}1&-1\\-1&1\end{smallmatrix}\right)$ on $\operatorname{span}\{e_i\otimes e_j,e_j\otimes e_i\}$, and it vanishes on the diagonal vectors $e_i\otimes e_i$.

\begin{prop}[Ground-cost quantum semidistance]\label{prop-qot-ground-cost-semidistance}
The matrix $C^{\mathrm Q}[\C_0]$ is positive semidefinite and its kernel is the symmetric tensor subspace of $\CC^n\otimes\CC^n$.
for all $A,B\in\mathbb H_n^{+,1}$,
\[
\mathcal W_{\C_0}^{\mathrm Q}(A,B)\geq0,
\qquad
\mathcal W_{\C_0}^{\mathrm Q}(A,B)=\mathcal W_{\C_0}^{\mathrm Q}(B,A),
\qquad
\mathcal W_{\C_0}^{\mathrm Q}(A,B)=0\ \Longleftrightarrow\ A=B.
\]
These properties do not, by themselves, include the triangle inequality.
\end{prop}
\begin{proof}
The vectors $(s_{i,j}^-)_{i<j}$ form an orthonormal basis of the antisymmetric tensor subspace. Since every coefficient in~\eqref{eq-qot-ground-cost-lift} is positive, the asserted positivity and kernel follow. The tensor-factor swap preserves the cost observable and exchanges the two partial traces, which proves symmetry of the value.

If $A=B=\sum_r\lambda_r u_ru_r^*$, then the symmetric purification $\psi_A=\sum_r\sqrt{\lambda_r}\,u_r\otimes u_r$ has partial traces $A$ and zero cost. Conversely, let a feasible $T$ have zero cost. The positive matrix $(C^{\mathrm Q}[\C_0])^{1/2}T(C^{\mathrm Q}[\C_0])^{1/2}$ then has zero trace, hence is zero. It follows that $(C^{\mathrm Q}[\C_0])^{1/2}T^{1/2}=0$, so $\operatorname{Ran}T\subset\ker C^{\mathrm Q}[\C_0]$, the symmetric tensor subspace. Decompose $T$ into rank-one projectors onto symmetric vectors. Under the identification $\CC^n\otimes\CC^n\simeq\CC^{n\times n}$, each such vector is a symmetric matrix $X$; the two partial traces of its rank-one projector are $XX^*$ and $X^\top\overline X$, which coincide. Summing these identities gives $A=B$.
\end{proof}
\begin{prop}[Convex classical restriction]\label{prop-qot-classical-restriction}
For $\a,\b\in\simplex_n$, the diagonal restriction in~\eqref{eq-qot-ground-cost-values} is
\begin{equation}\label{eq-qot-classical-square-root-program}
\mathcal W_{\C_0}^{\mathrm{cl}}(\a,\b)
=
\min_{\P\in\CouplingsD(\a,\b)}
\sum_{1\leq i<j\leq n}
\C_{0,i,j}\bigl(\sqrt{\P_{i,j}}-\sqrt{\P_{j,i}}\bigr)^2.
\end{equation}
The square roots mean that this is not a quadratic program in $\P$. It is nevertheless convex and admits the exact second-order cone formulation
\begin{equation}\label{eq-qot-classical-socp}
\min_{\substack{\P\in\CouplingsD(\a,\b),\ z_{i,j}\geq0\\
z_{i,j}^2\leq\P_{i,j}\P_{j,i}}}
\sum_{1\leq i<j\leq n}\C_{0,i,j}
\bigl(\P_{i,j}+\P_{j,i}-2z_{i,j}\bigr).
\end{equation}
For unit vectors $u,v\in\CC^n$, the value between the associated rank-one states is
\begin{equation}\label{eq-qot-ground-cost-pure-states}
\mathcal W_{\C_0}^{\mathrm Q}(uu^*,vv^*)
=
\sum_{1\leq i<j\leq n}\C_{0,i,j}|u_i v_j-u_j v_i|^2.
\end{equation}
$\mathcal W_{\C_0}^{\mathrm{cl}}(\delta_i,\delta_j)
=\mathcal W_{\C_0}^{\mathrm Q}(e_i e_i^*,e_j e_j^*)=\C_{0,i,j}$.
When $\C_{0,i,j}=1$ for $i\neq j$, formula~\eqref{eq-qot-ground-cost-pure-states} reduces to
$\mathcal W_{\C_0}^{\mathrm Q}(uu^*,vv^*)=1-|u^*v|^2$.
\end{prop}
\begin{proof}
Let $T$ be feasible between $\diag(\a)$ and $\diag(\b)$, and set $\P_{i,j}=T_{(i,j),(i,j)}$. The diagonal entries form a coupling $\P\in\CouplingsD(\a,\b)$. Since the cost only involves the two-dimensional principal blocks indexed by $(i,j)$ and $(j,i)$, positivity gives
$\Re T_{(i,j),(j,i)}\leq\sqrt{\P_{i,j}\P_{j,i}}$ and hence
\(\tr(C^{\mathrm Q}[\C_0]T) \geq \sum_{i<j}\C_{0,i,j} \bigl(\P_{i,j}+\P_{j,i}-2\sqrt{\P_{i,j}\P_{j,i}}\bigr).\)
Conversely, for any such $\P$, retain the diagonal entries $\P_{i,j}$ and set
$T_{(i,j),(j,i)}=T_{(j,i),(i,j)}=\sqrt{\P_{i,j}\P_{j,i}}$, with all other off-diagonal entries zero. The resulting matrix is a direct sum of positive rank-one $2\times2$ blocks and nonnegative $1\times1$ blocks, has the prescribed partial traces, and attains equality. This proves~\eqref{eq-qot-classical-square-root-program}.

The function $(u,v)\mapsto(\sqrt u-\sqrt v)^2=u+v-2\sqrt{uv}$ is jointly convex on $\RR_+^2$. Equivalently, the hypograph constraint $z^2\leq uv$, $z\geq0$, is a rotated second-order cone, and the minimum in~\eqref{eq-qot-classical-socp} saturates it.

A bipartite state with pure first marginal $uu^*$ is supported on $\operatorname{span}\{u\}\otimes\CC^n$; if its second marginal is also pure, equal to $vv^*$, the only feasible coupling is $(u\otimes v)(u\otimes v)^*$. Therefore~\eqref{eq-qot-ground-cost-lift} gives
\(\mathcal W_{\C_0}^{\mathrm Q}(uu^*,vv^*) =2\sum_{i<j}\C_{0,i,j}|(s_{i,j}^-)^*(u\otimes v)|^2 =\sum_{i<j}\C_{0,i,j}|u_i v_j-u_j v_i|^2.\)
Taking $u=e_i$ and $v=e_j$ proves the Dirac identity. For unit off-diagonal costs, the last assertion follows from the complex Lagrange identity
$\sum_{i<j}|u_i v_j-u_j v_i|^2=\norm{u}^2\norm{v}^2-|u^*v|^2$.
\end{proof}
\begin{defn}[Classical square-root transport]\label{def-classical-square-root-transport}
Let $\Xx$ be a Polish space, let $\alpha,\beta\in\Pp(\Xx)$, and let $c:\Xx\times\Xx\to[0,+\infty)$ be measurable and symmetric. For $\pi\in\Couplings(\alpha,\beta)$, let $\pi^\top$ be the image of $\pi$ by $(x,y)\mapsto(y,x)$. If $\lambda$ is any finite nonnegative measure dominating $\pi$ and $\pi^\top$, set
\begin{equation}\label{eq-classical-square-root-action}
\mathscr H_c(\pi,\pi^\top)
\eqdef
\frac12\int_{\Xx\times\Xx}c(x,y)
\left(
\sqrt{\frac{\d\pi}{\d\lambda}(x,y)}
-
\sqrt{\frac{\d\pi^\top}{\d\lambda}(x,y)}
\right)^2\d\lambda(x,y),
\end{equation}
where the integral may equal $+\infty$, and define
\begin{equation}\label{eq-classical-square-root-transport}
\mathcal W_c^{\mathrm{cl}}(\alpha,\beta)
\eqdef
\inf_{\pi\in\Couplings(\alpha,\beta)}\mathscr H_c(\pi,\pi^\top).
\end{equation}
\end{defn}
The definition is independent of $\lambda$ because its pointwise integrand is positively $1$-homogeneous in the two densities. The same perspective formula shows that $\pi\mapsto\mathscr H_c(\pi,\pi^\top)$ is convex. With $\phi_H(r)=(\sqrt r-1)^2$ as in Section~\ref{sec-phi-div},
\begin{equation}\label{eq-classical-square-root-weighted-hellinger}
\mathscr H_c(\pi,\pi^\top)
=\frac12\Divergm_{\phi_H}(c\pi\mid c\pi^\top).
\end{equation}
\begin{prop}[Wasserstein comparison and topology]\label{prop-classical-square-root-wasserstein-comparison}
Let $(\Xx,\dist)$ be a compact metric space and $1\leq p<+\infty$. Then $\mathcal W_{\dist^p}^{\mathrm{cl}}$ is symmetric, nonnegative and vanishes exactly when its two arguments agree, and
\(\mathcal W_{\dist^p}^{\mathrm{cl}}(\alpha,\beta) \leq \Wass_p(\alpha,\beta)^p.\)
Moreover, for sequences $(\alpha_k)_k,(\beta_k)_k\subset\Pp(\Xx)$, \(\mathcal W_{\dist^p}^{\mathrm{cl}}(\alpha_k,\beta_k)\longrightarrow0 \quad\Longleftrightarrow\quad \Wass_q(\alpha_k,\beta_k)\longrightarrow0\)
for one, and hence every, $1\leq q<+\infty$. For $p=2$, one has the sharper two-sided comparison
\begin{equation}\label{eq-classical-square-root-wasserstein-bounds}
\frac12\Wass_1(\alpha,\beta)^2
\leq
\mathcal W_{\dist^2}^{\mathrm{cl}}(\alpha,\beta)
\leq
\Wass_2(\alpha,\beta)^2.
\end{equation}
\end{prop}
\begin{proof}
Fix $\pi\in\Couplings(\alpha,\beta)$, take the symmetric dominating measure $\lambda=\pi+\pi^\top$, and write $r=\d\pi/\d\lambda$ and $s=\d\pi^\top/\d\lambda$. Symmetry of $\dist$ gives
\(\mathscr H_{\dist^p}(\pi,\pi^\top) = \int\dist(x,y)^p r(x,y)\d\lambda - \int\dist(x,y)^p\sqrt{r(x,y)s(x,y)}\d\lambda \leq \int\dist(x,y)^p\d\pi.\)
Taking the infimum over $\pi$ proves the upper bound. Since $\dist^p$ is bounded and continuous, the maps $\pi\mapsto\dist^p\pi$ and $\pi\mapsto\dist^p\pi^\top$ are weakly continuous. Equation~\eqref{eq-classical-square-root-weighted-hellinger} and Proposition~\ref{prop-basic-phi-divergence-properties} therefore show that the action is weakly lower semicontinuous. Compactness of the coupling set gives a minimizer. If its value is zero, then $r=s$ away from the diagonal, while every measure supported on the diagonal is invariant under transposition. Hence $\pi=\pi^\top$, so its two marginals agree. Conversely, the diagonal coupling has zero cost when $\alpha=\beta$. Symmetry follows by transposing couplings.

If $\Wass_q(\alpha_k,\beta_k)\to0$ for one finite $q\geq1$, compactness and Corollary~\ref{cor-topol-wass} give $\Wass_p(\alpha_k,\beta_k)\to0$, so the upper bound implies $\mathcal W_{\dist^p}^{\mathrm{cl}}(\alpha_k,\beta_k)\to0$. Conversely, suppose that the latter values tend to zero. By compactness, any subsequence admits a further subsequence such that $\alpha_k\rightharpoonup\alpha$, $\beta_k\rightharpoonup\beta$, and optimal couplings $\pi_k\rightharpoonup\pi$. Lower semicontinuity gives $\mathscr H_{\dist^p}(\pi,\pi^\top)=0$, hence $\alpha=\beta$. Thus every subsequential limit has coinciding marginals. Compactness and Corollary~\ref{cor-topol-wass} then give $\Wass_q(\alpha_k,\beta_k)\to0$ for every finite $q\geq1$.

For every $1$-Lipschitz function $f$, antisymmetry of $\pi-\pi^\top$ gives \(\int f\d(\alpha-\beta) = \frac12\int(f(x)-f(y))(\sqrt r-\sqrt s)(\sqrt r+\sqrt s)\d\lambda.\)
Cauchy--Schwarz, $|f(x)-f(y)|\leq\dist(x,y)$, and $(\sqrt r+\sqrt s)^2\leq2(r+s)$ yield
$|\int f\d(\alpha-\beta)|\leq\sqrt{2\mathscr H_{\dist^2}(\pi,\pi^\top)}$. Kantorovich--Rubinstein duality~\eqref{eq-w1-metric} and then the infimum over $\pi$ give the lower bound.
\end{proof}
\paragraph{Entropic regularization and Bregman iterations.}
\begin{defn}[von Neumann quantum entropy]\label{def-von-neumann-quantum-entropy}
For a density matrix or positive semidefinite matrix $T$, the shifted von Neumann entropy functional used here is
\[
H(T)=\tr\!\left(T(\log T-\Id)\right),
\qquad
\nabla H(T)=\log T,
\]
with the convention $0\log0=0$ on eigenvalues. This is the convex negative quantum entropy; on trace-one states it differs from the physical entropy $-\tr(T\log T)$ by a sign and an additive constant.
\end{defn}
\eqllead{For $\epsilon>0$ define}{eq-qot-entropic-primal}{
\mathrm{QOT}_C^\epsilon(A,B)
=
\min_{T\succeq0}
\left\{
\tr(CT)+\epsilon H(T)
:\,
\operatorname{Tr}_B(T)=A,
\operatorname{Tr}_A(T)=B
\right\}.
}
\begin{prop}[Entropic quantum OT duality]\label{prop-qot-entropic-duality}
Assume $A\succ0$, $B\succ0$ and $\epsilon>0$. Then~\eqref{eq-qot-entropic-primal} has a unique positive minimizer. Its dual is
\begin{equation}\label{eq-qot-entropic-dual}
\mathrm{QOT}_C^\epsilon(A,B)
=
\max_{F\in\mathbb H_n,\ G\in\mathbb H_m}
\left\{
\tr(FA)+\tr(GB)
-\epsilon\,
\tr\exp\!\left(\frac{F\otimes\Id_m+\Id_n\otimes G-C}{\epsilon}\right)
\right\}.
\end{equation}
At optimality, primal and dual variables are linked by the Gibbs formula
\begin{equation}\label{eq-qot-gibbs-coupling}
T_e(F,G)
=
\exp\!\left(\frac{F\otimes\Id_m+\Id_n\otimes G-C}{\epsilon}\right),
\end{equation}
with $\operatorname{Tr}_B(T_e)=A$ and $\operatorname{Tr}_A(T_e)=B$.
\end{prop}
\begin{proof}
The feasible set is compact and nonempty, and it contains the positive definite point $A\otimes B$. The trace entropy is strictly convex on positive semidefinite matrices, hence the regularized primal has a unique minimizer. This minimizer is positive definite: if it had a non-trivial kernel, moving toward $A\otimes B$ would have entropy directional derivative $-\infty$ at the boundary while changing the linear cost only at finite rate. Slater's condition justifies the Lagrange dual computation. The Fenchel identity
\(\sup_{T\succeq0}\ \tr(YT)-\epsilon H(T) = \epsilon\,\tr\exp(Y/\epsilon)\)
is the matrix analogue of the scalar exponential conjugacy. Applying it to the Lagrangian of~\eqref{eq-qot-entropic-primal}, with
$Y=F\otimes\Id_m+\Id_n\otimes G-C$, gives~\eqref{eq-qot-entropic-dual}. The stationarity condition of this Fenchel identity gives~\eqref{eq-qot-gibbs-coupling}; differentiating the dual objective with respect to $F$ and $G$ yields the two marginal equations.
\end{proof}
Writing $K=\exp(-C/\epsilon)$, the objective in~\eqref{eq-qot-entropic-primal} differs by a constant from $\epsilon$ times the quantum KL divergence
\[
D_H(T|K)
=
\tr\!\left(T(\log T-\log K)-T+K\right).
\]
The exact quantum analogue of Sinkhorn is an implicit alternating Bregman projection scheme onto the affine marginal sets
\(\mathcal M_A=\{T\succeq0:\operatorname{Tr}_B(T)=A\}, \qquad \mathcal M_B=\{T\succeq0:\operatorname{Tr}_A(T)=B\}.\)
\begin{prop}[Exact Bregman projections]\label{prop-qot-bregman-projections}
Assume $A,B\succ0$ and let $K=\exp(-C/\epsilon)$. The minimizer of~\eqref{eq-qot-entropic-primal} is equivalently the minimizer of $D_H(T|K)$ over $\mathcal M_A\cap\mathcal M_B$. Moreover, if a current positive definite matrix has Gibbs form $T_e(F,G)$, then its Bregman projection onto $\mathcal M_A$ has the form $T_e(F^+,G)$, where $F^+$ is chosen so that $\operatorname{Tr}_B T_e(F^+,G)=A$. The projection onto $\mathcal M_B$ is analogous. Thus, when each one-block marginal equation is solved exactly, alternating Bregman projections are equivalent to alternating block maximization of the dual~\eqref{eq-qot-entropic-dual}.
\end{prop}
\begin{proof}
Since $\log K=-C/\epsilon$, the identity
\(\tr(CT)+\epsilon H(T) = \epsilon D_H(T|K)-\epsilon\tr(K)\)
holds, so the primal minimizer is the constrained Bregman projection of $K$ up to an additive constant. For the projection of a positive definite matrix $S$ onto $\mathcal M_A$, the affine set contains the positive definite point $A\otimes\Id_m/m$; the entropy derivative $\log T-\log S$ is singular at the boundary, so the projection lies in the interior of the positive cone. Its Lagrangian first variation is
\(\log T-\log S-\Lambda\otimes\Id_m=0\)
for a Hermitian multiplier $\Lambda$. Hence $T=\exp(\log S+\Lambda\otimes\Id_m)$. If $S=T_e(F,G)$, this is again of the form $T_e(F+\epsilon\Lambda,G)$. The multiplier is fixed by the marginal equation $\operatorname{Tr}_B(T)=A$. The same argument applies to $\mathcal M_B$. Finally, the first-order optimality condition for maximizing~\eqref{eq-qot-entropic-dual} over one block is exactly the corresponding marginal equation, so the Bregman and block-dual views coincide.
\end{proof}
In the non-commutative case, however, the exact block equations
\(\operatorname{Tr}_B T_e(F,G)=A, \qquad \operatorname{Tr}_A T_e(F,G)=B\)
\paragraph{Gurvits scaling and quantum Sinkhorn.}
It replaces the true Gibbs coupling~\eqref{eq-qot-gibbs-coupling} by the symmetric factorization
\begin{equation}\label{eq-qot-symmetric-scaling}
T_s(F,G)
=
\exp\!\left(\frac{Z}{2\epsilon}\right)\exp(-C/\epsilon)\exp\!\left(\frac{Z}{2\epsilon}\right)
=
(U\otimes V)K(U\otimes V),
\qquad
Z=F\otimes\Id_m+\Id_n\otimes G,
\end{equation}
where $U=\exp(F/(2\epsilon))$, $V=\exp(G/(2\epsilon))$ and $K=\exp(-C/\epsilon)$. For two matrices $M,N$, their commutator is $[M,N]\eqdef MN-NM$.

To make the Choi identification unambiguous, use the product bases fixed above and associate with $K$ the positive linear map $\mathcal K:\mathbb H_m\to\mathbb H_n$ and its Hilbert--Schmidt adjoint $\mathcal K^\star:\mathbb H_n\to\mathbb H_m$ through
\begin{equation}\label{eq-qot-choi-map}
\mathcal K(X)
\eqdef
\operatorname{Tr}_B\!\bigl(K(\Id_n\otimes X)\bigr),
\qquad
\mathcal K^\star(Y)
\eqdef
\operatorname{Tr}_A\!\bigl(K(Y\otimes\Id_m)\bigr).
\end{equation}
Equivalently, \([\mathcal K(X)]_{i,k} = \sum_{j,\ell}K_{(i,j),(k,\ell)}X_{\ell,j}, \qquad [\mathcal K^\star(Y)]_{j,\ell} = \sum_{i,k}K_{(i,j),(k,\ell)}Y_{k,i}.\)
With this convention, the marginal equations for the symmetric coupling are exactly
\(U\,\mathcal K(V^2)\,U=A, \qquad V\,\mathcal K^\star(U^2)\,V=B.\)
They can be enforced by the explicit congruence normalizations
\begin{equation}\label{eq-qot-gurvits-updates}
\begin{aligned}
R_V&=\mathcal K(V^2),
&
U&\leftarrow
R_V^{-1/2}\bigl(R_V^{1/2} A R_V^{1/2}\bigr)^{1/2}R_V^{-1/2},
\\
S_U&=\mathcal K^\star(U^2),
&
V&\leftarrow
S_U^{-1/2}\bigl(S_U^{1/2} B S_U^{1/2}\bigr)^{1/2}S_U^{-1/2}.
\end{aligned}
\end{equation}
These inverse square roots are well-defined when $K\succ0$ and $U,V,A,B\succ0$. This is Gurvits/operator scaling with prescribed targets.

\subsection{Dynamic Time Warping}\label{sec-dynamic-time-warping}
\paragraph{Discrete variational problem.}
Let $x=(x_i)_{i=1}^n$ and $y=(y_j)_{j=1}^m$ be two sequences in a feature space $\mathcal Z$, and set $\C_{i,j}=c(x_i,y_j)$ for a nonnegative cost $c:\mathcal Z\times\mathcal Z\to\RR_+$. A warping path is a sequence $\omega=((i_\ell,j_\ell))_{\ell=1}^L$ that starts at $(1,1)$, ends at $(n,m)$, and whose increments $(i_{\ell+1}-i_\ell,j_{\ell+1}-j_\ell)$ belong to $\{(1,0),(0,1),(1,1)\}$.

\begin{defn}[Dynamic time warping]\label{def-dynamic-time-warping}
The DTW value associated with the feature cost $c$ is
\begin{equation}\label{eq-dtw-variational}
\mathrm{DTW}_c(x,y)
\eqdef
\min_{\omega\in\Omega_{n,m}}
\sum_{\ell=1}^L c(x_{i_\ell},y_{j_\ell})
=
\min_{\omega\in\Omega_{n,m}}\dotp{A_\omega}{\C}.
\end{equation}
\end{defn}
\paragraph{Dynamic programming.}
\begin{prop}[DTW dynamic-programming recurrence]\label{prop-dtw-dynamic-programming}
Set $D_{0,0}=0$, $D_{i,0}=D_{0,j}=+\infty$ for $i,j>0$, and, for $1\leq i\leq n$, $1\leq j\leq m$, define
\begin{equation}\label{eq-dtw-recurrence}
D_{i,j}
=
\C_{i,j}+\min\{D_{i-1,j},D_{i,j-1},D_{i-1,j-1}\}.
\end{equation}
Then $D_{n,m}=\mathrm{DTW}_c(x,y)$. Storing one minimizing predecessor at every cell and backtracking from $(n,m)$ recovers an optimal path. The computation takes $O(nm)$ time and $O(nm)$ memory; if only the value is required, two rows give $O(\min\{n,m\})$ memory.
\end{prop}
\begin{proof}
Every admissible path reaching $(i,j)$ enters it from exactly one of $(i-1,j)$, $(i,j-1)$ or $(i-1,j-1)$. Removing the last cell therefore leaves an admissible path to one of these predecessors, while appending $(i,j)$ to any predecessor path gives an admissible path to $(i,j)$. Minimizing over the three mutually exhaustive last steps proves~\eqref{eq-dtw-recurrence} by induction on $i+j$. The complexity follows because each of the $nm$ cells performs constant work.
\end{proof}
\paragraph{Continuous time warping.}
\begin{defn}[Continuous dynamic time warping]\label{def-continuous-dtw}
Let $\Gamma_\uparrow$ contain pairs $(\phi,\psi)$ of absolutely continuous, nondecreasing surjections of $[0,1]$ onto itself. Define
\begin{equation}\label{eq-continuous-dtw}
\mathrm{CDTW}_c(x,y)
\eqdef
\inf_{(\phi,\psi)\in\Gamma_\uparrow}
\int_0^1
c\bigl(x(\phi(s)),y(\psi(s))\bigr)
\bigl(\dot\phi(s)+\dot\psi(s)\bigr)\d s.
\end{equation}
\end{defn}
\paragraph{Soft-DTW and the Sinkhorn analogy.}
Fix $\epsilon>0$. The hard minimum in~\eqref{eq-dtw-recurrence} is nonsmooth when several paths tie. Soft-DTW replaces it with the log-sum-exp soft minimum,
\begin{equation}\label{eq-dtw-softmin}
\operatorname{softmin}_\epsilon(r_1,\ldots,r_q)
=
-\epsilon\log\!\left(\sum_{k=1}^q e^{-r_k/\epsilon}\right),
\end{equation}
and defines $D_{0,0}^\epsilon=0$, $D_{i,0}^\epsilon=D_{0,j}^\epsilon=+\infty$ for $i,j>0$, together with
\begin{equation}\label{eq-soft-dtw-recurrence}
D_{i,j}^\epsilon
=
\C_{i,j}
+
\operatorname{softmin}_\epsilon
\bigl(D_{i-1,j}^\epsilon,D_{i,j-1}^\epsilon,D_{i-1,j-1}^\epsilon\bigr),
\qquad
\mathrm{sDTW}_{c,\epsilon}(x,y)=D_{n,m}^\epsilon.
\end{equation}
Equivalently, it is the free energy of all monotone paths,
\begin{equation}\label{eq-soft-dtw-partition}
\mathrm{sDTW}_{c,\epsilon}(x,y)
=
-\epsilon\log
\sum_{\omega\in\Omega_{n,m}}
\exp\!\left(-\frac{\dotp{A_\omega}{\C}}{\epsilon}\right).
\end{equation}
The Gibbs variational identity gives
\begin{equation}\label{eq-soft-dtw-variational}
\mathrm{sDTW}_{c,\epsilon}(x,y)
=
\min_{q\in\Delta(\Omega_{n,m})}
\left\{
\sum_{\omega}q_\omega\dotp{A_\omega}{\C}
-\epsilon H(q)
\right\}.
\end{equation}
Its unique minimizer is the Gibbs law
\begin{equation}\label{eq-soft-dtw-gibbs-law}
\PP_\epsilon(\omega)
=
\frac{\exp(-\dotp{A_\omega}{\C}/\epsilon)}
{\sum_{\omega'}\exp(-\dotp{A_{\omega'}}{\C}/\epsilon)},
\qquad
E_\epsilon
\eqdef
\nabla_\C\mathrm{sDTW}_{c,\epsilon}(x,y).
\end{equation}
Moreover,
\begin{equation}\label{eq-soft-dtw-zero-temperature-bound}
\mathrm{DTW}_c(x,y)-\epsilon\log\abs{\Omega_{n,m}}
\leq
\mathrm{sDTW}_{c,\epsilon}(x,y)
\leq
\mathrm{DTW}_c(x,y),
\end{equation}
because the partition sum is bounded below by its largest term and above by $\abs{\Omega_{n,m}}$ times that term. Hence $\mathrm{sDTW}_{c,\epsilon}\to\mathrm{DTW}_c$ as $\epsilon\downarrow0$; at zero temperature one may define $\mathrm{sDTW}_{c,0}\eqdef\mathrm{DTW}_c$ by this one-sided extension, while the formulas above remain restricted to $\epsilon>0$. Differentiating the finite log-partition function gives
\begin{equation}\label{eq-soft-dtw-expected-alignment}
E_\epsilon
=
\EE_{\omega\sim\PP_\epsilon}[A_\omega].
\end{equation}
Thus $(E_\epsilon)_{i,j}$ is the probability that a Gibbs path visits cell $(i,j)$. The matrix $E_\epsilon$ is a diffuse expected alignment and converges, when the hard optimum is unique, to its path-incidence matrix.

In direct analogy with the Sinkhorn divergence of Section~\ref{sec-sinkhorn-div}, the soft-DTW divergence is
\begin{equation}\label{eq-soft-dtw-divergence}
\overline{\mathrm{sDTW}}_{c,\epsilon}(x,y)
=
\mathrm{sDTW}_{c,\epsilon}(x,y)
-\frac12\mathrm{sDTW}_{c,\epsilon}(x,x)
-\frac12\mathrm{sDTW}_{c,\epsilon}(y,y).
\end{equation}

\section{Dynamic Optimal Transport}
\label{sec-dynamic-optimal-transport}
\subsection{Evolutions over the Space of Measures}
\paragraph{Lagrangian and Eulerian descriptions.}
At the particle level, this advection is governed by
\begin{equation}
\frac{\d x(t)}{\d t} = v_t(x(t)), \label{eq:lagrangian-advection}
\end{equation}
In the Eulerian description, the same motion is written directly on the evolving measure: the particle ODE becomes the PDE
\begin{equation}
\frac{\partial \alpha_t}{\partial t} + \diverg(v_t \alpha_t) = 0. \label{eq:eulerian-advection}
\end{equation}
For general measures, and in particular for empirical measures, it is understood in the distributional sense: for every $\varphi\in C_c^1((0,1)\times\RR^d)$,
\begin{equation}\label{eq:eulerian-advection-weak}
\int_0^1\!\int_{\RR^d}
\left(\partial_t\varphi(t,x)+\dotp{v_t(x)}{\nabla_x\varphi(t,x)}\right)
\d\alpha_t(x)\d t
=
0.
\end{equation}
This equation is obtained from~\eqref{eq:eulerian-advection} by integration by parts. Hence, for smooth positive densities, the classical and weak formulations are equivalent; the weak formulation is useful because it still makes sense for discrete measures whose particles evolve according to~\eqref{eq:lagrangian-advection}.

\begin{defn}[Admissible continuity-equation evolution]\label{def-admissible-continuity-evolution}
Let $\Omega=\RR^d$, or let $\Omega\subset\RR^d$ be a bounded Lipschitz domain. A narrowly continuous curve $(\alpha_t)_{t\in[0,1]}$ in $\Pp(\overline\Omega)$ and a jointly Borel velocity field $v:(t,x)\mapsto v_t(x)$ are admissible between $\alpha_0$ and $\alpha_1$ if
\(\int_0^1\!\int_{\overline\Omega}\norm{v_t(x)}^2\d\alpha_t(x)\d t<+\infty\)
and, for every $\varphi\in C^1([0,1]\times\overline\Omega)$, with compact spatial support when $\Omega=\RR^d$,
\begin{equation}\label{eq-continuity-endpoint-no-flux}
\int_0^1\!\int_{\overline\Omega}
\left(\partial_t\varphi+\dotp{v_t}{\nabla_x\varphi}\right)\d\alpha_t\d t
+\int_{\overline\Omega}\varphi(0,\cdot)\d\alpha_0
-\int_{\overline\Omega}\varphi(1,\cdot)\d\alpha_1=0.
\end{equation}
On a bounded domain, allowing test functions that do not vanish on $\partial\Omega$ encodes the zero-normal-flux condition $(\alpha_t v_t)\cdot n=0$ weakly.
\end{defn}
\begin{prop}[Lagrangian flows solve the continuity equation]\label{prop-lagrangian-flow-continuity}
Let $(T_t)_{t\in[0,1]}$ be a $C^1$ family of diffeomorphisms of $\RR^d$ and define $\alpha_t=(T_t)_\sharp\alpha_0$. Assume that the derivatives below are integrable with respect to $\alpha_0$, and define the Eulerian velocity field by
\(v_t(T_t(y))=\partial_t T_t(y).\)
Then $(\alpha_t,v_t)$ solves the continuity equation in the weak sense~\eqref{eq:eulerian-advection-weak}.
In particular, if $\alpha_0=\frac1n\sum_i\delta_{x_i(0)}$ is empirical, then $\alpha_t=\frac1n\sum_i\delta_{x_i(t)}$ is empirical as well, with particle velocities $\dot x_i(t)=v_t(x_i(t))$.
\end{prop}
\begin{proof}
Let $\varphi\in C_c^1((0,1)\times\RR^d)$. Since $\alpha_t=(T_t)_\sharp\alpha_0$, \(\frac{\d}{\d t}\int \varphi(t,x)\d\alpha_t(x) = \frac{\d}{\d t}\int \varphi(t,T_t(y))\d\alpha_0(y).\)
The chain rule gives
\[
\frac{\d}{\d t}\int \varphi(t,T_t(y))\d\alpha_0(y)
=
\int
\left(
\partial_t\varphi(t,T_t(y))
+\dotp{\nabla_x\varphi(t,T_t(y))}{\partial_t T_t(y)}
\right)
\d\alpha_0(y).
\]
Using the definition of $v_t$ and the push-forward relation, this equals
\[
\int
\left(
\partial_t\varphi(t,x)+\dotp{\nabla_x\varphi(t,x)}{v_t(x)}
\right)
\d\alpha_t(x).
\]
Integrating in time and using the boundary values of $\varphi$ gives~\eqref{eq:eulerian-advection-weak}.
\end{proof}
\paragraph{From measure evolutions to vector fields.}
For a given evolution $(\alpha_t)_t$, there are typically infinitely many velocity fields $v_t$ satisfying
\begin{equation}\label{eq:inverse-flow}
\partial_t\alpha_t+\diverg(\alpha_t v_t)=0.
\end{equation}
The linear space of vector fields that leave a measure $\alpha$ invariant is \(\Hh_\alpha=\enscond{v\in L^2(\alpha;\RR^d)}{\diverg(\alpha v)=0\ \text{in distributions}}.\)
\paragraph{Dacorogna--Moser inversion.}
With a fixed convention for the inverse Laplacian, this gives formally
\begin{equation}\label{eq:dacorogna-moser}
v_t=-\frac{1}{\rho_t}\nabla\Delta^{-1}(\partial_t\rho_t),
\end{equation}
If $\alpha_i=\rho_i\,\d x$ are smooth positive densities with the same total mass on a bounded connected domain $\Omega$, set
\(\alpha_t=(1-t)\alpha_0+t\alpha_1=\rho_t\,\d x, \qquad \rho_t=(1-t)\rho_0+t\rho_1.\)
Choose a time-independent flux $w$ satisfying
\(\diverg w=\rho_0-\rho_1, \qquad w\cdot n=0\quad\hbox{on }\partial\Omega,\)
for instance $w=-\nabla\phi$ with $\Delta\phi=\rho_1-\rho_0$ and Neumann boundary condition. Then
\(v_t=\frac{w}{\rho_t}\)
satisfies $\partial_t\rho_t+\diverg(\rho_t v_t)=0$. The flow $\partial_t T_t=v_t\circ T_t$, $T_0=\Id$, therefore transports $\rho_0\d x$ onto $\rho_t\d x$, and $T_1$ solves the prescribed-Jacobian problem $\rho_1(T_1(x))\det(\nabla T_1(x))=\rho_0(x)$.

\paragraph{Least-square inversion and gradient structure.}
A more robust choice, used implicitly in flow matching, optimal transport and Wasserstein gradient flows, is to select among all admissible velocities the one with smallest kinetic energy:
\begin{equation}\label{eq:least-square-field}
\umin{v}
\frac12\int_0^1\!\int_{\RR^d}\norm{v_t(x)}^2\,\d\alpha_t(x)\d t
\quad
\text{subject to}
\quad
\partial_t\alpha_t+\diverg(\alpha_t v_t)=0.
\end{equation}
\begin{prop}[Least-square velocities are gradients]\label{prop-least-square-gradient-velocity}
Assume that $\alpha_t=\rho_t\,\d x$ is a smooth positive density curve, that $\partial_t\rho_t$ has zero integral, and that boundary terms vanish. The minimizer of~\eqref{eq:least-square-field}, if it exists, is a gradient field
\(v_t=\nabla\phi_t,\)
where $\phi_t$, unique up to an additive constant on each connected component, solves the weighted Poisson equation
\begin{equation}\label{eq:least-square-field-explicit}
-\diverg(\rho_t\nabla\phi_t)=\partial_t\rho_t,
\qquad
v_t=-\nabla\Delta_{\alpha_t}^{-1}(\partial_t\alpha_t),
\quad
\Delta_{\alpha_t}\phi=\diverg(\alpha_t\nabla\phi).
\end{equation}
\end{prop}
\begin{proof}
Introduce a Lagrange multiplier $\phi_t$ for the continuity equation. The constrained problem has the formal saddle formulation
\[
\umin{v}\umax{\phi}
\int_0^1
\left[
\frac12\int_{\RR^d}\norm{v_t(x)}^2\,\d\alpha_t(x)
+\int_{\RR^d}\phi_t(x)\big(\diverg(\alpha_t v_t)(x)+\partial_t\alpha_t(x)\big)\d x
\right]\d t.
\]
Integrating by parts in the divergence term gives, for each $t$,
\[
\int \left(\frac12\norm{v_t}^2-\dotp{\nabla\phi_t}{v_t}\right)\d\alpha_t
+
\int \phi_t\,\partial_t\alpha_t.
\]
The pointwise minimizer in $v_t$ is therefore $v_t=\nabla\phi_t$. Substituting this into the constraint $\partial_t\rho_t+\diverg(\rho_t v_t)=0$ gives the weighted Poisson equation in~\eqref{eq:least-square-field-explicit}. The inverse notation is just a shorthand for solving this equation on zero-mean right-hand sides, modulo additive constants.
\end{proof}
\subsection{Benamou--Brenier Dynamic Formulation of OT}
\label{sec-benamou-brenier-dynamic}
\paragraph{Benamou--Brenier formulation.}
\begin{defn}[Benamou--Brenier action]\label{def-benamou-brenier-action}
For an admissible pair $(\alpha_t,v_t)$, define
\(\operatorname{Act}_2(\alpha,v) \eqdef \int_0^1\!\int_{\RR^d}\norm{v_t(x)}^2\d\alpha_t(x)\d t.\)
\end{defn}
\begin{thm}[Benamou--Brenier]\label{thm-benamou-brenier}
For probability measures $\alpha_0,\alpha_1\in\Pp_2(\RR^d)$,
\begin{equation}\label{eq:benamou-brenier}
\Wass_2^2(\alpha_0,\alpha_1)
=
\inf_{(\alpha_t,v_t)}
\int_0^1\!\int_{\RR^d}\norm{v_t(x)}^2\,\d\alpha_t(x)\d t,
\end{equation}
where the infimum is over $(\alpha_t,v_t)$ solving $\partial_t\alpha_t+\nabla\!\cdot(\alpha_t v_t)=0$ with $\alpha_{t=0}=\alpha_0$ and $\alpha_{t=1}=\alpha_1$. If $\alpha_0$ has a density and $T$ is the optimal Monge map $T_\sharp\alpha_0=\alpha_1$, a minimizing curve is
\begin{equation}\label{eq:static-to-dynamic}
\alpha_t=((1-t)\Id+tT)_\sharp\alpha_0,
\qquad
v_t\big((1-t)x+tT(x)\big)=T(x)-x
\quad\text{for $\alpha_0$-a.e.\ $x$ and a.e.\ $t\in(0,1)$}.
\end{equation}
\end{thm}
\begin{proof}
For the inequality ``dynamic $\leq$ static'', assume first that a Monge map $T$ exists and define $(\alpha_t,v_t)$ by~\eqref{eq:static-to-dynamic}. Since the Lagrangian velocity $T(x)-x$ is independent of $t$,
\(\int_0^1\!\int\norm{v_t}^2\,\d\alpha_t\d t = \int\norm{T(x)-x}^2\,\d\alpha_0(x),\)
so the dynamic cost is no larger than the static Monge cost. Without a Monge map, the same construction is made with an optimal coupling $\pi$: sample $(X,Y)\sim\pi$ and move along the straight path $\gamma_{X,Y}(t)=(1-t)X+tY$. This path measure has action $\int\norm{x-y}^2\d\pi(x,y)$; projecting path velocities onto their conditional mean at time $t$ gives an admissible Eulerian velocity with no larger action, so the dynamic value is no larger than the Kantorovich value.

Conversely, for a smooth deterministic path, take the flow $T_t$ defined by $\dot T_t=v_t\circ T_t$ and $T_0=\Id$. Then $\alpha_t=(T_t)_\sharp\alpha_0$ and $(T_1)_\sharp\alpha_0=\alpha_1$. Jensen's inequality gives
\(\norm{T_1(x)-x}^2 \leq \int_0^1\norm{v_t(T_t(x))}^2\d t.\)
After integration with respect to $\alpha_0$, the Monge cost is bounded above by the dynamic action. For general finite-energy solutions of the continuity equation, the superposition principle lifts the curve to a probability measure on absolutely continuous paths; applying Jensen's inequality pathwise gives a coupling of the endpoints whose quadratic cost is no larger than the action. Thus the Kantorovich value is bounded above by the dynamic value.
\end{proof}
\paragraph{Convex momentum-based reformulation.}
Although~\eqref{eq:benamou-brenier} is not jointly convex in $(\alpha_t,v_t)$, it becomes convex after replacing velocities by momenta. Given a velocity \(v\in L^2(\alpha;\RR^d)\), define the momentum
\(\omega\eqdef \alpha v, \qquad \omega(A)=\int_A v(x)\,\d\alpha(x),\)
\begin{defn}[Quadratic perspective action]\label{def-quadratic-perspective-action}
\begin{equation}\label{eq-quadratic-perspective}
J(a,m)
\eqdef
\begin{cases}
\norm{m}^2/a, & a>0,\\
0, & a=0\ \text{and}\ m=0,\\
+\infty, & a=0\ \text{and}\ m\neq0,
\end{cases}
\qquad
(a,m)\in[0,+\infty)\times\RR^d.
\end{equation}
If \(\lambda\) dominates \(\alpha\) and \(|\omega|\), define
\begin{equation}\label{eq-measure-perspective-action}
\mathbb J(\alpha,\omega)
\eqdef
\int
J\left(
\frac{\d\alpha}{\d\lambda}(x),
\frac{\d \omega}{\d\lambda}(x)
\right)\d\lambda(x).
\end{equation}
\end{defn}
In particular,
\[
\mathbb J(\alpha,\omega)<+\infty
\quad\Longleftrightarrow\quad
\omega=v\alpha\ \text{with}\ v\in L^2(\alpha;\RR^d),
\qquad
\mathbb J(\alpha,\omega)=\int\norm{v}^2\,\d\alpha.
\]
The Benamou--Brenier problem can now be written with the linear continuity equation in the pair \((\alpha_t,\omega_t)\):
\begin{equation}\label{eq:benamou-brenier-convex}
\Wass_2^2(\alpha_0,\alpha_1)
=
\inf_{\substack{\partial_t\alpha_t+\diverg \omega_t=0\\
\alpha_{t=0}=\alpha_0,\ \alpha_{t=1}=\alpha_1}}
\int_0^1 \mathbb J(\alpha_t,\omega_t)\,\d t.
\end{equation}
In the absolutely continuous case \(\alpha_t=\rho_t\,\d x\) and \(\omega_t=m_t\,\d x\), this becomes
\[
\inf_{\substack{\partial_t\rho_t+\diverg m_t=0\\
\rho_{t=0}\d x=\alpha_0,\ \rho_{t=1}\d x=\alpha_1}}
\int_0^1\!\int_{\RR^d}
J(\rho_t(x),m_t(x))\,\d x\,\d t
=
\inf_{\substack{\partial_t\rho_t+\diverg m_t=0\\
\rho_{t=0}\d x=\alpha_0,\ \rho_{t=1}\d x=\alpha_1}}
\int_0^1\!\int_{\RR^d}
\frac{\norm{m_t(x)}^2}{\rho_t(x)}\,\d x\,\d t,
\]
\paragraph{Dual Hamilton--Jacobi formulation.}
The momentum formulation also has a useful dual. It turns the least-action problem into a Hamilton--Jacobi subsolution inequality for a scalar potential, with equality on the part of space-time actually visited by the optimal curve.

\begin{prop}[Dual Benamou--Brenier problem]\label{prop-benamou-brenier-dual}
Assume, for simplicity, that the densities are smooth, compactly supported, and that boundary terms vanish. Then the convex dynamic value~\eqref{eq:benamou-brenier-convex} has the dual formulation
\begin{equation}\label{eq:benamou-brenier-dual}
\Wass_2^2(\alpha_0,\alpha_1)
=
\sup_{\phi}
\enscond{
\int_{\RR^d}\phi_1\,\d\alpha_1
-
\int_{\RR^d}\phi_0\,\d\alpha_0
}{
\partial_t\phi_t+\frac14\norm{\nabla\phi_t}^2\leq 0
},
\end{equation}
where the supremum is over smooth test potentials $\phi:[0,1]\times\RR^d\to\RR$. In the general metric-measure statement, the same formula is understood after relaxation: $\phi$ is a Hamilton--Jacobi subsolution, for instance in the viscosity sense, and finite-dimensional discretizations follow directly from Fenchel--Rockafellar duality.
If $(\rho,m)$ and $\phi$ are smooth primal and dual optimizers, then
\begin{equation}\label{eq:bb-primal-dual-relation}
m_t=\frac{\rho_t}{2}\nabla\phi_t,
\qquad
\partial_t\phi_t+\frac14\norm{\nabla\phi_t}^2=0
\quad\text{on }\{\rho_t>0\}.
\end{equation}
Equivalently, the optimal Eulerian velocity is $v_t=m_t/\rho_t=\nabla\phi_t/2$.
\end{prop}
\begin{proof}
Let $(\rho,m)$ satisfy the continuity equation and let $\phi$ be smooth. Multiplying the constraint by $\phi$ and integrating by parts gives
\[
\int_{\RR^d}\phi_1\,\d\alpha_1
-
\int_{\RR^d}\phi_0\,\d\alpha_0
=
\int_0^1\!\int_{\RR^d}
\left(\rho_t\,\partial_t\phi_t+\dotp{m_t}{\nabla\phi_t}\right)\d x\d t.
\]
If $\partial_t\phi_t+\norm{\nabla\phi_t}^2/4\leq0$, Young's inequality yields, pointwise, \(\rho\,\partial_t\phi+\dotp{m}{\nabla\phi} \leq -\frac{\rho}{4}\norm{\nabla\phi}^2+\dotp{m}{\nabla\phi} \leq J(\rho,m),\)
with the perspective convention~\eqref{eq-quadratic-perspective}. Hence the dual objective of every feasible $\phi$ is no larger than the action of every feasible primal pair.

Conversely, introduce $\phi$ as a Lagrange multiplier for $\partial_t\rho+\diverg m=0$. After integration by parts, and up to the fixed endpoint contribution in~\eqref{eq:benamou-brenier-dual}, the pointwise minimization over $(\rho,m)$ contains
\(J(\rho,m)-\rho\,\partial_t\phi-\dotp{m}{\nabla\phi}.\)
For fixed $\rho>0$, its minimum over $m$ is attained at $m=\rho\nabla\phi/2$ and equals $-\rho(\partial_t\phi+\norm{\nabla\phi}^2/4)$. The recession cases in~\eqref{eq-quadratic-perspective} give the same conclusion when \(\rho=0\). Thus the infimum over $\rho\geq0$ is finite exactly when the Hamilton--Jacobi inequality holds, in which case the pointwise infimum is $0$. Fenchel--Rockafellar duality then gives no duality gap. Equality in the two pointwise inequalities gives~\eqref{eq:bb-primal-dual-relation}.
\end{proof}
If $\gamma$ is any smooth curve with $\gamma(0)=x$ and $\gamma(1)=y$, then
\(\frac{\d}{\d t}\phi_t(\gamma(t)) = \partial_t\phi_t(\gamma(t))+\dotp{\nabla\phi_t(\gamma(t))}{\dot\gamma(t)} \leq \norm{\dot\gamma(t)}^2,\)
where the last inequality uses $\partial_t\phi+\norm{\nabla\phi}^2/4\leq0$ and Young's inequality. Integrating and minimizing over curves gives \(\phi_1(y)-\phi_0(x)\leq\norm{x-y}^2.\)
Thus $(-\phi_0,\phi_1)$ is a feasible static Kantorovich dual pair for the quadratic cost. At optimality the inequality is saturated on the endpoint pairs connected by the primal characteristics, which is the Eulerian version of complementary slackness.

\paragraph{Proximal splitting.}
The convex momentum formulation also explains the original Benamou--Brenier solver. Suppressing discretization indices, write \(U=(\rho,m)\), where \(m\) is the density of the momentum measure. Using the perspective~\eqref{eq-quadratic-perspective}, set
\[
\mathcal F(U)=\int_0^1\!\int_{\RR^d}J(\rho_t(x),m_t(x))\,\d x\,\d t,
\qquad
\mathcal C=\enscond{(\rho,m)}{\partial_t\rho+\diverg m=0,\ \rho_0\d x=\alpha_0,\ \rho_1\d x=\alpha_1},
\]
The convex Benamou--Brenier problem is therefore
\(\min_U\; \mathcal F(U)+\mathcal G(U).\)
On the resulting Hilbert space of unknowns, or formally for an $L^2$-type product structure, the two proximal operators are
\[
\prox_{\tau\mathcal F}(\bar U)
=
\argmin_U
\frac12\norm{U-\bar U}_{\mathsf H}^2+\tau\mathcal F(U),
\qquad
\prox_{\tau\mathcal G}(\bar U)
=
\argmin_{U\in\mathcal C}
\frac12\norm{U-\bar U}_{\mathsf H}^2.
\]
For the standard product metric, the first map is local in $(t,x)$: it is the proximal operator of the convex perspective $J$. The second map is the orthogonal projection onto the affine set defined by the divergence equation and endpoint constraints. Douglas--Rachford then alternates these two simple operations:
\[
\begin{aligned}
U^{(\ell+1)} &= \prox_{\tau\mathcal F}(Z^{(\ell)}),\\
\widetilde U^{(\ell+1)} &= \prox_{\tau\mathcal G}(2U^{(\ell+1)}-Z^{(\ell)}),\\
Z^{(\ell+1)} &= Z^{(\ell)}+\widetilde U^{(\ell+1)}-U^{(\ell+1)}.
\end{aligned}
\]
\paragraph{Path-space formulation.}\label{rem-bb-path-space}
Let $\Ss=C([0,1];\RR^d)$ be the space of continuous paths endowed with the uniform topology. For $t\in[0,1]$ define the evaluation map
\(e_t:\Ss\to\RR^d, \qquad e_t(\gamma)=\gamma(t).\)
The Benamou--Brenier cost admits the equivalent formulation
\[
\Wass_2^2(\alpha_0,\alpha_1)
=
\inf_{M\in\Pp(\Ss)}
\enscond{
\int_{\Ss}\!\int_0^1\norm{\dot\gamma(t)}^2\d t\,\d M(\gamma)
}{
(e_0)_\sharp M=\alpha_0,\ (e_1)_\sharp M=\alpha_1
}.
\]
Furthermore, for a.e.\ $t$, denoting $\dot e_t(\gamma)=\dot\gamma(t)$ on absolutely continuous paths, the conditional law of the velocity is deterministic:
\((e_t,\dot e_t)_\sharp M^*(\d x,\d q) = \alpha_t(\d x)\delta_{v_t^*(x)}(\d q),\)
where $v_t^*$ is the optimal velocity field in the Benamou--Brenier formulation. Hence $M^*$ concentrates on straight-line geodesics and, for a.e.\ $t$, assigns exactly one direction at $\alpha_t$-a.e.\ spatial point.

\subsection{Path-Space Schr\"odinger Problem}
\label{sec-path-space-schrodinger}
\paragraph{From least-action paths to random bridges.}
Throughout this section, both endpoint measures live on the same Polish state space $\X$; this is the setting relevant to dynamic transport and Brownian bridges, whereas the static coupling problem~\eqref{eq-entropic-generic} can also compare measures on different spaces. Fix a compatible complete bounded metric $d_\X$ and equip $\Om=C([0,1];\X)$ with the uniform metric.

\begin{defn}[Path-space least action]\label{def-path-space-transport}
Let $\Aa:\Om\to[0,+\infty]$ be a lower-semicontinuous path action. Define
\begin{equation}\label{eq-path-space-ot}
\operatorname{PT}_{\Aa}(\al,\be)
\eqdef
\inf_{M\in\Pp(\Om)}
\enscond{
\int_\Om \Aa(\omega)\,\d M(\omega)
}{
(e_0)_\sharp M=\al,\ (e_1)_\sharp M=\be
}.
\end{equation}
\end{defn}
For the quadratic Wasserstein geometry on $\RR^d$, one takes
\[
\Aa(\omega)=
\begin{cases}
\displaystyle \int_0^1 \norm{\dot\omega_t}^2\,\d t,
& \text{if }\omega\text{ is absolutely continuous},\\[.4em]
+\infty, & \text{otherwise}.
\end{cases}
\]
The endpoint cost induced by the action is
\begin{equation}\label{eq-path-action-endpoint-cost}
c_\Aa(x,y)
\eqdef
\inf_{\omega\in\Om}
\enscond{\Aa(\omega)}{e_0(\omega)=x,\ e_1(\omega)=y}.
\end{equation}
\begin{prop}[Endpoint reduction of path-space transport]\label{prop-path-space-ot-endpoint-reduction}
Assume that, for every $\delta>0$, there is a universally measurable map $(x,y)\mapsto\omega_\delta^{x,y}$ with endpoints $(x,y)$ and
$\Aa(\omega_\delta^{x,y})\leq c_\Aa(x,y)+\delta$ whenever $c_\Aa(x,y)<+\infty$. Then~\eqref{eq-path-space-ot} has the same value as the Kantorovich problem
\(\inf_{\pi\in\Couplings(\al,\be)} \int_{\X\times\X} c_\Aa(x,y)\,\d\pi(x,y).\)
Moreover, if $\pi^\star$ is an optimal endpoint coupling and an exact universally measurable minimizing selection $(x,y)\mapsto\omega^{x,y}$ exists $\pi^\star$-almost everywhere, then
\(M^\star = \int_{\X\times\X}\delta_{\omega^{x,y}}\,\d\pi^\star(x,y)\)
is an optimal path law.
\end{prop}
\begin{proof}
Let $M$ be any feasible path law and set $\pi=(e_0,e_1)_\sharp M$. Then $\pi\in\Couplings(\al,\be)$ and, by the definition of $c_\Aa$, \(\int_\Om \Aa(\omega)\,\d M(\omega) \geq \int_{\X\times\X} c_\Aa(x,y)\,\d\pi(x,y).\)
This proves that the path-space value is at least the Kantorovich value. Conversely, given a coupling $\pi$ and a measurable selection $(x,y)\mapsto\omega^{x,y}$ with action arbitrarily close to $c_\Aa(x,y)$, the mixture of Dirac path laws
\(M=\int\delta_{\omega^{x,y}}\,\d\pi(x,y)\)
has endpoints $\al$ and $\be$ and action equal, up to the selected approximation error, to $\int c_\Aa\,\d\pi$. Optimizing over $\pi$ and letting the approximation error vanish proves equality. If exact minimizing paths are selected for an optimal $\pi^\star$, the displayed $M^\star$ is optimal.
\end{proof}
\paragraph{Entropic path-space problem.}
Let $\mathsf R^\epsilon\in\Pp(\Om)$ be a reference path law, for instance a Brownian or Langevin dynamics at noise level $\epsilon$. The corresponding entropy projection defines the dynamic Schr\"odinger problem.

\begin{defn}[Schr\"odinger bridge]\label{def-schrodinger-bridge}
\begin{equation}\label{eq-schrodinger-path-space}
\mathrm{SB}_\epsilon(\al,\be)
\eqdef
\inf_{M\in\Pp(\Om)}
\enscond{
\epsilon\KL(M|\mathsf R^\epsilon)
}{
(e_0)_\sharp M=\al,\ (e_1)_\sharp M=\be
}.
\end{equation}
\end{defn}
\paragraph{Stochastic-control interpretation.}
For Brownian references, the entropy projection has a direct control meaning. To fix constants, suppose that $\mathsf R^\epsilon$ is the law of
\(\d X_t = \sqrt{\epsilon}\,\d B_t, \qquad X_0\sim\al,\)
Under the standard finite-entropy assumptions, a path law $M$ with the same initial marginal and $M\ll \mathsf R^\epsilon$ can be represented, in the weak sense, by an adapted drift $u_t$ of finite energy,
\(\d X_t = u_t\,\d t+\sqrt{\epsilon}\,\d B_t, \qquad X_0\sim\al.\)
Girsanov's formula then gives
\[
\KL(M|\mathsf R^\epsilon)
=
\frac{1}{2\epsilon}\EE_M\int_0^1 \norm{u_t}^2\,\d t,
\qquad
\epsilon\KL(M|\mathsf R^\epsilon)=\frac12\EE_M\int_0^1 \norm{u_t}^2\,\d t.
\]
If the reference initial law is not $\al$, an additional term $\epsilon\KL(\al|(\mathsf R^\epsilon)_0)$ appears, but it is fixed by the endpoint constraint. Hence, up to this fixed endpoint cost, the Schr\"odinger bridge can be read as
\(\inf_u \enscond{ \frac12\EE\int_0^1\norm{u_t}^2\,\d t }{ \d X_t=u_t\,\d t+\sqrt{\epsilon}\,\d B_t,\quad X_0\sim\al,\ X_1\sim\be },\)
If $h_t$ denotes the positive space-time harmonic Schr\"odinger factor associated with the terminal constraint, then the optimal feedback drift has the Doob-transform form
\(u_t^\star(x)=\epsilon\nabla\log h_t(x).\)
\paragraph{Viscous Benamou--Brenier formulations.}
Formally, for smooth positive densities the controlled Fokker--Planck equation is \(\partial_t\rho_t+\diverg(\rho_t v_t) = \frac{\sigma}{2}\Delta\rho_t,\)
one minimizes, among curves joining the prescribed endpoint densities, the kinetic action
\(\int_0^1\!\int \frac12\norm{v_t(x)}^2\rho_t(x)\d x\d t.\)
Equivalently, one can absorb the diffusion into the velocity by writing
\(u_t=v_t-\frac{\sigma}{2}\nabla\log\rho_t, \qquad \partial_t\rho_t+\diverg(\rho_t u_t)=0.\)
Expanding $v_t=u_t+\frac{\sigma}{2}\nabla\log\rho_t$ and using the continuity equation for $(\rho_t,u_t)$ gives
\begin{align*}
\int_0^1\!\int
\frac12\norm{v_t}^2\rho_t\,\d x\,\d t
&=
\int_0^1\!\int
\left(
\frac12\norm{u_t}^2
+
\frac{\sigma^2}{8}\norm{\nabla\log\rho_t}^2
\right)\rho_t\,\d x\,\d t
\\
&\quad+
\frac{\sigma}{2}
\left[
\int\rho_1\log\rho_1\,\d x
-
\int\rho_0\log\rho_0\,\d x
\right].
\end{align*}
Since the entropy term depends only on the prescribed endpoints, the same minimizers are obtained from the modified Benamou--Brenier action
\[
\int_0^1\!\int
\left(
\frac12\norm{u_t(x)}^2
+
\frac{\sigma^2}{8}\norm{\nabla\log\rho_t(x)}^2
\right)
\rho_t(x)\d x\d t.
\]
Thus the Schr\"odinger bridge is a least-action interpolation with both transport kinetic energy and a Fisher-information penalty. If one instead writes the viscous equation with diffusion coefficient $\sigma\Delta\rho_t$, the same formula is obtained after replacing $\sigma$ above by $2\sigma$, so the Fisher coefficient becomes $\sigma^2/2$.

We write
\[
\mathsf R^\epsilon(\d\omega)
=
\int \mathsf R^{\epsilon,x,y}(\d\omega)\,
\mathsf R^\epsilon_{0,1}(\d x,\d y),
\qquad
\mathsf R^\epsilon_{0,1}\eqdef(e_0,e_1)_\sharp\mathsf R^\epsilon,
\]
where $\mathsf R^{\epsilon,x,y}$ is the reference bridge conditioned on endpoints $(x,y)$. Similarly, for any feasible $M$, write
\(M(\d\omega) = \int M^{x,y}(\d\omega)\,\pi(\d x,\d y), \qquad \pi\eqdef(e_0,e_1)_\sharp M.\)
\begin{prop}[Endpoint reduction of the Schr\"odinger problem]\label{prop-schrodinger-endpoint-reduction}
Assume that the regular conditional laws above exist and that the relative-entropy chain rule applies, with value $+\infty$ when absolute continuity fails. Then
\begin{equation}\label{eq-schrodinger-static-endpoint}
\mathrm{SB}_\epsilon(\al,\be)
=
\inf_{\pi\in\Couplings(\al,\be)}
\epsilon\KL(\pi|\mathsf R^\epsilon_{0,1}).
\end{equation}
For a fixed endpoint coupling $\pi$ with finite $\KL(\pi|\mathsf R^\epsilon_{0,1})$, the minimizing path law is the mixture of reference bridges
\begin{equation}\label{eq-schrodinger-bridge-mixture}
M^\pi
=
\int \mathsf R^{\epsilon,x,y}\,\d\pi(x,y).
\end{equation}
\end{prop}
\begin{proof}
If $M$ has finite relative entropy with respect to $\mathsf R^\epsilon$, then $\pi=(e_0,e_1)_\sharp M$ is necessarily absolutely continuous with respect to $\mathsf R^\epsilon_{0,1}$. The chain rule for relative entropy gives
\begin{equation}\label{eq-kl-chain-rule-path}
\KL(M|\mathsf R^\epsilon)
=
\KL(\pi|\mathsf R^\epsilon_{0,1})
+
\int_{\X\times\X}
\KL(M^{x,y}|\mathsf R^{\epsilon,x,y})
\,\d\pi(x,y).
\end{equation}
The second term is nonnegative and vanishes exactly when $M^{x,y}=\mathsf R^{\epsilon,x,y}$ for $\pi$-almost every $(x,y)$. Thus, once an endpoint coupling $\pi$ with finite $\KL(\pi|\mathsf R^\epsilon_{0,1})$ is fixed, the best path law is~\eqref{eq-schrodinger-bridge-mixture}, and the remaining minimization is precisely~\eqref{eq-schrodinger-static-endpoint}. If no such endpoint coupling exists, both sides are $+\infty$.
\end{proof}
\paragraph{Brownian bridges and Sinkhorn couplings.}
For $\X=\RR^d$, take $\mathsf R^\epsilon$ to be a Brownian reference dynamics. With the convention $\d X_t=\sqrt{\epsilon}\,\d B_t$ used above, its unit-time transition density is
\[
p_\epsilon(x,y)
=
(2\pi\epsilon)^{-d/2}
\exp\!\left(-\frac{\norm{x-y}^2}{2\epsilon}\right).
\]
Thus multiplying the endpoint entropy by $\epsilon$ produces the static cost $\norm{x-y}^2/2$; the convention $e^{-\norm{x-y}^2/\epsilon}$ used for quadratic Sinkhorn is recovered by replacing the Brownian noise variance by $\epsilon/2$. More generally, suppose that for two reference probability measures $\bar\alpha,\bar\beta$ the endpoint prior has the Gibbs form, with $0<Z_\epsilon<+\infty$,
\[
\mathsf R^\epsilon_{0,1}(\d x,\d y)
=
\frac1{Z_\epsilon}
\exp\!\left(-\frac{c(x,y)}{\epsilon}\right)
\bar\alpha(\d x)\bar\beta(\d y),
\]
and assume $\al\ll\bar\alpha$ and $\be\ll\bar\beta$. Then, whenever the terms are finite, the relative-entropy chain rule gives, for every $\pi\in\Couplings(\al,\be)$,
\[
\epsilon\KL(\pi|\mathsf R^\epsilon_{0,1})
=
\int c(x,y)\,\d\pi(x,y)
+
\epsilon\KL(\pi|\al\otimes\be)
+
\epsilon\KL(\al|\bar\alpha)
+
\epsilon\KL(\be|\bar\beta)
+
\epsilon\log Z_\epsilon.
\]
The last three terms are independent of $\pi$. For Brownian motion started from $\al$, one has $\mathsf R^\epsilon_{0,1}(\d x,\d y)=\al(\d x)p_\epsilon(x,y)\d y$. If $\be=b\,\d y$, the same calculation reads
\[
\epsilon\KL(\pi|\mathsf R^\epsilon_{0,1})
=
\frac12\int\norm{x-y}^2\d\pi
+
\epsilon\KL(\pi|\al\otimes\be)
+
\epsilon\int\log b\,\d\be
+
\frac{d\epsilon}{2}\log(2\pi\epsilon),
\]
under the present noise convention. Thus the required domination is $\be\ll\d y$ (and $\al\ll(\mathsf R^\epsilon)_0$ if the Brownian reference has another initial law); no absolute-continuity relation between $\al$ and $\be$ is involved. The optimal static coupling $\pi_\epsilon^\star$ is the endpoint law of the most likely controlled noisy dynamics, while the associated path law is
\(M_\epsilon^\star = \int \mathsf R^{\epsilon,x,y}\,\d\pi_\epsilon^\star(x,y).\)
\subsection{Generalized Dynamic Wasserstein Distances}
\label{sec-bb-extensions}
\label{sec-generalized-dynamic-wasserstein-distances}
The quadratic Benamou--Brenier formula is only one instance of a broader fixed-mass dynamic language. An arbitrary action first defines only a path value; metric properties require the symmetry, homogeneity and closure hypotheses isolated below.

\paragraph{Path actions.}
In the mass-preserving Euclidean setting, the basic input is an instantaneous action $\mathbb A(\alpha,w)$, where \(\alpha\) is the current measure and \(w\) is an admissible velocity representative. Its endpoint geometry is defined dynamically.

\begin{defn}[Generalized dynamic action value]\label{def-generalized-dynamic-action-distance}
Let $\mathbb A$ be a nonnegative extended-valued action on measure--velocity pairs. Define
\begin{equation}\label{eq-generalized-action-length-distance}
\mathsf E_{\mathbb A}(\alpha_0,\alpha_1)
\eqdef
\inf_{\alpha_t,v_t}
\left\{
\int_0^1 \mathbb A(\alpha_t,v_t)\,\d t
:
\partial_t\alpha_t+
\diverg(\alpha_t v_t)=0,\
\alpha_{t=0}=\alpha_0,\
\alpha_{t=1}=\alpha_1
\right\}.
\end{equation}
\end{defn}
Equivalently, one may quotient by velocity fields that induce the same first-order variation of the measure. The value~\eqref{eq-generalized-action-length-distance} need not be symmetric or separate measures, and its square root need not satisfy the triangle inequality.

\paragraph{Quadratic, or Riemannian, tangent actions.}
If the polarization of \(\mathbb A\) is represented by a positive self-adjoint operator
\(Q_\alpha:L^2(\alpha;\RR^d)\to L^2(\alpha;\RR^d)\),
\[
\mathbb A(\alpha;w,z)
=
\left\langle Q_\alpha w,z\right\rangle_{L^2(\alpha)},
\qquad
\mathbb A(\alpha,w)=\left\langle Q_\alpha w,w\right\rangle_{L^2(\alpha)},
\]
and if the quadratic form is nondegenerate after quotienting velocities that induce the same measure variation, then the associated quadratic path value is
\begin{equation}\label{eq-general-quadratic-tangent-action}
\mathsf D_Q^2(\alpha_0,\alpha_1)
=
\mathsf E_{\mathbb A}(\alpha_0,\alpha_1)
=
\inf_{\substack{\partial_t\alpha_t+\diverg(\alpha_t v_t)=0\\
\alpha_{t=0}=\alpha_0,\ \alpha_{t=1}=\alpha_1}}
\int_0^1
\left\langle Q_{\alpha_t}v_t,v_t\right\rangle_{L^2(\alpha_t)}
\d t.
\end{equation}
When, in addition, the dynamic problem is sequentially closed and attains finite infima as in Proposition~\ref{prop-homogeneous-dynamic-action-distance}, $\mathsf D_Q$ is an extended distance. The usual $\Wass_2$ geometry corresponds to $Q_\alpha=\Id$ in this simplified notation.

\paragraph{Local velocity actions.}
A velocity action is specified by a pointwise integrand
\(A:[0,+\infty)\times\RR^d\to[0,+\infty], \qquad (a,w)\mapsto A(a,w),\)
where $a\in\RR_+$ is a density value and $w\in\RR^d$ is a velocity value, and defines
\begin{equation}\label{eq-local-velocity-action}
\mathbb A(\alpha,w)
=
\int A\left(\frac{\d\alpha}{\d\lambda}(x),w(x)\right)\d\lambda(x).
\end{equation}
For a fixed reference $\lambda$, this covers density-dependent mobilities and congestion constraints. If, in addition, $A$ is positively $1$-homogeneous in its first variable, $A(\eta a,w)=\eta A(a,w)$ for $\eta\geq0$, then the same formula is intrinsic: replacing $\lambda$ by another dominating measure gives the same value. The usual Benamou--Brenier action is the model case
\begin{equation}\label{eq-bb-velocity-action}
A_2(a,w)=a\norm{w}^2,
\qquad
\mathbb A(\alpha,w)=\int\norm{w}^2\d\alpha.
\end{equation}
\paragraph{Homogeneous momentum actions.}\label{rem-generalized-bb}
When the local description is written with the same reference $\lambda$, so that $\alpha=a\lambda$ and $\omega=m\lambda$, the pointwise momentum perspective is
\begin{equation}\label{eq-general-momentum-perspective}
J_A(a,m)
\eqdef
\begin{cases}
A(a,m/a), & a>0,\\
0, & a=0\ \text{and}\ m=0,\\
+\infty, & a=0\ \text{and}\ m\neq0,
\end{cases}
\end{equation}
and the measure action relative to $\lambda$ is
\begin{equation}\label{eq-general-measure-momentum-action}
\mathbb J_{A,\lambda}(\alpha,\omega)
\eqdef
\int
J_A\!\left(
\frac{\d\alpha}{\d\lambda},
\frac{\d\omega}{\d\lambda}
\right)\d\lambda,
\end{equation}
For $A_2(a,w)=a\norm{w}^2$, one recovers the quadratic perspective
\(J_2(a,m)=\frac{\norm{m}^2}{a},\)
\begin{prop}[Concave mobilities give convex momentum actions]\label{prop-momentum-perspective-convexity}
Let \(I\subset[0,+\infty)\) be a convex interval, let \(\theta:I\to[0,+\infty)\) be concave, and let \(L:\RR^d\to[0,+\infty]\) be convex with \(L(0)=0\). Define, on the set where \(\theta(a)>0\),
\begin{equation}\label{eq-concave-mobility-perspective}
J_{\theta,L}(a,m)
\eqdef
\theta(a)L\!\left(\frac{m}{\theta(a)}\right),
\qquad
A_{\theta,L}(a,w)
\eqdef
J_{\theta,L}(a,aw)
=
\theta(a)L\!\left(\frac{aw}{\theta(a)}\right),
\end{equation}
and extend \(J_{\theta,L}\) to the boundary by lower semicontinuity. Then \(J_{\theta,L}\) is convex in \((a,m)\). In particular, this single construction contains
\(A(a,w)=aL(w) \quad\text{by taking }\theta(a)=a,\)
and the concave-mobility quadratic action
\(A(a,w)=\frac{a^2\norm{w}^2}{\theta(a)} \quad\text{by taking }L(u)=\norm u^2.\)
\end{prop}
\begin{proof}
The perspective \(P(s,m)=sL(m/s)\) of a convex function is convex on \(s>0\). Since \(L(0)=0\), it is also nonincreasing in the scalar \(s\) for fixed \(m\): if \(s_2\geq s_1>0\), then
\(L(m/s_2) = L\big((s_1/s_2)(m/s_1)\big) \leq (s_1/s_2)L(m/s_1),\)
hence \(P(s_2,m)\leq P(s_1,m)\). Let \(a_\zeta=(1-\zeta)a_0+\zeta a_1\) and \(m_\zeta=(1-\zeta)m_0+\zeta m_1\). By concavity of \(\theta\), \(\theta(a_\zeta) \geq (1-\zeta)\theta(a_0)+\zeta\theta(a_1).\)
Using the monotonicity in \(s\), and then the convexity of \(P\), gives
\[
J_{\theta,L}(a_\zeta,m_\zeta)
=
P(\theta(a_\zeta),m_\zeta)
\leq
P((1-\zeta)\theta(a_0)+\zeta\theta(a_1),m_\zeta)
\leq
(1-\zeta)J_{\theta,L}(a_0,m_0)+\zeta J_{\theta,L}(a_1,m_1).
\]
The boundary cases follow by the chosen lower-semicontinuous extension.
\end{proof}
\begin{prop}[Homogeneous dynamic actions define distances]\label{prop-homogeneous-dynamic-action-distance}
Fix a reference measure \(\lambda\), omitted from the notation only in the intrinsic case where \(\mathbb J_{A,\lambda}\) does not depend on \(\lambda\). Assume that the momentum perspective \(J_A\) defined in~\eqref{eq-general-momentum-perspective} is lower semicontinuous, convex in \((a,m)\), even in \(m\), and satisfies \(J_A(a,0)=0\) on its effective density domain. Assume moreover that, for some \(r>1\),
\(J_A(a,\xi m)=|\xi|^rJ_A(a,m) \qquad \text{for all admissible }a,\ \text{all }m,\ \text{and all }\xi\in\RR,\)
and that, for each admissible $a$, \(J_A(a,m)=0\) if and only if \(m=0\). Equivalently, the evenness, homogeneity and nondegeneracy assumptions translate, away from zero density, into
\(A(a,-w)=A(a,w), \qquad A(a,\xi w)=|\xi|^rA(a,w), \qquad A(a,w)=0 \Longleftrightarrow w=0 \qquad (a>0\text{ admissible}).\)
For a mass $M\geq0$, define the effective state space
\[
\mathcal S_{A,\lambda}^M
\eqdef
\left\{\alpha=a\lambda:\ \alpha(\mathcal X)=M,
J_A(a(x),0)=0\ \lambda\text{-a.e.}\right\}.
\]
For endpoints in $\mathcal S_{A,\lambda}^M$, define
\[
\mathsf D_{A,\lambda}(\alpha_0,\alpha_1)
\eqdef
\inf_{\substack{\partial_t\alpha_t+\diverg\omega_t=0\\
\alpha_{t=0}=\alpha_0,\ \alpha_{t=1}=\alpha_1}}
\left(
\int_0^1\mathbb J_{A,\lambda}(\alpha_t,\omega_t)\,\d t
\right)^{1/r}.
\]
Assume finally that the relaxed dynamic problem is sequentially closed and attains its infimum whenever the value is finite. Then \(\mathsf D_{A,\lambda}\) is an extended distance on $\mathcal S_{A,\lambda}^M$, and a genuine distance on each finite-action component: it is symmetric, satisfies the triangle inequality, and \(\mathsf D_{A,\lambda}(\alpha_0,\alpha_1)=0\) only when \(\alpha_0=\alpha_1\). In the intrinsic case, the reference can be omitted and the same conclusion holds on the corresponding effective fixed-mass state space; we write \(\mathsf D_A\).
\end{prop}
\begin{proof}
The constant curve \(\alpha_t=\alpha_0\), \(\omega_t=0\), gives \(\mathsf D_{A,\lambda}(\alpha_0,\alpha_0)=0\). Conversely, if \(\mathsf D_{A,\lambda}(\alpha_0,\alpha_1)=0\), attainment provides a relaxed zero-action minimizer, which satisfies
\(\int_0^1 \mathbb J_{A,\lambda}(\alpha_t,\omega_t)\,\d t=0.\)
Since \(\mathbb J_{A,\lambda}\geq0\), this implies \(\omega_t=0\) for a.e.\ \(t\). The continuity equation then gives \(\partial_t\alpha_t=0\) in the sense of distributions, so \(\alpha_0=\alpha_1\).

Symmetry follows by time reversal. If \((\alpha_t,\omega_t)_{t\in[0,1]}\) connects \(\alpha_0\) to \(\alpha_1\), then
\(\widetilde\alpha_t=\alpha_{1-t}, \qquad \widetilde\omega_t=-\omega_{1-t}\)
connects \(\alpha_1\) to \(\alpha_0\), satisfies the same continuity equation, and has the same action because \(J_A\) is even in \(m\).

Let \((\alpha^1,\omega^1)\) connect \(\alpha_0\) to \(\alpha_1\) with action \(E_1\), and let \((\alpha^2,\omega^2)\) connect \(\alpha_1\) to \(\alpha_2\) with action \(E_2\). For \(\tau\in(0,1)\), concatenate the two curves after the time changes \(s=t/\tau\) and \(s=(t-\tau)/(1-\tau)\):
\[
(\alpha_t,\omega_t)=
\begin{cases}
\bigl(\alpha^1_{t/\tau},\,\tau^{-1}\omega^1_{t/\tau}\bigr), & 0\leq t\leq \tau,\\[.2em]
\bigl(\alpha^2_{(t-\tau)/(1-\tau)},\,(1-\tau)^{-1}\omega^2_{(t-\tau)/(1-\tau)}\bigr), & \tau\leq t\leq1.
\end{cases}
\]
The rescaled curve is admissible from \(\alpha_0\) to \(\alpha_2\). By \(r\)-homogeneity in the momentum variable, its action is
\(\tau^{1-r}E_1+(1-\tau)^{1-r}E_2\).
Writing \(X=E_1^{1/r}\) and \(Y=E_2^{1/r}\), the optimal allocation is \(\tau=X/(X+Y)\), with the evident limiting convention if \(X+Y=0\), and gives
\[
\inf_{\tau\in(0,1)}
\left\{\tau^{1-r}E_1+(1-\tau)^{1-r}E_2\right\}
=
(X+Y)^r.
\]
Hence \(\mathsf D_{A,\lambda}(\alpha_0,\alpha_2)\leq E_1^{1/r}+E_2^{1/r}\). Taking \(E_1\) and \(E_2\) arbitrarily close to the two optimal action values yields
\(\mathsf D_{A,\lambda}(\alpha_0,\alpha_2) \leq \mathsf D_{A,\lambda}(\alpha_0,\alpha_1)+\mathsf D_{A,\lambda}(\alpha_1,\alpha_2),\)
with the usual extended-real convention when one of the distances is infinite.
\end{proof}
\paragraph{Concave-mobility actions.}
\begin{defn}[Concave-mobility dynamic distance]\label{def-concave-mobility-distance}
Let \(I\subset[0,+\infty)\) be a closed convex interval and let \(\theta:I\to[0,+\infty)\) be continuous and concave, with $\theta>0$ on the relative interior of $I$. Define the closed momentum perspective
\[
J_\theta(a,m)
\eqdef
\begin{cases}
\norm{m}^2/\theta(a), & a\in I,\ \theta(a)>0,\\
0, & a\in I,\ \theta(a)=0,\ m=0,\\
+\infty, & \text{otherwise}.
\end{cases}
\]
The corresponding velocity action is \(A_\theta(a,w)=J_\theta(a,aw), \qquad A_\theta(a,w)=\frac{a^2\norm{w}^2}{\theta(a)}\quad\text{when }\theta(a)>0.\)
For a fixed reference measure $\lambda$ and a mass $M\geq0$, set
\[
\mathcal S_{\theta,\lambda}^M
\eqdef
\left\{\alpha=a\lambda:\alpha(\mathcal X)=M,\ a(x)\in I\ \lambda\text{-a.e.}\right\},
\qquad
\mathsf W_{\theta,\lambda}(\alpha_0,\alpha_1)
\eqdef
\mathsf D_{A_\theta,\lambda}(\alpha_0,\alpha_1).
\]
\end{defn}
If \(\alpha=a\lambda\), the induced squared tangent action is \(\mathbb A_{\theta,\lambda}(\alpha,w) = \int A_\theta(a(x),w(x))\,\d\lambda(x) = \int \frac{a(x)^2}{\theta(a(x))}\norm{w(x)}^2\,\d\lambda(x),\)
where the integrand uses the closed convention in Definition~\ref{def-concave-mobility-distance}; the action is $+\infty$ if $\alpha\not\ll\lambda$ or $a(x)\notin I$ on a set of positive $\lambda$-measure. Equivalently, on the set where \(\theta(a(x))>0\),
\(\mathbb A_{\theta,\lambda}(\alpha,w) = \int \frac{a(x)}{\theta(a(x))}\norm{w(x)}^2\,\d\alpha(x).\)
Hence this is a local Riemannian case, in the sense of~\eqref{eq-general-quadratic-tangent-action}, whenever the multiplier \(a(x)/\theta(a(x))\) is finite and positive. The associated tensor is the multiplication operator
\begin{equation}\label{eq-concave-mobility-q-alpha}
\big(Q_{\theta,\lambda,\alpha}w\big)(x)
=
\frac{a(x)}{\theta(a(x))}w(x),
\qquad
\alpha=a\lambda,
\end{equation}
defined \(\alpha\)-a.e.\ on the set where \(\theta(a)>0\). Except for linear mobilities \(\theta(a)=ca\), and in particular the normalized case \(\theta(a)=a\) which recovers \(\Wass_2\), the pointwise velocity action \(A_\theta(a,w)\) is not positively \(1\)-homogeneous in the density variable \(a\).

\begin{prop}[Concave-mobility dynamic distances]\label{prop-concave-mobility-distance}
For a fixed reference measure \(\lambda\), and under the standard compactness hypotheses ensuring existence of relaxed minimizers, \(\mathsf W_{\theta,\lambda}=\mathsf D_{A_\theta,\lambda}\) is an extended distance on $\mathcal S_{\theta,\lambda}^M$ and a genuine metric on each finite-action component.
\end{prop}
\begin{proof}
Proposition~\ref{prop-momentum-perspective-convexity}, applied with \(L(u)=\norm u^2\), gives convexity on the positive-mobility region. Continuity of $\theta$, the explicit zero-mobility convention, and the barrier outside the closed interval $I$ give its lower-semicontinuous extension $J_\theta$. It is even in \(m\), satisfies \(J_\theta(a,0)=0\) on $I$, and obeys
\(J_\theta(a,\xi m)=|\xi|^2J_\theta(a,m)\).
Moreover \(J_\theta(a,m)=0\) if and only if \(m=0\), with the boundary convention used in its definition. For the fixed reference \(\lambda\), the hypotheses of Proposition~\ref{prop-homogeneous-dynamic-action-distance} therefore hold with \(r=2\); the existence assumption is supplied by the compactness hypothesis in the statement. That proposition gives symmetry, separation and the triangle inequality for \(\mathsf D_{A_\theta,\lambda}\), hence for \(\mathsf W_{\theta,\lambda}\).
\end{proof}
\paragraph{Dynamic spectral Wasserstein distances.}
\begin{defn}[Dynamic spectral Wasserstein distance]\label{def-dynamic-spectral-wasserstein}
Let $\gamma$ be a monotone spectral gauge on $\mathbb S_+^d$. Define
\begin{equation}\label{eq-spectral-tangent-action}
\mathbb A_\gamma(\alpha,v)
\eqdef
\gamma\!\left(\int v(x)v(x)^\top\d\alpha(x)\right).
\end{equation}
\eqllead{The associated distance is}{eq-dynamic-spectral-wasserstein}{
\mathsf W_{\gamma,\mathrm{dyn}}^2(\alpha_0,\alpha_1)
\eqdef
\mathsf E_{\mathbb A_\gamma}(\alpha_0,\alpha_1).
}
\end{defn}
In density--momentum variables, this corresponds to the measure action
\[
\mathbb J_\gamma(\alpha,\omega)
=
\gamma\!\left(\int
\left(\frac{\d\omega}{\d\alpha}\right)
\left(\frac{\d\omega}{\d\alpha}\right)^\top
\d\alpha\right),
\]
or, when \(\alpha=\rho\,\d x\) and \(\omega=m\,\d x\),
\[
\mathbb J_\gamma(\rho,m)
=
\gamma\!\left(\int \frac{m(x)m(x)^\top}{\rho(x)}\,\d x\right).
\]
In that case the velocity and momentum densities are \(A_{\mathrm{lin}}(a,w)=ca\norm w^2, \qquad J_{\mathrm{lin}}(a,m)=c\frac{\norm{m}^2}{a},\)
\begin{prop}[Static/dynamic equality for spectral Wasserstein]\label{prop-static-dynamic-spectral-wasserstein}
Let $\gamma$ be a monotone spectral gauge and let $\alpha_0,\alpha_1\in\Pp_2(\RR^d)$. Then
\(\mathsf W_{\gamma,\mathrm{dyn}}^2(\alpha_0,\alpha_1) = \Wass_\gamma^2(\alpha_0,\alpha_1).\)
The common value defines a distance on $\Pp_2(\RR^d)$; the triangle inequality is established by the robust representation in Proposition~\ref{prop-spectral-wasserstein-robust}.
\end{prop}
\begin{proof}
We first prove $\mathsf W_{\gamma,\mathrm{dyn}}^2\leq\Wass_\gamma^2$. Let $\pi\in\Couplings(\alpha_0,\alpha_1)$ and let $(X,Y)\sim\pi$. Set $Z_t=(1-t)X+tY$ and $\alpha_t=(Z_t)_\sharp\pi$. Define the Eulerian velocity as the conditional mean
\(v_t(z)\eqdef \EE[Y-X\mid Z_t=z]\).
Then $(\alpha_t,v_t)$ solves the continuity equation. If
\(M_\pi\eqdef\int(x-y)(x-y)^\top\d\pi(x,y), \qquad C_t\eqdef\int v_t(z)v_t(z)^\top\d\alpha_t(z),\)
conditional Jensen gives \(C_t\preceq M_\pi\). Indeed, for every direction \(u\in\RR^d\),
\[
u^\top C_tu
=
\EE\!\left[\EE[\dotp{u}{Y-X}\mid Z_t]^2\right]
\leq
\EE[(\dotp{u}{Y-X})^2]
=
u^\top M_\pi u.
\]
By the Loewner monotonicity of $\gamma$,
\(\int_0^1\gamma(C_t)\d t \leq \gamma(M_\pi)\).
Taking the infimum over $\pi$ gives the first inequality.

Conversely, let $(\alpha_t,v_t)$ be an admissible finite-action competitor. Since $\gamma$ is a finite positive gauge on the finite-dimensional cone $\mathbb S_+^d$, it is equivalent to the trace on this cone; hence finite spectral action implies the usual finite kinetic energy needed for the superposition principle. Thus there exists a probability law $\eta$ on absolutely continuous paths $\omega_\cdot$ such that $(e_t)_\sharp\eta=\alpha_t$ and $\dot\omega_t=v_t(\omega_t)$ for a.e.\ \(t\). Let $\pi=(e_0,e_1)_\sharp\eta$ be the induced endpoint coupling. With \(C_t=\int v_tv_t^\top\d\alpha_t\), Jensen's inequality along each path gives
\[
(\omega_1-\omega_0)(\omega_1-\omega_0)^\top
=
\left(\int_0^1\dot\omega_t\d t\right)
\left(\int_0^1\dot\omega_t\d t\right)^\top
\preceq
\int_0^1\dot\omega_t\dot\omega_t^\top\d t.
\]
After integration over $\eta$,
\(M_\pi\preceq \int_0^1 C_t\d t\).
Monotonicity of $\gamma$, followed by convexity, yields
\[
\gamma(M_\pi)
\leq
\gamma\!\left(\int_0^1 C_t\d t\right)
\leq
\int_0^1\gamma(C_t)\d t.
\]
Since $\pi$ is admissible for the static problem, $\Wass_\gamma^2(\alpha_0,\alpha_1)$ is bounded by the action of the dynamic competitor. Taking the infimum over all competitors gives the reverse inequality.

The crucial assumption is the monotonicity of $\gamma$: both parts of the proof produce Loewner-order comparisons between covariance matrices, and these comparisons imply inequalities between actions only for monotone gauges.
\end{proof}
\paragraph{Kernelized Benamou--Brenier distances.}\label{sec-kernelized-bb-distance}
\begin{defn}[Kernelized Benamou--Brenier distance]\label{def-kernelized-bb-distance}
Let $k$ be a Borel measurable positive definite kernel on $\RR^d$ with scalar RKHS $\RKHS_k$. Set
\(\RKHS_k^d\eqdef \RKHS_k\times\cdots\times\RKHS_k, \qquad \norm{v}_{\RKHS_k^d}^2 \eqdef \sum_{\ell=1}^d\norm{v_\ell}_{\RKHS_k}^2.\)
\begin{equation}\label{eq-kernelized-bb-distance}
\mathbb A_k(\alpha,v)
\eqdef
\norm{v}_{\RKHS_k^d}^2,
\qquad
\mathcal W_k^2(\alpha_0,\alpha_1)
\eqdef
\mathsf E_{\mathbb A_k}(\alpha_0,\alpha_1),
\end{equation}
where the infimum is restricted to strongly measurable $v\in L^2([0,1];\RKHS_k^d)$ satisfying the continuity equation.
\end{defn}
\begin{prop}[Kernelized dynamic distance]\label{prop-kernelized-bb-distance}
Assume that the kernel has uniformly bounded diagonal, \(\kappa_k^2\eqdef\sup_{x\in\RR^d}k(x,x)<+\infty.\)
Then the quantity \(\mathcal W_k\) defined by~\eqref{eq-kernelized-bb-distance} is an extended distance on probability measures. More precisely, it is symmetric, satisfies the triangle inequality, and \(\mathcal W_k(\alpha_0,\alpha_1)=0\) implies \(\alpha_0=\alpha_1\). It can take the value \(+\infty\) between different finite-action components.
\end{prop}
\begin{proof}
The Borel assumption makes every RKHS function measurable: kernel sections are Borel, their finite linear span is Borel, and the bounded evaluation estimate below turns RKHS-norm convergence into uniform convergence. The constant curve gives zero self-distance. To prove separation without assuming existence of a minimizing curve, use the RKHS evaluation bound
\(\norm{v(x)}\leq \kappa_k\norm{v}_{\RKHS_k^d}\).
For every $\varphi\in C_c^1(\RR^d)$ and every admissible curve of action $E$, the weak continuity equation and Cauchy--Schwarz give
\[
\begin{aligned}
\left|\int\varphi\,\d(\alpha_1-\alpha_0)\right|
&\leq
\int_0^1\!\int\norm{\nabla\varphi(x)}\,\norm{v_t(x)}\,\d\alpha_t(x)\d t\\
&\leq
\kappa_k\norm{\nabla\varphi}_\infty
\left(\int_0^1\norm{v_t}_{\RKHS_k^d}^2\d t\right)^{1/2}
=
\kappa_k\norm{\nabla\varphi}_\infty\sqrt E.
\end{aligned}
\]
Taking the infimum over admissible curves shows that $\mathcal W_k(\alpha_0,\alpha_1)=0$ forces equality of the integrals of every compactly supported smooth test function, hence $\alpha_0=\alpha_1$.

Symmetry follows by time reversal and by replacing \(v_t\) with \(-v_{1-t}\). For the triangle inequality, concatenate an almost optimal curve from \(\alpha_0\) to \(\alpha_1\) of action \(E_1\) with one from \(\alpha_1\) to \(\alpha_2\) of action \(E_2\). Compressing the first curve into a time interval of length \(\tau\) multiplies its action by \(1/\tau\), and compressing the second into length \(1-\tau\) multiplies its action by \(1/(1-\tau)\). Optimizing over \(\tau\in(0,1)\) gives \((\sqrt{E_1}+\sqrt{E_2})^2\), and taking infima yields the triangle inequality.
\end{proof}
One should read \(\mathcal W_k\) as an extended distance on finite-action components, not as a replacement for $\Wass_2$ on all of $\Pp_2(\RR^d)$.

\subsection{Nonlocal Wasserstein Distances}\label{sec-nonlocal-wasserstein-distances}
\begin{defn}[Logarithmic mean]\label{def-logarithmic-mean}
\begin{equation}\label{eq-logarithmic-mean}
\theta(a,b)\eqdef
\begin{cases}
\displaystyle\frac{a-b}{\log a-\log b}, & a,b>0\text{ and }a\neq b,\\[.4em]
a, & a=b,\\
0, & ab=0,
\end{cases}
\end{equation}
for $a,b\geq0$.
\end{defn}
For $a,b>0$, it satisfies \(\theta(a,b)(\log a-\log b)=a-b\); at vacuum, this relation is understood only as a limiting identity. This edge-wise chain rule identifies entropy-driven flows with the underlying reversible Markov or jump dynamics.

\paragraph{Continuum jump kernels.}
Let $(\mathcal X,d)$ be a Polish metric space, let $\mathfrak m$ be a Radon reference measure, and let $K(x,\cdot)$ be a nonnegative Borel kernel, possibly of infinite total mass. Assume that the pair measures introduced below are locally finite on the off-diagonal pair space $\mathcal G\eqdef\{(x,y)\in\mathcal X^2:x\neq y\}$ for the states under consideration. Define the jump measure $\mathsf J$ by
\begin{equation}\label{eq-nonlocal-jump-measure}
\int_{\mathcal G}\Phi(x,y)\,\mathsf J(\d x,\d y)
\eqdef
\int_{\mathcal X}\left(\int_{\mathcal X\setminus\{x\}}\Phi(x,y)K(x,\d y)\right)\mathfrak m(\d x).
\end{equation}
For a probability measure $\alpha$, introduce the two oriented edge measures
\begin{equation}\label{eq-nonlocal-oriented-edge-measures}
\alpha^1(\d x,\d y)\eqdef K(x,\d y)\alpha(\d x),
\qquad
\alpha^2(\d x,\d y)\eqdef K(y,\d x)\alpha(\d y).
\end{equation}
If $\nu$ is a signed Radon measure on $\mathcal G$ and $\varsigma$ dominates $\alpha^1$, $\alpha^2$, and $|\nu|$, write $r_1=\d\alpha^1/\d\varsigma$, $r_2=\d\alpha^2/\d\varsigma$, and $q=\d\nu/\d\varsigma$. The logarithmic mean is positively $1$-homogeneous, so the closed perspective below is independent of the choice of $\varsigma$.

\begin{defn}[Continuum nonlocal Wasserstein distance]\label{def-continuum-nonlocal-wasserstein}
Define the relaxed pair action by
\begin{equation}\label{eq-nonlocal-relaxed-action}
\mathbb A_K(\alpha,\nu)
\eqdef
\frac12\int_{\mathcal G}
\begin{cases}
\displaystyle\frac{|q|^2}{\theta(r_1,r_2)},&\theta(r_1,r_2)>0,\\[.35em]
0,&\theta(r_1,r_2)=0,\ q=0,\\
+\infty,&\theta(r_1,r_2)=0,\ q\neq0,
\end{cases}
\d\varsigma.
\end{equation}
An admissible pair consists of a narrowly continuous curve $(\alpha_t)_{t\in[0,1]}$ and a Borel family of antisymmetric signed flux measures $(\nu_t)_t$ such that
\(\int_0^1\!\int_{\mathcal G}(1\wedge d(x,y))\d|\nu_t|(x,y)\d t<+\infty\)
and, for every test function $\varphi$ that is $C^1$ in time, bounded Lipschitz in space, and compactly supported in space when needed,
\begin{equation}\label{eq-nonlocal-continuity-weak}
\int_0^1\!\int_{\mathcal X}\partial_t\varphi_t\d\alpha_t\d t
+\frac12\int_0^1\!\int_{\mathcal G}\bar\nabla\varphi_t\d\nu_t\d t
=
\int_{\mathcal X}\varphi_1\d\alpha_1-
\int_{\mathcal X}\varphi_0\d\alpha_0.
\end{equation}
The induced action distance is
\begin{equation}\label{eq-nonlocal-wasserstein-distance}
\mathcal W_K^2(\alpha_0,\alpha_1)
\eqdef
\inf_{(\alpha_t,\nu_t)}\int_0^1\mathbb A_K(\alpha_t,\nu_t)\d t.
\end{equation}
\end{defn}
For the density specialization, let $\alpha=\rho\mathfrak m$ and
\begin{equation}\label{eq-nonlocal-density-flux}
\nu(\d x,\d y)
=
v(x,y)\theta(\rho(x),\rho(y))\mathsf J(\d x,\d y),
\end{equation}
where $v$ is antisymmetric. Then $\alpha^1=\rho(x)\mathsf J$, $\alpha^2=\rho(y)\mathsf J$, and
\(\mathbb A_K(\alpha,\nu) = \frac12\int_{\mathcal G}|v(x,y)|^2\theta(\rho(x),\rho(y))\mathsf J(\d x,\d y).\)
Thus the earlier density--velocity formula is recovered, but a minimizing relaxed curve need not remain absolutely continuous with respect to $\mathfrak m$.

\begin{prop}[Metric properties of nonlocal transport]\label{prop-nonlocal-distance-properties}
Let $\mathcal X=\RR^d$. Assume that $K$ and $\mathfrak m$ satisfy Assumption~1.1 of: besides reversibility, the weighted kernels $(1\wedge\norm{x-y}^2)K(x,\d y)$ depend continuously on $x$ against bounded continuous tests and are uniformly integrable both near the diagonal and at infinity. Then $\mathcal W_K$ is an extended distance on $\Pp(\RR^d)$. Every finite-distance component is complete and geodesic, the infimum in~\eqref{eq-nonlocal-wasserstein-distance} is attained whenever it is finite, and a minimizer can be parametrized at constant speed.
\end{prop}
\begin{proof}
The density in~\eqref{eq-nonlocal-relaxed-action} is a convex lower-semicontinuous perspective and is $1$-homogeneous in $(r_1,r_2,q)$. Hence the action is well defined independently of $\varsigma$. The logarithmic mean satisfies Assumption~2.1 of. Proposition~3.4 there supplies compactness and closure of the measure--flux continuity equation, Proposition~4.3 supplies attainment and constant-speed parametrization, and Theorem~4.4 proves definiteness, completeness and the geodesic property. Symmetry follows by time reversal $(\alpha_t,\nu_t)\mapsto(\alpha_{1-t},-\nu_{1-t})$. Concatenating two curves and allocating time in proportion to the square roots of their actions gives the triangle inequality, exactly as in Proposition~\ref{prop-homogeneous-dynamic-action-distance}. The consequences for entropy dynamics and fractional PDE examples are developed in Section~\ref{sec-nonlocal-wasserstein-pdes}.
\end{proof}
For a fixed jump kernel, this geometry is genuinely nonlocal and does not coincide with ordinary $\Wass_2$. A local metric is recovered in a small-jump limit only if the kernel can transport through vacuum. Suppose that
\(M_2(\eta)\eqdef\int_{\RR^d}\norm{z}^2\eta(z)\,\d z\in(0,+\infty),\)
and, for some $c,r_0>0$ and $s\in(0,2)$,
\begin{equation}\label{eq-small-jump-singular-lower-bound}
\eta(z)\geq c\norm{z}^{-d-s}
\qquad\text{whenever }0<\norm{z}<r_0.
\end{equation}
\eqllead{and define}{eq-small-jump-kernel-scaling}{
K_\varepsilon(x,\d y)
\eqdef
\eta_\varepsilon(y-x)\,\d y,
\qquad
\eta_\varepsilon(z)
\eqdef
\varepsilon^{-d}\eta(z/\varepsilon).
}
Radiality implies symmetry and isotropy, and a change of variables gives
\begin{equation}\label{eq-small-jump-kernel-moments}
\int\norm{y-x}^2K_\varepsilon(x,\d y)
=
\varepsilon^2M_2(\eta),
\qquad
\int (y-x)(y-x)^\top K_\varepsilon(x,\d y)
=
\varepsilon^2\frac{M_2(\eta)}{d}\Id.
\end{equation}
To obtain a nontrivial local limit, one simultaneously accelerates the jump rate by $\varepsilon^{-2}$ and sets
\(\widehat K_\varepsilon \eqdef \frac{2d}{\varepsilon^2M_2(\eta)}K_\varepsilon, \qquad \int (y-x)(y-x)^\top\widehat K_\varepsilon(x,\d y)=2\Id.\)
Theorem~1.3 of therefore gives, for endpoints supported in a fixed compact set,
\[
\mathcal W_{\widehat K_\varepsilon}
=
\varepsilon\sqrt{\frac{M_2(\eta)}{2d}}\,
\mathcal W_{K_\varepsilon}
\underset{\varepsilon\to0}{\longrightarrow}
\Wass_2.
\]
If $B$ is invertible and
\(\eta_B(z)=|\det B|^{-1}\eta(B^{-1}z),\)
define
\[
K_{B,\varepsilon}(x,\d y)
\eqdef
\varepsilon^{-d}\eta_B\!\left(\frac{y-x}{\varepsilon}\right)\d y,
\qquad
\widehat K_{B,\varepsilon}
\eqdef
\frac{2d}{\varepsilon^2M_2(\eta)}K_{B,\varepsilon}.
\]
The change of variables $x=B\xi$, $y=B\zeta$, together with the $1$-homogeneity of the logarithmic mean, then gives
\[
\mathcal W_{\widehat K_{B,\varepsilon}}(\alpha_0,\alpha_1)
=
\mathcal W_{\widehat K_\varepsilon}(B^{-1}_\sharp\alpha_0,B^{-1}_\sharp\alpha_1)
\longrightarrow
\Wass_2(B^{-1}_\sharp\alpha_0,B^{-1}_\sharp\alpha_1).
\]
\paragraph{Discrete Wasserstein distances on Markov chains.}\label{sec-discrete-wasserstein-markov}
We write
\(\mathsf J_{i,j}\eqdef \pi_iK_{i,j}=\pi_jK_{j,i}\)
The transported object is a mass histogram
\[
\simplex_n\eqdef\left\{a\in\RR_+^n:\sum_i a_i=1\right\}.
\]
Relative densities only enter as auxiliary variables with respect to the invariant law: \(\rho_i(a)\eqdef \frac{a_i}{\pi_i}, \qquad a_i=\pi_i\rho_i(a).\)
The logarithmic mean \(\theta\) defined in~\eqref{eq-logarithmic-mean} is the mobility selected so that the entropy calculus later recovers exactly the Markov evolution; see Proposition~\ref{prop-discrete-markov-entropy-gradient}. For positive $a,b$, the identity \(\theta(a,b)(\log a-\log b)=a-b\), extended to boundary states by approximation, converts entropy gradients into density differences along graph edges. For a potential \(\psi\in\RR^n\), set
\(\bar\nabla\psi(i,j)\eqdef\psi_j-\psi_i.\)
\begin{defn}[Discrete Markov-chain Wasserstein distance]\label{def-discrete-markov-wasserstein}
\begin{equation}
\label{eq-discrete-markov-onsager}
(\mathcal K_\rho\psi)_i
\eqdef
\sum_j K_{i,j}\theta(\rho_i,\rho_j)(\psi_i-\psi_j),
\qquad
(\mathcal L_a\psi)_i
\eqdef
\pi_i(\mathcal K_{\rho(a)}\psi)_i.
\end{equation}
Its associated tangent action is
\begin{equation}
\label{eq-discrete-markov-action}
\mathbb A_K(a,\psi)
\eqdef
\frac12\sum_{i,j}\mathsf J_{i,j}\theta(\rho_i(a),\rho_j(a))
(\bar\nabla\psi(i,j))^2.
\end{equation}
\begin{equation}
\label{eq-discrete-markov-distance}
\mathcal W_K^2(a_0,a_1)
\eqdef
\inf_{a_t,\psi_t}
\int_0^1\mathbb A_K(a_t,\psi_t)\,\d t,
\qquad
\dot a_t+\mathcal L_{a_t}\psi_t=0,
\end{equation}
with endpoint conditions \(a_{t=0}=a_0\) and \(a_{t=1}=a_1\).
\end{defn}
\subsection{Dynamic Unbalanced Wasserstein Distances}
\label{sec-unbalanced-ot}
\paragraph{Balance equation and tangent variables.}
A representative quadratic balance equation is
\[
\partial_t\rho_t+\nabla\!\cdot m_t=s_t,
\qquad
\int_0^1\!\int
\left(\frac{\norm{m_t}^2}{\rho_t}+\kappa^2\frac{s_t^2}{\rho_t}\right)\d x\,\d t,
\]
\paragraph{Reaction--transport action.}
\begin{defn}[Dynamic Wasserstein--Fisher--Rao distance]\label{def-dynamic-wfr-distance}
\eqllead{For $\kappa>0$, define}{eq-wfr-velocity-action}{
A_\kappa(a,w,g)\eqdef a\bigl(\norm{w}^2+\kappa^2g^2\bigr),
}
\eqllead{and}{eq-wfr-momentum-perspective}{
J_\kappa(a,m,r)\eqdef
\begin{cases}
(\norm{m}^2+\kappa^2r^2)/a,&a>0,\\
0,&(a,m,r)=(0,0,0),\\
+\infty,&a=0,\ (m,r)\neq(0,0).
\end{cases}
}
If $\lambda$ dominates $\alpha$, $|\omega|$, and $|\sigma|$, set
\begin{equation}\label{eq-wfr-measure-action}
\mathbb J_\kappa(\alpha,\omega,\sigma)
\eqdef
\int J_\kappa\!\left(\frac{\d\alpha}{\d\lambda},
\frac{\d\omega}{\d\lambda},\frac{\d\sigma}{\d\lambda}\right)\d\lambda.
\end{equation}
The induced dynamic Wasserstein--Fisher--Rao distance is
\begin{equation}\label{eq-dynamic-unbalanced-ot}
\WFR_\kappa^2(\alpha_0,\alpha_1)
\eqdef
\inf_{\substack{\partial_t\alpha_t+\nabla\cdot\omega_t=\sigma_t\\
\alpha_{t=0}=\alpha_0,\ \alpha_{t=1}=\alpha_1}}
\int_0^1\mathbb J_\kappa(\alpha_t,\omega_t,\sigma_t)\d t.
\end{equation}
\end{defn}
\paragraph{Static and dynamic viewpoints.}
The dynamic balance-equation formula is not a separate object from static unbalanced OT: it is the least-action representation of the same cone distance. To make the normalization explicit, define on the cone $\mathfrak C[\RR^d]$ the squared cost
\begin{equation}\label{eq-wfr-scaled-cone-metric}
\Delta_\kappa\big((x,r),(y,s)\big)^2
\eqdef
4\kappa^2
\left[
r^2+s^2-2rs
\cos\!\left(
\frac{\norm{x-y}}{2\kappa}\wedge\frac{\pi}{2}
\right)
\right].
\end{equation}
The radii encode masses through the weighted projection $\mathsf P_2$ defined in Section~\ref{sec-unbalanced}. Accordingly, set
\begin{equation}\label{eq-wfr-scaled-cone-value}
\CW_\kappa(\alpha_0,\alpha_1)
\eqdef
\inf_{\substack{\gamma\in\Mm_+(\mathfrak C[\RR^d]^2)\\
\mathsf P_2\gamma_1=\alpha_0,\ \mathsf P_2\gamma_2=\alpha_1}}
\int
\Delta_\kappa\big((x,r),(y,s)\big)^2
\d\gamma(x,r,y,s).
\end{equation}
For $\kappa=1/2$, this is exactly the normalization of the Hellinger--Kantorovich cone cost used in Section~\ref{sec-unbalanced}.

Thus Definition~\ref{def-dynamic-wfr-distance} gives both the reaction--transport action and the induced dynamic WFR metric.

\begin{prop}[Static/dynamic equivalence for unbalanced OT]\label{prop-static-dynamic-unbalanced}
For nonnegative finite measures $\alpha_0,\alpha_1$ on $\RR^d$, the dynamic value~\eqref{eq-dynamic-unbalanced-ot} equals the static cone value~\eqref{eq-wfr-scaled-cone-value}. Hence $\WFR_\kappa=\CW_\kappa^{1/2}$ is the geodesic distance generated by the balance-equation least-action problem.
\end{prop}
\begin{proof}
The cone construction turns variation of mass into radial motion and spatial transport into angular motion on $\mathfrak C[\RR^d]$. The normalization can be checked on a smooth cone path. If its base position is $x_t$, its cone radius is $r_t$, and the projected mass is $a_t=r_t^2$, then its relative growth rate is $g_t=\dot a_t/a_t=2\dot r_t/r_t$. The infinitesimal energy induced by~\eqref{eq-wfr-scaled-cone-metric} is therefore
\[
4\kappa^2\dot r_t^2+r_t^2\norm{\dot x_t}^2
=
a_t\left(\norm{\dot x_t}^2+\kappa^2g_t^2\right),
\]
which is precisely the velocity action~\eqref{eq-wfr-velocity-action}.

The converse passage from an arbitrary Eulerian triple to a cone-valued dynamic plan is the substantive step. For $\kappa=1/2$, Theorems~4.3, 4.5 and 4.6 of provide respectively the dynamic-plan representation and the two action comparisons, and the proof of their Theorem~3.6 identifies the resulting distance with the cone metric. In their notation, the vector field is the chapter's velocity $w$, while the scalar field $\xi$ satisfies $g=4\xi$; hence their density $\norm w^2+4\xi^2$ is exactly $\norm w^2+\tfrac14g^2$. These theorems prove both inequalities for arbitrary finite measures, not only for smooth cone paths.

For general $\kappa$, rescale the base metric by $(2\kappa)^{-1}$ and the cone distance by $2\kappa$. The same dynamic-plan theorem then gives the cone cost~\eqref{eq-wfr-scaled-cone-metric} and the action $\norm w^2+\kappa^2g^2$. This proves equality of the two infima. The corresponding complete-separable-space formulation and relaxation are developed in; related normalizations appear in.
\end{proof}
\subsection{Variational Mean Field Games}
\label{sec-variational-mean-field-games}
\paragraph{From individual control to a population equilibrium.}
Given a candidate population density $(\rho_t)_{t\in[0,1]}$, a representative deterministic agent starting from $x\in\overline\Omega$ chooses an admissible path $\gamma:[0,1]\to\overline\Omega$ by minimizing
\begin{equation}\label{eq-mfg-individual-control}
\inf_{\gamma(0)=x}
\left\{
\int_0^1
\left(
\frac12\norm{\dot\gamma(t)}^2
+g\bigl(\rho_t(\gamma(t))\bigr)
\right)\d t
+\Psi(\gamma(1))
\right\}.
\end{equation}
\paragraph{A Benamou--Brenier planning problem.}
The consistency condition is usually coupled to an individual Hamilton--Jacobi--Bellman equation. A global variational reduction is available when the local coupling is a first variation. Let $G:[0,+\infty)\to[0,+\infty]$ be proper, closed and convex with $G(0)=0$; in the differentiable case, assume
\begin{equation}\label{eq-mfg-congestion-primitive}
g(r)=G'(r),
\end{equation}
\begin{defn}[Variational mean-field-game planning problem]\label{def-variational-mfg-planning}
\begin{equation}\label{eq-variational-mfg-velocity}
\inf_{(\rho_t,v_t)}
\left\{
\int_0^1\!\int_{\Omega}
\left(
\frac12\rho_t(x)\norm{v_t(x)}^2+G(\rho_t(x))
\right)\d x\,\d t
+\int_{\Omega}\Psi(x)\rho_1(x)\d x
\right\},
\end{equation}
\eqllead{subject to}{eq-variational-mfg-continuity}{
\partial_t\rho_t+\diverg(\rho_t v_t)=0,
\qquad
\rho_{t=0}\,\d x=\alpha_0,
\qquad
(\rho_t v_t)\cdot n=0\ \text{on }\partial\Omega.
}
\end{defn}
Unlike in~\eqref{eq:benamou-brenier}, the final measure is not prescribed. The terminal potential $\Psi$ selects desirable final states softly, while the integral of $G(\rho_t)$ penalizes crowded intermediate configurations.

\paragraph{Convex momentum formulation.}
To obtain the weakly lower-semicontinuous congestion functional, write the Lebesgue decomposition $\alpha=\rho\,\d x+\alpha^s$ and set
\begin{equation}\label{eq-mfg-congestion-functional}
\mathcal C_G(\alpha)
\eqdef
\int_\Omega G(\rho(x))\d x
+G^\infty\alpha^s(\overline\Omega),
\qquad
G^\infty\eqdef\lim_{r\to+\infty}\frac{G(r)}r.
\end{equation}
For superlinear congestion, $G^\infty=+\infty$ and finite energy forces $\alpha\ll\d x$; the recession term is needed for closedness when $G$ has only linear growth. Using the measure perspective action~\eqref{eq-measure-perspective-action}, the same problem becomes the generalized Benamou--Brenier formula
\begin{equation}\label{eq-variational-mfg-momentum}
\inf_{(\alpha_t,\omega_t)}
\left\{
\int_0^1
\left(
\frac12\mathbb J(\alpha_t,\omega_t)
+\mathcal C_G(\alpha_t)
\right)\d t
+\int_{\overline\Omega}\Psi\,\d\alpha_1
\right\},
\qquad
\begin{gathered}
\partial_t\alpha_t+\diverg\omega_t=0,
\qquad \alpha_{t=0}=\alpha_0,\\
\omega_t\cdot n=0\quad\text{on }\partial\Omega.
\end{gathered}
\end{equation}
Equivalently, in density--momentum variables, the running integrand is
\begin{equation}\label{eq-variational-mfg-density-momentum}
\frac12J(\rho_t,m_t)+G(\rho_t)
=
\frac{\norm{m_t}^2}{2\rho_t}+G(\rho_t),
\end{equation}
\paragraph{Optimality system and game interpretation.}
\begin{prop}[Potential MFG system]\label{prop-variational-mfg-system}
Assume that an optimizer of~\eqref{eq-variational-mfg-momentum} has a smooth positive density $\rho$, that $G$ is differentiable, and set $g=G'$. Then there is a value function $u$ such that
\begin{equation}\label{eq-variational-mfg-system}
\left\{
\begin{aligned}
-\partial_tu+\frac12\norm{\nabla u}^2&=g(\rho),
& u_{t=1}&=\Psi,\\
\partial_t\rho-\diverg(\rho\nabla u)&=0,
& \rho_{t=0}\,\d x&=\alpha_0,
\end{aligned}
\right.
\qquad
m=-\rho\nabla u,
\qquad
\rho\nabla u\cdot n=0\ \text{on }\partial\Omega.
\end{equation}
\end{prop}
\begin{proof}
Use $-u$ as the multiplier for $\partial_t\rho+\diverg m=0$. After integration by parts, the space--time Lagrangian density is \(\frac{\norm{m}^2}{2\rho}+G(\rho) +\rho\,\partial_tu+m\cdot\nabla u.\)
Stationarity in $m$ gives $m=-\rho\nabla u$. Stationarity in $\rho$, followed by this substitution, gives the first equation in~\eqref{eq-variational-mfg-system}. The free terminal density contributes $(\Psi-u_1)\rho_1$. Since admissible terminal variations preserve total mass, stationarity first gives $u_1-\Psi$ equal to a spatial constant. The multiplier $u$ is defined up to an additive constant, which we choose so that $u_1=\Psi$. Primal feasibility gives the second equation.
\end{proof}
\paragraph{Hard congestion through a bottleneck.}
A global hard crowd-capacity constraint is obtained from
\begin{equation}\label{eq-variational-mfg-hard-cap}
G_\kappa(r)=\iota_{[0,\kappa]}(r)
=
\begin{cases}
0,&0\leq r\leq\kappa,\\
+\infty,&\text{otherwise}.
\end{cases}
\end{equation}
In order to make the cap active while retaining a prescribed terminal preference, one may replace the linear terminal term in~\eqref{eq-variational-mfg-momentum} by the convex penalty
\begin{equation}\label{eq-variational-mfg-quadratic-terminal}
\Gamma_\eta(\rho_1)
=
\frac{\eta}{2}
\int_\Omega\abs{\rho_1(x)-\rho_\star(x)}^2\d x,
\end{equation}
where $\rho_\star$ is a desired terminal density.

The role of endpoint relaxation depends on the geometry. If two fixed endpoint densities on $\RR^d$ satisfy a spatially constant cap, their ordinary $\Wass_2$ geodesic satisfies the same cap by Proposition~\ref{prop-density-cap-geodesic-convex}, so the cap does not alter that geodesic.

\section{Wasserstein Gradient Flows}
\label{sec-wasserstein-gradient-flows}
\subsection{Minimizing Movements and Wasserstein Gradients}
\begin{defn}[JKO minimizing movement]\label{def-jko-minimizing-movement}
Given $\tau>0$ and $\alpha^{(0)}$, define recursively
\begin{equation}
\alpha^{(\ell+1)} \in \uargmin{\alpha\in\Pp_2(\RR^d)}
\left\{\frac{1}{2 \tau} \Wass_2(\alpha^{(\ell)}, \alpha)^2 + f(\alpha)\right\}. \label{eq:jko-discr}
\end{equation}
\end{defn}
\paragraph{Euclidean gradient flows.}
If we restrict \eqref{eq:jko-discr} to finite dimensions and assume $\alpha_t = \delta_{x(t)}$ and $\alpha = \delta_x$ (single Dirac measures), this matches the implicit Euler scheme:
\begin{equation*}
x(t+\tau) \in \uargmin{x}\left\{\frac{1}{2 \tau} \norm{x - x(t)}^2 + h(x)\right\},
\end{equation*}
where $h(x) = f(\delta_x)$. When $h$ is differentiable and the minimizer is unique, the solution is characterized by the implicit Euler formula
\(x(t+\tau) = (\Id + \tau \nabla h)^{-1}(x(t)).\)
In contrast, the explicit Euler scheme is: \(x(t+\tau) = (\Id - \tau \nabla h)(x(t)) = x(t) - \tau \nabla h(x(t)).\)
If, for instance, $\nabla h$ is globally Lipschitz, then for sufficiently small $\tau$ both schemes are well posed and their interpolants converge locally uniformly as $\tau\to0$ to the unique solution of
\begin{equation}
\dot{x}(t) = -\nabla h(x(t)). \label{eq:grad-flow-classical}
\end{equation}
\paragraph{Vertical derivatives (first- and higher-order variations).}
\begin{defn}[First variation]\label{def-first-variation}
Assume that \(f(\alpha)<+\infty\). A measurable function \(\delta f(\alpha)\) is a first variation of \(f\) at \(\alpha\) if, for every \(\beta\in\Pp(\RR^d)\) such that \(\delta f(\alpha)\) is integrable with respect to \(\lvert\beta-\alpha\rvert\) and the mixture \(\alpha_\vartheta\) remains in the domain of \(f\) for small \(\vartheta>0\),
\(f(\alpha_\vartheta) = f(\alpha) +\vartheta\int_{\RR^d}[\delta f(\alpha)](x) \d(\beta-\alpha)(x) +o(\vartheta) \qquad(\vartheta\downarrow0).\)
\end{defn}
It represents the directional linear form
\(Df(\alpha)[\eta] \eqdef \int_{\RR^d}[\delta f(\alpha)](x)\d\eta(x)\)
\begin{defn}[Higher-order vertical derivatives]\label{def-higher-order-vertical-derivatives}
Set \(D^0f=f\). The functional \(f\) admits vertical derivatives through order \(r\) if there are symmetric multilinear forms
\(D^j f(\alpha):\Mm_0(\Xx)^j\longrightarrow\RR, \qquad 1\leq j\leq r,\)
such that, along every mixture segment in \(\Pp(\Xx)\) and for fixed \(\eta_1,\ldots,\eta_{j-1}\in\Mm_0(\Xx)\),
\(\frac{\d}{\d\vartheta} D^{j-1}f(\alpha_\vartheta)[\eta_1,\ldots,\eta_{j-1}] = D^j f(\alpha_\vartheta) [\beta-\alpha,\eta_1,\ldots,\eta_{j-1}].\)
At the endpoints, derivatives are understood one-sided.
\end{defn}
When these multilinear forms have integral kernels, the notation
\(D^j f(\alpha)[\eta_1,\ldots,\eta_j] = \int_{\Xx^j} \frac{\delta^j f}{\delta\alpha^j} (\alpha,x_1,\ldots,x_j) \d\eta_1(x_1)\cdots\d\eta_j(x_j)\)
A canonical representative relative to the base measure \(\alpha\) is obtained by centering each variable, \(\int_{\Xx} \frac{\delta^j f}{\delta\alpha^j} (\alpha,x_1,\ldots,x_j)\,\d\alpha(x_i)=0 \qquad (1\leq i\leq j),\)
\paragraph{Horizontal derivative (Wasserstein gradient).}
\begin{defn}[Wasserstein gradient]\label{def-wasserstein-gradient}
Assume that $f$ admits a smooth first variation $\delta f(\alpha)$. Here smoothness means that $\nabla\delta f(\alpha)\in L^2(\alpha;\RR^d)$, that the transport chain rule holds along the admissible continuity-equation curves, and that the boundary conditions justify integration by parts. In the smooth formal calculus on $\Pp_2(\RR^d)$, the Wasserstein gradient of $f$ at $\alpha$ is the gradient vector field
\(\Wgrad f(\alpha) = \nabla_x \delta f(\alpha).\)
\end{defn}
The associated formal gradient flow is the continuity equation
\begin{equation}
\frac{\partial \alpha_t}{\partial t} + \diverg(-\Wgrad f(\alpha_t) \alpha_t) = 0. \label{eq:wassflow-pde}
\end{equation}
\begin{proposition}[Formal Wasserstein gradient]\label{prop-formal-wass-gradient}
Assume that $f$ admits a smooth first variation $\delta f(\alpha)$ and that $\alpha$ has a smooth positive density. Let $v$ be a sufficiently regular velocity field, compactly supported or satisfying a no-flux boundary condition. For perturbations generated by $(\Id+\tau v)_\sharp\alpha$, the differential of $f$ is
\(\frac{\d}{\d\tau}_{|\tau=0} f\big((\Id+\tau v)_\sharp\alpha\big) = \int \dotp{\nabla \delta f(\alpha)(x)}{v(x)}\d\alpha(x).\)
Hence, for the Riemannian metric $\norm{v}_{L^2(\alpha)}^2=\int \norm{v}^2\d\alpha$, the Wasserstein gradient is the vector field
\(\Wgrad f(\alpha)=\nabla \delta f(\alpha).\)
\end{proposition}
\begin{proof}
The push-forward expansion gives, in the sense of distributions, \((\Id+\tau v)_\sharp\alpha = \alpha-\tau\diverg(\alpha v)+o(\tau).\)
Using the definition of the first variation, \(f((\Id+\tau v)_\sharp\alpha) = f(\alpha)-\tau\int \delta f(\alpha)\diverg(\alpha v)\d x+o(\tau).\)
An integration by parts, with either compact support or vanishing boundary flux, gives \(-\int \delta f(\alpha)\diverg(\alpha v)\d x = \int \dotp{\nabla\delta f(\alpha)}{v}\d\alpha.\)
By definition of the Riesz representative for the $L^2(\alpha)$ metric, this representative is $\nabla\delta f(\alpha)$.
\end{proof}
\paragraph{From the JKO step to the velocity field.}
A first-order expansion of the JKO step explains why~\eqref{eq:wassflow-pde} uses the vector field $\Wgrad f(\alpha)$. This is a formal local argument: assume that $\alpha_t$ is absolutely continuous and that, for the nearby competitors selected by the JKO step, the optimal map from $\alpha_t$ can be written as $\Id+\tau v$. For these optimal displacement representatives,
\(\Wass_2^2\bigl(\alpha_t,(\Id+\tau v)_\sharp\alpha_t\bigr) = \tau^2\norm{v}_{L^2(\alpha_t)}^2.\)
The local JKO objective therefore reads \(\min_v \frac{\tau}{2}\norm{v}_{L^2(\alpha_t)}^2 +f\bigl((\Id+\tau v)_\sharp\alpha_t\bigr),\)
The two expansions
\begin{align*}
(\Id+\tau v)_\sharp\alpha_t
&=
\alpha_t-\tau\diverg(v\alpha_t)+o(\tau),\\
f\bigl((\Id+\tau v)_\sharp\alpha_t\bigr)
&=
f(\alpha_t)
+\tau\int\dotp{\nabla_x\delta f(\alpha_t)(x)}{v(x)}\d\alpha_t(x)
+o(\tau)
\end{align*}
give the first-order problem
\[
\min_v f(\alpha_t)
+\tau\int\left[
\frac12\norm{v(x)}^2
+\dotp{\Wgrad f(\alpha_t)(x)}{v(x)}
\right]\d\alpha_t(x)
+o(\tau).
\]
\paragraph{Metric derivative and curves of maximal slope.}
\begin{defn}[Metric derivative and absolutely continuous curve]\label{def-metric-derivative-ac-curve}
For \(1\leq p\leq+\infty\), a curve $x:[0,T]\to(\mathcal X,d)$ belongs to \(AC^p([0,T];\mathcal X)\) if some $g\in L^p(0,T)$ satisfies $d(x_s,x_t)\leq\int_s^t g(r)\d r$. Absolutely continuous means \(AC^1\). Its metric derivative is
\(|\dot x_t| \eqdef \lim_{h\to 0}\frac{d(x_{t+h},x_t)}{|h|}\)
for a.e.\ $t$.
\end{defn}
\begin{defn}[Metric slope and curve of maximal slope]\label{def-curve-maximal-slope}
For lower-semicontinuous $f:\mathcal X\to(-\infty,+\infty]$, define
\[
|\partial f|(x)
\eqdef
\begin{cases}
\limsup_{\substack{y\to x\\ y\neq x}}
\dfrac{( f(x)-f(y) )_+}{d(x,y)}, & f(x)<+\infty,\\[.4em]
+\infty, & f(x)=+\infty.
\end{cases}
\]
It is a strong upper gradient if $|f(y_t)-f(y_s)|\leq \int_s^t |\partial f|(y_r)|\dot y_r|\d r$ along every absolutely continuous $y$. When this property holds, a curve $x\in AC^2([0,T];\mathcal X)$ with $f(x_0)<+\infty$ is a curve of maximal slope for $f$ if $f\circ x$ has an absolutely continuous representative and
\(f(x_t) +\frac12\int_s^t |\dot x_r|^2\d r +\frac12\int_s^t |\partial f|^2(x_r)\d r \leq f(x_s), \qquad 0\leq s<t.\)
\end{defn}
\paragraph{Discrete evolutions.}
If $f$ admits a first variation whose spatial gradient is well defined at empirical measures, the characteristic form of~\eqref{eq:wassflow-pde} preserves the number and weights of the atoms. For equal weights, write $\alpha_t = \frac{1}{n} \sum_i \delta_{x_i(t)}$. The particles $X(t):= (x_i(t))_i$ evolve according to the coupled ODEs
\begin{equation}
\dot{x}_i(t) = -n\nabla_{x_i} F(X(t)), \label{eq:wassflows-particles}
\end{equation}
where $F(X):= f\left(\frac{1}{n} \sum_i \delta_{x_i}\right)$ and the factor $n$ comes from the empirical Wasserstein metric $\frac1n\sum_i \norm{\dot x_i}^2$.

\paragraph{Shannon neg-entropy.}
\begin{defn}[Score function]\label{def-score-function}
For $\alpha=\rho\,\d x$ with differentiable positive density $\rho$, its score is $\nabla\log\rho$. Relative to $\beta$, its relative score is $\nabla\log(\d\alpha/\d\beta)$ whenever this logarithmic density has a weak gradient.
\end{defn}
\paragraph{Interaction energies.}
In a similar spirit, to obtain nonlinear evolutions without requiring the measure to have density, one can consider
\begin{equation}
f(\alpha):= \iint k(x, y) \d \alpha(x) \d \alpha(y). \label{eq:quadratic-func}
\end{equation}
For a symmetric kernel $k$: \(\delta f(\alpha)(x) = 2 \int k(x, y) \d \alpha(y), \quad \Wgrad f(\alpha)(x) = 2 \int \nabla_x k(x, y) \d \alpha(y).\)
For $\alpha_0 = \frac{1}{n} \sum_i \delta_{x_i}$, the flow \eqref{eq:wassflow-pde} implies particles $(x_i(t))_i$ obey: \(\dot{x}_i(t) = -\frac{2}{n} \sum_j \nabla k(x_i(t), x_j(t)).\)
If $k$ is positive definite, or more generally conditionally positive definite on signed measures of zero total mass as for the energy-distance kernel $k(x,y)=-\norm{x-y}$, and one minimizes the squared kernel discrepancy to a teacher distribution $\beta$, then
\[
\norm{\alpha-\beta}_k^2
=
\iint k\d\alpha\d\alpha
-2\int\!\left(\int k(x,y)\d\beta(y)\right)\d\alpha(x)
+\mathrm{constant}.
\]
Thus MMD-type training energies are exactly an interaction energy plus a linear potential; the teacher distribution appears through the potential $x\mapsto-2\int k(x,y)\d\beta(y)$. The corresponding empirical Wasserstein gradient flow is
\(\dot x_i(t) = -\frac{2}{n}\sum_j\nabla_x k(x_i(t),x_j(t)) + 2\int\nabla_x k(x_i(t),y)\d\beta(y).\)
\paragraph{Stochastic particles and McKean--Vlasov limits.}
If the drift does not depend on the empirical measure, each particle evolves independently according to
\(\d X_t=b(X_t)\d t+\sqrt{2}\sigma\d B_t,\)
and the one-particle law $\alpha_t=\rho_t\d x$ directly satisfies the linear Fokker--Planck equation
\(\partial_t\rho_t=-\diverg(b\rho_t)+\sigma^2\Delta\rho_t.\)
The mean-field case is different: the drift is recomputed from the current empirical distribution of all particles,
\(\d X_i^n(t)=b(X_i^n(t),\al_t^n)\d t+\sqrt{2}\sigma\d B_i(t), \qquad \al_t^n=\frac1n\sum_{i=1}^n\delta_{X_i^n(t)}.\)
For finite $n$, the empirical law $\al_t^n$ is itself random. Under suitable Lipschitz, growth and chaotic-initialization assumptions, propagation of chaos states that finitely many particles become asymptotically independent as $n\to\infty$, all with the same deterministic law $\rho_t\d x$; equivalently, the empirical measure $\al_t^n$ converges in probability to this law. The limiting density solves the nonlinear Fokker--Planck, or McKean--Vlasov, equation
\(\partial_t \rho_t = -\diverg\big(b(x,\rho_t)\rho_t\big) + \sigma^2\Delta\rho_t.\)
When the interaction drift has variational form
\(b(x,\rho)=-\nabla \frac{\delta \mathcal E}{\delta \rho}(x),\)
this PDE is the Wasserstein gradient flow of the entropy-regularized energy
\(\mathcal E(\rho)+\sigma^2\int \rho\log\rho\,\d x.\)
For a prescribed target $\beta=\rho_\beta\d x$, take $b=\sigma^2\nabla\log\rho_\beta$. The SDE, the Fokker--Planck equation, and the deterministic continuity equation with score velocity are then three representations of the same gradient flow:
\begin{equation}\label{eq-relative-entropy-three-representations}
\partial_t\rho_t
=
\sigma^2\diverg\!\left(\rho_t\nabla\log\frac{\rho_t}{\rho_\beta}\right),
\qquad
v_t
=
\sigma^2\bigl(\nabla\log\rho_\beta-\nabla\log\rho_t\bigr).
\end{equation}
\paragraph{Gradient flows under density constraints.}
\begin{equation}\label{eq-density-constrained-energy}
f_{\rho_{\max}}(\rho\d x)=\int_\Omega h(x)\rho(x)\d x+
\iota_{\mathcal K_{\rho_{\max}}}(\rho\d x),
\qquad
\mathcal K_{\rho_{\max}}=
\enscond{\rho\d x\in\Pp_2(\Omega)}{0\leq\rho\leq\rho_{\max}\ \text{a.e.}},
\end{equation}
where $\Omega\subset\RR^d$ is a convex domain and $\mathcal K_{\rho_{\max}}$ is assumed nonempty; if $\Omega$ has finite volume, this requires $\rho_{\max}|\Omega|\geq1$. The indicator term has no classical first variation; it contributes the normal cone to $\mathcal K_{\rho_{\max}}$ in the Wasserstein first-order condition. The JKO step becomes
\begin{equation}\label{eq-density-constrained-jko}
\alpha^{(\ell+1)}
\in
\uargmin{\alpha\in\mathcal K_{\rho_{\max}}}
\frac{1}{2\tau}\Wass_2^2(\alpha^{(\ell)},\alpha)+\int_\Omega h\,\d\alpha.
\end{equation}
With no-flux boundary condition $\rho_t\nabla(h+p_t)\cdot n=0$ on $\partial\Omega$, the constrained gradient flow is the complementarity system
\begin{equation}\label{eq-density-constrained-pressure-flow}
\partial_t\rho_t=\diverg\left(\rho_t\nabla(h+p_t)\right),
\qquad
0\leq\rho_t\leq\rho_{\max},
\qquad
p_t\geq0,
\qquad
p_t(\rho_{\max}-\rho_t)=0.
\end{equation}
The pressure vanishes in the unsaturated region and acts only where $\rho_t=\rho_{\max}$, pushing mass away from congested zones so that the density cap remains satisfied.

\paragraph{Multi-species gradient flows.}
Several applications involve several densities living on the same physical space: chemical species, color channels, cell populations or competing phases. Let $\Omega\subset\RR^d$ be the physical domain and assume that the component masses $\mathbf m=(m_1,\ldots,m_p)$ are positive and fixed. A convenient notation is
\begin{equation}\label{eq-multispecies-space}
\Pp_{2,\mathbf m}(\Omega;\RR_+^p)=
\enscond{\alpha=(\alpha_1,\ldots,\alpha_p)}
{\alpha_i\in\Mm_+(\Omega),\ \alpha_i(\Omega)=m_i,\ \int_\Omega\norm{x}^2\d\alpha_i(x)<+\infty}.
\end{equation}
Since the components are finite measures rather than necessarily probabilities, set
\begin{equation}\label{eq-multispecies-product-metric}
\Wass_{2,\oplus}^2(\alpha,\eta)=
\sum_{i=1}^p m_i\,
\Wass_2^2\!\left(\frac{\alpha_i}{m_i},\frac{\eta_i}{m_i}\right).
\end{equation}
The corresponding JKO step is
\begin{equation}\label{eq-multispecies-jko}
\alpha^{(\ell+1)}
\in
\uargmin{\alpha\in\Pp_{2,\mathbf m}(\Omega;\RR_+^p)}
\frac{1}{2\tau}\Wass_{2,\oplus}^2(\alpha^{(\ell)},\alpha)+f(\alpha).
\end{equation}
For smooth densities and first variations $\phi_i=\delta f/\delta\rho_i$, this yields
\begin{equation}\label{eq-multispecies-pde}
\partial_t\rho_i=\diverg\left(\rho_i\nabla\phi_i\right),
\qquad
\phi_i=\frac{\delta f}{\delta\rho_i},
\qquad i=1,\ldots,p.
\end{equation}
Given a fixed reference measure $\beta=b\,\d x$ with density $b=\density{\beta}>0$ and total mass $\sum_i m_i$, set
\begin{equation}\label{eq-multispecies-composition-constraint}
\mathcal C_\beta=
\enscond{\alpha\in\Pp_{2,\mathbf m}(\Omega;\RR_+^p)}
{\sum_{i=1}^p \alpha_i=\beta}.
\end{equation}
The constrained JKO step is
\begin{equation}\label{eq-multispecies-constrained-jko}
\alpha^{(\ell+1)}
\in
\uargmin{\alpha\in\mathcal C_\beta}
\frac{1}{2\tau}\Wass_{2,\oplus}^2(\alpha^{(\ell)},\alpha)+f(\alpha).
\end{equation}
In the smooth-density regime, the formal constrained flow has a common pressure $\lambda_t$:
\begin{equation}\label{eq-multispecies-constrained-pde}
\partial_t\rho_i=\diverg\left(\rho_i\nabla(\phi_i-\lambda_t)\right),
\qquad
\diverg(b\nabla\lambda_t)=\diverg\left(\sum_{i=1}^p \rho_i\nabla\phi_i\right).
\end{equation}
For the separable Shannon entropy $f(\alpha)=\sum_i\int_\Omega \rho_i\log\rho_i\d x$, this becomes
\begin{equation}\label{eq-multispecies-modified-heat}
\partial_t\rho_i=\Delta\rho_i-\diverg(\rho_i\nabla\lambda_t),
\qquad
\diverg(b\nabla\lambda_t)=\Delta b.
\end{equation}
These equations are understood with periodic or no-flux boundary conditions. When $b$ is constant on a connected domain, the pressure can be chosen constant and the system reduces to independent heat equations which nevertheless preserve the pointwise sum $\sum_i\rho_i=b$.

\subsection{Geodesic Convexity and Convergence}
\label{sec-geodesic-convexity}
\paragraph{Geodesics and convexity.}
A constant-speed $\Wass_2$ geodesic between $\alpha_0$ and $\alpha_1$ is obtained, as in Definition~\ref{def-w2-geodesic-induced-by-plan}, from any optimal coupling $\pi^\star\in\Couplings(\alpha_0,\alpha_1)$ by the McCann interpolation
\(\alpha_t=((1-t)P_0+tP_1)_\sharp\pi^\star, \qquad t\in[0,1],\)
where $P_0(x,y)=x$ and $P_1(x,y)=y$. If the optimal plan is induced by a Brenier map $T$, this reduces to $((1-t)\Id+tT)_\sharp\alpha_0$.

\begin{defn}[Geodesic convexity]\label{def-geodesic-convexity}
A functional $f$ on $\Pp_2(\RR^d)$ is geodesically convex if for every $\Wass_2$ geodesic $(\alpha_t)_t$, \(f(\alpha_t)\leq (1-t)f(\alpha_0)+t f(\alpha_1).\)
It is $\lambda$-geodesically convex if the right-hand side is improved by $-\frac{\lambda}{2}t(1-t)\Wass_2^2(\alpha_0,\alpha_1)$.
\end{defn}
\begin{prop}[Geodesic convexity of linear and quadratic energies]\phantomsection\label{prop-basic-geodesic-convexity}
\begin{enumerate}
\item If $h$ is convex, then the linear energy $\alpha\mapsto\int h\d\alpha$ is geodesically convex; if $h$ is $\lambda$-strongly convex, it is $\lambda$-geodesically convex. Conversely, if this linear energy is geodesically convex for all Dirac endpoints, then $h$ is convex.
\item More generally, let \(k\geq2\) and let \(H_k:(\RR^d)^k\to\RR\cup\{+\infty\}\) be convex. Then the polynomial energy
\(\alpha\mapsto \int_{(\RR^d)^k} H_k(x_1,\ldots,x_k)\,\d\alpha(x_1)\cdots\d\alpha(x_k)\)
is geodesically convex whenever the integral is well defined. If $H_k$ is $\lambda$-strongly convex on the product space, this energy is $k\lambda$-geodesically convex. In particular, the quadratic interaction
\(\alpha\mapsto\frac12\iint W(x,y)\d\alpha(x)\d\alpha(y)\)
is geodesically convex when \(W\) is convex on \(\RR^d\times\RR^d\), and it is $\lambda$-geodesically convex when $W$ is $\lambda$-strongly convex there.
\end{enumerate}
\end{prop}
\begin{proof}
Let \(\pi\) be an optimal coupling between \(\alpha_0\) and \(\alpha_1\), and set \(X_t=(1-t)X_0+tX_1\) under \((X_0,X_1)\sim\pi\). Convexity of \(h\) gives \(h(X_t)\leq(1-t)h(X_0)+t h(X_1)\), and strong convexity gives the additional term \(-\frac{\lambda}{2}t(1-t)\norm{X_0-X_1}^2\). Integrating proves the first claim. Necessity follows by testing on Dirac masses: the geodesic between \(\delta_x\) and \(\delta_y\) is \(\delta_{(1-t)x+ty}\), so geodesic convexity of \(\int h\d\alpha\) gives the usual convexity inequality for \(h\).

For the polynomial statement, take \(k\) independent copies \((X_0^\ell,X_1^\ell)_{\ell=1}^k\) of the same optimal coupling and define \(X_t^\ell=(1-t)X_0^\ell+tX_1^\ell\). Convexity of \(H_k\) on the product space gives
\(H_k(X_t^1,\ldots,X_t^k) \leq (1-t)H_k(X_0^1,\ldots,X_0^k)+tH_k(X_1^1,\ldots,X_1^k).\)
Integrating over the product coupling proves the claim. If $H_k$ is $\lambda$-strongly convex, the additional term is
\(-\frac{\lambda}{2}t(1-t) \sum_{\ell=1}^k\EE\norm{X_0^\ell-X_1^\ell}^2 = -\frac{k\lambda}{2}t(1-t)\Wass_2^2(\alpha_0,\alpha_1),\)
which proves the stated constant. For the quadratic interaction, the prefactor $1/2$ cancels the factor $k=2$, leaving the constant $\lambda$.
\end{proof}
\begin{prop}[One-dimensional interactions and energy-distance MMD]\label{prop-1d-interaction-halfline-convex}
Let \(\varphi:\RR_+\to\RR\) be continuous with at most quadratic growth, and define on \(\Pp_2(\RR)\)
\(\mathcal I_\varphi(\alpha) \eqdef \frac12\iint \varphi(|x-y|)\d\alpha(x)\d\alpha(y).\)
Then \(\mathcal I_\varphi\) is geodesically convex for \(\Wass_2\) on \(\Pp_2(\RR)\) if and only if \(\varphi\) is convex on \(\RR_+\).
Fix also \(\beta\in\Pp_2(\RR)\). For the conditionally positive kernel \(k(x,y)=-|x-y|\) from Definition~\ref{def-positive-kernels}, the squared MMD of Definition~\ref{def-kernel-mmd-norm}, equivalently the squared energy distance,
\(\mathcal E_\beta(\alpha) \eqdef \MMD_k^2(\alpha,\beta) = \iint -|x-y|\,\d(\alpha-\beta)(x)\d(\alpha-\beta)(y),\)
is geodesically convex as a function of \(\alpha\).
\end{prop}
\begin{proof}
Let \(Q_0,Q_1\) be the quantile functions of \(\alpha_0,\alpha_1\), and let \(Q_t=(1-t)Q_0+tQ_1\) be the quantile function of the one-dimensional \(\Wass_2\) geodesic \(\alpha_t\). For \(r>s\), monotonicity gives \(Q_i(r)-Q_i(s)\geq0\) for \(i=0,1\), and hence
\(Q_t(r)-Q_t(s) = (1-t)(Q_0(r)-Q_0(s))+t(Q_1(r)-Q_1(s)).\)
Using the symmetry of the double integral and the quantile parametrization, one can write
\(\mathcal I_\varphi(\alpha_t) = \int_0^1\int_0^r \varphi(Q_t(r)-Q_t(s))\,\d s\,\d r.\)
Convexity of \(\varphi\) on \(\RR_+\) gives the desired convexity after integration.

Conversely, test geodesic convexity on two-point measures \(\alpha_i=\frac12(\delta_0+\delta_{a_i})\), with \(a_i\geq0\). Their monotone geodesic is \(\alpha_t=\frac12(\delta_0+\delta_{(1-t)a_0+t a_1})\). Since
\(\mathcal I_\varphi(\alpha_t) = \frac14\varphi(0)+\frac14\varphi((1-t)a_0+t a_1),\)
and the identical \(\frac14\varphi(0)\) terms cancel from the geodesic-convexity inequality, one obtains the ordinary convexity inequality for \(\varphi\) on \(\RR_+\).

For the energy distance, the function \(\varphi(r)=-r\) is affine on \(\RR_+\), so the self-interaction \(-\iint |x-y|\d\alpha(x)\d\alpha(y)\) is affine along one-dimensional Wasserstein geodesics even though \(z\mapsto-|z|\) is not convex on \(\RR\). Write \(Q_\beta\) for the quantile function of \(\beta\). Expanding the MMD and dropping the term depending only on \(\beta\) gives
\[
\mathcal E_\beta(\alpha)
=
2\int_0^1\!\!\int_0^1 |Q_\alpha(r)-Q_\beta(s)|\,\d r\,\d s
-
\int_0^1\!\!\int_0^1 |Q_\alpha(r)-Q_\alpha(s)|\,\d r\,\d s
+\mathrm{constant}.
\]
Here the constant is independent of \(\alpha\). The positive cross-distance term is convex in \(Q_\alpha\), since the absolute value is convex and \(Q_t\) is affine in \(t\). The self-distance term is affine along quantile geodesics because
\(\int_0^1\!\!\int_0^1 |Q_t(r)-Q_t(s)|\,\d r\,\d s = 2\int_0^1\int_0^r (Q_t(r)-Q_t(s))\,\d s\,\d r.\)
Thus \(\alpha\mapsto\mathcal E_\beta(\alpha)\) is geodesically convex.
\end{proof}
It should not be read as a general MMD statement: for a typical positive definite kernel \(k\), the decomposition
\(\MMD_k^2(\alpha,\beta) = \iint k\,\d\alpha\d\alpha -2\iint k\,\d\alpha\d\beta +\mathrm{constant}\)
\begin{thm}[McCann displacement convexity for internal energies]\label{thm-mccann-internal-energy}
Let \(g:[0,+\infty)\to\RR\cup\{+\infty\}\) be proper, lower-semicontinuous and convex, with \(g(0)=0\), set
\(g'_\infty\eqdef\lim_{r\to+\infty}\frac{g(r)}r\in\RR\cup\{+\infty\},\)
and define
\(\psi(r)=r^d g(r^{-d}),\qquad r>0.\)
If \(\psi\) is convex and nonincreasing on \((0,+\infty)\), then the relaxed internal energy
\(\overline{\mathcal U}_g(\alpha) \eqdef \int_{\RR^d}g(\rho(x))\,\d x +g'_\infty\alpha^\perp(\RR^d), \qquad \alpha=\rho\,\d x+\alpha^\perp,\)
is geodesically convex on its effective domain in \(\Pp_2(\RR^d)\). If, in addition,
\[
g^-(r)\leq C r^q
\qquad\text{for some }q\in\left(\frac{d}{d+2},1\right),
\]
then it never takes the value \(-\infty\) and is lower semicontinuous for \(\Wass_2\) convergence. On a bounded convex domain, lower semicontinuity holds without this tail condition. The convention is \(+\infty\cdot0=0\); when \(g'_\infty=+\infty\), finite energy forces \(\alpha^\perp=0\).
\end{thm}
\begin{proof}
It is enough to establish the inequality first for smooth positive compactly supported densities. Let \(T=\nabla\varphi\) be the Brenier map from \(\alpha_0=\rho_0\d x\) to \(\alpha_1=\rho_1\d x\), set \(T_t=(1-t)\Id+tT\), and denote
\(J_t(x)=\det\big((1-t)\Id+t\nabla T(x)\big)\).
The interpolation satisfies \(\alpha_t=(T_t)_\sharp\alpha_0\) and
\(\rho_t(T_t(x))J_t(x)=\rho_0(x)\).
By concavity of \(\det^{1/d}\) on positive semidefinite matrices, \(J_t(x)^{1/d}\geq (1-t)+t\det(\nabla T(x))^{1/d}.\)
Introduce \(s_t(x)=(J_t(x)/\rho_0(x))^{1/d}\). Since
\(\rho_1(T(x))\det\nabla T(x)=\rho_0(x)\), this gives
\(s_t(x)\geq (1-t)\rho_0(x)^{-1/d}+t\rho_1(T(x))^{-1/d}.\)
Using the change of variables \(y=T_t(x)\), one obtains
\(\int g(\rho_t(y))\,\d y =\int \rho_0(x)\,\psi(s_t(x))\,\d x.\)
Because \(\psi\) is nonincreasing and convex, \(\psi(s_t(x)) \leq (1-t)\psi(\rho_0(x)^{-1/d})+t\psi(\rho_1(T(x))^{-1/d}).\)
Multiplying by \(\rho_0\) and integrating gives
\(\overline{\mathcal U}_g(\alpha_t)\leq(1-t)\overline{\mathcal U}_g(\alpha_0)+t\overline{\mathcal U}_g(\alpha_1)\).
The recession term is the measure relaxation that accounts for concentration of mass, and approximation gives geodesic convexity on the effective domain. For lower semicontinuity, the positive part follows from the standard convex-integral relaxation. On \(\RR^d\), the displayed bound on \(g^-\), together with \(q>d/(d+2)\), makes the negative parts uniformly integrable on sets with bounded second moment; on bounded domains, tightness supplies this control directly. This yields the stated \(\Wass_2\)-lower-semicontinuity; see.
\end{proof}
\begin{prop}[Geodesic convexity of power divergences]\label{prop-power-divergences-geodesic-convexity}\label{prop-kl-gibbs-geodesic-convexity}
For \(m>0\), define
\[
\phi_m(r)=
\begin{cases}
\dfrac{r^m-mr+m-1}{m(m-1)}, & m\neq1,\\[2mm]
r\log r-r+1, & m=1.
\end{cases}
\]
Then the following geodesic-convexity statements hold.
\begin{enumerate}
\item If \(\Omega\subset\RR^d\) is a bounded convex domain, \(\beta=\density{\beta}\d x\) is uniform on \(\Omega\), and \(m\geq1-1/d\), then \(\alpha\mapsto\Divergm_{\phi_m}(\alpha|\beta)\) is geodesically convex on \(\Pp_2(\Omega)\).
\item If \(\beta=Z^{-1}e^{-V}\d x\) on \(\RR^d\), \(V\) is convex, and \(m\geq1\), then \(\alpha\mapsto\Divergm_{\phi_m}(\alpha|\beta)\) is geodesically convex on \((\Pp_2(\RR^d),\Wass_2)\).
\item In the KL case \(m=1\), if \(V\) is \(\lambda\)-strongly convex, then \(\alpha\mapsto\KL(\alpha|\beta)=\Divergm_{\phi_1}(\alpha|\beta)\) is \(\lambda\)-geodesically convex.
\end{enumerate}
The divergence uses the recession convention of Definition~\ref{def_divergence}. Thus singular mass has infinite cost for \(m\geq1\), whereas for \(0<m<1\) its cost is \(\alpha^\perp(\Omega)/(1-m)\).
\end{prop}
\begin{proof}
For the flat reference case, write \(\alpha=\rho\d x\) and let \(\density{\beta}=b\,\mathds{1}_\Omega\), where $b=1/|\Omega|$. Up to affine terms in \(\rho\), which only add constants on probability measures, the integrand \(b\,\phi_m(\rho/b)\) has non-affine part \(b^{1-m}\rho^m/(m(m-1))\) for \(m\neq1\), and \(\rho\log\rho\) for \(m=1\). For \(m=1\), McCann's criterion gives \(\psi(r)=-d\log r\). For \(m\neq1\), the transformed non-affine part is, up to the irrelevant factor \(b^{1-m}\),
\(\psi_m(r)=\frac{r^{d(1-m)}}{m(m-1)}\).
If \(m>1\), this function is convex and nonincreasing. If \(0<m<1\), it is convex and nonincreasing exactly when \(d(1-m)\leq1\), i.e.\ \(m\geq1-1/d\). The affine density term and the recession contribution combine with this relaxed internal energy to give precisely Definition~\ref{def_divergence}. This proves the first claim by Theorem~\ref{thm-mccann-internal-energy}.

Assume next that \(m>1\) and \(V\) is convex. Up to constants, \(\Divergm_{\phi_m}(\alpha|\beta) = \frac{Z^{m-1}}{m(m-1)} \int \rho(x)^m e^{(m-1)V(x)}\,\d x.\)
Let \(T=\nabla\varphi\) be the Brenier map from \(\alpha_0=\rho_0\d x\) to \(\alpha_1\), set \(T_t=(1-t)\Id+tT\), and write
\(J_t(x)=\det((1-t)\Id+t\nabla T(x))\). Along the geodesic \(\alpha_t=(T_t)_\sharp\alpha_0\),
\[
\int \rho_t^m e^{(m-1)V}\d x
=
\int \rho_0(x)^m
\exp\!\left((m-1)\big(V(T_t(x))-\log J_t(x)\big)\right)\d x.
\]
For each \(x\), the function \(t\mapsto V(T_t(x))\) is convex, and \(t\mapsto-\log J_t(x)\) is convex because \(\log\det\) is concave on positive definite matrices. Since \(m-1>0\), the exponential of their sum is convex in \(t\). Integrating gives geodesic convexity. The nonsmooth case follows by approximation.

Finally, for \(m=1\), \(\KL(\alpha|\beta) = \int \rho\log\rho\,\d x+\int V\d\alpha+\log Z.\)
The entropy term is \(0\)-geodesically convex by Theorem~\ref{thm-mccann-internal-energy}. The linear potential term is \(\lambda\)-geodesically convex when \(V\) is \(\lambda\)-strongly convex, by Proposition~\ref{prop-basic-geodesic-convexity}. Taking \(\lambda=0\) also covers the \(m=1\) case in the second claim.
\end{proof}
\paragraph{Geodesically convex constraints.}
\label{par-geodesically-convex-constraints}
\begin{proposition}[Geodesic convexity of density caps]\label{prop-density-cap-geodesic-convex}
Assume that $\Omega\subset\RR^d$ is convex and let $\rho_{\max}>0$. The set $\mathcal K_{\rho_{\max}}$ in~\eqref{eq-density-constrained-energy} is geodesically convex in $\Pp_2(\Omega)$: if $\alpha_0,\alpha_1\in\mathcal K_{\rho_{\max}}$, then the constant-speed $\Wass_2$ geodesic $(\alpha_t)_{t\in[0,1]}$ between them also belongs to $\mathcal K_{\rho_{\max}}$.
\end{proposition}
\begin{proof}
Assume first that the endpoints have smooth positive densities $\alpha_i=\rho_i\d x$ with $\rho_i\leq\rho_{\max}$, and let $T=\nabla\varphi$ be the Brenier map from $\alpha_0$ to $\alpha_1$. Since $\Omega$ is convex, $T_t=(1-t)\Id+tT$ maps $\Omega$ into $\Omega$. The interpolation is $\alpha_t=(T_t)_\sharp\alpha_0$ and
\(\rho_t(T_t(x))=\frac{\rho_0(x)}{\det((1-t)\Id+t\nabla T(x))}.\)
The concavity of $\det^{1/d}$ on positive semidefinite matrices and the identity $\rho_1(T(x))\det\nabla T(x)=\rho_0(x)$ give
\(\rho_t(T_t(x))^{-1/d} \geq (1-t)\rho_0(x)^{-1/d}+t\rho_1(T(x))^{-1/d} \geq\rho_{\max}^{-1/d}.\)
Hence $\rho_t\leq\rho_{\max}$ for all $t$. The general case follows by approximation and lower semicontinuity of the $L^\infty$ bound, as in McCann's displacement-convexity theory.
\end{proof}
\paragraph{Dissipation and finite metric length.}
For the classical gradient flow of a smooth function $h:\RR^d\to\RR$, $\dot x_t=-\nabla h(x_t)$, the chain rule gives \(\frac{\d}{\d t}h(x_t) = -\norm{\nabla h(x_t)}^2 = -\norm{\dot x_t}^2.\)
Thus, for $0\leq s<t$, Cauchy's inequality gives
\begin{equation}\label{eq-euclidean-gradient-flow-length}
\int_s^t \norm{\dot x_r}\,\d r
\leq
\sqrt{t-s}\left(\int_s^t\norm{\dot x_r}^2\d r\right)^{1/2}
=
\sqrt{t-s}\,\bigl(h(x_s)-h(x_t)\bigr)^{1/2}.
\end{equation}
\begin{prop}[Finite metric length from dissipation]\label{prop-gradient-flow-finite-length}
Let $(\alpha_r)_{r\in[0,T]}$ be an absolutely continuous curve in $(\Pp_2(\RR^d),\Wass_2)$. Assume that, for all $0\leq s<t\leq T$, it satisfies the exact energy-dissipation identity
\(f(\alpha_s)-f(\alpha_t) = \int_s^t |\dot\alpha_r|_{\Wass_2}^2\d r.\)
Then
\begin{equation}\label{eq-wasserstein-gradient-flow-length}
\operatorname{Length}_{\Wass_2}\bigl((\alpha_r)_{r\in[s,t]}\bigr)
=
\int_s^t |\dot\alpha_r|_{\Wass_2}\,\d r
\leq
\sqrt{t-s}\,\bigl(f(\alpha_s)-f(\alpha_t)\bigr)^{1/2}.
\end{equation}
In the smooth Wasserstein gradient-flow setting with velocity $v_r=-\Wgrad f(\alpha_r)$, the dissipation identity reads equivalently
\(f(\alpha_s)-f(\alpha_t) = \int_s^t |\dot\alpha_r|_{\Wass_2}^2\d r = \int_s^t\int \norm{v_r(x)}^2\d\alpha_r(x)\d r.\)
If one assumes only the metric energy-dissipation inequality for a curve of maximal slope, then the same length estimate holds with the right-hand side multiplied by $\sqrt2$.
\end{prop}
\begin{proof}
Cauchy's inequality gives
\[
\int_s^t |\dot\alpha_r|_{\Wass_2}\,\d r
\leq
\sqrt{t-s}
\left(\int_s^t |\dot\alpha_r|_{\Wass_2}^2\d r\right)^{1/2}.
\]
The exact energy-dissipation identity substitutes the last integral by $f(\alpha_s)-f(\alpha_t)$, which proves~\eqref{eq-wasserstein-gradient-flow-length}. In the metric theory of curves of maximal slope, the energy-dissipation inequality gives instead
\(\frac12\int_s^t |\dot\alpha_r|_{\Wass_2}^2\d r \leq f(\alpha_s)-f(\alpha_t),\)
and the same argument yields the additional factor $\sqrt2$.
\end{proof}
If \(m_T\leq f(\alpha_r)\) for \(0\leq r\leq T\), then for \(0\leq s<t\leq T\), \(\Wass_2(\alpha_s,\alpha_t) \leq \operatorname{Length}_{\Wass_2}\bigl((\alpha_r)_{r\in[s,t]}\bigr) \leq \sqrt{f(\alpha_0)-m_T}\,\sqrt{t-s},\)
with an additional factor \(\sqrt2\) under the energy-dissipation inequality. Thus an energy-dissipating trajectory is \(1/2\)-H\"older continuous as a curve in Wasserstein space on every time interval where the energy remains bounded from below.

\begin{prop}[Energy decay for convex Wasserstein flows]\label{prop-convex-wass-flow-rate}
Assume formally that \(f\) is geodesically convex, admits a smooth first variation, and has a minimizer \(\alpha^\star\). Let \((\alpha_t)_t\) be a smooth solution of the Wasserstein gradient flow
\(\partial_t\alpha_t+\diverg(\alpha_t v_t)=0, \qquad v_t=-\Wgrad f(\alpha_t).\)
Then
\(\frac{\d}{\d t} f(\alpha_t) = -\int \norm{\Wgrad f(\alpha_t)(x)}^2\,\d\alpha_t(x) \leq 0.\)
If \(T_t\) is the optimal map from \(\alpha_t\) to \(\alpha^\star\), then
\(f(\alpha_t)-f(\alpha^\star) \leq -\frac{\d}{\d t}\frac12 \Wass_2^2(\alpha_t,\alpha^\star),\)
and consequently
\(f(\alpha_t)-f(\alpha^\star) \leq \frac{\Wass_2^2(\alpha_0,\alpha^\star)}{2t}.\)
\end{prop}
\begin{proof}
The chain rule and Proposition~\ref{prop-formal-wass-gradient} give
\(\frac{\d}{\d t}f(\alpha_t) = \int \dotp{\Wgrad f(\alpha_t)(x)}{v_t(x)}\,\d\alpha_t(x) = -\int \norm{\Wgrad f(\alpha_t)(x)}^2\,\d\alpha_t(x).\)
Geodesic convexity along the geodesic
\(((1-s)\Id+sT_t)_\sharp\alpha_t\) gives
\(f(\alpha^\star)-f(\alpha_t) \geq \int \dotp{\Wgrad f(\alpha_t)(x)}{T_t(x)-x}\,\d\alpha_t(x).\)
Since \(v_t=-\Wgrad f(\alpha_t)\), this reads \(f(\alpha_t)-f(\alpha^\star) \leq \int \dotp{v_t(x)}{T_t(x)-x}\,\d\alpha_t(x).\)
\(\frac{\d}{\d t}\frac12\Wass_2^2(\alpha_t,\alpha^\star) = \int \dotp{x-T_t(x)}{v_t(x)}\,\d\alpha_t(x),\)
which proves the differential inequality. Integrating it from \(0\) to \(t\) and using the monotonicity of \(s\mapsto f(\alpha_s)\) gives
\[
t\bigl(f(\alpha_t)-f(\alpha^\star)\bigr)
\leq
\int_0^t \bigl(f(\alpha_s)-f(\alpha^\star)\bigr)\,\d s
\leq
\frac12\Wass_2^2(\alpha_0,\alpha^\star).
\]
\end{proof}
\paragraph{Convergence: the Wasserstein--PL viewpoint.}
The relevant quantity is the squared Wasserstein slope
\(|\partial f|^2(\alpha) \eqdef \int\norm{\Wgrad f(\alpha)(x)}^2\,\d\alpha(x)\)
\begin{defn}[Wasserstein--Polyak--{\L}ojasiewicz inequality]\label{def-wasserstein-pl}
Let \(f_\star\eqdef\inf_{\alpha\in\Pp_2(\RR^d)} f(\alpha)>-\infty\). The functional \(f\) satisfies the Wasserstein--PL inequality with constant \(\kappa>0\) if
\begin{equation}\label{eq-wasserstein-pl}
|\partial f|^2(\alpha)
\geq
2\kappa\bigl(f(\alpha)-f_\star\bigr)
\end{equation}
for every \(\alpha\) in the domain of \(f\). In the smooth Otto notation this reads \(\int\norm{\Wgrad f(\alpha)}^2\,\d\alpha \geq 2\kappa\bigl(f(\alpha)-f_\star\bigr).\)
\end{defn}
\begin{thm}[Wasserstein--PL convergence]\label{thm-wasserstein-pl-convergence}
Assume that \(f_\star>-\infty\), that \(f\) satisfies the Wasserstein--PL inequality~\eqref{eq-wasserstein-pl} with constant \(\kappa>0\), and that \((\alpha_t)_{t\geq0}\) is a global Wasserstein gradient flow satisfying the full energy-dissipation identity
\begin{equation}\label{eq-full-wasserstein-energy-dissipation}
\frac{\d}{\d t}f(\alpha_t)
=
-|\dot\alpha_t|^2
=
-|\partial f|^2(\alpha_t)
\qquad\text{for a.e.\ }t>0.
\end{equation}
Set \(E(t)\eqdef f(\alpha_t)-f_\star\). Then, for \(0\leq s\leq t\), \(E(t)\leq e^{-2\kappa(t-s)}E(s),\)
and the length of the flow segment satisfies \(\operatorname{Length}_{\Wass_2}\bigl((\alpha_r)_{r\in[s,t]}\bigr) \leq \sqrt{\frac{2}{\kappa}}\bigl(\sqrt{E(s)}-\sqrt{E(t)}\bigr).\)
If, in addition, \(f\) is lower semicontinuous for \(\Wass_2\)-convergence, then \(\alpha_t\) converges in \(\Wass_2\) to a minimizer \(\alpha_\infty\in\operatorname{Argmin} f\), and
\(\Wass_2(\alpha_t,\alpha_\infty) \leq \sqrt{\frac{2}{\kappa}}\sqrt{E(t)} \leq \sqrt{\frac{2}{\kappa}}\sqrt{E(0)}\,e^{-\kappa t}.\)
In particular, \(\operatorname{dist}_{\Wass_2}(\alpha_t,\operatorname{Argmin}f) \leq \sqrt{\frac{2}{\kappa}}\sqrt{E(0)}\,e^{-\kappa t}.\)
\end{thm}
\begin{proof}
The energy estimate is immediate from~\eqref{eq-full-wasserstein-energy-dissipation} and~\eqref{eq-wasserstein-pl}: \(E'(t) = -|\partial f|^2(\alpha_t) \leq -2\kappa E(t)\)
for almost every \(t\). Gronwall's lemma gives \(E(t)\leq e^{-2\kappa(t-s)}E(s)\).

The distance estimate uses the same two identities in a slightly different way. On the set where \(E(r)>0\), the PL inequality gives \(|\partial f|(\alpha_r)\geq \sqrt{2\kappa E(r)}.\)
Using the full dissipation identity and \(|\dot\alpha_r|=|\partial f|(\alpha_r)\),
\[
\begin{aligned}
\operatorname{Length}_{\Wass_2}\bigl((\alpha_r)_{r\in[s,t]}\bigr)
&=
\int_s^t|\dot\alpha_r|\,\d r
=
\int_s^t|\partial f|(\alpha_r)\,\d r \\
&=
\int_s^t
\frac{|\partial f|^2(\alpha_r)}{|\partial f|(\alpha_r)}\,\d r
\leq
\frac{1}{\sqrt{2\kappa}}
\int_s^t
\frac{|\partial f|^2(\alpha_r)}{\sqrt{E(r)}}\,\d r \\
&=
\frac{1}{\sqrt{2\kappa}}
\int_s^t
\frac{-E'(r)}{\sqrt{E(r)}}\,\d r
=
\sqrt{\frac{2}{\kappa}}\bigl(\sqrt{E(s)}-\sqrt{E(t)}\bigr).
\end{aligned}
\]
If \(E\) vanishes on a subinterval, then the flow is stationary there by~\eqref{eq-full-wasserstein-energy-dissipation}, so the same estimate follows by applying the argument on the part where \(E>0\). Since Wasserstein distance is bounded by metric length, the same upper bound controls \(\Wass_2(\alpha_s,\alpha_t)\).

Letting \(t\to+\infty\), the energy estimate gives \(E(t)\to0\), and the length estimate shows that the tail of \((\alpha_t)_t\) is Cauchy. Since \((\Pp_2(\RR^d),\Wass_2)\) is complete, there is a limit \(\alpha_\infty\). Lower semicontinuity yields
\(f(\alpha_\infty) \leq \liminf_{t\to+\infty}f(\alpha_t) = f_\star,\)
hence \(\alpha_\infty\in\operatorname{Argmin}f\). Sending the endpoint of the length estimate to \(+\infty\) with \(s=t\) fixed gives \(\Wass_2(\alpha_t,\alpha_\infty) \leq \sqrt{\frac{2}{\kappa}}\sqrt{E(t)}.\)
The distance-to-the-set bound follows because \(\alpha_\infty\) is one minimizer.
\end{proof}
\paragraph{Geodesic strong convexity.}
\begin{prop}[Strong geodesic convexity implies Wasserstein--PL]\label{prop-strong-geodesic-convexity-implies-pl}
Assume that \(f\) is \(\lambda\)-geodesically convex with \(\lambda>0\), that \(f_\star\) is attained at some \(\alpha^\star\), and that the Wasserstein slope is finite at \(\alpha\). Then
\(|\partial f|^2(\alpha) \geq 2\lambda\bigl(f(\alpha)-f_\star\bigr).\)
Thus positive Wasserstein strong convexity implies the Wasserstein--PL inequality with the same constant.
\end{prop}
\begin{proof}
Fix \(\alpha\) and a minimizer \(\alpha^\star\). Set
\(g=f(\alpha)-f_\star,\qquad r=\Wass_2(\alpha,\alpha^\star),\qquad a=|\partial f|(\alpha).\)
If \(r=0\), then \(\alpha=\alpha^\star\) and there is nothing to prove. Otherwise let \((\alpha_s)_{s\in[0,1]}\) be a constant-speed geodesic from \(\alpha\) to \(\alpha^\star\). By \(\lambda\)-geodesic convexity,
\(f(\alpha_s) \leq (1-s)f(\alpha)+s f_\star -\frac{\lambda}{2}s(1-s)r^2.\)
Therefore
\(f(\alpha)-f(\alpha_s) \geq s g+\frac{\lambda}{2}s(1-s)r^2.\)
Dividing by \(\Wass_2(\alpha,\alpha_s)=sr\) and letting \(s\downarrow0\) in the definition of the descending metric slope yields
\(a\geq \frac{g}{r}+\frac{\lambda}{2}r\).
Equivalently, \(g\leq ar-\lambda r^2/2\). Maximizing the right-hand side over \(r\geq0\) gives \(g\leq a^2/(2\lambda)\), which is the desired PL inequality.
\end{proof}
\paragraph{Wasserstein Kurdyka--{\L}ojasiewicz inequalities.}
\begin{defn}[Wasserstein Kurdyka--{\L}ojasiewicz inequality]\label{def-wasserstein-kl}
Let \(f_\star=\inf f\), let \(E(\alpha)=f(\alpha)-f_\star\), and fix \(c>0\), \(\theta\in[0,1)\). We say that \(f\) satisfies a power Wasserstein--KL inequality on a set \(\mathcal U\subset \Pp_2(\RR^d)\) if
\(|\partial f|(\alpha) \geq c\,E(\alpha)^\theta, \qquad \alpha\in\mathcal U, \quad 0<E(\alpha)<+\infty.\)
Equivalently, the desingularizing function
\(\psi(s)=\frac{s^{1-\theta}}{c(1-\theta)}\)
satisfies \(\psi'(E(\alpha))|\partial f|(\alpha)\geq 1\). The Wasserstein--PL inequality \eqref{eq-wasserstein-pl} is the special case \(\theta=1/2\) with \(c=\sqrt{2\kappa}\). The regime \(\theta>1/2\) gives sublinear rates, while \(\theta=1/2\) is exactly the exponential PL regime.
\end{defn}
\begin{thm}[Sublinear convergence under Wasserstein--KL]\label{thm-wasserstein-kl-sublinear}
Let \((\alpha_t)_{t\geq0}\) be a Wasserstein gradient flow of \(f\) satisfying the full energy-dissipation identity \eqref{eq-full-wasserstein-energy-dissipation}. Assume that \(f_\star> -\infty\), that \(E(t)=f(\alpha_t)-f_\star\) is finite, and that the trajectory stays in a region where the Wasserstein--KL inequality of Definition~\ref{def-wasserstein-kl} holds with \(\theta\in(1/2,1)\). If \(E(0)>0\), then
\[
E(t)
\leq
\left(E(0)^{1-2\theta}+c^2(2\theta-1)t\right)^{-1/(2\theta-1)}.
\]
Moreover, for every \(t\geq0\), \(\operatorname{Length}_{\Wass_2}\bigl((\alpha_s)_{s\geq t}\bigr) \leq \frac{E(t)^{1-\theta}}{c(1-\theta)}.\)
If \(f\) is lower semicontinuous on \((\Pp_2(\RR^d),\Wass_2)\), then \(\alpha_t\) converges to some \(\alpha_\infty\in\operatorname{Argmin}f\) and
\[
\Wass_2(\alpha_t,\alpha_\infty)
\leq
\frac{E(t)^{1-\theta}}{c(1-\theta)}
=
O\!\left(t^{-(1-\theta)/(2\theta-1)}\right).
\]
\end{thm}
\begin{proof}
If \(E\) reaches zero at some time, then the flow is already at a global minimizer. Indeed, the integrated energy-dissipation identity forces both \(|\dot\alpha_t|\) and \(|\partial f|(\alpha_t)\) to vanish afterwards. We therefore argue on intervals where \(E>0\). For a.e.\ \(t\),
\(-E'(t)=|\partial f|^2(\alpha_t)=|\dot\alpha_t|^2.\)
The Wasserstein--KL inequality yields
\(E'(t) \leq -c^2E(t)^{2\theta}\).
Since \(1-2\theta<0\), \(\frac{\d}{\d t}E(t)^{1-2\theta} = (1-2\theta)E(t)^{-2\theta}E'(t) \geq c^2(2\theta-1),\)
and integration from \(0\) to \(t\) gives the announced polynomial bound.

For the length estimate, use again \(|\dot\alpha_t|=|\partial f|(\alpha_t)\). For \(T>t\),
\[
\begin{aligned}
\int_t^T |\dot\alpha_s|\,\d s
&=
\int_t^T \frac{|\partial f|^2(\alpha_s)}{|\partial f|(\alpha_s)}\,\d s
\leq
\frac1c\int_t^T -E'(s)E(s)^{-\theta}\,\d s \\
&=
\frac{E(t)^{1-\theta}-E(T)^{1-\theta}}{c(1-\theta)}
\leq
\frac{E(t)^{1-\theta}}{c(1-\theta)}.
\end{aligned}
\]
Letting \(T\to+\infty\) shows that the tail has finite length. Hence \((\alpha_t)_t\) is Cauchy in \((\Pp_2(\RR^d),\Wass_2)\), which is complete, and converges to some \(\alpha_\infty\). Since the energy bound gives \(E(t)\to0\), lower semicontinuity gives \(f(\alpha_\infty)=f_\star\). The distance to \(\alpha_\infty\) is bounded by the remaining tail length, which proves the last estimate.
\end{proof}
\begin{prop}[Powers of PL energies are KL]\label{prop-powers-pl-are-kl}
Let \(g\geq0\) satisfy \(\inf g=0\) and the Wasserstein--PL inequality
\(|\partial g|^2(\alpha) \geq 2\kappa g(\alpha)\)
on a region where \(g>0\). Fix \(r\geq1\) and set \(f(\alpha)=g(\alpha)^r\). Whenever the metric-slope chain rule \(|\partial(g^r)|(\alpha)=r g(\alpha)^{r-1}|\partial g|(\alpha)\) holds, \(f\) satisfies the Wasserstein--KL inequality
\(|\partial f|(\alpha) \geq r\sqrt{2\kappa}\, f(\alpha)^{1-1/(2r)}.\)
Thus \(f\) has KL exponent \(\theta=1-1/(2r)\). For \(r=1\) this is PL; for \(r>1\), Theorem~\ref{thm-wasserstein-kl-sublinear} gives \(f(\alpha_t)=O(t^{-r/(r-1)})\) along the \(f\)-gradient flow, under its hypotheses.
\end{prop}
\begin{proof}
By the chain rule and the PL inequality for \(g\), \(|\partial f| = rg^{r-1}|\partial g| \geq r\sqrt{2\kappa}\,g^{r-1/2} = r\sqrt{2\kappa}\,f^{1-1/(2r)}.\)
The exponent and rate follow by substituting \(\theta=1-1/(2r)\) into Theorem~\ref{thm-wasserstein-kl-sublinear}.
\end{proof}
\paragraph{Convexity and curvature.}
The same language is not restricted to subsets of $\RR^d$. If $(\X,\dist,\mathfrak m)$ is a geodesic metric-measure space, $\Wass_2$ geodesics can be defined by transporting each pair of endpoints along metric geodesics, or more intrinsically by dynamical optimal plans on path space, as discussed in Section~\ref{sec-kantorovich-plan-interpolation}. Given a reference measure $\mathfrak m$, the entropy relative to $\mathfrak m$ is
\[
\mathrm{Ent}_{\mathfrak m}(\alpha)
\eqdef
\begin{cases}
\displaystyle \int_\X \rho\log\rho\,\d\mathfrak m,
& \text{if } \alpha=\rho\,\mathfrak m,\\
+\infty,
& \text{otherwise.}
\end{cases}
\]
\begin{thm}[Ricci curvature and entropy convexity]\label{thm-ricci-entropy-convexity}
Let $(M,g)$ be a smooth compact connected Riemannian manifold without boundary, and let $\mathfrak m=\mathrm{vol}_g$. For $\lambda\in\RR$, the lower Ricci bound $\mathrm{Ric}_g\geq\lambda g$ holds if and only if $\mathrm{Ent}_{\mathfrak m}$ is $\lambda$-geodesically convex on $(\Pp_2(M),\Wass_2)$.
\end{thm}
\subsection{Functional Inequalities via Optimal Transport}
Optimal transport proves inequalities by turning comparison into geometry. One transports a density by a monotone or Brenier map, interpolates along the resulting rays, and converts concavity of a Jacobian determinant or convexity of an energy into an integral estimate.

\paragraph{Brunn--Minkowski and isoperimetry.}
\begin{prop}[Brunn--Minkowski and Euclidean isoperimetry]\label{prop-ot-brunn-minkowski-isoperimetric}
Let $A,B\subset\RR^d$ be bounded Borel sets with positive Lebesgue measure. For $t\in[0,1]$, set
\((1-t)A+tB\eqdef\{(1-t)x+ty\;:\;x\in A,\ y\in B\}.\)
Then
\(\mathrm{Vol}((1-t)A+tB)^{1/d} \geq (1-t)\mathrm{Vol}(A)^{1/d} + t\,\mathrm{Vol}(B)^{1/d}.\)
\(\mathrm{Per}(A) \geq d\,\omega_d^{1/d}\mathrm{Vol}(A)^{(d-1)/d}.\)
The constant is sharp, with equality for Euclidean balls.
\end{prop}
\begin{proof}
We first give the argument for regular sets, the general Borel case following by approximation from inside and outside. Let $\alpha$ and $\beta$ be the uniform probability measures on $A$ and $B$, and let $T=\nabla\phi$ be the Brenier map with $T_\sharp\alpha=\beta$. The interpolating map $T_t=(1-t)\Id+tT$ sends $\alpha$ to a measure $\alpha_t$ supported in $(1-t)A+tB$. At points where $T$ is differentiable,
\(\det DT_t = \det((1-t)\Id+tDT)\).
Since $DT$ is symmetric positive semidefinite and $\det^{1/d}$ is concave on the positive semidefinite cone,
\(\det(DT_t)^{1/d} \geq (1-t)+t\det(DT)^{1/d}\).
The change-of-variables identity for the uniform measures gives $\det(DT)=\mathrm{Vol}(B)/\mathrm{Vol}(A)$ almost everywhere. If $\rho_t$ denotes the density of $\alpha_t$, a second application of the same formula gives
\(\rho_t(T_t(x))^{-1/d} = \mathrm{Vol}(A)^{1/d}\det(DT_t(x))^{1/d} \geq c_t,\)
where $c_t=(1-t)\mathrm{Vol}(A)^{1/d}+t\mathrm{Vol}(B)^{1/d}$. Thus $\rho_t\leq c_t^{-d}$ almost everywhere. Since $\alpha_t$ has total mass one and is supported in $(1-t)A+tB$, this gives $\mathrm{Vol}((1-t)A+tB)\geq c_t^d$, which is the displayed Brunn--Minkowski inequality.

For isoperimetry, use the homogeneity of Brunn--Minkowski and apply it to $A$ and $\epsilon B_1$. Writing $A+\epsilon B_1$ for the parallel set gives
\(\mathrm{Vol}(A+\epsilon B_1)^{1/d} \geq \mathrm{Vol}(A)^{1/d}+\epsilon\omega_d^{1/d}.\)
Using $\mathrm{Vol}(A+\epsilon B_1)=\mathrm{Vol}(A)+\epsilon\,\mathrm{Per}(A)+o(\epsilon)$ and differentiating at $\epsilon=0$ yields the claimed perimeter bound. Balls attain equality, so the constant is optimal.
\end{proof}
\paragraph{Pr\'ekopa--Leindler.}
\begin{prop}[Pr\'ekopa--Leindler inequality]\label{prop-ot-prekopa-leindler}
Let $u,v,w:\RR^d\to[0,+\infty)$ be integrable, and let $t\in[0,1]$. Assume that
\(w((1-t)x+ty)\geq u(x)^{1-t}v(y)^t \qquad \text{for all }x,y\in\RR^d.\)
Then
\[
\int_{\RR^d} w
\geq
\left(\int_{\RR^d}u\right)^{1-t}
\left(\int_{\RR^d}v\right)^t.
\]
\end{prop}
\begin{proof}
The endpoint cases are immediate, so fix $t\in(0,1)$. Set $U=\int u$ and $V=\int v$. If $U=0$ or $V=0$, the conclusion is trivial; hence assume $U,V>0$. Let $\alpha=(u/U)\d x$ and $\beta=(v/V)\d x$, and let $T$ be the Brenier map from $\alpha$ to $\beta$. Put $T_t=(1-t)\Id+tT$. On a smooth positive regularization, $T_t$ is the gradient of a strictly convex function and is therefore one-to-one almost everywhere. The area formula gives
\(\int w \geq \int w(T_t(x))\det(DT_t(x))\d x\).
For each eigenvalue $s\geq0$ of $DT$, weighted arithmetic--geometric mean gives $(1-t)+ts\geq s^t$. Hence
\(\det((1-t)\Id+tDT)\geq\det(DT)^t\).
\(\int w \geq \int u(x)^{1-t}v(T(x))^t\det(DT(x))^t\d x.\)
Since $T_\sharp\alpha=\beta$, the Monge--Amp\`ere identity gives
\(\det(DT(x))=\frac{u(x)V}{v(T(x))U}\)
almost everywhere. Substitution yields
\[
\int w
\geq
\left(\frac{V}{U}\right)^t\int u(x)\d x
=
U^{1-t}V^t.
\]
Approximation removes the smoothness and positivity assumptions.
\end{proof}
\paragraph{Logarithmic Sobolev inequalities.}
\begin{defn}[Relative Fisher information and logarithmic Sobolev inequality]\label{def-relative-fisher-lsi}
Let \(\beta\in\Pp(\RR^d)\) have a positive density. For a probability measure \(\alpha=h\beta\), define
\begin{equation}\label{eq-relative-fisher-information}
\mathcal I(\alpha|\beta)
\eqdef
4\int \norm{\nabla\sqrt h}^2\,\d\beta
=
\int \norm{\nabla\log h}^2\,\d\alpha,
\end{equation}
whenever $\sqrt h$ has a weak gradient and the integral is finite, and set \(\mathcal I(\alpha|\beta)=+\infty\) otherwise, including when \(\alpha\not\ll\beta\). The measure \(\beta\) satisfies a logarithmic Sobolev inequality with constant \(\lambda>0\) if
\begin{equation}\label{eq-general-log-sobolev}
\KL(\alpha|\beta)
\leq
\frac{1}{2\lambda}\mathcal I(\alpha|\beta)
\qquad
\text{for every probability measure }\alpha.
\end{equation}
\end{defn}
\begin{prop}[Curvature implies logarithmic Sobolev]\label{prop-curvature-log-sobolev}
Let \(\beta=Z^{-1}e^{-V}\d x\) be a smooth probability measure on \(\RR^d\). If \(\nabla^2V\succeq\lambda\Id\) for some \(\lambda>0\), then \(\beta\) satisfies \eqref{eq-general-log-sobolev} with constant \(\lambda\).
\end{prop}
\begin{proof}
Apply the HWI inequality of Proposition~\ref{prop-hwi-inequality}. Under the curvature assumption, for every absolutely continuous \(\alpha\),
\(\KL(\alpha|\beta) \leq \Wass_2(\alpha,\beta)\sqrt{\mathcal I(\alpha|\beta)} - \frac{\lambda}{2}\Wass_2^2(\alpha,\beta).\)
For fixed \(s=\sqrt{\mathcal I(\alpha|\beta)}\), maximize the right-hand side over \(r=\Wass_2(\alpha,\beta)\geq0\). The elementary inequality \(rs-\lambda r^2/2\leq s^2/(2\lambda)\) gives \eqref{eq-general-log-sobolev}. Approximation yields the standard nonsmooth statement. This is the Otto--Villani transport route from the Bakry--Emery curvature bound to logarithmic Sobolev; the original \(\Gamma_2\) calculus gives the same criterion directly.
\end{proof}
\begin{prop}[Logarithmic Sobolev inequality implies entropy decay]\label{prop-lsi-entropy-decay}
Let \(\beta=e^{-V}\d x/Z\) be a smooth probability measure on \(\RR^d\) satisfying the logarithmic Sobolev inequality \eqref{eq-general-log-sobolev} with constant \(\lambda>0\). Let \(\alpha_t\) solve the Wasserstein gradient flow of \(f(\alpha)=\KL(\alpha|\beta)\), namely
\(\partial_t\rho_t = \nabla\cdot(\rho_t\nabla V)+\Delta\rho_t, \qquad \alpha_t=\rho_t\d x.\)
If \(\KL(\alpha_0|\beta)<+\infty\), then
\(\KL(\alpha_t|\beta) \leq e^{-2\lambda t}\KL(\alpha_0|\beta).\)
Moreover, if \(\alpha_0,\beta\in\Pp_2(\RR^d)\), then
\(\Wass_2(\alpha_t,\beta) \leq \sqrt{\frac{2}{\lambda}\KL(\alpha_0|\beta)}\,e^{-\lambda t}.\)
\end{prop}
\begin{proof}
The first variation of \(f\) is \(\log(\rho/\rho_\beta)\), up to an additive constant. Along the smooth Wasserstein gradient flow,
\(\frac{\d}{\d t}\KL(\alpha_t|\beta) = -\int \norm{\nabla\log\frac{\d\alpha_t}{\d\beta}}^2\,\d\alpha_t = -\mathcal I(\alpha_t|\beta).\)
The logarithmic Sobolev inequality bounds the Fisher information below by \(2\lambda\KL(\alpha_t|\beta)\). Hence
\(\frac{\d}{\d t}\KL(\alpha_t|\beta) \leq -2\lambda\KL(\alpha_t|\beta),\)
and Gronwall's lemma gives the entropy estimate. The Wasserstein estimate follows from the length estimate in the Wasserstein--PL theorem: apply Theorem~\ref{thm-wasserstein-pl-convergence} to \(f=\KL(\cdot|\beta)\), whose unique minimizer is \(\beta\), and use the logarithmic Sobolev inequality as the corresponding Wasserstein--PL inequality.
\end{proof}
\paragraph{Gaussian logarithmic Sobolev inequality.}
\begin{prop}[Gaussian logarithmic Sobolev inequality]\label{prop-gaussian-log-sobolev-ot}
Let \(\gamma_d\) be the standard Gaussian measure on \(\RR^d\). If \(\rho>0\) is smooth, has sufficient decay, and \(\int \rho\,\d\gamma_d=1\), then
\(\mathrm{Ent}_{\gamma_d}(\rho) \eqdef \int \rho\log \rho\,\d\gamma_d \leq \frac12\int \frac{\norm{\nabla \rho}^2}{\rho}\,\d\gamma_d.\)
\end{prop}
\begin{proof}
Let \(\alpha=\rho\gamma_d\), and let \(T=\nabla\phi\) be the Brenier map with \(T_\sharp\alpha=\gamma_d\). Write \(T(x)=x+\xi(x)\). The Monge--Amp\`ere equation reads \(\rho(x)e^{-\norm{x}^2/2} = e^{-\norm{T(x)}^2/2}\det(DT(x)).\)
Thus
\(\log \rho(x) = -\dotp{x}{\xi(x)} -\frac12\norm{\xi(x)}^2 + \log\det(\Id+D\xi(x)).\)
Since \(DT\) is positive semidefinite, \(\log\det(\Id+D\xi)\leq\tr(D\xi)=\diverg\xi\). Multiplying by \(\rho\) and integrating with respect to \(\gamma_d\) gives
\(\mathrm{Ent}_{\gamma_d}(\rho) \leq \int \rho(\diverg\xi-\dotp{x}{\xi})\,\d\gamma_d -\frac12\int \rho\norm{\xi}^2\,\d\gamma_d.\)
Gaussian integration by parts gives \(\int \rho(\diverg\xi-\dotp{x}{\xi})\,\d\gamma_d = -\int\dotp{\nabla \rho}{\xi}\,\d\gamma_d.\)
Completing the square pointwise, \(-\dotp{\nabla \rho}{\xi}-\frac12 \rho\norm{\xi}^2 \leq \frac12\frac{\norm{\nabla \rho}^2}{\rho}.\)
This proves the inequality. The general Sobolev-class statement follows by approximation.
\end{proof}
For a translation tilt
\[
\rho_a(x)=\exp\!\left(\dotp{a}{x}-\frac12\norm{a}^2\right),
\qquad a\in\RR^d,
\]
A direct computation gives \(\int \rho_a\log \rho_a\,\d\gamma_d = \frac12\norm{a}^2, \qquad \frac12\int\frac{\norm{\nabla \rho_a}^2}{\rho_a}\,\d\gamma_d = \frac12\norm{a}^2.\)
Thus equality is attained along translations, and the sharp logarithmic Sobolev constant of the standard Gaussian is \(\lambda=1\). Proposition~\ref{prop-lsi-entropy-decay} then gives the classical exponential convergence of the Ornstein--Uhlenbeck flow.

\paragraph{Poincar\'e as a linearized Wasserstein--PL inequality.}
For \(|\epsilon|\) small enough, set
\(\beta_\epsilon=(1+\epsilon\xi)\beta, \qquad \beta_0=\beta.\)
Define the second-order energy and dissipation forms of \(f\) at \(\beta\) in the direction \(\xi\) by
\[
\mathsf H^f_\beta(\xi)
\eqdef
\lim_{\epsilon\to0}
\frac{2(f(\beta_\epsilon)-f(\beta))}{\epsilon^2},
\qquad
\mathsf D^f_\beta(\xi)
\eqdef
\lim_{\epsilon\to0}
\frac{|\partial f|^2(\beta_\epsilon)}{\epsilon^2},
\]
\begin{prop}[Linearizing a Wasserstein--PL inequality]\label{prop-linearized-pl-functional-inequality}
Let \(\beta\) be a minimizer of an energy \(f\), and assume that \(f\) satisfies \(|\partial f|^2(\alpha) \geq 2\kappa\bigl(f(\alpha)-f(\beta)\bigr)\)
in a neighborhood of \(\beta\). Let \(\xi\) be a smooth bounded perturbation with \(\int\xi\,\d\beta=0\), and set \(\beta_\epsilon=(1+\epsilon\xi)\beta\). If \(\mathsf H^f_\beta(\xi)\) and \(\mathsf D^f_\beta(\xi)\) exist, then
\(\kappa\,\mathsf H^f_\beta(\xi) \leq \mathsf D^f_\beta(\xi).\)
\end{prop}
\begin{proof}
Apply the PL inequality to \(\beta_\epsilon\): \(|\partial f|^2(\beta_\epsilon) \geq 2\kappa\bigl(f(\beta_\epsilon)-f(\beta)\bigr).\)
Divide by \(\epsilon^2\) and let \(\epsilon\to0\). The definitions of \(\mathsf H^f_\beta\) and \(\mathsf D^f_\beta\) give the claim.
\end{proof}
For \(f(\alpha)=\KL(\alpha|\beta)\), Proposition~\ref{prop-linearized-pl-functional-inequality} says precisely that logarithmic Sobolev linearizes to Poincar\'e. Indeed, for \(h_\epsilon=1+\epsilon\xi\),
\[
\KL(h_\epsilon\beta|\beta)
=
\frac{\epsilon^2}{2}\int\xi^2\,\d\beta+o(\epsilon^2),
\qquad
\mathcal I(h_\epsilon\beta|\beta)
=
\epsilon^2\int\norm{\nabla\xi}^2\,\d\beta+o(\epsilon^2).
\]
\begin{prop}[Poincar\'e from logarithmic Sobolev]\label{prop-poincare-linearized-wasserstein-pl}
Let \(\beta\) be a smooth probability measure on \(\RR^d\). Assume that \(\beta\) satisfies the logarithmic Sobolev inequality with constant \(\lambda>0\),
\[
\KL(\alpha|\beta)
\leq
\frac{1}{2\lambda}\,\mathcal I(\alpha|\beta),
\qquad
\mathcal I(\alpha|\beta)
\eqdef
\int\norm{\nabla\log\frac{\d\alpha}{\d\beta}}^2\,\d\alpha.
\]
Then, for every smooth \(\xi\) with \(\int\xi\,\d\beta=0\), \(\lambda\int\xi^2\,\d\beta \leq \int\norm{\nabla\xi}^2\,\d\beta.\)
In particular, if \(\d\beta=Z^{-1}e^{-V}\d x\) and \(\nabla^2 V\succeq \lambda\Id\), so that \(\beta\) is \(\lambda\)-strongly log-concave, then the conclusion holds by the Bakry--Emery criterion.
\end{prop}
\begin{proof}
For \(h_\epsilon=1+\epsilon\xi\), with \(|\epsilon|\) small enough, one has \(\beta_\epsilon=h_\epsilon\beta\) and
\(\KL(\beta_\epsilon|\beta) = \int h_\epsilon\log h_\epsilon\,\d\beta = \frac{\epsilon^2}{2}\int\xi^2\,\d\beta+o(\epsilon^2),\)
because \(\int\xi\,\d\beta=0\). The Fisher information expands as
\[
\mathcal I(\beta_\epsilon|\beta)
=
\int\norm{\nabla\log h_\epsilon}^2 h_\epsilon\,\d\beta
=
\epsilon^2\int\norm{\nabla\xi}^2\,\d\beta+o(\epsilon^2).
\]
Inserting these expansions in the logarithmic Sobolev inequality and letting \(\epsilon\to0\) gives the Poincar\'e inequality. The last assertion is the usual Bakry--Emery implication from \(\lambda\)-convexity of \(V\) to the logarithmic Sobolev inequality.
\end{proof}
If \(\d\beta=\rho_\beta\d x=Z^{-1}e^{-V}\d x\), the Poincar\'e inequality is equivalently a spectral gap for the \(\beta\)-weighted Laplacian
\(L_\beta\xi \eqdef \Delta\xi-\dotp{\nabla V}{\nabla\xi} = \rho_\beta^{-1}\nabla\cdot(\rho_\beta\nabla\xi).\)
Indeed \(L_\beta\) is symmetric in \(L^2(\beta)\) and
\(\int\norm{\nabla\xi}^2\,\d\beta = \langle \xi,-L_\beta\xi\rangle_{L^2(\beta)}.\)
For a mean-zero density perturbation \(\dot h\), define
\[
\norm{\dot h}_{-1,\beta}^2
\eqdef
\inf_{v}
\left\{
\int\norm v^2\,\d\beta
\;:\;
\dot h=-\rho_\beta^{-1}\nabla\cdot(\rho_\beta v)
\right\}.
\]
\paragraph{Talagrand's transport-entropy inequality.}
The distance estimate in Theorem~\ref{thm-wasserstein-pl-convergence} already reveals the general principle: when \(\KL(\cdot|\beta)\) satisfies logarithmic Sobolev with constant \(\lambda\), the entropy flow yields the transport-entropy bound \(\Wass_2^2(\alpha,\beta)\leq2\KL(\alpha|\beta)/\lambda\).

\begin{prop}[Gaussian \(T_2\) inequality]\label{prop-gaussian-talagrand-t2}
For every $\alpha\in\Pp_2(\RR^d)$, \(\frac12\Wass_2^2(\alpha,\gamma_d) \leq \KL(\alpha|\gamma_d),\)
with the convention that the right-hand side is \(+\infty\) when \(\alpha\not\ll\gamma_d\).
\end{prop}
\begin{proof}
If \(\KL(\alpha|\gamma_d)=+\infty\), the claim is immediate. Assume first that \(\alpha=\rho\gamma_d\) with \(\rho\) smooth and positive. Let $T=\nabla\phi$ be the Brenier map with $T_\sharp\gamma_d=\alpha$, and write $T(x)=x+\xi(x)$. The change of variables gives
\(\rho(T(x))e^{-\norm{T(x)}^2/2}\det(DT(x)) = e^{-\norm{x}^2/2}.\)
Hence
\(\log \rho(T(x)) = \dotp{x}{\xi(x)} + \frac12\norm{\xi(x)}^2 - \log\det(\Id+D\xi(x)).\)
Using again $\log\det(\Id+D\xi)\leq\diverg\xi$ and integrating with respect to $\gamma_d$, \(\KL(\alpha|\gamma_d) \geq \frac12\int\norm{\xi}^2\,\d\gamma_d + \int(\dotp{x}{\xi}-\diverg\xi)\,\d\gamma_d.\)
The last integral vanishes by Gaussian integration by parts. Since $T$ is optimal, \(\int\norm{\xi}^2\,\d\gamma_d=\Wass_2^2(\gamma_d,\alpha),\)
which proves the claim. Approximation gives the general case. Equality holds for translations \(\alpha=\Gaussian(a,\Id)\), so the constant is sharp.
\end{proof}
\paragraph{The HWI bridge.}
\begin{prop}[Otto--Villani HWI inequality]\label{prop-hwi-inequality}
Let $\beta=\rho_\beta\d x=e^{-V}\d x/Z$ be a smooth probability measure in $\Pp_2(\RR^d)$, and assume that $\nabla^2 V\succeq \lambda \Id$ for some $\lambda\in\RR$. For $\alpha\in\Pp_2(\RR^d)$, define the relative Fisher information by
\(\mathcal I(\alpha|\beta) \eqdef \int \norm{\nabla\log\frac{\rho_\alpha}{\rho_\beta}}^2\,\d\alpha,\)
when \(\alpha=\rho_\alpha\d x\) and the expression is finite, and set \(\mathcal I(\alpha|\beta)=+\infty\) otherwise. Then
\(\KL(\alpha|\beta) \leq \Wass_2(\alpha,\beta)\sqrt{\mathcal I(\alpha|\beta)} -\frac{\lambda}{2}\Wass_2^2(\alpha,\beta).\)
\end{prop}
\begin{proof}
If $\mathcal I(\alpha|\beta)=+\infty$, there is nothing to prove. Assume first that $\rho_\alpha$ is smooth and positive. Let $T$ be the Brenier map from $\alpha$ to $\beta$, set $v=T-\Id$, and let $\alpha_t=((1-t)\Id+tT)_\sharp\alpha$. By Theorem~\ref{thm-mccann-internal-energy} and Proposition~\ref{prop-basic-geodesic-convexity}, $f(\eta)=\KL(\eta|\beta)$ is $\lambda$-geodesically convex. Thus the one-dimensional function $e(t)=f(\alpha_t)$ satisfies
\(e(1)\geq e(0)+e'(0)+\frac{\lambda}{2}\Wass_2^2(\alpha,\beta).\)
Since $e(1)=f(\beta)=0$, it remains to compute the initial derivative. The continuity equation for the geodesic gives $\partial_t\alpha_t+\nabla\cdot(\alpha_t v_t)=0$ and $v_0=v$. Hence
\(e'(0) = \int \dotp{\nabla\log\frac{\rho_\alpha}{\rho_\beta}}{v}\,\d\alpha.\)
Therefore
\(\KL(\alpha|\beta) \leq -\int \dotp{\nabla\log\frac{\rho_\alpha}{\rho_\beta}}{T-\Id}\,\d\alpha -\frac{\lambda}{2}\Wass_2^2(\alpha,\beta).\)
Cauchy--Schwarz and $\int\norm{T-\Id}^2\d\alpha=\Wass_2^2(\alpha,\beta)$ give the announced bound. The general case follows by standard approximation and lower semicontinuity of entropy, Fisher information, and $\Wass_2$.
\end{proof}
When \(\lambda>0\), the HWI inequality is a bridge from curvature to PL\@. It is not itself a PL inequality because it still contains the distance term \(\Wass_2(\alpha,\beta)\). Optimizing the right-hand side, or simply using Young's inequality \(rs-\lambda r^2/2\leq s^2/(2\lambda)\), gives the logarithmic Sobolev estimate
\(\KL(\alpha|\beta) \leq \frac{1}{2\lambda}\mathcal I(\alpha|\beta).\)
\subsection{Training Two-Layer MLPs as Wasserstein Flows}
\label{sec-wasserstein-flows-mlp}
\paragraph{Wasserstein training of two-layer MLPs.}
A neuron is a particle
\(x=(u,v)\in\RR^d\times\RR^{d'},\)
where $u$ is the inner weight and $v$ is the outer vector weight. For a scalar nonlinearity $\sigma$, define the vector-valued feature
\(\psi(x,z)=v\,\sigma(\dotp{u}{z})\in\RR^{d'}.\)
The width-$n$ network and its mean-field version are \(G_X(z)=\frac1n\sum_{i=1}^n\psi(x_i,z), \qquad G_\alpha(z)=\int\psi(x,z)\d\alpha(x), \qquad \alpha=\frac1n\sum_i\delta_{x_i}.\)
Let $\zeta$ be a probability distribution on data-label pairs $(z,y)\in\RR^d\times\RR^{d'}$. The population risk is \(f(\alpha)=\int \ell(G_\alpha(z),y)\d\zeta(z,y),\)
The corresponding particle flow is
\[
\dot x_i=-n\nabla_{x_i}F(X),
\qquad
F(X)=f\!\left(\frac1n\sum_i\delta_{x_i}\right),
\]
which is the gradient flow of $F(X)=f(\alpha_X)$ for this Wasserstein particle metric, equivalently Euclidean gradient descent with the time scale multiplied by $n$. It gives a particle discretization of~\eqref{eq:wassflow-pde}.

Assume that $\ell$ is differentiable in its first variable. The first variation is
\begin{equation}\label{eq-mlp-first-variation-general}
\delta f(\alpha)(x)
=
\int
\dotp{\nabla_1\ell(G_\alpha(z),y)}{\psi(x,z)}
\d\zeta(z,y),
\end{equation}
and the Wasserstein gradient in parameter space is \(\Wgrad f(\alpha)(x) = \nabla_x\delta f(\alpha)(x) = \int [D_x\psi(x,z)]^\top\nabla_1\ell(G_\alpha(z),y) \d\zeta(z,y).\)
For the squared Euclidean loss $\ell(s,y)=\frac12\norm{s-y}^2$, the energy is the sum of a quadratic interaction and a linear potential:
\begin{equation}\label{eq-mlp-square-loss-quadratic-linear}
f(\alpha)
=
\frac12\iint k(x,x')\d\alpha(x)\d\alpha(x')
+
\int g(x)\d\alpha(x)
+\frac12\int\norm{y}^2\d\zeta(z,y),
\end{equation}
with
\begin{equation}\label{eq-mlp-kernel-linear-potential}
k(x,x')=\int\dotp{\psi(x,z)}{\psi(x',z)}\d\zeta(z,y),
\qquad
g(x)=-\int\dotp{y}{\psi(x,z)}\d\zeta(z,y).
\end{equation}
Thus
\(\delta f(\alpha)(x)=\int k(x,x')\d\alpha(x')+g(x), \qquad \Wgrad f(\alpha)(x)=\int\nabla_x k(x,x')\d\alpha(x')+\nabla_xg(x).\)
\paragraph{Classical convexity and stationarity.}
Consider an energy of the form
\(f(\alpha) = \frac12\iint k(x,x')\d\alpha(x)\d\alpha(x') + \int V(x)\d\alpha(x) +C,\)
Assume that the quadratic part is convex in the classical sense, namely convex for the affine structure of measures: \(Q((1-s)\alpha+s\beta)\leq (1-s)Q(\alpha)+sQ(\beta), \qquad Q(\alpha)=\frac12\iint k\d\alpha\d\alpha.\)
\begin{prop}[Affine convexity and stationary positive densities]\label{prop-classical-convex-stationary}
Let $\Omega\subset\RR^d$ be connected and open, let $f=Q+\int V\,\d\alpha+C$ be as above on $\Pp_2(\Omega)$, and assume that $Q$ is classically convex. Suppose that $\alpha_t\in AC^2_{\mathrm{loc}}([0,+\infty);\Pp_2(\Omega))$ is an energy-dissipation solution of the Wasserstein gradient flow of $f$, with $f(\alpha_0)<+\infty$, and converges in $\Wass_2$ to $\alpha_\infty=\rho_\infty\d x$, where $\rho_\infty>0$ almost everywhere on $\Omega$, and that $\inf_{t\geq0}f(\alpha_t)>-\infty$. Assume that $\delta f(\alpha_\infty)\in C^1(\Omega)$, that the first-variation formula and the affine subgradient inequality hold along every finite-energy segment from $\alpha_\infty$, and that the squared Wasserstein slope is lower semicontinuous along the convergence, in the sense that
\[
\int_\Omega\norm{\nabla\delta f(\alpha_\infty)}^2\d\alpha_\infty
\leq
\liminf_{n\to\infty}
\int_\Omega\norm{\nabla\delta f(\alpha_{t_n})}^2\d\alpha_{t_n}
\]
whenever $t_n\to\infty$. Then $\alpha_\infty$ is a global minimizer of $f$.
\end{prop}
\begin{proof}
The energy-dissipation inequality and the lower bound on $f$ along the convergent trajectory imply
\(\int_0^\infty \int_\Omega\norm{\nabla\delta f(\alpha_t)}^2\d\alpha_t\d t <+\infty.\)
Hence there exists $t_n\to\infty$ for which the inner integral tends to zero. The assumed lower semicontinuity gives \(\int_\Omega \norm{\nabla \delta f(\alpha_\infty)}^2\d\alpha_\infty=0.\)
Because $\rho_\infty>0$ almost everywhere and $\nabla\delta f(\alpha_\infty)$ is continuous, this gradient vanishes throughout $\Omega$. Connectedness then implies that $\delta f(\alpha_\infty)$ is constant.
\(\int_\Omega\delta f(\alpha_\infty)\d(\beta-\alpha_\infty)=0.\)
Classical convexity of $f$ in the affine variable $\alpha$ gives the subgradient inequality
\(f(\beta)\geq f(\alpha_\infty) + \int \delta f(\alpha_\infty)\d(\beta-\alpha_\infty) = f(\alpha_\infty).\)
The same conclusion is immediate when $f(\beta)=+\infty$, so no competitor in $\Pp_2(\Omega)$ has smaller energy. Strict positivity is essential to this argument: without it, stationarity controls the first variation only on the support explored by the limiting measure. For square-loss two-layer mean-field models, \eqref{eq-mlp-square-loss-quadratic-linear} is exactly of this quadratic-plus-linear form, and positive semidefiniteness of the induced kernel $k$ is the classical convexity assumption.
\end{proof}
\begin{prop}[Formal global optimality for two-homogeneous mean-field flows]\label{prop-formal-chizat-bach}
Assume that the feature is positively two-homogeneous in the neuron variable, \(\psi(\lambda x,z)=\lambda^2\psi(x,z) \qquad(\lambda>0),\)
and that $f(\alpha)=J(G_\alpha)$ with $J$ convex and differentiable as a functional of the predictor. Let $\alpha$ be a smooth stationary point of the Wasserstein flow, so that $\nabla_x\delta f(\alpha)(x)=0$ on $\supp(\alpha)$. Assume also full directional support: for every unit direction $\omega$, the support of $\alpha$ intersects the ray $\{\lambda\omega:\lambda>0\}$. Then $\alpha$ is a global minimizer of $f$ over the mean-field model class.
\end{prop}
\begin{proof}
Choose the first-variation representative
\(h_\alpha(x)=\delta f(\alpha)(x) = DJ(G_\alpha)\bigl[\psi(x,\cdot)\bigr].\)
When the predictor space is $L^2(\zeta)$ and $DJ(G_\alpha)$ has a Riesz representative, this derivative-dual pairing can equivalently be written as the $L^2(\zeta)$ inner product with that representative. The abstract notation above is all that is needed here.
Additive constants in first variations do not affect the Wasserstein gradient or first-order inequalities between probability measures; this normalization is useful because it inherits the homogeneity of $\psi$. By two-homogeneity of $\psi$, one has
\(h_\alpha(\lambda x)=\lambda^2h_\alpha(x), \qquad \lambda>0.\)
Taking \(x=0\) and any \(\lambda\neq1\) in this identity gives \(h_\alpha(0)=0\). We first record that stationarity forces \(h_\alpha\) to vanish on \(\supp(\alpha)\). Indeed, if \(x=r\omega\in\supp(\alpha)\) with \(r>0\) and \(\norm{\omega}=1\), then
\(0=\dotp{\nabla h_\alpha(r\omega)}{\omega} =\frac{\d}{\d s}h_\alpha(s\omega)\bigg|_{s=r} =2r h_\alpha(\omega),\)
so \(h_\alpha(\omega)=0\), and hence \(h_\alpha(x)=r^2h_\alpha(\omega)=0\). The case \(x=0\) has already been covered, so \(h_\alpha=0\) on \(\supp(\alpha)\) and \(\int h_\alpha\,\d\alpha=0\).

Suppose that a competitor \(\beta\) satisfies \(f(\beta)<f(\alpha)\). Since
\(G_\beta-G_\alpha = \int \psi(x,\cdot)\d(\beta-\alpha)(x),\)
convexity of \(J\) gives \(f(\beta)-f(\alpha) \geq \int h_\alpha(x)\d(\beta-\alpha)(x).\)
The left-hand side is negative, while \(\int h_\alpha\,\d\alpha=0\), so \(\int h_\alpha\,\d\beta<0\). Since \(h_\alpha(0)=0\), this strict negativity must occur away from the origin: there exists a nonzero point \(x=r\omega\) with \(r>0\), \(\norm{\omega}=1\), and \(h_\alpha(x)<0\). Equivalently \(h_\alpha(\omega)=h_\alpha(x)/r^2<0\). By full directional support, there is \(r_\omega>0\) such that \(r_\omega\omega\in\supp(\alpha)\). At this support point,
\(\dotp{\nabla h_\alpha(r_\omega\omega)}{\omega} = 2r_\omega h_\alpha(\omega) <0,\)
which contradicts the stationarity condition \(\nabla h_\alpha=0\) on \(\supp(\alpha)\). Thus no competitor has smaller risk.

The rigorous Chizat--Bach theorem replaces the full directional support assumption by propagation and overparameterization hypotheses ensuring that any negative first-variation direction is represented by the evolving support. The same radial-derivative contradiction, applied after this support-propagation step, then rules out non-optimal stationary limits.
\end{proof}
\subsection{Generalized Dynamic Wasserstein Flows}
\label{sec-generalized-dynamic-wasserstein-flows}
\paragraph{Generalized Wasserstein flows.}
\begin{defn}[Generalized action minimizing movement]\label{def-generalized-action-gradient-flow}
Let $d$ be a metric or extended metric generated by the action $\mathbb A$ through the path value in Definition~\ref{def-generalized-dynamic-action-distance}; for a quadratic action, $d^2=\mathsf E_{\mathbb A}$. Given \(\alpha^{(0)}\in\Pp(\RR^d)\) and \(\tau>0\), define
\begin{equation}\label{eq-generalized-jko-step}
\alpha^{(\ell+1)}
\in
\uargmin{\alpha\in\Pp(\RR^d)}
\left\{\frac{1}{2\tau}d^2(\alpha^{(\ell)},\alpha)+f(\alpha)\right\}.
\end{equation}
For each step size, let \(\bar\alpha_\tau(t)=\alpha^{(\ell)}\) on \(((\ell-1)\tau,\ell\tau]\). Any limit of \(\bar\alpha_{\tau_k}\) along a sequence \(\tau_k\downarrow0\), in a topology strong enough to identify the initial condition, is called a generalized minimizing movement of \(f\) for \(\mathbb A\). When such a limit is identified with a curve of maximal slope for the action metric, it is called the corresponding generalized action gradient flow.
\end{defn}
Compactness, coercivity and lower semicontinuity adapted to \(d\) can yield subsequential convergence of the interpolants in \eqref{eq-generalized-jko-step} to the minimizing movement just defined. A nonmetric discrepancy may still be inserted in the same variational update, but this metric theory does not then apply automatically. In the local mass-preserving metrics considered below, a smooth limiting curve is again represented by a continuity equation of the type introduced in \eqref{eq:eulerian-advection} and used for Wasserstein flows in \eqref{eq:wassflow-pde},
\begin{equation}\label{eq-generalized-flow-continuity-motivation}
\partial_t\alpha_t+\diverg(\alpha_t v_{\alpha_t})=0,
\end{equation}
The metric ingredient needed below is the squared-speed action $\mathbb A(\alpha,w)$ introduced in Section~\ref{sec-generalized-dynamic-wasserstein-distances}. Its path value is defined in~\eqref{eq-generalized-action-length-distance}; when the metric hypotheses of Proposition~\ref{prop-homogeneous-dynamic-action-distance} hold, the appropriate root generates the distance used by JKO. Others are built from a homogeneous local density but then squared at the metric level: for $\Wass_p$, the local densities are \(A_p(a,w)=a\norm w^p\) and \(J_p(a,m)=\norm m^p/a^{p-1}\), while the action paired with the JKO penalty is the squared Finsler speed
\[
\mathbb A_p(\alpha,w)=\left(\int\norm w^p\d\alpha\right)^{2/p}.
\]
In the \(\Wass_2\) case, all objects reduce to the familiar Benamou--Brenier ones: \(A_2(a,w)=a\norm w^2, \qquad J_2(a,m)=\frac{\norm{m}^2}{a}, \qquad \mathsf E_{\mathbb A}=\Wass_2^2.\)
\begin{defn}[Penalized minimization oracle]\label{def-generalized-action-steepest-direction}
Let \(\mathbb A(\alpha,\cdot)\) be the local squared-speed action at \(\alpha\). For a cotangent field \(u\) for which the pairing below is finite, usually \(u=\Wgrad f(\alpha)\), the penalized minimization oracle associated with this action is
\begin{equation}\label{eq-local-action-steepest-descent}
\operatorname{PMO}_{\mathbb A,\alpha}(u)
\in
\uargmin{w}
\left\{
\int \dotp{u(x)}{w(x)}\d\alpha(x)
+
\frac12\mathbb A(\alpha,w)
\right\},
\end{equation}
where the minimization is over admissible finite-action velocity representatives for the chosen geometry. In Hilbertian examples \(u\in L^2(\alpha;\RR^d)\); for \(\Wass_p\), the natural dual space is \(L^q(\alpha;\RR^d)\), with \(q=p/(p-1)\).
\end{defn}
For an energy $f$, the formal small-step expansion of~\eqref{eq-generalized-jko-step}, when the distance is generated by the squared-speed action \(\mathbb A\) and the JKO displacement admits an admissible finite-action tangent representative, selects the PMO direction
\begin{equation}\label{eq-generalized-gradient-from-action}
v_{\alpha}
=
\operatorname{PMO}_{\mathbb A,\alpha}(\Wgrad f(\alpha)).
\end{equation}
Equivalently, the formal evolution is
\begin{equation}\label{eq-local-action-gradient-flow-pde}
\partial_t\alpha_t
+
\diverg\!\left(
\alpha_t\,\operatorname{PMO}_{\mathbb A,\alpha_t}(\Wgrad f(\alpha_t))
\right)=0.
\end{equation}
Along smooth solutions of~\eqref{eq-local-action-gradient-flow-pde}, whenever the minimizer is attained and \(w\mapsto\mathbb A(\alpha,w)\) is \(2\)-homogeneous, one obtains the dissipation identity
\[
\frac{\d}{\d t}f(\alpha_t)
=
\left\langle \Wgrad f(\alpha_t),v_t\right\rangle_{\alpha_t}
=
-\mathbb A(\alpha_t,v_t),
\qquad
v_t=\operatorname{PMO}_{\mathbb A,\alpha_t}(\Wgrad f(\alpha_t)).
\]
For \(\Wass_2\), \(\mathbb A(\alpha,w)=\int\norm w^2\d\alpha\) and \(\operatorname{PMO}_{\mathbb A,\alpha}(u)=-u\). Hence~\eqref{eq-local-action-gradient-flow-pde} reduces to the standard Wasserstein gradient-flow equation
\[
\partial_t\alpha_t
=
\diverg\!\left(\alpha_t\Wgrad f(\alpha_t)\right),
\]
In the restricted Riemannian case of~\eqref{eq-general-quadratic-tangent-action}, where
\[
\mathbb A(\alpha,w)=\left\langle Q_\alpha w,w\right\rangle_{L^2(\alpha)},
\]
the PMO becomes the linear preconditioning rule
\begin{equation}\label{eq-generalized-gradient-q-preconditioner}
\operatorname{PMO}_{\mathbb A,\alpha}(u)
=
-Q_\alpha^{-1}u,
\qquad
v_\alpha=-Q_\alpha^{-1}\Wgrad f(\alpha).
\end{equation}
\[
\mathbb A_p(\alpha,w)=\left(\int\norm{w}^p\d\alpha\right)^{2/p}.
\]
Writing \(q=p/(p-1)\) and \(U=\norm{u}_{L^q(\alpha)}\), its PMO is the duality-map formula
\begin{equation}\label{eq-wp-pmo}
\operatorname{PMO}_{\mathbb A_p,\alpha}(u)(x)
=
-\,U^{2-q}\norm{u(x)}^{q-2}u(x),
\end{equation}
The examples below are organized by the same dictionary. Concave-mobility flows use $\mathbb A_{\theta,\lambda}(\alpha,v)$ from Section~\ref{rem-generalized-bb}; spectral flows use $\mathbb A_\gamma(\alpha,v)$ from~\eqref{eq-spectral-tangent-action}; kernelized Benamou--Brenier flows use $\mathbb A_k(\alpha,v)=\norm{v}_{\RKHS_k^d}^2$ with a restricted RKHS tangent space.

\paragraph{General mobility flows.}
Let \(\lambda\) be Lebesgue measure, write \(\alpha=\rho\lambda\), and let \(u=\nabla(\delta f/\delta\rho)\). For the velocity action \(\mathbb A_{\theta,\lambda}\), the PMO introduced in Definition~\ref{def-generalized-action-steepest-direction} is the pointwise minimizer of
\(\int \dotp{u}{w}\,\rho\,\d x + \frac12\int \frac{\rho}{\theta(\rho)}\norm{w}^2\rho\,\d x,\)
so, wherever \(\rho>0\) and \(\theta(\rho)>0\),
\begin{equation}\label{eq-concave-mobility-pmo}
\operatorname{PMO}_{\mathbb A_{\theta,\lambda},\rho\lambda}(u)
=
-\frac{\theta(\rho)}{\rho}u.
\end{equation}
For a smooth density \(\rho_t\) and an energy \(f\), the formal gradient flow associated with the pointwise momentum action \(J_\theta(a,m)=|m|^2/\theta(a)\) is
\begin{equation}\label{eq-general-mobility-gradient-flow}
\partial_t\rho_t
=
\nabla\!\cdot\!\left(\theta(\rho_t)\nabla\frac{\delta f}{\delta\rho}(\rho_t)\right).
\end{equation}
If
\(f(\rho)=\int U(\rho(x))\,\d x+\int V(x)\rho(x)\,\d x,\)
then \eqref{eq-general-mobility-gradient-flow} becomes
\[
\partial_t\rho
=
\nabla\!\cdot\!\left(\theta(\rho)\big(U''(\rho)\nabla\rho+\nabla V\big)\right).
\]
\paragraph{Dynamic spectral Wasserstein flows.}\label{sec-normalized-spectral-wasserstein-dynamics}
Specializing Definition~\ref{def-generalized-action-steepest-direction} to the spectral action \(\mathbb A_\gamma\), the descent direction associated with a field \(g\in L^2(\alpha;\RR^d)\) is simply \(\operatorname{PMO}_{\mathbb A_\gamma,\alpha}(g)\). The field \(g\) should be read as the Wasserstein gradient \(g=\Wgrad f(\alpha)=\nabla\delta f(\alpha)\).

\begin{prop}[Polar formula for the spectral direction]\label{prop-normalized-spectral-polar}
Let
\(S_\alpha(g)\eqdef\int g(x)g(x)^\top\d\alpha(x).\)
Assume that the inverse-trace problem admits a positive definite minimizer
\[
A_\alpha^\star
\in
\uargmin{A\in\mathcal B_\gamma\cap\mathbb S_{++}^d}
\tr\!\left(A^{-1}S_\alpha(g)\right).
\]
\eqllead{Then}{eq-normalized-spectral-direction-explicit}{
\operatorname{PMO}_{\mathbb A_\gamma,\alpha}(g)(x)
=
-(A_\alpha^\star)^{-1}g(x)
}
is the spectral PMO direction. If the minimizer lies on the boundary of $\mathcal B_\gamma$, the same formula is understood by a limiting argument along positive definite minimizers, or equivalently by replacing the inverse by the corresponding inverse on the range of $S_\alpha(g)$.
\end{prop}
\begin{proof}
Using the polar formula $\gamma(M)=\sup_{A\in\mathcal B_\gamma}\tr(AM)$, the PMO objective~\eqref{eq-local-action-steepest-descent} with \(\mathbb A=\mathbb A_\gamma\) can be written as
\[
\sup_{A\in\mathcal B_\gamma}
\left\{
\int \dotp{g}{v}\d\alpha
+
\frac12\int v^\top A v\d\alpha
\right\}.
\]
For a fixed $A\succ0$, the quadratic lower bound
\(\int \dotp{g}{v}\d\alpha+\frac12\int v^\top A v\d\alpha\)
is minimized by $v_A=-A^{-1}g$, with value $-\frac12\tr(A^{-1}S_\alpha(g))$. Let $v_\star=-(A_\alpha^\star)^{-1}g$ and
\(M_\star=\int v_\star(x)v_\star(x)^\top\d\alpha(x) = (A_\alpha^\star)^{-1}S_\alpha(g)(A_\alpha^\star)^{-1}.\)
Since $A_\alpha^\star$ minimizes $A\mapsto\tr(A^{-1}S_\alpha(g))$ over the convex polar set, its first-order optimality condition gives
\(\tr(AM_\star)\leq \tr(A_\alpha^\star M_\star) \qquad \text{for all } A\in\mathcal B_\gamma.\)
Thus $A_\alpha^\star$ is active in the polar representation of $\gamma(M_\star)$, and
\(\gamma(M_\star)=\tr(A_\alpha^\star M_\star) = \tr((A_\alpha^\star)^{-1}S_\alpha(g)).\)
For any $v$, the polar formula gives \(\int \dotp{g}{v}\d\alpha+\frac12\mathbb A_\gamma(\alpha,v) \geq \int \dotp{g}{v}\d\alpha+\frac12\int v^\top A_\alpha^\star v\d\alpha,\)
and the right-hand side is minimized at $v_\star$. The activity identity above shows that equality holds at $v_\star$, hence $v_\star$ is the spectral PMO direction. The boundary case follows by approximation with positive definite matrices in the polar set.
\end{proof}
For the trace gauge, $\mathcal B_\gamma=\enscond{A}{0\preceq A\preceq\Id}$ and the minimizer is $A_\alpha^\star=\Id$. For the operator gauge $\gamma(M)=\lambda_{\max}(M)$, one has $\mathcal B_\gamma=\enscond{A\succeq0}{\tr(A)\leq1}$; if $S_\alpha(g)\succ0$, then
\[
A_\alpha^\star
=
\frac{S_\alpha(g)^{1/2}}{\tr(S_\alpha(g)^{1/2})},
\qquad
\operatorname{PMO}_{\mathbb A_\gamma,\alpha}(g)(x)
=
-\tr(S_\alpha(g)^{1/2})\,S_\alpha(g)^{-1/2}g(x).
\]
Proposition~\ref{prop-static-dynamic-spectral-wasserstein} identifies the static spectral distance $\Wass_\gamma$ with the length distance generated by $\mathbb A_\gamma$. The infinitesimal descent direction is therefore defined directly by the general PMO problem~\eqref{eq-local-action-steepest-descent} specialized to this action. This action-based statement avoids an important overinterpretation: static/dynamic equality alone does not imply
\(\Wass_\gamma^2(\alpha,(\Id+\tau v)_\sharp\alpha)=\tau^2\mathbb A_\gamma(\alpha,v)+o(\tau^2)\)
\begin{prop}[Formal normalized spectral gradient flow]\label{prop-normalized-spectral-gradient-flow}
Assume that $f$ has a smooth first variation and that the spectral PMO is attained along a smooth curve. The formal gradient flow generated by the dynamic spectral action $\mathbb A_\gamma$ solves
\begin{equation}\label{eq-normalized-spectral-flow}
\partial_t\alpha_t
+
\diverg\!\left(
\alpha_t\,\operatorname{PMO}_{\mathbb A_\gamma,\alpha_t}(\Wgrad f(\alpha_t))
\right)=0.
\end{equation}
Equivalently, if $A_t^\star$ solves the inverse-trace problem for
\(S_t=\int \Wgrad f(\alpha_t)(x)\Wgrad f(\alpha_t)(x)^\top\d\alpha_t(x),\)
then the velocity is \(v_t(x)=-(A_t^\star)^{-1}\Wgrad f(\alpha_t)(x).\)
Moreover, along smooth solutions, \(\frac{\d}{\d t}f(\alpha_t) = -\mathbb A_\gamma(\alpha_t,v_t).\)
If the minimizing movements~\eqref{eq-generalized-jko-step} for $d=\Wass_\gamma$ converge to a smooth curve with this minimal-action metric derivative, that limit satisfies the same equation.
\end{prop}
\begin{proof}
The action-based steepest-descent rule~\eqref{eq-generalized-gradient-from-action} gives
\(
v_t=\operatorname{PMO}_{\mathbb A_\gamma,\alpha_t}(\Wgrad f(\alpha_t))
\).
Proposition~\ref{prop-normalized-spectral-polar} gives the explicit inverse-trace formula.
Finally, the variation formula gives $\frac{\d}{\d t}f(\alpha_t)=\int\dotp{\Wgrad f(\alpha_t)}{v_t}\d\alpha_t$. Since $v_t$ minimizes a $2$-homogeneous penalized objective, the one-dimensional rescaling $s\mapsto sv_t$ has derivative zero at $s=1$, hence
\(\int\dotp{\Wgrad f(\alpha_t)}{v_t}\d\alpha_t + \mathbb A_\gamma(\alpha_t,v_t) =0,\)
which proves the dissipation identity.
\end{proof}
\subsection{Nonlocal Wasserstein Flows}
\label{sec-nonlocal-wasserstein-flows}
\label{sec-nonlocal-wasserstein-pdes}
\paragraph{Nonlocal Wasserstein flows and fractional PDEs.}
\begin{prop}[Entropy gradient flow for reversible jump kernels]\label{prop-nonlocal-entropy-gradient-flow}
Let $\mathcal X=\RR^d$, let $\mathfrak m=\mathcal L^d$, and let $K(x,\d y)=\nu(\d y-x)$ be translation invariant. Assume the metric hypotheses of Proposition~\ref{prop-nonlocal-distance-properties} and the heat-kernel and moment conditions of Assumptions~5.5--5.6 in. Then the Markov semigroup generated by the closure of
\(L\psi(x) = \int\bigl(\psi(y)-\psi(x)\bigr)K(x,\d y)\)
is the gradient flow, for the distance \(\mathcal W_K\) defined in \eqref{eq-nonlocal-wasserstein-distance}, of the relative entropy
\(\mathrm{Ent}_{\mathfrak m}(\alpha) = \int \rho\log\rho\,\d\mathfrak m, \qquad \alpha=\rho\mathfrak m.\)
Equivalently, in weak form, the entropy flow is
\begin{equation}\label{eq-nonlocal-entropy-flow}
\partial_t\rho_t=L\rho_t.
\end{equation}
For singular kernels, \(L\) is understood through the closure of its Dirichlet form; for symmetric stable kernels, this agrees with the principal-value formula displayed below.
\end{prop}
\begin{proof}[Proof idea]
For positive arguments, the key identity is the logarithmic-mean relation
\(\theta(a,b)(\log a-\log b)=a-b\); boundary states follow by positive approximation.
The first variation of \(\mathrm{Ent}_{\mathfrak m}\) is \(\log\rho+1\). With the sign convention of \eqref{eq-nonlocal-continuity-weak}, the steepest descent velocity is
\(v(x,y) = -\bar\nabla(\log\rho+1)(x,y) = \log\rho(x)-\log\rho(y).\)
Hence \(\theta(\rho(x),\rho(y))v(x,y)=\rho(x)-\rho(y)\). Substituting this in \eqref{eq-nonlocal-continuity-weak} and using the symmetry of \(\mathsf J\) gives
\(\frac{\d}{\d t}\int \varphi\rho\,\d\mathfrak m = \int \varphi(x)\int(\rho(y)-\rho(x))K(x,\d y)\,\mathfrak m(\d x),\)
which is the weak form of \(\partial_t\rho=L\rho\). This calculation identifies the formal descent equation. Under the stated heat-kernel and moment assumptions, Theorem~5.11 of proves the corresponding evolution variational inequality, and hence the rigorous metric gradient-flow statement.
\end{proof}
On \(\mathcal X=\RR^d\), let
\(K_s(x,\d y) = c_{d,s}\frac{\d y}{|x-y|^{d+s}}, \qquad 0<s<2.\)
The associated generator is, up to the normalizing constant, \(L_s\psi(x) = \mathrm{p.v.}\int_{\RR^d} \bigl(\psi(y)-\psi(x)\bigr) \frac{c_{d,s}}{|x-y|^{d+s}}\,\d y = -(-\Delta)^{s/2}\psi(x),\)
Hence the \(\mathcal W_{K_s}\)-gradient flow of the entropy is
\begin{equation}\label{eq-fractional-heat-nonlocal-wasserstein}
\partial_t\rho_t
=
-(-\Delta)^{s/2}\rho_t.
\end{equation}
More generally, when a target-dependent jump kernel \(K_\beta\) satisfies detailed balance with a desired invariant measure \(\beta=\rho_\beta\mathfrak m\), the same construction can be applied to \(r_t=\d\alpha_t/\d\beta\). The gradient flow of \(\KL(\alpha|\beta)\) is then \(\partial_t r_t=L_\beta r_t\).

\paragraph{Discrete Wasserstein flows on Markov chains.}
For masses \(a_i=\pi_i\rho_i\), or equivalently relative densities \(\rho_i=a_i/\pi_i\) with respect to the invariant law \(\pi\), the entropy relative to \(\pi\) is
\(\operatorname{Ent}_\pi(\rho)\eqdef\sum_i\pi_i\rho_i\log\rho_i.\)
Applying the general minimizing movement~\eqref{eq-generalized-jko-step} with \(d=\mathcal W_K\) and \(f=\operatorname{Ent}_\pi\), the small-\(\tau\) first-order optimality condition gives
\(\frac{\rho^{(\ell+1)}-\rho^{(\ell)}}{\tau} =-\mathcal K_{\rho^{(\ell+1)}}\log\rho^{(\ell+1)}+o(1) =K\rho^{(\ell+1)}+o(1).\)
Here \((K\rho)_i=\sum_jK_{i,j}(\rho_j-\rho_i)\). Under smoothness and strict positivity, $\rho^{(\ell+1)}=\rho^{(\ell)}+O(\tau)$, so the last expression also equals $K\rho^{(\ell)}+O(\tau)$.

\begin{prop}[Entropy gradient flow of a reversible Markov chain]
\label{prop-discrete-markov-entropy-gradient}
Let \(K\) be reversible with invariant law \(\pi\). The gradient flow of \(\operatorname{Ent}_\pi\) for the metric \(\mathcal W_K\) is the forward equation of the Markov chain,
\begin{equation}
\label{eq-discrete-markov-gradient-flow}
\dot\rho_i(t)=\sum_jK_{i,j}\bigl(\rho_j(t)-\rho_i(t)\bigr).
\end{equation}
Equivalently, for the masses \(a_i(t)=\pi_i\rho_i(t)\), this is \(\dot a_i(t)=\sum_j(a_j(t)K_{j,i}-a_i(t)K_{i,j})\).
\end{prop}
\begin{proof}
The first variation of the entropy, with respect to the weighted pairing \(\sum_i\pi_i\xi_i\varphi_i\), is \(\log\rho_i+1\). Constants do not contribute to \(\mathcal K_\rho\), hence the metric gradient-flow equation is
\(\dot\rho=-\mathcal K_\rho\log\rho\).
For positive densities, using
\(\theta(a,b)(\log a-\log b)=a-b\),
one obtains, componentwise, \(\dot\rho_i =-\sum_jK_{i,j}\theta(\rho_i,\rho_j)(\log\rho_i-\log\rho_j) = \sum_jK_{i,j}(\rho_j-\rho_i),\)
which is the density form of the Kolmogorov forward equation. Boundary initial states follow by positive approximation, or by the instantaneous positivity of an irreducible finite chain. Multiplying by \(\pi_i\) and using detailed balance gives the equation for the probability masses.
\end{proof}
\subsection{Dynamic Unbalanced OT and WFR Flows}
\label{sec-dynamic-unbalanced-wfr-flows}
\label{sec-wfr-gradient-flows}
The generalized dynamic Wasserstein flows of Section~\ref{sec-generalized-dynamic-wasserstein-flows} are primarily mass-preserving: their tangent vectors are generated by continuity equations and are therefore velocity fields advecting the measure. Unbalanced transport changes the state space from probabilities to positive finite measures, where descent may both move and create or remove mass.

\begin{defn}[WFR minimizing movement]\label{def-wfr-jko-gradient-flow}
Let $f:\Mm_+(\RR^d)\to\RR\cup\{+\infty\}$ be an energy on positive finite measures and let $\WFR_\kappa$ be the Wasserstein--Fisher--Rao distance with growth scale $\kappa>0$, defined by the dynamic balance-equation formulation~\eqref{eq-dynamic-unbalanced-ot} and its static equivalence in Proposition~\ref{prop-static-dynamic-unbalanced}. The unbalanced JKO step is
\begin{equation}\label{eq-wfr-jko}
\alpha^{(\ell+1)}
\in
\uargmin{\alpha\in\Mm_+(\RR^d)}
\frac{1}{2\tau}\WFR_\kappa^2(\alpha^{(\ell)},\alpha)+f(\alpha).
\end{equation}
The continuous-time limit, when it exists, is called the $\WFR_\kappa$ gradient flow of $f$.
\end{defn}
Formally, a smooth curve $\alpha_t=\rho_t\d x$ has tangent vectors of the form
\(\partial_t\rho_t+\diverg(\rho_t v_t)=g_t\rho_t,\)
\begin{proposition}[Formal WFR gradient-flow PDE]\label{prop-wfr-gradient-flow-pde}
Assume that $f$ has a smooth first variation \(\phi_t=\delta f(\alpha_t)\) and that $\alpha_t=\rho_t\d x$ is smooth and positive. The formal $\WFR_\kappa$ gradient flow of $f$ is
\begin{equation}\label{eq-wfr-gradient-flow-pde}
\partial_t\rho_t
=
\diverg(\rho_t\nabla\phi_t)
-
\frac{1}{\kappa^2}\phi_t\rho_t.
\end{equation}
Along such a smooth solution,
\begin{equation}\label{eq-wfr-energy-dissipation}
\frac{\d}{\d t}f(\alpha_t)
=
-\int
\left(\norm{\nabla\phi_t}^2+\frac{1}{\kappa^2}\phi_t^2\right)\rho_t\,\d x
\leq 0.
\end{equation}
\end{proposition}
\begin{proof}
For an infinitesimal tangent pair $(v,g)$, the density variation is $\zeta=-\diverg(\rho v)+g\rho$. The first variation of $f$ is \(\int \phi\,\zeta = \int \dotp{\nabla\phi}{v}\rho\,\d x+\int \phi g\rho\,\d x,\)
after integration by parts. For the inner product
\(\dotp{(v,g)}{(\tilde v,\tilde g)}_\rho = \int \dotp{v}{\tilde v}\rho\,\d x+\kappa^2\int g\tilde g\rho\,\d x,\)
the Riesz representative is $(\nabla\phi,\phi/\kappa^2)$. Steepest descent therefore uses $(v,g)=(-\nabla\phi,-\phi/\kappa^2)$, which gives~\eqref{eq-wfr-gradient-flow-pde}.
Substituting this pair in the first-variation formula gives
\[
\frac{\d}{\d t}f(\alpha_t)
=
-\int
\left(\norm{\nabla\phi_t}^2+\frac{1}{\kappa^2}\phi_t^2\right)\rho_t\,\d x,
\]
which is~\eqref{eq-wfr-energy-dissipation}.
\end{proof}
\subsection{Conditional Wasserstein Training of Infinite ResNets}
\label{sec-conditional-wasserstein-resnets}
\paragraph{Finite-depth finite-width ResNets.}
With $L$ layers, step size $\tau=1/L$, widths $n_r$, and parameters $\theta_{r,i}\in\Theta$, a finite ResNet is the composition of the residual updates
\begin{equation}\label{eq-finite-resnet}
z_{r+1}
=
z_r+\tau G_{r,n_r}(z_r),
\qquad
G_{r,n_r}(z)
\eqdef
\frac1{n_r}\sum_{i=1}^{n_r}\psi(\theta_{r,i},z),
\qquad
z_0=z,\qquad r=0,\ldots,L-1.
\end{equation}
The associated population loss, with $n=(n_0,\ldots,n_{L-1})$, is \(f_{L,n}((\theta_{r,i})_{r,i}) = \int \ell(z_L^{\theta}(z),y)\d\zeta(z,y).\)
Equivalently, each layer carries an empirical neuron law
\(\alpha_r = \frac1{n_r}\sum_{i=1}^{n_r}\delta_{\theta_{r,i}}, \qquad G_{r,n_r}(z) = \int_\Theta\psi(\theta,z)\d\alpha_r(\theta).\)
Thus a finite ResNet is already a discrete conditional measure on depth and neuron space, \(\alpha^{L,n}(\d s,\d\theta) = \frac1L\sum_{r=0}^{L-1}\delta_{s_r}(\d s)\alpha_r(\d\theta), \qquad s_r=r/L,\)
\paragraph{Finite-depth mean-field limit.}
Each empirical law $\alpha_r$ is replaced by a probability measure in $\Pp_2(\Theta)$, and each residual block becomes a mean-field residual block
\(z_{r+1} = z_r+\tau G_{\alpha_r}(z_r), \qquad z_0=z, \qquad G_{\alpha_r}(z) \eqdef \int_\Theta\psi(\theta,z)\d\alpha_r(\theta).\)
The loss becomes a functional of the finite family of layer measures, \(f_L(\alpha_0,\ldots,\alpha_{L-1}) = \int \ell(z_L^\alpha(z),y)\d\zeta(z,y).\)
\paragraph{Continuous-depth conditional geometry.}
The state is therefore a conditional probability law
\(\alpha(\d s,\d\theta)=\alpha_s(\d\theta)\d s,\)
or, equivalently, a probability measure on $[0,1]\times\Theta$ whose first marginal is Lebesgue measure. The general conditional Wasserstein distance is defined in Section~\ref{sec-conditional-wasserstein-distances}. For two conditional neuron laws $\alpha(\d s,\d\theta)=\alpha_s(\d\theta)\d s$ and $\beta(\d s,\d\theta)=\beta_s(\d\theta)\d s$, the relevant distance is the depthwise quadratic Wasserstein metric
\begin{equation}\label{eq-resnet-depth-wasserstein}
\Wass_{2,\lambda}^2(\alpha,\beta)
=
\int_0^1\Wass_2^2(\alpha_s,\beta_s)\d s.
\end{equation}
\paragraph{Infinite-depth mean-field ResNet.}
The forward pass is
\begin{equation}\label{eq-mean-field-resnet-forward}
\partial_s z_s
=
G_{\alpha_s}(z_s)
\eqdef
\int_\Theta \psi(\theta,z_s)\d\alpha_s(\theta),
\qquad z_0=z.
\end{equation}
The population risk is \(f_{\mathrm{Res}}(\alpha) = \int \ell(z_1^\alpha(z),y)\d\zeta(z,y).\)
\paragraph{Conditional Wasserstein gradient flow.}
For signed perturbations $\xi(\d s,\d\theta)=\xi_s(\d\theta)\lambda(\d s)$ satisfying $\int_\Theta\d\xi_s=0$ for $\lambda$-a.e.\ $s$, we use the fiberwise first-variation convention
\[
\frac{\d}{\d\varepsilon}f(\alpha+\varepsilon\xi)\bigg|_{\varepsilon=0}
=
\int_S\int_\Theta \frac{\delta f}{\delta\alpha_s}(\alpha)(\theta)\d\xi_s(\theta)\d\lambda(s).
\]
The formal conditional Wasserstein gradient flow is the family of continuity equations
\begin{equation}\label{eq-conditional-resnet-gradient-flow}
\partial_t\alpha_{t,s}
+\diverg_\theta(\alpha_{t,s}v_{t,s})=0,
\qquad
v_{t,s}(\theta)
=
-\nabla_\theta\frac{\delta f}{\delta\alpha_s}(\alpha_t)(\theta),
\end{equation}
for $\lambda$-a.e.\ condition $s$, whenever the first variation is differentiable in the parameter variable. In the ResNet case one takes $f=f_{\mathrm{Res}}$ and $S=[0,1]$.

\paragraph{Particle discretization and layerwise matching.}
If two networks have the same depth grid and the same widths, and if their empirical layer laws are \(\alpha_r=\frac1{n_r}\sum_{i=1}^{n_r}\delta_{\theta_{r,i}}, \qquad \beta_r=\frac1{n_r}\sum_{i=1}^{n_r}\delta_{\theta'_{r,i}},\)
then their squared conditional distance is \(\Wass_{2,\lambda_L}^2(\alpha^{L,n},\beta^{L,n}) = \frac1L\sum_{r=0}^{L-1} \Wass_2^2(\alpha_r,\beta_r).\)
Keeping neuron labels fixed gives the labelled parameter distance
\(\frac1L\sum_{r=0}^{L-1}\frac1{n_r} \sum_{i=1}^{n_r}\norm{\theta_{r,i}-\theta'_{r,i}}^2,\)
The Riemannian metric represented by this normalization is \(\norm{\dot\theta}^2 = \frac1L\sum_{r=0}^{L-1}\frac1{n_r} \sum_{i=1}^{n_r}\norm{\dot\theta_{r,i}}^2.\)
\subsection{Second-Order Momentum Flows}
\label{sec-second-order-momentum-flows}
\paragraph{Finite-dimensional momentum.}
Let $F:\RR^{N}\to\RR$ be a smooth finite-dimensional objective. With step size $h>0$ and damping $\lambda_{\rm fric}>0$, one convenient velocity form is
\begin{equation}\label{eq-polyak-momentum}
S^{(\ell+1)}=(1-\lambda_{\rm fric} h)S^{(\ell)}-h\nabla F(X^{(\ell)}),
\qquad
X^{(\ell+1)}=X^{(\ell)}+hS^{(\ell+1)}.
\end{equation}
If $S^{(\ell)}$ approximates $\dot X(\ell h)$, this is a semi-implicit, or symplectic-Euler, discretization of the first-order system: the velocity update is explicit, while the position uses the newly updated velocity. The limiting system is
\begin{equation}\label{eq-heavy-ball-first-order-system}
\dot X(t)=S(t),
\qquad
\dot S(t)=-\lambda_{\rm fric} S(t)-\nabla F(X(t)),
\end{equation}
or equivalently the damped second-order equation
\begin{equation}\label{eq-heavy-ball-ode}
\ddot X(t)+\lambda_{\rm fric}\dot X(t)+\nabla F(X(t))=0.
\end{equation}
Nesterov's accelerated method uses a closely related extrapolation but evaluates the gradient at a look-ahead point,
\begin{equation}\label{eq-nesterov-momentum}
Y^{(\ell)}=X^{(\ell)}+\theta^{(\ell)}(X^{(\ell)}-X^{(\ell-1)}),
\qquad
X^{(\ell+1)}=Y^{(\ell)}-h\nabla F(Y^{(\ell)}),
\end{equation}
Under the nontrivial time scaling \(t=\ell\sqrt h\), a smooth interpolation with the compatible initial condition \(\dot X(0)=0\) formally converges, as \(h\to0\), to the time-dependent damping equation
\begin{equation}\label{eq-nesterov-ode}
\ddot X(t)+\frac{r}{t}\dot X(t)+\nabla F(X(t))=0.
\end{equation}
\paragraph{Empirical Wasserstein lift.}
For $X=(x_1,\ldots,x_n)\in(\RR^d)^n$, write
\begin{equation}\label{eq-empirical-momentum-lift}
\al_X \eqdef \frac1n\sum_{i=1}^n \delta_{x_i},
\qquad
F(X) \eqdef f(\al_X).
\end{equation}
The important point is that the configuration space is not used with the plain Euclidean metric, but with the empirical metric
\(\norm{\dot X}_n^2 \eqdef \frac1n\sum_{i=1}^n\norm{\dot x_i}^2,\)
Moving only $x_i$ to $x_i+\varepsilon \xi$ gives
\[
\left.\frac{\d}{\d\varepsilon}\right|_{\varepsilon=0}
F(x_1,\ldots,x_i+\varepsilon\xi,\ldots,x_n)
=
\frac1n\dotp{\nabla\delta f(\al_X)(x_i)}{\xi}.
\]
Thus, if $\nabla^{(n)}F$ denotes the gradient for the metric $\norm{\cdot}_n$, then
\begin{equation}\label{eq-empirical-lift-gradient}
(\nabla^{(n)}F(X))_i
=
n\nabla_{x_i}F(X)
=
\nabla\delta f(\al_X)(x_i)
=
\Wgrad f(\al_X)(x_i).
\end{equation}
The first-order Wasserstein particle descent is therefore $\dot x_i=v_{\al_X}(x_i)$ with
\begin{equation}\label{eq-momentum-wasserstein-velocity}
v_\al(x)\eqdef -\Wgrad f(\al)(x)=-\nabla\delta f(\al)(x).
\end{equation}
Applying the heavy-ball equation to the empirical metric gives the second-order Wasserstein momentum system
\begin{equation}\label{eq-wasserstein-momentum-particles}
\dot x_i(t)=s_i(t),
\qquad
\dot s_i(t)=-\lambda_{\rm fric} s_i(t)+v_{\al_t}(x_i(t)),
\qquad
\al_t=\frac1n\sum_{i=1}^n\delta_{x_i(t)}.
\end{equation}
Equivalently, \(\ddot x_i(t)+\lambda_{\rm fric}\dot x_i(t) = v_{\al_t}(x_i(t)) = -\Wgrad f(\al_t)(x_i(t)).\)
\begin{prop}[Discrete mechanical energy dissipation]\label{prop-second-order-energy-dissipation}
Assume that $f$ and the particle solution are smooth enough to justify differentiating the identities above along \eqref{eq-wasserstein-momentum-particles}. Define
\(\mathcal H_n(X,S) \eqdef F(X)+\frac1{2n}\sum_{i=1}^n\norm{s_i}^2, \qquad S=(s_1,\ldots,s_n).\)
Then along \eqref{eq-wasserstein-momentum-particles},
\begin{equation}\label{eq-second-order-energy-dissipation}
\frac{\d}{\d t}\mathcal H_n(X(t),S(t))
=
-\frac{\lambda_{\rm fric}}{n}\sum_{i=1}^n\norm{s_i(t)}^2
\leq 0.
\end{equation}
\end{prop}
\begin{proof}
By \eqref{eq-empirical-lift-gradient} and the definition $v_\al=-\Wgrad f(\al)$, one has $\nabla_{x_i}F(X)=-v_{\al_X}(x_i)/n$. Hence
\(\frac{\d}{\d t}F(X(t)) = \sum_{i=1}^n \dotp{\nabla_{x_i}F(X(t))}{s_i(t)} = -\frac1n\sum_{i=1}^n\dotp{v_{\al_t}(x_i(t))}{s_i(t)}.\)
On the other hand, \(\frac{\d}{\d t}\frac1{2n}\sum_i\norm{s_i(t)}^2 = \frac1n\sum_i\dotp{s_i(t)}{-\lambda_{\rm fric} s_i(t)+v_{\al_t}(x_i(t))}.\)
The force terms cancel, leaving \eqref{eq-second-order-energy-dissipation}.
\end{proof}
\paragraph{Phase-space formulation.}
\begin{prop}[Phase-space Liouville formulation]\label{prop-second-order-liouville}
Let $v_\al$ be smooth enough that \eqref{eq-wasserstein-momentum-particles} has a classical solution. Define the empirical phase-space measure
\(\eta_t^n \eqdef \frac1n\sum_{i=1}^n\delta_{(x_i(t),s_i(t))} \in \Pp(\RR^d_x\times\RR^d_s), \qquad \al_t^n=(\pi_x)_\sharp\eta_t^n.\)
Then $\eta_t^n$ solves, in the sense of distributions,
\begin{equation}\label{eq-second-order-liouville-empirical}
\partial_t\eta_t^n
+
\nabla_x\cdot(s\,\eta_t^n)
+
\nabla_s\cdot\bigl((-\lambda_{\rm fric} s+v_{\al_t^n}(x))\,\eta_t^n\bigr)
=0.
\end{equation}
If, as $n\to\infty$, the empirical phase-space laws converge to a smooth density $\eta_t(x,s)$ and the nonlinear fields converge accordingly, the limiting kinetic equation is
\begin{equation}\label{eq-second-order-liouville-density}
\partial_t\eta_t
+
\nabla_x\cdot(s\,\eta_t)
+
\nabla_s\cdot\bigl((-\lambda_{\rm fric} s+v_{\al_t}(x))\,\eta_t\bigr)
=0,
\qquad
\al_t=(\pi_x)_\sharp\eta_t.
\end{equation}
\end{prop}
\begin{proof}
Let $\varphi\in C_c^\infty(\RR^d_x\times\RR^d_s)$. Along the particle system,
\[
\frac{\d}{\d t}
\int \varphi\,\d\eta_t^n
=
\frac1n\sum_{i=1}^n
\left[
\dotp{\nabla_x\varphi(x_i,s_i)}{s_i}
+
\dotp{\nabla_s\varphi(x_i,s_i)}{-\lambda_{\rm fric} s_i+v_{\al_t^n}(x_i)}
\right].
\]
Rewriting the right-hand side as an integral against $\eta_t^n$ and integrating by parts gives \eqref{eq-second-order-liouville-empirical}. Passing formally to the mean-field limit gives \eqref{eq-second-order-liouville-density}.
\end{proof}
The numerical illustration below uses the energy-distance kernel
\(k(x,y)=-\norm{x-y}, \qquad f(\al)=\MMD_k^2(\al,\beta) = \iint k(x,y)\,\d(\al-\beta)(x)\d(\al-\beta)(y).\)
Away from the diagonal, its Wasserstein velocity is \(v_\al(x) = 2\int \frac{x-y}{\norm{x-y}}\,\d(\al-\beta)(y).\)
\paragraph{Entropy without an empirical energy.}
For the numerical illustration, we add the quadratic confinement $V(x)=\norm{x}^2/2$. In one dimension the overdamped equation becomes the Ornstein--Uhlenbeck Fokker--Planck equation
\(\partial_t\rho_t=\partial_{xx}\rho_t+\partial_x(x\rho_t),\)
whose stationary law is the centered Gaussian $\Gaussian(0,1)$. The Newton lift is discretized directly in phase space: for a density $\eta_t(x,s)$ with spatial marginal $\rho_t(x)=\int\eta_t(x,s)\,\d s$, it reads
\(\partial_t\eta_t+\partial_x(s\eta_t)+\partial_s\bigl((-\partial_x\log\rho_t(x)-x)\eta_t\bigr)=0.\)

\section{Generative Models via Transportation}
\label{sec-generative-models-transportation}
\subsection{Generative Models via Flow Matching}
\label{sec-generative-flow-matching}
\paragraph{Deterministic interpolants.}
\begin{defn}[Deterministic interpolant]\label{def-deterministic-interpolant}
Let \(\pi\in\Couplings(\alpha_0,\alpha_1)\). A deterministic interpolant is a jointly measurable family \(I_t:\RR^d\times\RR^d\to\RR^d\) such that \(t\mapsto I_t(x,y)\) is absolutely continuous for \(\pi\)-almost every \((x,y)\), its derivative is jointly measurable, \(\int_0^1\!\int_{\RR^d\times\RR^d}\norm{\partial_t I_t(x,y)}\,\d\pi(x,y)\d t<+\infty\), and \(I_0(x,y)=x\), \(I_1(x,y)=y\) for \(\pi\)-almost every \((x,y)\).
It induces \(\alpha_t=(I_t)_\sharp\pi\) and has pathwise velocity \(\partial_t I_t(x,y)\).
\end{defn}
More generally, $\alpha_t$ can be obtained by pushing a latent distribution $\pi \in \Pp(\RR^{d'})$ through a time-dependent map $I_t: \RR^{d'} \to \RR^d$; the latent dimension $d'$ may be larger than the data dimension $d$, i.e.
\begin{equation}\label{eq:interp-coupling}
\forall t \in [0,1], \quad \alpha_t:= (I_t)_\sharp \pi.
\end{equation}
\paragraph{Learning a flow field.}
\begin{prop}[Flow matching vector field]\label{prop-flow-matching-vector-field}
Assume that $t\mapsto I_t(u)$ is absolutely continuous for $\pi$-almost every $u$, that its derivative $(t,u)\mapsto\partial_t I_t(u)$ is jointly measurable, and that \(\int_0^1\!\int_{\RR^{d'}}\norm{\partial_t I_t(u)}^2\d\pi(u)\d t<+\infty\).
Let \(U\sim\pi\) and \(X_t=I_t(U)\), so that \(X_t\sim\alpha_t\).
For almost every $t$, consider the flow-matching problem over \(v_t\in L^2(\alpha_t;\RR^d)\)
\begin{equation}
\min_{v_t} \int_{\RR^{d'}} \norm{v_t(I_t(u)) - [\partial_t I_t](u)}^2 \, \d\pi(u). \label{eq:flow-matching}
\end{equation}
Its minimizer is characterized $\alpha_t$-almost everywhere by
\begin{equation}\label{eq:flow-match-conditional}
v_t(x)=\EE\bigl[\partial_t I_t(U)\mid X_t=x\bigr].
\end{equation}
There is a jointly measurable choice \((t,x)\mapsto v_t(x)\), obtained by disintegrating \(\d t\,\pi(\d u)\) under \((t,u)\mapsto(t,I_t(u))\).
\end{prop}
\begin{proof}
The push-forward identity reads \(\int_{\RR^d}\varphi(x)\,\d\alpha_t(x)=\int_{\RR^{d'}}\varphi(I_t(u))\,\d\pi(u)\) for every bounded continuous test function $\varphi$.
For almost every $t$, the minimizer in~\eqref{eq:flow-matching} is the orthogonal projection in $L^2(\pi;\RR^d)$ of the latent velocity $\partial_t I_t(u)$ onto the closed subspace of functions that depend on $u$ only through $I_t(u)$. This projection is the conditional expectation~\eqref{eq:flow-match-conditional}. Equivalently, disintegration of the joint measure \(\d t\,\pi(\d u)\) by the map \((t,u)\mapsto(t,I_t(u))\) supplies a jointly measurable version. Conditional Jensen gives \(\int\norm{v_t}^2\d\alpha_t\leq\int\norm{\partial_t I_t}^2\d\pi\),
so the Eulerian velocity is square-integrable in space--time. Its defining conditional-expectation identity is, for every \(m\in L^2(\alpha_t;\RR^d)\),
\begin{equation}\label{eq:v_t}
\int \dotp{m(x)}{v_t(x)} \, \d\alpha_t(x)
=
\int \dotp{m(I_t(u))}{[\partial_t I_t](u)} \, \d\pi(u).
\end{equation}

Testing first with \(\varphi\in C_c^\infty(\RR^d)\), and then integrating against compactly supported smooth functions of time, gives the distributional formulation on space--time. The fixed-time weak identity for
$\partial_t\alpha_t+\diverg(\alpha_t v_t)=0$ is that, for every $\varphi\in C_c^\infty(\RR^d)$,
\begin{equation}\label{eq:flow-matching-weak-target}
\frac{\d}{\d t}\int\varphi(x)\,\d\alpha_t(x)
-
\int\dotp{v_t(x)}{\nabla\varphi(x)}\,\d\alpha_t(x)
=0.
\end{equation}
By the pathwise chain rule, $t\mapsto\varphi(I_t(u))$ is absolutely continuous for $\pi$-almost every $u$. The space--time integrability assumption and Fubini's theorem justify integration of its derivative, so the push-forward identity gives, for almost every $t$,
\begin{equation}\label{eq:flow-matching-test-derivative}
\frac{\d}{\d t}\int \varphi(x)\d\alpha_t(x)
=
\int \dotp{\nabla\varphi(I_t(u))}{[\partial_t I_t](u)}\d\pi(u).
\end{equation}
On the other hand, applying~\eqref{eq:v_t} with $m=\nabla\varphi$ gives
\begin{equation}\label{eq:flow-matching-velocity-test}
\int\dotp{v_t(x)}{\nabla\varphi(x)}\,\d\alpha_t(x)
=
\int \dotp{\nabla\varphi(I_t(u))}{[\partial_t I_t](u)}\d\pi(u).
\end{equation}
Comparing~\eqref{eq:flow-matching-test-derivative} and~\eqref{eq:flow-matching-velocity-test} yields~\eqref{eq:flow-matching-weak-target}, which is the desired continuity equation.
\end{proof}
\paragraph{Sampling algorithm.}
Once the field $v_t$ has been estimated by the tractable regression~\eqref{eq:flow-matching}, generation amounts to integrating its characteristics. For a sufficiently regular version of the exact population field, let $T_t$ denote the flow map solving
\begin{equation}\label{eq-flow-matching-sampling-map}
\partial_tT_t(x)=v_t(T_t(x)),
\qquad T_0(x)=x.
\end{equation}
\paragraph{Connection with diffusion models.}\label{par-diffusion-model-connection}
In the special case where $I_t(x,y)=(1-t)x+ty$ is a linear interpolation and $\pi = \alpha \otimes \beta$, the curve $\alpha_t$ is a convolution of rescaled versions of $\alpha_0$ and $\alpha_1$. The flow-matching problem~\eqref{eq:flow-matching} becomes
\(\min_{(v_t)_t} \int_0^1\!\int_{\RR^{d} \times \RR^d} \norm{v_t( (1-t)x+t y ) - (y-x) }^2 \, \d\alpha_0(x) \d\alpha_1(y)\,\d t.\)
\begin{proposition}[Tweedie identity]\label{prop:Tweedie}
Let $W$ be a random vector in $\RR^{d}$ with law $\beta\in\Pp_1(\RR^d)$.
For $\sigma>0$, let $Z\sim\Gaussian(0,\Id)$ be independent of $W$ and observe
\(Y=W+\sigma Z.\)
Denote by $\rho_\sigma$ the smooth positive density of \(\beta_\sigma=\beta * \Gaussian(0,\sigma^{2}\Id)\), which is the law of $Y$.
Then the conditional mean admits the following everywhere-defined version: \(\EE\bigl[W\mid Y=y\bigr] =y+\sigma^{2}\nabla \log \rho_\sigma(y) \qquad\text{for all } y \in \RR^{d}.\)
\end{proposition}
\begin{proof}
Let $\varphi_\sigma$ be the $\Gaussian(0,\sigma^{2}\Id)$ density. Bayes' rule gives \(\EE[W\mid Y=y]=\rho_\sigma(y)^{-1}\int_{\RR^{d}} w\varphi_\sigma(y-w)\,\d\beta(w)\).
Differentiating the Gaussian convolution under the integral sign and using
$
\nabla_y\varphi_\sigma(y-w)
= -\sigma^{-2}(y-w)\,\varphi_\sigma(y-w)
$
yields
\(\nabla\rho_\sigma(y) = \int \nabla_y\varphi_\sigma(y-w)\,\d\beta(w) = -\sigma^{-2}\Bigl(y-\EE[W\mid Y=y]\Bigr)\,\rho_\sigma(y).\)
Rearranging finishes the proof.
\end{proof}
\begin{proposition}[Gaussian-endpoint flow-matching field]\label{prop:flow}
Let \(\alpha\in\Pp_2(\RR^d)\), let $X_0\sim\alpha$, and let $Z\sim\Gaussian(0,\Id)$ be independent.
For $t\in(0,1)$ set \(X_t=(1-t)X_0+tZ\) and \(\alpha_t=\operatorname{Law}(X_t)=\rho_t\,\d x\). The minimizer $v^\star:\RR^d\times(0,1)\to\RR^d$ of the flow-matching regression~\eqref{eq:flow-matching} is
\(v^\star(x,t) = -\frac{1}{1-t}\,x \;-\; \frac{t}{1-t}\,\nabla\log\rho_t(x) \qquad (x\in\RR^{d},\;t\in(0,1)).\)
In particular, for each $t\in(0,1)$ this field is the gradient field \(v^\star(\cdot,t)=-\nabla\bigl(\norm{\cdot}^2/[2(1-t)]+t\log\rho_t/(1-t)\bigr)\).
\end{proposition}
\begin{proof}
Fix $t\in(0,1)$ and write $W=(1-t)X_0$ and $\sigma=t$, so that
$X_t=W+\sigma Z$ matches the setting of Proposition~\ref{prop:Tweedie}.
Conditional expectations satisfy
$
v^\star(x,t)
= \EE[Z-X_0\mid X_t=x]
= \frac{1}{t}\,\EE[X_t-W\mid X_t=x]
-\,\frac{1}{1-t}\,\EE[W\mid X_t=x].
$
Applying Proposition~\ref{prop:Tweedie} to $\EE[W\mid X_t=x]$ and
noting $\EE[Z\mid X_t=x]=-t\nabla\log\rho_t(x)$
gives the claimed formula.
\end{proof}
\paragraph{When is the induced map optimal?}
Integrating the learned velocity gives a deterministic map from $\alpha_0$ to $\alpha_1$, but this map is not automatically the Brenier optimal map. It is optimal only in special cases where the accumulated flow remains the gradient of a convex potential.

\begin{prop}[Gaussian flow matching and optimality]\label{prop-gaussian-flow-matching-optimality}
Let $\Sigma_0,\Sigma_1\succ0$ and let $X_0\sim\Gaussian(0,\Sigma_0)$ and $X_1\sim\Gaussian(0,\Sigma_1)$ be independent. Consider the linear flow-matching interpolation
\(X_t=(1-t)X_0+tX_1, \qquad \alpha_t=\operatorname{Law}(X_t)=\Gaussian(0,\Sigma_t),\)
where $\Sigma_t=(1-t)^2\Sigma_0+t^2\Sigma_1$. Then the exact flow-matching velocity is affine, $v_t(x)=A_tx$, with $A_t=\bigl(t\Sigma_1-(1-t)\Sigma_0\bigr)\Sigma_t^{-1}$.
The induced flow map $T_t^{\rm FM}$ from $\alpha_0$ to $\alpha_t$ is
\begin{equation}\label{eq-gaussian-product-fm-map}
T_t^{\rm FM}
=
\Sigma_0^{1/2}
\Bigl((1-t)^2\Id+t^2\Sigma_0^{-1/2}\Sigma_1\Sigma_0^{-1/2}\Bigr)^{1/2}
\Sigma_0^{-1/2}.
\end{equation}
\eqllead{In particular,}{eq-gaussian-product-fm-terminal-map}{
T_1^{\rm FM}
=
\Sigma_0^{1/2}
\bigl(\Sigma_0^{-1/2}\Sigma_1\Sigma_0^{-1/2}\bigr)^{1/2}
\Sigma_0^{-1/2}.
}
This terminal map coincides with the quadratic optimal transport map
\begin{equation}\label{eq-gaussian-brenier-map-comparison}
T^{\rm OT}
=
\Sigma_0^{-1/2}
\bigl(\Sigma_0^{1/2}\Sigma_1\Sigma_0^{1/2}\bigr)^{1/2}
\Sigma_0^{-1/2}
\end{equation}
if and only if $\Sigma_0\Sigma_1=\Sigma_1\Sigma_0$.
\end{prop}
\begin{proof}
The conditional-expectation formula gives
\(v_t(x)=\EE[X_1-X_0\mid X_t=x]\).
Since all variables are jointly Gaussian, this conditional expectation is linear and
\(v_t(x) = \operatorname{Cov}(X_1-X_0,X_t)\operatorname{Cov}(X_t)^{-1}x = \bigl(t\Sigma_1-(1-t)\Sigma_0\bigr)\Sigma_t^{-1}x,\)
which proves the stated velocity formula. To solve the characteristic equation, whiten the source by setting \(C=\Sigma_0^{-1/2}\Sigma_1\Sigma_0^{-1/2}\) and \(\widetilde X_t=\Sigma_0^{-1/2}X_t\).
In these coordinates the source covariance is $\Id$ and
\(\widetilde\Sigma_t=(1-t)^2\Id+t^2C\).
Because $\Id$ and $C$ commute, the affine flow map in whitened coordinates is simply $\widetilde T_t=\widetilde\Sigma_t^{1/2}$. Indeed,
$\frac{\d}{\d t}\widetilde\Sigma_t^{1/2}=\bigl(tC-(1-t)\Id\bigr)\widetilde\Sigma_t^{-1/2}$, which is exactly the equation $\dot{\widetilde T}_t=\widetilde A_t\widetilde T_t$ with $\widetilde T_0=\Id$. Returning to the original coordinates gives~\eqref{eq-gaussian-product-fm-map}, and $t=1$ gives~\eqref{eq-gaussian-product-fm-terminal-map}.

Both $T_1^{\rm FM}$ and $T^{\rm OT}$ push $\Gaussian(0,\Sigma_0)$ to $\Gaussian(0,\Sigma_1)$. The Brenier map between nondegenerate Gaussians is the unique symmetric positive definite linear map with this property. Hence $T_1^{\rm FM}=T^{\rm OT}$ if and only if $T_1^{\rm FM}$ is symmetric positive definite. The map $T_1^{\rm FM}$ is similar to $C^{1/2}$, so if it is symmetric then it is automatically positive definite. Since $C^{1/2}$ is symmetric positive definite, \((T_1^{\rm FM})^\top=\Sigma_0^{-1/2}C^{1/2}\Sigma_0^{1/2}\).
Thus symmetry of $T_1^{\rm FM}$ is equivalent to
$\Sigma_0 C^{1/2}=C^{1/2}\Sigma_0$, hence to $\Sigma_0 C=C\Sigma_0$ by functional calculus. Multiplying this identity on the left and right by $\Sigma_0^{1/2}$ gives $\Sigma_0\Sigma_1=\Sigma_1\Sigma_0$. Conversely, if $\Sigma_0$ and $\Sigma_1$ commute, they are orthogonally co-diagonalizable, and both~\eqref{eq-gaussian-product-fm-terminal-map} and~\eqref{eq-gaussian-brenier-map-comparison} reduce in that basis to the diagonal map with entries $\sqrt{\lambda_{1,k}/\lambda_{0,k}}$. This proves the equivalence.
\end{proof}
\paragraph{Variations on the interpolant.}
There is first a harmless ambiguity: a monotone reparametrization $X_t=(1-\lambda(t))X_0+\lambda(t)X_1$ of the linear bridge only changes the speed of the flow,
\(v_t(x)=\lambda'(t)\,v^{\rm lin}_{\lambda(t)}(x), \qquad v^{\rm lin}_{r}(x)=\EE[X_1-X_0\mid (1-r)X_0+rX_1=x].\)
\begin{prop}[Static Gaussian interpolants and Ornstein--Uhlenbeck noising]\label{prop-static-ou-equivalence}
Let $\alpha\in\Pp_2(\RR^d)$ and $\sigma>0$. Let $X_0\sim\alpha$ and $Z\sim\Gaussian(0,\sigma^2\Id)$ be independent, and set $X_t=a_tX_0+b_tZ$ on an open time interval where the differentiable scalar functions satisfy $a_t>0$ and $b_t>0$. If $\rho_t$ is the density of $X_t$, then its flow-matching velocity is
\begin{equation}\label{eq-scalar-gaussian-interpolant-velocity}
v_t(x)=\frac{\dot a_t}{a_t}x+
\left(\frac{\dot a_t b_t^2}{a_t}-\dot b_t b_t\right)\sigma^2\nabla\log \rho_t(x).
\end{equation}
For the variance-preserving coefficients $a_\tau=e^{-\tau}$ and $b_\tau=\sqrt{1-e^{-2\tau}}$, with $\tau>0$, these marginals are exactly those of the Ornstein--Uhlenbeck diffusion
\begin{equation}\label{eq-ou-forward-diffusion}
\d X_\tau^{\rm OU}=-X_\tau^{\rm OU}\d\tau+\sqrt{2}\,\sigma\d B_\tau,
\qquad X_0^{\rm OU}\sim\alpha,
\end{equation}
and~\eqref{eq-scalar-gaussian-interpolant-velocity} reduces to $v_\tau(x)=-x-\sigma^2\nabla\log \rho_\tau(x)$. Moreover, the finite-time coefficients $a_t=\cos(\pi t/2)$ and $b_t=\sin(\pi t/2)$ describe the same marginal curve after the time change $\tau=-\log\cos(\pi t/2)$.
\end{prop}
\begin{proof}
Tweedie's identity gives $\EE[X_0\mid X_t=x]=a_t^{-1}(x+b_t^2\sigma^2\nabla\log \rho_t(x))$ and $\EE[Z\mid X_t=x]=-b_t\sigma^2\nabla\log \rho_t(x)$. Substitution in $v_t(x)=\dot a_t\EE[X_0\mid X_t=x]+\dot b_t\EE[Z\mid X_t=x]$ proves~\eqref{eq-scalar-gaussian-interpolant-velocity}. The explicit solution of~\eqref{eq-ou-forward-diffusion} is $X_\tau^{\rm OU}=e^{-\tau}X_0^{\rm OU}+\sqrt{2}\,\sigma\int_0^\tau e^{-(\tau-s)}\d B_s$. The stochastic integral is independent of $X_0^{\rm OU}$ and Gaussian with covariance $\sigma^2(1-e^{-2\tau})\Id$, proving equality of the one-time marginals. Substituting the OU coefficients in~\eqref{eq-scalar-gaussian-interpolant-velocity} gives the stated velocity, and the trigonometric parametrization follows from $e^{-\tau}=\cos(\pi t/2)$.
\end{proof}
\paragraph{Stochastic interpolants.}
\begin{defn}[Stochastic interpolant]\label{def-stochastic-interpolant}
Let $(X_0,X_1)\sim\pi\in\Couplings(\alpha_0,\alpha_1)$, let $Z\sim\Gaussian(0,\Id)$ be independent of $(X_0,X_1)$, and let $I_t$ be a deterministic interpolant for this coupling. Given an absolutely continuous noise schedule $\gamma:[0,1]\to\RR$ satisfying $\gamma(0)=\gamma(1)=0$, the stochastic interpolant is
\begin{equation}\label{eq-static-noise-stochastic-interpolant}
X_t=I_t(X_0,X_1)+\gamma(t)Z,
\qquad \alpha_t=\operatorname{Law}(X_t).
\end{equation}
Assume in addition that $\EE\int_0^1\norm{\partial_tI_t(X_0,X_1)+\dot\gamma(t)Z}^2\,\d t<+\infty$.
\end{defn}
The same draw of $Z$ is used for every $t$: conditionally on $(X_0,X_1,Z)$, the path remains absolutely continuous. Thus~\eqref{eq-static-noise-stochastic-interpolant} is the push-forward construction~\eqref{eq:interp-coupling} with latent variable $u=(X_0,X_1,Z)$. Its Eulerian velocity is
\(v_t(x)=\EE\bigl[\partial_tI_t(X_0,X_1)+\dot\gamma(t)Z\mid X_t=x\bigr],\)
\subsection{One-Step Generative Models}
\label{sec-one-step-generative-models}
\paragraph{Non-parametric sampling via gradient flow.}
Let $\beta$ be the target law and let $\mathcal D$ be a discrepancy that vanishes only on the diagonal. Define
\begin{equation}\label{eq-one-step-discrepancy-energy}
\mathcal E_\beta(\alpha)\eqdef\mathcal D(\alpha,\beta).
\end{equation}
A direct sampling strategy follows the formal $\Wass_2$ gradient flow
\begin{equation}\label{eq-nonparametric-discrepancy-flow}
\partial_t\alpha_t+\diverg(\alpha_t v_t)=0,
\qquad
v_t=-\Wgrad\mathcal E_\beta(\alpha_t)
=-\nabla\delta_\alpha\mathcal E_\beta(\alpha_t).
\end{equation}
If this flow is globally well posed and converges to a global minimizer of $\mathcal E_\beta$, then its limit is $\beta$. These are substantial assumptions, not automatic consequences of the definition of a discrepancy. If $\alpha_0=n^{-1}\sum_{i=1}^n\delta_{X_i(0)}$ and the dynamics preserves empirical measures, then~\eqref{eq-nonparametric-discrepancy-flow} becomes the particle system
\begin{equation}\label{eq-nonparametric-particle-flow}
\dot X_i(t)=v_t(X_i(t))
=-\Wgrad\mathcal E_\beta(\alpha_t)(X_i(t)),
\qquad
\alpha_t=\frac1n\sum_{i=1}^n\delta_{X_i(t)}.
\end{equation}
If $T_\theta$ denotes the monotone transport from $(P_\theta)_\sharp\alpha$ to $(P_\theta)_\sharp\beta$, then the formal descent velocity is
\begin{equation}\label{eq-sliced-wasserstein-flow-velocity}
v_{\SW}[\alpha,\beta](x)
=
\int_{\Sphere^{d-1}}
\big(T_\theta(P_\theta x)-P_\theta x\big)\,\theta\,\d\sigma(\theta).
\end{equation}
\paragraph{One-step models via parameter-domain discrepancy flows.}
The most direct one-step strategy is to train the generator parameters themselves by descending the discrepancy~\eqref{eq-one-step-discrepancy-energy}. Let $\zeta$ be a simple latent distribution and let $f_\theta:\Zz\to\Xx$ be a neural generator. Define
\begin{equation}\label{eq-one-step-parameter-discrepancy}
H_\beta(\theta)
\eqdef
\mathcal E_\beta\big((f_\theta)_\sharp\zeta\big)
=
\mathcal D\big((f_\theta)_\sharp\zeta,\beta\big).
\end{equation}
The resulting training dynamics is the ordinary Euclidean gradient flow in parameter space,
\begin{equation}\label{eq-one-step-parameter-flow}
\dot\theta_t=-\nabla_\theta H_\beta(\theta_t).
\end{equation}
Set
\(\alpha_t\eqdef(f_{\theta_t})_\sharp\zeta, \qquad X_t=f_{\theta_t}(Z), \qquad Z\sim\zeta.\)
\begin{prop}[Velocity induced by a parameter path]\label{prop-parameter-path-continuity}
Assume that $\Xx=\RR^d$, that $t\mapsto\theta_t$ is continuously differentiable on $[0,T]$, and that, for $\zeta$-almost every $z$, the map $\theta\mapsto f_\theta(z)$ is continuously differentiable. Assume also that $(t,z)\mapsto\partial_\theta f_{\theta_t}(z)$ is jointly measurable and that
\(\int_0^T\!\int_\Zz \norm{\partial_\theta f_{\theta_t}(z)\dot\theta_t} \,\d\zeta(z)\d t<+\infty.\)
Let $\eta_{t,x}$ be a jointly measurable disintegration of $\d t\,\zeta(\d z)$ under $(t,z)\mapsto(t,f_{\theta_t}(z))$. Then the Eulerian velocity
\begin{equation}\label{eq-one-step-induced-velocity}
v_t(x)=\int \partial_\theta f_{\theta_t}(z)\dot\theta_t\,\d\eta_{t,x}(z)
\end{equation}
satisfies, in the distributional sense on $(0,T)\times\RR^d$,
\begin{equation}\label{eq-one-step-induced-continuity}
\partial_t\alpha_t+\diverg(\alpha_t v_t)=0.
\end{equation}
If $\theta_t$ follows~\eqref{eq-one-step-parameter-flow}, then $\dot\theta_t=-\nabla_\theta H_\beta(\theta_t)$ in~\eqref{eq-one-step-induced-velocity}.
\end{prop}
\begin{proof}
The chain rule gives $\dot X_t=\partial_\theta f_{\theta_t}(Z)\dot\theta_t$. For every bounded measurable $\psi$, disintegration means
\[
\int \psi(z)\,\d\zeta(z)=\int_\Xx\left(\int \psi(z)\,\d\eta_{t,x}(z)\right)\d\alpha_t(x)
\]
for almost every $t$.
\(\frac{\d}{\d t}\int \varphi(x)\,\d\alpha_t(x) = \int \dotp{\nabla\varphi(x)}{v_t(x)}\,\d\alpha_t(x),\)
which is the weak form of~\eqref{eq-one-step-induced-continuity}.
\end{proof}
The velocity~\eqref{eq-one-step-induced-velocity} depends on the parameter value $\theta_t$ and on the latent disintegration, not only on the measure $\alpha_t$. If the parametrization is non-identifiable, two different parameters can generate the same measure while inducing different velocities. Unlike standard stochastic-gradient constructions, this substitution is not generally unbiased for nonlinear discrepancies:
\[
\EE\left[\nabla_\theta H_{\hat\beta_n}(\theta)\right]
\neq
\nabla_\theta H_\beta(\theta).
\]
\paragraph{One-step models from Wasserstein flows.}
Let $\zeta$ be a simple latent distribution and let $\alpha_\theta=(G_\theta)_\sharp\zeta$ be the model distribution. Typical choices for $\mathcal E_\beta$ include a smoothed relative entropy, an MMD/IPM loss, or the debiased Sinkhorn divergence $\bar\MK_\c^\epsilon(\alpha,\beta)$ introduced in Section~\ref{sec-sinkhorn-div}. Writing its formal Wasserstein descent as
\begin{equation}\label{eq-one-step-wgf}
\partial_t\al_t+\diverg(\al_t v_t)=0,
\qquad
v_t(x)=-\Wgrad\mathcal E_\beta(\al_t)(x)
=-\nabla\delta_\alpha \mathcal E_\beta(\al_t)(x),
\end{equation}
one fits, at each training time $t$, a parametric residual field $U_{\eta_t}$ along the current model distribution:
\begin{equation}\label{eq-one-step-l2-fit}
\min_\eta \int
\norm{U_\eta(x)-v_t(x)}^2
\,\d\al_t(x).
\end{equation}
In a particle or generator implementation, the learned residual is then used to update the current generator by
\[
\alpha_{\theta}^{+}
=
(\Id+\tau U_{\eta_t})_\sharp \alpha_\theta,
\qquad\text{or equivalently}\qquad
G_\theta^{+}(z)=G_\theta(z)+\tau U_{\eta_t}(G_\theta(z)).
\]
\paragraph{Stein variational gradient descent.}\label{sec-svgd-generative-flow}
\begin{defn}[SVGD vector field]\label{def-svgd-vector-field}
Let $\beta=\rho_\beta\d x$, where $\rho_\beta$ is positive and continuously differentiable. Let $k$ be a continuously differentiable positive-definite kernel with scalar RKHS $\RKHS_k$, and assume that its first-derivative evaluations are bounded on $\RKHS_k$. If the $\RKHS_k^d$-valued map $y\mapsto k(y,\cdot)\nabla\log\rho_\beta(y)+\nabla_yk(y,\cdot)$ is Bochner integrable under $\alpha$, the SVGD field is the unique minimizer
\begin{equation}\label{eq-svgd-local-problem}
v_\alpha^{\mathrm{SVGD}}
\eqdef
\uargmin{v\in\RKHS_k^d}
\left\{
-\int\!\bigl(\dotp{\nabla\log\rho_\beta(x)}{v(x)}+\diverg v(x)\bigr)\,\d\alpha(x)
+\frac12\norm{v}_{\RKHS_k^d}^2
\right\}.
\end{equation}
By the reproducing property, it has the explicit form
\begin{equation}\label{eq-svgd-velocity}
v_{\alpha}^{\mathrm{SVGD}}(x)
=
\int
\Big(k(y,x)\nabla\log\rho_\beta(y)+\nabla_y k(y,x)\Big)
\d\alpha(y).
\end{equation}
\end{defn}
For smooth positive $\alpha$, under the usual Stein boundary condition, the linear term in~\eqref{eq-svgd-local-problem} is the directional derivative of $\KL(\cdot|\beta)$ along $(\Id+s v)_\sharp\alpha$. Unlike that density-level expression, the right-hand side of~\eqref{eq-svgd-velocity} remains meaningful for empirical $\alpha$.
\eqllead{Thus}{eq-svgd-mean-field}{
\partial_t\alpha_t+\diverg\bigl(\alpha_t v_{\alpha_t}^{\mathrm{SVGD}}\bigr)=0.
}
For particles $\alpha_n^{(\ell)}=n^{-1}\sum_i\delta_{X_i^{(\ell)}}$, this becomes the explicit update
\begin{equation}\label{eq-svgd-particle-update}
X_i^{(\ell+1)}
=
X_i^{(\ell)}
+
\tau\,\frac{1}{n}\sum_{j=1}^n
\Big(k(X_j^{(\ell)},X_i^{(\ell)})\nabla\log\rho_\beta(X_j^{(\ell)})
+
\nabla_1k(X_j^{(\ell)},X_i^{(\ell)})\Big).
\end{equation}
\paragraph{Drifting flows.}
A simple continuous construction uses a strictly positive kernel $K_\epsilon(x,y)$ for which the following integrals are finite and defines, for every probability measure $\nu$,
\begin{equation}\label{eq-normalized-kernel-drift}
B_\epsilon[\nu](x)
\eqdef
\frac{\int (y-x)K_\epsilon(x,y)\,\d\nu(y)}
{\int K_\epsilon(x,y)\,\d\nu(y)}.
\end{equation}
For the Gaussian kernel $K_\epsilon(x,y)=\exp(-\norm{x-y}^2/(2\epsilon))$, this normalized field is a score of a smoothed density:
\begin{equation}\label{eq-normalized-kernel-score}
B_\epsilon[\nu](x)
=
\epsilon\nabla\log\!\left(\int K_\epsilon(x,y)\,\d\nu(y)\right).
\end{equation}
The drifting velocity is
\begin{equation}\label{eq-cross-minus-self-drift}
v_t(x)=B_\epsilon[\beta](x)-B_\epsilon[\al_t](x)
=
\epsilon\nabla\log
\frac{\int K_\epsilon(x,y)\,\d\beta(y)}
{\int K_\epsilon(x,y)\,\d\al_t(y)}.
\end{equation}
The first term pulls samples toward data, while the second acts as a repulsive correction before coalescence. If all empirical particles occupy the same point, their self displacement is zero and they all receive the same target drift, so this deterministic field cannot recreate diversity.

Sampling follows the usual continuity equation $\partial_t\al_t+\diverg(\al_t v_t)=0$ from a prescribed simple law $\al_{t=0}=\al_0$; equivalently, particles solve $\dot X_i(t)=v_t(X_i(t))$. A direct implementation repeatedly evaluates~\eqref{eq-cross-minus-self-drift}.

\subsection{Moment Measures}
\label{sec-moment-measures}
\begin{defn}[Moment measure of a convex function]\label{def-moment-measure}
Let $u:\RR^d\to\RR\cup\{+\infty\}$ be a proper lower-semicontinuous convex function such that $0<Z_u\eqdef \int_{\RR^d} e^{-u(x)}\,\d x <+\infty$. The normalized log-concave measure associated with $u$ is $\eta_u \eqdef Z_u^{-1} e^{-u(x)}\,\d x$.
These assumptions force the effective domain of $u$ to have nonempty interior, so convex differentiability implies that $\nabla u$ exists $\eta_u$-almost everywhere. The moment measure of $u$ is
\(\mathfrak M(u)\eqdef (\nabla u)_\sharp \eta_u.\)
\end{defn}
Formally, if $u$ is smooth and $e^{-u}$ decays fast enough for the boundary term to vanish, then
\(\int y\,\d\mathfrak M(u)(y) = Z_u^{-1}\int \nabla u(x)e^{-u(x)}\,\d x = -Z_u^{-1}\int \nabla(e^{-u(x)})\,\d x =0.\)
Thus moment measures are necessarily centered. The nonsmooth theory uses essentially continuous convex functions: these are lower-semicontinuous convex functions whose set of discontinuity points has zero $\mathcal H^{d-1}$ measure.

\begin{thm}[Cordero--Erausquin--Klartag]\label{thm-moment-measure-characterization}
Let $\al\in\Pp_1(\RR^d)$ be a probability measure. There exists an essentially continuous convex function $u$ such that $\al=\mathfrak M(u)$ if and only if
\(\int y\,\d\al(y)=0 \qquad\text{and}\qquad \supp(\al)\text{ is not contained in a hyperplane.}\)
Under these conditions, $u$ is unique up to translations of the argument and additive constants.
\end{thm}
\paragraph{Optimal-transport variational formulation.}
For a centered $\al\in\Pp_1(\RR^d)$, define the maximal-correlation transport functional, with values in $\RR\cup\{+\infty\}$,
\begin{equation}\label{eq-moment-max-correlation}
\mathsf{MC}_\al(\eta)
\eqdef
\sup_{\pi\in\Couplings(\eta,\al)}
\int_{\RR^d\times\RR^d} x\cdot y\,\d\pi(x,y).
\end{equation}
By Kantorovich duality for the scalar-product cost,
\begin{equation}\label{eq-moment-max-correlation-dual}
\mathsf{MC}_\al(\eta)
=
\inf_{v\ \mathrm{convex}}
\left\{
\int v(x)\,\d\eta(x)+\int v^*(y)\,\d\al(y)
\right\},
\end{equation}
where the infimum is over convex functions for which both integrals are well defined and $v^*$ is the Legendre transform. If $\eta,\al\in\Pp_2(\RR^d)$, then
\begin{equation}\label{eq-moment-correlation-w2}
\mathsf{MC}_\al(\eta)
=
\frac12\int \norm{x}^2\,\d\eta(x)
+
\frac12\int \norm{y}^2\,\d\al(y)
-
\frac12\Wass_2^2(\eta,\al).
\end{equation}
The variational problem attached to a centered target $\al$ is
\begin{equation}\label{eq-moment-measure-variational}
\min_{\eta\in\Pp_1(\RR^d)}
\left\{
\mathcal H(\eta)+\mathsf{MC}_\al(\eta)
\right\},
\qquad
\mathcal H(r\,\d x)\eqdef \int r(x)\log r(x)\,\d x,
\end{equation}
where $\mathcal H$ is the relative entropy with respect to Lebesgue measure, in the sigma-finite-reference convention corresponding to the KL divergence of Definition~\ref{def-measure-relative-entropy}; set $\mathcal H(\eta)=+\infty$ when $\eta$ is not absolutely continuous.

\begin{prop}[Variational characterization of moment measures]\label{prop-moment-hidden-convexity}
Let $\al\in\Pp_1(\RR^d)$ be centered and not supported on a hyperplane. Then the minimization problem~\eqref{eq-moment-measure-variational} admits a solution, unique up to translations. Every minimizer is a probability measure with log-concave density $\eta=Z_u^{-1}e^{-u}\,\d x$, where $u$ is convex, essentially continuous, and satisfies the moment-measure equation $\al=(\nabla u)_\sharp \eta$.
Conversely, any essentially continuous convex $u$ satisfying this equation yields a global minimizer. If $\al\in\Pp_2(\RR^d)$, the objective $\eta\mapsto\mathcal H(\eta)+\mathsf{MC}_\al(\eta)$ is displacement convex along $\Wass_2$ geodesics of absolutely continuous measures in $\Pp_2(\RR^d)$ whenever the terms are finite. When merely $\al\in\Pp_1(\RR^d)$, the maximal-correlation term remains convex along such finite-second-moment geodesics by approximation; existence and uniqueness on the full $\Pp_1$ domain follow from the lower-semicontinuity argument of Santambrogio.
\end{prop}
\begin{proof}
The existence proof has two ingredients. First, since $\al$ is centered, translating $\eta$ does not change either $\mathcal H(\eta)$ or $\mathsf{MC}_\al(\eta)$, so one can center a minimizing sequence. Second, the assumption that $\al$ is not supported on a hyperplane gives a coercive lower bound of the form $\mathsf{MC}_\al(\eta)\geq c_\al\int\norm{x}\,\d\eta(x)$ for centered absolutely continuous $\eta$, with $c_\al>0$. Together with the lower-semicontinuity estimates for entropy and maximal correlation, this yields a minimizer.

Let $\eta=r\,\d x$ be a minimizer and let $u$ be a convex optimizer in the dual formula~\eqref{eq-moment-max-correlation-dual}. Keeping $u$ fixed and varying $\eta$ in~\eqref{eq-moment-measure-variational} gives the Euler equation
\(\log r(x)+1+u(x)=\mathrm{constant} \qquad\text{on }\{r>0\},\)
so $\eta=Z_u^{-1}e^{-u}\d x$. The optimality condition for the scalar-product transport problem says that an optimal coupling is supported on $\{(x,y):y\in\partial u(x)\}$. Since $\eta$ is absolutely continuous and $u$ is convex, $\partial u(x)=\{\nabla u(x)\}$ for $\eta$-almost every $x$, hence $\al=(\nabla u)_\sharp\eta$.

Conversely, add a constant to \(u\) so that \(\eta=e^{-u}\d x\), and first let \(\nu\in\Pp_2(\RR^d)\) be compactly supported. Write \(T=\nabla v\) for the Brenier map from \(\eta\) to \(\nu\), and \(\eta_t=((1-t)\Id+tT)_\sharp\eta\). Denote the lower right derivative of the objective along this geodesic by
\[
D^+\!\left(\mathcal H+\mathsf{MC}_\al\right)(\eta;\nu)
\eqdef
\liminf_{t\downarrow0}
\frac{(\mathcal H+\mathsf{MC}_\al)(\eta_t)-(\mathcal H+\mathsf{MC}_\al)(\eta)}{t}.
\]
The nonsmooth one-sided estimate of is
\[
D^+\!\left(\mathcal H+\mathsf{MC}_\al\right)(\eta;\nu)
\geq
-\int(\Delta^{\mathrm{ac}}v-d)\,\d\eta
+\int\dotp{\nabla v-x}{\nabla u}\,\d\eta,
\]
where \(\Delta^{\mathrm{ac}}v\) is the Lebesgue density of the absolutely continuous part of the distributional Laplacian of the convex potential \(v\). Essential continuity gives the boundary inequality
\(\int\dotp{x}{\nabla u}\,\d\eta\leq d\), and therefore
\[
D^+\!\left(\mathcal H+\mathsf{MC}_\al\right)(\eta;\nu)
\geq
-\int\Delta^{\mathrm{ac}}v\,\d\eta
+\int\dotp{\nabla v}{\nabla u}\,\d\eta.
\]
The distributional Hessian of the convex function \(v\) is a positive matrix-valued measure, so
\(\Delta v=\Delta^{\mathrm{ac}}v\,\d x+\Delta^s v\) with \(\Delta^s v\geq0\). Since \(\nabla u\,e^{-u}=-\nabla(e^{-u})\), cutoff integration by parts yields
\[
-\int\Delta^{\mathrm{ac}}v\,e^{-u}\d x
-\int\dotp{\nabla v}{\nabla(e^{-u})}\d x
\geq
-\int e^{-u}\,\d(\Delta v)
-\int\dotp{\nabla v}{\nabla(e^{-u})}\d x
=0.
\]
Here boundedness of \(\nabla v\), integrability of \(\nabla(e^{-u})\), and the decay of the nondegenerate log-concave density make the cutoff boundary term vanish. Thus every such one-sided derivative is nonnegative. Displacement convexity gives global minimality against compactly supported competitors, and the approximation theorem in extends it to the full \(\Pp_1\) domain. The equality case in the displacement-convexity argument yields uniqueness modulo translations; this is precisely the unavoidable invariance because translating \(\eta\) does not change \(\mathsf{MC}_\al\) when \(\al\) is centered.

The entropy term is displacement convex by McCann's theorem, recalled in Theorem~\ref{thm-mccann-internal-energy}. If $\al$ has finite second moment, identity~\eqref{eq-moment-correlation-w2} writes $\mathsf{MC}_\al$ as the sum of the $1$-convex moment term $\eta\mapsto\frac12\int\norm{x}^2\,\d\eta$ and the $(-1)$-convex term $\eta\mapsto-\frac12\Wass_2^2(\eta,\al)$, hence $\mathsf{MC}_\al$ is displacement convex. For a target with only a finite first moment, Santambrogio obtains the same convexity along $\Pp_2$ geodesics by approximation and proves the full variational characterization by lower semicontinuity.
\end{proof}
\paragraph{Conjugate moment measures for generation.}
Their existence theorem assumes that the target \(\beta\) is absolutely continuous and supported on a compact convex set; under these hypotheses they propose the conjugate factorization
\begin{equation}\label{eq-conjugate-moment-measure}
\beta
=
(\nabla w^*)_\sharp
\left(Z_w^{-1}e^{-w(z)}\,\d z\right).
\end{equation}
\subsection{Evolution in Depth of Transformers}
\label{sec-transformer-depth-evolution}
\paragraph{Attention as a context-dependent velocity.}
With residual scale \(1/T\), it acts as
\begin{equation}\label{eq-transformer-residual-attention-step}
x_i
\longmapsto
x_i+\frac1T\operatorname{Att}^{(n)}_\theta(x_1,\ldots,x_n)_i
=
x_i+\frac1T
\sum_{j=1}^n
\frac{e^{\dotp{Qx_i}{Kx_j}}Vx_j}
{\sum_{\ell=1}^n e^{\dotp{Qx_i}{Kx_\ell}}}.
\end{equation}
\paragraph{Token measure evolution.}
Let \(x_i^{(\ell)}\) be token \(i\) after layer \(\ell\), and let \(\al^{(\ell)}=n^{-1}\sum_i\de_{x_i^{(\ell)}}\) be the corresponding empirical law. With \(\tau=1/T\), equation~\eqref{eq-transformer-residual-attention-step} is exactly
\begin{equation}\label{eq-transformer-residual-measure-step}
\al^{(\ell+1)}
=
\left(
\Id+\tau\Gamma_{\theta^{(\ell)}}[\al^{(\ell)}]
\right)_\sharp
\al^{(\ell)},
\qquad
\ell=0,\ldots,T-1.
\end{equation}
As \(T\to+\infty\), a time interpolation of~\eqref{eq-transformer-residual-measure-step} converges formally to the conservation equation
\begin{equation}\label{eq-transformer-mean-field-pde}
\partial_t\al_t
+
\diverg\!\left(\al_t\Gamma_{\theta_t}[\al_t]\right)
=
0.
\end{equation}
\paragraph{\texorpdfstring{$L^2$ attention and mean shift.}{L^2 attention and mean shift.}}
Take, for simplicity, the same token space for queries, keys and values, and set
\[
K_\epsilon(x,y)=\exp(-\norm{x-y}^2/(2\epsilon)),
\qquad
\rho_\epsilon[\alpha](x)=\int K_\epsilon(x,y)\d\alpha(y),
\qquad
m_\epsilon[\alpha](x)
=
\frac{\int yK_\epsilon(x,y)\d\alpha(y)}
{\rho_\epsilon[\alpha](x)}.
\]
This gives
\begin{equation}\label{eq-l2-attention-mean-shift}
M_\epsilon[\alpha](x)
\eqdef
m_\epsilon[\alpha](x)-x
=
\frac{\int (y-x)K_\epsilon(x,y)\d\alpha(y)}
{\rho_\epsilon[\alpha](x)}
=
\epsilon\nabla\log\rho_\epsilon[\alpha](x)
\end{equation}
Its damped update is
\[
x_i^{(\ell+1)}
=
(1-\tau)x_i^{(\ell)}+\tau m_\epsilon[\alpha^{(\ell)}](x_i^{(\ell)})
=
x_i^{(\ell)}+\tau M_\epsilon[\alpha^{(\ell)}](x_i^{(\ell)})
\]
which is an explicit Euler step of the continuous-time mean-shift equation
\(\partial_t\alpha_t+\diverg\bigl(\alpha_tM_\epsilon[\alpha_t]\bigr)=0.\)
\paragraph{Gradient structure and limitations.}
When the token space has dimension $d$ and the query/key space has dimension $r$, take $Q,K\in\RR^{r\times d}$ and $V\in\RR^{d\times d}$. If $V=Q^\top K$, the field $\Gamma_\theta[\alpha]$ is a gradient vector field in the token variable. Indeed, define the log-partition potential
\(\Phi_\alpha(x) = \int \exp(\dotp{Qx}{Ky})\d\alpha(y), \qquad U_\alpha(x)=\log\Phi_\alpha(x).\)
Then
\(\nabla_x U_\alpha(x) = \frac{\int Q^\top K y\,\exp(\dotp{Qx}{Ky})\d\alpha(y)} {\int \exp(\dotp{Qx}{Kz})\d\alpha(z)} = \Gamma_\theta[\alpha](x).\)
Indeed, the natural scalar candidate $f_{\rm att}(\alpha)=\int U_\alpha(x)\,\d\alpha(x)$ has first variation
\(\delta f_{\rm att}(\alpha)(z) = U_\alpha(z) + \int \frac{\exp(\dotp{Qx}{Kz})}{\Phi_\alpha(x)} \,\d\alpha(x),\)
up to an additive constant. The second term is the response of every query normalization to a perturbation of the key distribution, and its spatial gradient is absent from $\Gamma_\theta[\alpha]$.

\subsection{Flows over the Gaussian Manifold}
\label{sec-gaussian-closure-transport-dynamics}
\paragraph{Gaussianity preservation.}
\begin{prop}[Affine velocities preserve Gaussianity]\label{prop-gaussian-affine-closure}
Let $\alpha_0=\Gaussian(m_0,\Sigma_0)$, with $\Sigma_0\succ0$. Let $b_t\in\RR^d$ and $A_t\in\RR^{d\times d}$ be locally integrable on a time interval, and let $(m_t,\Sigma_t)$ solve
\(\dot m_t=b_t, \qquad \dot\Sigma_t=A_t\Sigma_t+\Sigma_tA_t^\top, \qquad (m_{t=0},\Sigma_{t=0})=(m_0,\Sigma_0).\)
Then, as long as this matrix ODE is defined, $\Sigma_t\succ0$ and
\(\alpha_t=\Gaussian(m_t,\Sigma_t)\)
is the solution of the continuity equation
\(\partial_t\alpha_t+\diverg(\alpha_t v_t)=0, \qquad v_t(x)=b_t+A_t(x-m_t).\)
In particular, if a smooth functional $f$ has a Wasserstein gradient on Gaussian measures of the affine form
\(\Wgrad f(\Gaussian(m,\Sigma))(x) = b_f(m,\Sigma)+A_f(m,\Sigma)(x-m),\)
with $A_f(m,\Sigma)$ symmetric, then the displayed moment ODE below produces a Gaussian distributional solution of the full Wasserstein gradient-flow equation for $f$:
\(\dot m_t=h_f(m_t,\Sigma_t), \qquad \dot\Sigma_t=H_f(m_t,\Sigma_t),\)
where \(h_f(m,\Sigma)=-b_f(m,\Sigma), \qquad H_f(m,\Sigma) = -\bigl(A_f(m,\Sigma)\Sigma+\Sigma A_f(m,\Sigma)\bigr).\)
Conversely, fix a non-degenerate Gaussian \(\Gaussian(m,\Sigma)\). Suppose that a finite-energy Wasserstein tangent field \(V_{m,\Sigma}\) is represented by a gradient in \(L^2(\Gaussian(m,\Sigma))\) and is tangent to the Gaussian manifold at this Gaussian. Then there exist \(b(m,\Sigma)\) and a symmetric matrix \(A(m,\Sigma)\) such that
\(V_{m,\Sigma}(x)=b(m,\Sigma)+A(m,\Sigma)(x-m) \qquad \Gaussian(m,\Sigma)\text{-a.e.}\)
Without the Wasserstein gradient representative, this converse is false because one may add velocity fields with zero \(\Gaussian(m,\Sigma)\)-weighted divergence.
Finally, any smooth Gaussian curve with positive definite covariance can be generated by an affine velocity. If one wants the velocity to be a Wasserstein tangent gradient, one chooses the unique symmetric solution of the Lyapunov equation
\(A_t\Sigma_t+\Sigma_t A_t=\dot\Sigma_t.\)
\end{prop}
\begin{proof}
Let \(X_t\) follow the characteristic ODE \(\dot X_t=b_t+A_t(X_t-m_t)\) with \(X_0\sim\Gaussian(m_0,\Sigma_0)\). Since \(\dot m_t=b_t\), the centered variable \(\tilde X_t=X_t-m_t\) solves the homogeneous linear ODE \(\dot{\tilde X}_t=A_t\tilde X_t\). If \(\Phi_t\) is the fundamental matrix \(\dot\Phi_t=A_t\Phi_t\), \(\Phi_{t=0}=\Id\), then
\(X_t=m_t+\Phi_t(X_0-m_0), \qquad \Sigma_t=\Phi_t\Sigma_0\Phi_t^\top.\)
Hence \(X_t\) is Gaussian and \(\Sigma_t\succ0\), and
\(\dot\Sigma_t = \frac{\d}{\d t}\EE\bigl(\tilde X_t\tilde X_t^\top\bigr) = A_t\Sigma_t+\Sigma_t A_t^\top.\)
This proves Gaussian preservation and the moment ODE\@. The Wasserstein gradient-flow statement follows by inserting the descent velocity $v_t=-\Wgrad f(\alpha_t)$.

For the converse, fix $\alpha=\Gaussian(m,\Sigma)$ and denote its density by $\rho$. Tangency to the Gaussian manifold means that the density variation $-\diverg(\rho V)$ is generated by some moment variation $(\dot m,\dot\Sigma)$, with $\dot\Sigma$ symmetric. Set $b=\dot m$, and let $A=A^\top$ be the unique solution of
\(A\Sigma+\Sigma A=\dot\Sigma\).
By the first part of the proposition, the affine gradient field $V_0(x)=b+A(x-m)$ generates exactly the same infinitesimal Gaussian variation. Hence
\(\diverg\bigl(\rho(V-V_0)\bigr)=0\)
in the distributional sense. Both $V$ and $V_0$ belong to the $L^2(\alpha)$ closure of gradient fields. The weighted Helmholtz decomposition therefore makes $V-V_0$ simultaneously a tangent gradient and orthogonal to every tangent gradient, so $V=V_0$ in $L^2(\alpha)$. Equivalently, for a smooth representative $V-V_0=\nabla\psi$, integration by parts gives $\int\norm{\nabla\psi}^2\d\alpha=0$. This proves that the Wasserstein tangent representative is affine. The qualification is essential: without selecting the gradient representative, one can add a nonzero field with zero $\alpha$-weighted divergence without changing the Gaussian curve.

For a prescribed smooth Gaussian curve, set $b_t=\dot m_t$ and choose any matrix $A_t$ satisfying $A_t\Sigma_t+\Sigma_tA_t^\top=\dot\Sigma_t$. Since $\Sigma_t$ is positive definite, the Lyapunov map $A\mapsto A\Sigma_t+\Sigma_t A$ is invertible on symmetric matrices, which gives the unique symmetric choice when a gradient velocity is required. In that case $v_t$ is the gradient of the quadratic potential $x\mapsto \dotp{b_t}{x}+\dotp{A_t(x-m_t)}{x-m_t}/2$.
\end{proof}
\paragraph{Gaussian-preserving gradient flows.}
\begin{prop}[Gaussian closure catalogue]\label{prop-centered-gaussian-covariance-catalogue}
Let \(\gamma=\Gaussian(\bar m,\bar\Sigma)\) on \(\RR^d\), with \(\bar\Sigma\succ0\), and let the initial condition be \(\alpha_0=\Gaussian(m_0,\Sigma_0)\), with \(\Sigma_0\succ0\). For each functional displayed below, evaluating its ambient Wasserstein gradient on a Gaussian gives an affine field.
\(\alpha_t=\Gaussian(m_t,\Sigma_t)\)
is an exact distributional solution of the corresponding ambient gradient-flow PDE, and its mean and covariance satisfy
\(\dot{m}_t=h(m_t,\Sigma_t), \qquad \dot{\Sigma}_t=H(m_t,\Sigma_t).\)
Write \(\delta_m=m-\bar m\), \(A=\bar\Sigma^{-1}\), and
\(M_{\Sigma,\bar\Sigma} \eqdef \Sigma^{-1/2}\bigl(\Sigma^{1/2}\bar\Sigma\Sigma^{1/2}\bigr)^{1/2}\Sigma^{-1/2}.\)
With the normalizations displayed in the first column, the mean vector field \(h\) and covariance vector field \(H\) are listed in the following table. If that PDE is unique in the selected solution class, this Gaussian solution is the Wasserstein flow from \(\alpha_0\). Gradients with respect to \(\Sigma\) are symmetric Frobenius gradients on the cone of covariance matrices.
\begin{center}
\begingroup
\small
\renewcommand{\arraystretch}{1.55}
\textbf{Exact Gaussian solutions of ambient Wasserstein gradient equations.}\par\smallskip
\endgroup
\end{center}
Here, in the MMD row, \(R=\Sigma+m\,m^\top-\bar\Sigma-\bar m\,\bar m^\top.\)
The debiased Sinkhorn row uses the notation of Corollary~\ref{cor-gaussian-sinkhorn-divergence}: for Gaussian inputs,
\(\bar\MK_{\norm{\cdot-\cdot}^2}^{\epsilon}(\alpha,\gamma) = \norm{\delta_m}^2+\Bb_\epsilon(\Sigma,\bar\Sigma)^2,\)
with the closed-form covariance gradient
\[
G_\epsilon(\Sigma,\bar\Sigma)
=
\tau_\epsilon(\Sigma)
-
\bar\Sigma^{1/2}
\tau_\epsilon\bigl(B_{\Sigma,\bar\Sigma}^{1/2}\bigr)
B_{\Sigma,\bar\Sigma}^{-1/2}
\bar\Sigma^{1/2},
\qquad
B_{\Sigma,\bar\Sigma}=\bar\Sigma^{1/2}\Sigma\bar\Sigma^{1/2}.
\]
Here \(\tau_\epsilon\) is the scalar function
\(\tau_\epsilon(r) \eqdef \frac{\sqrt{\epsilon^2+16r^2}-\epsilon}{4r}, \qquad r>0,\)
applied to positive matrices by spectral calculus. Equivalently, for \(M\succ0\), \(\tau_\epsilon(M) = \bigl(\sqrt{\epsilon^2 \Id+16M^2}-\epsilon \Id\bigr)(4M)^{-1}.\)
With this convention, \(G_\epsilon=\nabla_\Sigma \Bb_\epsilon(\Sigma,\bar\Sigma)^2\).
In the sliced row, \(\sigma\) is the normalized spherical measure on \(\Sphere^{d-1}\), and
\[
G_{\mathrm{sw}}(\Sigma,\bar\Sigma)
=
\int_{\Sphere^{d-1}}
\left(
1-\sqrt{\frac{\theta^\top\bar\Sigma\theta}{\theta^\top\Sigma\theta}}
\right)
\theta\theta^\top\,\d\sigma(\theta).
\]
Here
\[
\mathcal I(\alpha|\gamma)
=
\int \left|\nabla\log\frac{\rho(x)}{\rho_\gamma(x)}\right|^2\rho(x)\,\d x
=
\int
\left|\nabla\log\rho(x)+A(x-\bar m)\right|^2\rho(x)\,\d x
\qquad(\alpha=\rho\,\d x),
\]
where $\rho_\gamma$ is the density of $\gamma$.
\end{prop}
\begin{proof}
For the moment-functional row, the formula is exact on the full Wasserstein space. Indeed, write $m_\alpha=\int x\,\d\alpha(x)$ and $\Sigma_\alpha=\int (x-m_\alpha)(x-m_\alpha)^\top\d\alpha(x)$. If \(\eta\) is a signed perturbation with zero total mass, then $\d m_\alpha[\eta]=\int x\,\d\eta(x)$ and $\d\Sigma_\alpha[\eta]=\int (x-m_\alpha)(x-m_\alpha)^\top\d\eta(x)$. Thus, up to an irrelevant additive constant, the first variation is $\delta f(\alpha)(x)=\dotp{\nabla_m g}{x}+\dotp{\nabla_\Sigma g}{(x-m_\alpha)(x-m_\alpha)^\top}$,
and its spatial gradient is the affine field
\(\nabla_x\delta f(\alpha)(x) = \nabla_m g+2\nabla_\Sigma g\,(x-m_\alpha).\)

For the quadratic potential row, the ambient first variation is \(\delta f(\alpha)(x)=\frac12x^\top Bx+\dotp{\ell}{x}, \qquad \nabla_x\delta f(\alpha)(x)=B m+\ell+B(x-m),\)
which gives the displayed affine velocity. For the quadratic interaction row, the ambient first variation is $\delta f(\alpha)(x)=\frac12\int (x-y)^\top G(x-y)\d\alpha(y)$, with $\nabla_x\delta f(\alpha)(x)=G(x-m_\alpha)$.
These first variations are quadratic in \(x\), hence the affine Gaussian flow coincides with the full Wasserstein gradient flow.

For the KL row, write $\alpha=\rho\,\d x$ and let $\rho_\gamma$ denote the density of $\gamma$. The ambient first variation is \(\log(\rho/\rho_\gamma)\) up to an additive constant. If \(\alpha=\Gaussian(m,\Sigma)\), then
\(\nabla\log(\rho/\rho_\gamma)(x) = -\Sigma^{-1}(x-m)+A(x-\bar m),\)
so the descent velocity is
\(v(x)=-A\delta_m+(\Sigma^{-1}-A)(x-m)\).
Proposition~\ref{prop-gaussian-affine-closure} gives \(h(m,\Sigma)=-A\delta_m\) and \(H(m,\Sigma)=2\Id-\Sigma A-A\Sigma\). For the Fisher row, writing \(u=\log(\rho/\rho_\gamma)\), the ambient first variation of \(\frac12\mathcal I(\alpha|\gamma)\) is, up to constants,
\(-\Delta u-\frac12|\nabla u|^2-\dotp{\nabla u}{\nabla\log\rho_\gamma}.\)
For Gaussian \(\rho\) and Gaussian \(\gamma\), this is quadratic in \(x\). Multiplying by two for \(\mathcal I(\alpha|\gamma)\) gives the descent velocity
\(v(x)=-2A^2\delta_m+2(\Sigma^{-2}-A^2)(x-m).\)
Hence the Wasserstein velocity is affine and the Gaussian family is closed, with the displayed \(h\) and \(H\). Although the ambient Fisher flow is a fourth-order PDE for a generic density, it is therefore an exact finite-dimensional closure in this quadratic-reference case.

For \(\Wass_2^2(\alpha,\gamma)\), the Brenier map from \(\Gaussian(m,\Sigma)\) to \(\gamma\) is
\(T(x)=\bar m+M_{\Sigma,\bar\Sigma}(x-m)\).
The descent velocity for the unhalved squared distance is \(2(T-\Id)\), which gives the mean and covariance equations through Proposition~\ref{prop-gaussian-affine-closure}. For the MMD row, \(k(x,y)=\dotp{x}{y}^2\) identifies the kernel mean embedding with the raw second moment. Thus
\(\MMD_k^2(\alpha,\gamma) = \norm{\Sigma+m\,m^\top-\bar\Sigma-\bar m\,\bar m^\top}_{\mathrm F}^2 = \norm{R}_{\mathrm F}^2,\)
and the ambient first variation is \(2x^\top R x\), whose descent velocity is \(-4Rx\).

For the Sinkhorn row, Corollary~\ref{cor-gaussian-sinkhorn-divergence} gives the closed Gaussian formula as a squared mean displacement plus the smoothed Bures term \(\Bb_\epsilon^2\). This is not only a formula on the Gaussian submanifold. The first variation of the entropic value with respect to the first marginal is a Sinkhorn dual potential, and for Gaussian marginals the cross and self dual potentials are quadratic. Their debiased combination is therefore quadratic, so its spatial gradient is affine. This proves existence of the displayed Gaussian ambient solution; identifying it with every weak solution requires uniqueness in the chosen class. Since the restriction of the functional to Gaussians is
\((m,\Sigma)\mapsto \norm{m-\bar m}^2+\Bb_\epsilon(\Sigma,\bar\Sigma)^2,\)
the quadratic coefficient is \(G_\epsilon=\nabla_\Sigma\Bb_\epsilon(\Sigma,\bar\Sigma)^2\), and the moment-functional computation gives the displayed \(h\) and \(H\). To compute this gradient, set \(B_{\Sigma,\bar\Sigma}=\bar\Sigma^{1/2}\Sigma\bar\Sigma^{1/2}\). Since \(\psi_\epsilon'(r)=-2\tau_\epsilon(r)\), differentiating
\(\tr\psi_\epsilon(B_{\Sigma,\bar\Sigma}^{1/2})\) gives
\(-\bar\Sigma^{1/2} \tau_\epsilon\bigl(B_{\Sigma,\bar\Sigma}^{1/2}\bigr) B_{\Sigma,\bar\Sigma}^{-1/2} \bar\Sigma^{1/2},\)
whereas differentiating the self term
\(-\frac12\tr\psi_\epsilon(\Sigma)\) gives \(\tau_\epsilon(\Sigma)\). This proves the displayed closed form for \(G_\epsilon\). When \(\bar\Sigma=\Id\), this covariance contribution is a spectral function of \(\Sigma\). If \(\lambda\) is an eigenvalue of \(\Sigma\), its covariance-eigenvalue derivative is
\(4\sqrt{\lambda+\epsilon^2/16}-4\sqrt{\lambda^2+\epsilon^2/16},\)
which is \(\dot\lambda\) in the displayed ODE. The centered spatial affine-velocity coefficient in that eigendirection is \(\dot\lambda/(2\lambda)\).

For sliced Wasserstein, the exact ambient Wasserstein gradient is obtained by averaging the one-dimensional gradients. Each projection is Gaussian:
\[
(P_\theta)_\sharp\alpha=\Gaussian(\dotp{\theta}{m},\theta^\top\Sigma\theta),
\qquad
(P_\theta)_\sharp\gamma=\Gaussian(\dotp{\theta}{\bar m},\theta^\top\bar\Sigma\theta).
\]
Hence
\[
\SW_2^2(\alpha,\gamma)
=
\int_{\Sphere^{d-1}}
\left[
\dotp{\theta}{\delta_m}^2
+
\left(\sqrt{\theta^\top\Sigma\theta}
-\sqrt{\theta^\top\bar\Sigma\theta}\right)^2
\right]\d\sigma(\theta).
\]
If \(T_\theta\) is the monotone transport from \((P_\theta)_\sharp\alpha\) to \((P_\theta)_\sharp\gamma\), then the descent velocity for the unhalved sliced objective is
\(2\int_{\Sphere^{d-1}}\bigl(T_\theta(P_\theta x)-P_\theta x\bigr)\theta\,\d\sigma(\theta).\)
For Gaussian marginals, \(T_\theta\) is affine. Thus this velocity is affine in \(x\), and Proposition~\ref{prop-gaussian-affine-closure} applies. The spherical identity \(\int\theta\theta^\top\d\sigma(\theta)=\Id/d\) gives the mean equation, and differentiating the covariance term gives \(G_{\mathrm{sw}}\).
\end{proof}
\paragraph{One-dimensional case.}
Fix a Gaussian target $\gamma=\Gaussian(\bar m,\bar\sigma^2)$, with $\bar\sigma>0$, and write $\alpha=\Gaussian(m,\sigma^2)$, with $\sigma>0$. By Proposition~\ref{prop-gaussian-w2-bures}, the parametrization $(m,\sigma)\mapsto\Gaussian(m,\sigma^2)$ identifies the one-dimensional Gaussian family with the open Euclidean half-plane. Hence every Gaussian-preserving functional in Proposition~\ref{prop-centered-gaussian-covariance-catalogue} is represented by the scalar restriction
\(g_f(m,\sigma)\eqdef f\bigl(\Gaussian(m,\sigma^2)\bigr).\)
Set $\delta=m-\bar m$. For the Sinkhorn row below, with $\epsilon>0$, abbreviate
\[
\psi_\epsilon(r)
\eqdef
-\frac{\sqrt{\epsilon^2+16r^2}-\epsilon}{2}
-\frac{\epsilon}{2}
\log\!\left(
\frac{2\epsilon}{\sqrt{\epsilon^2+16r^2}+\epsilon}
\right),
\qquad r>0.
\]
\paragraph{Constrained evolution on the Gaussian manifold.}
Let
\(\mathcal G=\{\Gaussian(m,\Sigma):m\in\RR^d,\ \Sigma\succ0\}\)
be the Gaussian submanifold of $\Pp_2(\RR^d)$. The Wasserstein gradient of a functional constrained to a smooth submanifold $\mathcal M\subset\Pp_2$ is defined as the Riesz representative of the differential restricted to tangent velocities of $\mathcal M$. Equivalently, it is the small-step limit of the constrained JKO scheme
\(\alpha^{(\ell+1)}\in \uargmin{\alpha\in\mathcal M} \frac{1}{2\tau}\Wass_2^2(\alpha,\alpha^{(\ell)})+f(\alpha).\)
For $\mathcal M=\mathcal G$, tangent velocities are affine gradient fields $v(x)=b+A(x-m)$ with $A=A^\top$. The constrained gradient is therefore the $L^2(\Gaussian(m,\Sigma))$ projection of the ambient Wasserstein gradient onto this finite-dimensional affine space, whenever the ambient gradient exists.

\begin{prop}[Gaussian-constrained Wasserstein gradients]\label{prop-gaussian-gradient-bullet-list}
Let $f$ be a smooth functional and assume that its restriction to nondegenerate Gaussian measures can be written as \(f(\Gaussian(m,\Sigma))=F(m,\Sigma).\)
Then the Wasserstein gradient constrained to the Gaussian family is the affine vector field
\(v_F(x) = \nabla_m F(m,\Sigma) + 2\nabla_\Sigma F(m,\Sigma)(x-m),\)
where $\nabla_\Sigma F$ denotes the symmetric matrix derivative. Equivalently, $v_F$ is the $L^2(\Gaussian(m,\Sigma))$ projection of the ambient Wasserstein gradient onto affine gradient fields, whenever the ambient gradient exists. Hence the gradient descent flow constrained to Gaussian measures satisfies
\begin{equation}\label{eq-gaussian-wgf-closure}
\dot m_t=-\nabla_m F(m_t,\Sigma_t),
\qquad
\dot\Sigma_t=-2\bigl(\Sigma_t\nabla_\Sigma F(m_t,\Sigma_t)+\nabla_\Sigma F(m_t,\Sigma_t)\Sigma_t\bigr),
\end{equation}
and the descent velocity is affine.
\end{prop}
\begin{proof}
Test the functional along a Gaussian tangent vector, represented by an affine gradient field
\(v(x)=b+A(x-m)\)
with $A$ symmetric. The induced first-order variations are $\dot m=b$ and $\dot\Sigma=A\Sigma+\Sigma A$. Therefore
\[
\d F(m,\Sigma)[b,A\Sigma+\Sigma A]
=
\dotp{\nabla_m F}{b}
+
\tr\!\left(\nabla_\Sigma F(A\Sigma+\Sigma A)\right).
\]
Since $A$, $\Sigma$ and $\nabla_\Sigma F$ are symmetric, the second term equals
\(2\,\tr(\nabla_\Sigma F\,A\Sigma) = \int \dotp{2\nabla_\Sigma F(x-m)}{A(x-m)}\d\Gaussian(m,\Sigma)(x).\)
Together with the mean term, this gives \(\d F(m,\Sigma)[\dot m,\dot\Sigma] = \int \dotp{v_F(x)}{v(x)}\d\Gaussian(m,\Sigma)(x)\)
for all affine gradient fields $v$. This identifies the constrained Wasserstein gradient in the induced $L^2(\alpha)$ metric, or equivalently the projection of the ambient gradient when it exists. Substituting the descent velocity $-v_F$ in Proposition~\ref{prop-gaussian-affine-closure} gives~\eqref{eq-gaussian-wgf-closure}.
\end{proof}
\paragraph{Contractive Gaussian projection.}
\begin{defn}[Gaussian projection]\label{def-gaussian-projection}
For \(\alpha\in\Pp_2(\RR^d)\), its Gaussian projection is \(\GaussProj(\alpha) \eqdef \Gaussian(m_\alpha,\Sigma_\alpha),\)
the Gaussian with the same mean and covariance as \(\alpha\).
\end{defn}
\begin{thm}[Gelbrich theorem]\label{thm-gelbrich-projection}
For every $\al,\be\in\Pp_2(\RR^d)$, \(\Wass_2^2\bigl(\GaussProj(\al),\GaussProj(\be)\bigr) = \norm{m_\al-m_\be}^2+\Bb^2(\Sigma_\al,\Sigma_\be) \leq \Wass_2^2(\al,\be).\)
\end{thm}
\begin{proof}
Take any coupling $(X,Y)$ of $\al$ and $\be$, center the variables, and write $C=\EE[(X-m_\al)(Y-m_\be)^\top]$. In the positive definite case, positivity of the block covariance matrix implies the factorization $C=\Sigma_\al^{1/2}K\Sigma_\be^{1/2}$ with $\norm{K}_{\mathrm{op}}\leq1$, and therefore, by operator/nuclear norm duality, $\tr C\leq \tr\bigl((\Sigma_\al^{1/2}\Sigma_\be\Sigma_\al^{1/2})^{1/2}\bigr)$.
The semidefinite case follows by adding $\eta\Id$ to both covariance matrices and letting $\eta\downarrow0$.
Expanding $\EE\norm{X-Y}^2$ gives the lower bound
\(\EE\norm{X-Y}^2 \geq \norm{m_\al-m_\be}^2+\Bb^2(\Sigma_\al,\Sigma_\be).\)
Taking the infimum over couplings proves the inequality, while equality for Gaussian laws is Proposition~\ref{prop-gaussian-w2-bures}.
\end{proof}
\begin{thm}[Hugo Lavenant Gaussian-preservation criterion]\label{thm-lavenant-gaussian-preserving-jko}
Let $f:\Pp_2(\RR^d)\to(-\infty,+\infty]$ satisfy
\(f\bigl(\GaussProj(\al)\bigr)\leq f(\al) \qquad\forall\al\in\Pp_2(\RR^d),\)
with $\GaussProj$ defined in Definition~\ref{def-gaussian-projection}. If $\gamma$ is Gaussian and $\alpha^\star$ minimizes the JKO step
\(\eta\mapsto f(\eta)+\frac1{2\tau}\Wass_2^2(\gamma,\eta),\)
then $\GaussProj(\alpha^\star)$ is also a minimizer. If this JKO minimizer is unique, it is Gaussian.
\end{thm}
\begin{proof}
For the JKO claim, $\GaussProj(\gamma)=\gamma$ because $\gamma$ is Gaussian. Hence, for any competitor $\eta$,
\(f\bigl(\GaussProj(\eta)\bigr)+\frac1{2\tau}\Wass_2^2\bigl(\gamma,\GaussProj(\eta)\bigr) \leq f(\eta)+\frac1{2\tau}\Wass_2^2(\gamma,\eta).\)
Applying this to a minimizer $\eta=\alpha^\star$ shows that $\GaussProj(\alpha^\star)$ is again a minimizer. Uniqueness forces $\alpha^\star=\GaussProj(\alpha^\star)$.
\end{proof}
\paragraph{Moment closure beyond Gaussianity.}
If $f(\alpha)=g(m_\alpha,\Sigma_\alpha)$ and $G=\nabla_\Sigma g(m_\alpha,\Sigma_\alpha)$ denotes the symmetric Frobenius gradient, then the Wasserstein descent velocity is
\(v_\alpha(x)=-\nabla_m g(m_\alpha,\Sigma_\alpha)-2G(x-m_\alpha)\).
Consequently, every sufficiently regular flow starting from an arbitrary $\alpha_0\in\Pp_2(\RR^d)$ satisfies
\begin{equation}\label{eq-distribution-free-mean-covariance-closure}
\dot m_t=-\nabla_m g(m_t,\Sigma_t),
\qquad
\dot\Sigma_t
=
-2\bigl(\Sigma_t\nabla_\Sigma g(m_t,\Sigma_t)
+\nabla_\Sigma g(m_t,\Sigma_t)\Sigma_t\bigr).
\end{equation}
For Gaussian initial data, this affine velocity also preserves Gaussianity; for general initial data, the law need not become Gaussian, but its mean and covariance still follow the same closed vector field $(h,H)$ listed in that proposition.

\begin{prop}[Scalar moment closure and eikonal characterization]\label{prop-scalar-moment-closure}
Let $\Omega\subset\RR^d$ be connected and open, let $\varphi\in C^2(\Omega)$, and set
\(m_\varphi(\alpha)\eqdef\int_\Omega\varphi(x)\d\alpha(x)\)
for probability measures such that $\varphi$ and $\norm{\nabla\varphi}^2$ are integrable. For every smooth $g:\RR\to\RR$, the Wasserstein gradient flow of $f_g(\alpha)=g(m_\varphi(\alpha))$ satisfies, whenever the differentiation is justified,
\(\dot m_t = -g'(m_t)\int_\Omega\norm{\nabla\varphi(x)}^2\d\alpha_t(x), \qquad m_t=m_\varphi(\alpha_t).\)
This equation is autonomous for every $g$, meaning that its right-hand side depends on $\alpha_t$ only through $m_t$, if and only if there are constants $a,b\in\RR$ such that
\begin{equation}\label{eq-scalar-moment-affine-carre}
\norm{\nabla\varphi(x)}^2=a+b\varphi(x)
\qquad (x\in\Omega).
\end{equation}
In that case, $\dot m_t=-(a+b m_t)g'(m_t)$.
Moreover, on every connected open set $U\subset\{x:\nabla\varphi(x)\neq0\}$, the restriction of~\eqref{eq-scalar-moment-affine-carre} to $U$ is equivalent to the existence of an eikonal coordinate $r\in C^2(U)$ and constants $c,\kappa,\lambda\in\RR$ such that
\begin{equation}\label{eq-scalar-moment-eikonal-form}
\norm{\nabla r}=1,
\qquad
\varphi=c+\kappa r+\frac{\lambda}{2}r^2.
\end{equation}
Under~\eqref{eq-scalar-moment-affine-carre}, one may take $\lambda=b/2$.
\end{prop}
\begin{proof}
The first variation is $\delta f_g(\alpha)(x)=g'(m_\varphi(\alpha))\varphi(x)$, so the continuity equation gives the displayed evolution of $m_t$. Condition~\eqref{eq-scalar-moment-affine-carre} then proves sufficiency.

For necessity, it is enough to take $g(s)=s$ and write $\psi=\norm{\nabla\varphi}^2$. Autonomy for Dirac laws first shows that $\psi=q\circ\varphi$ for some function $q$ on $\varphi(\Omega)$. Because $\Omega$ is connected, $\varphi(\Omega)$ is an interval. Comparing a Dirac law with any two-point mixture having the same $\varphi$-moment shows that
\(q\bigl((1-\zeta)s_0+\zeta s_1\bigr) = (1-\zeta)q(s_0)+\zeta q(s_1) \qquad (0\leq\zeta\leq1).\)
Hence $q$ is affine, which is precisely~\eqref{eq-scalar-moment-affine-carre}.

On a regular set $U$, the right-hand side of~\eqref{eq-scalar-moment-affine-carre} is positive. Fix $s_0\in\varphi(U)$ and define $r(x)=\int_{s_0}^{\varphi(x)}(a+bs)^{-1/2}\d s$.
Then $\norm{\nabla r}=1$. Inverting this change of variables gives~\eqref{eq-scalar-moment-eikonal-form} with $c=s_0$, $\kappa=\sqrt{a+bs_0}$, and $\lambda=b/2$. Conversely,~\eqref{eq-scalar-moment-eikonal-form} gives
$\norm{\nabla\varphi}^2=\kappa^2-2\lambda c+2\lambda\varphi$.
\end{proof}


\end{document}
